\documentclass[11pt]{article}
\usepackage[T1]{fontenc}
\usepackage[utf8]{inputenc}
\usepackage{lmodern,amsmath,amssymb,amsthm,mathrsfs,mathtools,booktabs,longtable,array}
\usepackage[margin=1in]{geometry}
\usepackage[colorlinks=true,urlcolor=blue,linkcolor=blue]{hyperref}
\newtheorem{theorem}{Theorem}[section]
\newtheorem{proposition}[theorem]{Proposition}
\newtheorem{lemma}[theorem]{Lemma}
\newtheorem{corollary}[theorem]{Corollary}
\theoremstyle{remark}\newtheorem{remark}[theorem]{Remark}
\DeclareMathOperator{\CH}{CH}
\DeclareMathOperator{\Alg}{Alg}
\DeclareMathOperator{\Pic}{Pic}
\DeclareMathOperator{\Alb}{Alb}
\DeclareMathOperator{\NS}{NS}
\DeclareMathOperator{\End}{End}
\DeclareMathOperator{\Hom}{Hom}
\DeclareMathOperator{\Nm}{Nm}
\newcommand{\Q}{\mathbf Q}\newcommand{\Z}{\mathbf Z}\newcommand{\F}{\mathbf F}
\newcommand{\tun}[1]{#1}
\setlength{\emergencystretch}{2em}
\title{C44\_POS: primitive signs, special axes, and the 1964 programme}
\author{Research continuation of JAC1964 and SIGN\_NEXT}
\date{22 November 2026}
\begin{document}
\maketitle
\begin{abstract}
Final research record at the scheduled stopping point. The general fixed-axis sign and the Fermat
Abel--Jacobi injection remain unresolved. New checked results include
an exact axis primitive decomposition, lifts of every embedded smooth
axis (without a cycle-lifting assertion), and an unconditional sign
theorem whenever at most one potential Tate direction remains after
a known positive algebraic span. An extending orthogonal-block criterion
now also proves the sign on actual closed affine-$S_3$ families with
three residual potential Tate directions, including families with a
known positive span of rank six. A totally real coefficient-field
criterion further gives actual dihedral families with five residual
potential directions and a known positive span of rank ten, including
characteristic three. A quaternion-packet criterion now controls residual imaginary-quadratic
planes, with actual cyclic families in every characteristic. Its finite-place
comparison does not assume algebraicity of the remaining Tate directions.
The universal gate outside these verified correspondence criteria remains open. Other
checked findings include closed supersingular axes with their entire
positive kernels controlled, intrinsically nonliftable surface centres
whose selected primitive classes are now excluded from every abstract smooth
projective lift of ramification index $e\leq5$ at $p=3$ and $e\leq3p$ at the
constructed primes $p\geq7$, $p\equiv3\pmod4$. The corresponding conditional
centre theorem at $p=5$ excludes $e\leq10$ when its nonliftable centre is supplied.
For larger ramification,
class lifting is reduced to an explicit finite initial support-lifting
obstruction. Other results include lifts of every finite
curve cover as a map, special-cover sign tests, exact integral
arithmetic obstructions, a characteristic-$p$ detector/division theorem,
and geometric extensions of the primitive incidence polarization.
The algebraically trivial Fermat obstruction is reduced to specified
nonzero universal cycles from excess Picard quotients; several entire
classes of singular ample image surfaces are eliminated in rational
Chow. The projected geometric 5-primary torsion vanishes, and the known
nonzero class has bounded geometric 5-divisibility; the exact residual
completed Chow term and the arithmetic finite-generation gate remain.
Full Abel--Jacobi injectivity remains open.
The proposed supersingular Fermat test case already satisfies the
required Hodge standard conjecture; its corrected restriction and
weighted Gysin conditions are verified explicitly. A separate
nonsupersingular degree-seven Fermat calculation supplies explicit
new special cycles, all-prime Tate and positivity for the fourfold,
and numerical and 13-adic positivity for its square.
Selected closed nonabelian axes also satisfy the Hodge standard
conjecture on all their powers for positive product polarizations,
strong Tate in every degree, and the corrected scan-050 conditions;
on the fourfold itself the sign holds for every ample polarization.
The historical prior-unknown status is not asserted. An explicit
full-monodromy family excludes uniform abelian realization of its whole
variable local system, without excluding isolated special-fibre
correspondences.
Every use of 032 or Milne's 2026 report remains conditional on those
unrefereed preprints. The separate Milne 2022 transfer assertion is
exposed wherever used. Whole-proof reading events and finite
computations are distinguished in \texttt{CHECK\_EVENTS.tsv}.
\end{abstract}
\tableofcontents
\part{Checked findings and explicitly conditional deductions}
\section{Normalization, scope, and inherited results}
The starting isometry is JAC1964 H9-T7, equation (H9.18):
\[
 A_X\oplus A_B\simeq\CH^m(Y)_{\mathrm{hom},\Q}/\mathcal H_0,
 \qquad Q((z,w),(z',w'))=(-1)^m\int_Xzz'+(-1)^{m-1}\int_Bww'.
\]
Here $A_B=\ker\kappa_*$ on the algebraic middle classes of the axis,
and $h_BA_B=0$. This is an inherited statement, not a new result of
this run. For an arbitrary abelian ambient, the ambient sign is conditional
on the unrefereed 032 Corollary 8.4 and Milne2026 report.
For the explicit elliptic-power families, A42b supplies it unconditionally. Weak semidefiniteness and strict positive definiteness
are distinguished throughout; no numerical/homological equivalence
on an arbitrary axis is silently assumed.

Geometric generic parameter points are taken over algebraic closures
of function fields. They are not closed points over
$\overline{\mathbf F}_p$. A monodromy vanishing assertion at the
former does not settle a sign at every latter point.

Every finding retains its stable C44 identifier. A displayed
proposition that assumes an additional unproved input is a checked
conditional implication, not a proof of that input. Prior statements
retain their JAC1964 or SIGN\_NEXT identifiers and are not claimed
as findings of this continuation. The governing source is the attached
\texttt{PROMPT(20261121-231316).md}.

% BEGIN work/axis/axis_contribution.tex
% C44_POS /root/axis_independent; GPT-6 Astra Pro; effort setting not exposed.
% Mathematical deductions for root review; no claim of formal verification.
\section{The exact axis primitive space and the limits of the proposed transfers}

\paragraph{Notation and carried input.}
Put $r=m-1$, so that $\dim X=2r+2$ and $\dim B=2r$.
Write $L_X$ and $L_B$ for multiplication by $h$ and $h_B$.
All cycle groups have rational coefficients, and $\operatorname{Alg}^j$
means the image of the rational cycle-class map for one fixed
$\ell\ne p$. No equality between numerical and homological equivalence
on $B$ is assumed. The carried statement JAC1964 H9-T7, equation
(H9.18), is the isometry
\[
 A_X\oplus A_B\ \simeq\ \mathcal H/\mathcal H_0,
 \qquad
 Q((z,w),(z',w'))=(-1)^m\int_Xzz'+(-1)^r\int_Bww'.
\]
Its proof establishes $L_Bw=0$ for $w\in A_B$, not merely
$L_B^2w=0$. The factorization $B\hookrightarrow Y\hookrightarrow X$
through a smooth pencil member and the injectivity of the latter
Gysin map are essential to that stronger conclusion. We use this
checked input without reproving its height isometry.

\subsection*{C44-A01. The axis gate is the entire remaining primitive sign on the axis}

\textbf{Proposition.}
There is an orthogonal direct sum of rational algebraic classes
\[
 \operatorname{Alg}^{r}(B)\cap P^{2r}(B)
 =\kappa^*\bigl(\operatorname{Alg}^{r}(X)\cap P^{2r}(X)\bigr)
       \mathbin{\perp} A_B,                                  \tag{C44.A01.1}
\]
where $P^{2r}(B)=\ker L_B$ and $P^{2r}(X)=\ker L_X^3$.
On the first summand the signed middle form is
\[
   (-1)^r\int_B\kappa^*a\,\kappa^*a'
       =(-1)^r\int_Xh^2aa'.                                 \tag{C44.A01.2}
\]
For every $0\le j<r$, restriction identifies the primitive algebraic
classes of codimension $j$ on $X$ and $B$, including their signed
Lefschetz pairings.
Consequently, \emph{conditional on the unrefereed 032/Milne input
asserting the Hodge standard conjecture for $X$}, strict positivity
on $A_B$ is equivalent to the full Hodge standard conjecture $I(B)$
for $h_B$. Weak nonnegativity on $A_B$ is equivalent to the version
of $I(B)$ with every positive-definite assertion replaced by
positive-semidefinite. These last two statements must not be
identified unless nondegeneracy on algebraic classes is additionally
known.

\textbf{Proof.}
Hard Lefschetz makes $L_X^2:H^{2r}(X)\longrightarrow H^{2r+4}(X)$
an isomorphism. Its inverse $\beta$ is induced by an algebraic
correspondence because the Lefschetz standard conjecture is known
for abelian varieties, without 032. The self-intersection and
projection formulas give
$\kappa_*\kappa^*=L_X^2$. Hence
$e=\kappa^*\beta\kappa_*$ is a projector and every algebraic middle
class $x$ of $B$ has a unique decomposition
\[
 x=\kappa^*a+w,\qquad a=\beta\kappa_*x,\qquad w\in A_B.       \tag{C44.A01.3}
\]
Both components are rational algebraic classes. The summands are
orthogonal because
$\int_B\kappa^*a\,w=\int_Xa\,\kappa_*w=0$.
H9-T7 gives $L_Bw=0$. The Gysin map
\[
 \kappa_*:H^{2r+2}(B)\longrightarrow H^{2r+6}(X)
\]
is an isomorphism, being dual to weak Lefschetz in degree $2r-2$.
Therefore $L_B\kappa^*a=0$ if and only if
$\kappa_*L_B\kappa^*a=L_X^3a=0$. This proves (C44.A01.1).
The projection formula proves (C44.A01.2).

For $j<r$, weak Lefschetz gives an isomorphism
$\kappa^*:H^{2j}(X)\longrightarrow H^{2j}(B)$.
Its inverse on this degree is the algebraic operator
\[
 (L_X^{2r+2-2j})^{-1}\kappa_*L_B^{2r-2j};                   \tag{C44.A01.4}
\]
indeed its composition with $\kappa^*$ is the identity.
Thus this is an isomorphism on rational algebraic classes too.
For $1\le j<r$, the map $\kappa_*$ in the degree of
$L_B^{2r-2j+1}\kappa^*a$ is an isomorphism by dual weak
Lefschetz in degree $2j-2$. Thus the primitive condition on $B$
is equivalent to
\[
 L_X^{2r-2j+3}a=0,
\]
which is precisely the primitive condition on $X$. For $j=0$
both primitive spaces are $H^0$, so the same assertion holds.
Finally
\[
 (-1)^j\int_B h_B^{2r-2j}\kappa^*a\,\kappa^*a'
 =(-1)^j\int_Xh^{2r+2-2j}aa'.
\]
The conditional positivity assertions now follow from the
orthogonal direct sum and the signed Hodge standard conjecture on
$X$. No positivity assertion on the new summand has been proved.
\hfill$\square$

\paragraph{Programme location and scope.}
This is a sharpened identification of JAC1964 QH9-axis, and concerns
cote 44, Th.~5, scan 004, with the primitive margin on scan 008.
In cote 16 terminology it isolates $I(B)$; the algebraic operator
$\beta$ uses $C_G(X)=B(X)$, not $I(B)$. It does not close
QH9-axis and does not upgrade its requested semidefinite sign to
strict positivity.

\subsection*{C44-A02. Every graph transfer to an abelian target annihilates the axis classes}

\textbf{Proposition.}
Let $D$ be an abelian variety and $f:B\longrightarrow D$ any
morphism. Then $f_*w=0$ in cohomology for every $w\in A_B$.
The same holds for a rational map, interpreted through its unique
extension. Consequently $A_B$ is orthogonal to every cohomology
class pulled back by such a morphism in the complementary degree.
This is unconditional. It does not rule out correspondences on
$B\times D$ which are not assembled from these graph maps.

\textbf{Proof.}
Choose a geometric base point $b$ on $B$. Let
$a_B:B\longrightarrow\operatorname{Alb}(B)$ be the pointed Albanese
map and let $u:\operatorname{Alb}(B)\longrightarrow X$ be induced by
$x\mapsto\kappa(x)-\kappa(b)$.
Weak Lefschetz in degree one identifies $H^1(X)$ with $H^1(B)$.
The identification of degree-one cohomology with the dual of the
Albanese Tate module therefore makes $u$ an isogeny.
After subtracting the constant $f(b)$, the morphism $f$ factors as
$v a_B$. Since $u$ is invertible in the isogeny category, there are
an integer $N>0$ and a homomorphism $t:X\longrightarrow D$ such that
$t u=[N]_D v$. Thus
\[
 [N]_D\circ f=t\circ\kappa+\text{a constant}.
\]
Translations on an abelian variety act trivially on cohomology.
Functoriality of proper pushforward gives
\[
 [N]_{D,*}f_*w=t_*\kappa_*w=0.
\]
Multiplication by $N$ induces an invertible map on rational
$\ell$-adic cohomology, so $f_*w=0$. The orthogonality assertion is
projection formula. A rational map from a smooth variety to an
abelian variety extends by Weil's extension theorem; one exact
modern locator is Milne, \emph{Abelian Varieties}, course notes
version 2.00, I, Theorem 3.2, printed p.~17 (PDF p.~23).\hfill$\square$

\paragraph{Programme location and scope.}
This specifies a limit of the proposed curve-test attack on
QH9-axis. It concerns the incidence higher Picard framework of
cote 44, Th.~1--4, scan 002, and the incidence divisor on scans
008--009. It complements JAC-P19 (no abelian rational domination)
and the already proved annihilation calculation in H9-T13, rather
than replacing either. An algebraic correspondence shifts
cohomological degree by an even number, so $H^1(C)$ by itself
cannot map into $H^{2r}(B)$ through an algebraic correspondence.
The elliptic tensor construction $w\boxtimes P_E$ resolves that
parity issue but its ambient pushforward is zero, exactly as
recorded in H9-T13. The missing input is a genuinely new
pairing-compatible correspondence, not a graph of another
Albanese map.

\subsection*{C44-A03. The finite-field test is absence of potential Tate roots, not one exceptional slope}

\textbf{Proposition.}
Suppose the axis and the embedding are defined over $\mathbf F_q$.
Let
\[
 V_B=\ker\bigl(\kappa_*:H^{2r}(B_{\overline{\mathbf F}_q},
 \mathbf Q_\ell)\longrightarrow H^{2r+4}(X_{\overline{\mathbf F}_q},
 \mathbf Q_\ell)(2)\bigr)
\]
be the untwisted primitive variable cohomology, with the twists in
the target understood by Gysin compatibility. Write $F$ for
geometric Frobenius and $P_B(t)=\det(t-F\mid V_B)$.
Every algebraic axis class is contained in
\[
 T_B=\bigcup_{a\ge1}\ker(F^a-q^{ra}\mid V_B).                \tag{C44.A03.1}
\]
In particular all its normalized Newton slopes are $r$.
If $P_B$ has no root $q^r\zeta$ for a root of unity $\zeta$,
then $A_B=0$.
For $D=\deg P_B>0$, a finite exact sufficient test is
\[
 \gcd\bigl(P_B(t),\ t^N-q^{rN}\bigr)=1,
 \qquad N=\operatorname{lcm}(1,2,\ldots,2D^2).               \tag{C44.A03.2}
\]
Here $P_B\in\mathbf Q[t]$ is obtained by dividing the middle
Frobenius polynomial of $B$ by that of the restricted ambient
summand. No Tate conjecture and no Frobenius semisimplicity
conjecture is needed for this vanishing criterion.

\textbf{Proof.}
A rational algebraic cycle on $B_{\overline{\mathbf F}_q}$ is
defined over a finite extension $\mathbf F_{q^a}$. Its codimension
$r$ cycle class in untwisted cohomology satisfies $F^a w=q^{ra}w$.
This proves (C44.A03.1). In crystalline realization, the same
cycle is a Frobenius eigenclass of slope $r$; equivalently every
root $q^r\zeta$ has normalized valuation $r$ at every prime over
$p$.
If the polynomial has no such root, every kernel in (C44.A03.1)
is zero, which proves the asserted vanishing.

The splitting (C44.A01.3) is Frobenius-compatible after the
polarization is defined over $\mathbf F_q$; its image is isomorphic
to $H^{2r}(X)$. Thus the usual rational Frobenius polynomials of
$X$ and $B$ give $P_B\in\mathbf Q[t]$ without asserting
rationality of a chosen algebraic basis of $A_B$.
If $q^r\zeta$ is a root and $n$ is the order of $\zeta$, the
cyclotomic polynomial of $\zeta$ divides
$q^{-rD}P_B(q^rt)$, so $\varphi(n)\le D$.
For $n=\prod p^a$,
\[
 \frac{\varphi(n)^2}{n}=\prod_{p^a\parallel n}p^{a-2}(p-1)^2
 \ge\frac12:
\]
each factor for $p\ge3$ is at least one, and the sole possible
factor less than one is $1/2$ for $2^1$. Consequently $n\le2D^2$,
so $n\mid N$. This proves (C44.A03.2), including its equivalence
to absence of roots of this form. For $D=0$ the conclusion is
immediate.\hfill$\square$

\paragraph{Why H9-T5 does not supply this test.}
JAC1964 H9-T5 is a theorem about the \emph{odd} variable local
system $R^{2m-1}$ of the pencil and its variable incidence abelian
factor. Its irreducibility hypothesis turns a nonzero injection
$H^1(D)(-(m-1))$ into an isomorphism; one slope outside
$[m-1,m]$ then contradicts that isomorphism. The even-degree
fixed-axis space in (C44.A03.1) is a different representation.
At a closed point the arithmetic Galois group is procyclic, and
an isolated algebraic class need not extend to a flat algebraic
class on the generic axis family. No family-monodromy
irreducibility statement can therefore be substituted into
(C44.A03.1) without a separate spreading theorem.
An out-of-range slope in some other direct summand is compatible
with a slope-$r$ Tate summand. The correct stronger sufficient
Newton condition here is absence of \emph{every} slope-$r$
subspace, while the cyclotomic criterion above is sharper.
This repairs the proposed use of H9-T5 without producing a
witness for every special axis.

\paragraph{Supersingular ambient varieties do not resolve this distinction.}
If $X$ is supersingular, $H^{2r}(X)$ has only slope $r$ because
$H^1(X)$ has slope $1/2$ and the cohomology algebra is exterior.
That identifies the slopes of the ambient summand in
(C44.A01.3), not the slopes or cycle classes of $V_B$.
A nonliftable algebraic class pulled back from $X$ lies in the
ambient summand and is removed by the projector $1-e$. Thus
ambient nonliftability alone gives no nonliftable element of
$A_B$.

\subsection*{C44-A05. The embedded axis itself always has a characteristic-zero lift}

\textbf{Proposition.}
Let $k=\overline{\mathbf F}_p$, let $X/k$ be an abelian variety,
let $L$ be ample, and let $s_0,s_1\in H^0(X,L)$ have a smooth
codimension-two common zero scheme $B$. There is a mixed-characteristic
discrete valuation ring $R$, with finite residue field contained in $k$,
an abelian scheme $\mathcal X/R$, a relatively ample line bundle
$\mathcal L$, and two sections whose common zero scheme
$\mathcal B/R$ is smooth and projective, such that the geometric
special fibre is the given $(X,L,s_0,s_1,B)$. In particular this
applies to every axis in QH9-axis, including supersingular ambient
abelian varieties. The proposition asserts no lifting of a specified
middle-dimensional algebraic class on $B$.

\textbf{Proof.}
Descend the given data to a finite subfield of $k$. The polarized
lifting theorem of Mumford--Norman--Oort gives an abelian scheme
with polarization lifting $(X,\lambda_L)$ over an integral
mixed-characteristic base. An exact statement is Oort,
\emph{Did Earlier Thoughts Inspire Grothendieck?}, Theorem (8.3).1,
p.~29. By finite-presentation descent we may use an integral base
$S$ of finite type over $\mathbf Z$, with a closed point of finite
residue field giving the prescribed geometric special fibre. The
scheme, group laws, and polarization all descend; passing to a
sufficiently large finite-type model preserves smoothness and
properness. This is the usual limit construction, as in Stacks
Project, Lemma 32.10.1, tag 01ZM, Lemma 32.8.9, tag 0C0C,
and Lemma 32.13.1, tag 081F.

The chosen line bundle also lifts after an \'etale neighborhood of
the special point: Olsson, \emph{Compactifications of moduli of
abelian varieties: an introduction}, Lemma 2.1.8 and its proof,
pp.~299--300. Its proof preserves the chosen special-fibre line
bundle. Alternatively, after choosing a representative of the
lifted polarization, the difference from $L$ lies in the dual
abelian variety and lifts after a smooth base change. Take an
integral component through the chosen point and still dominating
$S$; this is possible by flatness of the \'etale neighborhood.
Write $\pi:\mathcal X\to S$ for the resulting family and
$\mathcal L$ for its ample line bundle. For an ample line bundle
on an abelian scheme,
\[
 R^i\pi_*\mathcal L=0\quad(i>0),\qquad
 E=\pi_*\mathcal L\ \text{is locally free and commutes with base change}.
\]
The exact family statement is Olsson, \S2.1.4, p.~299.

The total space $W$ of the vector bundle $E\oplus E$ carries a
universal pair of sections on $\mathcal X\times_S W$. It is
integral of characteristic zero and of finite type over
$\mathbf Z$. The given pair determines a closed point $b\in W$
after a finite residue-field extension. Let $\mathcal Z$ denote
the universal common zero scheme. On $\mathcal Z$, failure of the
two relative differentials to have rank two is closed. Its image
in $W$ is closed because $\mathcal Z\to W$ is proper. The
complement $U$ contains $b$, since $B$ is smooth of codimension
two. The relative Jacobian criterion makes $\mathcal Z_U\to U$
smooth of relative dimension $\dim X-2$, and it is projective.
Thus $U$ is an integral finite-type mixed-characteristic
parameter scheme carrying a smooth projective lift of the
specified geometric fibre.

For completeness one can pass to an arithmetic trait without
changing that geometric fibre. Let $A=\mathcal O_{U,b}$, with
maximal ideal $\mathfrak m$, and put $d=\dim A$. This is an
excellent local domain, $0\ne p\in\mathfrak m$, and its residue
field is finite. Choose $x_1,\ldots,x_{d-1}$ so that
$(p,x_1,\ldots,x_{d-1})$ is $\mathfrak m$-primary. Set
$I=(x_1,\ldots,x_{d-1})$. Some minimal prime $P$ of $I$ does
not contain $p$: otherwise $p\in\sqrt I$ would give
$\sqrt I=\mathfrak m$, contrary to Krull's height theorem.
The local domain $A/P$ has dimension one, since $p$ is nonzero
and generates an ideal with maximal radical in that quotient.
Its normalization is finite by excellence; localization at a
point above its closed point is a discrete valuation ring $R$
with characteristic-zero fraction field and finite residue field.
Pulling back the family along
$\operatorname{Spec}R\to\operatorname{Spec}A\to U$ proves the
claim. The residue extension embeds in $k$, so the geometric
special fibre is the prescribed one.\hfill$\square$

\paragraph{Programme location and remaining obstruction.}
This completes the geometric lifting substep of QH9-axis(a),
in the schemes-over-$\mathbf Z$ framework. The unresolved issue
is surjectivity, for some such lift, of specialization on the
specified rational algebraic middle classes. Smooth proper base
change is an isomorphism on cohomology; it is not a lifting
theorem for cycles. The sign transfer to classes which do lift
is exactly the already settled H9-T8a specialization argument.
Cote 44, Th.~5, scan 004, and the primitive margin on scan 008
therefore retain the same cycle-level gate.

\section{Unresolved axis steps after the ordered attacks}

\paragraph{Lifting.}
For any specified smooth projective mixed-characteristic lift
of $(B,h_B)$ and any rational algebraic class on the geometric
generic fibre specializing to $w\in A_B$, compatibility of
cycle classes with smooth proper specialization carries
$L_Bw=0$ back to the same primitive relation on the generic
fibre. The specialization is an isomorphism on $\ell$-adic
cohomology, so a nonzero $w$ has nonzero generic class, and the
intersection number is unchanged. H9-T8a then gives
$(-1)^r\int_Bw^2>0$. This is the carried lifting argument,
not a surjectivity theorem for Chow specialization.

No proof that every $w\in A_B$ lifts with $B$ is obtained.
Nor is a specific cohomological class in $A_B$ proved
nonliftable in this contribution. Attempts using nonliftable
ambient Tate classes fail because their projection to $A_B$
is zero. Agugliaro, \emph{Examples for the standard conjecture
of Hodge type}, Doc. Math. 31 (2026), Proposition 5.12,
pp.~876--877, supplies a rank-two nonliftable Tate block for
CM lifts of certain abelian varieties; Remark 3.5, p.~870,
explains the passage to arbitrary abelian lifts. Even after
032 supplies algebraicity, these remain ambient examples.
They are not examples for the axis kernel.

\paragraph{Correspondences.}
C44-A02 and JAC1964 H9-T13 eliminate graph/ambient transfers.
A general correspondence from $B$ to an abelian variety could
still act nontrivially, but a splitting preserving the signed
geometric pairing has not been constructed. Merely splitting
an abstract Tate class from $h(B)$ would not identify that
pairing with a positive form.

\paragraph{Newton slopes.}
C44-A03 gives the correct finite-field vanishing test.
No Frobenius polynomial satisfying that test has been exhibited
for every special axis, and axes with $A_B\ne0$ necessarily
fail it. A slope witness in the odd pencil variation does not
settle the even fixed-axis sign.

\paragraph{Current gate.}
The unresolved assertion remains weak nonnegativity of
$(-1)^r\int_Bw^2$ for every algebraic $w\in\ker\kappa_*$.
Conditional on 032/Milne, C44-A01 identifies it with the weak
Hodge-standard assertion for $B$. No wrong-sign cycle is
constructed. No main question is marked closed by these
route boundaries.

% END work/axis/axis_contribution.tex

% BEGIN work/axis/axis_correspondence_obstruction.tex
% C44_POS /root/axis_independent; GPT-6 Astra Pro; effort not exposed.
\subsection*{C44-A02b. Divisor-generated correspondence tests annihilate the axis kernel}

\textbf{Proposition.}
Let $\kappa:B\hookrightarrow X$ be a smooth Lefschetz axis in an
abelian $(2r+2)$-fold, $r\ge2$, and let $A$ be any abelian variety.
Every rational divisor cohomology class on $B\times A$ is the
restriction of a rational divisor cohomology class on $X\times A$.
Hence every class
in the algebra generated by these divisors is an ambient
restriction. Any algebraic correspondence from $B$ to $A$ whose
cohomology class is such a restriction acts as zero on $A_B$.
In particular this holds for every correspondence class in the
divisor-generated algebra of $B\times A$.

\textbf{Proof.}
Choose $b_0\in B$ and the origin of $A$. For a line bundle $D$
on $B\times A$, remove its restrictions from the two factors.
The remaining normalized line bundle is a family of
algebraically trivial line bundles on $A$, parametrized by $B$,
and therefore is the pullback of the Poincar\'e bundle along
$f\times1_A$, for a morphism $f:B\to\widehat A$ with
$f(b_0)=0$. This is the defining universal property of the dual
abelian variety.

The map $\operatorname{Alb}(B)\to X$ induced by
$b\mapsto\kappa(b)-\kappa(b_0)$ is an isogeny, as used and
proved cohomologically in C44-A02. After multiplying by a
positive integer $N$, the map $[N]f$ therefore extends to a
morphism $X\to\widehat A$. Pullback of the Poincar\'e bundle
gives a line bundle on $X\times A$ restricting to the $N$-th
tensor power of the normalized part of $D$.
The factor from $A$ already extends. The factor from $B$ has
a rational algebraic divisor-class extension to $X$ by the
low-degree algebraic inverse in C44-A01, with $j=1$.
These three extensions prove the assertion for rational
cohomology classes of divisors. Taking products proves the
assertion about the divisor-generated algebra.

For the vanishing, write $\rho=\kappa\times1_A$ and
$[\Gamma]=\rho^*\gamma$ for an ambient cohomology class
$\gamma$ of the appropriate degree. For $w\in A_B$, the
projection formula gives
\[
 \begin{split}
 \Gamma_*w
 &=p_{A,*}\rho_*\bigl(p_B^*w\,\rho^*\gamma\bigr)\\
 &=p_{A,*}\bigl(p_X^*(\kappa_*w)\,\gamma\bigr)=0.
 \end{split}
\]
Here the first $p_A$ on the right is the projection from
$X\times A$. This calculation in fact proves the stated
vanishing for every ambient-restriction correspondence, without
the divisor-generation hypothesis.\hfill$\square$

\paragraph{Scope.}
This is a constraint on proposed curve tests in the manner of
cote~44 Th.~4 (scan~002), when used to attack the axis term of
Th.~5 (scan~004). It extends the graph-map obstruction of
C44-A02. It does not exclude general correspondences. In
particular, an extending injection of the kind in C44-A08 must
use correspondence classes involving the new middle
cohomology of $B$; ambient divisor and Poincar\'e operations
alone cannot supply it. It does not assert that any such new
correspondence is nonalgebraic or nonexistent. The universal
Poincar\'e property used above is Milne, \emph{Abelian Varieties},
course notes v.~2.00 (2008), I \S8, printed pp.~36--37
(definition of the dual abelian variety and universal sheaf);
the other inputs are C44-A01 and
C44-A02, already checked in this run.

% END work/axis/axis_correspondence_obstruction.tex

% BEGIN work/axis/axis_forced_tate_root.tex
% C44_POS /root/axis_independent; GPT-6 Astra Pro; effort not exposed.
\subsection*{C44-A03b. A forced central Frobenius root in the first open dimension}

\textbf{Theorem.}
Let $X/\mathbf F_q$ be an abelian variety of dimension $2r+2$,
$r\ge2$, with an ample line bundle $L$ of class $h$. Let
$B\subset X$ be a smooth complete intersection of two members
of $|L|$, defined over $\mathbf F_q$, and put
\[
 V_B=\ker\bigl(\kappa_*:H^{2r}(B,\mathbf Q_\ell)
       \longrightarrow H^{2r+4}(X,\mathbf Q_\ell)(2)\bigr),
 \qquad d=\int_Xh^{2r+2}.
\]
Then
\[
 \begin{split}
 \dim V_B
 &=(2r+1)d+2\binom{4r+3}{2r-1}-\binom{4r+4}{2r}\\
 &=(2r+1)d-\frac1{r+1}\binom{4r+4}{2r}.
 \end{split}
 \tag{C44.A03b.1}
\]
This dimension is odd exactly when $r=2^s-2$ for an integer
$s\ge2$. In these dimensions the Frobenius on $V_B$ has an
eigenvalue $+q^r$ or $-q^r$. In particular the potential Tate
eigenspace is nonzero after an extension of degree at most two,
and $V_B$ has a Newton slope $r$.

For an abelian sixfold, $r=2$, formula (C44.A03b.1) is
\[
                  \dim V_B=5d-165,
\]
which is always odd. Thus the no-potential-Tate-root criterion
C44-A03 can never prove $A_B=0$ for any axis in the first open
dimension. The same obstruction occurs for ambient dimensions
$6,14,30,62,\ldots$.

\textbf{Proof.}
The tangent sequence of the smooth complete intersection gives
\[
 c(T_B)=(1+h|_B)^{-2},\qquad
 \chi(B)=\int_Bc_{2r}(T_B)=(2r+1)d.
\]
Here the Euler/Chern identity can be read entirely in Weil
cohomology: the self-intersection of the diagonal is both the
alternating sum of the Betti numbers, by K\"unneth and Poincar\'e
duality, and the degree of the top Chern class of its normal
bundle $T_B$, by self-intersection (Fulton, \emph{Intersection
Theory}, \S6.3). No characteristic-zero comparison is needed.

Weak Lefschetz gives $b_i(B)=b_i(X)$ for $i<2r$, and Poincar\'e
duality determines the higher Betti numbers from those lower
ones. Since $b_i(X)=\binom{4r+4}{i}$, we obtain
\[
 b_{2r}(B)
 =\chi(B)-2\sum_{i=0}^{2r-1}(-1)^i\binom{4r+4}{i}
 =(2r+1)d+2\binom{4r+3}{2r-1}.
\]
The restriction of $H^{2r}(X)$ is injective and splits off
orthogonally, by $\kappa_*\kappa^*=L_h^2$ and hard Lefschetz.
Subtracting its dimension proves the first equality of
(C44.A03b.1); the second is a binomial identity.
The exterior-algebra description of abelian cohomology is
Milne, \emph{Abelian Varieties}, course notes v.~2.00 (2008),
I Theorem~12.1 (printed pp.~55--56).

Riemann--Roch on $X$ gives $d=(2r+2)!\chi(L)$, so $d$ is even.
Reducing the first formula modulo two gives
\[
 \dim V_B\equiv\binom{4r+4}{2r}
                 \equiv\binom{2r+2}{r}\pmod2.
\]
The second congruence follows from $(1+x)^2=1+x^2$ over
$\mathbf F_2$. More generally, expanding
$(1+x)^n=\prod_{j:\,n_j=1}(1+x^{2^j})$ in binary shows that
$\binom{2r+2}{r}$ is odd exactly when adding $r$ and $r+2$
has no binary carries, equivalently $r\mathbin{\&}(r+2)=0$.
This requires $r$ to be even. Writing $r=2k$ reduces it to
$k\mathbin{\&}(k+1)=0$, which says that $k$ consists solely
of trailing binary ones: $k=2^{s-1}-1$.
Thus the stated parity criterion holds.

The middle intersection form restricts nondegenerately to $V_B$
by the same orthogonal splitting. Frobenius $F$ is a similitude
with multiplier $q^{2r}$, so $q^{-r}F$ is an orthogonal
transformation of $V_B$. Over an algebraic closure of the
coefficient field its eigenvalues occur in reciprocal pairs,
counted with multiplicity. In odd dimension at least one
eigenvalue is fixed by inversion, hence is $+1$ or $-1$.
This argument does not assume semisimplicity. An eigenvalue
always has a nonzero actual eigenspace, and its square equals
one. Consequently $F^2$ has eigenvalue $q^{2r}$ on a nonzero
vector in $V_B$, as claimed. The same characteristic polynomial
occurs in crystalline cohomology: apply Katz--Messing,
\emph{Some consequences of the Riemann hypothesis for varieties
over finite fields}, Invent.~Math.~23 (1974), Theorem~1(i) and
the preceding observation that crystalline cohomology satisfies
its hypotheses, p.~74, to $B$ and $X$, and take their middle
quotient. The corresponding root has normalized $p$-adic
valuation $r$, so it is also a forced central Newton
slope.\hfill$\square$

\paragraph{What this does and does not settle.}
No algebraicity of the forced eigenvector is asserted. If the
Tate conjecture for codimension-$r$ cycles on $B$ holds over
the relevant finite extension, the algebraic variable projector
$1-\kappa^*(L_h^2)^{-1}\kappa_*$ shows that $A_B\ne0$.
That is a further hypothesis about $B$, not a consequence of
032 on the ambient abelian variety. The theorem is compatible
with geometric-generic vanishing in JAC1964 JAC-P09: the
geometric generic parameter is not a finite-field closed point.

For route~(c) of QH9-axis, the theorem rules out the full
fixed-axis variable-part vanishing certificate in these
dimensions. A slope outside the interval in H9-T5 can coexist
with this forced central root. It remains possible to isolate
a smaller factor by an independently justified projector, but
no such universal factor is supplied here. In cote~44 terms,
this constrains the proposed Newton route to Th.~5, scan~004,
and its primitive axis summand; it does not decide its sign.

% END work/axis/axis_forced_tate_root.tex

% BEGIN work/axis/root_explicit_axes.tex
\section{An explicit nonzero axis subspace in every characteristic}

\subsection*{C44-A04. Axes containing fibres of an abelian product}
\begin{theorem}[C44-A04]\label{thm:product-axis}
Let $r=m-1\ge2$, let $S$ and $T$ be abelian varieties of dimensions
$r$ and $r+2$, and put $X=S\times T$. Choose very ample classes
$s$ and $t$ on $S$ and $T$ and let $h=s+t$ on $X$. Let
$B\subset X$ be a smooth complete intersection of two divisors of
class $h$. Suppose that it contains the distinct fibres
$S_i=S\times\{t_i\}$ for $1\le i\le N$. Put
\[
 d_S=\int_Ss^r,\qquad d_T=\int_Tt^{r+2},\qquad
 a=\frac1{d_T}\sum_{j=0}^r(-1)^j(j+1)s^jt^{r-j},
 \qquad w_i=[S_i]_B-\kappa^*a.
 \tag{C44.A04.1}
\]
Then $w_i\in A_B$ and the exact signed Gram matrix is
\[
 \left((-1)^r\int_B w_iw_j\right)_{i,j}
       =(r+1)d_S\left(I_N-\frac1{d_T}{\bf1}{\bf1}^{t}\right).
 \tag{C44.A04.2}
\]
Smoothness of $B$ forces $N\le d_T$. Hence this matrix is positive
semidefinite in every characteristic, without 032. If $N<d_T$ it is
positive definite, and the $w_i$ are linearly independent rational
algebraic classes. In particular a single contained fibre supplies a
nonzero axis class of the required sign.

Smooth such axes containing a prescribed fibre exist after replacing
$s,t$ by sufficiently divisible positive multiples. The two defining
sections may be chosen to give a genuine Lefschetz pencil. Thus the
construction includes products of supersingular elliptic curves and
gives special axes with $A_B\ne0$ over $\overline{\F}_p$ for every $p$.
The assertion concerns the displayed subspace, not all of $A_B$.
\end{theorem}
\begin{proof}
All computations are in rational Weil cohomology, with their usual
cycle representatives. Since $s^{r+1}=0$ and $t^{r+3}=0$, the finite
identity $(1+z)^2\sum_{j=0}^r(-1)^j(j+1)z^j\equiv1\pmod{z^{r+1}}$
gives
\[
                      h^2a=t^{r+2}/d_T=[S_i]_X.
\]
Thus $\kappa_*w_i=0$ by self-intersection, and $w_i\in A_B$.
The normal bundle of $S_i$ in $X$ is trivial of rank $r+2$.
Because $B$ is smooth along $S_i$, its normal sequence is exact:
\[
 0\longrightarrow N_{S_i/B}\longrightarrow\mathcal O_S^{r+2}
   \longrightarrow L_S\oplus L_S\longrightarrow0.
\]
Here $c_1(L_S)=s$. Consequently
$c(N_{S_i/B})=(1+s)^{-2}$ and
\[
                 \int_B[S_i]^2=(-1)^r(r+1)d_S.
\]
Distinct fibres are disjoint, so their intersection is zero.
Restriction of $a$ to any $S_i$ has degree
$(-1)^r(r+1)d_S/d_T$; the same is true of
$\int_B(\kappa^*a)^2=\int_Xa h^2a$.
Expanding each $w_iw_j$ therefore gives (C44.A04.2).

We prove the bound on the number of contained fibres without using a
positivity conjecture. Write the two defining sections in a basis of
$H^0(S,L_S)$ as
$F_\nu=\sum_\alpha\sigma_\alpha\tau_{\nu\alpha}$,
where $\tau_{\nu\alpha}\in H^0(T,L_T)$ and $\nu=1,2$.
Let $Z\subset T$ be the common zero scheme of all coefficients
$\tau_{\nu\alpha}$. Its geometric points are exactly the fibres
contained in $B$. Suppose $T_zZ$ has positive dimension at a point
$z\in Z$. Differentiating the coefficients shows that the normal
surjection along $S\times\{z\}$ factors through
\[
       \mathcal O_S^{r+1}\longrightarrow L_S\oplus L_S
\]
(or through a trivial bundle of still smaller rank). Its kernel
would have rank at most $r-1$ but total Chern class $(1+s)^{-2}$.
Its $r$th Chern class would therefore be both zero, by rank, and
$(-1)^r(r+1)s^r$, whose degree is nonzero. This is impossible.
Thus $Z$ has zero tangent space at every geometric point; it is a
finite reduced scheme. Choose $r+2$ general linear combinations of
the coefficient sections. Since their common base scheme is finite,
successive dimension cuts yield a zero-dimensional complete
intersection containing $Z$. Its degree is $d_T$, so
$N\le\#Z\le d_T$.
The eigenvalues of the matrix in parentheses in (C44.A04.2) are
$1$ with multiplicity $N-1$ and $1-N/d_T$ with multiplicity one.
This proves all the asserted signs and, for $N<d_T$, independence.

For existence take the prescribed smooth centre $S_0=S\times\{0\}$.
For sufficiently positive $L_S,L_T$, the linear system
$H^0(X,I_{S_0}\otimes(L_S\boxtimes L_T))$ separates second jets
at each point off $S_0$, separates the first jets at each pair of
distinct points off $S_0$, and maps onto the conormal first jets
along $S_0$. These are uniform assertions for all such points:
apply relative Serre vanishing to the universal length-bounded
jet subschemes and the fixed coherent ideal $I_{S_0}$ (or use
high powers of product sections). A common sufficiently high
power satisfies all three conditions.
The latter assertion also follows directly from product sections:
sections of $L_S$ generate its fibres and sections of $I_0L_T$
generate the cotangent space of $T$ at $0$ after taking a high power.
For two general sections, the locus on $S_0$ where their normal
derivatives have rank at most one has expected codimension $r+1$;
the evaluation to the two normal covectors is surjective at each
point, so the incidence dimension estimate, with $\dim S_0=r$,
makes this locus empty. Away from $S_0$, jet separation and
Bertini give a smooth codimension-two intersection. Here are the
dimension counts for the pencil condition. Write $n=\dim X=2r+2$.
At a specified point off $S_0$, singularity imposes $n+1$
independent first-jet equations on a section. Degeneracy of its
quadratic part imposes an additional proper equation, by the
second-jet surjectivity. Letting the point vary gives a bad locus
of codimension at least two in the projective linear system.
Singularities at two specified distinct points impose $2(n+1)$
independent equations; their parameter space has dimension $2n$,
again giving codimension at least two. For singularity along the
imposed centre, the normal covector must vanish; this imposes
$r+2$ conditions at a specified point of $S_0$, and $\dim S_0=r$.
Its incidence image has codimension at least two as well. A general
projective line avoids these bad loci. Smoothness of its axis
also prevents the line from being tangent to the discriminant at
a single ordinary quadratic point: such tangency says that the
other pencil section vanishes at that point, putting a singular
member's zero differential into the axis's defining pair.
Thus the line gives a Lefschetz pencil, with distinct singular
values and no singularity on $S_0$. Since $n$ is even, the
nondegenerate quadratic forms are a nonempty open also in
characteristic two. These are the ordinary Lefschetz local models.
All conditions used are nonempty Zariski-open conditions over the
algebraically closed field. Over $\overline{\F}_p$ the chosen
sections are defined over a finite extension. This proves existence.
\end{proof}

\paragraph{C44-A04a: a concrete sixfold normalization.}
Take $X=E^6$, $S=E^2$, $T=E^4$ and
$h=d\sum_{i=1}^{6}[0]_i$, with $d$ sufficiently large for the
construction. For an axis containing $E^2\times\{0\}$,
\[
 d_S=2d^2,\qquad d_T=24d^4,\qquad
 \int_Bw^2=6d^2-\frac1{4d^2}>0.
\]
The integral cycle $\widetilde w=d_T[S]_B-
\kappa^*(t^2-2st+3s^2)$ has
\[
        \int_B\widetilde w^2=144d^6(24d^4-1)>0.
\]
By H9-T7 the corresponding height classes have these exact heights
in its normalization. One may choose $E$ supersingular in every
characteristic. These are positive special-axis classes, not a
classification of nonliftable classes and not a wrong-sign example.

\paragraph{Archive and source location.}
The height interpretation is JAC1964 H9-T7, (H9.18), for cote 44,
Th.~5, scan 004. The primitive qualification is the margin on scan
008. The normal-bundle and self-intersection calculations use Fulton,
\emph{Intersection Theory}, Section 6.3, self-intersection formula,
and the codimension-two blowup normalization already checked in
H9-T7. The existence argument uses the ordinary jet construction
underlying SGA~7~II, XVII, 1.1--1.2 (pp.~213--214), 2.5
(p.~217), 3.2.1--3.2.4 (p.~219), 3.3 (p.~221), 3.7
(p.~234), and 4.2--4.3 (pp.~236--239); its imposed-centre
incidence argument is provided here, not attributed as a separate
statement to SGA. In cote 16 language the calculation verifies $I$
on a specified algebraic primitive subspace. It does not establish
$I(B)$ on its complement or close QH9-axis.

% END work/axis/root_explicit_axes.tex

% BEGIN work/axis/axis_boundary_relation.tex
% C44_POS /root/axis_independent; GPT-6 Astra Pro; effort not exposed.
\subsection*{C44-A04b. The equality case of the product-fibre Gram bound}

\textbf{Proposition.}
In C44-A04, suppose the smooth axis $B\subset S\times T$ contains
$N=d_T$ distinct fibres $S_i=S\times\{t_i\}$, where
$\dim S=r$, $\dim T=r+2$, $s=c_1(L_S)$,
$t=c_1(L_T)$ and $d_T=\int_Tt^{r+2}$. Put
\[
 a=\frac1{d_T}\sum_{j=0}^r(-1)^j(j+1)s^jt^{r-j},
 \qquad w_i=[S_i]_B-\kappa^*a.
\]
Then
\[
 \sum_{i=1}^{d_T}w_i=0
 \tag{C44.A04b.1}
\]
already in the rational Chow group of $B$, for these displayed
representatives. This is the only linear relation in cohomology.
Thus the signed H9-T7 form is positive definite on the actual
span of all the $w_i$, of dimension $d_T-1$. Together with
C44-A04, the form is positive definite on the span of any
collection of contained product fibres, including the boundary
case $N=d_T$.

\textbf{Proof.}
Let $W\subset H^0(T,L_T)$ be the span of the coefficient sections
of the two equations of $B$ in a basis of $H^0(S,L_S)$.
C44-A04 proves that their common zero scheme $Z\subset T$ is
finite and reduced and that its length is at most $d_T$.
It contains the $d_T$ given points, so these are exactly $Z$.
Choose $g=r+2$ general sections $t_1,\ldots,t_g\in W$.
They form a proper zero-dimensional complete intersection on
$T$, of length $d_T$, containing $Z$ scheme-theoretically.
Hence their common zero scheme equals $Z$.

The Koszul resolution of $I_ZL_T$ has terms
$\bigwedge^j k^g\otimes L_T^{1-j}$ for $1\le j\le g$.
For $j\ge2$, Serre duality and ample-line-bundle vanishing on an
abelian variety give
$H^i(T,L_T^{1-j})=0$ for $i<g$. Consequently its degree-zero
cohomology gives
\[
 H^0(T,I_ZL_T)=\langle t_1,\ldots,t_g\rangle.
 \tag{C44.A04b.2}
\]
The same sections are linearly independent, since otherwise
$g-1$ equations would have a zero-dimensional common zero scheme
on a $g$-fold, contrary to the height theorem. Every coefficient
of the equations of $B$ belongs to (C44.A04b.2). Thus these two
equations can be written as $U(t_1,\ldots,t_g)^t$, with a matrix
of sections on $S$ giving
\[
 U:\mathcal O_S^{g}\longrightarrow L_S^{\oplus2}.
\]
This map is surjective everywhere. Indeed, at any point of the
reduced complete intersection $Z$, the differentials of the
$t_j$ form a basis of the conormal space, and smoothness of $B$
along the corresponding entire fibre $S\times\{t_i\}$ says
precisely that $U$ has rank two everywhere on $S$.

Let $K=\ker U$, a rank-$r$ bundle with
$c(K)=(1+s)^{-2}$. On $B$, the vector of the $t_j$ is a section
of $K\boxtimes L_T$. Its zero scheme is exactly $S\times Z$.
Since it has codimension $r$ in the smooth variety $B$, this is
a regular section. The top Chern class formula therefore gives
\[
 \begin{split}
 \sum_{i=1}^{d_T}[S_i]_B
 &=c_r(K\boxtimes L_T)\\
 &=\kappa^*\sum_{j=0}^r(-1)^j(j+1)s^jt^{r-j}
 =d_T\kappa^*a.
 \end{split}
\]
This proves (C44.A04b.1) in Chow. By C44-A04 the signed Gram
matrix is $(r+1)d_S(I-J/d_T)$, whose nullspace at $N=d_T$ is
exactly the constant vector. Any cohomological relation is in
this nullspace. The exhibited relation is therefore the only
one, and the induced form on the quotient is positive definite.
For fewer than $d_T$ fibres the matrix itself is positive
definite by C44-A04.\hfill$\square$

\paragraph{Inputs and scope.}
The ample vanishing and duality used in (C44.A04b.2) are the
standard abelian-variety index theorem; see Milne,
\emph{Abelian Varieties}, course notes v.~2.00 (2008),
I Theorem~11.1(c), printed p.~54. The regular-section formula is
Fulton, \emph{Intersection Theory}, \S14.1. The apparent isotropic
vector in the boundary Gram matrix therefore cannot supply a
nonliftable nonzero axis class. This sharpens the explicit
positive subspace of cote~44 Th.~5, scan~004, without deciding
QH9-axis on its complement.

% END work/axis/axis_boundary_relation.tex

% BEGIN work/axis/axis_asymptotic_centres.tex
% C44_POS /root/axis_independent; GPT-6 Astra Pro; effort not exposed.
\section{Prescribed smooth centres: an unconditional asymptotic axis sign}

\subsection*{C44-A06. Fixed smooth centres become positive after increasing the polarization}

\textbf{Theorem.}
Let $X$ be an abelian variety of dimension $2r+2$, $r\ge2$, over
an algebraically closed field of arbitrary characteristic. Fix an
ample line bundle $L_0$ with class $h_0$, and pairwise disjoint
smooth closed subvarieties $Z_1,\ldots,Z_s\subset X$ of dimension
$r$. For every sufficiently large integer $N$ there is a Lefschetz
pencil in $|L_0^{\otimes N}|$ whose smooth axis $B_N$ contains every
$Z_i$. Put $\kappa:B_N\hookrightarrow X$ and
\[
 a_i=(L_{h_0}^{2})^{-1}[Z_i]_X,
 \qquad w_i=[Z_i]_{B_N}-N^{-2}\kappa^*a_i.
 \tag{C44.A06.1}
\]
These are rational algebraic classes and $w_i\in A_{B_N}$. Define
\[
 \begin{split}
 P_i(N)&=\sum_{j=0}^r(-1)^j(r-j+1)N^{r-j}
       \int_{Z_i}c_j(N_{Z_i/X})h_0^{r-j},\\
 G_{ij}&=\int_X a_i a_jh_0^2.
 \end{split}
 \tag{C44.A06.2}
\]
The exact signed Gram matrix is
\[
 \left((-1)^r\int_{B_N}w_iw_j\right)_{i,j}
       =\operatorname{diag}(P_i(N))-(-1)^rN^{-2}G.
 \tag{C44.A06.3}
\]
In particular it is positive definite for all sufficiently large
$N$, independently of 032. More explicitly, let
\[
 \begin{split}
 d_i&=(r+1)\int_{Z_i}h_0^r>0,\qquad d_{\min}=\min_i d_i,\\
 S_i&=\sum_{j=1}^r\left|(r-j+1)
             \int_{Z_i}c_j(N_{Z_i/X})h_0^{r-j}\right|,\\
 R&=\max_i\sum_j|G_{ij}|.
 \end{split}
\]
Whenever a smooth axis containing these centres exists, the matrix
is positive definite provided
\[
 N\ge1,\qquad N\ge\max_i\frac{2S_i}{d_i},
 \qquad N^{r+2}>\frac{2R}{d_{\min}}.
 \tag{C44.A06.4}
\]
Thus the $w_i$ are nonzero and linearly independent in that range.

\textbf{Proof.}
The inverse Lefschetz operator in (C44.A06.1) is algebraic because
$B(X)$ is known for abelian varieties, without the Hodge standard
conjecture. Since $[B_N]=N^2h_0^2$, self-intersection gives
$\kappa_*\kappa^*=N^2L_{h_0}^2$, so $\kappa_*w_i=0$.
The normal sequence of the smooth inclusions $Z_i\subset B_N\subset X$
is
\[
 0\longrightarrow N_{Z_i/B_N}\longrightarrow N_{Z_i/X}
 \longrightarrow (L_0^{\otimes N}|_{Z_i})^{\oplus2}
 \longrightarrow0.
\]
Therefore
$c(N_{Z_i/B_N})=c(N_{Z_i/X})(1+Nh_0)^{-2}$.
The self-intersection formula and the coefficient identity
$(1+x)^{-2}=\sum_{k\ge0}(-1)^k(k+1)x^k$ show that
$(-1)^r\int_{B_N}[Z_i]^2=P_i(N)$.
For distinct $i,j$ the intersection is zero because the centres
are disjoint. Projection formula gives
\[
 \int_{B_N}[Z_i]\kappa^*a_j=G_{ij},\qquad
 \int_{B_N}\kappa^*a_i\kappa^*a_j=N^2G_{ij}.
\]
Expanding (C44.A06.1) proves (C44.A06.3), with no asymptotic
approximation.

For $N\ge1$ the lower-degree terms in $P_i$ have total absolute
value at most $S_iN^{r-1}$. The first two inequalities in
(C44.A06.4) imply $P_i(N)\ge\tfrac12d_iN^r$.
For a real vector $x=(x_i)$, symmetry of $G$ and
$2|x_ix_j|\le x_i^2+x_j^2$ give
$|x^tGx|\le R\sum_i x_i^2$. Hence the right side of
(C44.A06.3), evaluated at $x$, is at least
\[
 \left(\tfrac12d_{\min}N^r-RN^{-2}\right)\sum_i x_i^2,
\]
which is positive for $x\ne0$ under the last inequality.

It remains to justify existence with the prescribed centres.
Let $Z=\coprod_iZ_i$ and $V_N=H^0(X,I_ZL_0^{\otimes N})$.
For $N$ large, Serre vanishing for $I_Z^2L_0^{\otimes N}$ and
global generation of $(I_Z/I_Z^2)L_0^{\otimes N}$ make the
normal first-jet map from $V_N$ onto the conormal fibre at each
point of $Z$ surjective. Off $Z$, $V_N$ separates second jets at
one point and simultaneous first jets at any two distinct points.
To see uniformity, first choose $a$ for which $I_ZL_0^{\otimes a}$
is globally generated. At any one or two points off $Z$, a general
section of this sheaf is nonzero at all those points. Multiplying
it by sections of a sufficiently high power of $L_0$, which
separates the indicated finite jets, gives the required maps from
$V_N$. High powers of a very ample bundle separate these jets;
handling the finitely many residue classes of $N$ gives the claim
for every sufficiently large $N$.

Set $n=2r+2$. For a pair of sections, rank at most one for the
two normal covectors at a fixed point of $Z$ has codimension
$r+1$ in the pair parameter space. Since $\dim Z=r$, a general
pair has rank two everywhere on $Z$. At a point off $Z$, the
conditions that both section values vanish and their two first
derivatives have rank at most one have codimension
$2+(n-1)=n+1$. Varying the point therefore still leaves a proper
bad locus. A general pair consequently defines a smooth
codimension-two axis.

The second-jet and two-point first-jet surjectivity also give the
Lefschetz property. A singular section at a fixed point off $Z$
imposes $n+1$ independent linear conditions. A degenerate
quadratic part imposes at least one more; after varying the point,
these sections have codimension at least two in $\mathbf P(V_N)$.
Sections singular at two distinct points also have codimension at
least two, since the fixed-pair conditions have codimension
$2(n+1)$ and the pair of points has dimension $2n$.
A general pencil avoids the closures of both loci. Along $Z$ no
member is singular, because the two normal covectors are
independent. Smoothness of the axis also excludes singularities
of a member elsewhere on the axis. Thus every singular member
has just one ordinary quadratic singularity. In characteristic
two an invertible Hessian is possible here because $n$ is even.
These are the local incidence arguments of SGA~7~II, XVII,
1.1--1.2, 3.2.1--3.2.4, 3.3, 3.7, and 4.2--4.3, with the
imposed-centre conditions verified above. The generic member is
smooth by the same first-jet count. All the required opens in the
pencil parameter space are nonempty, so the construction exists.
Over $\overline{\mathbf F}_p$ a chosen pencil is defined over a
finite extension.\hfill$\square$

\paragraph{Scope and programme location.}
This gives a new unconditional positive subspace of the precise
H9-T7 axis height form, in cote~44 Th.~5, scan~004, with the
primitive qualification of scan~008. It allows prescribed smooth
centres whose own lifting is not known, but it does not prove that
the projected classes are nonliftable. Its threshold depends on
the chosen centres, and it changes the axis and the polarization.
It supplies neither a uniform bound over all cycles in one fixed
$B$ nor positivity on the rest of $A_{B_N}$. It therefore does not
close QH9-axis or $I(B)$ in cote~16 terminology. The geometric
existence argument has the same characteristic-two scope as
C44-A04. All normal and self-intersection operations are those of
Fulton, \emph{Intersection Theory}, \S6.3.

% END work/axis/axis_asymptotic_centres.tex

% BEGIN work/axis/axis_product_curves.tex
% C44_POS /root/axis_independent; GPT-6 Astra Pro; effort not exposed.
\section{A fixed-polarization positive axis class from a product of curves}

\subsection*{C44-A07. Products of smooth curves in abelian threefold factors}

\textbf{Theorem.}
Let $A,A'$ be abelian threefolds over an algebraically closed field of
arbitrary characteristic, with ample line bundles $L,M$ of classes
$s,t$. Put $X=A\times A'$, $h=s+t$, and let $B\subset X$ be a smooth
complete intersection of two sections of $L\boxtimes M$.
Suppose that $B$ contains $Z=C\times D$, where $C\subset A$ and
$D\subset A'$ are smooth connected curves. Write
\[
 a=\int_Cs,\quad b=\int_Dt,\quad
 V=\int_As^3,\quad W=\int_{A'}t^3,\quad
 u=a/V,\quad v=b/W.
\]
Then $0<u,v\le1$, and the rational algebraic class
\[
 w=[Z]_B-\kappa^*(L_h^2)^{-1}[Z]_X\in A_B
 \tag{C44.A07.1}
\]
satisfies the unconditional strict lower bound
\[
 \int_Bw^2\ \ge\ ab(2u+2v-3uv)>0.
 \tag{C44.A07.2}
\]
In particular this class is nonzero. When $B$ is a Lefschetz axis,
this is the positive H9-T7 axis sign ($m=3$) on the line spanned by
$w$. Neither 032 nor a lifting of the embedded curves is used.

\paragraph{1. Smoothness forces two independent conormal directions on each curve.}
Set $K_C=H^0(A,I_CL)$ and $K_D=H^0(A',I_DM)$. By the degree-zero
K\"unneth formula and elementary linear algebra, the sections of
$L\boxtimes M$ vanishing on $Z$ form
\[
 K_C\otimes H^0(A',M)+H^0(A,L)\otimes K_D.
 \tag{C44.A07.3}
\]
Let $e_C$ be the generic rank of the evaluation of conormal derivatives
\[
 K_C\otimes\mathcal O_C\longrightarrow N_{C/A}^{\vee}\otimes L|_C,
\]
and define $e_D$ similarly. The smooth normal sequence for $Z\subset
B\subset X$ gives a surjection
\[
 N_{C/A}\boxplus N_{D/A'}\longrightarrow
 (L|_C\boxtimes M|_D)^{\oplus2}.
 \tag{C44.A07.4}
\]
Its first and second summands factor through the derivatives in
(C44.A07.3). If $e_C=0$, restrict (C44.A07.4) to $C\times\{y\}$.
It would be a surjection from a trivial rank-two bundle to
$(L|_C)^{\oplus2}$, up to a constant one-dimensional factor. Its
determinant cannot be invertible because $2a>0$. Thus $e_C>0$,
and similarly $e_D>0$.

Suppose $e_C=1$. Saturate the derivative image in
$N_{C/A}^{\vee}\otimes L|_C$; the saturation is a line subbundle
$P_C$. It has a nonzero global section, so $\deg P_C\ge0$.
Put $Q_C=L|_C\otimes P_C^{-1}$ and $q=\deg Q_C\le a$.
All the derivatives in the first summand of (C44.A07.4) factor
through the line-bundle quotient $N_{C/A}\twoheadrightarrow Q_C$.
If also $e_D=1$, a rank-two quotient
$Q_C\boxplus Q_D$ would map surjectively to the rank-two target in
(C44.A07.4). Taking determinants and restricting to $C\times\{y\}$
would give $q=2a$, a contradiction. Hence $e_D=2$.

Put $c=\deg N_{C/A}=2g(C)-2$ and
$d=\deg N_{D/A'}=2g(D)-2$. The equality $e_D=2$ supplies a nonzero
section of $\det(N_{D/A'}^{\vee}\otimes M|_D)$, so $d\le2b$.
Now (C44.A07.4) factors through a surjection from the rank-three
bundle $Q_C\boxplus N_{D/A'}$. Its kernel is a line bundle. The
degree of the second Chern class of the formal quotient by
$(L|_C\boxtimes M|_D)^{\oplus2}$ must therefore vanish. It is
\[
 qd-2(ad+bq)+6ab=(q-2a)(d-2b)+2ab>0,
\]
since $q\le a$ and $d\le2b$. This is impossible. Consequently
$e_C=e_D=2$, and taking the two determinants gives
\[
 c\le2a,\qquad d\le2b.
 \tag{C44.A07.5}
\]
All saturation arguments here are on smooth curves, so the saturated
subsheaves and the corresponding quotients are vector bundles.

\paragraph{2. A curve whose conormal derivatives have rank two is bounded by a complete intersection.}
Remove the common fixed divisor $H$ from the linear system $K_C$.
It cannot contain $C$: at the generic point of a curve contained
in a common divisor, every derivative of a section divisible by its
local equation lies in at most one conormal direction. Thus the
residual system in $L(-H)$ still vanishes on $C$ and has no fixed
divisor. Two general residual divisors meet properly and contain
$C$ in their effective intersection cycle. If $s_0=s-[H]$, this
shows that $s_0^2-[C]$ is the class of an effective curve cycle.
The class
\[
 s^2-s_0^2=[H](s+s_0)
\]
is effective as a cohomology class: both factors have effective
divisor representatives, and translating one of them generally
makes their intersection proper without changing its class.
Therefore
\[
 s^2-[C]\text{ and }t^2-[D]
 \quad\text{are effective curve classes.}
 \tag{C44.A07.6}
\]
Intersecting with $s,t$ gives $0<u,v\le1$.

\paragraph{3. The precise Fourier bound for effective curves.}
We record the needed form and its characteristic-independent
justification. For a polarized abelian threefold $(A,s)$ let
\[
 \mathcal F_s(x)=\chi(L)^{-1}\varphi_L^*
       p_{\widehat A,*}(p_A^*x\,\exp(c_1(\mathcal P))),
 \qquad T_s(x)=-\mathcal F_s(x)
\]
on degree-four cohomology; $\mathcal P$ is the normalized Poincar\'e
bundle on $A\times\widehat A$. Changing its sign makes no difference
in this degree. The normalized Fourier formulas are
\[
 \mathcal F_s(s^2)=-2s,\qquad
 \mathcal F_s(s\alpha_0)=\alpha_0\quad(s^2\alpha_0=0).
 \tag{C44.A07.7}
\]
These follow directly from Beauville, \emph{The action of
$\mathrm{SL}_2$ on abelian varieties}, JRMS 25 (2010),
Corollary~5.4, p.~261, with $g=3$, and the normalization in
\S4.1, p.~260. That paper works over an algebraically closed field
of arbitrary characteristic (\S1.1).

For every effective curve cycle $F$ on $A$, the class $T_s([F])$
is a nef rational divisor. Indeed it suffices to normalize an
integral component $\nu:\widetilde F\to A$ and use the restriction
homomorphism $\rho:\widehat A\to J(\widetilde F)$. The Poincar\'e
identity on a curve gives
\[
 -\frac12 p_{\widehat A,*}
       \bigl([F]c_1(\mathcal P)^2\bigr)=\rho^*\theta_{\widetilde F}.
 \tag{C44.A07.8}
\]
Thus $T_s([F])=\chi(L)^{-1}\varphi_L^*\rho^*\theta_{\widetilde F}$,
which is nef. For clarity, (C44.A07.8) uses only the canonical
principal polarization of a Jacobian, not complex Hodge positivity.
In a symplectic basis $e_i,f_i$ of $H^1(\widetilde F)$ with
$\int e_if_j=\delta_{ij}$, the universal Poincar\'e tensor is
$\sum_i(e_i\otimes v_i-f_i\otimes u_i)$ and the theta class is
$\sum_i u_iv_i$. Its square pushed along the curve is
$-2\sum_i u_iv_i$, proving the formula. The identification of these
tensors in every characteristic is the autoduality statement of
Milne, \emph{Jacobian Varieties} (2021 corrected reprint of the
1986 chapter), Theorem~6.6, Lemmas~6.8--6.9 and Summary~6.11,
pp.~21--24. Equivalently, pullback to the square of the curve is
the diagonal minus the two base-point fibres, as in Summary~6.11.

Put $\alpha=L_s^{-1}[C]=us+\alpha_0$, with $s^2\alpha_0=0$,
and $Q_\alpha=\int_A s\alpha_0^2$. The inverse is algebraic by the
known Lefschetz standard conjecture for abelian varieties.
Equations (C44.A07.6)--(C44.A07.8) show that
\[
 T_s([C])=2us-\alpha_0,
 \qquad 0\le T_s([C])\le2s
 \tag{C44.A07.9}
\]
in the nef order. We claim that
\[
 0\le -Q_\alpha\le 2u(1-u)V.
 \tag{C44.A07.10}
\]
Here is an algebraic proof of the bound, including the reality
needed to use (C44.A07.9). For a rational divisor class $E$ put
$f=\varphi_L^{-1}\varphi_E\in\operatorname{End}^0(A)$. This is
self-adjoint for the Rosati involution. The positive definite
Rosati pairing $(x,y)=\operatorname{Tr}(xy^\dagger)$ makes left
multiplication by $f$ self-adjoint: cyclicity of the trace gives
$(fx,y)=(x,fy)$. It follows that the minimal polynomial of $f$
has only real roots. Rosati positivity here is the classical
theorem, for example Conrad, \emph{Abelian varieties: geometry,
parameter spaces, and arithmetic} (2005), Theorem~5.2.4,
pp.~45--46; it is unrelated to 032.

Riemann--Roch and the degree polynomial of an endomorphism give
\[
 \det(T+f\mid V_\ell A)
   =\left(\frac{\int_A(Ts+E)^3}{V}\right)^2.
 \tag{C44.A07.11}
\]
Exact locators are Milne, \emph{Abelian Varieties}, course notes
v.~2.00 (2008), Proposition~10.20 (printed pp.~50--51) and
Theorem~11.1(a),(b) (printed p.~54). In particular the monic cubic
on the right before squaring has three real roots, counted with
multiplicity. Write it as $\prod_{i=1}^3(T+\mu_i)$.
For $E=T_s([C])$, both $E$ and $2s-E$ are nef. The mixed
intersection coefficients of a nef divisor with $s$ are
nonnegative, so their cubics are strictly positive at every
positive real $T$. Applied to $f$ and $2-f$, this shows
$0\le\mu_i\le2$.

Expanding (C44.A07.11) using $E=2us-\alpha_0$ gives
\[
 \sum_i\mu_i=6u,\qquad
 \sum_i\mu_i^2=12u^2-6Q_\alpha/V.
\]
The lower bound $\sum_i\mu_i^2\ge(\sum_i\mu_i)^2/3$ and the
upper bound $\sum_i\mu_i^2\le2\sum_i\mu_i$ prove
(C44.A07.10). The identical argument for
$\beta=L_t^{-1}[D]=vt+\beta_0$ gives
\[
 0\le -Q_\beta\le2v(1-v)W,
 \qquad Q_\beta=\int_{A'}t\beta_0^2.
 \tag{C44.A07.12}
\]

\paragraph{4. Exact self-intersection and ambient subtraction.}
The normal sequence gives
$c(N_{Z/B})=c(N_{Z/X})(1+s+t)^{-2}$. Since
$N_{Z/X}=N_{C/A}\boxplus N_{D/A'}$, its degree-two part gives
\[
 \int_B[Z]^2=6ab-2(ad+bc)+cd
              =(c-2a)(d-2b)+2ab\ \ge2ab.
 \tag{C44.A07.13}
\]
This is the ordinary self-intersection formula (Fulton,
\emph{Intersection Theory}, \S6.3), with (C44.A07.5).

Using $s^2\alpha_0=t^2\beta_0=0$, direct multiplication gives
\[
 \begin{split}
 a_Z:=(L_{s+t}^2)^{-1}[C][D]
 &=\tfrac12\alpha_0\beta_0+v s\alpha_0+u t\beta_0\\
 &\quad+uv\bigl(st-\tfrac12(s^2+t^2)\bigr).
 \end{split}
 \tag{C44.A07.14}
\]
Indeed multiplication by $(s+t)^2$ gives
$(us^2+s\alpha_0)(vt^2+t\beta_0)$ term by term.
Consequently
\[
 G:=\int_Xa_Z[C][D]
       =\tfrac12Q_\alpha Q_\beta+u^2v^2VW.
 \tag{C44.A07.15}
\]
The self-intersection identity $\kappa_*\kappa^*=L_h^2$ gives
$\kappa_*w=0$ and
$\int_Bw^2=\int_B[Z]^2-G$. Equations (C44.A07.10) and
(C44.A07.12), together with $ab=uvVW$, imply
\[
 \frac{G}{ab}\le uv+2(1-u)(1-v).
\]
Subtracting this from (C44.A07.13) proves (C44.A07.2).
Finally $2u+2v-3uv\ge u+v-uv>0$ when $0<u,v\le1$.
\hfill$\square$

\paragraph{Scope and programme location.}
The theorem is valid for supersingular factors and for curves
whose embedded support is not known to lift. It does not claim
that $w$ is nonliftable. It establishes the axis height sign on
one explicit nonzero line at the original polarization, not on
the span of all such lines or on all of $A_B$. In cote~44 it
addresses the remaining axis summand in Th.~5 (scan~004), with
the primitive qualification of the letter on scan~008. In
cote~16 terms it supplies a particular positive primitive
algebraic class toward $I(B)$, without settling that condition
or its product version. It strengthens the calculation of a
candidate class; QH9-axis itself remains open in this work.

% END work/axis/axis_product_curves.tex

% BEGIN work/axis/axis_product_grid.tex
% C44_POS /root; GPT-6 Astra Pro; effort not exposed.
\section{A semipositive span at one fixed product polarization}

\begin{theorem}[C44-A09: a grid of product curves]
Let $A,A'$ be abelian threefolds over an algebraically closed field
of arbitrary characteristic, with ample classes $s=c_1(L)$ and
$t=c_1(M)$. Let $B\subset A\times A'$ be a smooth complete
intersection of two members of $|L\boxtimes M|$.
Let $C_1,\ldots,C_e\subset A$ and $D_1,\ldots,D_f\subset A'$
be two nonempty families of smooth connected curves, pairwise
disjoint within each family. Suppose that $B$ contains every
$Z_{ij}=C_i\times D_j$. Put
\[
 \begin{gathered}
 V=\int_As^3,\qquad W=\int_{A'}t^3,\qquad
 a_i=\int_{C_i}s,\quad b_j=\int_{D_j}t,\\
 u_i=a_i/V,\qquad v_j=b_j/W,\qquad
 U_0=\sum_i u_i,\quad V_0=\sum_jv_j,\\
 w_{ij}=[Z_{ij}]_B-\kappa^*(L_{s+t}^2)^{-1}[Z_{ij}]_{A\times A'}.
 \end{gathered}
 \tag{C44.A09.1}
\]
Then $0<U_0,V_0\le1$ and $w_{ij}\in A_B$.
The intersection form $\int_Bww'$ is positive semidefinite on
the entire rational span of the $w_{ij}$. If $e=1$ or $f=1$, its
Gram matrix in these generators is positive definite, so the
generators are independent. All these statements are unconditional.

More precisely, for any real coefficients $x_{ij}$, define
\[
 R_i=\sum_jv_jx_{ij},\qquad
 S_j=\sum_i u_ix_{ij},\qquad
 z=\sum_{i,j}u_iv_jx_{ij}.
\]
The exact axis form satisfies the explicit lower bound
\[
 \begin{split}
 \int_B\left(\sum_{i,j}x_{ij}w_{ij}\right)^2
 \ge VW\left(2\sum_i u_iR_i^2+2\sum_jv_jS_j^2-3z^2\right)\ge0.
 \end{split}
 \tag{C44.A09.2}
\]
A vector of zero norm must therefore have every weighted row sum
$R_i$ and every weighted column sum $S_j$ equal to zero. No assertion
that every such vector has zero norm, or is a cohomological relation,
is made.
\end{theorem}

\begin{proof}
Write $C=\coprod_iC_i$, $D=\coprod_jD_j$ and
\[
 K_C=H^0(A,I_CL),\qquad K_D=H^0(A',I_DM).
\]
The restriction kernel on $C\times D$ is
$K_C\otimes H^0(M)+H^0(L)\otimes K_D$, by the same tensor-kernel
identity as (C44.A07.3). At every pair $C_i\times D_j$, the
normal surjection associated with the smooth axis factors
through the conormal derivatives of these two global spaces.
The rank-zero and rank-one contradictions in part~1 of C44-A07
apply to these spaces themselves: their proof used only this
factorization, saturation on a smooth curve, and the surjectivity
of the normal map. It did not require that all sections vanishing
on the individual component be available. Hence $K_C$ has
conormal rank two generically on each $C_i$, and $K_D$ has rank
two generically on each $D_j$.

Removing the common fixed divisor of $K_C$, no component $C_i$
can be contained in that divisor, since that would leave at most
one conormal direction. Two general residual divisors meet
properly and contain every $C_i$. Restoring the fixed divisor
by the general-translation intersection argument of C44-A07
proves
\[
          s^2-\sum_i[C_i]\ \text{and}\
          t^2-\sum_j[D_j]\quad\text{are effective curve classes}.
 \tag{C44.A09.3}
\]
In particular $0<U_0,V_0\le1$.
Writing $c_i=2g(C_i)-2$ and $d_j=2g(D_j)-2$, the conormal
rank-two determinants also give $c_i\le2a_i$, $d_j\le2b_j$.
The normal-bundle formula consequently yields
\[
 T_{ij}:=\int_B[Z_{ij}]^2
       =(c_i-2a_i)(d_j-2b_j)+2a_ib_j\ge2a_ib_j.
 \tag{C44.A09.4}
\]
Distinct $Z_{ij}$ are disjoint.

\paragraph{A simultaneous Fourier bound.}
Decompose
\[
 L_s^{-1}[C_i]=u_is+\alpha_i,\quad s^2\alpha_i=0,
 \qquad
 L_t^{-1}[D_j]=v_jt+\beta_j,\quad t^2\beta_j=0.
\]
Let
\[
 Q^C_{ik}=\int_As\alpha_i\alpha_k,\qquad
 Q^D_{jl}=\int_{A'}t\beta_j\beta_l,
 \qquad D_u=\operatorname{diag}(u_i),\quad
 D_v=\operatorname{diag}(v_j).
\]
We claim the inequalities of real symmetric matrices
\[
 0\le -Q^C\le2V(D_u-uu^t),\qquad
 0\le -Q^D\le2W(D_v-vv^t).
 \tag{C44.A09.5}
\]
Here $u=(u_i)$ and $v=(v_j)$ are column vectors. Matrix
inequalities mean nonnegative quadratic forms, not entrywise
inequalities.

Use the normalized Fourier map $T_s$ of C44-A07. Each
$T_s([C_i])=2u_is-\alpha_i$ is nef, and (C44.A09.3) gives
$\sum_iT_s([C_i])\le2s$ in the nef order. Let $E_i$ be the
Rosati-self-adjoint endomorphism corresponding to this divisor.
As proved in C44-A07 by Rosati positivity and Riemann--Roch,
nef divisors correspond here to positive self-adjoint operators
in the faithful left-regular representation of
$\operatorname{End}^0(A)\otimes\mathbf R$. To make the
representation comparison explicit, the Tate representation is
faithful, so every root of the minimal polynomial of a rational
self-adjoint endomorphism occurs among the roots of its degree
polynomial. The Riemann--Roch and nef arguments of C44-A07 make
all these roots nonnegative. This is exactly positivity in the
left-regular Rosati representation, and extends to real
self-adjoint elements by continuity. Thus
\[
                        E_i\ge0,\qquad \sum_iE_i\le2I.
\]
For arbitrary real $x_i$, elementary operator Cauchy--Schwarz gives
\[
             \left(\sum_i x_iE_i\right)^2
                         \le2\sum_i x_i^2E_i.
 \tag{C44.A09.6}
\]
For example, let $R$ be the block-row operator with entries
$E_i^{1/2}$ and let $D$ act by the scalar $x_i$ on its $i$-th
summand. Then $RR^*=\sum_iE_i\le2I$, so $R^*R\le2I$, and
\[
 (RDR^*)^2=RD(R^*R)DR^*\le2RD^2R^*.
\]
No commutativity among the $E_i$ is needed.

Let $\tau=\tfrac12\operatorname{Tr}_{V_\ell A}$ on the rational
endomorphism algebra, extended to the real algebra. It is
positive on positive self-adjoint elements: such an element
has a self-adjoint square root in the real finite-dimensional
algebra, and Rosati trace positivity applies to its square.
Also $\tau(I)=3$ and $\tau(E_i)=6u_i$. For
$E(x)=\sum_i x_iE_i$ and $u(x)=\sum_i x_iu_i$, the degree
polynomial calculation (C44.A07.11) gives
\[
 \tau(E(x)^2)=12u(x)^2-\frac6Vx^tQ^Cx.
 \tag{C44.A09.7}
\]
It is a polynomial identity, so is valid for real as well as
rational $x_i$. Applying $\tau$ to (C44.A09.6) proves the
upper bound in (C44.A09.5). The lower bound follows from
$\tau(E(x)^2)\ge\tau(E(x))^2/\tau(I)$, which is
Cauchy--Schwarz for the positive Rosati trace pairing. The
second matrix bound is identical.

\paragraph{The complete Gram matrix and its lower bound.}
The exact inverse and ambient subtraction in C44-A07, applied
bilinearly to each pair, give
\[
 G_{ij,kl}=\tfrac12 Q^C_{ik}Q^D_{jl}
                              +VWu_iu_kv_jv_l,
 \qquad
 Q_B=\operatorname{diag}(T_{ij})-G.
 \tag{C44.A09.8}
\]
The first identity also follows by multiplying the displayed
inverse (C44.A07.14) for $C_i\times D_j$ against
$[C_k][D_l]$. The remaining terms have too high a degree on one
threefold factor and vanish. The second identity uses disjointness
of distinct centres and the same projection formula as C44-A07.

The tensor product of positive semidefinite forms preserves
inequalities. Applying (C44.A09.5) therefore shows
\[
 \begin{split}
 G&\le2VW(D_u-uu^t)\otimes(D_v-vv^t)
                           +VWuu^t\otimes vv^t,\\
 Q_B&\ge VW\left(2D_u\otimes vv^t+
                  2uu^t\otimes D_v-3uu^t\otimes vv^t\right),
 \end{split}
 \tag{C44.A09.9}
\]
where (C44.A09.4) was used in the second line. Evaluation on
$x=(x_{ij})$ is the first inequality of (C44.A09.2).
Its right-hand quadratic form is nonnegative because it equals
\[
 2\sum_i u_i(R_i-z/U_0)^2+
 2\sum_j v_j(S_j-z/V_0)^2+
       (2/U_0+2/V_0-3)z^2.
 \tag{C44.A09.10}
\]
The last coefficient is at least one. All weights are strictly
positive. Thus equality in this lower form forces $z=0$ and
all $R_i,S_j$ zero. If either family has a single component,
the weighted row and column sums determine every $x_{ij}$;
hence the lower form is positive definite. This proves all
claims, including independence in that case.
\end{proof}

\paragraph{Programme consequence and limitation.}
Unlike an assertion that several individual classes have
positive squares, C44-A09 controls every linear combination in
the displayed span at the original polarization. For a
Lefschetz axis it is exactly the H9-T7 sign required for this
part of the height in cote~44 Th.~5, scan~004, with the
primitive qualification on scan~008. It is an unconditional
special-axis result in every characteristic, including
supersingular ambient factors. It neither asserts that all of
$A_B$ is generated by these product curves nor settles the
universal $I(B)$ of cote~16. The proved criterion requires the
entire product grid, not an arbitrary collection of unrelated
product curves. Its intersection and Fourier inputs are the
fully specified classical inputs of C44-A07; no 032 consequence
is used.

% END work/axis/axis_product_grid.tex

% BEGIN work/axis/axis_grid_strict.tex
% C44_POS /root/axis_independent; GPT-6 Astra Pro; effort not exposed.
\subsection*{C44-A09a. A generating curve makes the entire grid strictly positive}

\textbf{Theorem.}
In the setting and notation of C44-A09, suppose at least one
$C_i$ generates $A$ as an abelian variety, or at least one
$D_j$ generates $A'$. Here a curve generates if no translate of
a proper abelian subvariety contains it. Then the complete
Gram matrix of the classes $w_{ij}$ is positive definite.
In particular all the $w_{ij}$ are independent. This holds in
every characteristic without 032. It applies in particular
whenever either threefold factor is geometrically simple.

\textbf{Proof.}
Suppose $C_{i_0}$ generates $A$. The induced homomorphism
$J(C_{i_0})\to A$ is surjective, by the universal property of
the Jacobian (Milne, \emph{Jacobian Varieties}, Proposition~6.1,
printed p.~19 in the 2021 corrected reprint). Its dual
$\rho:\widehat A\to J(C_{i_0})$ has finite kernel, so is finite
onto its image. The explicit formula (C44.A07.8) shows that
$T_s([C_{i_0}])=\chi(L)^{-1}\varphi_L^*\rho^*\theta_{C_{i_0}}$
is ample as a rational divisor class. Consequently the
corresponding Rosati-self-adjoint operator $E_{i_0}$ of C44-A09
is strictly positive and invertible. For example, the nef
argument there makes its spectrum nonnegative; its ample top
self-intersection is nonzero, so the Riemann--Roch degree
polynomial excludes eigenvalue zero. Tate faithfulness makes
this conclusion valid in the left-regular representation as
well as in the Tate representation.

Keep the notation of C44-A09 and define
\[
 A_C=-Q^C,\quad H_C=2V(D_u-uu^t),\quad M_C=H_C-A_C,
 \qquad
 A_D=-Q^D,\quad H_D=2W(D_v-vv^t),\quad M_D=H_D-A_D.
\]
All six matrices are positive semidefinite by (C44.A09.5).
We first prove that $M_C$ is positive definite on
$u^\perp=\{x:\sum_i u_ix_i=0\}$.

Use the faithful real left-regular representation with its
positive Rosati inner product. For real $x=(x_i)$ let
$R=(E_1^{1/2},\ldots,E_e^{1/2})$, let $D$ be the block-diagonal
operator with scalar blocks $x_i$, and put
$F=RDR^*=\sum_i x_iE_i$. The exact trace calculation in
C44-A09 gives
\[
 x^tM_Cx=\frac V6\tau\bigl(2\sum_i x_i^2E_i-F^2\bigr),
 \qquad \tau=\tfrac12\operatorname{Tr}_{V_\ell A}.
 \tag{C44.A09a.1}
\]
The operator inside the trace is
\[
 S=RD(2I-R^*R)DR^*\ge0.
\]
The functional $\tau$ is faithful on positive elements: if a
nonzero positive element is written as the square of its
self-adjoint square root in the finite-dimensional real
algebra, this follows from positive definiteness of the
Rosati trace pairing. Therefore equality in (C44.A09a.1)
implies $S=0$. It follows that
$(2I-R^*R)^{1/2}DR^*=0$, and hence
\[
                    2DR^*=R^*RDR^*=R^*F.
 \tag{C44.A09a.2}
\]
In the $i$-th component this says
$2x_iE_i^{1/2}=E_i^{1/2}F$.
For $i=i_0$, invertibility gives $F=2x_{i_0}I$.
Substituting into the other components gives
$(x_i-x_{i_0})E_i^{1/2}=0$ for every $i$.
Each $E_i$ is nonzero because $\tau(E_i)=6u_i>0$.
All $x_i$ are therefore equal. Such a vector in $u^\perp$
must vanish, since $\sum_i u_i>0$. This proves the asserted
strict positivity of $M_C$ on $u^\perp$.

Let $Q_B$ be the complete grid Gram matrix and let
\[
 K=VW\bigl(2D_u\otimes vv^t+2uu^t\otimes D_v
                         -3uu^t\otimes vv^t\bigr)
\]
be its lower bound in (C44.A09.9). The same calculation, retaining
the slack in the tensor inequality, gives
\[
 \begin{split}
 Q_B&\ge K+\tfrac12(H_C\otimes H_D-A_C\otimes A_D)\\
    &=K+\tfrac12(M_C\otimes H_D+A_C\otimes M_D)\\
    &\ge K+\tfrac12M_C\otimes H_D.
 \end{split}
 \tag{C44.A09a.3}
\]
By (C44.A09.10), $K$ is positive semidefinite and its kernel
consists precisely of arrays with zero weighted row and
column sums. As a tensor subspace this is
$u^\perp\otimes v^\perp$.
On $v^\perp$ one has
$y^tH_Dy=2W\sum_jv_jy_j^2>0$ for $y\ne0$.
Thus $M_C\otimes H_D$ is positive definite on the kernel
of $K$. Inequality (C44.A09a.3) proves that $Q_B$ is positive
definite everywhere. If a $D_j$ generates $A'$, interchange
the factors. Finally every positive-dimensional curve in a
geometrically simple abelian variety generates it up to
translation, proving the last assertion.\hfill$\square$

\paragraph{Programme scope.}
This proves strict positivity for a full family of special-axis
classes at one fixed polarization, strengthening C44-A09 and
the single-class result C44-A07. It controls the corresponding
entire summand of the H9-T7 height for cote~44 Th.~5, scan~004.
It makes no claim that the grid exhausts $A_B$, and supplies
neither a universal correspondence injection nor an individual
nonliftable class on the axis.

% END work/axis/axis_grid_strict.tex

% BEGIN work/axis/axis_supersingular_nonliftable.tex
% C44_POS /root/axis_independent; GPT-6 Astra Pro; effort not exposed.
\section{An explicit supersingular nonliftable class outside the axis kernel}
\subsection*{C44-A10. Individual class nonliftability}

\textbf{Theorem.}
For $E/\overline{\mathbf F}_3$ given by $y^2=x^3-x$, there is an
explicit integral divisor-product cycle $z\in\operatorname{CH}^2(E^4)$
whose nonzero cohomology class cannot specialize from an algebraic
cycle in any characteristic-zero abelian-scheme lift of $E^4$.
If $a$ is its pullback to $X=E^6$ and $\kappa:B\hookrightarrow X$
is a smooth complete intersection of two ample divisors of class $h$,
then $\kappa^*a$ cannot lift even under an arbitrary abstract smooth
projective lift of $B$. Nevertheless,
\[
 \kappa_*\kappa^*a=h^2a\ne0,\qquad
 (1-\kappa^*(L_h^2)^{-1}\kappa_*)\kappa^*a=0.                 \tag{A10.1}
\]
Thus this is individual class nonliftability, but gives no
nonliftable class in $A_B$.

\paragraph{Explicit divisors and cycle.}
Choose $\eta\in\mathbf F_9$ with $\eta^2=-1$ and define
$\iota(x,y)=(-x,\eta y)$, $\pi(x,y)=(x^3,y^3)$.
Then $\iota^2=-1$, $\pi\iota=-\iota\pi$, and $\pi^2=-3$:
there are four points over $\mathbf F_3$, so the Frobenius
polynomial is $T^2+3$, and Cayley--Hamilton plus Tate-module
faithfulness gives the identity. Thus $[3]=-\pi^2$ is purely
inseparable and $E$ is supersingular.
The degree/characteristic-polynomial statements are in Milne,
\emph{Abelian Varieties}, v.~2.00, I Proposition 10.20,
printed pp.~50--52; the point-count identity is in II, \S1.

Choose a basis of $H^1(E,\mathbf F_5)$ with pullback matrices
\[
 I=\iota^*=\begin{pmatrix}2&0\\0&3\end{pmatrix},\qquad
 P=\pi^*=\begin{pmatrix}0&1\\2&0\end{pmatrix}.                 \tag{A10.2}
\]
The invertible $\pi$ interchanges the two distinct $\iota$
eigenspaces, and rescaling using $\pi^2=2$ gives this form.
Do not assume Weil normalization: write $[0]=\lambda e_0e_1$
with $\lambda\in\mathbf F_5^\times$ and use the same basis on all
four factors. Put $D_i=\operatorname{pr}_i^*[0]$ and, for $i<j$,
\[
 R_{ij,u}=(\operatorname{pr}_i+u\operatorname{pr}_j)^*[0]
                    -D_i-(\deg u)D_j,\quad
 u=1,\iota,\pi,\pi\circ\iota,\quad \deg u=1,1,3,3.            \tag{A10.3}
\]
These are integral divisors defined over $\mathbf F_9$.
For $a=2i-2,b=a+1,c=2j-2,d=c+1$, their normalized modulo-five
cross coordinates in the order $(ac,ad,bc,bd)$ are respectively
\[
 (0,1,4,0),\ (0,3,3,0),\ (1,0,0,3),\ (2,0,0,4).
\]
Indeed they are the cross terms in
\[
 (e_a+M_{11}e_c+M_{21}e_d)(e_b+M_{12}e_c+M_{22}e_d),
\]
where $M=1,I,P,IP$. The last matrix is $IP$ because pullback
reverses $\pi\circ\iota$.
Define integral divisors, for each pair of factors, by
\[
 \begin{array}{ll}
 F_{ac}=3R_{ij,\pi}+4R_{ij,\pi\circ\iota},&
 F_{ad}=3R_{ij,1}+R_{ij,\iota},\\
 F_{bc}=2R_{ij,1}+R_{ij,\iota},&
 F_{bd}=R_{ij,\pi}+2R_{ij,\pi\circ\iota},
 \end{array}
 \qquad F_{2i-2,2i-1}=D_i.                                  \tag{A10.4}
\]
Direct multiplication of the four cross columns shows
$\operatorname{cl}_5(F_{ab})=\lambda e_ae_b$.
Equivalently these are the representatives in $\{0,1,2,3,4\}$
of the inverse of the $28$-by-$28$ divisor coefficient matrix,
ordered by the four $D_i$, lexicographic pairs $i<j$, and the
four displayed endomorphisms.

Let $c_{abcd}$, in increasing lexicographic four-tuple order
on $\{0,\ldots,7\}$, be the following 70 integers read rowwise:
\[
\begin{array}{rrrrrrrrrr}
4&4&1&3&4&0&0&3&0&1\\
3&3&3&1&2&2&3&4&2&4\\
3&2&2&1&3&4&2&4&3&4\\
1&4&2&4&1&3&1&4&3&2\\
0&0&3&2&3&1&1&4&2&3\\
2&1&4&4&1&3&4&1&3&0\\
4&2&2&3&1&3&4&4&2&2
\end{array}
\]
Set $z=\sum_{a<b<c<d}c_{abcd}F_{ab}F_{cd}$.
Its modulo-five class is $\lambda^2w$, where
$w=\sum c_{abcd}e_ae_be_ce_d$. This completely specifies the
actual integral cycle, without requiring any conjectural classes.

\paragraph{Exact certificate for the actual rational tensor.}
For $w\in\bigwedge^4V$ let
$\mathcal D_w:\operatorname{End}(V)\to\bigwedge^4V$ send
$T$ to its derivative exterior action $T\cdot w$.
Its kernel is the infinitesimal stabilizer.
Order columns by $E_{ab}$, $a,b=0,\ldots,7$ lexicographically,
and rows by increasing lexicographic four-tuples $J$.
Its entries are
\[
 (\mathcal D_w)_{J,(a,b)}=
 \begin{cases}
 c_J,&a=b,\ a\in J,\\
 (-1)^{\operatorname{pos}_J(a)+\operatorname{pos}_K(b)}c_K,
 &a\ne b,\ a\in J,\ b\notin J,\ K=(J\setminus\{a\})\cup\{b\},\\
 0,&\text{otherwise}.
 \end{cases}                                                \tag{A10.5}
\]
Positions are zero-based, and both tuples are sorted.
The first 64 rows, with entries represented in $\{0,1,2,3,4\}$,
have exact integer determinant
\[
 26723729698312903479574054159209851065213568
                       \equiv3\pmod5.                       \tag{A10.6}
\]
The supplied \texttt{supersingular\_class\_certificate.py}
constructs the divisor matrix and tensor; its
\texttt{supersingular\_class\_certificate.json} supplies both
the divisor inverse and the full inverse of this minor modulo
five. Both inverse products were checked.
A separate target-row implementation of (A10.5) and exact Bareiss
elimination reproduced (A10.6). The root agent independently
rebuilt the matrix by source-column exterior replacement and
reproduced the same exact determinant. Random search selected
the fixed displayed vector; verification of it is exact.

Milne, loc.\ cit., I Theorem 12.1, printed pp.~55--56, gives
$H^*(E^4,\mathbf Z_5)=\bigwedge^*H^1(E^4,\mathbf Z_5)$,
torsion free. Lift the chosen modulo-five basis arbitrarily to a
$\mathbf Z_5$ basis. The derivative matrix of the actual integral
class of $z$ reduces to $\mathcal D_{\lambda^2w}$.
Its indicated minor is congruent to
$(\lambda^2)^{64}\cdot3=\lambda^{128}\cdot3=3$ modulo five.
Consequently the actual $\mathbf Q_5$ tensor has zero infinitesimal
stabilizer. Its integral coordinates need not equal the selected
residue representatives: higher $5$-adic terms cannot change
this unit-minor conclusion.

\paragraph{Characteristic-zero obstruction.}
If the class lifts algebraically to a characteristic-zero abelian
fourfold, spread the generic fibre and cycle to a finitely generated
field and embed it into $\mathbf C$. Smooth proper base change and
Betti--\'etale comparison preserve $H^1$, exterior products and
$\mathcal D$. Thus its rational Betti tensor also has zero
infinitesimal stabilizer after scalar extension.
But an algebraic degree-four class is $(2,2)$.
The real Hodge complex structure $J$ on $H^1$ acts by $+i$ on
$H^{1,0}$ and $-i$ on $H^{0,1}$, satisfies $J^2=-1$, and
annihilates a $(2,2)$ exterior tensor. It is a nonzero stabilizer,
contradiction. Although $J$ need not be rational, the derivative
matrix of a rational Betti tensor is rational: a nonzero kernel
over $\mathbf R$ implies rank deficiency over $\mathbf Q$ and
hence over $\mathbf Q_5$.
For the Hodge statements see Deligne, \emph{Hodge cycles on abelian
varieties}, notes by Milne, revised 1 October 2018, Introduction
printed p.~2 and \S3, Hodge structures, printed p.~23.
No Hodge conjecture is used.

\paragraph{Pullback and arbitrary abstract axis lifts.}
The support $S(t)$ of an exterior tensor $t$ is its least supporting
subspace, equivalently the span of its contractions by three
covectors. It is intrinsic, equivariant and compatible with field
extension. The certified tensor on an eight-dimensional space has
full support, since a nonzero endomorphism vanishing on a proper
support would annihilate it.
Hence $a$ on $E^6$ has intrinsic support of dimension eight and
zero infinitesimal stabilizer on that support.
Any $(2,2)$ lift would have Hodge-circle-invariant support;
$J$ restricted to it still squares to $-1$ and annihilates the
tensor, contradiction. No preservation of the selected $E^4$
factor by the lift is assumed.

Weak Lefschetz gives $H^1(E^6)\simeq H^1(B)$ and
$H^4(E^6)\hookrightarrow H^4(B)$. Therefore the intrinsic cup map
\[
                     \bigwedge^4H^1(B)\longrightarrow H^4(B)
\]
is injective, and its rank and image transport through any abstract
smooth proper lift. If $\kappa^*a$ lifted to an algebraic class
$\gamma$, its $\mathbf Q_5$ realization would remain in this image.
After spreading to a complex fibre, rational Betti $\gamma$ lies
in the rational cup-image too: membership can be tested after
scalar extension to $\mathbf Q_5$.
Its unique rational preimage is $(2,2)$, because cup product is
an injective morphism of Hodge structures. The preimage retains
its eight-dimensional support with zero infinitesimal stabilizer.
The same Hodge-circle contradiction excludes the abstract lift.
Finally hard Lefschetz gives
$L_h^2:H^4(E^6)\simeq H^8(E^6)$, and the nonzero $a$ satisfies
(A10.1) by self-intersection. \hfill$\square$

\paragraph{Transfer obstructions.}
For an axis of dimension $2r$ in an abelian $(2r+2)$-fold,
every $w\in\ker\kappa_*\subset H^{2r}(B)$ satisfies
\[
                   w\,\kappa^*\alpha=0\quad
                         (\alpha\in H^{>0}(X)).
\]
Its Gysin image is zero by projection formula; Gysin above the
middle is injective by dual weak Lefschetz. Products beyond top
degree vanish automatically. Thus the displayed $H^1$ tensor
cannot be recovered from an axis-kernel class by ambient cup
products.
Likewise, if $i:S\hookrightarrow B$ has codimension $c<r$ and
$\eta=i^*\kappa^*\theta$ with $\theta\in H^{2r-2c}(X)$, then
$i_*\eta=\kappa^*\theta\,[S]_B$ is ambient, since weak Lefschetz
makes $[S]_B\in H^{2c}(B)$ ambient. Its variable projection is zero.

\paragraph{Sources, scope and genuine checks.}
Primary URLs:
\texttt{https://jmilne.org/math/CourseNotes/AV.pdf};
\texttt{https://www.jmilne.org/math/Documents/Deligne82.pdf}.
Smooth proper base change is SGA4, Expos\'e XVI, Th\'eor\`eme 2.2,
together with Betti--\'etale comparison in characteristic zero.
C44-A05 lifts geometry; this theorem shows a single algebraic
middle class need not lift. It lies wholly in the ambient summand,
where the sign is covered by the conditional 032 input.
It gives neither a nonliftable $A_B$ class nor a wrong-sign class;
QH9-axis remains open.

The author checked the whole proof on 22 November 2026
(Europe/Berlin), including the actual integral tensor and all
lifting quantifiers. The root independently read all three
complete proof messages and passed the whole proof at 00:49:30 UTC
that date, also reproducing the exact determinant independently.
Both agents are GPT-6 Astra Pro; reasoning effort is not exposed.
The closed divisor formulas (A10.4) were subsequently checked by
hand and exact modulo-five JavaScript. None of these events is
a claimed verification of QH9-axis.

% END work/axis/axis_supersingular_nonliftable.tex

% BEGIN work/axis/axis_product_ci_matching.tex
% C44_POS /root/axis_independent; GPT-6 Astra Pro; effort not exposed.
\section{A sharp fixed-polarization bound for products of complete-intersection curves}
\subsection*{C44-A11. Two is the maximum for smooth complete intersections}

\textbf{Theorem.}
Let $S,T$ be abelian threefolds, $L,M$ ample line bundles of
classes $s,t$, and $B\subset S\times T$ a smooth complete
intersection of two sections of $L\boxtimes M$. Put
$V=\int_Ss^3$, $W=\int_Tt^3$.
Suppose it contains $Z_i=C_i\times D_i$, $1\le i\le N$,
where the $C_i$ are pairwise disjoint smooth scheme-theoretic
complete intersections of two divisors in $|L|$, and the $D_i$
are pairwise disjoint smooth scheme-theoretic complete
intersections of two divisors in $|M|$. Then $N\le2$.
Set
\[
 a=st-\tfrac12(s^2+t^2),\qquad w_i=[Z_i]_B-\kappa^*a.           \tag{A11.1}
\]
Then $w_i\in A_B$ and
\[
                         (w_iw_j)_B=VW(2\delta_{ij}-1).       \tag{A11.2}
\]
At $N=2$, $w_1+w_2=0$ in rational cohomology, and the actual
span is a positive line. For a Lefschetz pencil in the usual
one-node-per-singular-member sense the stronger bound $N\le1$
holds (C44-A11b below).

\paragraph{1. Section spaces and finite maps.}
Put $A=H^0(S,L)$, $F=H^0(T,M)$ and
$U_i=H^0(\mathcal I_{C_i}L)$, $V_i=H^0(\mathcal I_{D_i}M)$.
Each has dimension two by the twisted Koszul sequence
\[
 0\longrightarrow L^{-1}\longrightarrow\mathcal O_S^2
                     \longrightarrow\mathcal I_{C_i}L\longrightarrow0.
\]
Here $H^0(L^{-1})=H^1(L^{-1})=0$, using Serre duality and
ample-line-bundle cohomology vanishing on abelian varieties;
the same argument applies to $M$.
For $i\ne j$, $U_i\cap U_j=0$. Otherwise the two curves'
combined defining equations would use at most three sections
of $L$. Such sections have a common zero: a nowhere-zero
section of $L^{\oplus q}$, $q\le3$, would imply $s^q=0$,
contradicting $\int_Ss^3>0$. A common zero belongs to
$C_i\cap C_j$, impossible. Similarly $V_i\cap V_j=0$.

Let $f_0,f_1\in A\otimes F$ define $B$. Vanishing on the products
gives $f_\nu\in U_i\otimes F+A\otimes V_i$, exactly the kernel
of the tensor product of the two restriction maps.
If $N\ge2$, set $A_0=U_1\oplus U_2$, $F_0=V_1\oplus V_2$.
Direct-sum complements give
\[
 (U_1\otimes F+A\otimes V_1)\cap(U_2\otimes F+A\otimes V_2)
                         =U_1\otimes V_2+U_2\otimes V_1.     \tag{A11.3}
\]
Thus both equations lie in the four-by-four tensor space
$A_0\otimes F_0$.
The four sections in $A_0$ have common zero scheme
$C_1\cap C_2=\varnothing$, and give
$\phi:S\to\mathbf P(A_0^\vee)\simeq\mathbf P^3$ with
$\phi^*\mathcal O(1)=L$. Ampleness excludes positive-dimensional
fibres, so $\phi$ is finite and surjective, of degree $V$.
Likewise $\psi:T\to\mathbf P^3$ is finite surjective of degree
$W$. Both are flat by the Cohen--Macaulay/regular-local
criterion for finite morphisms between smooth equidimensional
varieties.

\paragraph{2. A rank-three Chern obstruction.}
Every further $U_i$ lies in $A_0$, and every $V_i$ lies in $F_0$.
Suppose $\dim(U_i\cap A_0)\le1$. Restricting $f_\nu$ to
$S\times D_i$ places its terms in
$(U_i\cap A_0)\otimes H^0(D_i,M|_{D_i})$.
The derivative in the $S$ direction along $Z_i$ therefore uses
at most one constant conormal generator. Terms whose coefficients
vanish on $D_i$ have zero derivative in that direction.
Since $N_{C_i/S}\simeq(L|_{C_i})^{\oplus2}$, the full normal
surjection required by smoothness of $B$ factors through
\[
 (L|_{C_i})\boxtimes\mathcal O_{D_i}\ \oplus\
 \mathcal O_{C_i}\boxtimes(M|_{D_i})^{\oplus2}
                 \longrightarrow (L\boxtimes M)|_{Z_i}^{\oplus2}.
                                                               \tag{A11.4}
\]
If the first block is zero, add an unused first summand.
This rank-three surjection has a line-bundle kernel.
For $x=c_1(L|_{C_i})$, $y=c_1(M|_{D_i})$, one has $x^2=y^2=0$,
and the virtual second Chern class of the kernel is
\[
                    \left[\frac{(1+x)(1+y)^2}{(1+x+y)^2}\right]_2
                           =2xy.
\]
Its degree is $2VW>0$, contradicting vanishing of a line bundle's
second Chern class. Thus $U_i\subset A_0$; interchanging factors
gives $V_i\subset F_0$.
Each $C_i$ is consequently the scheme-theoretic inverse image
of a line $\ell_i\subset\mathbf P^3$, and each $D_i$ is the
inverse image of a line $m_i$. Surjectivity shows that each
family of lines is pairwise disjoint.

\paragraph{3. Classification of the smooth bilinear pencil.}
The equations define $B_0\subset\mathbf P^3\times\mathbf P^3$
of bidegrees $(1,1),(1,1)$ with
$B=(\phi\times\psi)^{-1}(B_0)$ scheme-theoretically.
The codimension is two, since its finite pullback has codimension
two. Moreover $B_0$ is smooth: at a singular point the two
gradients are dependent, and remain so after composition with
any differential. Every point has a preimage, which would then
be singular in $B$. Each $\ell_i\times m_i$ lies in $B_0$.

Write its equations as $x^{\mathsf t}M_0y=x^{\mathsf t}M_1y=0$.
Smoothness forces the determinant of this $4$-by-$4$ pencil to
be nonzero with four distinct projective roots, in every
characteristic. For a pencil member of rank at most two, its
left and right kernels have dimension at least two. Choose
nonzero $x,y$ in them with $x^{\mathsf t}M'y=0$ for an
independent member $M'$. This gives a point of $B_0$ with
dependent gradients. For a rank-three member, the two kernel
lines yield a singular point exactly when
$x^{\mathsf t}M'y=0$. By the adjugate formula this is precisely
vanishing of the determinant's derivative at that root.
If the determinant were identically zero, use a rank-three
member and the same derivative argument, or the lower-rank
case. Thus neither an identically zero determinant nor a
multiple root is possible.

After choosing an invertible member and changing coordinates,
\[
 M_0=1,\qquad M_1=\operatorname{diag}(\lambda_1,\ldots,\lambda_4),
                      \quad\lambda_i\ne\lambda_j\ (i\ne j).   \tag{A11.5}
\]
A product $\mathbf P(U)\times\mathbf P(V)$ of lines lies in
$B_0$ exactly when $V=U^\perp$ for $M_0$ and $U$ is invariant
under $M_1^{\mathsf t}=M_1$. Distinct eigenvalues make every
invariant two-plane a span of two eigenvectors. Hence all such
products are
\[
 Z_I^0=\mathbf P(\langle e_i:i\in I\rangle)
          \times\mathbf P(\langle e_j:j\notin I\rangle),
                        \qquad |I|=2.                        \tag{A11.6}
\]
Two distinct products are disjoint exactly when their two-subsets
are complementary. Otherwise their subsets and their complements
each share one index, producing an intersection point.
A pairwise-disjoint collection thus has at most two members:
$N\le2$.

\paragraph{4. Gram matrix and actual boundary relation.}
In the Chow ring of $S\times T$, $[Z_i]=s^2t^2$ and
$(s+t)^2(st-\tfrac12(s^2+t^2))=s^2t^2$ because $s^4=t^4=0$.
Thus $w_i$ lies in the Gysin kernel. The normal sequence gives
\[
 [Z_i]_B^2=\int_{Z_i}
       \left[\frac{(1+x)^2(1+y)^2}{(1+x+y)^2}\right]_2=2VW.
\]
Distinct products are disjoint and the ambient projection term
is $\int_{S\times T}a[Z_i]=VW$, proving (A11.2).

At $N=2$, zero norm alone cannot establish a zero class.
On $B_0$ put $a_0=xy-\tfrac12(x^2+y^2)$ and
$w_I^0=[Z_I^0]_{B_0}-a_0|_{B_0}$.
The same calculation with $V=W=1$ gives
$(w_I^0+w_{I^c}^0)^2=0$ and puts the sum in the Gysin kernel.
It is Lefschetz primitive: Gysin in degree six is injective
by dual weak Lefschetz in degree two.
Lift the diagonal model and all coordinate products to a
mixed-characteristic discrete valuation ring by lifting the
distinct $\lambda_i$ with their differences invertible.
The same matrix criterion gives a smooth projective model.
On a complex fibre the sum is an algebraic primitive $(2,2)$
class of square zero, hence zero by Hodge--Riemann positivity.
If needed, spread the finite data to a finitely generated
characteristic-zero field before embedding into $\mathbf C$.
Smooth proper base change proves
$w_I^0+w_{I^c}^0=0$ in the original rational cohomology.
This lifting is only of the explicitly diagonal bilinear
variety, not of the given abelian varieties or curve collections.

The finite flat map $f_B:B\to B_0$ pulls the two products back
to $Z_1,Z_2$ and $x,y$ back to $s,t$. Thus $w_1+w_2=0$.
Each $w_i^2=VW>0$, so the actual span is a positive line;
for $N=1$ positivity is immediate. \hfill$\square$

\subsection*{C44-A11a. The six products give a positive three-dimensional subspace}

Under the $N=2$ reduction define, for all six two-subsets,
\[
 \widetilde Z_I=f_B^*[Z_I^0],\qquad
 \widetilde w_I=\widetilde Z_I-\kappa^*a=f_B^*w_I^0.
\]
They are algebraic cycle classes even if additional inverse-image
curves are singular. Their ambient classes are $s^2t^2$, so the
$\widetilde w_I$ lie in the Gysin kernel.
For distinct noncomplementary $I,J$, write
$I\cap J=\{i\}$, $I^c\cap J^c=\{j\}$.
The products meet in the unique point $(e_i,e_j)$.
Linearizing the two equations gives
$dy_i+dx_j=0$ and $\lambda_i\,dy_i+\lambda_j\,dx_j=0$,
hence $dy_i=dx_j=0$.
The four remaining tangent coordinates on $B_0$ split into
complementary pairs used by the two surfaces. Their intersection
is transverse of multiplicity one. Together with self-intersection
two and complementary disjointness, subtraction of the ambient
term one gives
\[
 (w_I^0w_J^0)=
 \begin{cases}
 1,&J=I,\\
 -1,&J=I^c,\\
 0,&\text{otherwise}.
 \end{cases}                                                \tag{A11a.1}
\]
Each complementary sum vanishes by the explicit lifting argument.
Thus $w_{12}^0,w_{13}^0,w_{14}^0$ form an orthonormal basis of
the span of the six classes. Projection formula for the finite
degree-$VW$ map yields
\[
             (\widetilde w_{1i}\widetilde w_{1j})_B
                   =VW\delta_{ij},\qquad i,j=2,3,4.          \tag{A11a.2}
\]
This is an unconditional positive three-dimensional subspace for
smooth complete intersections. It is not asserted as an H9
Lefschetz-axis example, because of the next obstruction.

\subsection*{C44-A11b. The Lefschetz hypothesis forces $N\le1$}

Suppose every singular pencil member has exactly one ordinary
quadratic singularity. For $N\ge2$, the reduction above contradicts
that requirement.
First the pencil is unique for its scheme-theoretic axis in the
given linear system. If $H=L\boxtimes M$, twisted Koszul gives
\[
 0\longrightarrow H^{-1}\longrightarrow\mathcal O_{S\times T}^2
                    \longrightarrow\mathcal I_BH\longrightarrow0.
\]
Serre duality and ample cohomology vanishing give
$H^0(H^{-1})=H^1(H^{-1})=0$, so
$H^0(\mathcal I_BH)=\langle f_0,f_1\rangle$.
There is no alternative $H$-pencil with the same axis.

For each $i$, the target member
$\sum_j(\lambda_j-\lambda_i)x_jy_j=0$ has its unique singular
point $p_i=(e_i,e_i)$ outside $B_0$, since $\sum_jx_jy_j=1$ there.
In affine coordinates $x_i=y_i=1$ its equation is exactly
\[
                    q_i=\sum_{j\ne i}(\lambda_j-\lambda_i)x_jy_j,
                                                               \tag{A11b.1}
\]
a nonsingular quadratic form on six variables in every
characteristic.
Every preimage under $g=\phi\times\psi$ is a singular point of
the pulled-back member, with quadratic tangent form $q_i\circ dg$.
If $dg$ is not injective, a nonzero vector in its kernel is both
isotropic and polar-radical for this quadratic form. Its
projective quadric is singular; hence the hypersurface does not
have an ordinary quadratic singularity. This argument includes
characteristic two.

If all preimages instead have ordinary quadratic singularities,
$dg$ is invertible at each. The finite flat map is then \'etale
there, so the fibre consists of $VW$ distinct geometric points.
That one member has at least $VW$ nodes. Riemann--Roch makes
$V,W$ positive multiples of $3!$, so $VW>1$.
Thus either a forbidden nonordinary singularity occurs or more
than one node lies in the same member. These points are outside
the axis; blowing it up does not change the contradiction.
Therefore a genuine Lefschetz axis satisfies $N\le1$.
\hfill$\square$

\paragraph{Sources, scope and genuine checks.}
The cohomology and Riemann--Roch inputs are Milne,
\emph{Abelian Varieties}, v.~2.00, I Theorem 11.1(c),(b),
printed p.~54:
\texttt{https://jmilne.org/math/CourseNotes/AV.pdf}.
Nonzero sections exist here, so the unique nonzero cohomology
degree in part (c) is zero; Serre duality gives the negative
vanishings. Flatness uses Stacks Project, Lemma 10.128.1,
Tag 00R4: \texttt{https://stacks.math.columbia.edu/tag/00R4}.
The block identity, Chern obstruction, pencil classification,
coordinate intersections and singularity-count argument are
proved explicitly above.

Both curve families separately must be pairwise disjoint, and
each curve must be cut out by two divisors of the given
polarization. Arbitrary product-cycle collections are not covered.
The $N=2$ three-dimensional positive subspace is a smooth-CI
result; it cannot supply the claimed configuration in the usual
H9 Lefschetz setup. QH9-axis remains open.

The author checked the entire proof including both addenda on
22 November 2026 (Europe/Berlin). The root independently read
all three complete A11 proof parts and both complete addenda,
recording whole-proof PASS at 00:58:45 UTC that date.
Checks included the rank-three Chern obstruction, all-characteristic
matrix criterion, actual boundary relation, transverse products,
and strict Lefschetz qualification. Both agents are GPT-6 Astra Pro;
reasoning effort is not exposed.

% END work/axis/axis_product_ci_matching.tex

% BEGIN work/axis/relative_form_transfer.tex
% C44_POS /root; GPT-6 Astra Pro; effort not exposed.
\section{A relative correspondence criterion that retains the geometric form}

\begin{theorem}[C44-A08: positive-spectrum transfer along an extending correspondence]
Let $R$ be a discrete valuation ring of mixed characteristic $(0,p)$,
with fraction field admitting an embedding into $\mathbf C$.
Fix a prime $\ell\ne p$ for the special-fibre cohomology.
Let $\mathscr X/R$ be smooth projective of relative dimension $n$,
with a relatively ample class $h$, and let $0\le 2r\le n$.
Let $V$ be a relative algebraic direct summand of
$P^{2r}(\mathscr X)$, self-dual for
\[
 q_X(x,y)=(-1)^r\int_{\mathscr X}xyh^{n-2r}.
 \tag{C44.A08.1}
\]
Assume its projector and the inverse of this restricted pairing
are represented by algebraic correspondences over $R$ at the
level of Weil cohomology. Suppose there are abelian schemes
$\mathscr A_\nu/R$ and algebraic correspondences over $R$ giving
an injection on every geometric fibre
\[
 \alpha:V\longrightarrow
       \bigoplus_\nu H^{2s_\nu}(\mathscr A_\nu)(s_\nu-r).
 \tag{C44.A08.2}
\]
A relative split Chow-motive injection is sufficient for this
hypothesis; no such injection is assumed to exist in general.

Then, \emph{conditional on the unrefereed 032/Milne 2026
preprints}, the actual signed geometric form (C44.A08.1) is
positive definite on the rational algebraic classes of $V$ in
the geometric special fibre. In particular, no cycle-lifting
hypothesis is required for those special algebraic classes.
The assertion does not use Milne 2022, Proposition~4.2.
\end{theorem}

\begin{proof}
Choose relatively ample symmetric classes on the abelian schemes,
after a finite extension of the trait if necessary. The algebraic
Lefschetz operators of abelian varieties can be taken relatively.
Explicitly, on a relatively polarized abelian scheme of positive
dimension $g$, with $\theta=c_1(L)$ and $D=\int\theta^g/g!$, the cycles
\[
 L=\Delta_*\theta,\qquad
 \Lambda=\frac{\delta^*(\theta^{g-1})}{D(g-1)!},
 \qquad \delta(x,y)=y-x,
\]
give the raising and lowering operators fibrewise. Their
$\mathfrak{sl}_2$ identities and the usual rational polynomial
formulas give relative cohomological primitive and Lefschetz
projectors. Zero-dimensional factors are treated as Tate summands.
On every even-degree target factor choose the
orthogonal direct sum of its signed primitive Lefschetz pairings,
with positive rational coefficients on the summands, and include
the displayed Tate shifts. Write $Q_A$ for the resulting direct-sum
pairing. It is represented by an algebraic correspondence over
$R$. With the common real normalization of Tate factors, its
complex realization satisfies
\[
                  Q_A(Cu,\overline u)>0\quad(u\ne0),
 \tag{C44.A08.3}
\]
where $C$ is the Weil operator. This is the classical
Hodge--Riemann theorem on each primitive summand, extended by the
orthogonal Lefschetz decomposition. On algebraic classes in the
special fibre the same $Q_A$ is positive definite, conditional
on 032 Corollary~8.4 and Milne's 2026 report. Indeed its primitive
projections are algebraic and each signed primitive form has
that positivity. This is precisely where the unrefereed input
enters.

Pull back this pairing and put
\[
                  q_1=\alpha^*Q_A,\qquad
                  c=q_X^{-1}q_1\in\operatorname{End}(V).
 \tag{C44.A08.4}
\]
The assumed algebraic inverse and the transpose of $\alpha$
make $c$ an algebraic correspondence over the trait, at the
cohomological level required here. The characteristic polynomial
of its Betti action is a polynomial $P(T)\in\mathbf Q[T]$.
We show that every root of $P$ is a strictly positive real number.

On the complex generic fibre both $q_X$ and $q_1$ are
polarizations in the following explicit sense: their associated
Hermitian forms
\[
 H_X(x,y)=q_X(Cx,\overline y),\qquad
 H_1(x,y)=q_1(Cx,\overline y)
\]
are positive definite. For $q_X$ this is the primitive
Hodge--Riemann relation; for $q_1$ it follows from the Hodge
compatibility and injectivity of $\alpha$ and (C44.A08.3).
The algebraic correspondence $c$ preserves Hodge types and the
real structure, hence commutes with $C$. Symmetry of the two
bilinear forms gives
\[
 H_X(cx,y)=H_X(x,cy),\qquad H_X(cx,x)=H_1(x,x)>0
 \quad(x\ne0).
\]
Consequently $c$ is self-adjoint and positive for $H_X$.
It is diagonalizable with strictly positive real eigenvalues,
as claimed. This argument uses the common Weil operator; the
indefinite geometric bilinear form on all Betti classes is not
mistaken for a positive form.

Smooth proper base change and compatibility of cycle classes
with specialization identify the characteristic polynomial of
$c$ in the two geometric fibres. Thus $P(c)=0$ in the special
fibre, and $P(0)\ne0$. Cayley--Hamilton expresses $c^{-1}$ as a
polynomial in $c$ with rational coefficients. Both $c$ and
$c^{-1}$ therefore preserve special rational algebraic classes.

For complete precision no unproved rational dimension assertion
about the space of algebraic classes is needed. For any such
class $z$, let
\[
 U_z=\operatorname{span}_{\mathbf Q}
       \{z,cz,\ldots,c^{d-1}z\},\qquad d=\deg P.
\]
It is finite-dimensional, consists of algebraic classes, and is
stable under $c$ and $c^{-1}$. On $U_z$, $q_1$ is positive definite:
$\alpha z\ne0$ for $z\ne0$, and $Q_A$ is positive on algebraic
classes by the conditional input. Moreover,
\[
 q_1(cx,y)=q_X(c^2x,y)=q_X(cx,cy)=q_1(x,cy).
\]
Thus $c$ is $q_1$-self-adjoint on $U_z$. Its eigenvalues are among
the positive roots of $P$, so $c^{-1}$ is positive for $q_1$.
Finally
\[
                   q_X(z,z)=q_1(c^{-1}z,z)>0
                   \qquad(z\ne0).
 \tag{C44.A08.5}
\]
Since $z$ was arbitrary, this proves the assertion on all special
algebraic classes. The newly constructed inverse $c^{-1}$ is an
actual rational polynomial in $c$; algebraicity of $q_X^{-1}$ was
an explicit hypothesis. No square root,
abstract comparison isometry, or unproved positivity transfer
is inserted.
\end{proof}

\begin{corollary}[C44-A08a: a sufficient correspondence condition for a fixed axis]
Let $B$ be a smooth axis in an abelian $(2r+2)$-fold as in
QH9-axis. Suppose an embedded mixed-characteristic lift from
C44-A05 carries an extending algebraic injection of its variable
middle cohomology into a sum of Tate-shifted even cohomology of
abelian schemes. Then the H9-T7 axis form on $A_B$ is positive
definite, \emph{conditional on the unrefereed 032/Milne 2026
preprints}.
\end{corollary}

\begin{proof}
The variable projector in middle cohomology is
\[
       p_{\rm var}=1-\kappa^*(L_h^2)^{-1}\kappa_*.
 \tag{C44.A08.6}
\]
It is algebraic relatively, since the ambient abelian Lefschetz
inverse is algebraic relatively. It is self-adjoint for the
middle intersection pairing. Its image is $\ker\kappa_*$,
which is primitive by C44-A01; the restriction of the middle
pairing to this image is nondegenerate. Its inverse as a duality
morphism is represented by the transpose self-adjoint projector
and the middle Poincare duality correspondence. Hence all the
hypotheses of C44-A08 hold, and (C44.A08.1) is exactly
$(-1)^r\int_Bww'$, the second summand of H9-T7.
\end{proof}

\paragraph{Exact remaining gate and Grothendieck's frameworks.}
This is a criterion, not an assertion that every axis admits
such a relative abelian correspondence. C44-A05 lifts the
embedded geometry but supplies no injection (C44.A08.2).
C44-A02 excludes graph maps as a source for a nonzero variable
injection. An arbitrary characteristic-$p$ correspondence that
does not extend to a characteristic-zero fibre also does not
meet the criterion. Thus QH9-axis is not closed. The distinction
is the same geometric-polarization issue exposed by the prior
JAC-P20 transfer audit, but the present hypothesis permits an
explicit proof via a positive spectrum. It gives a sufficient
condition for cote~44 Th.~5, scan~004, with the primitive
qualification of scan~008, and proves the relevant $I$ summand
in the terminology of cote~16 scans~049--050.

The cohomological tools are Weil cohomology, Poincare duality,
Lefschetz decomposition, and smooth proper base change (SGA~4,
XVI, Corollary~2.2), together with algebraic correspondences and
Cayley--Hamilton. Relative Fourier/Lefschetz operators on an
abelian scheme are supplied by the same Poincare-bundle
correspondences as Beauville, \emph{The action of SL(2) on abelian
varieties}, JRMS 25 (2010), Theorem~3.3, pp.~258--259, and
Theorem~4.2, p.~260;
their identities are fibrewise identities in Weil cohomology.
The characteristic-zero sign is the Hodge--Riemann bilinear
relation (Voisin, \emph{Hodge Theory and Complex Algebraic
Geometry I}, Theorem~6.32). The only conjectural positivity
input is the expressly conditional 032/Milne assertion on the
abelian special fibre.

% END work/axis/relative_form_transfer.tex

% BEGIN work/axis/axis_frobenius_selection.tex
\section{Selecting closed fibres with exactly the prescribed algebraic part}
\label{sec:c44-a12}

\begin{proposition}[C44-A12: even orthogonal complement and closed-point selection]
Let \(U/\mathbf F_q\) be smooth, geometrically connected and
quasiprojective, and let \(f:\mathcal B\to U\) be smooth projective
of relative dimension \(2r\). Fix \(\ell\ne\operatorname{char}\mathbf F_q\).
Let \(\mathcal V\subset R^{2r}f_*\mathbf Q_\ell(r)\) be a lisse
orthogonal direct summand for the signed middle pairing
\(Q(x,y)=(-1)^r\int xy\).
Suppose that rational relative cycle classes
\(z_1,\ldots,z_s\) lie in \(\mathcal V\), are linearly independent
and have constant nonsingular rational Gram matrix \(G\).
Write
\[
 \mathcal V=\mathcal K\perp\mathcal W,\qquad
 \mathcal K=\sum_i\mathbf Q_\ell z_i .
\]
Assume that \(d=\operatorname{rank}\mathcal W>0\) is even, that the
geometric monodromy closure on \(\mathcal W\) contains
\(\operatorname{SO}(\mathcal W,Q)\), and that every normalized
closed-point Frobenius characteristic polynomial on \(\mathcal W\)
lies in \(\mathbf Q[T]\).
Then a Zariski-dense set of closed points \(u\) has
\[
 \Alg^r(\mathcal B_{\bar u})\cap\mathcal V_{\bar u}
       =\sum_i\mathbf Q z_{i,\bar u}.                    \tag{A12.1}
\]
In particular, if \(G\) is positive definite, the entire algebraic
part of this summand is positive definite at those closed points.
No Tate conjecture, cycle-lifting assertion or 032 input is used.
\end{proposition}

\begin{proof}
Put \(N=\operatorname{lcm}(1,\ldots,2d^2)\).
The elementary cyclotomic degree bound used in C44-A03 says that a
root of unity which is a root of a rational degree-\(d\) polynomial
has order at most \(2d^2\). Thus a Frobenius polynomial here has a
root-of-unity root if and only if
\(\det(F^N-1)=0\), without a semisimplicity assumption.

On the even orthogonal group's identity component the regular function
\[
                     P(g)=\det(g^N-1)
\]
is not identically zero. After passing to an algebraic closure,
a maximal torus has eigenvalues
\(t_1,t_1^{-1},\ldots,t_{d/2},t_{d/2}^{-1}\), and one may choose
every \(t_i^N\ne1\).
The geometric monodromy image therefore contains some element
\(\gamma\) with \(P(\gamma)\ne0\).

Let \(\Gamma\) be the compact arithmetic monodromy image on
\(\mathcal W\). The subset \(P\ne0\) in \(\Gamma\) is open for
the \(\ell\)-adic topology and invariant under conjugation.
Choose an open normal subgroup \(K\) of \(\Gamma\) such that
\(\gamma K\) lies in this subset. The union of the conjugates of
\(\gamma K\) still lies there. It is the inverse image of a
nonempty conjugacy class in the finite quotient \(\Gamma/K\).

The finite quotient is realized by a finite etale Galois cover of
\(U\). Finite-field Chebotarev supplies a closed point whose
Frobenius image lies in that class. For precision concerning the
constant field, the selected class lies in the geometric subgroup
because \(\gamma\) was chosen geometrically. One may first extend
the constants so that the finite quotient cover has a geometrically
connected component and then apply Chebotarev to that component.
Equivalently, one takes extension degrees in the appropriate
constant-field congruence class in the general theorem.
The resulting Frobenius has \(P(F)\ne0\), hence has no
root-of-unity eigenvalue on \(\mathcal W\).
If one uses a rational point over an extension rather than a closed
point over the original field, its Frobenius is a positive power of
the latter. A root-of-unity eigenvalue of the latter would also
give one for that power, so the conclusion descends to its
underlying closed point.

The same argument applies after replacing \(U\) by any nonempty
open subset: its fundamental group surjects onto that of \(U\),
so the monodromy image does not shrink. This proves the claimed
Zariski density.

Now let \(z\) be an algebraic rational cycle class in
\(\mathcal V_{\bar u}\). Its orthogonal projection to \(\mathcal K\)
has coefficients
\[
  (a_1,\ldots,a_s)^{\mathsf t}
       =G^{-1}\bigl(Q(z,z_1),\ldots,Q(z,z_s)\bigr)^{\mathsf t}
       \in\mathbf Q^s .                                  \tag{A12.2}
\]
The rationality follows from intersection numbers of rational
algebraic cycles; it is not merely a statement about
\(\mathbf Q_\ell\)-linear spans.
Thus \(z-\sum_i a_i z_i\) is still an algebraic class and lies
in \(\mathcal W_{\bar u}\).
Every algebraic cycle over \(\overline{\mathbf F}_q\) is defined
over a finite field extension. After the Tate twist its class
is fixed by a positive power of Frobenius.
There are no such nonzero vectors on \(\mathcal W_{\bar u}\),
because its Frobenius has no root-of-unity eigenvalues.
This proves (A12.1), and the sign assertion follows from \(G\).
\end{proof}

\paragraph{Exact Chebotarev source.}
The finite-quotient existence statement, including the constant-field
coset, is Alexei Entin, \emph{Monodromy of hyperplane sections of curves
and decomposition statistics over finite fields}, Theorem 3 in
\S4, author version arXiv:1805.05454v3, pp.~9--11;
published in \emph{IMRN} 2021, pp.~10409--10441,
doi:10.1093/imrn/rnz120.
Only its nonempty-class conclusion over sufficiently large finite
extensions is used. No numerical error coefficient from the
Chebotarev sketch in Stacks, Tag 03W3, is used.

\begin{corollary}[C44-A12a: application to a family of axes]
Suppose a family of axes has constant abelian ambient variety,
with its variable middle summand \(\mathcal V=\ker\kappa_*\).
If a known positive cycle span has even orthogonal complement
with the monodromy hypothesis of C44-A12, then a Zariski-dense
set of closed axes has exactly that algebraic variable span.
The rational-polynomial hypothesis holds whenever the full
middle and ambient Frobenius polynomials are rational:
the polynomial on the complement is their quotient, followed
by removing the factors supplied by the known cycle classes.
\end{corollary}

\begin{proof}
Weak Lefschetz and hard Lefschetz give the orthogonal splitting
of the middle cohomology into the ambient image and the variable
kernel. The ambient image has the ambient Frobenius polynomial.
The classes \(z_i\) have normalized Frobenius eigenvalue one,
so their factor is \((T-1)^s\).
The quotient polynomial is rational. C44-A12 applies.
\end{proof}

\paragraph{Relation to the main question.}
For cote 44, Th.~5, scan 004 and the primitive margin on scan 008,
this is a sufficient construction principle for complete control
of \(A_B\) at actual finite-field points, rather than merely at
a geometric generic parameter point. It does not assert the
monodromy hypothesis for an arbitrary fixed axis.
The odd-rank obstruction of C44-A03b explains why the known
cycle must first be removed in the abelian-sixfold case.

% END work/axis/axis_frobenius_selection.tex

% BEGIN work/scan_w/closed_rank_one_axis.tex
% C44_POS /root/scan_w; complete draft, independent review requested.
\section{Closed axes with a prescribed nonzero rank-one algebraic kernel}
\label{sec:c44-a13}

\begin{theorem}[C44-A13: closed rank-one axes in a supersingular abelian sixfold]
Let $E/\mathbf Q$ be the elliptic curve $y^2=x^3-x$, with origin $o$,
and put
\[
 A=E^6,\qquad H=\sum_{i=1}^6h_i,\qquad
 h_i=\operatorname{pr}_i^*[o],\qquad
 Y=E^2\times\{o\}^4.
\]
For a sufficiently large integer $t$, there is a one-parameter family
of genuine smooth axes
\[
 B_u=D_0\cap D_u\subset A,\qquad D_0,D_u\in|tH|,
 \qquad Y\subset B_u,
 \tag{A13.1}
\]
defined on a nonempty open $U\subset\mathbf P^1_{\mathbf Q}$.
After excluding finitely many primes (depending on this model and a
fixed coefficient prime $\ell$), there are infinitely many primes
$p\equiv3\pmod4$ for which the following holds.
The reduction $A_p=E_p^6$ is supersingular, and a Zariski-dense set
of closed points $u$ of $U_{\mathbf F_p}$ has
\[
 A_{B_u}:=\Alg^2(B_{\bar u})\cap
 \ker\bigl(\kappa_*:H^4(B_{\bar u})\longrightarrow H^8(A_{\bar u})\bigr)
       =\mathbf Q w_u,
 \tag{A13.2}
\]
where $w_u$ is the explicitly algebraic primitive class
\[
 \begin{split}
 \alpha_0={}&\frac14h_1h_2
 -\frac1{12}\sum_{i\in\{1,2\},\,j\in\{3,4,5,6\}}h_ih_j
 +\frac1{12}\sum_{3\leq i<j\leq6}h_ih_j,\\
 w_u={}&[Y]_{B_u}-t^{-2}\kappa^*\alpha_0.
 \end{split}
 \tag{A13.3}
\]
Its actual geometric intersection form is
\[
                    \int_{B_u}w_u^2
                       =6t^2-\frac1{4t^2}>0.
 \tag{A13.4}
\]
Thus the required sign on the entire nonzero $A_{B_u}$ is
unconditional at these closed finite-field axes. Their full variable
cohomology is nonzero, of rank $3600t^6-165$.
The original pencil generated by $D_0,D_u$ may be required to be a
Lefschetz pencil. No Tate conjecture or unrefereed 032 result is used.
The construction does not assert that the entire variable motive
is of abelian type.
\end{theorem}

\begin{proof}
\emph{The family and its two pencils.}
Let $\pi:Z=\operatorname{Bl}_Y A\to A$ and let $F$ be the exceptional
divisor. For $t\gg0$, the line bundle
$L'=\pi^*\mathcal O_A(tH)\otimes\mathcal O_Z(-F)$ is very ample.
Its sections are exactly the sections of $\mathcal I_Y(tH)$.
Choose a general member $\widetilde D_0$ of $|L'|$.
Its image $D_0\subset A$ is smooth and contains $Y$, and
$\widetilde D_0=\operatorname{Bl}_Y D_0$.
Indeed the normal derivative of a section along $Y$ is a section
of $\mathcal O_Y(tH)^{\oplus4}$. For $t$ sufficiently large these
derivatives are generated by the chosen section space; their common
zero locus has expected codimension four, exceeding $\dim Y=2$.
Away from $Y$, Bertini applies. Likewise two general sections give
a $2\times4$ matrix of normal derivatives, whose rank-one locus
has expected codimension three. It misses $Y$, so their original
intersection $B$ is smooth along $Y$ as well as away from it.
These statements can equivalently be obtained by applying the
usual jet-generation bounds to $\mathcal I_Y(tH)$.

Embed $Z$ by $L'$. The induced embedding of $\widetilde D_0$ admits
Lefschetz pencils in characteristic zero: SGA~7~II, XVII,
Theorem~2.5.2, printed p.~217. Choose a general pencil of these
hyperplane sections of $\widetilde D_0$. It gives a family
\[
 \widetilde B_u=\widetilde D_0\cap\widetilde D_u
                  =\operatorname{Bl}_Y B_u.
 \tag{A13.5}
\]
The restriction map from all sections on $Z$ to the complete
section space on $\widetilde D_0$ need not be surjective; the
induced projective embedding is all that is used.
Remove the finitely many auxiliary-pencil critical values and
any additional values for which the original $B_u$ is not smooth.
The open that remains is nonempty by the preceding general-pair
argument. We may also require that the line spanned by
$D_0,D_u$ give a Lefschetz pencil on $A$. This is an open nonempty
condition on general pairs in this system. It follows alternatively
from the Lefschetz condition for the corresponding pencil on $Z$:
smoothness of $B_u$ along $Y$ makes every member of that pencil
smooth along $Y$, so its critical points are away from the
exceptional divisor and coincide with the original critical points.
For a general fixed $\widetilde D_0$, the open set of Lefschetz
lines through that member is nonempty. Choose the auxiliary
pencil generically so that its general member gives such a line;
then remove the finitely many parameters outside this open.
The auxiliary pencil varying the axis and the original pencil
having that axis are different pencils.

All the required nonempty open conditions are defined over
$\mathbf Q$. Choose the first member and the auxiliary pencil
over $\mathbf Q$ in these opens. This proves the asserted
number-field model of (A13.1).

\emph{The primitive class and its norm.}
Write $h=tH$ and $\kappa:B\hookrightarrow A$.
The algebraic ambient projection in degree four is
$\kappa^*(L_h^2)^{-1}\kappa_*$; it is defined since
$\kappa_*\kappa^*=L_h^2:H^4(A)\to H^8(A)$ is an isomorphism.
For the particular class $[Y]=h_3h_4h_5h_6$, direct multiplication,
using $h_i^2=0$, gives
\[
                  H^2\alpha_0=[Y].
 \tag{A13.6}
\]
For clarity, the coefficients of a four-index monomial are
$12(1/12)=1$ when all indices are in $\{3,4,5,6\}$,
$2(3(-1/12)+3(1/12))=0$ when precisely one is in $\{1,2\}$,
and $2(1/4+4(-1/12)+1/12)=0$ when both are.
Consequently $(L_h^2)^{-1}[Y]=t^{-2}\alpha_0$, proving
$\kappa_*w=0$ and (A13.3).

The normal bundle sequence is
\[
 0\longrightarrow N_{Y/B}\longrightarrow
 \mathcal O_Y^{\oplus4}\longrightarrow
 \mathcal O_Y(h)^{\oplus2}\longrightarrow0.
 \tag{A13.7}
\]
Thus $c_2(N_{Y/B})=3h^2|_Y$ and
$[Y]_B^2=3\int_Yh^2=6t^2$.
Orthogonal projection subtracts the ambient norm:
\[
 \int_Bw^2
 =\int_B[Y]_B^2-\int_Yt^{-2}\alpha_0
 =6t^2-\frac1{4t^2}.
\]
This computation is valid over every good characteristic, without
positivity assumptions. In particular $w$ is nonzero and
nonisotropic.

\emph{The exact vanishing system.}
Let $q:\widetilde B\to B$ be (A13.5), and write $j$ and $p$ for
the inclusion and projection of its exceptional divisor
$\mathbf P(N_{Y/B})$. The blowup formula gives
\[
 \begin{split}
 H^4(Z)&=H^4(A)\oplus H^2(Y)(-1)\oplus\mathbf Q_\ell(-2),\\
 H^4(\widetilde D_0)&\simeq H^4(Z),\\
 H^4(\widetilde B)&=q^*H^4(B)\perp j_*p^*H^2(Y)(-1).
 \end{split}
 \tag{A13.8}
\]
The middle isomorphism is weak Lefschetz for the ample
fivefold $\widetilde D_0\subset Z$. The last decomposition is
orthogonal: the mixed pairing is zero by the projection formula,
and the exceptional pairing is the negative of the middle pairing
on $H^2(Y)$. Pullback $q^*$ preserves the middle pairing.
The additive blowup formula is SGA~7~II, XVIII,
Theorem~2.2 and (2.2.2), printed pp.~260--261; the stated maps
and pairings follow directly from the projective-bundle formula
and $j^*j_*=-\xi$, where $\xi=c_1\mathcal O_{\mathbf P(N)}(1)$.

Restriction from $Z$ sends its $H^2(Y)(-1)$ summand onto the
exceptional summand in (A13.8). The remaining exceptional
generator restricts to $j_*\xi$. The excess-intersection formula
for the rank-two normal bundle says
\[
 q^*[Y]_B=j_*(\xi+p^*c_1(N_{Y/B})).
 \tag{A13.9}
\]
Therefore the image of $H^4(\widetilde D_0)$ in
$H^4(\widetilde B)$ is exactly
\[
 j_*p^*H^2(Y)(-1)
  +q^*\bigl(\kappa^*H^4(A)+\mathbf Q_\ell[Y]_B\bigr).
 \tag{A13.10}
\]
There is no assertion that a restriction of an ambient
correspondence detects variable cohomology. The additional
class $[Y]_B$ enters through the exceptional divisor and has
nonzero primitive component $w$.

Let $\mathcal V=\ker\kappa_*$ in the family and put
$\mathcal W=\mathcal V\cap w^\perp$.
Equations (A13.8)--(A13.10) identify $q^*\mathcal W$, isometrically,
with the vanishing cohomology of the auxiliary Lefschetz pencil
on $\widetilde D_0$. In particular this is an identification of
local systems and pairings, not just a dimension comparison.
All maps are algebraic and all these identities persist in the
relative model at good primes.

\emph{Rank and characteristic-zero monodromy.}
Since $H^6/6!=1$, one has $h^6=720t^6$.
From $c(T_B)=(1+h)^{-2}|_B$, the topological Euler number is
$5h^6=3600t^6$. Weak Lefschetz and Poincare duality give
\[
 b_4(B)=3600t^6+330,\qquad
 \operatorname{rank}\mathcal V=3600t^6-165,\qquad
 d:=\operatorname{rank}\mathcal W=3600t^6-166.
 \tag{A13.11}
\]
In particular $d>0$ is even. To exclude finite monodromy in
characteristic zero, the Koszul resolution gives
$\chi(\mathcal O_B)=(2^6-2)t^6=62t^6$.
The groups $H^i(B,\mathcal O_B)$ for $i<4$ have dimensions
$1,6,15,20$, so
\[
 h^{4,0}(B)=62t^6+10,\qquad
 h^{4,0}(\mathcal W)=62t^6-5>0.
 \tag{A13.12}
\]
Removing $w$ changes only type $(2,2)$, and the ambient
$H^{4,0}(A)$ has dimension $15$.

By SGA~7~II, XIX, Proposition~3.4, printed p.~334, finite
monodromy of an even-dimensional Lefschetz pencil over
$\mathbf C$ is equivalent to its vanishing cohomology being
purely of type $(2,2)$. Thus the monodromy here is infinite.
Deligne, \emph{Weil II}, Theorem~4.4.1 and Lemma~4.4.2s,
printed pp.~227--228, then give the full orthogonal geometric
monodromy group on $\mathcal W$.
Removing the additional finitely many parameter values used
above does not reduce this image, since the fundamental group
of a smaller nonempty open curve surjects onto that of the
larger open curve.

\emph{Preserving that monodromy at good primes.}
Spread this one fixed auxiliary pencil, its smooth proper
families, and its finitely many omitted values over
$\mathbf Z[1/M]$. Increase $M$ so that the total discriminant
and other omitted values form a finite etale divisor in the
relative $\mathbf P^1$, and so that all the stated smoothness
and Lefschetz conditions hold. This is possible because there
are only finitely many points and bad reduction conditions
in this fixed characteristic-zero model.

Fix $\ell$. The local system $\mathcal W$ has a stable
$\mathbf Z_\ell$ lattice: intersect it with the torsion-free
integral cohomology lattice. Enlarge the finite exceptional
set of primes to include $2,\ell$ and the prime divisors of
$|\operatorname{GL}_d(\mathbf F_\ell)|$.
For any remaining $p$, every finite congruence quotient of its
compact integral monodromy image has order prime to $p$,
since the kernel of reduction to $\operatorname{GL}_d(\mathbf
F_\ell)$ is pro-$\ell$. Its image is therefore pro-prime-to-$p$.

Apply the prime-to-$p$ specialization isomorphism for a smooth
proper relative curve minus a relative normal-crossing divisor.
In the present case the curve is $\mathbf P^1$ and the divisor
is finite etale. After strict henselization at $p$, SGA~1,
XIII, Corollary~2.9, construction~2.10 and the proof of
Corollary~2.12 identify the prime-to-$p$ geometric fundamental
groups in characteristic zero and $p$, compatibly with the
local system extending over the trait. Hence their monodromy
images on $\mathcal W$ are conjugate. In particular the
characteristic-$p$ geometric closure is still the full
orthogonal group. This explicit argument avoids inferring
infinite monodromy in characteristic $p$ from a Hodge number.
The source's comparison is on pp.~289--292 of the revised
edition (original pp.~390--394). SGA~7~II, XIX, \S3.6,
printed p.~335, uses the same relative-pencil comparison;
the prime exclusions above give the stated sufficient version
without any additional tame-specialization assumption.

There are infinitely many primes $p\equiv3\pmod4$ outside
this finite set. For such $p$, the character sum for
$y^2=x^3-x$ cancels under $x\mapsto-x$, so its trace is zero
and its Frobenius polynomial is $T^2+p$. Its Newton slopes
are both $1/2$, proving supersingularity of $E_p$ and $A_p$.

\emph{Closed-point selection and the whole algebraic kernel.}
The relative class $w$ has the constant positive rational
Gram matrix (A13.4). The middle Frobenius polynomial of each
smooth $B_u$ and that of the constant ambient $A$ are rational;
their quotient, followed by division by the factor $T-1$
of the normalized class $w$, is the rational Frobenius
polynomial of $\mathcal W$.
Apply C44-A12 and C44-A12a, using the even rank in (A13.11)
and the just-proved orthogonal monodromy. A Zariski-dense set
of closed points has no root-of-unity normalized Frobenius
eigenvalue on $\mathcal W$. Every algebraic class has finite
Frobenius orbit, so its $\mathcal W$ component vanishes.
The rational Gram projection from A12 gives the equality
with $\mathbf Qw$, not merely with $\mathbf Q_\ell w$.
Equations (A13.2) and (A13.4) follow.

Finally, the selected closed parameter lifts over an unramified
mixed-characteristic trait in the fixed open parameter model.
The subvariety $Y$ and every divisor in the polynomial
expression (A13.3) lift along with it. Thus all of $A_{B_u}$
lifts in this example. No nonliftable class or general
fixed-axis sign theorem is being claimed.
\end{proof}

\paragraph{Primary sources actually read.}
The original SGA~7~II scan, including XVII~2.5.2,
XVIII~2.2, XVIII~6.6--6.7, XIX~1.4 and XIX~3.4--3.6, is
\url{https://publications.ias.edu/sites/default/files/Number12.pdf}.
The original \emph{Weil II} proof of 4.4.1 and 4.4.2s is in
\url{https://publications.ias.edu/sites/default/files/Number40.pdf}.
SGA~1 XIII~2.9--2.12 and its specialization proof were read in
the revised primary edition
\url{https://arxiv.org/pdf/math/0206203}.
For the finite-field selection, C44-A12 uses Entin's Theorem~3,
whose whole primary proof was independently read in
\url{https://arxiv.org/pdf/1805.05454v3}, pp.~9--11.
This result supplies a new nonzero closed-axis family for cote~44,
Th.~5, scans~004 and~008, while retaining the general axis question.

% END work/scan_w/closed_rank_one_axis.tex

% BEGIN work/scan_w/closed_odd_rank_axes.tex
% C44_POS /root/scan_w; separate extension preserving the checked A13 file.
\section{An odd number of prescribed primitive axis classes}
\label{sec:c44-a13b}

\begin{theorem}[C44-A13b: closed axes with a prescribed odd-rank kernel]
Let $N\geq1$ be odd. Retain the elliptic curve $E$, the sixfold
$A=E^6$, the product polarization $H$, and the class $\alpha_0$
of C44-A13. There is a number field $K$ and a collection of
distinct points $a_1,\ldots,a_N\in E^4(K)$ such that the
pairwise disjoint abelian surfaces
\[
                     Y_i=E^2\times\{a_i\}\subset A_K
\]
have the following property. For all sufficiently large integers
$t$, a nonempty open $U\subset\mathbf P^1_K$ carries smooth axes
$B_u=D_0\cap D_u$, with $D_0,D_u\in|tH|$ and $Y_i\subset B_u$
for every $i$, each the axis of a genuine Lefschetz pencil.
There is a finite set $S$ of rational primes, depending on these
chosen data and a coefficient prime $\ell$, such that, for every
$p\equiv3\pmod4$ outside $S$ and every place $\mathfrak p$ of $K$
above $p$, the reduction of this family has a Zariski-dense set
of closed points $u$ for which
\[
 \begin{split}
 A_{B_u}
 &=\Alg^2(B_{\bar u})\cap\ker(\kappa_*:H^4(B_{\bar u})
                                      \longrightarrow H^8(A_{\bar u}))\\
 &=\bigoplus_{i=1}^N\mathbf Q w_{i,u},\qquad
 w_{i,u}=[Y_i]_{B_u}-t^{-2}\kappa^*\alpha_0.
 \end{split}
 \tag{A13b.1}
\]
The ambient abelian sixfold is supersingular. On this entire
algebraic variable kernel the actual geometric intersection
matrix is
\[
 G_N=6t^2 I_N-\frac1{4t^2}J_N,
 \tag{A13b.2}
\]
where $J_N$ is the matrix with every entry equal to one. In
particular the matrix is positive definite once $24t^4>N$.
Every class in (A13b.1) has a relative algebraic lift in the
chosen incidence model over a mixed-characteristic trait.
The assertions are unconditional; neither Tate nor 032 is used.
\end{theorem}

\begin{proof}
Choose $N$ distinct torsion points of $E^4(\overline{\mathbf Q})$
and a number field $K$ over which they are all rational. This
is only a convenient way to choose the centres; no torsion
property enters their intersection or monodromy calculations.
Translations act trivially on cohomology, so
\[
       [Y_i]_A=h_3h_4h_5h_6=H^2\alpha_0
                 \quad\hbox{for every }i.
 \tag{A13b.3}
\]
All centres have trivial rank-four normal bundle in $A$.

Let $Y=\coprod_iY_i$, $Z=\operatorname{Bl}_Y A$, and let $F$
be the sum of its exceptional divisors. For $t\gg0$,
$L'=\pi^*\mathcal O_A(tH)\otimes\mathcal O_Z(-F)$ is very
ample and the normal derivatives of its section space are
generated along each $Y_i$. The normal derivative of one
section is a row of four entries on a surface; its zero
locus has expected codimension four. For two sections,
rank at most one has expected codimension three. The finite
disjoint union of centres therefore introduces no zero or
rank-drop point for general sections. Bertini off $Y$ gives
smooth original divisors and smooth original intersections.
Choose a general $\widetilde D_0\in|L'|$ and a general
auxiliary hyperplane pencil on its induced projective embedding.
Exactly as in A13, this gives
\[
       \widetilde B_u=\operatorname{Bl}_Y B_u
                       \subset\widetilde D_0.
 \tag{A13b.4}
\]
There are finitely many parameters to remove to make $B_u$
smooth. For the fixed general $\widetilde D_0$, the Lefschetz
lines through it form a nonempty open set of lines in $|L'|$.
The auxiliary pencil can be chosen so that its general
$\widetilde D_u$ determines a line in this open; remove the
finitely many exceptional parameters. The two normal rows
are independent at every point of every centre, so all
original pencil members are smooth along $Y$. Consequently
their critical points are off the exceptional locus and the
original pencil generated by $D_0,D_u$ is Lefschetz as well.
These are nonempty open conditions over $K$, so the choices
may be made over $K$. This construction uses the induced
embedding of $\widetilde D_0$, and does not require surjectivity
onto its complete section space.

By (A13b.3), every $w_i$ belongs to $\ker\kappa_*$. For each
centre the normal bundle sequence is
\[
 0\longrightarrow N_{Y_i/B}\longrightarrow\mathcal O_{Y_i}^4
       \longrightarrow\mathcal O_{Y_i}(tH)^2\longrightarrow0.
\]
Hence $[Y_i]_B^2=6t^2$. Distinct $Y_i$ are disjoint, so
$[Y_i]_B[Y_j]_B=0$ for $i\ne j$. The restriction of the
divisor polynomial $\alpha_0$ to any translate $Y_i$ has
degree $1/4$, and $H^2\alpha_0=[Y_i]_A$. Orthogonal ambient
projection thus gives
\[
           \int_Bw_iw_j=6t^2\delta_{ij}-\frac1{4t^2}.
\]
The eigenvalues of this matrix are $6t^2$ on the sum-zero
subspace and $6t^2-N/(4t^2)$ on the all-ones line. Taking
$t$ still larger so that $24t^4>N$ proves (A13b.2) and the
linear independence of these $N$ classes in every good fibre.

The vanishing-system identification in A13 extends as an
actual orthogonal identification, as follows. The blowup
formula and weak Lefschetz give
\[
 H^4(\widetilde D_0)\simeq H^4(Z)
   =H^4(A)\oplus\bigoplus_{i=1}^N
                  \bigl(H^2(Y_i)(-1)\oplus\mathbf Q_\ell(-2)\bigr).
 \tag{A13b.5}
\]
For $q:\widetilde B\to B$, let $j_i,p_i,\xi_i$ denote the
maps and tautological class on its $i$th exceptional divisor.
Then
\[
 H^4(\widetilde B)=q^*H^4(B)\perp
              \bigoplus_{i=1}^N j_{i*}p_i^*H^2(Y_i)(-1).
 \tag{A13b.6}
\]
Each exceptional summand has pairing minus the pairing on
$H^2(Y_i)$; different exceptional summands are orthogonal
because their supports are disjoint. The excess formula is
$q^*[Y_i]_B=j_{i*}(\xi_i+p_i^*c_1N_{Y_i/B})$. Thus the image
of (A13b.5) in (A13b.6) is exactly the exceptional sum plus
\[
             q^*\bigl(\kappa^*H^4(A)+\sum_i\mathbf Q_\ell[Y_i]_B\bigr).
\]
The ambient pairing is nondegenerate by hard Lefschetz, and
the pairing on the additional primitive span is $G_N$.
Writing $\mathcal V=\ker\kappa_*$, the vanishing cohomology
of the auxiliary pencil is therefore isometric to
\[
 q^*\mathcal W_N,
 \qquad
 \mathcal W_N=\mathcal V\cap\bigcap_{i=1}^Nw_i^\perp.
 \tag{A13b.7}
\]
All maps are relative algebraic maps. This identifies local
systems and their forms, not just their dimensions.

The topology and Koszul calculation for a smooth intersection
of two members of $|tH|$ are unchanged by the prescribed
centres. Since only type $(2,2)$ classes have been removed,
\[
 d_N:=\operatorname{rank}\mathcal W_N=3600t^6-165-N,
 \qquad h^{4,0}(\mathcal W_N)=62t^6-5.
 \tag{A13b.8}
\]
For our sufficiently large $t$, the rank is positive and even,
because $N$ is odd, and the displayed Hodge number is positive.
SGA~7~II, XIX, Proposition~3.4 excludes finite characteristic-zero
vanishing monodromy; Deligne, \emph{Weil II}, 4.4.1 and 4.4.2s
then give full orthogonal geometric monodromy on $\mathcal W_N$.
These are precisely the primary monodromy results used in A13,
applied to the odd-dimensional $\widetilde D_0$ in (A13b.4).

Spread this fixed data over $\mathcal O_K[1/M]$. Enlarge the
finite set of excluded rational primes to include those below
any bad reduction place of the elliptic curve, centres, blowups,
chosen sections or smooth families, any collision of the centre
sections, and any failure of the omitted parameter divisor to
be finite etale. Include also $2$, the fixed coefficient prime
$\ell$, and all prime divisors of $|\operatorname{GL}_{d_N}
(\mathbf F_\ell)|$. There are finitely many such primes.
A stable integral lattice in $\mathcal W_N$ has compact image
whose finite congruence quotients are all prime to every
remaining residue characteristic $p$. The prime-to-$p$
specialization isomorphism of SGA~1, XIII, 2.9--2.12, on
the relative $\mathbf P^1$ minus its finite etale boundary
identifies the geometric images in characteristic zero and
at every remaining place $\mathfrak p$. This is the same
integral-image argument as in A13 and gives full orthogonal
monodromy at every such place. It does not infer
characteristic-$p$ monodromy from the Hodge number alone.

Infinitely many rational primes $p\equiv3\pmod4$ remain.
At every place of $K$ above one of them, the reduction of $E$
is supersingular: its model over $\mathbf F_p$ has Frobenius
polynomial $T^2+p$, and finite residue-field extension preserves
supersingularity. Each $w_i$ is a rational relative algebraic
class fixed by normalized Frobenius. The normalized Frobenius
polynomial on $\mathcal W_N$ is rational, obtained by dividing
the polynomial of the full middle cohomology by that of the
ambient cohomology and by $(T-1)^N$. C44-A12 and A12a now
apply with the nonsingular rational Gram matrix $G_N$ and
the even rank $d_N$. They give a Zariski-dense set of closed
points with no root-of-unity Frobenius eigenvalue on
$\mathcal W_N$. All algebraic classes have finite Frobenius
orbit, so the whole algebraic kernel equals the known span.
The coefficients in this span are rational because they are
$G_N^{-1}$ times the vector of rational intersection numbers.
This proves the direct-sum equality in (A13b.1).

Finally, a selected closed parameter lifts after an unramified
extension of the completed local trait of the parameter model.
All $Y_i$, the divisors $h_j$, the chosen incidence variety,
and the polynomial expression for $w_i$ lift over this trait.
Thus every rational cycle class in (A13b.1) has an algebraic
lift. The assertion concerns these classes and does not claim
that every individual representing cycle deforms. In
particular these examples still do not produce nonliftable
primitive classes.
\end{proof}

\paragraph{Parity and scope.}
The oddness of $N$ is exactly what makes the complementary
orthogonal rank even. For even $N$, an odd-dimensional
orthogonal transformation always has an eigenvalue $1$ or
$-1$, so the root-of-unity-free open used in A12 is empty.
This construction therefore gives no analogous exact-span
conclusion for even $N$. The result is an existence theorem
for dense sets of closed points in explicit incidence models;
it does not specify one numerical parameter, does not settle
the arbitrary-axis sign, and does not assert that the full
variable motive is of abelian type.

\paragraph{Source locators.}
The proofs of the three monodromy inputs and the blowup maps
are those read and located in C44-A13: SGA~7~II, XVII~2.5.2
(p.~217), XVIII~2.2 (pp.~260--261), XIX~3.4 (p.~334);
\emph{Weil II}, 4.4.1--4.4.2s (pp.~227--229); SGA~1,
XIII~2.9--2.12 (revised pp.~289--292). Closed-point selection
is the separately checked C44-A12/A12a, using Entin, Theorem~3,
Section~4 (author arXiv:1805.05454v3, pp.~9--11).

% END work/scan_w/closed_odd_rank_axes.tex

% BEGIN work/axis/nonliftable_surface_centres.tex
% C44_POS /root/axis_independent; GPT-6 Astra Pro; effort not exposed.
\section{Smooth nonliftable centres and the primitive cancellation}
\subsection*{C44-A14. A smooth surface with an intrinsic lifting obstruction}

\textbf{Theorem.}
Let $k=\overline{\mathbf F}_3$ and $E/k$ be the supersingular elliptic
curve $y^2=x^3-x$. There is a smooth connected surface
$i:S\hookrightarrow A=E^4$, defined over a finite subfield of $k$,
with the following properties:
\[
 h^1(S,\mathcal O_S)=4,\qquad
 i^*:H^1(A,\mathbf Q_5)\xrightarrow{\sim}H^1(S,\mathbf Q_5).
                                                               \tag{A14.1}
\]
The intrinsic top cup form
\[
 f_S:\bigwedge^4 H^1(S,\mathbf Q_5)\longrightarrow
                 H^4(S,\mathbf Q_5)\simeq\mathbf Q_5(-2)
                                                               \tag{A14.2}
\]
has zero infinitesimal stabilizer in
$\operatorname{End}(H^1(S,\mathbf Q_5))$. Consequently $S$ admits
no abstract smooth projective lift to characteristic zero, even
after extension of the residue field.

Embed $S'=S\times\{0\}^2$ in $X=E^6$. For every sufficiently high
power of an ample line bundle on $X$, there are genuine Lefschetz
axes $B\subset X$ containing $S'$. For each such axis, the
uncorrected algebraic class
\[
                   \gamma=[S']_B\in H^4(B,\mathbf Q_5(2))
                                                               \tag{A14.3}
\]
cannot lift to an algebraic class under any abstract smooth
projective characteristic-zero lift of $B$. The primitive class
\[
 w=\gamma-\kappa^*(L_h^2)^{-1}[S']_X\in A_B,
 \qquad h=c_1(L),\quad B\in|L|\cap|L|,                         \tag{A14.4}
\]
is nonzero for all sufficiently high powers, by C44-A06.
The argument below does \emph{not} prove that $w$ is nonliftable.
Its known obstruction vanishes exactly under this projection.

\paragraph{1. The tensor input, including the lifting quantifier.}
Use the explicit integral divisor-product cycle
\[
 z=\sum_{a<b<c<d}c_{abcd}F_{ab}F_{cd}\in\operatorname{CH}^2(A)
\]
of C44-A10. The integers $c_{abcd}$ and the integral divisors
$F_{ab}$ are specified there. Its modulo-five cohomology tensor
is $\lambda^2u_0$ in an eight-dimensional exterior algebra, where
$\lambda\in\mathbf F_5^\times$ and the derivative map
\[
 \mathcal D_{u_0}:\operatorname{End}(\mathbf F_5^8)
                     \longrightarrow\bigwedge^4\mathbf F_5^8,
 \qquad T\longmapsto T\cdot u_0
                                                               \tag{A14.5}
\]
has a $64$-by-$64$ minor with determinant $3$ modulo five.
This is the explicit certificate in
\texttt{supersingular\_class\_certificate.json}; its integer
determinant and two independent reconstructions are recorded in
C44-A10. We use the actual residue class of the integral cycle,
not a claim about unspecified higher $5$-adic coordinates.

Every integral degree-four class whose reduction equals
$\operatorname{cl}_5(z)$ therefore has zero infinitesimal
stabilizer over $\mathbf Q_5$. Indeed, lift the selected
modulo-five basis to a $\mathbf Z_5$ basis. Milne,
\emph{Abelian Varieties}, I Theorem 12.1, printed pp.~55--56,
identifies integral cohomology with the exterior algebra and
gives torsion freeness. The same minor of the actual derivative
matrix is a $5$-adic unit. Changing its higher $5$-adic entries
cannot destroy this conclusion.

\paragraph{2. A split ample bundle with the required second Chern class.}
For divisor classes $x,y$ modulo five,
\[
                  xy=3\big((x+y)^2-x^2-y^2\big).
                                                               \tag{A14.6}
\]
The sum of a line bundle $M$ and its inverse has
$c_1(M\oplus M^{-1})=0$ and $c_2(M\oplus M^{-1})=-c_1(M)^2$.
Thus a direct sum of such pairs realizes
$\operatorname{cl}_5(z)$ as its second Chern class. For complete
explicitness, for each summand $cF_{ab}F_{cd}$ put
$x=c_1(\mathcal O(F_{ab}))$ and $y=c_1(\mathcal O(F_{cd}))$.
Take the following numbers of inverse pairs, using the least
representatives in $\{0,1,2,3,4\}$:
\[
 \begin{array}{c|c}
 \text{line bundle in the inverse pair}&\text{number of pairs}\\ \hline
 \mathcal O(F_{ab}+F_{cd})&-3c\bmod5\\
 \mathcal O(F_{ab})&3c\bmod5\\
 \mathcal O(F_{cd})&3c\bmod5 .
 \end{array}
                                                               \tag{A14.7}
\]
Each pair has zero first Chern class, so different pairs have no
cross term in the second Chern class. Equations (A14.6)--(A14.7)
give the asserted residue class.

Fix an ample line bundle $H$ on $A$. Replace each line summand
$M_j$ by
\[
                        L_j=M_j\otimes H^{\otimes5n_j}
\]
for a sufficiently large positive integer $n_j$. Every $L_j$
can then be chosen very ample. These changes preserve all Chern
classes modulo five. Write
\[
 \mathcal E=\bigoplus_{j=1}^R L_j,\qquad
 c_2(\mathcal E)\equiv\operatorname{cl}_5(z)\pmod5 .             \tag{A14.8}
\]
Here $R\ge2$ is a finite even integer. Each summand is very ample,
so its sections prescribe arbitrary first jets at any one point
of the smooth variety $A$.

\paragraph{3. A smooth determinantal surface, without an unproved
transversality assumption.}
Choose a general homomorphism
\[
                         \psi:\mathcal O_A^{R-1}\longrightarrow
                                      \mathcal E .
                                                               \tag{A14.9}
\]
All its entries vary independently in the corresponding spaces
of global sections. At a fixed point of $A$, evaluation and
first-jet evaluation are surjective. The space of $R$-by-$(R-1)$
matrices of rank at most $R-3$ has codimension
$3\cdot2=6>4$. The incidence dimension estimate consequently
excludes this rank locus for a general $\psi$.

The matrix stratum of rank exactly $R-2$ is smooth of codimension
two. At a point of that stratum the normal space is
$\operatorname{Hom}(\ker\psi,\operatorname{coker}\psi)$, of
dimension two. The map from the four-dimensional tangent space
of $A$ to this normal space can be prescribed arbitrarily by
the first jets of the entries of $\psi$. Failure of its
surjectivity is the rank-at-most-one locus for a $2$-by-$4$
matrix, which has codimension three. Including the two
conditions on the value gives a bad incidence of codimension
five. Varying the point of the fourfold still leaves a proper
bad locus in the parameter space. Together with the preceding
rank estimate, this proves that a general $\psi$ has a smooth
rank-drop locus
\[
                    S=\{x:\operatorname{rank}\psi_x=R-2\}
                                                               \tag{A14.10}
\]
of pure dimension two, if nonempty. These dimension counts use
first jets and hold in characteristic three without generic
smoothness for a possibly inseparable map.

The locus is nonempty. Otherwise $\psi$ is everywhere of rank
$R-1$ and its cokernel is a line bundle, whence
$c_2(\mathcal E)=0$. This contradicts (A14.8), since the certified
modulo-five tensor is nonzero.

For clarity, its cycle class can be computed directly. On
$A\times\mathbf P^{R-2}$ the incidence
\[
                    \{(x,[v]):\psi_x(v)=0\}
                                                               \tag{A14.11}
\]
is the zero locus of a section of
$\operatorname{pr}_A^*\mathcal E\otimes
 \operatorname{pr}_{\mathbf P}^*\mathcal O(1)$.
There is no deeper rank stratum, so (A14.11) is isomorphic to
$S$: the kernel line is unique and varies algebraically.
It is a regular zero locus of codimension $R$. If
$\xi=c_1(\mathcal O(1))$, pushforward of its top Chern class
takes the coefficient of $\xi^{R-2}$ in
$\prod_{j=1}^R(c_1(L_j)+\xi)$. It follows that
\[
                    [S]_A=c_2(\mathcal E).
                                                               \tag{A14.12}
\]
This proves the class formula directly, without relying on a
generality assertion in a degeneracy-locus formula.

\paragraph{4. Connectedness and coherent irregularity.}
The maximal-minor complex gives an exact sequence
\[
 0\longrightarrow
       \big((\det\mathcal E)^{-1}\big)^{\oplus(R-1)}
   \longrightarrow \mathcal E\otimes(\det\mathcal E)^{-1}
   \longrightarrow I_S\longrightarrow0.                      \tag{A14.13}
\]
One can verify exactness locally here. Where an $(R-2)$-minor
is invertible, row and column operations reduce $\psi$ to an
identity block and a single two-entry column $(a,b)$.
The rank-drop ideal is the regular sequence $(a,b)$, and
(A14.13) reduces to its Koszul resolution plus split acyclic
summands. Away from $S$ a maximal minor is invertible and the
complex is split exact. This covers $A$.

Both left terms are direct sums of inverses of ample line
bundles: the summands of
$\mathcal E\otimes(\det\mathcal E)^{-1}$ are
$(\bigotimes_{k\ne j}L_k)^{-1}$. On an abelian fourfold an
anti-ample line bundle has zero cohomology in every degree
strictly less than four. This follows from Serre duality,
$\omega_A\simeq\mathcal O_A$, and ample vanishing; see Milne,
loc.\ cit., I Theorem 11.1(c), printed p.~54.
The long exact sequence of (A14.13) gives
\[
          H^0(A,I_S)=H^1(A,I_S)=H^2(A,I_S)=0.
                                                               \tag{A14.14}
\]
The sequence $0\to I_S\to\mathcal O_A\to\mathcal O_S\to0$
therefore gives
\[
             H^0(S,\mathcal O_S)=k,\qquad
             H^1(S,\mathcal O_S)\simeq H^1(A,\mathcal O_A).
                                                               \tag{A14.15}
\]
Thus $S$ is connected and $h^1(\mathcal O_S)=4$.
As a smooth connected variety over an algebraically closed
field, it is irreducible.

\paragraph{5. The intrinsic cup form and the necessary
\'etale comparison.}
Set $V=H^1(A,\mathbf Q_5)$, of dimension eight, and
$u=[S]_A\in\bigwedge^4V$, suppressing Tate twists after choosing
a nonzero generator. Equations (A14.8) and (A14.12), together
with the unit-minor certificate, imply
\[
                         \ker\mathcal D_u=0.
                                                               \tag{A14.16}
\]
Let $\operatorname{vol}:\bigwedge^8V\to\mathbf Q_5$ be the
trace pairing, and define
$f_u(\alpha)=\operatorname{vol}(u\wedge\alpha)$ for
$\alpha\in\bigwedge^4V$.
This covariant four-form also has zero infinitesimal stabilizer.
Indeed the derivative of the contragredient action is
\[
                  T\cdot f_u=f_{\,T\cdot u-(\operatorname{tr}T)u}.
                                                               \tag{A14.17}
\]
If it is zero, nondegeneracy of the exterior trace pairing gives
$T\cdot u=(\operatorname{tr}T)u$. Scalar multiplication by the
identity acts on $\bigwedge^4V$ by four. Hence
$T-(\operatorname{tr}T)I/4$ annihilates $u$ and is zero by
(A14.16). Taking traces in dimension eight gives
$\operatorname{tr}T=2\operatorname{tr}T$, and therefore $T=0$.

In particular the contraction radical of $f_u$ is zero.
Otherwise a nonzero endomorphism whose image is a nonzero line
in this radical would annihilate $f_u$.
Projection formula identifies $f_u$ with
\[
     (v_1,\ldots,v_4)\longmapsto
                 \int_S i^*v_1\,i^*v_2\,i^*v_3\,i^*v_4 .
                                                               \tag{A14.18}
\]
Every element of $\ker i^*:H^1(A,\mathbf Q_5)\to H^1(S,\mathbf Q_5)$
would thus belong to this radical. Consequently $i^*$ is injective.

Coherent irregularity alone would not justify an \'{e}tale
isomorphism. The additional required bound is
\[
 \dim H^1(S,\mathbf Q_5)
       =2\dim\operatorname{Pic}^0_{S/k,\mathrm{red}}
       \le2\dim T_0\operatorname{Pic}_{S/k}
       =2h^1(S,\mathcal O_S)=8.                              \tag{A14.19}
\]
For the equality, the Kummer sequence identifies the rational
$5$-adic Tate module of the Picard variety with
$H^1(S,\mathbf Q_5(1))$; the N\'eron--Severi torsion disappears
after tensoring with $\mathbf Q_5$. This is explicitly stated
in SGA~1, Expos\'e XI, Corollary 6.6 and the paragraph following
it, revised p.~235, original pp.~305--306.
The tangent-space identity is recalled in XI, Remarks 6.11,
revised p.~237, original p.~308. Dimension of a group scheme
at its identity is at most the dimension of its tangent space.
The injection and (A14.19) prove (A14.1), and (A14.18) then
proves that the intrinsic form (A14.2) has zero infinitesimal
stabilizer.

\paragraph{6. Abstract nonliftability of the surface.}
Suppose $S$ were the special fibre of a smooth projective
mixed-characteristic lift. Smooth proper base change preserves
the rational cohomology and cup tensor. Spread the generic fibre
to a finitely generated characteristic-zero field and embed that
field into $\mathbf C$. Betti--\'etale comparison identifies
the complex fibre's rational top cup form after scalar extension
with (A14.2). Its infinitesimal stabilizer therefore remains
zero. This is a rank condition on a rational derivative matrix,
so it can be checked over $\mathbf Q_5$ or over $\mathbf R$.

On the real $H^1$ of a smooth projective complex surface, the
Hodge circle has a nonzero infinitesimal generator $J$ with
$J^2=-I$. The top cohomology $H^4$ has type $(2,2)$, so the
circle acts trivially on it. Compatibility of cup product with
Hodge structures implies that $J$ annihilates the top fourfold
cup form. This contradicts its zero infinitesimal stabilizer.
No embedded lift, preservation of the selected abelian
fourfold, or Hodge conjecture is assumed.

All choices involved finitely many equations over $k$.
The nonempty parameter opens used above have $k$-points; each
chosen point and all defining data lie in a finite extension of
$\mathbf F_3$. The surface therefore descends to a finite
field. This proves the first part of the theorem.

\paragraph{7. The class on an axis and the exact cancellation.}
Embed $S'=S\times\{0\}^2$ in $X=E^6$. The imposed-centre
construction in C44-A06 gives, for every sufficiently high
power of an ample polarization, a genuine Lefschetz pencil
whose smooth axis $B$ contains $S'$. Weak Lefschetz identifies
$H^1(B,\mathbf Q_5)$ with
\[
       H^1(X,\mathbf Q_5)=V\oplus U,\qquad\dim V=8,\quad\dim U=4.
\]
For the class $\gamma=[S']_B$ consider the intrinsic form
\[
 f_\gamma(v_1,\ldots,v_4)=\int_B\gamma v_1v_2v_3v_4.
                                                               \tag{A14.20}
\]
It is the pullback of $f_S$ from $V$, because restriction to
$S'$ kills the last two elliptic factors. Its radical is
exactly $U$, and its induced form on the quotient by the radical
has zero infinitesimal stabilizer.

If $\gamma$ lifted to an algebraic class in an arbitrary
abstract smooth projective lift of $B$, the generic rational
Betti class would have type $(2,2)$. Cup product and top trace
would make the transported form (A14.20) Hodge-circle invariant.
Its radical is intrinsic and hence stable under that circle.
The quotient of $H^1$ by the radical inherits a weight-one
Hodge structure. Its real infinitesimal generator still
satisfies $J^2=-I$ and is nonzero, while annihilating the induced
four-form. This contradicts the certified zero stabilizer.
Thus the \emph{cohomology class} $\gamma$ is nonliftable;
the claim is stronger than failure to lift its chosen support.

Hard Lefschetz on $X$ and algebraicity of its inverse give
(A14.4), a rational algebraic element of the exact kernel
$A_B=\ker\kappa_*$. C44-A06 makes it nonzero and positive for
all sufficiently high powers. However, for every
$w\in\ker\kappa_*$ and every four $v_i\in H^1(X)$, projection
formula gives
\[
       \int_B w\,\kappa^*v_1\kappa^*v_2\kappa^*v_3\kappa^*v_4
          =\int_X\kappa_*w\,v_1v_2v_3v_4=0.                  \tag{A14.21}
\]
Every $H^1(B)$ class comes from $X$. Consequently the intrinsic
four-form of the projected class $w$ is identically zero.
The entire obstruction (A14.20) is carried by the subtracted
ambient class.

The same cancellation occurs for differences of translates.
Take two distinct points of the last $E^2$ and the disjoint
surfaces $S'_1,S'_2$. Their ambient cycle classes agree because
translation acts trivially on cohomology. A sufficiently ample
axis contains both, and
$[S'_1]_B-[S'_2]_B\in A_B$ is nonzero by the positive Gram
matrix of C44-A06. Both individual supports and both
uncorrected classes are nonliftable, but their four-forms agree
and cancel. Nonliftability of the supports does not exclude a
different characteristic-zero cycle representing the projected
or difference cohomology class. No such exclusion is established
here. \hfill$\square$

\paragraph{Sources and scope.}
The explicit tensor and arithmetic certificate are C44-A10,
with files \texttt{axis\_supersingular\_nonliftable.tex} and
\texttt{supersingular\_class\_certificate.json}.
Primary sources used above are:
\[
 \begin{gathered}
 \texttt{https://jmilne.org/math/CourseNotes/AV.pdf},\\
 \texttt{https://arxiv.org/pdf/math/0206203},\\
 \texttt{https://www.jmilne.org/math/Documents/Deligne82.pdf}.
 \end{gathered}
\]
The Milne locators are I Theorems 11.1(c) and 12.1,
printed pp.~54--56. The SGA~1 locators are XI 6.6, revised
p.~235/original pp.~305--306, and XI 6.11, revised p.~237/original
p.~308. The Hodge-circle description and the necessary Hodge
type of algebraic classes are Deligne's revised notes,
Introduction p.~2 and \S3 p.~23. Smooth proper base change is
SGA~4 XVI Theorem 2.2. The rank-stratum, Koszul and incidence
top-Chern arguments have been given in full above.

This construction certifies a smooth abstractly nonliftable
centre and an individually nonliftable class on a liftable
axis. It does not produce a certified nonliftable class in
$A_B$, and it does not settle the universal QH9-axis sign.

% END work/axis/nonliftable_surface_centres.tex

% BEGIN work/scan_w/axis_one_tate_direction.tex
% C44_POS /root/scan_w. Root's determinant idea; complete proof and all-ell repair.
\section{One remaining potential Tate direction determines the axis sign}
\label{sec:c44-a15}

\begin{theorem}[C44-A15: a positive span and one potential Tate direction]
Let $r\geq1$ and let $X/\mathbf F_q$ be an abelian variety of dimension $2r+2$,
let $L$ be an ample line bundle with $h=c_1(L)$, and let
$\kappa:B\hookrightarrow X$ be a smooth complete intersection
of two members of $|L|$. In particular this applies to
the smooth axes in QH9-axis. For each prime $\ell\ne p$, set
\[
 V_\ell=\ker\bigl(\kappa_*:H^{2r}(B_{\overline{\mathbf F}_q},
             \mathbf Q_\ell(r))\longrightarrow
        H^{2r+4}(X_{\overline{\mathbf F}_q},\mathbf Q_\ell(r+2))\bigr),
 \qquad q_B(x,y)=(-1)^r\int_Bxy.
 \tag{A15.1}
\]
Fix an initial coefficient prime $\ell_0\ne p$ for the definition
of $\Alg^r(B)$ as the image of rational cycles, and put
$A_B=\Alg^r(B)\cap V_{\ell_0}$.
Suppose a finite-dimensional rational subspace $K\subset A_B$
has a positive definite Gram matrix $G$ for $q_B$; the case
$K=0$, with $\det G=1$, is allowed.

There is a common algebraic projection of cycle representatives
onto $V_\ell$ for every $\ell\ne p$. Use it on a rational basis
of $K$, and extend the finite ground field so that the resulting
cycles are all defined over it. Their span $K_\ell\subset V_\ell$
has the same Gram matrix $G$ in every realization. Write
\[
       W_\ell=K_\ell^\perp\cap V_\ell,
 \qquad
       \Phi_W(T)=\det(T-F\mid W_\ell)\in\mathbf Q[T],
 \tag{A15.2}
\]
where $F$ is normalized Frobenius on the Tate-twisted cohomology.
The polynomial is independent of $\ell$. Let $\tau$ be the total
\emph{algebraic multiplicity} of its roots that are roots of
unity. If
\[
                              \tau\leq1,
 \tag{A15.3}
\]
then $q_B$ is positive definite on all of $A_B$.
Moreover $A_B=K$ or $A_B=K\perp\mathbf Qz$ for a class $z$ of
strictly positive square. When $\tau=0$, only $A_B=K$ occurs.
The given cycles spanning $K$ need not lift to characteristic zero.
No Tate conjecture, Frobenius semisimplicity, standard conjecture
$D(B)$, or unrefereed 032 result is assumed.
\end{theorem}

\begin{proof}
All determinants below are ordinary Gram determinants modulo
squares, with no factor $(-1)^{d(d-1)/2}$.

\emph{One algebraic projector for all coefficient primes.}
The self-intersection formula and hard Lefschetz give
\[
 \kappa_*\kappa^*=L_h^2:
       H^{2r}(X)\xrightarrow{\sim}H^{2r+4}(X)(2).
 \tag{A15.4}
\]
There is a rational algebraic correspondence $\Gamma$ on $X$
which realizes $(L_h^2)^{-1}$ simultaneously for every $\ell\ne p$.
Indeed replace $h$ by the symmetric rational divisor
$(h+[-1]^*h)/2$, which has the same cohomology class, and apply
K\"unnemann's hard Lefschetz isomorphism for Chow motives with
$g=2r+2$ and $i=2r$. Its inverse, composed with the indicated
Chow--K\"unneth projectors, is one correspondence $\Gamma$.
The source is \emph{Inventiones Mathematicae} 113 (1993),
Theorem~5.2, printed p.~96, with the primitive decomposition
in Theorem~5.1 on the same page. Only the absolute theorem over
the finite ground field is needed here. No application of its
relative formulation to a mixed-characteristic base is being made.

Consequently the correspondence
\[
                 P=\Delta_B-\kappa^*\Gamma\kappa_*
 \tag{A15.5}
\]
acts on middle cohomology as the orthogonal projector onto
$V_\ell$ for every $\ell\ne p$. These statements concern its
cohomological action; no idempotence claim on all Chow groups
is needed. The ambient summand is nondegenerate, since its
pairing is $(-1)^r\int_X ab h^2$, nondegenerate by hard
Lefschetz. Thus $V_\ell$ is itself nondegenerate.

Let $C_1,\ldots,C_k$ be cycle representatives for a rational
basis of $K$. Replace them by $P_*C_1,\ldots,P_*C_k$ and retain
the names. The initial $\ell_0$-classes are unchanged. Their
intersection numbers, hence their rational Gram matrix $G$,
are unchanged as well. Now their classes lie in $V_\ell$ in
\emph{every} realization, and $\det G\ne0$ makes them independent
in every realization. This explicit replacement is important:
vanishing of $\kappa_*C_i$ in one realization alone would not
justify vanishing in all realizations without another argument.
After a finite ground-field extension all these actual rational
cycles and correspondences are defined over that field, so
normalized Frobenius fixes $K_\ell$ pointwise.

The characteristic polynomial of $V_\ell$ is the quotient of
the normalized middle characteristic polynomial of $B$ by that
of $H^{2r}(X)$; the latter is the ambient summand. These
polynomials are rational, independent of $\ell$, and have all
complex roots on the unit circle, by Deligne, \emph{Weil I},
Theorem~1.6, printed p.~276. Dividing further by $(T-1)^k$
gives $\Phi_W$. Passing to a finite field extension raises
eigenvalues to a common positive power and does not change
their total root-of-unity multiplicity.

\emph{A positive rational determinant for the full variable space.}
Use the embedded polarized characteristic-zero lift of $B\subset X$
from C44-A05. Choose a complex embedding of a field of definition
of its generic fibre. Let $V_{\rm B}$ be the rational Betti
kernel of its Gysin map and let $q_{\rm B}$ be the signed
intersection form in (A15.1). Smooth proper base change and
Betti--etale comparison identify
\[
           (V_\ell,q_B)\simeq
              (V_{\rm B},q_{\rm B})\otimes_{\mathbf Q}\mathbf Q_\ell
                \quad\hbox{for every }\ell\ne p.
 \tag{A15.6}
\]
One lift is used for all these comparisons. The kernel and its
pairing are geometric, so this step does not require the cycles
$C_i$ or the special-fibre correspondence $\Gamma$ to lift.

The Betti kernel is primitive for $h|_B$. Indeed
$\kappa_*(h|_B\,v)=h\kappa_*v=0$, and Gysin on
$H^{2r+2}(B)$ is injective, being dual to the weak Lefschetz
isomorphism on $H^{2r-2}$. Hence $h|_B\,v=0$.
The Hodge--Riemann relations make $q_{\rm B}$ positive on
type $(r,r)$ and give opposite signs on successive pairs of
off-diagonal Hodge types. Every negative real block comes
from a conjugate pair of distinct Hodge types and has even
dimension. The negative index is therefore even. It follows
that, in any rational Betti basis,
\[
                       D:=\det(q_{\rm B})>0.
 \tag{A15.7}
\]
In particular $D$ represents the determinant squareclass of
$q_B|_{V_\ell}$ simultaneously in every $\mathbf Q_\ell$.
This only uses the sign parity on the full variable space;
it does not declare $q_{\rm B}$ positive definite there.

\emph{Removing the known span by rational intersection numbers.}
For any cycle representing $x\in A_B$, first replace it by
its image under $P$, then subtract its orthogonal projection
onto the $C_i$. The coefficients of that projection are
\[
     G^{-1}\bigl(q_B(P_*x,C_1),\ldots,q_B(P_*x,C_k)\bigr)^{\rm t},
 \tag{A15.8}
\]
where $x$ denotes the chosen cycle in this formula. They are
rational intersection numbers. The resulting rational cycle
$Z$ has its class in $W_\ell$ for every $\ell\ne p$ and
represents exactly the residual component of $x$ at $\ell_0$.
Thus the same representative can be used in every realization.

Every algebraic cycle over $\overline{\mathbf F}_q$ is defined
over a finite extension. Its normalized Frobenius orbit is
therefore finite. Its class in $W_\ell$ is annihilated by
$F^e-1$ for some $e>0$, so it belongs to the sum of the
root-of-unity eigenspaces. If $\tau=0$, every residual class
vanishes, giving $A_B=K$. Assume henceforth that $\tau=1$ and
that some residual class $z=\operatorname{cl}_{\ell_0}(Z)$ is
nonzero; otherwise the same conclusion holds.

\emph{The simple potential Tate line and its complement.}
Because $\Phi_W$ is rational and has just one root of unity,
counted with multiplicity, that root is $1$ or $-1$ and is
simple. Replace $F$ by $F^2$ and call the resulting isometry
$S$. In every realization its eigenvalue $1$ is simple and
it has no other root-of-unity eigenvalue. Its one-dimensional
$1$-eigenspace $T_\ell$ is nondegenerate: the other generalized
eigenspaces are orthogonal to it, so degeneracy on this line
would contradict nondegeneracy of $W_\ell$. We have
\[
              W_\ell=T_\ell\perp U_\ell,
 \tag{A15.9}
\]
with $S|_{U_\ell}$ having no eigenvalues $1$ or $-1$.
In particular $\dim U_\ell$ is even, since eigenvalues of an
orthogonal transformation occur in reciprocal pairs away
from $1$ and $-1$.

At $\ell_0$, the class $z$ spans $T_{\ell_0}$. Therefore
\[
                           a:=q_B(Z,Z)\in\mathbf Q^*
 \tag{A15.10}
\]
is nonzero. Since it is an actual rational intersection number,
the class of this same cycle $Z$ is nonzero in every realization.
By (A15.8) it belongs to $W_\ell$ in each realization, and by
its finite Frobenius orbit it spans $T_\ell$. Thus $a$ is the
Gram determinant of $T_\ell$ in the basis supplied by $Z$
for every $\ell\ne p$. This proves the required simultaneous
nonvanishing and kernel membership without $D(B)$.

\emph{The non-Tate determinant from a skew-adjoint operator.}
Let
\[
  \psi(T)=\det(T-S\mid U_\ell)\in\mathbf Q[T],\qquad
  A=S-S^{-1}\quad\hbox{on }U_\ell.
\]
The polynomial $\psi$ is independent of $\ell$: the
characteristic polynomial after squaring is determined by
$\Phi_W$, and one divides out its simple factor $T-1$.
The operator $A$ is invertible and skew-adjoint for $q_B$.
Consequently $b(u,v)=q_B(Au,v)$ is a nondegenerate alternating
form. Its matrix determinant is a Pfaffian square, giving
\[
                   \det(q_B|_{U_\ell})\equiv\det A
                      \pmod{(\mathbf Q_\ell^*)^2}.
 \tag{A15.11}
\]
Indeed $\det b=\det A\,\det(q_B|_{U_\ell})$ is a square,
and an element and its inverse have the same squareclass.
Reciprocal eigenvalue pairing gives $\det(S|_{U_\ell})=1$.
As the dimension is even, it follows that
\[
       \delta:=\det A=\psi(1)\psi(-1)\in\mathbf Q_{>0}.
 \tag{A15.12}
\]
To verify the strict sign, all roots of $\psi$ have complex
absolute value one and avoid $1,-1$. Pairing each root with
its complex conjugate gives positive factors
$|1-\lambda|^2$ and $|1+\lambda|^2$ in these two evaluations.
If $U_\ell=0$, both the determinant and the empty product
are $1$. The calculation uses characteristic polynomials
only, and requires no semisimplicity on $U_\ell$.

\emph{The missing-$p$ squareclass argument fixes the real sign.}
Combining (A15.6), the orthogonal decomposition by $K_\ell$,
(A15.9), and (A15.11), for every prime $\ell\ne p$ we obtain
\[
            a\,\det G\,\delta/D\in(\mathbf Q_\ell^*)^2.
 \tag{A15.13}
\]
The map
\[
 \mathbf Q^*/(\mathbf Q^*)^2\longrightarrow
        \prod_{\ell\ne p}\mathbf Q_\ell^*/(\mathbf Q_\ell^*)^2
 \tag{A15.14}
\]
is injective. One proof is that a nontrivial quadratic extension
of $\mathbf Q$ has an inert prime outside any fixed finite
set. The exact squareclass implication is also explicitly
proved and cited in Ancona--Marmora, \emph{Commentarii
Mathematici Helvetici} 100 (2025), proof of Corollary~4.4,
printed p.~629, referring to Ireland--Rosen, second edition,
\S5.2, Theorem~3. Applying it to (A15.13) gives
\[
                     a=\frac{D}{\det G\,\delta}\,c^2
                       \quad\hbox{for some }c\in\mathbf Q^*.
 \tag{A15.15}
\]
All three rational factors in the denominator and numerator
on the right are positive. Hence $a>0$.

Every nonzero algebraic residual direction has just been
shown positive. If one such class $z$ exists, every other
residual class is its rational multiple: the multiple is
$q_B(z',z)/q_B(z,z)\in\mathbf Q$, and equality follows in
the one-dimensional nondegenerate $T_{\ell_0}$. Thus
$A_B=K\perp\mathbf Qz$; otherwise $A_B=K$. Positivity
on the whole space follows in either case.
\end{proof}

\begin{corollary}[C44-A15a: the simple rank-one gate]
For any fixed axis in the theorem, if the normalized variable
Frobenius polynomial has at most one root of unity counted
with algebraic multiplicity, its entire axis form has the
required positive sign unconditionally. No algebraic
generator of the possible Tate line needs to be exhibited.
If that line contains a nonzero algebraic class, the argument
proves its positive square; it does not assert its algebraicity.
\end{corollary}

\begin{proof}
Take $K=0$ in C44-A15. The convention $\det G=1$ gives
exactly the asserted conclusion.
\end{proof}

\paragraph{A sharper remaining gate in the abelian sixfold case.}
For $\dim X=6$ and $\dim B=4$, Gauss--Bonnet, weak Lefschetz
and Poincare duality give
\[
        \dim V=5\int_Xh^6-165=3600\chi(X,\mathcal O_X(h))-165,
 \tag{A15.16}
\]
which is odd. All non-root-of-unity eigenvalues of its
orthogonal Frobenius occur in reciprocal pairs. Thus the
total root-of-unity multiplicity is odd. C44-A15a settles
the unavoidable multiplicity-one case: the fixed-axis sign
question not covered by this test now requires multiplicity
at least three. More generally, after a known positive span
of dimension $k$ has been removed, the residual multiplicity
has parity $\dim V-k$. An even-dimensional known span in a
sixfold axis can therefore be completed by this theorem
when its residual multiplicity is one. This sharpens a
cohomological sufficient test; it does not prove that every
axis satisfies the multiplicity bound.

\paragraph{Programme locator and conditional height consequence.}
This result concerns cote~44, Theorem~5, scan~004, its primitive
margin on scan~008, and the exact H9-T7 form on QH9-axis.
It sharpens route~(c) from exclusion of every potential Tate
direction to a sufficient sign test allowing one such residual
direction after a known positive span. The axis sign proved
here is unconditional. By the decomposition in C44-A01, the
corresponding full-height and $I(B)$ consequences still use
the abelian positivity input conditional on the unrefereed
032/Milne2026 results. Neither those conditional consequences
nor this sufficient test settle the arbitrary-axis case.

\paragraph{Primary source checks and scope.}
K\"unnemann's actual scanned Theorems~5.1--5.2 on printed
p.~96 and their surrounding construction were read at
\url{https://gdz.sub.uni-goettingen.de/download/pdf/PPN356556735_0113/LOG_0011.pdf}.
The base-scheme convention on p.~86 was also checked: the
quoted relative formulation is over a field, and only its
absolute special-fibre case is used above. Deligne's actual
Theorem~1.6, printed p.~276, was read at
\url{https://www.numdam.org/item/PMIHES_1974__43__273_0.pdf}.
The complete proof of Ancona--Marmora Corollary~4.4, printed
p.~629, and the signature discussion in Proposition~4.5,
p.~630, were read at
\url{https://ems.press/content/serial-article-files/51595?nt=1}.
The theorem uses only the elementary squareclass consequence
of that paper, not its CM or supersingularity hypotheses.
The geometric lift is C44-A05; no lifting of an individual
test cycle or of the known positive span is used. A bound on
the dimension of algebraic classes alone would not replace
(A15.3), which concerns the whole potential Tate part of the
Frobenius characteristic polynomial.

% END work/scan_w/axis_one_tate_direction.tex

% BEGIN work/axis/constrained_moment_monodromy.tex
% C44_POS /root/axis_independent; GPT-6 Astra Pro; effort not exposed.
\section{Constrained axes: fourth moments without lifting the centres}
\label{sec:c44-a18}

\begin{theorem}[C44-A18: orthogonal monodromy with prescribed centres]
Let $k$ be a finite field of odd characteristic, let $X/k$ be an
abelian variety of dimension $2r+2$, $r\geq2$, and let $H$ be a
very ample line bundle with $h=c_1(H)$. Fix pairwise disjoint
smooth geometrically connected closed subvarieties
$Z_1,\ldots,Z_s\subset X$, all of dimension $r$, and put
$Z=\coprod_i Z_i$; the empty collection is allowed.
For all sufficiently large integers $t$, after a finite extension
of $k$ there is a smooth member
$D_0\in|I_Z\otimes H^{\otimes t}|$ such that the following holds.

There is a nonempty open subset $\Omega$ of the affine space
\[
             E_t=H^0(X,I_Z\otimes H^{\otimes t})
\]
whose points $a$ give smooth complete intersections
$B_a=D_0\cap(a=0)$, and for which the original pencil spanned
by $D_0$ and $(a=0)$ is a genuine Lefschetz pencil on $X$.
For $\ell\ne\operatorname{char}k$, let
\[
 \begin{split}
 \mathcal V_a&=\ker\bigl(\kappa_{a*}:H^{2r}(B_{\bar a},
          \mathbf Q_\ell(r))\longrightarrow
                   H^{2r+4}(X_{\bar k},\mathbf Q_\ell(r+2))\bigr),\\
 b_i&=(L_h^2)^{-1}[Z_i]_X,\qquad
 w_i=[Z_i]_{B_a}-t^{-2}\kappa_a^*b_i,\\
 \mathcal K&=\sum_i\mathbf Q_\ell w_i,
 \qquad \mathcal W=\mathcal K^\perp\cap\mathcal V,
 \qquad Q(x,y)=(-1)^r\int_{B_a}xy.
 \end{split}                                                    \tag{A18.1}
\]
The classes $w_i$ are rational algebraic relative classes, their
constant signed Gram matrix is positive definite, and
$\mathcal V=\mathcal K\perp\mathcal W$ is a splitting of lisse
sheaves. The geometric monodromy closure on $\mathcal W$ is the
full orthogonal group, provided $\operatorname{rank}\mathcal W>8$;
this rank condition also holds for all sufficiently large $t$.
The varieties $Z_i$ and their cycle classes are not assumed to
lift to characteristic zero. No 032 input is used.
\end{theorem}

\begin{proof}
We distinguish the integer $s$ counting the centres, a section
$a\in E_t$, and the cardinality $Q_0$ of a finite extension field.

\emph{The family and the nondegenerate known summand.}
C44-A06 supplies a nonempty open set of pairs of sections with
smooth common zero locus containing $Z$, whose original pencils
are Lefschetz. Its imposed-centre jet proof gives this for every
sufficiently large $t$. Projection of this open set onto the first
section space contains a nonempty open. Choose $D_0$ in this open,
after extending the finite field, and take the corresponding
nonempty open $\Omega\subset E_t$ of second sections. Thus $D_0$
is smooth, and the assertion about the original pencil is part of
the chosen open condition.

The algebraic inverse Lefschetz correspondence on the abelian
variety supplies the classes $b_i$ simultaneously for all
coefficient primes. Self-intersection and hard Lefschetz identify
the ambient image in the middle cohomology of $B_a$ as a
nondegenerate constant summand. Its orthogonal complement is
$\mathcal V$. The cycle $Z_i\times\Omega$ inside the universal
axis gives a relative cycle class; the ambient correction in
(A18.1) is also relative. Hence every $w_i$ is a section of
$\mathcal V$, fixed by normalized arithmetic Frobenius.

The exact Gram formula of C44-A06 is constant in this family and
positive definite for large $t$. In particular the $w_i$ are
independent and $\mathcal K$ is nondegenerate. It follows that
$\mathcal W$ is a lisse nondegenerate orthogonal direct summand.
Nothing here assumes that the stated classes exhaust the fixed
part. That will follow from the monodromy calculation.

In the calculation below we use the untwisted sheaves, denoted by
$\mathcal V^{\rm un}$ and $\mathcal W^{\rm un}$, of weight $2r$.
The $s$ known lines in $\mathcal K^{\rm un}$ have eigenvalue
$Q_0^r$ over a field of cardinality $Q_0$.

\emph{Four-point evaluation outside the centres.}
Choose an integer $c$ such that $I_Z\otimes H^{\otimes c}$ is
globally generated. For $t\geq c+3$, evaluation from $E_t$ is
surjective at any set of at most four distinct geometric points
of $X\setminus Z$. Here is a direct argument, also over a finite
extension $k'/k$ with $Q_0=\#k'>4$. At each such point the
evaluation of $H^0(I_ZH^c)\otimes k'$ is nonzero. The union of at
most four proper linear hyperplanes does not cover this vector
space, so there is a section $e$ nonzero at all these points.
For each selected point $x_i$, very ampleness gives a hyperplane
section through each $x_j$, $j\ne i$, but not through $x_i$.
Their product, padded to degree $t-c$ by a hyperplane section
nonzero at all the points, vanishes at the other points and not
at $x_i$. Multiplying these products by $e$ proves simultaneous
surjectivity. This also proves the assertion geometrically, and
surjectivity descends to $k'$.

Set
\[
                     U=D_0\setminus Z.
\]
It is smooth and geometrically connected of dimension $2r+1$.
The same four-point evaluation statement holds for points of $U$,
because it was proved using sections on $X$. No completeness
claim about the induced linear system on $D_0$ is required.

\emph{A trace function with the centre contribution removed.}
For every $a\in E_t(k')$, including singular or zero sections,
define the rational number
\[
 T(a,k')=\#\bigl(U\cap(a=0)\bigr)(k')-\frac{\#U(k')}{Q_0}.
                                                               \tag{A18.2}
\]
For $a\in\Omega(k')$ we claim the uniform estimate
\[
 T(a,k')=\operatorname{Tr}
       (F_{k',a}\mid\mathcal W^{\rm un}_{\bar a})
                   +O(Q_0^{r-1/2}).                            \tag{A18.3}
\]
The implied constant depends on the fixed geometry, not on $k'$
or $a$ in the smooth open.

Indeed $B_a$ is a smooth ample divisor of the smooth
$(2r+1)$-fold $D_0$. The trace formula, weak Lefschetz in degrees
below $2r$, and its Poincare-dual Gysin statement in degrees above
$2r$, give
\[
 \#B_a(k')-\frac{\#D_0(k')}{Q_0}
       =\operatorname{Tr}(F_{k',a}\mid\mathcal V^{\rm un}_{\bar a})
                       +O(Q_0^{r-1/2}).                       \tag{A18.4}
\]
For clarity, the terms in degrees greater than $2r$ on $B_a$
cancel the corresponding terms of degrees greater than $2r+2$
on $D_0$, divided by $Q_0$. In the middle, the restriction image
of $H^{2r}(D_0)$ has the same Frobenius polynomial as
$H^{2r+2}(D_0)(1)$, by hard Lefschetz; its removal leaves the
variable summand. The remaining terms have weights at most
$2r-1$. Moreover $H^{2r}(X)\simeq H^{2r}(D_0)$ by weak
Lefschetz, so this variable summand is exactly $\mathcal V$ in
(A18.1). The ranks of the low-degree groups are constant and
Deligne's weight bound supplies the stated uniform error.
This is the argument of Katz's Key Lemma 2.2.6, with $D_0$ as
the ambient variety.

Since the $Z_i$ are smooth geometrically connected projective
$r$-folds, the trace formula and the Weil bounds give
\[
                      \#Z(k')=sQ_0^r+O(Q_0^{r-1/2}).           \tag{A18.5}
\]
Subtracting $(1-Q_0^{-1})\#Z(k')$ from (A18.4) gives (A18.2).
On the cohomology side, removing $\mathcal K^{\rm un}$ subtracts
exactly $sQ_0^r$. Since $Q_0^{r-1}=O(Q_0^{r-1/2})$, equations
(A18.4)--(A18.5) prove (A18.3).

\emph{The exact fourth moment.}
Write $M=\#U(k')$, $v=Q_0^{-1}$ and let $a$ be uniformly random
in the finite vector space $E_t(k')$. For $x\in U(k')$, put
$I_x(a)=1$ if $a(x)=0$, and $I_x(a)=0$ otherwise. Four-point
evaluation implies that these variables are independent for
every set of at most four distinct points, and each is a
Bernoulli variable of parameter $v$. Thus
$T(a,k')=\sum_x(I_x(a)-v)$, and elementary expansion gives the
exact identity
\[
 \begin{split}
 \frac{1}{\#E_t(k')}\sum_{a\in E_t(k')} T(a,k')^4
   &=3M(M-1)\sigma^4+M\mu_4,\\
 \sigma^2&=v(1-v),\\
 \mu_4&=v(1-v)^4+(1-v)v^4.
 \end{split}                                                   \tag{A18.6}
\]
Only coincidence patterns $(2,2)$ and $(4)$ survive: a singleton
has expectation zero. As $U$ is a dense open in the smooth
geometrically connected $(2r+1)$-fold $D_0$,
$M=Q_0^{2r+1}+O(Q_0^{2r+1/2})$. Consequently
\[
 Q_0^{-4r}\frac{1}{\#E_t(k')}
       \sum_{a\in E_t(k')}T(a,k')^4
                   =3+O(Q_0^{-1/2}).                         \tag{A18.7}
\]
For example, the second term in (A18.6), after this normalization,
is $O(Q_0^{-2r})$.

Let $G$ be the geometric monodromy closure of $\mathcal W$ and
$W$ its fibre after extension of the coefficient field to an
algebraic closure. The moment is
\[
        M_4(G,W)=\dim\bigl(W^{\otimes2}\otimes
                                  (W^\vee)^{\otimes2}\bigr)^G.
\]
Purity makes $G$ reductive. Katz's Lemma 2.1.4 identifies this
integer with the limsup of the normalized fourth trace moments
over the smooth parameter open. Equation (A18.3) changes those
moments by $o(1)$: both the trace and $T$ on this open are
$O(Q_0^r)$, and their difference is $O(Q_0^{r-1/2})$.
Since fourth powers are nonnegative, adding all the omitted
sections can only increase the sum. Therefore (A18.7) gives
\[
                         M_4(G,W)\leq3.                       \tag{A18.8}
\]
This is precisely the use of Katz's Second Variant Lemma 2.1.6.
No estimate on traces of singular members is used or needed.
Because $G\subset O(W,Q)$ and $\dim W>8$, the three standard
orthogonal contraction invariants are independent. Hence
$M_4(G,W)\geq3$, and equality holds.

\emph{A nonzero reflection in the residual summand.}
For large $t$ the imposed-centre jet argument also produces a
section whose restriction to $D_0$ has exactly one ordinary
quadratic singularity at a chosen point $x\in U$, and is smooth
everywhere else, including along $Z$.
Here one prescribes the value and first derivatives to vanish at
$x$, and chooses a nondegenerate quadratic term on $T_xD_0$;
odd characteristic is used here. The system of sections with
this fixed second jet has arbitrary first jets at a second
point outside $Z\cup\{x\}$ for large $t$. A further singular
point therefore imposes $2r+2$ conditions while varying in
dimension $2r+1$, so a general section has no such point.
Along $Z_i$, its normal first derivative is a section of the
rank-$(r+1)$ bundle
$N^*_{Z_i/D_0}\otimes H^t|_{Z_i}$. This bundle is generated by
the constrained sections even after the prescribed jet at $x$;
its zero condition has codimension $r+1>r=\dim Z_i$, so a
general section has no singular point there either.
The required uniform generation follows as in C44-A06 by
twisting the ideal and using bounded jets; equivalently,
hyperplanes vanishing to order three at $x$ and nonzero at a
second point preserve the fixed jet while generating the
indicated first jets. Thus these dimension counts apply to
actual linear systems.

Choose a transverse one-parameter deformation of this nodal
member, by a section nonzero at $x$. Its total space is regular
at the node, and proper over the local trait; its punctured
generic point can be chosen in $\Omega$, since $\Omega$ is
dense open. SGA 7 II, Expose XV, Theorem 3.4(i),(iii), printed
pp.~195--196, gives a global vanishing cycle $\delta$ with
\[
                       \int\delta^2=2(-1)^r.                 \tag{A18.9}
\]
In particular $\delta\ne0$. The same theorem says that the
quadratic inertia character is surjective when the total space
is regular. An inertia element acts by the Picard--Lefschetz
reflection in $\delta$.

Every ambient class and every $w_i$ is fixed by this local
monodromy. The reflection formula and $\delta\ne0$ then imply
that $\delta$ is orthogonal to both the ambient image and
$\mathcal K$. Thus $\delta\in W$, and this is a nontrivial
reflection on $W$, not merely on the entire middle cohomology.
Katz's Theorem 1.4.6 now applies to (A18.8) and proves
$G=O(W,Q)$.

Finally the rank of the variable cohomology tends to infinity
with $t$: C44-A03b gives
\[
 \operatorname{rank}\mathcal V
  =(2r+1)t^{2r+2}\int_Xh^{2r+2}
       +2\binom{4r+3}{2r-1}-\binom{4r+4}{2r}.
                                                               \tag{A18.10}
\]
Subtracting the fixed integer $s$ proves the remaining rank
assertion. Restriction to any further nonempty parameter open
does not change geometric monodromy, since its fundamental
group surjects onto that of the original smooth open.
\end{proof}

\begin{corollary}[C44-A18a: closed axes through arbitrary smooth centres]
With the hypotheses and sufficiently large $t$ of C44-A18,
there is a Zariski-dense set of closed parameters $a\in\Omega$
for which the signed axis form is positive definite on the
\emph{whole} rational algebraic variable space $A_{B_a}$.
Put $d=\operatorname{rank}\mathcal W$.
If $d$ is even, then
\[
                         A_{B_a}=\sum_i\mathbf Q w_i.
                                                               \tag{A18.11}
\]
If $d$ is odd, the right side is either all of $A_{B_a}$ or has
codimension one, and the additional line has positive square.
No lifting of the centres, Tate conjecture or 032 input is used.
\end{corollary}

\begin{proof}
The normalized Frobenius polynomial of $\mathcal W$ is rational:
it is the middle polynomial of $B_a$ divided by the ambient
polynomial and by $(T-1)^s$. In even rank C44-A12 selects a dense
set of closed fibres with no root of unity on $\mathcal W$,
which gives (A18.11). In odd rank C44-A16 selects a dense set
with total root-of-unity multiplicity one. The positive Gram
matrix of C44-A06 and the unconditional determinant criterion
C44-A15 then give precisely the stated alternative and sign.
These statements apply to the original finite-field axes, not
only to geometric generic parameter points.
\end{proof}

\paragraph{Application to the nonliftable centres and the remaining gate.}
Take the intrinsically nonliftable surface $S$ of C44-A14,
embedded as $S\times\{0\}^2\subset E^6$ in characteristic three.
There is just one prescribed centre. Formula (A18.10) makes the
residual rank even, so C44-A18a yields actual closed genuine
Lefschetz axes with $A_B=\mathbf Q w_S$ and $Q(w_S,w_S)>0$.
For every such axis the uncorrected class $[S]_B$ remains
nonliftable under any smooth projective characteristic-zero
lift, by C44-A14. This does \emph{not} certify that $w_S$ itself
is nonliftable: the ambient correction cancels the known
fourfold cup tensor exactly. Thus the construction removes the
monodromy/lift-of-centre obstruction, but the individual
primitive-class lifting question remains open.

\paragraph{Primary source locators and Grothendieck's programme.}
The moment machinery used above is N.~Katz,
\emph{Larsen's alternative, moments, and the monodromy of
Lefschetz pencils}, in \emph{Contributions to Automorphic Forms,
Geometry, and Number Theory} (2004), pp.~521--560:
Theorem 1.4.6 and its proof on p.~535; Lemmas 2.1.4--2.1.6
on pp.~540--541; Key Lemma 2.2.6 with complete proof on
pp.~544--545; the four-point calculation in 2.2.7--2.2.13 on
pp.~545--550. Author PDF:
\url{https://web.math.princeton.edu/~nmk/LAMMLF.pdf}.
The exact global norm and nontrivial inertia assertions are
SGA 7 II, Expose XV, Theorem 3.4(i),(iii), pp.~195--196;
the original-pencil existence uses the imposed-centre extension
proved in C44-A06, with SGA 7 II, Expose XVII, 2.4 and
3.7--4.3 as its ordinary-pencil input.

This settles the signed axis summand of JAC1964 H9-T7 for the
closed axes selected in C44-A18a. In cote 44, Th.~5, scan 004,
it gives the remaining axis sign for these actual special
pencils, on the primitive domain specified by the margin on
scan 008. Any assembly of the full height with the separate
abelian primitive term still uses the unrefereed 032/Milne 2026
input for that separate term. The result does not prove the
sign for every fixed axis and supplies no wrong-sign class.
In the SGA and Weil-cohomology framework, its new point is a
full orthogonal local system over the constrained parameter
space, constructed without a characteristic-zero model of its
prescribed centres.

% END work/axis/constrained_moment_monodromy.tex

% BEGIN work/axis/axis_realization_kernel.tex
% C44_POS /root/axis_independent. Author contribution; independent review pending.
\section{A common cycle-class kernel and a legitimate crystalline lifting test}
\label{sec:c44-a18b}

\begin{proposition}[C44-A18b: realization independence after one potential Tate direction]
Let $X/\mathbf F_q$ be an abelian variety of dimension $2r+2$,
and let $\kappa:B\hookrightarrow X$ be a smooth intersection of two
divisors of the same ample class $h$. Let $\Gamma$ be an absolute
Chow correspondence realizing the inverse of
$L_h^2:H^{2r}(X)\longrightarrow H^{2r+4}(X)(2)$, and put
\[
 P=\Delta_B-\kappa^*\Gamma\kappa_*,\qquad
 \mathscr C=P_*\mathrm{CH}^r(B_{\overline{\mathbf F}_q})_{\mathbf Q},
 \qquad q_B(z,z')=(-1)^r\deg(zz').
 \tag{A18b.1}
\]
The same correspondence is used in every realization. Choose actual
cycles $C_1,\ldots,C_s\in\mathscr C$ with nonsingular rational Gram
matrix $G=(q_B(C_i,C_j))$; $s=0$ is allowed. After extending the finite
ground field, assume these cycles and the correspondences are defined
over it. Let $K_\ell$ be their span and
\[
 V_\ell=\ker\kappa_*\subset H^{2r}(B,\mathbf Q_\ell(r)),\qquad
 W_\ell=V_\ell\cap K_\ell^\perp.
\]
Suppose the total algebraic multiplicity of root-of-unity roots of
the normalized Frobenius polynomial $\Phi_W$ is at most one.
Then the kernels, on the rational group $\mathscr C$, of every
$\ell$-adic cycle-class map ($\ell\ne p$) and of the rational
crystalline cycle-class map coincide.

Here a crystalline class over $\overline{\mathbf F}_q$ is tested after
any finite field of definition and extension of its Witt fraction
field; faithful scalar extension makes the vanishing condition
independent of that choice. The assertion is about projected cycles,
and does not assert realization independence for all cycles on $B$.
No positivity assumption on $G$, no Tate conjecture and no result
conditional on 032 are needed.
\end{proposition}

\begin{proof}
K\"unnemann's absolute Chow-motivic hard Lefschetz theorem
(\emph{Inventiones Mathematicae} 113 (1993), Theorem~5.2,
printed p.~96; Theorem~5.1 gives the decomposition used there)
supplies one $\Gamma$. If necessary replace the divisor by its
symmetric rational symmetrization, which has the same cohomology
class. Because this is a motivic inverse, it acts as the required
inverse in rational crystalline cohomology as well as in every
$\ell$-adic realization. We use only the absolute theorem over a
field, not a relative theorem over a mixed-characteristic base.
The self-intersection formula $\kappa_*\kappa^*=L_h^2$ then shows
that $P$ is the orthogonal cohomological projector onto $V$ in
each realization. Its action on Chow groups need not be idempotent.
Both the ambient summand and its orthogonal complement $V$ are
nondegenerate, by hard Lefschetz and Poincar\'e duality.

Intersection numbers are rational and agree with cup-product
traces in each realization. Thus the $C_i$ have the same matrix
$G$, are independent in each realization, and give nondegenerate
subspaces $K$. The characteristic polynomial on $V$ is the quotient
of the normalized middle polynomial of $B$ by that on
$H^{2r}(X)$. Katz--Messing, \emph{Some consequences of the Riemann
hypothesis for varieties over finite fields}, \emph{Inventiones
Mathematicae} 23 (1974), Theorem~1, printed p.~74, identifies
these characteristic polynomials in the reasonable cohomology
theories, including rational crystalline cohomology. Dividing by
$(T-1)^s$ gives the same rational polynomial $\Phi_W$ in all of them.
For crystalline cohomology over $\mathbf F_{q'}$, normalized
Frobenius means $q'^{-r}$ times the $q'$-power crystalline Frobenius,
which is linear over $W(\mathbf F_{q'})[1/p]$ and is an isometry.
Passing to a finite extension raises all normalized eigenvalues
to a common power and preserves their total root-of-unity
multiplicity.

We record the needed linear algebra without a semisimplicity
assumption. Let $F$ be an isometry of a nondegenerate symmetric
space over a characteristic-zero field, with rational
characteristic polynomial having at most one root of unity,
counted with algebraic multiplicity. If such a root exists it is
$1$ or $-1$ and is simple. Its eigenline $T$ is nondegenerate:
the generalized eigenspaces for every other eigenvalue are
orthogonal to $T$, since an isometry pairs only reciprocal
generalized eigenvalues. Were $T$ degenerate, its nonzero vector
would be orthogonal to the whole space. Moreover every vector
annihilated by some $F^e-1$ belongs to $T$; if there is no
root-of-unity root, every such vector is zero. In particular a
finite-orbit vector of square zero is zero.

Let $Z\in\mathscr C$, and first suppose its class vanishes in one
of the realizations under consideration. The cycle-theoretic
equalities
\[
                  q_B(Z,C_i)=0\quad(1\leq i\leq s),
                  \qquad q_B(Z,Z)=0
 \tag{A18b.2}
\]
follow in $\mathbf Q$. In every other realization, $Z$ lies in
$V$ because it is an actual projected cycle, and it lies in
$K^\perp$ by (A18b.2). Choose a finite field of definition for $Z$
as well. Its normalized Frobenius orbit is finite, and its square
is zero by (A18b.2). The preceding linear algebra on $W$ forces
its class to vanish. This proves all kernel equalities at once.
\end{proof}

\begin{corollary}[C44-A18c: a necessary filtered test for a class, not just a representative]
Keep the assumptions of C44-A18b. Let $R$ be a complete discrete
valuation ring of mixed characteristic with finite residue field,
and suppose that, after an allowed finite residue-field extension,
there is a smooth projective lift
\[
          \mathcal B\lhook\joinrel\longrightarrow\mathcal X
                       \longrightarrow\operatorname{Spec}R
 \tag{A18c.1}
\]
of the polarized embedded complete intersection $B\subset X$,
where $\mathcal X$ is an abelian scheme and $\mathcal B$ is cut
out by two relative divisors of its lifted line bundle. Write
$K_R=\operatorname{Frac}(R)$ and $K_0=W(k)[1/p]$. For
$Z\in\mathscr C$, suppose that its class in one $\ell$-adic
realization is the specialization of the class of a rational
algebraic cycle on $\mathcal B_{\overline K_R}$. Then, under
crystalline--de Rham comparison for this lift,
\[
 \operatorname{cl}_{\mathrm{cris}}(Z)\otimes1
       \in \operatorname{Fil}^r
                H^{2r}_{\mathrm{dR}}(\mathcal B_{K_R}/K_R).
 \tag{A18c.2}
\]
Finite extensions of $R$ may be made to define the generic cycle;
membership in the displayed filtration descends by faithful
field extension. The conclusion is only a necessary condition
for this specified embedded lift. It gives neither a converse
nor a lift-independent failure of (A18c.2).
\end{corollary}

\begin{proof}
After a finite extension let $\beta$ be a generic-fibre cycle
whose class specializes to $\operatorname{cl}_\ell(Z)$.
On the generic abelian variety choose an algebraic inverse
$\Gamma_\eta$ of $L_h^2$ and form the analogous algebraic
projector $P_\eta$. Its $\ell$-adic action is the orthogonal
projector onto the generic Gysin kernel. Smooth proper base
change, together with compatibility of Gysin and cup products,
identifies that projector with the special cohomological
projector $P$. Consequently $\beta'=P_{\eta*}\beta$ has the
same specialization of its $\ell$-adic class as $\beta$.

Let $\gamma=\operatorname{sp}(\beta')$ be the specialization
as an actual rational Chow cycle. The specialized correspondence
$\operatorname{sp}(\Gamma_\eta)$ acts as the inverse of $L_h^2$
also in crystalline cohomology: crystalline--de Rham comparison
intertwines its action with the generic de Rham inverse.
That inverse is unique. Thus $\operatorname{cl}_{\mathrm{cris}}
(\gamma)$ belongs to $V_{\mathrm{cris}}$, and applying the
original special-fibre $P$ leaves it unchanged. No lifting of the
original special-fibre correspondence $\Gamma$ has been assumed.

The actual cycle $P_*\gamma-Z$ belongs to $\mathscr C$ and has
zero $\ell$-adic class. C44-A18b gives
\[
 \operatorname{cl}_{\mathrm{cris}}(Z)
   =\operatorname{cl}_{\mathrm{cris}}(P_*\gamma)
   =\operatorname{cl}_{\mathrm{cris}}(\gamma).
 \tag{A18c.3}
\]
The right side corresponds to the de Rham class of $\beta'$,
which belongs to $\operatorname{Fil}^r$ because $\beta'$ is an
algebraic codimension-$r$ cycle. This proves (A18c.2).

For the comparison invoked here, see Berthelot--Ogus,
\emph{$F$-isocrystals and de Rham cohomology. I},
\emph{Inventiones Mathematicae} 72 (1983), Theorem~2.4,
printed pp.~169--170, Corollary~2.5 and base-change
Proposition~2.7, printed p.~171. Corollary~3.7 and
Remark~3.7.1, printed pp.~178--179, give compatibility with
the Chern classes of vector bundles over the same possibly
ramified trait. This also gives the cycle compatibility used
above: take the horizontal closure of each generic integral
cycle, resolve its structure sheaf by vector bundles on the
regular projective total space, and use the degree-$r$ leading
term of its Chern character. Flatness of the horizontal closure
over the trait identifies its fibre cycle with Chow
specialization. The identical argument on the product gives
compatibility for specialized correspondences.
Compatibility with cycle
classes and specialized correspondences is the standard cycle
compatibility; Emerton, \emph{A $p$-adic variational Hodge
conjecture and modular forms with complex multiplication},
\S2, (2.1) and the paragraph immediately before Conjecture~2.2,
printed pp.~4--5, explicitly states the cycle compatibility,
and Remark~2.5, printed pp.~6--7, the correspondence compatibility.
For complete ramified traits, the same comparison is stated in
Lieblich--Olsson, \emph{Fourier--Mukai partners of K3 surfaces in
positive characteristic}, \emph{Annales scientifiques de
l'\'Ecole normale sup\'erieure} 48 (2015), \S2.7, printed p.~1007,
with cup-product and trace compatibility. The cycle compatibility
comes from the Chern-class construction; it is not a variational
Hodge conjecture or an assertion of its converse.
\end{proof}

\paragraph{Exact scope for the programme.}
For the closed axes selected in C44-A18a, these results identify
the crystalline class of a projected algebraic $\ell$-adic class
independently of its chosen projected representative. In particular,
for the axes containing the nonliftable smooth surface of C44-A14,
it is legitimate to test the crystalline filtration of
$w_S=P_*[S]$ along any specified embedded lift. What has not been
proved is that this class fails the filtration test for every
possible lift. Nonliftability of $S$ itself, or of the uncorrected
class $[S]$, supplies no such assertion: the ambient cup tensor
cancels under $P$, as proved in C44-A14a. These results repair a
possible class-versus-representative gap in route (a) of
JAC1964 QH9-axis; they do not close that row or prove a new
universal sign assertion. In Grothendieck's language they concern
the Weil realizations of the primitive axis correspondence in
cote~44, Th.~5, scan~004 and the primitive warning on scans~008--009,
without assuming condition $D(B)$ for all of $B$.

% END work/axis/axis_realization_kernel.tex

% BEGIN work/scan_w/equivariant_axis_blocks.tex
% C44_POS /root/scan_w. No universal lifting of a special Tate projector asserted.
\section{Several potential Tate directions separated by geometric symmetries}
\label{sec:c44-a19}

\begin{theorem}[C44-A19: the orthogonal block criterion]
Let $B\subset X/\mathbf F_q$ be the smooth intersection of two
members of one ample line bundle as in C44-A15, with
$\dim X=2r+2$. Suppose an embedded polarized lift
$\mathcal B\subset\mathcal X$ over a mixed-characteristic trait
has a characteristic-zero generic fibre admitting a complex
embedding. On its relative variable middle cohomology suppose
there are rational algebraic correspondences $e_1,\ldots,e_b$
whose realizations are pairwise orthogonal self-adjoint
idempotents summing to the identity. In particular they
preserve this variable cohomology on both fibres and in every
coefficient realization. Put
\[
 V_{i,\ell}=e_iV_\ell,\qquad
          q=(-1)^r\int_B(-)(-).
 \tag{A19.1}
\]
Assume that the normalized Frobenius polynomial on every
$V_{i,\ell}$ is rational and independent of $\ell\ne p$.
Let $K\subset A_B$ be a known positive rational algebraic
subspace stable under the $e_i$; $K=0$ is allowed. Use the
common cycle projection described below to realize
$K_i=e_iK$ in each $V_{i,\ell}$, and let $\tau_i$ be the total
root-of-unity multiplicity on $V_{i,\ell}\cap K_{i,\ell}^\perp$.
If
\[
                            \tau_i\leq1\quad(1\leq i\leq b),
 \tag{A19.2}
\]
then $q$ is positive definite on the whole rational algebraic
axis space $A_B$. Its rational dimension satisfies
\[
     \dim K\leq\dim A_B\leq\dim K+
                              \#\{i:\tau_i=1\}.
 \tag{A19.3}
\]
Every additional nonzero algebraic direction has positive
square; the existence or liftability of those directions is
not asserted. No Tate conjecture, Frobenius semisimplicity,
standard conjecture $D(B)$, or 032 input is used.
\end{theorem}

\begin{proof}
All Gram determinants use the ordinary convention of A15.
The relative algebraic correspondences give rational Hodge
endomorphisms on the Betti variable space of the generic
fibre. Their images $V_{i,\mathrm B}$ are orthogonal
nondegenerate rational Hodge substructures. They are primitive
for the inherited ample class: the entire variable space is
primitive by the dual weak-Lefschetz argument of A15.
On each image the Hodge--Riemann relations make the signed
form positive on type $(r,r)$, while every negative real
summand arises from a conjugate pair of off-diagonal Hodge
types. Its negative index is therefore even. Consequently
\[
              D_i:=\det(q|V_{i,\mathrm B})>0.
 \tag{A19.4}
\]
For a zero-dimensional image take $D_i=1$. Smooth proper
base change and Betti--etale comparison for these same
relative correspondences identify this determinant
squareclass with that on $V_{i,\ell}$ for every $\ell\ne p$.
The special algebraic classes themselves need not lift.

Here is the promised correction of cycle representatives.
Use the common special-fibre algebraic projector
\[
                  P=1-\kappa^*(L_h^2)^{-1}\kappa_*
 \tag{A19.5}
\]
from A15, constructed using K\"unnemann's absolute Chow
Lefschetz inverse. Apply $e_iP$ to chosen representatives
of $K$. At the initial coefficient prime these images span
$K_i$. Choose a rational basis from these images. Its
positive rational Gram matrix $G_i$ makes the chosen
cycles independent in every realization; their classes
are in $V_{i,\ell}$ in every realization. This choice
does not infer realization-independent vanishing from
vanishing at one prime. After a finite extension all
these cycles and correspondences are defined, and their
normalized Frobenius action on $K_{i,\ell}$ is trivial.
The residual polynomial is the given block polynomial
divided by $(T-1)^{\dim K_i}$, hence remains rational and
independent of $\ell$. Field extension preserves $\tau_i$.

For an arbitrary class in $A_B$, first apply $P$ to its
cycle representative, then apply each $e_i$, and subtract
its orthogonal projection to the chosen $K_i$ using
$G_i^{-1}$ and rational intersection numbers. This gives
actual rational cycles in the residual block for every
$\ell\ne p$, whose initial classes sum to the given
class modulo $K$. Distinct blocks remain orthogonal.

The determinant proof of A15 now applies separately to
each block, as can also be seen directly. An algebraic
residual class has finite Frobenius orbit. If $\tau_i=0$
it is zero. If $\tau_i=1$, the unique cyclotomic root is
simple and is $1$ or $-1$. For $S=F^2$ its eigenspace
$T_{i,\ell}$ at $1$ is a nondegenerate line; write its
orthogonal complement as $U_{i,\ell}$. A nonzero residual
cycle $z$ at the initial prime has
$a=q(z,z)\in\mathbf Q^*$, so its realization is nonzero
and spans $T_{i,\ell}$ at every prime.

On $U_{i,\ell}$, the operator $S-S^{-1}$ is invertible
and skew-adjoint. The determinant of the resulting
alternating pairing is a square. If $\psi_i$ is the
common rational polynomial of $S$ on $U_{i,\ell}$, this
gives, exactly as in A15,
\[
 \det(q|U_{i,\ell})\equiv
      \delta_i:=\psi_i(1)\psi_i(-1)>0
                    \pmod{(\mathbf Q_\ell^*)^2}.
 \tag{A19.6}
\]
Positivity of $\delta_i$ follows by pairing conjugate
unit-circle eigenvalues, none equal to $\pm1$.
Consequently
\[
                 a\det(G_i)\delta_i/D_i
                  \in(\mathbf Q_\ell^*)^2
                         \quad\hbox{for every }\ell\ne p.
 \tag{A19.7}
\]
The missing-$p$ rational squareclass injection used in
A15 makes this a rational square. The other displayed
factors are positive, so $a>0$. Pairing any further
residual cycle with $z$ gives its rational coefficient
on the same line. Each block therefore adds at most
one positive rational direction. Orthogonality proves
both the sign and (A19.3).
\end{proof}

\begin{corollary}[C44-A19a: a tame elementary two-group criterion]
Let $R$ be a mixed-characteristic discrete valuation ring
with finite residue field $k$ of odd characteristic, and
let $(\mathcal X,\mathcal L)/R$ be a polarized abelian
scheme of relative dimension $2r+2$. Suppose the constant
group $G=(\mathbf Z/2\mathbf Z)^s$ acts on $\mathcal X$
and $\mathcal L$ is given an actual $G$-linearization.
Let $B\subset X_k$ be smooth and be cut out by two
simultaneous eigenvectors
\[
       s_0\in H^0(X_k,L)_{\chi_0},\qquad
       s_1\in H^0(X_k,L)_{\chi_1}.
 \tag{A19.8}
\]
They may have the same character. Assume a complex
embedding of the generic field as in A19.

Then the actual embedded $G$-axis lifts. Its rational
orthogonal character projectors are
\[
        e_\chi=\frac1{|G|}\sum_{g\in G}\chi(g)g^*,
             \qquad\chi:G\longrightarrow\{1,-1\}.
 \tag{A19.9}
\]
Their normalized Frobenius polynomials on variable
cohomology are rational and independent of $\ell\ne p$.
For any known positive $G$-stable rational algebraic
subspace $K$, condition (A19.2) on these character
blocks implies positivity on all of $A_B$.
In particular distinct character blocks can each
contribute one unknown potential Tate direction.
No extension of a special-fibre Frobenius projector
is part of this hypothesis or conclusion.
\end{corollary}

\begin{proof}
The ample special line bundle $L$ has a nonzero section.
Its polarization has finite kernel, so the index theorem
for line bundles on abelian varieties gives
$H^j(X_k,L)=0$ for $j>0$. Upper semicontinuity gives the
same vanishing on the generic fibre. Coherent base change
therefore makes $H^0(\mathcal X,\mathcal L)$ a finite
free $R$-module with surjective reduction onto
$H^0(X_k,L)$. These are Milne's Theorems 13.3(c) and
4.2(c),(e), author notes pp.~25 and 6--7, respectively.

Because $2$ is invertible in $R$, the operators (A19.9)
also act on this module. Choose arbitrary lifts of
$s_0,s_1$ and apply $e_{\chi_0},e_{\chi_1}$, respectively.
The resulting sections are actual eigensection lifts
with their prescribed reductions. Their common zero
locus $\mathcal B$ is $G$-stable. Along its special fibre
the two relative differentials have full rank, since
$B$ is smooth of codimension two. Thus $\mathcal B/R$
is smooth along the special fibre. The nonsmooth locus
is closed and proper over $R$; if nonempty its image
would contain the closed point. It is therefore empty.
This proves smoothness of the entire embedded lift.

The graph correspondences of the $G$-action are relative
algebraic correspondences. They commute with restriction,
Gysin and multiplication by $h$, and so preserve the
variable middle cohomology. An automorphism's adjoint
for cup product is its inverse. Here every element is
its own inverse, so (A19.9) are self-adjoint,
pairwise orthogonal rational idempotents summing to one.
No restriction $\ell\nmid|G|$ is needed: the coefficient
fields have characteristic zero even when $\ell=2$.

We prove degreewise rationality of the equivariant
Frobenius traces, rather than infer it from the total
characteristic polynomial. All group elements are
defined over $k$, hence commute with Frobenius. For
each $g\in G$ and each positive integer $n$, base change
$B$ to $k_n=\mathbf F_{q^n}$ and descend its base change
to $k_{2n}$ by the cocycle taking the nontrivial Galois
element to $g$. Effective projective Galois descent
gives a smooth projective twist $B^{(g,n)}/k_n$.
Its Frobenius on cohomology is $g^*F^n$ (the inverse
convention gives the same operator since $g=g^{-1}$).
Deligne's Weil I, Theorem 1.6, gives rational,
coefficient-independent characteristic polynomials
degree by degree for this twist. In particular
\[
       \operatorname{Tr}(g^*F^n\mid H^j(B_{\bar k}))
 \tag{A19.10}
\]
is rational and independent of $\ell$. The same argument
applies to $X$. Twisting is legitimate as a projective
scheme: finite Galois descent is effective here, and
the given linearized ample line bundle descends with
the cocycle. See Stacks, Section 39.25, Lemma 39.25.3
in the version read for this proof.

Weak Lefschetz and hard Lefschetz identify the ambient
middle image with $H^{2r}(X)$ in a $G$- and
Frobenius-equivariant way. Subtract its traces in
(A19.10) and divide by the Tate-normalizing factor.
Taking the rational average (A19.9) now gives all
power traces of Frobenius on each $e_\chi V_\ell$.
Newton identities recover a rational characteristic
polynomial independent of $\ell$. Its degree is also
common, by the equivariant smooth proper comparison
with the generic Betti fibre. The hypotheses of A19
are now established, including the two requirements
that are not automatic for an arbitrary special
Frobenius eigenspace: an actual extending self-dual
algebraic block and a common rational polynomial.
\end{proof}

\begin{corollary}[C44-A19b: finite-group matrix idempotents]
In A19a replace the elementary two-group by any finite
constant group $G$ with $p\nmid|G|$. Assume the given
$G$-action and linearized ample bundle extend over $R$.
Let $U_R$ be a free rank-two $R$-representation of $G$,
and suppose a $G$-equivariant map
$U_k\longrightarrow H^0(X_k,L)$ has a smooth
codimension-two common zero locus $B$.
Then this actual embedded axis and its $G$-action lift.
Every set of pairwise orthogonal self-adjoint rational
idempotents in $\mathbf Q[G]$ summing to one satisfies
the geometric and rational-polynomial hypotheses of A19.
Self-adjointness here is for the group-algebra involution
$g^*=g^{-1}$. The idempotents need not be central.

For $G=S_3$ and a transposition $\sigma$, one such set is
\[
 \begin{aligned}
 e_1&=\tfrac16\sum_{g\in S_3}g,&
 e_{\mathrm{sgn}}&=\tfrac16\sum_{g\in S_3}\mathrm{sgn}(g)g,\\
 e_{\mathrm{std}}&=1-e_1-e_{\mathrm{sgn}},&
 e_+&=\tfrac12e_{\mathrm{std}}(1+\sigma),\qquad
 e_- =\tfrac12e_{\mathrm{std}}(1-\sigma).
 \end{aligned}
 \tag{A19.11}
\]
The four orthogonal blocks are $e_1,e_{\mathrm{sgn}},e_+,e_-$.
The last two have the same Frobenius polynomial, that of
the standard multiplicity module. If the trivial block
has no potential Tate direction and the sign and standard
multiplicity modules each have at most one, then the
axis sign is positive even though its total potential
Tate multiplicity can be three.
\end{corollary}

\begin{proof}
The coherent base-change argument of A19a is unchanged.
Lift the given map $U_k\to H^0(X_k,L)$ arbitrarily to
an $R$-linear map. Average it by
\[
                 \varphi\longmapsto
       \frac1{|G|}\sum_{g\in G}g\varphi g^{-1}.
 \tag{A19.12}
\]
Since $|G|$ is a unit in $R$, its reduction is the
original equivariant map. The same Jacobian and properness
argument shows that its common zero locus is a smooth
embedded lift. This does not assume that an arbitrary
residue-field representation lifts: $U_R$ is given.

Every group-algebra element defines a relative graph
correspondence, and the adjoint involution is $g\mapsto g^{-1}$.
Thus the stated algebraic identities give the required
orthogonal blocks in every realization. For rationality,
the twisting argument (A19.10) works for every finite-order
$g$: use the cyclic extension of degree $\operatorname{ord}(g)$
over $k_n$, choosing the cocycle or its inverse according
to the Frobenius convention. Hence all
$\operatorname{Tr}(g^*F^n\mid H^j)$ are rational and
independent of $\ell$. For $e=\sum a_gg$ with $a_g\in\mathbf Q$,
\[
       \operatorname{Tr}(F^n\mid eV)
                    =\operatorname{Tr}(eF^n\mid V)
                    =\sum_g a_g\operatorname{Tr}(gF^n\mid V),
 \tag{A19.13}
\]
because $e$ commutes with Frobenius and is idempotent.
Newton identities again give a common rational polynomial.
Centrality of $e$ is unnecessary.

The formulas (A19.11) are self-adjoint orthogonal
idempotents by the elementary character decomposition
of $\mathbf Q[S_3]\simeq\mathbf Q\oplus\mathbf Q\oplus M_2(\mathbf Q)$.
On the standard representation
$\{(x_1,x_2,x_3)\in\mathbf Q^3:\sum x_i=0\}$ with
its usual positive dot product, the two $\sigma=(12)$
eigenlines are spanned by $(1,1,-2)$ and $(1,-1,0)$;
their squared lengths are $6$ and $2$. They tensor with
the same multiplicity module, on which Frobenius acts.
The geometric pairing on the isotypic part is this
dot product tensor a rational symmetric pairing on
that module. Thus the two block polynomials agree,
and their relative scalar factors are positive.
Applying A19 to the four blocks proves the assertion.
\end{proof}

\paragraph{An explicit source of the geometric hypotheses.}
For odd residue characteristic take the good-reduction
model of $E:y^2=x^3-x$, put $\mathcal X=\mathcal E^6$,
and use $\mathcal L=\boxtimes_{i=1}^6\mathcal O_{\mathcal E}(t[0])$.
Coordinate sign changes give a constant elementary
two-group and its natural linearization on this line
bundle. Every smooth complete intersection of two
simultaneous eigensections in its special fibre
therefore has the lift and character correspondences
required in A19a. For instance, for even $t=2d$,
the invariant subsystem is the pullback of
$|\mathcal O(d,\ldots,d)|$ on $(\mathbf P^1)^6$ under
the coordinate double covers; a general pair transverse
to the finitely many branch strata gives a smooth
complete intersection for sufficiently large $d$.
This describes actual geometry, not a postulated
relative projector. It does not by itself verify
(A19.2) for a specified finite-field member.

\paragraph{Scope of the additional test.}
A19 can settle a residual block of total potential
Tate multiplicity two or more when its directions
are separated by these extending rational orthogonal
projectors and every block satisfies (A19.2).
It does not settle two potential directions in one
character block. No explicit fixed-axis member of
total multiplicity at least three satisfying all of
these conditions is certified here. In particular
one must not declare such a member solely from the
existence of sign symmetries. A pencil spanned by
invariant sections can have simultaneous orbits of
nodes; the displayed smooth complete intersections
are not being asserted to be axes of genuine
Lefschetz pencils without a separate check.

For cote 44, Theorem 5, scan 004 and the primitive
qualification on scan 008, the criterion controls
the same H9-T7 axis form. The positivity conclusion
is unconditional under its stated geometric and
Frobenius tests. As in A15, assembly with the separate
abelian term in the full height remains conditional
on the unrefereed 032/Milne2026 input. The arithmetic
rank-two obstruction of A17 is not removed for an
axis lacking an extending orthogonal decomposition.

% END work/scan_w/equivariant_axis_blocks.tex

% BEGIN work/scan_w/affine_s3_axis_geometry.tex
% C44_POS /root/scan_w. Geometry and parity input for C44-A21.
\section{An affine symmetric-group family of genuine Lefschetz axes}
\label{sec:c44-a20}

\begin{theorem}[C44-A20: the affine $S_3$ axis family]
Let $k$ be a finite field of characteristic $p>3$, let $E/k$
be an elliptic curve, and, after a finite extension of $k$,
choose a nonzero point $T\in E[3](k)$. On
\[
       X=E^2\times E^2\times E^2,
       \qquad(a,b;u,v)\in X,
\]
put
\[
 \begin{aligned}
 R(P,Q)&=(Q,-P-Q),& S(P,Q)&=(Q,P),\\
 r(a,b;u,v)&=(a+T,b;Ru,Rv),&
 s(a,b;u,v)&=(-a,-b;Su,Sv).
 \end{aligned}                                                   \tag{A20.1}
\]
These generate a faithful action of $G=S_3$, with
$r^3=s^2=1$ and $srs=r^{-1}$. Its linear tangent
representation is $2\varepsilon\oplus2\rho$, where
$\varepsilon$ is the sign representation and $\rho$
the rational two-dimensional standard representation.
The order-three elements act freely. The loci fixed by
the three involutions are pairwise disjoint, each a
union of sixteen translates of an abelian surface.

Choose the product very ample bundle
$M=\boxtimes_{i=1}^6\mathcal O_E(3[0])$ on $X$ and put
\[
               H=\left(\bigotimes_{g\in G}g^*M\right)^{\otimes2}.
 \tag{A20.2}
\]
It has a natural $G$-linearization, and every stabilizer
acts trivially on its fibres. Allow either $Z=\varnothing$
or the following prescribed centre. Choose a point
$c\in E^2\times E^2$ with a free orbit for the linear
$S_3$-action on the last four factors, and put
\[
                  Z=\coprod_{g\in G}Z_g,
             \qquad Z_g=E^2\times\{gc\}.
 \tag{A20.3}
\]
These six surfaces are pairwise disjoint and avoid all
fixed loci. All these choices are defined after a finite
extension of $k$.

For every sufficiently large $t$, let $L=H^t$ and
\[
        \mathcal E_t=\operatorname{Hom}_G
                  (\rho,H^0(X,I_ZL)).
 \tag{A20.4}
\]
There is a nonempty open $\Omega\subset\mathcal E_t$
whose parameters give smooth fourfolds $B\subset X$,
the common zero loci of the two component sections.
Their original two-dimensional systems are genuine
Lefschetz pencils: every singular member has exactly
one ordinary quadratic singularity, outside the smooth
axis, and the generic member is smooth.
The group acts faithfully on their parameter line.

Let $h=c_1(L)$, let $F$ be one component of $X^s$, and set
\[
                 d_6=\int_Xh^6,\qquad d_2=\int_Fh^2.
\]
For the middle variable cohomology $V$, the multiplicities
of $\mathbf1,\varepsilon,\rho$ are exactly
\[
 \begin{split}
 m_1&=\frac{5d_6}{6}-8d_2-12,\\
 m_\varepsilon&=\frac{5d_6}{6}+8d_2-39,\\
 m_\rho&=\frac{5d_6}{3}-57.
 \end{split}                                                     \tag{A20.5}
\]
Their parities are respectively even, odd, odd, and
all exceed eight for large $t$.
With the six centres the positive known primitive span
$K$ is the regular rational representation of $G$;
the residual multiplicity ranks have parities odd,
even, odd and still exceed eight for large $t$.

The family also admits, in its parameter-space boundary,
a member with exactly one free orbit of six ordinary
nodes, and a member with exactly one involution-fixed
orbit of three ordinary nodes. All nodes can be chosen
outside $Z$, with no other singularities. Each admits
a transverse one-parameter smoothing within $\mathcal E_t$
with regular total space at every node and geometric
generic parameter in $\Omega$. The stabilizer of a node
in the orbit of three acts by its nontrivial character
on that node's vanishing-cycle line.

Finally the action, the linearized bundle, the prescribed
centres if present, and each of these smooth axes admit
a common mixed-characteristic lift, carrying their
actual $G$-action and their two-dimensional generator
representation. No algebraic class in the remaining
variable cohomology is asserted to lift.
\end{theorem}

\begin{proof}
\emph{The affine action and its fixed loci.}
The matrices $R,S$ satisfy $R^3=S^2=1$ and $SRS=R^{-1}$.
The translations in (A20.1) satisfy the same relation,
since $3T=0$ and $s$ negates the first coordinate.
An $r$-fixed point would have $a+T=a$, which is impossible.
The same is true of $r^2$.
The $s$-fixed locus is
\[
          E[2]^2\times\{(P,P):P\in E\}
                         \times\{(Q,Q):Q\in E\}.
 \tag{A20.6}
\]
It has sixteen components, all abelian surfaces.
The other two involutions are conjugate to $s$ and
have the same description up to an automorphism.
A point fixed by two distinct involutions would be
fixed by an element of order three. Hence their fixed
loci are disjoint, and these are all points with
nontrivial stabilizer. At every involution-fixed point
the tangent eigen-dimensions are $(2,4)$ for $(+1,-1)$.
Translation does not change tangent representations.

The norm tensor in (A20.2) has the standard permutation
linearization. Squaring makes its fibre character at
each order-two stabilizer trivial. There are no other
nontrivial stabilizers. It is very ample: all factors
are very ample, hence globally generated, and their
tensor product is very ample. A point $c$ outside the
finitely many proper fixed loci in $E^4$ has the required
free orbit. The surfaces (A20.3) are distinct parallel
fibres, hence disjoint. No point of $Z_g$ can be fixed
by an involution, because $gc$ is not fixed by it.

\emph{Compatible jets in the parameter space.}
For large $t$, sections of $I_ZH^t$ have arbitrary
bounded jets at any fixed number of points outside $Z$,
and generate the normal first jets along $Z$.
The bounds can be chosen uniformly in the points.
One elementary construction first chooses $c_0$ with
$I_ZH^{c_0}$ globally generated. A section of this
sheaf can be chosen nonzero at any prescribed bounded
set outside $Z$. Multiply it by products of hyperplanes
to interpolate the required values and derivatives;
the degree needed is bounded by the number and orders
of the jets. For normal jets along $Z$, use the
generators of $(I_Z/I_Z^2)H^{c_0}$, and the same
hyperplane interpolation. Increasing $c_0$ gives
these generators by Serre vanishing. Extra hyperplane
factors vanishing to order three at an already fixed
orbit preserve its second jets while providing all
the stated jets at another orbit.

Apply this construction to whole $G$-orbits, of length
at most six, and to two such orbits when needed.
Specify compatible translated jets, then average a
lift in $H^0(I_ZH^t)\otimes\rho^*$ by $|G|^{-1}\sum g$.
Because $p>3$, averaging is exact. Thus evaluation
from (A20.4) surjects onto every compatible jet target.
At a free point there are no restrictions on the two
component jets. At a point fixed by $s$, use an
$s$-eigenbasis of $\rho$ and local tangent coordinates
with two plus and four minus directions. The allowed
jets have the form
\[
 \begin{array}{c|c|c}
       &\text{value}&\text{first derivative}\\ \hline
 f_+&\text{arbitrary}&\text{arbitrary in the two plus directions}\\
 f_-&0&\text{arbitrary in the four minus directions}.
 \end{array}                                                    \tag{A20.7}
\]
Second derivatives have the corresponding even or odd
parity. These are exactly the compatible targets;
they are not extra assumptions on a generic section.

\emph{Smoothness of the axis and its fixed curves.}
At a free point off $Z$, singularity of the common
zero locus imposes two value conditions and rank at
most one on a $2$ by $6$ derivative matrix. The latter
has codimension five, so the total codimension is seven,
greater than the six-dimensional point space.
At an involution-fixed point, the common zero condition
is $f_+=0$. The two derivative rows belong to different
eigenspaces, so rank drops precisely if one row is zero.
The two bad conditions have codimensions $1+2=3$ and
$1+4=5$, both greater than the fixed-locus dimension two.
Along a prescribed centre, both values and tangential
derivatives vanish. The two normal rows are arbitrary
in a four-dimensional conormal fibre; rank at most one
has codimension three, greater than the centre dimension
two. These incidence counts show that the common zero
locus of a general parameter is smooth everywhere.
They also show that on each component $F$ of $X^s$
the fixed locus $B^s$ is the smooth divisor $(f_+|F=0)$
of class $h|F$. Thus $B^s$ is a union of sixteen smooth
ample curves. The same holds for the other involutions.

\emph{Why the original pencil is genuinely Lefschetz.}
The standard representation has faithful projective
$S_3$-action. There are only eight geometric points of
$\mathbf P(\rho)$ with nontrivial stabilizer: the two
order-three eigenlines and the six involution eigenlines.
For each of these finitely many lines $u$, a general
member $D_u=(\varphi(u)=0)$ is smooth.
At a free point this follows from arbitrary first jets:
a singularity imposes seven conditions while the point
varies in dimension six. At a point fixed by an
involution, an invariant eigenmember has arbitrary
value and two plus derivatives, hence three conditions
over a surface. An anti-invariant eigenmember vanishes
there identically, but has four arbitrary minus
derivatives, hence four conditions over that surface.
For any other line the two components of $u$ in (A20.7)
are nonzero, and its value and all six derivatives
are arbitrary. Along $Z$ independence of the two
normal rows makes every member smooth. Thus all eight
exceptional members can simultaneously be required
smooth on a nonempty open of $\mathcal E_t$.

Now vary a nonexceptional $u$. A member singular at a
free point imposes seven conditions; a degenerate
quadratic part imposes one more. The pair $(x,u)$ has
dimension seven, so a general parameter has no such
degenerate singularity. A singularity at a point with
involution stabilizer imposes seven conditions by
(A20.7), while $(x,u)$ has dimension three, and is absent.
Two singular points in distinct free $G$-orbits impose
fourteen independent first-jet conditions. Their points
and the parameter $u$ vary in dimension thirteen, so
a general pencil avoids them. If the second point is
$gx$ in the same orbit, singularity there is equivalent
to singularity of $\varphi(g^{-1}u)$ at $x$. Since $u$
has trivial projective stabilizer, these are two
independent components of $\varphi$, so again fourteen
conditions, now over a seven-dimensional parameter
space $(x,u)$. This case is absent as well.

The smooth axis excludes a singular point on the base
locus. The nondegenerate Hessian and the nonzero
parameter derivative at each remaining singularity
make the critical pairs isolated and reduced. They
form a finite scheme, because $X\times\mathbf P^1$ is
proper. The generic pencil member is therefore smooth,
and each singular member has just one ordinary node.
There are in fact singular members: the universal
first jet is a section of
$J^1(L)\boxtimes\mathcal O_{\mathbf P^1}(1)$ of rank
seven on $X\times\mathbf P^1$. Since $\Omega_X^1$ is
trivial, its top Chern number is
\[
            \int_{X\times\mathbf P^1}(h+\eta)^7
                        =7\int_Xh^6>0.
 \tag{A20.8}
\]
The transverse zeros just described have this length.
All required good conditions are open and their bad
incidences have the smaller dimensions verified above.
Their intersection is the nonempty open $\Omega$.

\emph{The two boundary degenerations.}
At a chosen free point off $Z$, prescribe the two
values to vanish, their derivative matrix to have
rank one, and the remaining equation after eliminating
the first to have a nondegenerate quadratic part in
five variables. The compatible orbit jets prescribe
exactly six ordinary nodes. At a chosen point $x$ fixed
by $s$, outside $Z$, prescribe
\[
      f_+(x)=0,\qquad df_+(x)=0,\qquad df_-(x)\ne0,
 \tag{A20.9}
\]
and a nondegenerate quadratic part for $f_+$ on
$\ker df_-$, with two plus and three minus variables.
This is an allowed even quadratic form. It gives one
ordinary node at $x$ and at each of its two translates,
and no additional prescribed singularities.

In either case the compatible second-jet evaluation
is surjective. Keeping these jets fixed, the remaining
point and normal-jet incidences have exactly the same
strict dimension inequalities as above away from the
chosen orbit. A general choice is consequently smooth
away from that orbit, including along $Z$. Choose an
equivariant deformation changing the surviving nodal
equation's constant term; at a fixed node this is the
allowed $f_+$ value. Its total space is regular at
every node. A general such one-parameter direction has
its generic parameter in the dense open $\Omega$.

For completeness the stabilizer character can be
calculated algebraically. Near the fixed node eliminate
the odd equation using its nonzero odd derivative.
The remaining smooth five-dimensional total space has
an invariant Morse function with tangent eigen-dimensions
$(2,3)$. An equivariant etale Morse change of coordinates
reduces it to an invariant nondegenerate quadratic form.
Here is a direct construction in odd characteristic:
choose eigen-coordinates by averaging; write the
function as $x^tA(x)x$ with symmetric $A$ and average
so $A(sx)=sA(x)s$. With $A_0=A(0)$, the matrix
$A_0^{-1}A(x)$ has a unique square root near the identity
on an etale neighbourhood, since the derivative of
$P\mapsto P^2$ there is multiplication by two.
Uniqueness makes this root $A_0$-self-adjoint and
$s$-equivariant. The change $y=P(x)x$ is etale and
puts the function into the form $y^tA_0y$.

Compactify the smooth affine quadratic fibre by one
extra coordinate. Its dimension is four and $s$ acts
on the six quadratic coordinates with determinant
$(-1)^3=-1$. It therefore negates the volume element
in the centre of the even Clifford algebra and
exchanges its two primitive idempotents. These label
the two families of maximal generators of the quadric.
Their cohomology-class difference generates primitive
middle cohomology and is negated by $s$. The affine
quadric calculation identifies this line with the
local vanishing line, so its stabilizer character is
nontrivial. Exact primary locators are SGA 7 II,
XII, 1.4--1.5, pp.~65--66; Proposition 1.12, p.~69;
Theorem 3.3, pp.~75--77; 3.5--3.7, pp.~78--81;
and XV, 2.2.3--2.2.5, pp.~182--185. XV, Theorem 3.4,
pp.~195--196 supplies the global nonzero norm and
surjective quadratic inertia for the regular smoothing.
Thus the two orbit representations of vanishing lines
are respectively $\mathbf Q[G]$ and
$\operatorname{Ind}_{\langle s\rangle}^G\varepsilon|_{\langle s\rangle}$.

\emph{Exact equivariant Euler calculation.}
The tangent bundle of $X$ is trivial and
$N_{B/X}=L|B\oplus L|B$. Consequently
\[
                  \chi(B)=\int_Bc_4(T_B)=5d_6.
 \tag{A20.10}
\]
For a finite-order automorphism of order prime to $p$,
its alternating cohomological trace equals the Euler
characteristic of its fixed locus. This is
Deligne--Lusztig, \emph{Annals of Mathematics} 103
(1976), Theorem 3.2, printed p.~119, in the tame case.
Since $r$ is free, its trace is zero. Each of the
sixteen fixed curves of $s$ is a smooth ample divisor
on an abelian surface, so has Euler characteristic
$-\int_Fh^2=-d_2$. All the components $F$ are translates
and have the same intersection number. Thus the
equivariant Euler character on the three conjugacy
classes $(1,s,r)$ is
\[
                            (5d_6,-16d_2,0).
 \tag{A20.11}
\]
Its multiplicities are
$5d_6/6-8d_2$, $5d_6/6+8d_2$, $5d_6/3$.

Translations act trivially on $H^1(X)$, whose rational
$G$-representation is $4\varepsilon\oplus4\rho$.
The exterior-power generating characters are
\[
 (1+z)^{12},\qquad(1+z)^4(1-z)^8,\qquad(1+z^3)^4.
 \tag{A20.12}
\]
The degree-four coefficients are $(495,-17,0)$;
their irreducible multiplicities are $(74,91,165)$.
The alternating sum of the coefficients in degrees
zero through three is $(-165,-5,-3)$, whose
multiplicities are $(-31,-26,-54)$.
Weak Lefschetz identifies the lower cohomology of
$B$ with that of $X$, and Poincare duality identifies
the upper cohomology with their duals. All the rational
$S_3$ representations are self-dual. Subtracting twice
this lower alternating character and the ambient
$H^4(X)$ character from (A20.11) gives (A20.5).
The sum of dimensions is $5d_6-165$, as in A03b.

Riemann--Roch on the abelian sixfold gives
$d_6=720\chi(L)$. Hence $5d_6/6$ and $5d_6/3$ are
even integers, and $8d_2$ is even. This proves the
claimed parities. As $L=H^t$, the leading terms are
positive multiples of $t^6$, while $d_2$ grows only
as $t^2$. Each multiplicity therefore exceeds eight
for all sufficiently large $t$.

\emph{The optional six positive classes.}
Put $h_0=c_1(H)$, let $A_0=(L_{h_0}^2)^{-1}[Z_1]_X$
be the common rational algebraic ambient correction,
and set
\[
       w_g=[Z_g]_B-t^{-2}\kappa^*A_0,\qquad
       d=\int_{Z_1}h_0^2>0,\qquad
       c_0=\int_XA_0^2h_0^2.
 \tag{A20.13}
\]
All $Z_g$ have the same ambient cohomology class. Their
normal bundles in $X$ are trivial of rank four, so the
normal exact sequence gives
$c_2(N_{Z_g/B})=3t^2h_0^2|Z_g$. The disjoint-centre
Gram calculation of A06 becomes the exact matrix
\[
                      3t^2d\,I_6-t^{-2}c_0J_6.
 \tag{A20.14}
\]
It is positive definite for large $t$; for example
$3t^4d>6|c_0|$ suffices. Thus the six classes are
independent. The common ambient correction is
$G$-invariant, so $G$ permutes them regularly.
Removing $K$ subtracts multiplicities $(1,1,2)$
from (A20.5), proving the residual parity assertion.

\emph{The actual equivariant lift.}
Lift the elliptic curve over a mixed-characteristic
henselian trait. Since $p>3$, its finite etale
three-torsion lifts $T$ after a finite extension;
the integer matrices in (A20.1) then give the same
affine action on the abelian scheme. The specified
product bundle $M$ extends using the zero section of
the elliptic scheme. Apply the relative norm and square
to that actual extension; its reduction is exactly
the prescribed $H$ in (A20.2), including its linearization.
When centres are present, lift
$c$ by smoothness of the abelian scheme; its free
orbit and the six disjoint fibre centres remain so
over the trait. Properness excludes a new collision
whose closure would meet the special fibre.
For large $t$, relative Serre vanishing for the
centre ideal gives reduction-surjectivity of
$H^0(I_ZH^t)$. With no centres this also follows
from the abelian line-bundle index theorem.
Average an arbitrary lift of the map (A20.4) by
$(1/6)\sum_{g\in G}g\varphi g^{-1}$. It is an actual
equivariant lift of the two generators. Smoothness
of the special complete intersection and properness
then give a smooth lift exactly as in A19b.

For a fully explicit ambient model one may use
$E:y^2=x^3-x$, over a number field containing its
three-torsion and then localize or complete at a
prime above any $p>3$. At primes $p\equiv3\pmod4$
the ambient $E^6$ is supersingular. The argument above
works over all $p>3$, after the indicated finite
field and trait extensions, without relying on
supersingularity or on an extension of any unknown
special algebraic cycle.
\end{proof}

\paragraph{Role in the programme.}
A20 supplies the geometric hypotheses, ranks and
two reflection configurations for A21; it does not
assert their monodromy conclusion without that proof.
Together with A19b and A21, the odd sign and standard
multiplicity modules can account for three residual
potential Tate directions. With the optional centres
there is already a nonzero positive regular span of
rank six, but the residual multiplicity is still
three; A15 alone cannot settle that residual space.
These are genuine original Lefschetz pencils in
the same ample system, so their axes apply to
cote 44, Theorem 5, scan 004 and its primitive
qualification in scan 008. No assertion of the
universal fixed-axis sign or of algebraicity of
the three remaining potential directions is made.

% END work/scan_w/affine_s3_axis_geometry.tex

% BEGIN work/axis/s3_multiplicity_monodromy.tex
% C44_POS /root; GPT-6 Astra Pro; effort not exposed.
% A19, A20 and this entire proof have passed independent whole-proof readings.
\section{Three residual Tate directions separated by an actual symmetric-group action}
\label{sec:c44-a21}

The affine action and its smooth-pencil existence assertions in the following
statement are the explicit geometric construction C44-A20. Its checked
equivariant jet argument supplies the geometric input for the moment calculation.

\begin{theorem}[C44-A21: monodromy on the three rational multiplicity spaces]
Let $k$ be a finite field of characteristic $p>3$ and let $G=S_3$ act on
$X=E^6$ by the affine action of C44-A20.  Its action on the tangent space has
representation $2\varepsilon\oplus2\mathrm{Std}$, where $\varepsilon$ is the
sign representation and $\mathrm{Std}$ the rational two-dimensional standard
representation.  An element of order three acts freely; the loci fixed by the
three involutions are disjoint unions of abelian surfaces.  Fix a sufficiently
large power $L$ of the $G$-linearized ample bundle in that construction.

Allow either no prescribed centre, or the single free $G$-orbit of six
abelian-surface fibres $Z_g$ specified there.  Write $Z$ for their union in
the second case.  Let
\[
 \mathcal E=H^0(X,I_Z\otimes L\otimes\mathrm{Std})^G,
 \qquad
 \Omega\subset\mathcal E                                       \tag{A21.1}
\]
be the nonempty smooth open of C44-A20: a parameter $a$ gives a smooth axis
$B_a$ cut out by its two component sections, and their original pencil on
$X$ is a genuine Lefschetz pencil.  In the case without centres, set $K=0$.
With the six centres, let $K$ be their primitive projected cycle span.  Its
signed Gram matrix is positive definite by C44-A06; it is a constant copy
of the regular rational representation of $G$.

Let $V$ be the middle Gysin kernel, with its pairing $q(x,y)=\int_{B_a}xy$
and twist $(2)$, and put $W=K^\perp\cap V$.  Over every coefficient field
$\mathbf Q_\ell$, $\ell\ne p$, decompose
\[
 W=(\mathbf1\otimes M_1)\perp(\varepsilon\otimes M_\varepsilon)
                      \perp(\mathrm{Std}\otimes M_{\mathrm{Std}}).
                                                                  \tag{A21.2}
\]
The multiplicity sheaves are nondegenerate orthogonal lisse sheaves.
For $L$ sufficiently large their ranks exceed eight, and each has full
orthogonal geometric monodromy.  This is an unconditional assertion; it
uses no Tate conjecture, cycle lifting, or 032 input.
\end{theorem}

\begin{proof}
We give the moment argument directly over the finite field.  All schemes,
centres and group elements are defined over $k$, after the finite extension
allowed in C44-A20.  For the calculation, remove the twist $(2)$ and give
the multiplicity sheaves their natural weight four.  We write $Q$ for the
cardinality of a finite extension $k'/k$.

The splitting (A21.2) and its forms are actual lisse constructions.
The finite group algebra is semisimple over every $\mathbf Q_\ell$,
including $\ell=2,3$.  Its three rational irreducible representations are
absolutely irreducible and self-dual with symmetric forms: use norm one
on the one-dimensional representations and the usual dot product on
$\{(x_1,x_2,x_3):\sum x_i=0\}$ for the standard representation.
The $G$-invariant nondegenerate pairing on $W$ is orthogonal between
distinct isotypic summands, and on each summand it is the chosen form
on $\rho$ tensor a symmetric form on $\operatorname{Hom}_G(\rho,W)$.
Schur's lemma and nondegeneracy show that the latter form is itself
nondegenerate.  These operations commute with the family and Frobenius.
The fixed positive Gram projector onto $K$, when present, also commutes
with $G$ and the family.  This justifies all the orthogonal multiplicity
sheaves used below.

\paragraph{The free quotient and four-point evaluation.}
Remove from $X$ the loci fixed by nonidentity elements and also $Z$, if
present, and call the resulting smooth open $U$.  The order-three elements
have no fixed points, so the only other removed sets are surfaces.  The
quotient
\[
                            Y=U/G                               \tag{A21.3}
\]
is smooth and geometrically connected of dimension six, and $U\to Y$ is
a geometrically connected finite etale $G$-torsor.  The equivariant vector
bundle $L\otimes\mathrm{Std}$ descends to a rank-two vector bundle on $Y$.
A point $y\in Y(k')$ has a Frobenius conjugacy class $h_y\subset G$.

For every set of at most four distinct points of $Y(k')$, evaluation of
$\mathcal E\otimes k'$ onto the corresponding four descended vector-bundle
fibres is surjective.  To see this uniformly, their inverse images in $U$
are at most twenty-four distinct geometric points.  For $L$ sufficiently
large, sections of $I_ZL$ evaluate surjectively at every such collection.
One may first twist until $I_Z$ is globally generated and then multiply
by products of hyperplanes separating at most twenty-four points, as in
C44-A18.  The bound is independent of the positions of these points.
Tensor by $\mathrm{Std}$ and average over $G$.  Since $p>3$, taking
$G$-invariants is exact, and evaluation remains surjective onto the
invariants of the orbit-fibre target.  These invariants are precisely the
descended fibres.  The map is defined over $k'$; surjectivity after
algebraic closure implies surjectivity over $k'$.

For a uniformly random $a\in\mathcal E(k')$, let $I_y(a)$ be one if the
section vanishes at $y$, and zero otherwise.  Thus $I_y$ is Bernoulli with
parameter $v=Q^{-2}$, and any collection of at most four distinct such
variables is independent.  This remains true when the fibre of the
$G$-torsor at $y$ is nonsplit: the descended vector-bundle fibre is still
a two-dimensional $k'$-vector space.

\paragraph{The character-weighted trace function.}
Let $\rho$ be any of the three rational irreducible representations
$\mathbf1,\varepsilon,\mathrm{Std}$, with character $\chi_\rho$ and dimension
$d_\rho=1,1,2$, respectively.  Define, for every section, including singular
ones,
\[
 T_\rho(a,k')=\sum_{y\in Y(k')}\chi_\rho(h_y)(I_y(a)-Q^{-2}).
                                                                  \tag{A21.4}
\]
For smooth parameters in $\Omega$ this has the uniform approximation
\[
 T_\rho(a,k')=\operatorname{Tr}(F_{k',a}\mid M_\rho^{\mathrm{un}})
                                  +O(Q^{3/2}).                   \tag{A21.5}
\]
Here and below the constants may depend on the fixed large $L$, but not
on $k'$ or on the parameter in the smooth open.

We verify this approximation with the ambient contribution included.
For a $G$-stable scheme $T$, put $N_g(T)=\#\operatorname{Fix}(gF_{k'};T)$.
On a free $G$-orbit stable under Frobenius, the character projector gives
\[
 \frac1{6}\sum_{g\in G}\chi_\rho(g^{-1})N_g(T)
       =\sum_{y\in(T/G)(k')}\chi_\rho(h_y).                       \tag{A21.6}
\]
Indeed identify an orbit with $G$: Frobenius acts by a right translation,
and the sum over all fixed pairs $(g,x)$ is six times the character of
that translation.  All the characters here are real and invariant under
inversion, so the choice of inverse Frobenius convention gives the same
identity.  Applying (A21.6) to the zero locus in $U$ and to $U$ identifies
(A21.4) with the character average of
$N_g(B_a\cap U)-Q^{-2}N_g(U)$.

Weak Lefschetz for the two ample sections, its dual Gysin assertion, and
hard Lefschetz on $X$ give, for each $g$,
\[
 N_g(B_a)-Q^{-2}N_g(X)
       =\operatorname{Tr}(gF_{k',a}\mid V^{\mathrm{un}})
                                      +O(Q^{3/2}).               \tag{A21.7}
\]
To check the degrees: $H^i(B_a)$ for $i>4$ identifies by Gysin with
$H^{i+4}(X)(2)$, so its contribution cancels against the corresponding
term of $Q^{-2}N_g(X)$.  In degree four, the ambient image of $H^4(X)$
has the same Frobenius and $G$-action as $H^8(X)(2)$ by $L_h^2$; its
removal leaves $V^{\mathrm{un}}$.  All remaining terms have absolute
value $O(Q^{3/2})$.  The finite-order action of $g$ does not affect this
bound, and the smooth family has fixed cohomological ranks.  Equivalently
one may apply the ordinary trace formula to the finite-field twist by
$g$, which is smooth projective.

Each removed involution-fixed surface meets $B_a$ in a smooth curve in
the chosen open.  Its contribution to $N_g(B_a)$ is $O(Q)$, whereas that
to $Q^{-2}N_g(X)$ is $O(1)$.  All such estimates are uniform.  Their
removal therefore fits within the error in (A21.7).

If the six centres are present, their quotient is one abelian surface
$Z_0$, and its torsor is split because all six components are defined
individually over $k$.  Their subtraction from the character average is
\[
 d_\rho(1-Q^{-2})\#Z_0(k')
                         =d_\rho Q^2+O(Q^{3/2}).                 \tag{A21.8}
\]
This is exactly the leading trace of the $d_\rho$ constant weight-four
lines in the $\rho$-multiplicity of the regular representation $K$.
Subtracting them leaves $M_\rho$.  With no centres this step is empty.
Finally, averaging the right side of (A21.7) against $\chi_\rho/6$
extracts the multiplicity trace, with no extra factor $d_\rho$:
character orthogonality gives
$\frac16\sum_g\chi_\rho(g^{-1})\operatorname{Tr}(gF\mid V)
=\operatorname{Tr}(F\mid\operatorname{Hom}_G(\rho,V))$.
This proves (A21.5).

\paragraph{An exact fourth moment for each multiplicity sheaf.}
Put $w_y=\chi_\rho(h_y)$ and $S_j=\sum_yw_y^j$.  The finite-etale
character sheaf on $Y$ and the compact-support trace formula imply
\[
                        S_2=Q^6+O(Q^{11/2}),\qquad S_4=O(Q^6).
                                                                  \tag{A21.9}
\]
Here the leading coefficient of $S_2$ is
$\dim(\rho\otimes\rho)^G=1$; the connected geometric torsor has full
geometric group $G$, and every representation is absolutely irreducible
and orthogonal.  All other compact-support cohomology has weight at most
eleven.  The second bound also follows immediately from boundedness of
the three integer characters and $\#Y(k')=O(Q^6)$.  This is the usual
finite-torsor form of the Weil-II weight estimate, not a supposition of
equidistribution without an error term.

Four-wise independence, applied to (A21.4), gives the exact identity
\[
 \begin{split}
 \frac1{\#\mathcal E(k')}\sum_a T_\rho(a,k')^4
     &=3v^2(1-v)^2(S_2^2-S_4)+\mu_4S_4,\\
 \mu_4&=v(1-v)^4+(1-v)v^4,
 \qquad v=Q^{-2}.
 \end{split}                                                     \tag{A21.10}
\]
The terms with a singleton index vanish; the remaining index patterns
are the two pairs and the quadruple.  The character values may be
negative, but they occur only through their squares or fourth powers in
this formula.  By (A21.9),
\[
 Q^{-8}\frac1{\#\mathcal E(k')}\sum_aT_\rho(a,k')^4
                                  =3+O(Q^{-1/2}).                 \tag{A21.11}
\]
For example, the fourth-diagonal contribution is $O(Q^{-4})$ after
normalization.

The geometric monodromy group of $M_\rho$ is reductive, by purity.
On the smooth open, (A21.5) differs from the true trace by $O(Q^{3/2})$,
while both functions are $O(Q^2)$.  Thus normalized fourth moments differ
by $o(1)$.  Katz's Lemmas 2.1.4--2.1.6, printed pp.~540--541, identify
the limsup on that open with the fourth monodromy moment and bound it
above by (A21.11), since adding the omitted nonnegative fourth powers
can only increase the sum.  Hence
\[
                  M_4(G_\rho,M_\rho)\leq3.                      \tag{A21.12}
\]
The three independent orthogonal contraction tensors give the reverse
inequality when the rank exceeds eight.  Therefore the moment is
exactly three.  No trace bound on singular fibres is required.

\paragraph{Reflections on the actual multiplicity spaces.}
The equivariant jet construction in C44-A20 gives a degeneration with
one free orbit of six ordinary double points, away from every fixed
locus and centre, and no other singularity.  A transverse one-parameter
smoothing has regular total space at all six nodes.  SGA~7~II,
Expose~XV, Theorem~3.4(i),(iii), pp.~195--196, supplies mutually
orthogonal vanishing cycles of square two, with one inertia element
acting as minus one on their span and as one on its orthogonal
complement.  The span is the regular representation of $G$.
Consequently this inertia acts as a reflection on the one-dimensional
vanishing part of $M_1$.

There is also a degeneration with one orbit of three nodes, each fixed
by an involution, and no other singularity.  At one node the even
component of the section has zero value and zero first derivatives
along the fixed surface, while the odd component has nonzero normal
derivative.  Eliminating the latter equation leaves an ordinary
quadratic singularity in five variables, with involution eigenspaces
of dimensions two and three.  The involution acts by its determinant,
namely minus one, on the local vanishing line, as established in the
local quadratic calculation of C44-A20.  Thus the three vanishing
cycles form
\[
                  \operatorname{Ind}_{C_2}^{S_3}\varepsilon
                          =\varepsilon\oplus\mathrm{Std}.       \tag{A21.13}
\]
The same regular smoothing and SGA~7 formula make inertia act as minus
one on this span.  It is a reflection on both $M_\varepsilon$ and
$M_{\mathrm{Std}}$.

These cycles lie in $W$, not merely in the full middle cohomology.
Every ambient class and every prescribed cycle is fixed by inertia.
The nonzero orthogonal node cycles are linearly independent; the
Picard--Lefschetz formula therefore forces each of them to be orthogonal
to the ambient image and to $K$.  The involution character and the
regular-orbit decompositions above hence apply inside $W$ itself.

Each multiplicity monodromy is orthogonal, reductive, of rank greater
than eight, has fourth moment three, and contains a reflection.
Katz's Theorem~1.4.6, printed p.~535, now gives the full orthogonal group
in all three cases.  The rank growth follows from C44-A20's equivariant
Euler calculation.  Passing to the smaller open imposing the original
Lefschetz-pencil condition does not change geometric monodromy.
\end{proof}

\begin{corollary}[C44-A21a: actual closed axes with three residual Tate directions]
For each sufficiently large polarization in C44-A20 there is a
Zariski-dense set of closed parameters in $\Omega$ for which the entire
rational algebraic axis space $A_B$ is positive definite.  Without the
six centres, the total root-of-unity multiplicity of normalized variable
Frobenius is exactly three.  With the six centres, their known positive
span has dimension six and the total root-of-unity multiplicity on its
orthogonal complement is exactly three; the full variable multiplicity
is nine.  In particular the latter examples have $\dim_{\mathbf Q}A_B\geq6$.
All these assertions are unconditional and concern genuine original
Lefschetz pencils over finite fields.
\end{corollary}

\begin{proof}
The three rational multiplicity Frobenius polynomials are independent of
$\ell$.  Indeed the traces of $gF^n$ are rational and independent of
$\ell$ by the finite-field twists used in C44-A19, and character averaging
followed by Newton identities gives the assertion.  The known regular
cycle summand, if present, is removed by its rational Gram projector.

Let $H$ be the joint geometric monodromy closure in the product of the
three orthogonal groups.  C44-A21 shows that each projection is onto;
therefore each projection of $H^0$ is onto the corresponding special
orthogonal group.  For every even-dimensional factor, the condition of
having no root of unity is detected by a nonzero regular polynomial on
that special orthogonal group; for every odd-dimensional factor, the
condition of having a unique simple root of unity is likewise detected
by the polynomial of C44-A16.  One common exponent suffices: if $D$ is
the largest of the three ranks, use
$M=\operatorname{lcm}(1,\ldots,2D^2)$, since a root of unity of degree at
most $D$ has order at most $2D^2$.

Each inverse image of this nonempty open is a nonempty open in the
connected group $H^0$.  Their intersection is nonempty.  The finite
quotient and constant-field Chebotarev argument of C44-A12 and A16 thus
selects a Zariski-dense set of closed parameters meeting all three
conditions simultaneously.  Independence of the three monodromy factors
is not assumed; surjectivity of the separate projections suffices.

The parity computation of C44-A20 is
\[
 \begin{array}{c|ccc}
 &\operatorname{rank}M_1&\operatorname{rank}M_\varepsilon
                                  &\operatorname{rank}M_{\mathrm{Std}}\\\hline
 Z=\varnothing&0&1&1\\
 \#\{Z_g\}=6&1&0&1
 \end{array}\pmod2.                                               \tag{A21.14}
\]
Thus the residual potential Tate multiplicity on $W$ is, respectively,
$0+1+2\cdot1=3$ or $1+0+2\cdot1=3$.

Choose an involution $s\in G$.  The four rational self-adjoint,
mutually orthogonal idempotents
\[
 e_1,\quad e_\varepsilon,\quad
       e_{\mathrm{Std}}\frac{1+s}{2},\quad
       e_{\mathrm{Std}}\frac{1-s}{2}                              \tag{A21.15}
\]
sum to one.  They are actual correspondences extending with the affine
action and the equivariant axis over the mixed-characteristic model of
C44-A20.  Their images on $W$ have Frobenius polynomials respectively
those of $M_1,M_\varepsilon,M_{\mathrm{Std}},M_{\mathrm{Std}}$.
Every such image has at most one residual potential Tate direction.
The known cycles split into the same orthogonal blocks, with positive
Gram matrices.  C44-A19 applies its common-projector determinant proof
to these four extending blocks, so their full algebraic subspaces are
positive.  Their orthogonal sum is all of $A_B$.

The six known classes, when imposed, are independent by their positive
definite Gram matrix, proving the last dimension assertion.  No
algebraicity is asserted for any of the three additional potential Tate
directions, and no Tate conjecture is needed for positivity of those
which are algebraic.
\end{proof}

\paragraph{Scope, sources and the main gate.}
These axes have three residual root-of-unity directions even after the
explicit known positive span is removed.  They therefore fall outside
the multiplicity-one sufficient condition C44-A15; the new reason for
their sign is the four actual extending rational projectors.  The
construction removes a nonempty further class of axes from the original
QH9-axis gate.  It does not prove that every axis has such projectors,
prove a general Tate conjecture, or exhibit a wrong-sign cycle.

The monodromy inputs are Katz, \emph{Larsen's alternative, moments, and
the monodromy of Lefschetz pencils} (2004), Theorem~1.4.6 p.~535 and
Lemmas~2.1.4--2.1.6 pp.~540--541, and SGA~7~II, Expose~XV,
Theorem~3.4 pp.~195--196.  The finite-torsor estimate (A21.9) is the
compact-support trace formula with Deligne's Weil-II weight theorem,
\emph{La conjecture de Weil II}, Theorem~3.3.1 and its compact-support
Corollary~3.3.4, printed p.~206, applied to the finite-monodromy weight-zero character
sheaf.  The torsor averaging and exact fourth-moment identities are
proved above rather than borrowed as an unstated equivariant theorem.

In Grothendieck's language this supplies the signed axis term of
cote~44, Th.~5, scan~004, for these closed special pencils, with the
primitive margin of scan~008 enforced.  The separate ambient abelian
term in H9-T7, and hence any assembled full-height conclusion using it,
remain conditional on the unrefereed 032/Milne 2026 preprints.  The
present axis sign itself is unconditional.  Its setting is the actual
relative correspondence algebra and SGA~7 monodromy, not an abstract
choice of a positive form on a numerical-motive summand.

% END work/axis/s3_multiplicity_monodromy.tex

% BEGIN work/axis/axis_realization_blocks.tex
% C44_POS /root/axis_independent. Extends the common-kernel test, not a liftability theorem.
\section{The crystalline test with several symmetry-separated Tate directions}
\label{sec:c44-a18d}

\begin{corollary}[C44-A18d: a common kernel for finite-group blocks]
Let $B\subset X/\mathbf F_q$ and $P$, $\mathscr C$ be as in
C44-A18b. Let a finite constant group $G$ act on this embedded
complete intersection, preserving the ample class. Suppose
$e_1,\ldots,e_b\in\mathbf Q[G]$ are self-adjoint pairwise
orthogonal idempotents summing to one, where the involution
on the group algebra sends $g$ to $g^{-1}$. Let $K$ be a
$G$-stable rational algebraic subspace of the variable
cohomology at one initial coefficient prime, and suppose
its rational intersection form is nondegenerate and each
$K_i=e_iK$ is nondegenerate. A positive definite $K$, including
$K=0$, satisfies this assumption.

Use actual cycles $P_*e_{i*}P_*C$ to represent each $K_i$,
choose a rational basis in every block, and extend the finite
field to define all these cycles. Suppose that in each
orthogonal residual block
\[
       W_{i,\ell}=e_iV_\ell\cap K_{i,\ell}^{\perp}
\]
the total normalized Frobenius root-of-unity multiplicity
is at most one. Then all rational crystalline and
$\ell$-adic cycle-class kernels coincide on $\mathscr C$.
The necessary filtration test of C44-A18c therefore holds
for a projected class along every specified embedded
polarized lift of $B\subset X$.

The finite-group action and its projectors need not extend
along the lift used for that necessary test. Their role in
the present assertion is on the special fibre. In particular
the assertion applies to the C44-A21 axes, including those
with residual potential Tate multiplicity three. It does
not assert that a projected class on any such axis is
nonliftable.
\end{corollary}

\begin{proof}
The group action commutes with Gysin, restriction and
$L_h^2$, and hence with its unique inverse in every
cohomological realization. Thus it commutes with the
cohomological action of $P$. All the stated idempotent
and adjoint identities hold in rational crystalline
cohomology as well as in every $\ell$-adic realization.
The actual projected basis cycles have the same rational
nondegenerate Gram matrices in all these realizations;
they supply common nondegenerate subspaces $K_i$ and
orthogonal residual spaces.

The Frobenius polynomials of the blocks are common rational
polynomials also in crystalline cohomology. Indeed, for
each finite-order $g\in G$ and each $n>0$, the finite-field
twist by $g$ is smooth projective and its Frobenius acts
as $gF^n$, with the inverse cocycle chosen if the Frobenius
convention requires it. Katz--Messing, Theorem~1,
printed p.~74, applies to the twist and identifies its
degreewise characteristic polynomials in crystalline and
$\ell$-adic cohomology. Subtract the ambient traces, then
take the rational linear combination specified by $e_i$.
Since $e_i$ commutes with Frobenius,
\[
  \operatorname{Tr}(F^n\mid e_iV)
       =\operatorname{Tr}(e_iF^n\mid V).
\]
Newton identities give the same rational block polynomial
in all realizations. Removing $(T-1)^{\dim K_i}$ preserves
this assertion. This is the finite-twist argument of
C44-A19b with Katz--Messing also included; no independence
conjecture for cycle classes has been used.

If $Z\in\mathscr C$ has zero class in one realization,
set $Z_i=P_*e_{i*}Z$. Its class in that realization is
zero, so all its pairings with the actual basis cycles
of $K_i$ and its self-pairing vanish in $\mathbf Q$.
In every other realization its class lies in $W_i$,
has finite normalized Frobenius orbit, and has square
zero. The simple potential-Tate-line argument of
C44-A18b forces that class to vanish. Summing over $i$
gives $\operatorname{cl}(Z)=0$ in every realization,
because $\sum_i P e_i$ acts as the identity on $V$.
This proves the common-kernel assertion.

The proof of C44-A18c uses only this common-kernel
assertion on the special projected cycle group, and
an algebraic Lefschetz inverse on the generic abelian
variety. It does not use an extension of the special
finite-group action. The same proof therefore gives
the stated filtration test along any embedded polarized
lift under consideration.
\end{proof}

\paragraph{Scope in Grothendieck's programme.}
For cote~44, Th.~5, scan~004 and the primitive instruction
on scans~008--009, this supplies a representative-independent
crystalline necessary condition on the actual axis cycle
classes in the new three-direction families. The sign of
those families is separately proved by C44-A19--A21.
Failure of the filtration condition for every possible
lift remains unproved, so route (a) of QH9-axis is not
closed by this corollary.

% END work/axis/axis_realization_blocks.tex

% BEGIN work/axis/axis_primitive_semiregularity.tex
% C44_POS /root/axis_independent. Coherent calculation only; arithmetic lifting remains open.
\section{Primitive coherent semiregularity of a prescribed smooth surface}
\label{sec:c44-a22}

\begin{proposition}[C44-A22: surjectivity from the coherent Gysin kernel]
Let $X$ be an abelian sixfold over an algebraically closed field
of arbitrary characteristic, let $H$ be ample, and let
$S\subset X$ be a fixed smooth projective surface. For every
sufficiently large integer $t$, and for every smooth complete
intersection $\kappa:B\hookrightarrow X$ of two members of
$|H^t|$ containing $S$, put $N=N_{S/B}$ and
\[
 V^{3,1}_{\mathrm{coh}}
  =\ker\left(\kappa_*:
          H^1(B,\Omega_B^3)\longrightarrow H^3(X,\Omega_X^5)\right).
 \tag{A22.1}
\]
Here $\kappa_*$ is the coherent Hodge-cohomology Gysin map.
Restriction followed by the natural conormal exterior quotient
induces a surjection
\[
 V^{3,1}_{\mathrm{coh}}
       \longrightarrow H^1(S,N^*\otimes K_S).
 \tag{A22.2}
\]
Consequently its dual is an injection
\[
 H^1(S,N)\lhook\joinrel\longrightarrow
                     (V^{3,1}_{\mathrm{coh}})^*.
 \tag{A22.3}
\]
This is the coherent primitive component of the usual dual
description of the semiregularity map. No degeneration of a
Hodge--de Rham spectral sequence or mixed-characteristic
obstruction theorem is assumed in this statement.
The bound on $t$ depends on $(X,H,S)$ and is uniform in
the smooth complete intersection containing $S$.
\end{proposition}

\begin{proof}
Write $L=H^t$, $I=I_{S/B}$, and $F=\Omega_B^3$.
Choose a nonzero invariant volume form on $X$; its only role
is to trivialize $K_X$ and identify $\Omega_X^5$ with $T_X$.
Adjunction gives $K_B=L^2|B$. The tangent sequence, tensored
with $K_B$, gives
\[
 0\longrightarrow F\longrightarrow (L^2|B)^{\oplus6}
              \longrightarrow (L^3|B)^{\oplus2}
              \longrightarrow0.
 \tag{A22.4}
\]
Its terms are locally free. Tensoring with $I$ is therefore
still exact.

\emph{The required vanishing is uniform in $B$.}
The twisted Koszul resolutions on $X$ are
\[
 \begin{split}
 0&\longrightarrow\mathcal O_X\longrightarrow L^{\oplus2}
                   \longrightarrow L^2
                   \longrightarrow\kappa_*(L^2|B)\longrightarrow0,\\
 0&\longrightarrow L\longrightarrow (L^2)^{\oplus2}
                   \longrightarrow L^3
                   \longrightarrow\kappa_*(L^3|B)\longrightarrow0.
 \end{split}
 \tag{A22.5}
\]
Ample line bundles on an abelian variety have no higher
coherent cohomology, in every characteristic. Thus
\[
 H^j(B,L^3|B)=0\ (j>0),\qquad
 H^j(B,L^2|B)\simeq H^{j+2}(X,\mathcal O_X)\ (j>0).
 \tag{A22.6}
\]
For $t$ sufficiently large, Serre vanishing on the fixed
$S$ and fixed ideal $I_{S/X}$ gives
\[
 \begin{gathered}
 H^j(S,L^a|S)=0\quad(j>0,\ a=2,3),\qquad
 H^1(X,I_{S/X}L^a)=0\quad(a=2,3),\\
 H^2\bigl(S,\Omega_S^1\otimes K_S^{-1}\otimes L^2|S\bigr)=0.
 \end{gathered}
 \tag{A22.7}
\]
In particular $H^0(B,L^a|B)\to H^0(S,L^a|S)$ is
surjective: already the sections from $H^0(X,L^a)$
surject onto the latter space. The restriction exact
sequences and (A22.6)--(A22.7) now give
\[
 \begin{split}
 H^1(B,IL^3)&=H^2(B,IL^3)=0,\\
 H^j(B,IL^2)&\xrightarrow{\sim}H^j(B,L^2|B)
                                  \quad(j=1,2).
 \end{split}
 \tag{A22.8}
\]
Here and below a line bundle next to $I$ means its
restriction to $B$.

\emph{Restriction of the differential forms is surjective.}
Apply cohomology to (A22.4) and to its tensor product
with $I$. By (A22.6) and (A22.8), the commutative map
\[
 H^2(B,IF)\longrightarrow H^2(B,F)
 \tag{A22.9}
\]
identifies with the isomorphism
$H^2(B,IL^2)^{\oplus6}\to H^2(B,L^2|B)^{\oplus6}$.
It is therefore an isomorphism. The restriction sequence
$0\to IF\to F\to F|S\to0$ implies that
\[
                  H^1(B,F)\longrightarrow H^1(S,F|S)
                     \quad\hbox{is surjective}.
 \tag{A22.10}
\]

The rank-two conormal sequence for $S\subset B$, on
taking third exterior powers, has the exact quotient
\[
 0\longrightarrow\det(N^*)\otimes\Omega_S^1
       \longrightarrow F|S\longrightarrow N^*\otimes K_S
       \longrightarrow0.
 \tag{A22.11}
\]
By adjunction its left term is
$L^2|S\otimes K_S^{-1}\otimes\Omega_S^1$.
The last vanishing in (A22.7) proves that
$H^1(S,F|S)\to H^1(S,N^*\otimes K_S)$ is surjective.
Together with (A22.10), this already proves surjectivity
from the full $H^1(B,F)$.

\emph{The ambient quotient can be killed without changing the target.}
Let
\[
  \theta:H^1(B,F)\longrightarrow H^1(B,L^2|B)^{\oplus6}
 \tag{A22.12}
\]
be the map in (A22.4). By (A22.8), the image of
$H^1(B,IF)\to H^1(B,F)$ maps onto the entire target of
$\theta$: in the corresponding ideal-tensored sequence
$H^1(B,IF)\to H^1(B,IL^2)^{\oplus6}$ is surjective,
and the right vertical map is an isomorphism. Every
class from $H^1(B,IF)$ restricts to zero on $S$.

For precision, the coherent Gysin map in (A22.1) is
$\theta$ followed by the double Koszul boundary
\[
 H^1(B,L^2|B)\otimes H^0(X,\Omega_X^5)
        \xrightarrow{\sim} H^3(X,\Omega_X^5).
 \tag{A22.13}
\]
The boundary is an isomorphism by (A22.5)--(A22.6).
This identification follows directly from the local
formula, which also fixes the compatibility needed
here. If $B$ is locally cut out by $f_0,f_1$, the
adjunction injection in (A22.4) sends a lifted
$3$-form $\widetilde\omega$ to
$\widetilde\omega\wedge df_0\wedge df_1$ in
$\Omega_X^5\otimes L^2|B$. The double boundary sends it to
the local-cohomology class represented by
\[
       \frac{\widetilde\omega\wedge df_0\wedge df_1}{f_0f_1}.
 \tag{A22.14}
\]
This is exactly the coherent Gysin construction:
Stacks Project, Section~50.23, Tag~0G82, Remark~50.23.1
and Remark~50.23.7. A simultaneous change of orientation
would multiply both maps by the same sign and would
not change their kernel.

Now take a target class in $H^1(S,N^*\otimes K_S)$,
and lift it to $v\in H^1(B,F)$ using (A22.10)--(A22.11).
Choose $u$ in the image of $H^1(B,IF)$ with
$\theta(u)=\theta(v)$, as is possible by (A22.12).
Then $v-u$ is in the Gysin kernel by (A22.13), and
has the same image on $S$ as $v$. This proves (A22.2).
Serre duality on the smooth surface identifies
$H^1(S,N^*\otimes K_S)^*$ with $H^1(S,N)$ and gives
(A22.3).
\end{proof}

\paragraph{What the proposition does and does not supply.}
The dual-of-restriction description used here is the algebraic
description of semiregularity; compare Dan--Kaur,
\emph{Semi-regular varieties and variational Hodge conjecture},
\emph{Comptes Rendus Math\'ematique} 354 (2016), \S2.1,
printed p.~298. The proof above supplies its own global
restriction argument and does not appeal to a
characteristic-zero variational theorem. The ample-vanishing
input is Milne, \emph{Abelian Varieties}, Chapter~I,
Theorem~11.1(c), printed p.~54 in the author notes.

This calculation applies to the smooth centres in C44-A14
and to the high-degree genuine axes containing them whose
existence is proved in C44-A06 and whose closed fibres are
selected in C44-A18. It shows that the coherent obstruction
map remains injective after its target is restricted to the
dual of the coherent Gysin kernel; discarding the ambient
quotient loses no coherent normal obstruction in this range.

An additional theorem would be needed to conclude that a
rational crystalline primitive class lying in
$\operatorname{Fil}^2$ along an arbitrary possibly ramified
mixed-characteristic lift forces an all-orders lift of the
given surface. Such a theorem has not been proved here.
In particular no passage from a rational filtered class to
an integral deformation of this particular support is being
made. C44-A18b--A18d repair the separate issue of changing
cycle representatives; they do not supply this obstruction
compatibility or a converse to the filtration test.

For cote~44, Th.~5, scan~004 and the primitive warning on
scans~008--009, the proposition isolates a precise remaining
deformation step in route (a) of QH9-axis. It neither proves
that the primitive class of the nonliftable surface is itself
nonliftable nor settles the universal axis sign.

% END work/axis/axis_primitive_semiregularity.tex

% BEGIN work/scan_w/higher_dimensional_s3_axes.tex
% C44_POS /root/scan_w; GPT-6 Astra Pro; reasoning effort not exposed.
% Dimensional extension of the independently checked A19--A21 mechanism.
\section{The symmetric-group construction in every higher even dimension}
\label{sec:c44-a23}

The following is a dimensional generalization of the checked construction
C44-A19--A21.  It is not a separately staged progress-clock event.  We give
the changed geometry, signs, moment powers and representation parities;
the common-projector determinant theorem remains exactly C44-A19.

\begin{theorem}[C44-A23: higher-dimensional symmetry-separated axis signs]
Let $r\geq2$, let $k$ be a finite field of characteristic $p>3$, and let
$E/k$ be an elliptic curve.  After a finite extension choose
$0\ne T\in E[3](k)$ and put
\[
 n=2r+2,\qquad a=2r-2,\qquad
 X=E_{\rm sign}^{\,a}\times E_{\rm std}^{\,2}\times E_{\rm std}^{\,2}.
 \tag{A23.1}
\]
On each standard pair use
$R(P,Q)=(Q,-P-Q)$ and $S(P,Q)=(Q,P)$.
Let the order-three generator $\rho$ translate the first sign coordinate
by $T$, act trivially on the other sign coordinates, and act by $R$ on
both standard pairs.  Let the involution $\sigma$ negate every sign
coordinate and act by $S$ on both standard pairs.  These define an
affine $G=S_3$ action.

Set
\[
 M=\mathop{\boxtimes}_{j=1}^{n}\mathcal O_E(3[0]),\qquad
 H=\left(\bigotimes_{g\in G}g^*M\right)^{\otimes2},\qquad L=H^t.
 \tag{A23.2}
\]
For all sufficiently large $t$, the equivariant standard section pairs
have a nonempty smooth parameter open whose members cut out smooth
axes $B\subset X$ of dimension $2r$, and whose original pencils are
genuine Lefschetz pencils.  One may either impose no centre, or impose
the following six centres.  Let $S_0=E^r$ be the first $r$ sign factors,
and choose a point $c$ in the remaining factors with free orbit under
the induced linear $G$ action.  After another finite extension put
\[
                       Z_g=S_0\times\{gc\}\quad(g\in G).
 \tag{A23.3}
\]
These are six disjoint abelian $r$-folds.  Their primitive projected
classes span a positive definite constant regular representation $K$;
without centres set $K=0$.

Write $V=\ker\kappa_*$ for the variable middle cohomology of the axis,
with its twist $(r)$ and signed pairing
$q_B(x,y)=(-1)^r\int_Bxy$, and put $W=V\cap K^\perp$.
The three rational multiplicity sheaves of $W$ for
$\mathbf1,\varepsilon,\mathrm{Std}$ have ranks greater than eight for
large $t$, and each has full orthogonal geometric monodromy.

Define bits $\alpha_r,\beta_r\in\{0,1\}$ by
\[
 (\alpha_r,\beta_r)=
 \begin{cases}
 (0,0),&r\text{ odd},\\
 \left(\binom{r-1}{r/2-2},\binom{r-1}{r/2-1}\right)\bmod2,
       &r\text{ even},
 \end{cases}
 \tag{A23.4}
\]
where an out-of-range binomial coefficient is zero.
There is a Zariski-dense set of closed parameters for which the entire
rational algebraic axis space $A_B$ is positive definite and the total
root-of-unity multiplicity of normalized Frobenius on $W$ is exactly
\[
 \tau(W)=
 \begin{cases}
 \alpha_r+3\beta_r,&K=0,\\
 2-\alpha_r+\beta_r,&\dim K=6.
 \end{cases}
 \tag{A23.5}
\]
With centres, the total multiplicity on $V$ is $6+\tau(W)$ and
$\dim_{\mathbf Q}A_B\geq6$.  These conclusions are unconditional.
They do not assert algebraicity of any residual potential Tate line.
\end{theorem}

\begin{proof}
We verify every dimensional change to C44-A20--A21.

\paragraph{The action, fixed loci, and actual equivariant lifts.}
The identities $R^3=S^2=1$ and $SRS=R^{-1}$ are unchanged.  Negation
conjugates translation by $T$ to translation by $-T$, so the affine
generators satisfy the same relations.  Both order-three elements act
freely.  The fixed locus of $\sigma$ has
\[
                 2^{2a}=2^{4r-4}
 \tag{A23.6}
\]
components, each an abelian surface: the sign coordinates are two-torsion
and both standard pairs lie on their diagonals.  Fixed loci of distinct
involutions are disjoint, since an intersection would contain a point
fixed by an order-three element.  At a fixed point the tangent
eigenspace dimensions are $2$ and $2r$.

The specified norm-square $H$ is very ample, has a natural
$G$ linearization, and has trivial stabilizer character on every fixed
surface.  The last assertion follows from squaring the order-two fibre
character.  The standard representation has one even and one odd
component at an involution.  Thus, at such a point, the even section
has arbitrary value and two arbitrary tangent derivatives, whereas
the odd section has zero value and $2r$ arbitrary normal derivatives.
The corresponding second jets have exactly the invariant parity.

There are at least $r$ sign factors, because $2r-2\geq r$.
The translation by $T$ is inside $S_0$, so the quotient by these factors
carries the stated linear action.  Its free locus is nonempty because
it contains the two standard factors.  Consequently the six fibres
in (A23.3) are disjoint, form a regular orbit, and avoid all involution
fixed loci.  They may all be defined over a finite extension.

Lift $E$ to a smooth elliptic scheme over a mixed-characteristic trait,
and lift its prime-to-$p$ three-torsion section after an allowed finite
extension.  The integer matrices and affine translation extend.  The
\emph{specified} product bundle $M$ extends by the zero section, and
the same norm and square extend $H$.  The point $c$ lifts by smoothness
after completion; its six distinct fibres remain disjoint.  For large
$t$, relative Serre vanishing and averaging by $1/6$ lift the actual
standard section pair, including the prescribed vanishing if the
centres are imposed.  Properness of the nonsmooth locus then gives a
smooth embedded polarized lift of every chosen smooth special axis.
This is precisely the lift needed for C44-A19; it is not a claim that
an arbitrary polarization on a supersingular product extends.

\paragraph{Smooth axes and genuine original pencils.}
The uniform compatible-jet argument of C44-A20 works in dimension $n$:
first globally generate the fixed ideal by an ample twist, then use
products of separating hyperplanes for a bounded number of points and
jets, and average.  Along the smooth prescribed centres, the same
Serre argument applies to their conormal jets.  The bounds depend on
the fixed data, but are uniform in points and orbit configurations.

At a free point the bad complete-intersection jet has two zero values
and a rank-at-most-one $2$-by-$n$ derivative matrix.  Its codimension
is $2+(n-1)=n+1>n$.  On a fixed surface, the bad even derivative has
codimension $1+2=3>2$, and the bad odd derivative, together with the
zero even value, has codimension $1+2r>2$.  Along each $r$-dimensional
centre, the two rows in its $(r+2)$-dimensional normal space fail to
have rank two in codimension $r+1>r$.  These estimates prove the
nonempty smooth-axis open and ensure that every pencil member is
smooth along the centres.

The faithful projective standard action on the pencil has eight
exceptional parameters: two order-three eigenlines and the six
involution eigenlines.  They can all be chosen to give smooth members.
At a free point a critical section has codimension $n+1>n$.
For an involution eigenmember the critical-point codimensions on its
fixed surface are $3$ for the even member and $2r$ for the odd member;
both exceed two.  A member with neither eigencharacter has all $n$
derivatives available, so has critical-point codimension $n+1$ there.
These arguments cover all eight members, even at fixed loci of the
other involutions.

At a nonexceptional parameter, a degenerate critical point imposes
$n+2$ conditions on the $(n+1)$-dimensional point--parameter space.
Two critical points in distinct free orbits impose $2(n+1)$ conditions
on a space of dimension $2n+1$.  If they lie in one free orbit, the
parameter and its translate are distinct; the two independent section
components then both have zero first jet at one point, again imposing
$2(n+1)$ conditions, now on dimension $n+1$.  A fixed-point critical
pair imposes $n+1$ conditions on dimension three.  Each of these bad
loci is avoided.  Smoothness of the base prevents a critical point on
the base; a nondegenerate Hessian and the nonzero pencil derivative
make the remaining critical scheme reduced and finite.  Properness
gives finiteness without any unaccounted boundary points.

This scheme is nonempty.  The first-jet bundle on $X$ has an extension
of $L$ by $L^{\oplus n}$, so the top Chern number for the universal
pencil on $X\times\mathbf P^1$ is
\[
 \int_{X\times\mathbf P^1}(h+\eta)^{n+1}
                =(n+1)\int_Xh^n>0.
 \tag{A23.7}
\]
Thus every singular member has exactly one ordinary double point,
and the original pencil is a genuine Lefschetz pencil.

\paragraph{The two node orbits and their signed norms.}
At a free orbit one may prescribe a rank-one complete-intersection
jet and a nondegenerate remaining quadratic form in $n-1=2r+1$
variables.  This gives exactly one orbit of six nodes.  At an
involution-fixed point, prescribe zero value and derivative for the
even component and a nonzero odd derivative.  Elimination of one odd
variable leaves an invariant quadratic form with eigenspace dimensions
\[
                              (2,2r-1).
 \tag{A23.8}
\]
A nondegenerate such form gives exactly one orbit of three nodes.
In both cases the same bounded-jet independence at other points and
the preceding strict codimension inequalities exclude every other
singularity.  Varying the remaining constant term gives a transverse
smoothing regular at all the nodes; its general member lies in the
smooth-axis open and may also be required to lie in the genuine-pencil
open.

The equivariant etale Morse and even-quadric calculation in C44-A20
applies in every odd number $2r+1$ of quadratic variables.  The
involution determinant in (A23.8) is $(-1)^{2r-1}=-1$; after
compactification it exchanges the two generator families of the
even-dimensional quadric, hence negates its primitive line.  The
three-node vanishing representation is therefore again
$\operatorname{Ind}_{C_2}^{S_3}\varepsilon=\varepsilon\oplus\mathrm{Std}$.
The free orbit gives the regular representation.  SGA~7 supplies
pairwise orthogonal node cycles with ordinary cup square
$2(-1)^r$, hence \emph{signed} square two for $q_B$.  A common quadratic
tame inertia acts by minus one on their span.  It gives reflections
on the trivial and sign multiplicities from the free orbit and on the
standard multiplicity from the three-node orbit.  Ambient classes and
prescribed centre classes are fixed by inertia, so these node cycles
lie in the residual $W$, by their nonzero norms and independence.

\paragraph{The prescribed positive regular span.}
Write $h_0=c_1(H)$ and
\[
 d=\int_{Z_1}h_0^r>0,\qquad
 A_0=(L_{h_0}^2)^{-1}[Z_1]_X,\qquad
 c_0=\int_XA_0^2h_0^2.
 \tag{A23.9}
\]
The inverse is represented by the absolute abelian Chow correspondence
of C44-A15.  All six parallel fibres have the same ambient class.
Their primitive cycles are
$w_g=[Z_g]_B-t^{-2}\kappa^*A_0$.
Their normal bundle in $X$ is trivial, and the two defining divisors
give
$c(N_{Z_g/B})=(1+th_0|Z_g)^{-2}$.
Thus their signed Gram matrix is
\[
        (r+1)t^r d\,I_6-(-1)^r t^{-2}c_0J_6.
 \tag{A23.10}
\]
It is positive definite, for example when
$(r+1)t^{r+2}d>6|c_0|$.  Hence the $w_g$ are independent and form a
constant regular $G$ representation in every realization.  This also
checks the sign change when $r$ is odd.

\paragraph{Fourth moments in the changed dimension.}
Remove the fixed surfaces and the centres from $X$ and take the free
quotient $Y$.  It is a smooth geometrically connected $n$-fold with
connected $G$-torsor.  Four quotient points still require at most
twenty-four geometric points upstairs.  Equivariant evaluation and
descent therefore give the same four-wise independent Bernoulli
variables of parameter $v=Q^{-2}$ as in C44-A21.

For each character $\chi$ define
$T_\chi=\sum_{y\in Y(\mathbf F_Q)}\chi(h_y)(I_y-Q^{-2})$.
The trace identity in C44-A21, now in degree $2r$, is
\[
 T_\chi=\operatorname{Tr}(F\mid M_\chi^{\rm un})
                                  +O(Q^{r-1/2}).
 \tag{A23.11}
\]
Indeed the terms above degree $2r$ cancel by Gysin against the ambient
terms shifted by four, the ambient middle image cancels by $L_h^2$,
and every remaining term has weight at most $2r-1$.  The removed fixed
surfaces meet $B$ in curves, costing only $O(Q)$.  A centre quotient is
one $r$-fold with split torsor, so its subtraction is
$\dim(\chi)Q^r+O(Q^{r-1/2})$, exactly the regular $K$ contribution.
The character average still has factor $1/6$ and no extra dimension
factor.

The finite-character local system gives
$S_2=Q^n+O(Q^{n-1/2})$ and $S_4=O(Q^n)$.
The exact Bernoulli fourth moment remains
\[
                 3v^2(1-v)^2(S_2^2-S_4)+\mu_4S_4.
 \tag{A23.12}
\]
After division by $Q^{4r}$ it tends to three; the diagonal contribution
is $O(Q^{-2r})$.  The trace error in (A23.11) tends to zero after
normalization by $Q^r$.  The nonnegative fourth powers on the singular
parameter locus give the same upper bound on the smooth open.  Purity,
the three orthogonal contraction tensors, the displayed reflections
and Katz's rank-greater-than-eight theorem yield full orthogonal
monodromy in each multiplicity space.  Their ranks grow like a positive
multiple of $t^n$, as the next calculation also shows.

\paragraph{An exact representation-ring parity calculation.}
The equivariant Euler characters of the smooth axis at the three
conjugacy classes $1,\sigma,\rho$ are
\[
 \left((2r+1)\int_Xh^n,
              -2^{4r-4}\int_Fh^2,0\right),
 \tag{A23.13}
\]
where $F$ is any fixed-surface component.  The first entry follows
from $c(T_B)=(1+h)^{-2}$; the second is the Euler characteristic of
the fixed smooth curves; the third is zero because $\rho$ acts freely.
The primary tame fixed-point trace theorem used in C44-A20 applies
unchanged.  All multiplicities of this Euler representation are even:
$\int_Xh^n=n!\chi(L)$, $n\geq6$, and the factor $2^{4r-4}$ gives
the required divisibility in the trivial and sign multiplicities.
Their leading $t^n$ coefficients are positive.

By weak Lefschetz and duality, removing the ambient cohomology changes
this Euler representation by twice the alternating lower cohomology
and by $H^{2r}(X)$.  Consequently the multiplicity parities on $V$
are precisely those on $H^{2r}(X)$.
In the rational representation ring modulo two, write $U=\mathrm{Std}$.
The relations $U^2=1+\varepsilon+U$ and $\varepsilon U=U$ give
$U^4=1+\varepsilon+U$ modulo two.  Since
$H^1(X)=4(r-1)\varepsilon+4U$ as a $G$ representation, its full exterior
algebra has generating series
\[
 \lambda_z H^1(X)
   =(1+\varepsilon z)^{4(r-1)}(1+Uz+\varepsilon z^2)^4
   \equiv(1+z^4)^{r-1}
             \bigl(1+(1+\varepsilon+U)z^4+z^8\bigr)\pmod2.
 \tag{A23.14}
\]
For odd $r$, the coefficient of $z^{2r}$ is zero.  For $r=2s$ it has
sign and standard coefficients $\binom{r-1}{s-1}$; its trivial
coefficient is
\[
 \binom{r-1}{s}+\binom{r-1}{s-1}+\binom{r-1}{s-2}
       \equiv\binom{r-1}{s-2}\pmod2,
 \tag{A23.15}
\]
because $\binom{2s}{s}$ is even for $s\geq1$.
This proves (A23.4): the three multiplicity parities on $V$ are
$(\alpha_r,\beta_r,\beta_r)$.
Removing the regular $K$ subtracts multiplicities $(1,1,2)$, leaving
$(1-\alpha_r,1-\beta_r,\beta_r)$.

\paragraph{Simultaneous Frobenius selection and the full algebraic sign.}
The rational multiplicity Frobenius polynomials are common across
coefficient primes by the finite-order twists of C44-A19.
Each projection of the connected joint geometric monodromy group is
the appropriate special orthogonal group.  The even-rank no-root
selector of C44-A12 and odd-rank simple-root selector of C44-A16 have
nonempty inverse-image opens in this one connected group.  Their
intersection is nonempty.  The same constant-field finite-quotient
Chebotarev argument as in C44-A21 therefore gives a Zariski-dense set
of closed parameters at which all three multiplicities are minimal
at once.  No independence of the three factors is asserted or needed.

The parity list gives (A23.5), remembering that the standard
multiplicity contributes twice to $W$.  The four actual relative
self-adjoint projectors
$e_1,e_\varepsilon,e_U(1+\sigma)/2,e_U(1-\sigma)/2$
each have residual potential Tate multiplicity at most one.
Their embedded polarized lifts have already been constructed, and
the positive $K$ splits into positive orthogonal pieces.
C44-A19 proves the sign on each complete algebraic block, so on their
sum $A_B$.  The positive Gram matrix proves the six-dimensional
lower bound when centres are imposed.  This proves the theorem.
\end{proof}

\paragraph{Concrete dimensions and scope.}
For $r=2$ this recovers C44-A21.  For every odd $r\geq3$, the six-centre
construction has residual multiplicity exactly two and full potential
Tate multiplicity eight.  For $r=2^j\geq4$, all binomial coefficients
in (A23.4) are odd: modulo two,
$(1+x)^{2^j-1}=1+x+\cdots+x^{2^j-1}$.
Thus the construction without centres has multiplicity four, while
the six-centre construction has residual multiplicity two and full
multiplicity eight.  All these entire algebraic axis forms are
positive, even though the unsplit residual block has more than one
potential Tate direction.

The arguments are valid for every $E$ in characteristic $p>3$.
Taking $E:y^2=x^3-x$ and $p\equiv3\pmod4$ gives the same supersingular
ambient examples as C44-A20, after finite extensions.  No claim is made
that any residual direction contains an algebraic cycle, or that these
known centres or their primitive classes are nonliftable.

The primary inputs and their scope are unchanged: the absolute abelian
Chow Lefschetz inverse of K\"unnemann, Theorem~5.2 p.~96; tame
fixed-point traces of Deligne--Lusztig, Theorem~3.2 p.~119; SGA~7~II,
Expose~XII, Theorem~3.3 and the subsequent affine-quadric comparison,
and Expose~XV, Theorem~3.4 pp.~195--196; Katz's Theorem~1.4.6 and
Lemmas~2.1.4--2.1.6; and the finite-character weight estimate of
Weil~II, Theorem~3.3.1 pp.~204--205 and Corollary~3.3.4 p.~206.
These theorems apply in the indicated arbitrary dimensions; no
additional equivariant monodromy theorem or mixed-characteristic
cycle-lifting converse is being imported.

For cote~44, Th.~5, scan~004 and the primitive instruction on scans
008--009, this extends the same sufficient family for QH9-axis to
higher-dimensional genuine pencils.  The axis sign is unconditional;
the separate ambient abelian term in a full height conclusion remains
subject to the unrefereed 032/Milne~2026 assumptions already recorded
in the programme.

% END work/scan_w/higher_dimensional_s3_axes.tex

% BEGIN work/scan_w/totally_real_axis_blocks.tex
% C44_POS /root/scan_w; GPT-6 Astra Pro; reasoning effort not exposed.
\section{A totally real coefficient field can separate several Tate directions}
\label{sec:c44-a24}

\begin{theorem}[C44-A24: one residual direction over a totally real field]
Let $B\subset X/\mathbf F_q$ be a smooth intersection of two members of
one ample line bundle, with $\dim X=2r+2$, and set
\[
 V_\ell=\ker\kappa_*\subset H^{2r}(B,\mathbf Q_\ell(r)),\qquad
 q_B(x,y)=(-1)^r\int_Bxy.
 \tag{A24.1}
\]
Assume the polarized embedded complete intersection has a smooth
mixed-characteristic lift with a complex generic fibre as in C44-A19.
Let $e$ be an actual relative rational algebraic correspondence whose
actions on the variable cohomologies are self-adjoint idempotents, and
put $U=eV$.  Suppose a totally real number field $F$ acts on $U$ through
actual relative rational algebraic correspondences, with identity $e$,
and every element of $F$ is self-adjoint for $q_B$.  The field and
projector identities are required on both fibres and in the coefficient
realizations; no assertion that they hold as equalities of Chow cycles
is necessary.

Assume the normalized Frobenius on $U$ has a common polynomial
$\Phi_F(T)\in F[T]$: on every component $U_v$ for a finite place
$v\mid\ell\ne p$, its characteristic polynomial is the image of
$\Phi_F$ in $F_v[T]$.  The $F_v$-ranks are the common $F$-rank of the
Betti space.  All special correspondences are defined after an allowed
finite extension of the finite ground field.

Let $K$ be a known positive rational algebraic subspace of the initial
realization of $U$, stable under $F$; $K=0$ is allowed.  Use the common
cycle representatives constructed below to realize it everywhere.
Suppose the residual polynomial on $K^\perp\cap U$ has at most one
root of unity, counted with algebraic multiplicity, over any, equivalently
every, embedding of $F$ into an algebraic closure.  Then $q_B$ is
positive definite on the complete rational algebraic subspace of $U$.
If $d=[F:\mathbf Q]$ and $\tau_F$ is that residual multiplicity, then
\[
 \dim_{\mathbf Q}K\ \leq\
 \dim_{\mathbf Q}\operatorname{Alg}(U)
       \ \leq\ \dim_{\mathbf Q}K+d\tau_F.
 \tag{A24.2}
\]
When the residual algebraic part is nonzero it is an $F$-line and its
$F$-valued form is totally positive.  The cycles on that line need not
lift.  No Tate conjecture, standard conjecture $D(B)$, Frobenius
semisimplicity, CM action, or 032 input is used.

The same assertion applies to several orthogonal relative blocks with
possibly different totally real fields.  If they sum to $V$, positivity
on their complete algebraic subspaces proves positivity on all of $A_B$.
\end{theorem}

\begin{proof}
We keep the field $F$ distinct from the normalized Frobenius operator
$\mathsf F$.  After a finite extension all chosen cycles and
correspondences are defined, so $\mathsf F$ commutes with $F$.

\paragraph{The actual $F$-valued form and its cycle values.}
Choose a rational basis $f_1,\ldots,f_d$ of $F$ and its trace-dual basis
$f_1^\vee,\ldots,f_d^\vee$.  On the Betti space define
\[
             b(x,y)=\sum_{j=1}^d q_B(f_jx,y)f_j^\vee.
 \tag{A24.3}
\]
Then $\operatorname{Tr}_{F/\mathbf Q}(a b(x,y))=q_B(ax,y)$ for
every $a\in F$.  Nondegeneracy of the trace pairing, commutativity of
$F$, and self-adjointness imply that $b$ is symmetric and $F$-bilinear,
and that $q_B=\operatorname{Tr}_{F/\mathbf Q}b$.
It is nondegenerate because $q_B$ is nondegenerate on $U$.
The same construction over $F\otimes\mathbf Q_\ell$ gives its local
forms $b_v$; their trace is the corresponding rational pairing.

For actual rational cycles in the block, (A24.3) lies in $F$ itself:
each coefficient is a rational intersection number of actual cycles.
These elements of $F$ have exactly the same images in every $F_v$.
This gives more information than the rational trace of a norm; no
realization independence of arbitrary vanishing classes is assumed.

To put all representatives in the block, first use the common absolute
special-fibre projector
\[
       P=1-\kappa^*(L_h^2)^{-1}\kappa_*
 \tag{A24.4}
\]
from C44-A15 and then apply $e$.  These are actual rational algebraic
correspondences.  The field operations in (A24.3), and below, mean
applications of chosen rational correspondences realizing the
specified elements of $F$, with $P$ and $e$ inserted if necessary.
Their cohomological identities hold everywhere by hypothesis.  Chow
idempotence is not required.

Choose an $F$-basis $C_1,\ldots,C_s$ of the known space $K$, represented
by these actual projected cycles, and put
\[
                         G=(b(C_i,C_j))\in M_s(F).
 \tag{A24.5}
\]
It is nonsingular.  Positivity of the rational form on $K$ makes every
real conjugate of $G$ positive definite, so $\det G$ is totally positive.
The same matrix in every $F_v$ makes the actual cycles $F_v$-independent
and gives common nondegenerate subspaces $K_v$ of rank $s$.
After defining the finitely many $f_jC_i$ over a finite field, their
normalized Frobenius is the identity.  The common residual polynomial
is therefore $\Phi_F(T)/(T-1)^s$.

For any additional algebraic cycle in $U$, subtract its $b$-orthogonal
projection onto $K$, using $G^{-1}$ and the elements of $F$ in (A24.3).
This operation is a finite rational combination of actual cycle
correspondences.  It gives one actual rational cycle whose class lies
in $K_v^\perp\cap U_v$ in every realization, and represents the intended
residual class initially.  This is the field-valued version of the
common-representative repair in C44-A19.

\paragraph{A totally positive Betti determinant at every embedding.}
For a real embedding $\sigma:F\hookrightarrow\mathbf R$, let
$U_\sigma=U_{\mathrm B}\otimes_{F,\sigma}\mathbf R$.
The relative algebraic correspondences preserve Hodge types.  The
space $U_\sigma$ is therefore a real Hodge substructure of primitive
middle cohomology.  Self-adjointness makes the distinct real-embedding
summands orthogonal, and the restriction of $q_B$ is $\sigma b$.
The Hodge--Riemann relations are positive on type $(r,r)$; all negative
real directions come in conjugate off-diagonal pairs.  Thus the
negative index of $\sigma b$ is even for \emph{each} $\sigma$, not
only for their rational trace form.
In any $F$-basis of the Betti space this proves that
\[
                            D=\det_F b_{\mathrm B}\in F^*
              \quad\hbox{is totally positive}.
 \tag{A24.6}
\]
This is an ordinary $F$-determinant; it is not the determinant of the
rational trace form, and no discriminant-of-$F$ factor is omitted.
For a zero-dimensional block the conclusion is empty.

Smooth proper base change and Betti--etale comparison respect these
same relative correspondences.  They are consequently $F_v$-linear
and preserve $b_v$, which is characterized by its trace identities.
It follows that $D$ gives the determinant squareclass of $b_v$ for
every finite $v\nmid p$.  The known special cycle space $K$ does not
need to extend to the Betti fibre for this assertion.

\paragraph{One potential Tate line over $F$, with common cycle values.}
The root-of-unity condition is unchanged under any embedding of the
coefficient field: an extension of that embedding to the algebraic
closure permutes roots and preserves roots of unity and their
multiplicities.  If the multiplicity is one, its root is $1$ or $-1$.
Indeed the polynomial at a real embedding has real coefficients,
so any nonreal root of unity would have a second conjugate.  The
corresponding equality at $1$ or $-1$ is then an equality in $F$,
and its simple multiplicity is common to every place.

If $\tau_F=0$, every algebraic residual class is zero: its normalized
Frobenius orbit is finite.  Suppose instead $\tau_F=1$ and set
$S=\mathsf F^2$ on each residual local space.  Its eigenline $T_v$ at
one is simple and nondegenerate.  For completeness, an isometry pairs
only reciprocal generalized eigenspaces.  If the simple line at one
were degenerate, it would be orthogonal to the whole space.  Every
finite-orbit vector lies on this line, without any semisimplicity
assumption on the remaining generalized eigenspaces.

Let $z$ be an actual residual cycle nonzero at the initial coefficient
prime.  It is nonzero in at least one $F_v$ component, and the
nondegenerate line there gives
\[
                              a=b(z,z)\in F^*.
 \tag{A24.7}
\]
Because this is one nonzero element of $F$, its image is nonzero at
\emph{every} finite place.  Therefore the same actual cycle is nonzero
and spans $T_v$ at every $v\nmid p$.  This is the needed repair of the
cross-place nonvanishing issue; it is stronger than just knowing that
the rational self-intersection has a given sign or trace.

\paragraph{The non-Tate determinant and the missing finite places.}
Let $N_v=T_v^\perp$ inside the residual space.  On $N_v$ the isometry
$S$ has no eigenvalue $1$ or $-1$.  Its common $F$-polynomial is denoted
$\psi(T)$.  It is obtained from the common residual polynomial by
squaring eigenvalues and removing the simple factor $T-1$.
The dimension of $N_v$ is even, and $\det S=1$ there, because all
eigenvalues occur in reciprocal pairs.

The operator $A=S-S^{-1}$ is invertible and skew-adjoint for $b_v$.
The form $b_v(Ax,y)$ is therefore nondegenerate alternating.  Its
determinant is a square.  Consequently
\[
 \det b_v|N_v\equiv \det A
       =\delta:=\psi(1)\psi(-1)
                              \pmod{(F_v^*)^2}.
 \tag{A24.8}
\]
The equality for $\det A$ uses the even dimension and $\det S=1$;
it includes the empty complement, with $\delta=1$.
The common rational Frobenius polynomial on the underlying block is
the field norm of $\Phi_F$.  Purity therefore puts every root of every
real conjugate of $\psi$ on the unit circle.  None is $1$ or $-1$,
so conjugate reciprocal pairs show that $\delta$ is totally positive.

The orthogonal decomposition into $K_v$, $T_v$ and $N_v$ gives
\[
              \frac{a\det(G)\delta}{D}\in(F_v^*)^2
                             \qquad(v\nmid p).
 \tag{A24.9}
\]
The map
\[
 F^*/(F^*)^2\longrightarrow
             \prod_{v\notin\mathcal S}F_v^*/(F_v^*)^2
 \tag{A24.10}
\]
is injective for any finite set $\mathcal S$ of finite places.  Here
is the exact proof needed.  If $c$ is a nonsquare in $F$, the quadratic
extension $F(\sqrt c)/F$ is nontrivial.  Number-field Chebotarev gives
infinitely many unramified primes with nontrivial Frobenius, and one
can choose such a prime outside $\mathcal S$.  At that prime the
extension is a quadratic field extension, so $c$ is not a local square.
Apply this with $\mathcal S=\{v:v\mid p\}$ to (A24.9).

The ratio in (A24.9) is thus a square in $F^*$.  All its other factors
are totally positive by (A24.5), (A24.6) and (A24.8).  Hence $a$ is
totally positive.  The form on the rational space $Fz$ is
\[
                          q_B(fz,fz)
                            =\operatorname{Tr}_{F/\mathbf Q}(a f^2),
 \tag{A24.11}
\]
which is positive definite.

Finally, if $y$ is any other actual residual cycle, then
$b(y,z)/a\in F$.  At every place both $y$ and $z$ lie on the same
line and $z$ is nonzero, so
$y=(b(y,z)/a)z$ cohomologically in the initial realization and in all
the other realizations.  Thus the complete residual algebraic space
has $F$-dimension at most one.  Adding the known positive space $K$
proves positivity and (A24.2).  Orthogonality proves the stated
multiple-block version.
\end{proof}

\begin{corollary}[C44-A24a: the polynomial hypothesis from a finite symmetry algebra]
Suppose the relative action and embedded equivariant lift come from
a finite constant group $G$ as in C44-A19b.  Let $e\in\mathbf Q[G]$
be self-adjoint, and suppose the corner $e\mathbf Q[G]e$ contains a
totally real field $F$, with unit $e$, fixed pointwise by the adjoint
involution $g\mapsto g^{-1}$.  Then the common $F$-valued Frobenius
polynomial required in C44-A24 exists.  In particular that theorem
applies once its known-positive-space and residual multiplicity tests
are satisfied.
\end{corollary}

\begin{proof}
The relative Betti $F$-module is finite-dimensional and free, and
comparison gives its common rank in every $F_v$ realization.
For $f_j$ in a rational basis of $F$, the element $f_j=e f_j e$ is a
rational linear combination of group elements.  The finite-field
twists in C44-A19b give rational, coefficient-independent numbers
\[
             t_{j,n}=\operatorname{Tr}_{\mathbf Q_\ell}
                       (f_j\mathsf F^n\mid U_\ell)
                         \qquad(n\geq1).
 \tag{A24.12}
\]
In detail the trace on $U$ equals the trace of the same group-algebra
element times Frobenius on $V$; each group-element trace is obtained
degreewise from its smooth projective twist, and the ambient traces
are subtracted.  Centrality of $e$ or of $f_j$ in the whole group
algebra is unnecessary.

Trace transitivity for a free $F\otimes\mathbf Q_\ell$-module gives
\[
      t_n:=\sum_jt_{j,n}f_j^\vee\in F,
       \qquad t_n\mapsto\operatorname{Tr}_{F_v}
                              (\mathsf F^n\mid U_v)
                                  \quad\hbox{for every }v\mid\ell.
 \tag{A24.13}
\]
Indeed pairing both sides with every $f_j$ under the field trace
recovers (A24.12), and the trace pairing is nondegenerate.
Newton identities in characteristic zero now reconstruct one
polynomial $\Phi_F\in F[T]$ with the required images.  This is an
actual reconstruction of the $F$-coefficients, rather than an
inference from rationality of the norm polynomial alone.
\end{proof}

\paragraph{The dihedral algebra giving a nonsplit rational example of the criterion.}
For $G=D_5=\langle\rho,\sigma:\rho^5=\sigma^2=1,
\sigma\rho\sigma=\rho^{-1}\rangle$, put
\[
 e_1=\frac1{10}\sum_g g,\qquad
 e_\varepsilon=\frac1{10}\sum_g\varepsilon(g)g,\qquad
 e_0=1-e_1-e_\varepsilon,\qquad
 e_\pm=\frac12e_0(1\pm\sigma).
 \tag{A24.14}
\]
These are the four self-adjoint orthogonal rational projectors.
The nontrivial rational simple component is
\[
              e_0\mathbf Q[D_5]\simeq M_2(F),
                         \qquad F=\mathbf Q(\sqrt5).
 \tag{A24.15}
\]
One direct verification uses $\alpha^2+\alpha-1=0$ and the matrices
\[
       \rho\longmapsto
          \begin{pmatrix}0&-1\\1&\alpha\end{pmatrix},\qquad
       \sigma\longmapsto
          \begin{pmatrix}0&1\\1&0\end{pmatrix}.
 \tag{A24.16}
\]
They satisfy the dihedral relations.  The two real embeddings of $F$
give the two distinct absolutely irreducible two-dimensional
representations; together with $1,\varepsilon$ their squared degrees
sum to ten.  These explicit split representations verify (A24.15).
The central element $\rho+\rho^{-1}$ acts as $\alpha$ and is
self-adjoint.  Thus
$e_\pm\mathbf Q[D_5]e_\pm\simeq F$ is fixed pointwise by the adjoint
involution.

If a geometric axis supplies this actual relative symmetry algebra,
each of the two $e_\pm$ blocks can have residual $F$-rank one, hence
residual rational potential Tate dimension two.  C44-A24 still proves
its complete algebraic sign.  The available rational group-algebra
projectors cannot split either corner further, since a field has no
nontrivial idempotents.  Consequently this case is not obtained by
applying the rational multiplicity-one test C44-A19 to the supplied
group-algebra projectors.  No nonexistence of other geometric
projectors is asserted, and no smooth-pencil or closed-point example
is inferred from this algebra calculation alone.

\paragraph{Sources and scope.}
The geometric comparison and common-cycle inverse are the checked
inputs of C44-A15 and A19.  The new number-field squareclass argument
uses only the quadratic case of Chebotarev: Milne, \emph{Class Field
Theory}, Chapter~V, Theorem~3.23 and Corollary~3.24, printed p.~164,
or \emph{Algebraic Number Theory}, Theorem~8.31, printed p.~147.
The finite-place deletion is proved explicitly in (A24.10); the
rational-only statement in Ancona--Marmora is not being cited as a
number-field theorem.  The trace-dual reconstruction, every-embedding
determinant sign and dihedral algebra calculation are proved above.

For cote~44, Th.~5, scan~004 and the primitive warning on scan~008,
this supplies an additional sufficient test for the actual QH9-axis
form.  It requires an extending self-adjoint field of correspondences
and a verified common field-valued polynomial.  It does not promote
a special Frobenius eigenspace to a relative factor.  The axis sign
is unconditional under these explicit hypotheses; the separate
ambient abelian contribution to a full height theorem retains the
032/Milne~2026 dependency already stated in the programme.

% END work/scan_w/totally_real_axis_blocks.tex

% BEGIN work/axis/affine_d5_axes.tex
% C44_POS /root; GPT-6 Astra Pro; effort not exposed.
% Complete candidate proof; promote only after whole-proof checks.
\section{A dihedral family with five residual potential Tate directions}
\label{sec:c44-a25}

\begin{theorem}[C44-A25: an actual totally real coefficient-field application]
Let $k$ be a finite field of odd characteristic $p\ne5$, and let
$E/k$ be an elliptic curve. After finite extensions there are genuine
Lefschetz pencils on $X\simeq E^6$, with smooth axes $B$ of dimension
four, for which the following hold on a Zariski-dense set of closed
parameters. Write $V=\ker\kappa_*$ with twist $(2)$ and its ordinary
intersection pairing.

One can arrange either no prescribed cycle space $K$, or a positive
constant rational algebraic space $K$ of dimension ten. In both cases
the root-of-unity multiplicity of normalized Frobenius on
$W=V\cap K^\perp$ is exactly five. The entire rational algebraic axis
space $A_B$ is nevertheless positive definite. With centres,
\[
                 \dim_{\mathbf Q}A_B\geq10,\qquad \tau(V)=15.
                 \tag{A25.1}
\]
The construction carries an actual extending action of the dihedral
group of order ten. Its two nontrivial minimal rational matrix blocks
have coefficient field $F=\mathbf Q(\sqrt5)$ and residual potential
Tate dimension one over $F$, hence two over $\mathbf Q$.
C44-A24 proves their signs. The result is unconditional; it does not
assert algebraicity or liftability of any of the five residual
potential directions.
\end{theorem}

\begin{proof}
We give the geometry and all changes to the checked affine-$S_3$
construction C44-A20--A21. No monodromy or smooth-pencil assertion is
inferred merely from the group algebra.

\paragraph{An integral affine dihedral action.}
Index five coordinates by $\mathbf Z/5\mathbf Z$ and set
\[
 A=\ker\left(\Sigma:E^5\longrightarrow E\right),\qquad
 X=E^2\times A.
 \tag{A25.2}
\]
The map $(z_1,z_2,z_3,z_4)\mapsto
(-\sum_{i=1}^4z_i,z_1,z_2,z_3,z_4)$ identifies $A$ with $E^4$.
Let $R$ cyclically permute the five coordinates and let $S$ send the
coordinate with index $j$ to the coordinate with index $-j$.
Choose $0\ne T\in E[5](k)$ after a finite extension and define
\[
 \rho(a,b;z)=(a+T,b;Rz),\qquad
 \sigma(a,b;z)=(-a,-b;Sz).
 \tag{A25.3}
\]
These satisfy $\rho^5=\sigma^2=1$ and
$\sigma\rho\sigma=\rho^{-1}$. Every nonidentity rotation is free,
because its translation on the first factor is nonzero.
The fixed locus of $S$ in $A$ is exactly
\[
                  (-2u-2v,u,v,v,u),\qquad (u,v)\in E^2.
 \tag{A25.4}
\]
It is connected even in characteristic three. Thus $X^\sigma$ is a
union of sixteen abelian surfaces, with tangent eigenspaces of
dimensions two and four. The five reflection fixed loci are pairwise
disjoint: a common fixed point would be fixed by a nonidentity
rotation. These observations use only $p\ne2,5$.

Use the displayed coordinates $X\simeq E^6$ to define the particular
product bundle $M=\boxtimes_{j=1}^6\mathcal O_E(3[0])$, and put
\[
 H=\left(\bigotimes_{g\in D_5}g^*M\right)^{\otimes2},\qquad L=H^t.
 \tag{A25.5}
\]
It is very ample with its natural group linearization. Its fibre
character on every reflection fixed surface is trivial, because the
norm has been squared. We use this specified extending polarization,
not an arbitrary divisor on a supersingular product.

After another finite extension choose a root
$\alpha^2+\alpha-1=0$ in the residue field. The two-dimensional
representation $U$ with matrices
\[
 R_U=\begin{pmatrix}0&-1\\1&\alpha\end{pmatrix},\qquad
 S_U=\begin{pmatrix}0&1\\1&0\end{pmatrix}
 \tag{A25.6}
\]
is absolutely irreducible and faithful on its projective line.
Its rotations have eigenvalue ratio of order five, and its
reflections have two distinct eigenvalues. The representation
extends over an unramified mixed-characteristic ring containing the
lifted root $\alpha$, since the discriminant five is a unit.

For the ten-centre version choose $c\in A$ with free $D_5$ orbit and
put
\[
                     Z_g=E^2\times\{gc\}\qquad(g\in D_5).
 \tag{A25.7}
\]
The free locus is nonempty; each nonidentity fixed locus in $A$ is
proper. These ten surfaces are disjoint and avoid all fixed loci in
$X$. Their union is denoted $Z$; in the other version let $Z$ be empty.
The section parameter space is
\[
                 \mathcal E_t=\operatorname{Hom}_{D_5}
                                  (U,H^0(X,I_ZL)).
 \tag{A25.8}
\]
Its common zero locus is the axis and its original two-dimensional
section space defines the pencil.

\paragraph{Smooth axes and genuine original pencils.}
The compatible-jet proof of C44-A20 applies with ten-point free orbits
and five-point reflection orbits. Here is why no new restriction on
jets is introduced. The order ten is invertible, so averaging projects
the usual evaluation maps onto their equivariant parts. For large
$t$, products of separating hyperplanes give simultaneous prescribed
jets at any bounded number of distinct points; the bound is uniform
as the points vary. Vanishing along the fixed union $Z$ is handled by
first globally generating its ideal and by relative Serre vanishing
for its conormal jets. On a free orbit the two section components
therefore have arbitrary prescribed jets at one representative. At a
reflection fixed point they have exactly the even and odd jets for
the tangent decomposition $(2,4)$, as the fibre of $L$ is trivial.

At a free point, failure of axis smoothness has two zero values and
a rank-at-most-one $2$-by-$6$ derivative matrix. Its codimension is
seven, greater than six. At a fixed surface, the even component must
vanish; its zero derivative has codimension three in all, and zero
odd derivative has codimension five in all. Both exceed the dimension
two of the fixed locus. Along a centre, rank failure of the two rows
in its four-dimensional normal space has codimension three, again
greater than two. Thus a nonempty parameter open has smooth axes,
and every pencil member is smooth along the centres and along the
smooth base of the pencil.

There are exactly twelve exceptional parameters on the projective
pencil line: the two common eigenlines of the rotations and the ten
eigenlines of the five reflections. A nontrivial stabilizer is cyclic
of order five or two. Distinct reflection eigenline sets cannot meet,
since the product of two reflections is a rotation and a reflection
interchanges the two rotation eigenlines. All twelve members can be
required to be smooth. At a free point a fixed member has critical
jet codimension seven. At a reflection fixed surface an eigenmember
has critical codimension three if even and four if odd. Both exceed
two. A member not an eigenmember of that reflection has all six
first derivatives available and critical codimension seven there.
This also treats fixed loci of other reflections at an exceptional
parameter.

For the other parameters, a degenerate critical point imposes eight
conditions on the seven-dimensional point--parameter incidence.
Two critical points in distinct free orbits impose fourteen conditions
on dimension thirteen. For two distinct points in the same free
orbit, the parameter and its translate are distinct; the two section
components then have zero first jet at one representative, imposing
fourteen conditions on dimension seven. A fixed-point critical pair
has codimension seven on a three-dimensional incidence. The compatible
jet maps above justify these counts simultaneously with the imposed
centres. Thus a general parameter has only nondegenerate critical
points, no two in the same member. The smooth base excludes critical
points there. The nondegenerate Hessian and nonzero parameter derivative
make the critical scheme reduced; properness makes it finite.

It is nonempty. On $X\times\mathbf P^1$ the first-jet bundle tensored
by $\mathcal O(1)$ has top Chern number
\[
                   \int_{X\times\mathbf P^1}(h+\eta)^7
                             =7\int_Xh^6>0,
 \tag{A25.9}
\]
where $h=c_1(L)$. The section given by the pencil has its zero scheme
exactly at these critical pairs. This proves the claimed genuine
Lefschetz-pencil open, denoted $\Omega$.

The exact action, polarization and section pair lift. Lift $E$ and
its prime-to-$p$ torsion point, the coordinate maps defining $A$, the
specified product bundle and the norm-square. Lift $c$ by smoothness;
the centres remain disjoint. Relative Serre vanishing lifts the given
section pair, and averaging by $1/10$ makes it equivariant for the
lift of (A25.6). The nonsmooth locus is closed in the proper family;
its empty special fibre makes the generic axis smooth. These are
actual extending correspondences and embedded lifts of the kind
required in C44-A24.

\paragraph{Positive regular centres and exact character ranks.}
Write $h_0=c_1(H)$, $d=\int_{Z_1}h_0^2$, and
$A_0=(L_{h_0}^2)^{-1}[Z_1]_X$, using the absolute abelian Chow
inverse from C44-A15. Put $c_0=\int_X A_0^2h_0^2$.
All centre classes have the same ambient class. Their projected
classes $w_g=[Z_g]_B-t^{-2}\kappa^*A_0$ have Gram matrix
\[
                         3t^2d I_{10}-t^{-2}c_0 J_{10}.
 \tag{A25.10}
\]
Indeed the normal bundle of a centre in $X$ is trivial, whereas its
normal bundle in $B$ has total Chern class $(1+th_0)^{-2}$; disjoint
centres have zero mutual intersections before projection. The matrix
is positive definite whenever $3t^4d>10|c_0|$. Thus their span $K$ is
a constant regular rational representation of dimension ten in every
realization. The same construction supplies its actual cycle basis.

Let $d_6=\int_Xh^6$ and $d_2=\int_{F_0}h^2$, where $F_0$ is any
reflection fixed surface. In the four
conjugacy classes $(1,\sigma,\rho^{\pm1},\rho^{\pm2})$ the equivariant
Euler characteristic of $B$ is
\[
                         (5d_6,-16d_2,0,0).
 \tag{A25.11}
\]
The first entry follows from $c(T_B)=(1+h)^{-2}$. The fixed locus of
a reflection on $B$ is a smooth divisor in each of its sixteen fixed
abelian surfaces, hence has Euler characteristic $-16d_2$. Rotations
have no fixed points. The tame fixed-point trace theorem used in
C44-A20 gives precisely these alternating cohomology traces.

On $H^1(X)$ the exterior-character generating polynomials are
\[
 \begin{gathered}
 (1+z)^{12},\qquad (1+z)^4(1-z)^8,\\
 (1+z)^2(1+z^5)^2,\qquad (1+z)^2(1+z^5)^2.
 \end{gathered}
 \tag{A25.12}
\]
For a rotation the augmentation representation contributes
$(1+z^5)/(1+z)$, twice in cohomology, while the sign factors contribute
$(1+z)^4$. It follows that $H^4(X)$ has character
$(495,-17,0,0)$ and absolute multiplicities $(41,58,99,99)$.
The alternating sum of degrees zero through three has character
$(-165,-5,0,0)$ and multiplicities $(-19,-14,-33,-33)$.
Weak Lefschetz, duality and the ambient middle image therefore give
\[
 \begin{aligned}
 m_1&=d_6/2-8d_2-3,\\
 m_\varepsilon&=d_6/2+8d_2-30,\\
 m_{U_1}=m_{U_2}&=d_6-33.
 \end{aligned}
 \tag{A25.13}
\]
Here $U_1,U_2$ are the two real conjugate two-dimensional irreducibles.
For example their characters at the rotation classes are respectively
$(\alpha,-1-\alpha)$ and $(-1-\alpha,\alpha)$; their reflection
characters are zero. The displayed multiplicities follow directly
by character inner products, using class sizes $1,5,2,2$.
Since $d_6=720\chi(L)$, their parities are $(1,0,1,1)$, and all ranks
exceed eight for sufficiently large $t$. Removing the regular $K$
subtracts $(1,1,2,2)$, so the residual parities are $(0,1,1,1)$.
The standard rational isotype is a module over $M_2(\mathbf Q(\sqrt5))$,
with the two conjugate multiplicity spaces coming from one module over
that coefficient field.

\paragraph{Full orthogonal monodromy from real character moments.}
Remove all fixed surfaces and all centres from $X$ and take the free
quotient $Y$. It is a smooth geometrically connected sixfold with a
connected $D_5$-torsor. For a rational point $y$ over $\mathbf F_Q$,
write $g_y$ for the conjugacy class of its torsor Frobenius. A section
in (A25.8) descends to a section of a rank-two bundle on $Y$.
Simultaneous evaluation at four quotient points uses at most forty
geometric points upstairs. The uniform jet argument therefore gives
four-wise independent zero indicators $I_y$ with probability $v=Q^{-2}$,
including nonsplit torsor fibres.

For each of the four real irreducible characters $\chi$, put
\[
              T_\chi=\sum_{y\in Y(\mathbf F_Q)}
                                  \chi(g_y)(I_y-Q^{-2}).
 \tag{A25.14}
\]
Choose either real embedding when $\chi$ is field-valued. The elementary
orbit identity underlying the character trace is
\[
 \frac1{10}\sum_{g\in D_5}\chi(g^{-1})\#\operatorname{Fix}(gF;T)
                       =\sum_{y\in(T/D_5)(\mathbf F_Q)}\chi(g_y)
 \tag{A25.15}
\]
for a free finite set or variety $T$, interpreted pointwise. It follows
by counting the elements conjugating $g_y$ to $g^{-1}$ in its ten-point
orbit. This is a multiplicity trace: there is no extra factor $\dim\chi$.

Apply (A25.15) to $B$ minus its fixed and prescribed loci, and subtract
$Q^{-2}$ times the corresponding ambient trace. As in C44-A21, Gysin
cancels degrees above four against ambient degrees shifted by four;
$L_h^2$ cancels the ambient middle image. All remaining terms have
weight at most three. Fixed surfaces cost only $O(Q)$, since their
intersection with $B$ is a curve. In the centre case the torsor on
$Z\to Z/D_5$ is split, with quotient one $E^2$; its subtraction is
$\dim(\chi)(1-Q^{-2})\#E^2(\mathbf F_Q)
=\dim(\chi)Q^2+O(Q^{3/2})$, precisely the constant regular $K$ term.
On the smooth parameter open this gives the uniform estimate
\[
                  T_\chi=\operatorname{Tr}(F\mid M_\chi^{\rm un})
                                     +O(Q^{3/2})
 \tag{A25.16}
\]
where $M_\chi$ is the residual multiplicity sheaf. On the complement
we retain the real point-count expression $T_\chi$ from (A25.14),
without asserting any lisse-sheaf trace interpretation there.

The finite-monodromy character sheaf and orthogonality give
\[
       S_2:=\sum_y\chi(g_y)^2=Q^6+O(Q^{11/2}),\qquad
       S_4:=\sum_y\chi(g_y)^4=O(Q^6).
 \tag{A25.17}
\]
The leading coefficient is one, because the invariant space of the
square of an irreducible real character has dimension one and
arithmetic Frobenius fixes that canonical invariant. This applies to
each real embedding of the standard character, not just to their sum.
Four-wise independence gives the exact fourth moment
\[
 \operatorname{Avg}_{\mathcal E_t(\mathbf F_Q)}T_\chi^4
       =3v^2(1-v)^2(S_2^2-S_4)
            +\bigl(v(1-v)^4+(1-v)v^4\bigr)S_4.
 \tag{A25.18}
\]
Dividing by $Q^8$ gives limit three. The weights and all point-count
expressions are real in either real embedding, so their fourth powers
are nonnegative even off $\Omega$. The same extension argument and
Katz moment lemmas used in C44-A21 bound the smooth-open monodromy
fourth moment by three. The three independent orthogonal contraction
tensors give the reverse inequality. Thus the fourth moment is three
in each multiplicity sheaf. Self-duality gives a symmetric
nondegenerate form on each multiplicity space, since all four
irreducibles are of real orthogonal type. Purity supplies the
reductivity needed in Katz's criterion.

For completeness the required reflections also occur in this actual
family. At a free orbit prescribe a rank-one derivative and a
nondegenerate remaining quadratic form in five variables. This gives
exactly ten nodes. At a reflection fixed point set the even value and
its derivative to zero, keep the odd derivative nonzero, and impose
a nondegenerate invariant quadratic form after eliminating that odd
variable. Its tangent signs are $(2,3)$, so this gives exactly five
nodes. Independent compatible jets and the strict inequalities above
exclude all other singularities. Varying the constant even or remaining
component gives a regular transverse smoothing, with general member
in $\Omega$.

The same equivariant etale Morse calculation as in C44-A20 applies
because only two must be inverted. The determinant of the reflection
on the five remaining variables is minus one; on the compactified
even quadric it interchanges the two generator families and acts by
minus one on the vanishing line. Consequently the ten-node orbit is
the regular representation, and the five-node orbit is
\[
                  \operatorname{Ind}_{C_2}^{D_5}(\mathrm{sgn})
                                   =\varepsilon\oplus U_1\oplus U_2.
 \tag{A25.19}
\]
The node cycles are orthogonal and each has square two. They are
independent and lie in $W$, because ambient and centre classes are
fixed by inertia. Tame inertia gives reflections on the trivial and
sign multiplicities from the free orbit, and on each standard
multiplicity from (A25.19). Katz's rank-greater-than-eight, fourth-moment
and reflection theorem therefore gives full orthogonal geometric
monodromy in every multiplicity factor.

\paragraph{Closed-point selection with the correct coefficient degree.}
The common field-valued polynomials and actual field action are those
of C44-A24a and its explicit dihedral calculation. Choose an auxiliary
coefficient prime different from $p$ which splits
$\mathbf Q(\sqrt5)$. Let $D$ be the total rational rank of $W$ and use
\[
                          N=\operatorname{lcm}(1,2,\ldots,2D^2).
 \tag{A25.20}
\]
This bounds the orders of every possible root of unity in the rational
Frobenius polynomial. In particular it uses rational degree, not the
possibly smaller degree of an $F$-polynomial.

Let $\mathcal H$ be the joint geometric monodromy group in the four
absolute orthogonal multiplicity factors. Every projection is onto;
therefore every projection of its connected component is onto the
corresponding special orthogonal group. In each even-rank factor the
no-root selector of C44-A12 is a nonempty open, and in each odd-rank
factor the simple-root selector of C44-A16 is a nonempty open. Their
inverse images intersect in the connected irreducible group
$\mathcal H^0$. The same finite-quotient Chebotarev argument, after a
constant-field extension, selects a Zariski-dense set of closed
parameters with all multiplicities minimal simultaneously. No product
independence of the factors is assumed.

The parities in (A25.13) give total multiplicity
$1+2+2=5$ with no centres and $1+2+2=5$ after removing the regular
centre space, now with the one-dimensional contribution in the sign
factor. All ten centre directions are Frobenius-fixed, giving total
multiplicity fifteen in the latter case.

Finally apply the four actual self-adjoint projectors of (A24.14).
The trivial and sign blocks have at most one residual rational Tate
direction. Each of the two standard matrix blocks has at most one
residual direction over $\mathbf Q(\sqrt5)$, hence rational residual
dimension two. The known space $K$ is stable under the group algebra
and decomposes orthogonally with positive restrictions. The embedded
lifts, common field polynomials and self-adjoint field actions have
all been supplied. C44-A24 proves positivity on the whole algebraic
subspace of every block, hence on all of $A_B$. The Gram matrix gives
the lower bound ten. This proves the theorem.
\end{proof}

\paragraph{Exact checks, primary inputs and scope.}
The accompanying script \texttt{verify\_d5\_data.py} checks the four
self-adjoint orthogonal rational projectors, the polynomial
$(\rho+\rho^{-1})^2+(\rho+\rho^{-1})-1$ on the standard block,
the exterior-character multiplicities, and enumerated field-valued
Bernoulli fourth moments by exact rational arithmetic. These are
arithmetic certificates, not substitutes for the geometric proof.

The primary geometric inputs have the same precise locators as in
C44-A20--A21: Deligne--Lusztig, Theorem~3.2 p.~119 for tame fixed-point
Euler traces; SGA~7~II, Expose~XII, Theorem~3.3 pp.~75--77 and the
affine-quadric comparison, and Expose~XV, Theorem~3.4 pp.~195--196 for
the ordinary double point pairing and local monodromy; Katz,
Theorem~1.4.6 and Lemmas~2.1.4--2.1.6 for orthogonal monodromy from
moments and reflection; Weil~II, Theorem~3.3.1 pp.~204--205 and
Corollary~3.3.4 p.~206 for the finite-character estimates. The quotient,
actual lifting, common-polynomial and number-field sign arguments are
the proved C44-A19, A24 and A24a.

This family is available also in characteristic three, and in the
supersingular case whenever such an $E$ is chosen. It proves the
QH9-axis sign on these actual closed families, beyond the supplied
rational-block multiplicity-one test. It does not prove the universal
fixed-axis sign or exclude other geometric projectors. For cote~44,
Th.~5, scan~004 and the primitive warning on scans~008--009, the axis
term is unconditional. The separate ambient abelian term in the full
height formula remains conditional on the unrefereed 032/Milne~2026
inputs; those preprints are not used in this axis theorem.

% END work/axis/affine_d5_axes.tex

% BEGIN work/scan_w/unbounded_dihedral_axis_blocks.tex
% C44_POS /root/scan_w; GPT-6 Astra Pro; reasoning effort not exposed.
% One dimensional/field-degree generalization of A24--A25; no separate clock event.
\section{Unbounded residual multiplicities across dimensions}
\label{sec:c44-a26}

This is one dimensional and coefficient-field generalization of the
checked C44-A24--A25 mechanism, not a separate progress-clock event.
The following choice of dimensions allows both an unbounded totally
real field degree and an explicit nonzero known positive cycle space.

\begin{theorem}[C44-A26: dihedral axes with unbounded residual potential rank]
Fix an odd characteristic $p$.  For every odd prime $\nu\geq5$ with
$\nu\ne p$, put
\[
 d=(\nu-1)/2,\quad r=\nu-2,\quad a=\nu-1,\quad n=2\nu-2,
 \qquad \epsilon_\nu=(-1)^d.
 \tag{A26.1}
\]
Over a finite extension of a finite field of characteristic $p$, there
are smooth axes $B$ of dimension $2r$ in $X\simeq E^n$, cut out by
two members of one ample line bundle, whose original pencils are
genuine Lefschetz pencils.  For each sufficiently large degree the
parameter family contains a Zariski-dense set of closed members with
positive definite entire rational algebraic axis space $A_B$.

One may arrange either of the following two versions.  Without a
prescribed centre space the total root-of-unity multiplicity of the
normalized variable Frobenius is exactly $\nu-\epsilon_\nu$.
With prescribed centres there is a constant positive algebraic
subspace $K\subset A_B$ of dimension $2\nu$, and
\[
 \tau(V\cap K^\perp)=\nu+\epsilon_\nu,\qquad
 \tau(V)=3\nu+\epsilon_\nu,\qquad
                  \dim_{\mathbf Q}A_B\geq2\nu.
 \tag{A26.2}
\]
The sign is for $q_B(x,y)=(-1)^r\int_Bxy$.
These are unconditional assertions.  The residual directions are
potential Tate directions; their algebraicity is not asserted.
The residual multiplicities are unbounded as the dimension grows,
even with the characteristic $p$ fixed.  No unbounded-rank assertion
in a fixed ambient dimension is being made.
\end{theorem}

\begin{proof}
We prove all changes from C44-A25.

\paragraph{Integral geometry and the actual field of correspondences.}
Let $E$ be any elliptic curve over the finite field and set
\[
 A=\ker\left(\Sigma:E^\nu\longrightarrow E\right)\simeq E^{\nu-1},
 \qquad X=E_{\rm sign}^{a}\times A.
 \tag{A26.3}
\]
Index the augmentation coordinates by $\mathbf Z/\nu\mathbf Z$.
Let $R$ rotate these coordinates and $S$ send index $j$ to $-j$.
After a finite extension choose an order-$\nu$ point $T$ on $E$.
The generator $\rho$ translates the first sign coordinate by $T$
and acts by $R$ on $A$; the generator $\sigma$ negates every sign
coordinate and acts by $S$ on $A$.
They satisfy the dihedral relations, giving a group of order $2\nu$.
All nonidentity rotations act freely.  The reflection fixed locus
in $A$ is connected and isomorphic to $E^d$: its paired coordinates
are $u_1,\ldots,u_d$, and its remaining coordinate is
$-2\sum_j u_j$.  Hence each reflection fixes $2^{2a}$ disjoint
abelian $d$-folds, and distinct reflection fixed loci are disjoint.
The tangent signs there have dimensions
\[
                            (d,a+d)=(d,3d).
 \tag{A26.4}
\]

Using the explicit product coordinates on $X$, choose
$M=\boxtimes_{i=1}^n\mathcal O_E(3[0])$ and
$H=(\bigotimes_{g\in D_\nu}g^*M)^{\otimes2}$, with its natural
linearization.  It is very ample and its reflection stabilizer
characters are trivial.  Set $L=H^t$.

After a finite extension let the residue field contain a primitive
$\nu$th root $\zeta$, and set $\alpha=\zeta+\zeta^{-1}$.
Use the rank-two generator representation
\[
 \rho\longmapsto\begin{pmatrix}0&-1\\1&\alpha\end{pmatrix},\qquad
 \sigma\longmapsto\begin{pmatrix}0&1\\1&0\end{pmatrix}.
 \tag{A26.5}
\]
Its rotation eigenvalue ratio has order $\nu$, and reflections
interchange the distinct rotation eigenlines.  It is absolutely
irreducible and projectively faithful in every $p\ne2,\nu$.
The root of unity lifts over an unramified mixed-characteristic
extension because $\nu$ is invertible.  Thus the actual representation,
elliptic variety, affine action and specified norm-square polarization
all lift.

The characteristic-zero coefficient field is
\[
              F_\nu=\mathbf Q(\zeta_\nu+\zeta_\nu^{-1}),
                  \qquad [F_\nu:\mathbf Q]=d.
 \tag{A26.6}
\]
For clarity, irreducibility of $1+x+\cdots+x^{\nu-1}$ follows from
Eisenstein after replacing $x$ by $x+1$; adjoining the real sum has
index two in the cyclotomic field.  The matrices (A26.5) over $F_\nu$
give the $d$ distinct real two-dimensional irreducibles $U_j$ under
its real embeddings.  Together with $1,\varepsilon$, their squared
dimensions sum to $2\nu$, so they give the explicit split algebra
\[
               \mathbf Q[D_\nu]\simeq
                      \mathbf Q\oplus\mathbf Q\oplus M_2(F_\nu).
 \tag{A26.7}
\]
The sum $\rho+\rho^{-1}$ is central and self-adjoint and generates
the last centre.  If $e_0$ is its central isotypic idempotent, the
self-adjoint idempotents $e_\pm=e_0(1\pm\sigma)/2$ have corners
isomorphic to $F_\nu$, fixed pointwise by the adjoint involution.
These are the actual relative field blocks of C44-A24, not merely
special Frobenius eigenspaces.

\paragraph{Optional centres, compatible jets and original pencils.}
There are $a=r+1$ sign factors.  Let $S_0=E^r$ be the first $r$ of
them and choose a point $c$ with free orbit in the complementary
$E\times A$.  The translation by $T$ is inside $S_0$, so this
complement carries the linear dihedral action and has a nonempty
free locus.  The $2\nu$ fibres
\[
                            Z_g=S_0\times\{gc\}
 \tag{A26.8}
\]
are disjoint, avoid every fixed locus, and form a regular orbit.
Define them all over a finite extension, and let $Z$ be their union,
or the empty set in the version without centres.
The parameter space is $\operatorname{Hom}_{D_\nu}
(U,H^0(X,I_ZL))$, where $U$ is (A26.5).

All compatible jets at a bounded number of orbits are available for
large $t$, by the global-generation, separating-hyperplane and
averaging argument of C44-A25.  Here the orbit sizes are $2\nu$ and
$\nu$, and averaging divides by $2\nu$, a unit.  The conormal-jet
argument along the fixed smooth centres is unchanged.

At a free point a singular axis imposes codimension $n+1>n$.
On a fixed $d$-fold the bad even derivative imposes $1+d>d$
conditions including the even value; the bad odd derivative imposes
$1+3d>d$.  Along an $r$-fold centre the normal space in $X$ has rank
$r+2$, and a two-row rank failure has codimension $r+1>r$.
Thus a nonempty open has smooth axes and all pencil members smooth
along the centres and the smooth base.

The projective generator representation has $2+2\nu$ exceptional
parameters: two rotation eigenlines and the two eigenlines of each
reflection.  They are distinct by the same product-of-reflections
argument as in C44-A25.  At a fixed locus the critical codimensions
for an even and an odd eigenmember are $d+1$ and $3d$, respectively;
both exceed $d$.  A member not an eigenmember has all $n$ derivatives
available.  At a free point a fixed member has critical codimension
$n+1$.  All exceptional members can therefore be made smooth.

For the other parameters, a degenerate critical point has codimension
$n+2$ on dimension $n+1$; two points in distinct free orbits have
codimension $2(n+1)$ on dimension $2n+1$; and two points in one free
orbit give the same codimension on dimension $n+1$, since their two
parameters are distinct.  A fixed-point critical pair has codimension
$n+1$ on dimension $d+1$.  These strict inequalities, and the smooth
base, leave a proper reduced finite critical scheme with no two
critical points in one member.  Its length is the positive number
\[
                            (n+1)\int_Xh^n.
 \tag{A26.9}
\]
This is the first-jet top Chern calculation of C44-A23 in dimension
$n$.  The original pencil is therefore genuinely Lefschetz.

Lift $c$ by smoothness, so its orbit centres lift disjointly.  Relative
Serre vanishing and averaging lift each actual equivariant section
pair.  The proper nonsmooth locus has empty special fibre and is
empty.  Thus every selected smooth special axis has the actual
embedded polarized equivariant lift required in C44-A24.

\paragraph{The known Gram matrix and the two reflection configurations.}
Put $h_0=c_1(H)$, $b_0=\int_{Z_1}h_0^r>0$,
$A_0=(L_{h_0}^2)^{-1}[Z_1]_X$, and
$c_0=\int_XA_0^2h_0^2$.  The inverse is the absolute abelian Chow
inverse of C44-A15.  All centres are parallel fibres.  Their actual
primitive classes have signed Gram matrix
\[
            (r+1)t^r b_0 I_{2\nu}-(-1)^r t^{-2}c_0J_{2\nu}.
 \tag{A26.10}
\]
Indeed $N_{Z_g/X}$ is trivial and
$c(N_{Z_g/B})=(1+th_0|Z_g)^{-2}$.  The matrix is positive for
$(r+1)t^{r+2}b_0>2\nu|c_0|$.  Its span is a constant regular
rational representation $K$ of dimension $2\nu$.

Prescribing compatible two-jets gives a degeneration with exactly
one free orbit of $2\nu$ nodes, and another with exactly one
reflection-fixed orbit of $\nu$ nodes, all off the centres.  At the
latter, eliminate a nonzero odd derivative.  The remaining quadratic
variables have signs
\[
                              (d,3d-1).
 \tag{A26.11}
\]
Both invariant quadratic blocks can be nondegenerate in odd
characteristic.  The other-point jet inequalities exclude all other
singularities, and varying the remaining constant term yields a
regular transverse smoothing with general member in the required
open.

The determinant action on the local vanishing line, proved by the
equivariant quadric calculation in C44-A20, is
$(-1)^{3d-1}=(-1)^{d+1}$.  Thus the $\nu$-node representation is
\[
 \operatorname{Ind}_{C_2}^{D_\nu}(\mathrm{sgn})
       =\varepsilon\oplus\bigoplus_{j=1}^{d}U_j\quad(d\text{ even}),
 \qquad
 \operatorname{Ind}_{C_2}^{D_\nu}(1)
       =1\oplus\bigoplus_{j=1}^{d}U_j\quad(d\text{ odd}).
 \tag{A26.12}
\]
The free orbit is the regular representation.  Each node has ordinary
square $2(-1)^r$, hence signed square two; different node cycles are
orthogonal.  They lie in $W=V\cap K^\perp$ because ambient and centre
classes are inertia-fixed.  The unique quadratic tame inertia gives
reflections on $1,\varepsilon$ from the free orbit and on every $U_j$
multiplicity from (A26.12).

\paragraph{The moments retain the full orthogonal groups.}
Remove the fixed loci and the centres and take the free quotient $Y$.
It is a smooth geometrically connected $n$-fold with connected
dihedral torsor.  Four quotient points lift to at most $8\nu$
geometric points, so rank-two section evaluation gives four-wise
independent Bernoulli indicators of parameter $Q^{-2}$.
For each of the $d+2$ real characters, the weighted centred sum
$T_\chi$ has, on the smooth parameter open,
\[
              T_\chi=\operatorname{Tr}(F\mid M_\chi^{\rm un})
                                      +O(Q^{r-1/2}).
 \tag{A26.13}
\]
The character average has factor $1/(2\nu)$ and no dimension factor.
Gysin and $L_h^2$ cancel the ambient contributions as in C44-A23.
The removed fixed $d$-folds meet the axis in $(d-1)$-folds, costing
$O(Q^{d-1})$, which is smaller than the displayed error since
$d\leq r$.  The centre quotient is one $r$-fold with split torsor;
its leading subtraction $\dim(\chi)Q^r$ removes exactly the regular
$K$ contribution.

The square-character local system has one invariant fixed by
arithmetic Frobenius.  Consequently
$S_2=Q^n+O(Q^{n-1/2})$ and $S_4=O(Q^n)$.
The exact centred fourth moment is the formula of C44-A23 with
$v=Q^{-2}$.  Dividing by $Q^{4r}$ gives limit three, with diagonal
error $O(Q^{-2r})$.  Every character value is real in each embedding
of $F_\nu$, so nonsmooth parameters still contribute nonnegative
fourth powers.  Purity, the orthogonal contraction tensors, the
reflections just exhibited, and Katz's criterion give full orthogonal
geometric monodromy in all $d+2$ absolute multiplicity factors once
their ranks exceed eight.  The following rank calculation proves
that growth.

\paragraph{Exact multiplicity parities.}
For a fixed-locus component $F_0$, the equivariant Euler traces are
\[
 E_1=(2r+1)\int_Xh^n,\qquad
 E_\sigma=(-1)^{d-1}2^{2a}\int_{F_0}h^d,\qquad
 E_{\rho^j}=0\quad(1\leq j<\nu).
 \tag{A26.14}
\]
The fixed components are parallel translates and have the same
intersection degree.  Tame fixed-point traces give the displayed
Euler values.  The scalar multiplicities are
$E_1/(2\nu)\pm E_\sigma/2$, and every $U_j$ multiplicity is
$E_1/\nu$.  All are even, because
$\int_Xh^n=n!\chi(L)$, $n\geq\nu+1$, and $a\geq2$.
They have positive leading coefficients of degree $n$ in $t$;
the fixed-locus contribution has lower degree $d$.

Twice the lower ambient cohomology makes no parity change.  Thus
the multiplicity parities of $V$ are those of $H^{2r}(X)$.
Let $D$ be the latter's dimension and $c$ its trace at any nonidentity
rotation.  The exterior-character polynomials are
\[
 (1+z)^{4\nu-4},\qquad
 (1+z)^{\nu-1}(1-z)^{3\nu-3},\qquad
 (1+z)^{2\nu-4}(1+z^\nu)^2,
 \tag{A26.15}
\]
at identity, reflection and rotation, respectively.  The last formula
is the same for every nonidentity rotation because $\nu$ is prime.
In degree $2r=2\nu-4$ it gives
\[
       D\equiv0\pmod2,\qquad
       c=1+2\binom{2\nu-4}{\nu-4}\equiv1\pmod2.
 \tag{A26.16}
\]
For the first parity, reduce
$(1+z)^{4\nu-4}$ to $(1+z^4)^{\nu-1}$ modulo two; the required
degree is two modulo four.
Since $\sum_{j=1}^{\nu-1}\chi_{U_i}(\rho^j)=-2$ and reflection
characters vanish, each standard multiplicity is
\[
                              m_U=(D-c)/\nu,
 \tag{A26.17}
\]
and is odd.

The negative eigenspace of a reflection on $H^{2r}(X)$ has even
dimension.  Its binomial expression sums terms with an odd number
of factors from the $3\nu-3$ negative eigenvectors and an odd number
from the $\nu-1$ positive eigenvectors.  Both ambient numbers are
even, so both odd-index binomial coefficients are even; each product
is divisible by four.  On restriction to a reflection, every $U_i$
contributes one negative dimension.  Hence the sign multiplicity
has parity $d$, and the trivial multiplicity has the same parity
because $D$ is even.  If $b=d\bmod2$, the full parity list is
\[
                 (b,b,\underbrace{1,\ldots,1}_{d}).
 \tag{A26.18}
\]
Removing regular $K$ subtracts $(1,1,2,\ldots,2)$, giving
$(1-b,1-b,1,\ldots,1)$.

\paragraph{Selection and the complete sign.}
Choose an auxiliary coefficient prime splitting $F_\nu$ and different
from $p$.  The $d+2$ absolute factors are then available over that
coefficient field.  Their common rational or $F_\nu$-valued
Frobenius polynomials follow from C44-A24a.
Use a root-order exponent bounded by the \emph{total rational rank}
of $W$, as in C44-A25, so no coefficient-field degree is lost.
The separate projections of the connected joint geometric group
onto the special orthogonal factors are surjective.  Their finitely
many nonempty minimum-Tate selector opens therefore intersect in
that one connected group.  The constant-field finite-quotient
Chebotarev argument gives dense closed parameters attaining every
minimum simultaneously, without assuming product independence.

From (A26.18), their total residual multiplicities are
$2b+2d=\nu-\epsilon_\nu$ without centres and
$2(1-b)+2d=\nu+\epsilon_\nu$ with centres.
The $2\nu$ known directions are Frobenius-fixed, giving (A26.2).
The two scalar projectors have at most one residual rational Tate
direction each.  The two self-adjoint field corners $e_+,e_-$ each
have at most one residual direction over $F_\nu$.  Their actual
extending actions, lifts, common polynomials and positive known
subspaces satisfy C44-A24.  That theorem proves positivity on every
complete algebraic block, and therefore on $A_B$.
The Gram matrix proves its stated nonzero dimension bound.

There are infinitely many primes $\nu\ne2,p$, so the residual
multiplicities are unbounded.  All dimension increases and all
field extensions are explicit in this construction.  This proves
the theorem.
\end{proof}

\paragraph{Primary input and exact scope.}
The compatible-jet, tame fixed-point, local quadric, SGA~7 node-pairing,
Katz moment, finite-character weight and finite-field Chebotarev
inputs are exactly those checked for C44-A23--A25, now used in the
displayed dimensions.  The trace-dual field form and number-field
squareclass step remain C44-A24.  Formulae (A26.14)--(A26.18) prove
the new parity calculation directly.  No unverified equivariant
monodromy theorem or general algebraic-cycle lifting theorem has
been added.

The axes and their entire algebraic signs are actual closed examples,
not a formal choice of forms on abstract motive summands.  The known
centre span is nonzero, but no algebraicity or nonliftability of the
additional potential directions is asserted.  For cote~44, Th.~5,
scan~004 and its primitive warning, this extends the same sufficient
family in QH9-axis across dimensions; it does not settle the universal
fixed-axis question.  The axis sign is unconditional, while a full
height conclusion retains the separate ambient 032/Milne~2026 input.

% END work/scan_w/unbounded_dihedral_axis_blocks.tex

% BEGIN work/axis/axis_hodge_flatness.tex
% C44_POS /root/axis_independent. A preparatory arithmetic input, not a lifting converse.
\section{Integral Hodge flatness for the embedded axis lifts}
\label{sec:c44-a22a}

\begin{proposition}[C44-A22a: Hodge cohomology of every embedded lift]
Let $R$ be a mixed-characteristic discrete valuation ring, let
$\mathcal X/R$ be an abelian scheme of relative dimension $g=d+2$,
and let $\mathcal L$ be a relatively ample line bundle. Suppose
$\mathcal B\subset\mathcal X$ is a smooth relative complete
intersection of two sections of $\mathcal L$, of relative dimension
$d\geq1$. Then all modules
\[
 H^q(\mathcal B,\Omega^p_{\mathcal B/R})
 \tag{A22a.1}
\]
are finite free over $R$ and commute with arbitrary base change.
The relative Hodge--de Rham spectral sequence degenerates at its
first page, its abutment is finite free, and the Hodge filtration
has finite free successive quotients. These statements hold without
an unramifiedness assumption on $R$ or a lower bound on its residue
characteristic.

More precisely, on either geometric fibre $B$ of this family,
\[
 h^q(B,\Omega_B^p)=
 \begin{cases}
 \binom gp\binom gq,&p+q<d,\\[2pt]
 \binom g{d-p}\binom g{d-q},&p+q>d.
 \end{cases}
 \tag{A22a.2}
\]
For each $p$, its only remaining Hodge number is determined by
$\chi(B,\Omega_B^p)$ and is therefore constant in the family.
In particular every embedded lift used after C44-A22 satisfies
the free-Hodge-terms and degeneration hypotheses appearing in
integral crystalline obstruction comparisons such as Pan's
Proposition~4.1. The proposition does not extend that graph-cycle
comparison to general cycles and does not supply a cycle lift.
\end{proposition}

\begin{proof}
We first calculate on a geometric fibre, in arbitrary
characteristic, and write $X,B,L$ for the corresponding objects.
The canonical bundle of $B$ is $L^2|B$ by adjunction.
For every integer $a\geq3$, the twisted Koszul resolution
\[
 0\longrightarrow L^{a-2}\longrightarrow
        (L^{a-1})^{\oplus2}\longrightarrow L^a
        \longrightarrow\kappa_*(L^a|B)\longrightarrow0
 \tag{A22a.3}
\]
and the ample-line-bundle vanishing on an abelian variety give
$H^j(B,L^a|B)=0$ for $j>0$. Serre duality consequently gives
\[
                H^q(B,L^{-j}|B)=0
                     \qquad(j\geq1,\ q<d).
 \tag{A22a.4}
\]
The untwisted Koszul resolution on $X$, using
$H^q(X,L^{-a})=0$ for $a>0$ and $q<g$, gives
\[
              H^q(B,\mathcal O_B)=H^q(X,\mathcal O_X)
                        \qquad(q<d).
 \tag{A22a.5}
\]

The conormal sequence is
\[
 0\longrightarrow E=(L^{-1}|B)^{\oplus2}
   \longrightarrow F=\Omega_X^1|B\simeq\mathcal O_B^{\oplus g}
   \longrightarrow\Omega_B^1\longrightarrow0.
 \tag{A22a.6}
\]
The case $p=0$ is already (A22a.5). For $1\leq p\leq d$
the divided-power exterior resolution of (A22a.6) is
\[
 0\longrightarrow\Gamma^pE\longrightarrow
   \Gamma^{p-1}E\otimes F\longrightarrow\cdots\longrightarrow
   E\otimes\bigwedge^{p-1}F\longrightarrow
   \bigwedge^p F\longrightarrow\Omega_B^p\longrightarrow0.
 \tag{A22a.7}
\]
Here divided powers, rather than an identification with symmetric
powers involving factorials, make the resolution valid in every
characteristic. For completeness, its exactness can be checked
locally after splitting $F=E\oplus\Omega_B^1$. Decomposition by
the exterior degree in $\Omega_B^1$ reduces it to the duals of
the positive homogeneous pieces of the ordinary Koszul resolution
of the zero section in the polynomial algebra on $E^*$. Those
pieces are finite split exact complexes of free modules, so their
duals are exact over every base ring. The term of exterior degree
$p$ in the quotient supplies the last term in (A22a.7).

For $j\geq1$ the term
$\Gamma^j E\otimes\bigwedge^{p-j}F$ is a finite direct sum of
$L^{-j}|B$. It has no cohomology below degree $d$ by (A22a.4).
The hypercohomology spectral sequence of (A22a.7), putting its
$j$th negative term in degree $-j$, therefore yields
\[
 H^q(B,\Omega_B^p)\simeq
 H^q\bigl(B,\textstyle\bigwedge^p F\bigr)
                    \qquad(p+q<d).
 \tag{A22a.8}
\]
Indeed every negative-term contribution has total degree at least
$d-p$; its possible incoming or outgoing differentials cannot
affect a total degree below $d-p$. Combining (A22a.5) with
$\dim H^q(X,\mathcal O_X)=\binom gq$ proves the first line of
(A22a.2). Serre duality identifies the dual of
$H^q(B,\Omega_B^p)$ with
$H^{d-q}(B,\Omega_B^{d-p})$ and proves its second line.

Return to the relative family. Each
$\Omega^p_{\mathcal B/R}$ is a vector bundle, flat over $R$.
Its Euler characteristic is constant in the proper flat family.
For fixed $p$, all fibre cohomology dimensions except possibly
$q=d-p$ are already identical by (A22a.2). The Euler characteristic
therefore forces that remaining dimension also to be identical.

We spell out the integral base-change consequence. The derived
global sections of this vector bundle are represented by a bounded
complex $C$ of finite free $R$-modules and commute with derived
base change. For a finite $R$-module $H^q(C)$, let $t_q$ be the
number of its nonzero torsion invariant factors. The universal
coefficient sequence over the DVR gives
\[
 \dim_k H^q(C\otimes_R k)
  =\dim_{\operatorname{Frac}(R)}H^q(C\otimes_R\operatorname{Frac}(R))
                       +t_q+t_{q+1}.
 \tag{A22a.9}
\]
Constancy of every fibre dimension gives $t_q=0$ for all $q$.
Thus all $H^q(C)$ are finite free. A bounded finite free complex
with free cohomology splits in the derived category over a DVR;
alternatively, the universal coefficient sequence has no Tor
term for any base change. This proves (A22a.1) and its asserted
base-change property.

The relative Hodge spectral sequence now has finite free first
page. After tensoring with the fraction field, it degenerates at
the first page by Deligne--Illusie, Corollary~2.7, printed p.~257,
which applies to every smooth proper scheme over a characteristic-zero field.
Its first differential is therefore zero, since its target is
torsion free. The second page is consequently the same free
module as the first, so the same argument kills the second
differential; induction kills every differential. The finite
filtration on each de Rham cohomology group has these free
graded pieces and hence has a finite free abutment and finite
free successive quotients, as claimed.
\end{proof}

\paragraph{Primary inputs and the remaining arithmetic distinction.}
The only ample vanishing used is Milne, \emph{Abelian Varieties},
Chapter~I, Theorem~11.1(c), printed p.~54 of the 2008 notes,
\url{https://www.jmilne.org/math/CourseNotes/AV.pdf}.
The perfect-complex and derived-base-change input is the Stacks
Project, Section~36.32, Tag~0BDM, Lemmas~36.32.1--36.32.2 and
their cited Lemma~36.30.4; the DVR argument above proves the
additional freeness step explicitly. The exterior resolution is
verified locally in the proof and does not require a
characteristic-zero exterior-power formula.
For the generic-fibre degeneration, the exact primary locator is
Deligne--Illusie, \emph{Rel\`evements modulo $p^2$ et d\'ecomposition
du complexe de de Rham}, \emph{Inventiones} 89 (1987),
Corollary~2.7 and its proof, printed p.~257,
\url{https://publications.ias.edu/sites/default/files/Number57.pdf}.

Pan, \emph{$p$-adic deformations of graph cycles},
\url{https://arxiv.org/pdf/1610.03836}, Assumptions~4.1 and
Proposition~4.1, printed pp.~12--13, concern an isomorphism
between two already smooth lifted schemes over a square-zero
divided-power thickening. His proof uses the inverse of the
isomorphism on differential forms explicitly in \S5,
printed pp.~20--23. The freeness just proved supplies his
Hodge-cohomology hypotheses for our ambient axes; it does not
remove that isomorphism hypothesis, produce a lift of a surface
support, or turn rational vanishing of an obstruction into
integral vanishing.

% END work/axis/axis_hodge_flatness.tex

% BEGIN work/axis/axis_integral_primitive_tests.tex
% C44_POS /root/axis_independent. No arithmetic surface-lifting theorem is asserted.
\section{Primitive integral test forms without inverse-Lefschetz denominators}
\label{sec:c44-a22b}

\begin{proposition}[C44-A22b: a saturated lattice of primitive test forms]
Let $R$ be a mixed-characteristic discrete valuation ring with fraction
field $K$ and residue field $k$. Let $\mathcal X/R$ be an abelian
sixfold scheme with relatively ample line bundle $\mathcal L$.
Suppose a smooth complete intersection $\mathcal B\subset\mathcal X$
of two sections of $\mathcal L$ factors as
\[
            \mathcal B\subset\mathcal D\subset\mathcal X,
 \tag{A22b.1}
\]
where $\mathcal D/R$ is a smooth member of the same pencil.
For an embedded lift of a genuine Lefschetz axis this can always be
arranged after an allowed finite extension, by choosing a smooth
special pencil member and lifting its two coefficients.

Use the Hodge filtrations on relative de Rham cohomology and put
\[
 \mathscr T_R=
 \ker\left(F^3H^4_{\mathrm{dR}}(\mathcal B/R)
       \xrightarrow{\kappa_*}F^5H^8_{\mathrm{dR}}(\mathcal X/R)\right).
 \tag{A22b.2}
\]
The Gysin map in (A22b.2) is integral. It is surjective onto the
displayed target, its kernel $\mathscr T_R$ is finite free, and
these statements commute with base change. Moreover, reduction
followed by the Hodge graded projection gives a surjection
\[
 \mathscr T_R\otimes_R k\longrightarrow
 \ker\left(H^1(B,\Omega_B^3)
                  \xrightarrow{\kappa_*}H^3(X,\Omega_X^5)\right).
 \tag{A22b.3}
\]

If $\beta\in H^4_{\mathrm{dR}}(\mathcal B/R)\otimes_R K$ and
$a\in H^4_{\mathrm{dR}}(\mathcal X/R)\otimes_R K$ satisfy
\[
                     \beta-\kappa^*a\in F^2H^4_{\mathrm{dR}}(\mathcal B_K/K),
 \tag{A22b.4}
\]
then $\beta$ pairs to zero with every element of $\mathscr T_R$.
If $\beta$ is integral, this is an equality over $R$, and therefore
holds modulo every power of its maximal ideal. No integrality or
reduction modulo $p$ of an inverse Lefschetz projector is needed.
\end{proposition}

\begin{proof}
For a smooth Cartier divisor, the logarithmic residue sequence
\[
 0\longrightarrow\Omega^\bullet_{Y/R}
   \longrightarrow\Omega^\bullet_{Y/R}(\log Z)
   \longrightarrow\Omega^\bullet_{Z/R}[-1]\longrightarrow0
 \tag{A22b.5}
\]
gives the integral Gysin map, with Hodge shift one. Apply it to
each inclusion in (A22b.1) and compose. The resulting map has
cohomological shift four and Hodge shift two. Its associated
graded maps are the coherent Gysin maps; locally both are given
by wedging with the two differentials and taking the double
boundary, namely the local representative
$\omega\wedge df_0\wedge df_1/(f_0f_1)$. The construction and
its projection formula commute with base change. Choices of
the sign of each Cartier Gysin cause no issue: fix the signs
to agree with the positive cycle class; the two-divisor
composition has the usual codimension-two normalization.

By C44-A22a, all Hodge pieces of $\mathcal B$ are finite free
with base change. The same facts hold for the abelian scheme
$\mathcal X$. We check surjectivity on the only two relevant
graded pieces, entirely over $R$.

Choose an invariant volume form on $\mathcal X$ to trivialize
its canonical bundle. On the top graded piece the Gysin is
\[
 H^0(\mathcal B,K_{\mathcal B/R})
   =H^0(\mathcal B,\mathcal L^2|\mathcal B)
       \longrightarrow H^2(\mathcal X,\mathcal O_{\mathcal X})
       =H^2(\mathcal X,\Omega^6_{\mathcal X/R}).
 \tag{A22b.6}
\]
It is the double boundary in
\[
 0\to\mathcal O_{\mathcal X}\to\mathcal L^{\oplus2}
   \to\mathcal L^2\to\kappa_*(\mathcal L^2|\mathcal B)\to0.
 \tag{A22b.7}
\]
Relative ample vanishing on the abelian scheme shows that this
boundary is surjective. The same computation after any base
change proves the corresponding base-change assertion.

On the other graded piece, the relative tangent sequence and
adjunction give
\[
 0\to\Omega^3_{\mathcal B/R}\to
       (\mathcal L^2|\mathcal B)^{\oplus6}\to
       (\mathcal L^3|\mathcal B)^{\oplus2}\to0.
 \tag{A22b.8}
\]
The twisted Koszul computations used in C44-A22 give
\[
 H^1(\mathcal B,\mathcal L^3|\mathcal B)=0,
 \qquad
 H^1(\mathcal B,\mathcal L^2|\mathcal B)
          \simeq H^3(\mathcal X,\mathcal O_{\mathcal X}).
 \tag{A22b.9}
\]
Consequently the map
\[
 H^1(\mathcal B,\Omega^3_{\mathcal B/R})
            \longrightarrow H^3(\mathcal X,\Omega^5_{\mathcal X/R})
 \tag{A22b.10}
\]
is surjective. Its identification with the Gysin follows from
the same local double-boundary formula as C44-A22. All the
groups in (A22b.6) and (A22b.10) are finite free, and their
formation and the surjections commute with base change.

The filtration on the source of (A22b.2) has just its $F^4$
piece and its degree-three quotient. On the target, $F^7=0$
and the pieces are $F^6$ and the degree-five quotient.
Surjectivity of (A22b.6) and (A22b.10) therefore gives
surjectivity in (A22b.2). Its kernel is finite free. More
explicitly, if $C_4$ and $C_3$ denote the respective kernels
of (A22b.6) and (A22b.10), there is an exact sequence
\[
                  0\longrightarrow C_4\longrightarrow
                   \mathscr T_R\longrightarrow C_3\longrightarrow0.
 \tag{A22b.11}
\]
To lift an element of $C_3$, first lift it to $F^3$; its
Gysin lies in $F^6$ and can be cancelled using (A22b.6).
All three modules in (A22b.11) are free. Tensoring by $k$
and using the established base change gives (A22b.3), and
the same proof works for every base change.

Finally, write $w=\beta-\kappa^*a$ and take
$\omega\in\mathscr T_R$. The product of $w\in F^2$ and
$\omega\in F^3$ is in $F^5H^8_{\mathrm{dR}}(\mathcal B_K/K)=0$,
because $\mathcal B_K$ has dimension four. Projection formula
also gives
\[
 \int_{\mathcal B_K}\kappa^*a\,\omega
             =\int_{\mathcal X_K}a\,\kappa_*\omega=0.
 \tag{A22b.12}
\]
Thus $\beta$ annihilates $\mathscr T_R$. If $\beta$ is
integral, its cup product and trace are in $R$ and vanish
in $K$, so they vanish in $R$. This proves the last assertion.
\end{proof}

\paragraph{The exact surface-obstruction consequence, and what remains missing.}
Suppose the special fibre contains a smooth surface $S$ to which
C44-A22 applies. Composing (A22b.3) with its surjection gives
\[
 \mathscr T_R\otimes_R k\twoheadrightarrow
                         H^1(S,N_{S/B}^*\otimes K_S).
 \tag{A22b.13}
\]
Duality consequently embeds $H^1(S,N_{S/B})$ into
$\operatorname{Hom}_k(\mathscr T_R\otimes_R k,k)$.
This statement is entirely integral linear algebra and coherent
cohomology.

For an unramified lift over $W(k)$, the integral crystalline
class of the codimension-two smooth surface gives an integral
$\beta$ under crystalline--de Rham comparison. Its primitive
rational projection satisfies (A22b.4) whenever that primitive
class is in $F^2$. Hence its periods against (A22b.13) vanish
integrally, at every Witt truncation. This removes denominators
in the inverse Lefschetz operator from this particular period
test, including at residue characteristic three.

It is still necessary to prove that these vanishing periods
annihilate the \emph{actual} normal obstruction at each successive
divided-power thickening for an arbitrary lci surface. Pan's
graph-isomorphism theorem alone does not state that assertion.
For an arbitrary ramified lift, even the comparison of this
integral period lattice with the successive support classes
requires additional justification. Neither missing implication
is asserted here. In particular the existence of a filtered
multiple $p^Nw$ is not being used to erase a torsion obstruction.

\paragraph{Primary locators.}
The logarithmic construction and projection formula are Stacks
Project, Section~50.15, Tag~0FMU, Lemmas~50.15.2, 50.15.4,
50.15.5 and 50.15.7. The local coherent Gysin identification is
Section~50.23, Tag~0G82, Remarks~50.23.1 and 50.23.7.
The integral unramified comparison is the usual crystalline
Poincare lemma for a smooth lift; it is the identification used
explicitly in Pan, \emph{$p$-adic deformations of graph cycles},
Definition~1.6 and Assumptions~4.1, and his cited Berthelot--Ogus,
\emph{Notes on crystalline cohomology}, Theorem~7.1.

% END work/axis/axis_integral_primitive_tests.tex

% BEGIN work/axis/axis_pd_surface_obstruction.tex
% C44_POS /root/axis_independent. Entire-proof author and two independent checks passed.
\section{A square-zero divided-power obstruction calculation for surfaces}
\label{sec:c44-a27}

\paragraph{Scope.}
Step~3 identifies the canonical crystalline class by divided-power
descent; the local Cartan calculation alone would not suffice.
The application is restricted to unramified embedded lifts;
arbitrary ramified or abstract lifts are not covered.

\begin{lemma}[C44-A27: the square-zero PD surface calculation]
Let $k$ be a perfect field of characteristic $p>2$, put $W=W(k)$,
and write $R_m=W/p^{m+1}$ for $m\geq1$. Let $\mathcal Y_m/R_m$
be a smooth proper fourfold and let
\[
 S_{m-1}\subset\mathcal Y_{m-1}
       =\mathcal Y_m\otimes_{R_m}R_{m-1}
 \tag{A27.1}
\]
be a smooth closed relative surface. Write $Y$ and $S$ for the special
fibres, $N=N_{S/Y}$, and $I=p^mR_m\simeq k$.
Let $o(S_{m-1})\in H^1(S,N)\otimes_k I$ be the obstruction to
lifting (A27.1) as a closed flat subscheme of $\mathcal Y_m$.

Let $c_m$ be the canonical crystalline fundamental class of $S$,
viewed in $H^4_{\mathrm{dR}}(\mathcal Y_m/R_m)$ by comparison.
For every
$\omega\in F^3H^4_{\mathrm{dR}}(\mathcal Y_m/R_m)$, choose a preimage
in the hypercohomology of the truncated complex $\Omega^{\geq3}$,
and let $\bar\omega_3\in H^1(Y,\Omega_Y^3)$ be the image of that
chosen preimage in differential-form degree three, modulo $p$.
(In the application C44-A22a supplies degeneration and free Hodge
pieces.) For every such choice one has
\[
 \int_{\mathcal Y_m}c_m\,\omega
   =\pm\left\langle\sigma(o(S_{m-1})),\bar\omega_3\right\rangle
        \quad\hbox{in }I\subset R_m.
 \tag{A27.2}
\]
Here $\sigma:H^1(S,N)\to H^3(Y,\Omega_Y^1)$ is the coherent
semiregularity map, with the normalizations of C44-A22, and the
bracket is Serre duality. The single sign depends only on the
total-complex convention and is immaterial for vanishing.
\end{lemma}

\begin{proof}
\emph{Step 1: local lifts and the actual normal obstruction.}
Take a sufficiently fine affine cover of $Y$, refined on intersections
as necessary. Smoothness of (A27.1) and of its ambient space permits
etale local coordinates in which its ideal has two generators. Lifting
these coordinates gives smooth local closed lifts
$S_i\subset U_i\subset\mathcal Y_m$. Their ideals $J_i$ are generated
by a regular sequence $(f_{i1},f_{i2})$ that is part of relative etale
coordinates.

Two such lifts on an overlap differ by an element of
$N\otimes_k I$. Locally lift that element through
$T_Y|S\twoheadrightarrow N$ to an ambient derivation $D_{ij}$
with values in $I\mathcal O_{\mathcal Y_m}$. Thus
$\phi_{ij}=1+D_{ij}$ is an $R_m$-automorphism that carries one
local ideal to the other and reduces to the identity on
$\mathcal Y_{m-1}$. These maps can be chosen on an affine
refinement of the Cech nerve; equivalently one can use the
usual affine Cech hypercover. Nothing requires a global lift
of an individual normal section to a vector field.

The normal images $v_{ij}$ of $D_{ij}$ form a Cech one-cocycle.
Its class is exactly $o(S_{m-1})$: changing the local lifts
changes this cocycle by a normal coboundary, and its vanishing
allows the corrected ideals to glue. This is the usual
embedded deformation obstruction, not an obstruction to a
chosen resolution or to an abstract deformation of $S$.
Since $I^2=0$, composition of these automorphisms adds their
derivations. Consequently
$D_{ij}+D_{jk}-D_{ik}$ has zero normal image on triple overlaps.

Use local cohomology with support in the underlying closed set $S$.
For every locally free $\Omega^a_{\mathcal Y_m/R_m}$ its local
cohomology sheaves vanish in all support degrees except degree two:
locally the support is cut out by the regular sequence
$(f_{i1},f_{i2})$. The lower vanishing is the regular-sequence
depth assertion; the upper vanishing follows from the two-generator
localization complex. Thus de Rham cohomology with these supports
can be computed by
\[
 \mathcal H^2_S(\Omega^\bullet_{\mathcal Y_m/R_m})[-2].
 \tag{A27.3}
\]
In this complex the local fundamental class is
\[
 \nu_i=\frac{df_{i1}\wedge df_{i2}}{f_{i1}f_{i2}}
       \in\mathcal H^2_S(\Omega^2_{\mathcal Y_m/R_m})(U_i).
 \tag{A27.4}
\]
The expression denotes the generalized fraction, including its
orientation convention. It is independent of the regular generators
and is closed for the de Rham differential. This is also the
de Rham representative of the crystalline local Gysin class:
for one smooth Cartier equation the first Chern/Gysin class is
$d\log f$, and composing the two smooth Cartier Gysins gives
(A27.4). This construction is integral.
More precisely the local chain-level Gysin is
\[
 \Omega^a_{S_i/R_m}\longrightarrow
       \mathcal H^2_S(\Omega^{a+2}_{U_i/R_m}),\qquad
 \alpha\longmapsto
       \frac{df_{i1}\wedge df_{i2}\wedge\widetilde\alpha}
                       {f_{i1}f_{i2}}.
 \tag{A27.4a}
\]
A different lift of $\alpha$ differs by a term divisible by
one of the $f_{ia}$, or containing one of their differentials.
The former has only one denominator and the latter has a
repeated wedge factor, so both are zero. This map commutes
with the differentials and is natural for isomorphisms of
smooth pairs. Its value on $1$ is exactly $\nu_i$.

\emph{Step 2: the explicit Cartan cocycle.}
All differentials are relative to $R_m$, so $dI=0$. For a
derivation $D$ with values in $I$, the identity
\[
 (1+D)^*-1=d\iota_D+\iota_Dd
 \tag{A27.5}
\]
holds on the localized de Rham complexes, hence on (A27.3).
It follows directly on functions and differentials, and products
of two $I$-valued terms vanish. Put
\[
 \eta_{ij}=\iota_{D_{ij}}\nu_i
   =\frac{D_{ij}(f_{i1})\,df_{i2}
             -D_{ij}(f_{i2})\,df_{i1}}{f_{i1}f_{i2}}.
 \tag{A27.6}
\]
Up to the fixed Cech sign, $\delta\nu=d\eta$.

If $D$ has zero normal image, then
$D(f_{ia})\in I(f_{i1},f_{i2})$. In (A27.6) each resulting
summand has at most one of the two denominators, and therefore
is zero in $\mathcal H^2_S$. Thus $\eta_{ij}$ depends only on
$v_{ij}$. On a triple overlap replacing $\nu_j$ by $\nu_i$
inside contraction changes nothing, because the difference
already has coefficients in $I$. The normal cocycle condition
therefore gives $\delta\eta=0$.

Consequently $(\nu_i,\pm\eta_{ij})$ is a total cocycle of degree
four in (A27.3), with no component of differential-form degree
zero. Changing local lifts by normal sections changes this
cocycle by a total boundary, by the same calculation. The map
\[
 N\longrightarrow\mathcal H^2_S(\Omega_Y^1),\qquad
 v\longmapsto\frac{v(f_1)df_2-v(f_2)df_1}{f_1f_2}
 \tag{A27.7}
\]
followed by the local-to-global map on $H^1$ is the coherent
semiregularity map. Its Serre dual is the conormal exterior
map used in C44-A22. In particular $[\eta]$ has image
$\pm\sigma(o(S_{m-1}))$ in $H^3(Y,\Omega_Y^1)\otimes I$.

\emph{Step 3: identification with the canonical crystalline class.}
We give the PD comparison detail here because merely constructing
the preceding total cocycle would not identify it with $c_m$.
Crystalline--de Rham comparison is computed using the Cech nerve
of smooth local liftings and the divided-power envelopes of
the special-fibre diagonals over $R_m$, compatibly with the
canonical divided powers on $(p)$. On a common etale coordinate
chart $x_1,\ldots,x_4$, the diagonal ideal is generated by
$p$ and the differences $z_a=x_a'-x_a$; its PD envelope is
the PD polynomial algebra in the $z_a$, with the prescribed
base divided powers. The maps under comparison both reduce
to the identity on the special fibre, and their coordinate
differences in the present situation lie in the top ideal $I$.
The two projection maps become
chain homotopic on de Rham complexes by the PD Poincare
contraction. For one PD variable its contraction is
\[
             z^{[b]}dz\longmapsto z^{[b+1]},
 \tag{A27.8}
\]
and for several variables one tensors these contractions in
an ordered coordinate sequence. These are the same PD
polynomial contractions that prove the comparison theorem;
they commute with localization and hence with the complex
computing supports in (A27.3).

For the top Witt ideal $I$, all divided powers of order at
least two are zero. Indeed, for $a\in W$ and $j\geq2$,
\[
 v_p\bigl((p^ma)^j/j!\bigr)
       \geq mj-v_p(j!)\geq m+1,
 \tag{A27.9}
\]
because $p>2$ and $v_p(j!)\leq(j-1)/(p-1)$.
Thus evaluating a PD diagonal at
$z_a=D(x_a)\in I$ kills every term of PD degree at least two.
The degree-one part of (A27.8), evaluated on a differential
form, is precisely
\[
 \sum_a D(x_a)\,\iota_{\partial/\partial x_a}
                       =\iota_D.
 \tag{A27.10}
\]
This proves that the comparison homotopy between identity and
$1+D$ is the Cartan homotopy in (A27.5), with no division by
$p$ and no unspecified rational scalar. The calculation is
independent of the chosen etale coordinates: both sides
transform by the derivation rule, or one can use the PD
Taylor formula on a common etale refinement.

For completeness, higher comparison corrections contribute
nothing here. In the tensor contractions on a double or
higher diagonal, a two-step correction contains two
coordinate differences, or one divided power of order at
least two; after evaluation these belong to $I^2$ or
$\gamma_{\geq2}(I)$. Both are zero. A discrepancy among the
chosen overlap automorphisms on a triple intersection is
tangent to the local surface, and its single contraction
with the fundamental fraction is zero by Step~2. Hence no
additional degree-zero-form term survives in the descent
of this particular fundamental class.

Now compute the crystalline Gysin image of $1$ using the
local smooth lifts $S_i\subset U_i$ in this Cech comparison.
On each single chart its image is (A27.4), by the two Cartier
normalizations in Step~1. On an overlap, naturality of
Gysin for the lifting $\phi_{ij}$ of the identity on the
special pair compares these two representatives by the
PD diagonal homotopy just computed. Its contribution is
exactly (A27.6). The unit on the source is unchanged, and
the higher comparison contributions vanish as above.
Equivalently, apply the comparison homotopy to the chain-level
Gysin (A27.4a): the source contraction kills $1$, while the
target contraction gives $\eta_{ij}$. The homotopies are
additive on these fractions, $H_{D+D'}=H_D+H_{D'}$; every
composite cross-term has coefficients in $I^2$. This
explains the cocycle compatibility without choosing any
global splitting of a de Rham filtration or complex.
Cohomological descent therefore sends the canonical
crystalline fundamental class to $(\nu_i,\pm\eta_{ij})$.
This is an identification of the canonical class, not a
choice of a de Rham class with the same reduction. The
construction can equivalently use the already given lift
$S_{m-1}$: reduction modulo the nilpotent base ideal gives
the same crystalline PD envelopes and the same class.

\emph{Step 4: pairing with the test form.}
Represent $\omega$ by a Cech cocycle in $\Omega^{\geq3}$.
In total degree four its only components have bidegrees
$(1,3)$ and $(0,4)$. Wedging either with $\nu_i$, of form
degree two, gives zero on the relative fourfold. Wedging
$\eta_{ij}$ with the form-degree-four component also gives
zero. The sole surviving term is $\eta\wedge\omega_3$.
Its coefficients lie in $I$, so it depends only on
$\bar\omega_3$. The support-to-global trace identifies
this product with Serre duality between (A27.7) applied to
the obstruction and $\bar\omega_3$. This is (A27.2).
\end{proof}

\begin{corollary}[C44-A27a: a restricted primitive nonliftability criterion]
Let $S\subset X$ be a smooth surface in an abelian sixfold over
a perfect field $k$ of characteristic $p>2$. Choose $t$ sufficiently
large for C44-A22, and let $B\subset X$ be a smooth axis of a
genuine Lefschetz pencil in $|H^t|$ containing $S$. Suppose $S$
is geometrically nonliftable: no base extension of $S$ to an
algebraically closed field admits a smooth projective lift to
characteristic zero.

For every smooth projective embedded lift
\[
        \mathcal B\subset\mathcal X\quad\hbox{over }W(k)
 \tag{A27.11}
\]
as a two-divisor complete intersection with its given polarization,
the crystalline primitive class
\[
 w_{\mathrm{cris}}
   =[S]_{\mathrm{cris}}-
       \kappa^*(L_{H^t}^2)^{-1}\kappa_*[S]_{\mathrm{cris}}
 \tag{A27.12}
\]
does not belong to $F^2H^4_{\mathrm{dR}}(\mathcal B_{K}/K)$
under crystalline--de Rham comparison, where $K=W(k)[1/p]$.
The same assertion holds after replacing $k$ by any perfect
extension and considering an unramified lift over its Witt ring.
\end{corollary}

\begin{proof}
Suppose (A27.12) is filtered. Choose a smooth special member
of the given pencil, extending coefficients after an unramified
residue-field extension if needed, and lift the coefficients.
Properness makes the resulting relative member smooth, so
C44-A22b supplies the integral test lattice $\mathscr T_W$.
The integral crystalline class $\beta=[S]_{\mathrm{cris}}$
annihilates this lattice exactly over $W$, despite possible
denominators in the ambient inverse Lefschetz operator.

Inductively suppose $S$ has been lifted as a smooth surface
$S_{m-1}\subset\mathcal B_{m-1}$. Apply C44-A27 to the
reductions of every element of $\mathscr T_W$. Its left side
vanishes. By C44-A22b these elements map surjectively onto
the coherent primitive tests, and C44-A22 says that the
pairings against those tests inject $H^1(S,N_{S/B})$.
Thus the actual embedded obstruction is zero, so a flat
closed lift $S_m\subset\mathcal B_m$ exists. Such a lift is
smooth: flatness and the smooth special fibre give the
infinitesimal smoothness criterion. Starting with $S$, this
constructs compatible smooth embedded lifts at every order.

Formal GAGA for the projective $W$-scheme $\mathcal B$
algebraizes the compatible ideal sheaves to a closed
$\mathcal S\subset\mathcal B$. It is flat, and its smooth
locus contains the special fibre. The nonsmooth locus is
closed and proper over $W$; if nonempty its image would
meet the closed point, a contradiction. Hence $\mathcal S$
is a smooth projective characteristic-zero lift of $S$,
contrary to hypothesis.
\end{proof}

\paragraph{Class-level application.}
The surface of C44-A14 lies in $E^4$ in characteristic three;
embed it as $S\times\{0\}^2\subset E^6$. Choose $t$ large
enough for both C44-A22 and C44-A18. The latter supplies
closed genuine Lefschetz axes containing $S$ for which
$A_B=\mathbf Qw$, with $w^2>0$, and the common-realization
kernel result C44-A18b applies. These axes have unramified
embedded lifts: lift $E:y^2=x^3-x$ and the product polarization
over the Witt ring and lift the two sections using abelian
ample vanishing.

If the $\ell$-adic class $w$ were the specialization of an
algebraic generic-fibre class along any such unramified
embedded lift, C44-A18c would force (A27.12) to lie in
$F^2$. The corollary forbids this. This is therefore
an actual nonliftability statement about a nonzero primitive
cohomology class, with the stated restricted lifting
quantifier; it does not merely say that the chosen support
fails to lift. It uses the common-kernel theorem to handle
a different special representative of the same $\ell$-adic
class. The positive norm also shows that this example is
not a counterexample to the desired axis sign.
This exclusion also covers an algebraic cycle on the geometric
generic fibre of an unramified embedded model: even if the cycle
requires a ramified finite field extension of its generic field,
membership of the fixed rational crystalline class in $F^2$
descends faithfully along that field extension. This does not
exclude a new ramified deformation of the model itself.

\paragraph{Limits retained explicitly.}
There is no assertion here about an arbitrary ramified
mixed-characteristic trait, or about an abstract lift of
$B$ without the prescribed complete-intersection embedding
and polarization. The top Witt ideal has the particular
trivial PD structure used in (A27.9); the integral lattice
comparison and the first deformation steps need further
work for a ramified trait. Multiplying the primitive class
by a power of $p$ does not remove these missing steps.

\paragraph{Primary comparison and normalization locators.}
Bhatt--de Jong, \emph{Crystalline cohomology and de Rham
cohomology}, Lemma~2.8 and Remark~2.9, printed p.~4, give
the PD Taylor comparisons and independence of etale charts;
Lemma~2.13, p.~5, gives the PD polynomial Poincare comparison;
Remark~2.19, p.~7, gives its finite-level version; and
Theorem~3.6 with its proof, pp.~8--9, gives canonical global
Cech comparison. Corollary~2.10 identifies reduction
modulo $p$ through the identical PD envelopes.
These statements are not being cited as a surface
semiregularity theorem; Steps~2--4 provide that calculation.
The primary PDF is
\url{https://arxiv.org/pdf/1110.5001}.

For the local Gysin normalization, Stacks Project,
Section~50.23, Tag~0G82, Remarks~50.23.1 and 50.23.7,
gives the generalized-fraction formula and its independence.
Shiho, \emph{On logarithmic Hodge--Witt cohomology of
regular schemes}, Lemma~3.6 and Remark~3.7, printed
pp.~587--588, explicitly gives the integral smooth-Cartier
de Rham formula over $\mathbf Z/p^n$ and its relation to
the crystalline/Hodge--Witt Gysin, including the sign
convention correction; it cites Berthelot, Chapter~VI,
Proposition~3.1.3 for the de Rham calculation.
The primary PDF is
\url{https://www.ms.u-tokyo.ac.jp/journal/pdf/jms140404.pdf}.
For the canonical crystalline normalization itself, Gros,
\emph{Classes de Chern et classes de cycles en cohomologie
de Hodge--Witt logarithmique}, Chapter~II, Section~3.4,
printed pp.~43--44, diagrams~(3.4.4)--(3.4.8), compares the
crystalline Gysin with the local de Rham and Hodge--Witt
Gysins. Chapter~II, Definition~4.1.7 and the paragraph
following it, printed pp.~50--51, identify the resulting
smooth-cycle class with Berthelot's crystalline class.
The primary PDF is
\url{https://www.numdam.org/item/10.24033/msmf.322.pdf}.

% END work/axis/axis_pd_surface_obstruction.tex

% BEGIN work/axis/axis_abstract_unramified_lifts.tex
% C44_POS /root. Author and independent entire-proof readings passed.
\section{Recovering the Albanese embedding of an abstract axis lift}
\label{sec:c44-a29}

\paragraph{Scope.}
This result extends the unramified conclusion of C44-A27 to
abstract smooth projective lifts. It does not assert that an abstract
deformation remains a two-divisor complete intersection, and it does
not address the integral transport problem for a ramified deformation.

\begin{lemma}[C44-A29a: Hodge flatness and the Albanese embedding]
Let $k=\overline{\mathbf F}_p$,
let $X/k$ be an abelian sixfold, and let $B\subset X$ be a smooth
axis of a genuine Lefschetz pencil of ample class $h$.
Let $R$ be a complete mixed-characteristic discrete valuation ring
with residue field $k$, and let $\mathcal B/R$ be any smooth
projective lift of the abstract variety $B$.

All $H^q(\mathcal B,\Omega^a_{\mathcal B/R})$ are finite free,
commute with arbitrary base change, and the relative Hodge--de Rham
spectral sequence degenerates at $E_1$. There is an abelian scheme
$\mathcal X/R$ and a closed immersion
\[
                 a:\mathcal B\hookrightarrow\mathcal X
\tag{A29.1}
\]
whose special fibre is the original immersion $\kappa:B\hookrightarrow X$
up to a translation and an identified isomorphism of its target.
\end{lemma}

\begin{proof}
First choose one embedded characteristic-zero lift of the polarized
axis, which exists by C44-A05. It may use a different complete trait
and may be ramified. C44-A22a applied to that lift proves the numerical
identity
\[
              \sum_{a+q=n}h^q(B,\Omega_B^a)=b_n(B)
              \qquad\hbox{for every }n.
\tag{A29.2}
\]
Indeed its free degenerate de Rham cohomology has generic dimension
$b_n(B)$ by characteristic-zero comparison and smooth proper
$\ell$-adic base change. Thus (A29.2) is an identity concerning the
fixed special fibre, independent of a new lift.

For the given abstract lift, upper semicontinuity gives
$h^q(B,\Omega_B^a)\geq h^q(\mathcal B_K,\Omega_{\mathcal B_K}^a)$
for each pair. The sum of the generic numbers with $a+q=n$ is again
$b_n(B)$, by characteristic-zero Hodge--de Rham degeneration and
smooth proper base change. Equality of the sums in (A29.2) forces
equality in every individual inequality. The perfect-complex
base-change argument of C44-A22a now gives freeness and arbitrary
base change for every Hodge group. Explicitly, for a fixed $a$,
if $t_q$ counts the elementary torsion summands of
$H^q(\mathcal B,\Omega^a)$, the universal coefficient sequence gives
\[
 h^q(B,\Omega_B^a)=h^q(\mathcal B_K,\Omega^a)+t_q+t_{q+1}.
\tag{A29.3}
\]
Every $t_q$ is consequently zero. All differentials on the free
$E_1$ page vanish because they vanish after inverting the uniformizer;
the same argument applies successively to every page. The resulting
filtration has free quotients and commutes with base change.

We next check the Picard hypotheses integrally. For the two smooth
ample inclusions $B\subset D\subset X$, restriction gives
\[
                  \operatorname{Pic}(X)\simeq\operatorname{Pic}(B).
\tag{A29.4}
\]
Here one may choose a smooth member $D$ of the given pencil. Apply
SGA~2, XII, Corollary~3.6 to each inclusion. The required
$H^1$ and $H^2$ vanishing for every negative ample twist on $D$
and $B$ follows from the abelian negative-twist vanishing and their
one-term and two-term Koszul resolutions. All schemes are regular
and the dimensions are six, five and four, so the parafactorial
and dimension hypotheses are satisfied.

The same Koszul calculation gives
\[
 H^i(X,\mathcal O_X)\simeq H^i(B,\mathcal O_B),\qquad i=1,2.
\tag{A29.5}
\]
In particular $h^1(\mathcal O_B)=6$. Weak Lefschetz gives $b_1(B)=12$,
so the reduced Picard variety has dimension six; its tangent space
also has dimension six by (A29.5), hence $\operatorname{Pic}^0_{B/k}$
is smooth. Restriction from $X^\vee$ is an isomorphism: (A29.4)
gives injectivity on geometric points, (A29.5) gives an isomorphism
on tangent spaces, and its image has the full dimension of the
connected proper target. The kernel is therefore a finite etale
group scheme with only the identity geometric point, and is trivial.

The relative Picard scheme of the smooth projective family
$\mathcal B/R$ exists and is separated. Its generic identity component
is smooth of dimension six, by characteristic-zero Hodge theory and
$b_1=12$; its special identity component has just been identified with
$X^\vee$. Kleiman's Proposition~5.20 in \emph{The Picard scheme}
applies over the reduced base $\operatorname{Spec}R$: the identity
components form a smooth proper group scheme $\mathcal P/R$,
and its formation has these prescribed fibres. Thus $\mathcal P$
is an abelian scheme.

Choose a point of $B(k)$ and lift it to a section of $\mathcal B/R$,
using smoothness and completeness. The rigidified Poincare bundle
on $\mathcal B\times_R\mathcal P$ then gives the relative Albanese
map to the dual abelian scheme $\mathcal X=\mathcal P^\vee$.
This construction commutes with passage to either fibre. On the
special fibre its pullback on Picard varieties is the identity
under (A29.4), so its universal property identifies it with the
given $\kappa$ after translating the chosen base point to zero.
This is the usual Poincare construction of the Albanese, not an
appeal to a characteristic-zero-only relative Albanese theorem.

Finally this relative map is a closed immersion. It is proper, and
its non-quasi-finite locus is closed in the proper $R$-scheme
$\mathcal B$ and misses the special fibre; hence that locus is empty.
A proper quasi-finite map is finite. The cokernel of
$\mathcal O_{\mathcal X}\to a_*\mathcal O_{\mathcal B}$ is coherent;
its reduction is zero because the special map is a closed immersion.
Its support is proper over $R$ and misses the special fibre, so is
empty by Nakayama and properness. The finite map is therefore a
closed immersion, as claimed. Since source and target are smooth
over $R$, it is a regular immersion of relative codimension two.
\end{proof}

\begin{lemma}[C44-A29b: primitive integral tests for the recovered embedding]
For (A29.1), put
\[
 \mathscr T_R=\ker\left(F^3H^4_{\mathrm{dR}}(\mathcal B/R)
       \xrightarrow{a_*}F^5H^8_{\mathrm{dR}}(\mathcal X/R)\right).
\tag{A29.6}
\]
The map is integral and surjective, its kernel is finite free and
commutes with base change, and reduction maps the kernel onto
\[
 C_3=\ker\left(H^1(B,\Omega_B^3)
                         \xrightarrow{\kappa_*}H^3(X,\Omega_X^5)\right).
\tag{A29.7}
\]
An integral class $\beta$ such that $\beta-a^*c\in F^2$ rationally,
for some rational ambient class $c$, pairs to zero with every
element of $\mathscr T_R$, integrally and at every truncation.
\end{lemma}

\begin{proof}
There is no need for a global factorization into two Cartier
divisors. For a smooth regular immersion of codimension two, the
local Gysin map is the natural chain map already used in (A27.4a):
\[
 \alpha\longmapsto
       \frac{df_1\wedge df_2\wedge\widetilde\alpha}{f_1f_2}
       \in\mathcal H^2_{\mathcal B}
                    (\Omega^{a+2}_{\mathcal X/R})
       \quad(\alpha\in\Omega^a_{\mathcal B/R}).
\tag{A29.8}
\]
Changing the lifted form adds terms divisible by an equation or
containing its differential, which vanish in the fraction. Changing
regular generators gives the same fraction by the determinant
transformation rule. These maps therefore glue, commute with the
de Rham differentials, and give an integral Gysin with cohomological
shift four and Hodge shift two. The localization-complex construction
commutes with base change and its cup-product formula is the
projection formula. Its associated graded maps reduce to the
coherent Gysin maps of the original special embedding.

All Hodge groups of both schemes are free with base change, by
A29a and the corresponding abelian-scheme statement. Modulo the
uniformizer, the two graded maps
\[
 H^0(\mathcal B,\Omega^4)\longrightarrow H^2(\mathcal X,\Omega^6),
 \qquad
 H^1(\mathcal B,\Omega^3)\longrightarrow H^3(\mathcal X,\Omega^5)
\tag{A29.9}
\]
are surjective by the special-fibre calculations (A22b.6)--(A22b.10).
Nakayama therefore makes both maps surjective over $R$. Their
kernels are free and the surjections remain exact after every base
change. The filtered two-step argument (A22b.11) now proves (A29.6)
and (A29.7): lift a degree-three kernel element, and cancel its
degree-six Gysin by an element of $F^4$ using the first map in
(A29.9). This uses only the special complete intersection, not a
complete-intersection presentation of the deformation.

For the last claim, $(\beta-a^*c)\omega=0$ for
$\omega\in\mathscr T_R$ because $F^2F^3\subset F^5H^8=0$ on a
fourfold. Also $\int a^*c\,\omega=\int c\,a_*\omega=0$.
These are rational equalities; the left side for integral $\beta$
lies in the torsion-free ring $R$, so vanishes integrally.
\end{proof}

\begin{theorem}[C44-A29: primitive nonliftability for every unramified abstract lift]
Assume $p>2$ and the hypotheses on the surface $S\subset X$ and
the sufficiently high axis $B$ in C44-A27a, with
$k=\overline{\mathbf F}_p$. For every abstract smooth projective lift $\mathcal B/W(k)$,
the fixed crystalline class
\[
 w_{\mathrm{cris}}=[S]_{\mathrm{cris}}
       -\kappa^*(L_h^2)^{-1}\kappa_*[S]_{\mathrm{cris}}
\tag{A29.10}
\]
does not lie in $F^2H^4_{\mathrm{dR}}(\mathcal B_K/K)$.
Consequently, for the closed A18 axes containing the A14 surface,
their nonzero primitive $\ell$-adic class cannot be the specialization
of an algebraic class on the geometric generic fibre of any
unramified abstract lift. The class has positive norm and is not a
counterexample to the axis sign.
\end{theorem}

\begin{proof}
Apply A29a to recover the smooth Albanese embedding over $W(k)$
and A29b to obtain its integral primitive tests. The special fibre
is the same $\kappa$, so C44-A22 still makes pairing against these
tests injective on $H^1(S,N_{S/B})$.
Let $\beta=[S]_{\mathrm{cris}}$, which is integral under the
unramified crystalline--de Rham comparison. If (A29.10) is filtered,
the projection formula in A29b makes all its surface periods zero
integrally. At every step $W/p^{m+1}\to W/p^m$, C44-A27 identifies
these periods with the actual embedded obstruction. The same
induction and formal GAGA argument as in A27a would give a smooth
projective characteristic-zero lift of the geometrically nonliftable
$S$. This contradiction proves the crystalline assertion. No
complete-intersection presentation of $\mathcal B$ has entered.

We supply the class-level comparison, because the generic fibre
of (A29.1) was not asserted to remain a complete intersection.
On degree-four ambient cohomology the actual relative map
\[
 M=a_*a^*:H^4(\mathcal X)\longrightarrow H^8(\mathcal X)(2)
\tag{A29.11}
\]
has special fibre $L_h^2$, hence is an isomorphism in the relevant
rational realizations. On the generic abelian variety it has an
algebraic inverse. Indeed choose any polarization there and let
$J:H^8\to H^4$ be its algebraic inverse Lefschetz operator. The
endomorphism $JM$ is invertible; its characteristic polynomial is
rational by Betti realization and has nonzero constant term. Its
inverse is a rational polynomial in $JM$. Thus
\[
                      M^{-1}=(JM)^{-1}J
\tag{A29.12}
\]
is an actual rational algebraic correspondence. All cycles involved
may be defined over a finite extension of the generic field.
Taking their closures and specializing identifies (A29.12)
cohomologically with $(L_h^2)^{-1}$; the inverse identities
specialize in every realization. Equivalently, apply unramified
comparison to the fixed model and then extend its generic field.

The algebraic generic operator
\[
                         P_K=1-a^*M^{-1}a_*
\tag{A29.13}
\]
therefore specializes cohomologically to the fixed axis projector
$P_0=1-\kappa^*(L_h^2)^{-1}\kappa_*$, both etale and crystalline.
Suppose an algebraic generic class specializes etale to $w$ and
let $z$ be a special cycle representative obtained by specialization.
Then $P_0z-w$ is a cycle in the projected group with zero etale
class. The common-realization kernel theorem C44-A18b for the
selected A18 axis gives
$P_{0,\mathrm{cris}}\operatorname{cl}_{\mathrm{cris}}(z)
 =w_{\mathrm{cris}}$.
On the other hand, applying (A29.13) to the generic algebraic class
gives an algebraic class, hence a class in $F^2$, whose crystalline
specialization is precisely this vector. This contradicts the
first assertion. If the generic cycle or correspondence required
a finite field extension, filtration membership of the fixed vector
descends faithfully to $K$. This proves the stated geometric-generic
class exclusion for every unramified abstract model.
\end{proof}

\paragraph{Exact sources and remaining scope.}
The Picard restriction input is SGA~2, Expose~XII,
Corollary~3.6, original printed p.~153, recomposed edition p.~121:
\url{https://arxiv.org/pdf/math/0511279}.
The relative Picard result is Kleiman, \emph{The Picard scheme},
Proposition~5.20 and its proof, author pp.~47--48;
Theorem~4.8, author pp.~27--28, supplies representability for a projective family.
The Poincare bundle, abelian duality and Albanese construction are
Exercise~4.3, author p.~27 with its answer on p.~65, and
Remarks~5.24--5.25, author pp.~50--51;
the same rigidified-bundle construction over the present base was
spelled out above:
\url{https://arxiv.org/pdf/math/0504020}.
Perfect-complex base change and generic degeneration are the exact
inputs already checked in A22a. The integral local Gysin construction
and canonical crystalline obstruction comparison are those checked
in A27; A29b extends the test-lattice argument by Nakayama.

The preparation A29a--A29b works for ramified abstract lifts too.
The final obstruction induction uses the unramified integral
comparison and the top Witt divided powers. A genuinely ramified
deformation still requires a comparison with its successive support
classes; nothing here removes that remaining gate. No assumption
that every abstract deformation is a two-divisor complete
intersection is made or needed.

% END work/axis/axis_abstract_unramified_lifts.tex

% BEGIN work/axis/axis_ramified_finite_order.tex
% C44_POS /root/axis_independent. Author and independent whole-proof readings passed.
\section{The finite initial range in a ramified surface-lifting problem}
\label{sec:c44-a31}

\paragraph{Review status and scope.}
This checked proof extends the obstruction calculation of
C44-A27 after an explicitly prescribed initial surface lift.  It does not
assert that this initial lift exists on an arbitrary ramified model.  The
conclusion is stronger than merely clearing a denominator in the primitive
projector, and its additional hypothesis is correspondingly geometric.

\begin{theorem}[C44-A31: an explicit finite initial gate]
Let $p>2$, let $k=\overline{\mathbf F}_p$, and let
$S\subset B\subset X$ satisfy C44-A27a: $S$ is a smooth geometrically
nonliftable projective surface in an abelian sixfold, and $B$ is a
sufficiently ample smooth axis containing $S$ for which the primitive
semiregularity assertion C44-A22 holds.  Let $R$ be a complete
mixed-characteristic discrete valuation ring with residue field $k$,
uniformizer $\pi$, and absolute ramification index $e=v_\pi(p)$.
Let $\mathcal B/R$ be any smooth projective lift of the abstract variety
$B$, and put
\[
             n_0=\left\lfloor\frac{e}{p-1}\right\rfloor+1,
 \qquad B_n=\mathcal B\otimes_R R/(\pi^n).
\tag{A31.1}
\]
Suppose there is a smooth closed surface
\[
                S_{n_0}\subset B_{n_0}
\tag{A31.2}
\]
over $R/(\pi^{n_0})$, with special fibre the specified immersion
$S\subset B$.  Then the fixed rational crystalline primitive class
\[
 w=[S]-\kappa^*(L_h^2)^{-1}\kappa_*[S]
\tag{A31.3}
\]
does not belong to $F^2H^4_{\mathrm{dR}}(\mathcal B_K/K)$ under the
Berthelot--Ogus comparison.

In particular the assertion holds for every such model with $e<p-1$,
since then $n_0=1$ and (A31.2) is the original surface.  For the closed
A18 axes containing the A14 surface, under hypothesis (A31.2) the
nonzero primitive etale class also cannot specialize from an algebraic
class on the geometric generic fibre of $\mathcal B$.
\end{theorem}

\begin{proof}
We first specify exactly which integral cycle class is used.  The ideal
$J=(\pi^{n_0})\subset R$ has its canonical divided powers
$\gamma_j(x)=x^j/j!$, which are $\pi$-adically nilpotent because
$n_0>e/(p-1)$.  This is Berthelot--Ogus, Lemma~3.9.  Moreover
$n_0\leq e$ for $p>2$, so $A_0=R/(\pi^{n_0})$ is a local Artin
$k$-algebra of characteristic $p$.
Here the coefficient $k$-algebra structure is induced by
$W(k)\to R$: since $p\in(\pi^{n_0})$, this map factors through
$W(k)/p=k$.

The crystalline site of $B_{n_0}$ over the PD base $(R,J)$ has the
canonical integral comparison
\[
 H^4_{\mathrm{cris}}(B_{n_0}/R)
        \simeq H^4_{\mathrm{dR}}(\mathcal B/R).
\tag{A31.4}
\]
Here and below the left side means the compatible $p$-adic limit;
one can equivalently work at each quotient and then pass to the limit.
The smooth regular immersion (A31.2) has its integral crystalline
Gysin class.  Denote its image under (A31.4) by
\[
                   \beta'\in H^4_{\mathrm{dR}}(\mathcal B/R).
\tag{A31.5}
\]
The construction is the natural local generalized-fraction Gysin of
(A27.4a), now over this PD base.  In particular it is a specified
integral class, not an arbitrarily chosen lift of a reduced class.

We check the rational identity
\[
                  \beta'\otimes K=[S]_{\mathrm{cris}}\otimes K
\tag{A31.6}
\]
under the fixed Berthelot--Ogus comparison.  Apply their
Theorem~2.1 and Corollary~2.2 over
$\operatorname{Spec}k\subset\operatorname{Spec}A_0
 \subset\operatorname{Spf}R$.
Both $B_{n_0}$ and $S_{n_0}$ are smooth deformations of their special
fibres.  The resulting rational comparisons are compatible with the
pullback by the immersion, by (2.2.4), and with cup products: they
are composites of Frobenius pullbacks and their inverses.  They
also preserve trace.  To see the normalization without choosing a
splitting, restrict an ample bundle to either smooth proper scheme.
The comparison preserves its first Chern class by Proposition~3.5;
its top power has nonzero integral degree, and the trace on that
power is this degree on both sides.  Indeed, take a sufficiently
large power of the ample bundle on the given smooth projective
$A_0$-scheme.  Serre vanishing on each layer of the finite
nilpotent ideal filtration makes restriction of its sections to
the special fibre surjective.  Choose general sections there
whose common zeros are transverse and zero-dimensional, and
lift these sections to $A_0$.  The relative Jacobian criterion
then makes their common zero scheme finite etale over $A_0$;
its local residue trace is its degree.  This requires no lift
of $S_{n_0}$ beyond $A_0$.  Top cohomology on each connected
component is one-dimensional after inverting $p$, by the same
comparison with the special fibre.  This proves trace compatibility.
The adjunction formula for the natural Gysin now gives (A31.6):
for every degree-four class $x$ on the ambient fourfold,
\[
       \int_{B_{n_0}} [S_{n_0}]\,x
          =\int_{S_{n_0}}i^*x
          =\int_S i_0^*x_0
          =\int_B[S]x_0.
\tag{A31.7}
\]
The perfect rational Poincare pairing on $B$ identifies the two
classes.  To identify the comparison with the usual one, first
use the first triangle of Berthelot--Ogus, Remark~2.6(i), with
$n=n_0$ and $n'=e$.  It identifies the integral comparisons for
$B_{n_0}$ and $B_e$ with de Rham cohomology.  Their functorial
rational comparisons are compatible with this restriction by
Remark~2.3(i).  The second diagram of Remark~2.6(i), applied at
$n=e$, then identifies the result with the usual comparison
defined using the $(p)$-thickening.  This verifies (A31.6) for the fixed comparison,
rather than for a new transport selected to fit $S_{n_0}$.

Assume, for contradiction, that (A31.3) is filtered.  The preparatory
lemmas A29a--A29b apply to the present possibly ramified abstract
model: they recover the relative Albanese embedding and its integral
primitive test lattice $\mathscr T_R$.  All its relevant cohomology
and Hodge quotients are free.  By (A31.6), the projection formula and
$F^2F^3=0$ on a fourfold give
\[
                  \int_{\mathcal B/R}\beta'\omega=0
                         \quad(\omega\in\mathscr T_R).
\tag{A31.8}
\]
These equalities hold integrally, since their left sides are in the
torsion-free ring $R$.  No integral assertion about
$(L_h^2)^{-1}$ is made.

We now extend the prescribed surface successively.  Suppose
$m\geq n_0$ and a smooth embedded lift $S_m\subset B_m$ reducing
to $S_{n_0}$ has been constructed.  In
$R/(\pi^{m+1})\to R/(\pi^m)$ the top ideal
$I_m=(\pi^m)/(\pi^{m+1})$ has square zero, is annihilated by the
maximal ideal, and inherits \emph{trivial higher divided powers}
from the canonical PD ideal $J$.  Indeed, for $j\geq2$,
\[
 v_\pi\!\left(\frac{\pi^{mj}}{j!}\right)
      =mj-e v_p(j!)
      \geq m+(j-1)\left(m-\frac{e}{p-1}\right)
      >m.
\tag{A31.9}
\]
The valuation is an integer, so the strict final inequality implies
that it is at least $m+1$.  The quotients
used here inherit the PD structure, since the same inequality
makes their defining ideals stable under all divided powers.

The canonical local calculation of A27 therefore applies over this
PD base with $I_m$ in place of the top Witt ideal.  More explicitly,
choose local smooth lifts of $S_m$ to $B_{m+1}$.  Their overlap
automorphisms $1+D_{ij}$ have coefficients in $I_m$.  The natural
chain Gysin sends $1$ to
$df_1\wedge df_2/(f_1f_2)$; its PD comparison homotopy evaluates to
ordinary contraction $\iota_{D_{ij}}$, since (A31.9) kills every
divided power of degree at least two.  All composite corrections
have coefficients in $I_m^2=0$.  Thus the degree-one overlap term
is exactly
\[
 \frac{D_{ij}f_1\,df_2-D_{ij}f_2\,df_1}{f_1f_2},
\tag{A31.10}
\]
the semiregularity image of the actual normal obstruction.

The class in this calculation is the reduction of the fixed
$\beta'$.  This last compatibility is essential: the local
crystalline Gysin is natural for PD base change, and comparison
using any finite-level lift of the same pair sends its source
unit to the displayed fraction.  The PD comparison diagrams of
Remark~2.6(i), at finite level, commute for successive truncations.
Consequently computing the descent with $S_m$ instead of
$S_{n_0}$ does not alter the original class; the source comparison
homotopy annihilates $1$.  This is exactly the canonical-descent
argument of A27, with the original PD base $(R,J)$ retained.

Pairing with the primitive tests and using (A31.8) therefore shows
that the semiregularity image of the obstruction to extending
$S_m$ is zero.  Reduction of $\mathscr T_R$ surjects onto the
coherent primitive test space by A29b.  The injectivity statement
of A22 forces the obstruction itself to vanish.  A smooth
embedded $S_{m+1}$ exists.  Starting with (A31.2), induction gives
a compatible formal smooth surface in the projective $\mathcal B$.
Formal GAGA algebraizes its ideal.  The resulting closed subscheme
is smooth over $R$, since its formal completion is smooth and the
closed nonsmooth locus of a proper scheme cannot miss the special
fibre.  Its generic fibre is a smooth projective characteristic-zero
lift of $S$, contradicting geometric nonliftability.

For the final assertion use the class-level argument of A29: the
recovered Albanese embedding has an algebraic generic primitive
projector obtained by an abelian inverse Lefschetz map and
Cayley--Hamilton.  Its specialization is the fixed projector in
every realization.  The common-kernel lemma A18b for the selected
axis identifies the crystalline specialization of a putative
algebraic generic primitive class with (A31.3).  Algebraicity
would place that class in $F^2$, contradicting what was just
proved.  A finite extension used to define the generic cycle does
not affect the argument: filtration membership of this fixed
vector descends to $K$.
\end{proof}

\paragraph{What the remaining ramified gate says.}
If an arbitrary ramified model has the filtered primitive class,
then the particular surface cannot lift smoothly as an embedded
surface even through $R/(\pi^{n_0})$.  For a selected A18 axis,
the same failure is therefore necessary for any model on whose
geometric generic fibre the primitive etale class lifts algebraically.
When $p=3$ and $e=2$,
$n_0=2$; the unresolved part is already the first equicharacteristic
infinitesimal deformation.  It is not removed by assigning a
different divided-power structure to that first square-zero ideal.
For example, in $R=\mathbf Z_3[\pi]/(\pi^2-3)$ the canonical
structure has
$\gamma_3(\pi)=\pi/2$, nonzero modulo $\pi^2$, whereas an
artificial trivial PD structure has zero third divided power.
These define different integral transports.  The proof above
starts only where the \emph{inherited canonical} structure becomes
nilpotent and its successive top restrictions become trivial.

\paragraph{Primary-source locators.}
P.~Berthelot and A.~Ogus, \emph{$F$-isocrystals and de Rham
cohomology. I}, Invent.~Math.~72 (1983), 159--199:
Theorem~2.1 and complete Frobenius construction, pp.~167--168;
Corollary~2.2 and Remark~2.3, pp.~168--169; Theorem~2.4 and its
complete proof, pp.~169--170; Remark~2.6(i), p.~171;
Proposition~3.5, Corollaries~3.6--3.7 and Remark~3.7.1,
pp.~178--179; Lemmas~3.9--3.10 and Theorem~3.8 proof,
pp.~179--180.  The original diagrams on pp.~171, 178 and 180
were visually inspected in this run.
\url{https://gdz.sub.uni-goettingen.de/download/pdf/PPN356556735_0072/LOG_0013.pdf}.
The local Gysin normalization and PD diagonal comparison are the
primary inputs already read and checked in C44-A27, not an
additional ramified semiregularity theorem imported from the
literature.

% END work/axis/axis_ramified_finite_order.tex

% BEGIN work/axis/axis_ramified_translation_correction.tex
% C44_POS /root/axis_independent. Scope lemma; no separate progress-clock event.
\section{An explicit first ramified correction in codimension two}
\label{sec:c44-a31a}

\begin{lemma}[C44-A31a: ordinary contraction misses a canonical PD term]
Let $p>2$, let $k$ be a perfect field of characteristic $p$, and set
\[
 R=W(k)[\pi]/(\pi^{p-1}-p),\qquad A=R/(\pi^2),
 \qquad I=(\pi)/(\pi^2)\subset A.
\tag{A31a.1}
\]
Equip $(\pi)\subset R$ and its quotient $I$ with the canonical
divided powers.  On the smooth affine plane over $A$, compare the
two regular pairs $(x,y)$ and $(x+\pi,y)$.  For their codimension-two
generalized-fraction Gysin forms, the canonical one-variable PD
comparison homotopy has degree-one term
\[
 \eta_{\mathrm{PD}}
       =\log(1+\pi/x)\frac{dy}{y}
       =\frac{\pi\,dy}{xy}+\frac{\pi\,dy}{x^p y}
       \quad\pmod{\pi^2}.
\tag{A31a.2}
\]
Ordinary square-zero Cartan contraction gives only
$\eta_{\mathrm{lin}}=\pi\,dy/(xy)$.
Their difference is a nonzero closed, nonexact element of
\[
 \mathcal H^2_{(x,y)}
     (\Omega^\bullet_{A[x,y]/A})\otimes_A I
\tag{A31a.3}
\]
in differential-form degree one.  Equivalently it gives a
nonzero local supported de Rham class of total degree three.
The same calculation holds after adjoining two further smooth
coordinates, so it applies to a local codimension-two surface
inside a smooth fourfold.
\end{lemma}

\begin{proof}
The ring $R$ is a complete DVR with absolute ramification
$e=p-1$.  Berthelot--Ogus, Lemma~3.9, puts the ideal $(\pi)$
exactly at the boundary $1=e/(p-1)$: its canonical divided
powers exist, but are not $\pi$-adically nilpotent.  Explicitly,
\[
 v_\pi(\gamma_j(\pi))
       =j-(p-1)v_p(j!)=s_p(j),
\tag{A31a.4}
\]
where $s_p(j)$ is the sum of the base-$p$ digits of $j$.
For $j=p^a$ the valuation is always one.  The quotient by
$(\pi^2)$ inherits the structure, since $(\pi^2)$ is PD-stable;
the same valuation estimate with $2j$ instead of $j$ makes
every divided power of an element of $(\pi^2)$ remain in that
ideal.  In particular $I^2=0$ does not imply trivial higher
divided powers on $I$.

Nevertheless the PD logarithm converges.  In the $\pi$-adically
completed localization in which $x$ is invertible it is
\[
 \log(1+\pi/x)
  =\sum_{j\ge1}(-1)^{j-1}(j-1)!\gamma_j(\pi/x)
  =\sum_{j\ge1}\frac{(-1)^{j-1}\pi^j}{j x^j}.
\tag{A31a.5}
\]
The coefficient valuation is $j-(p-1)v_p(j)$, which tends to
infinity.  It equals one exactly for $j=1$ and $j=p$.
Indeed, writing $j=p^a u$ with $p\nmid u$, the case $a=0$
is immediate; for $a=1$ equality requires $u=1$; for $a\ge2$
one has $p^a-a(p-1)>1$.  Every other coefficient vanishes
modulo $\pi^2$.  The $p$th coefficient is positive because
$p$ is odd and equals $\pi^p/p=\pi$.  This proves the second
equality in (A31a.2).  At each finite $p$-torsion level the
PD logarithm is in fact a finite sum because the factorial
coefficients eventually vanish; convergence does not require
PD-nilpotence of $I$.

The first equality in (A31a.2) is the standard PD contraction
evaluated along the displacement $x\mapsto x+\pi$, with the
second coordinate fixed.  Before evaluation, its primitive
for $d\log(x+z)-d\log x$ is the PD logarithm
$\log(1+z/x)$.  Taking the product with $d\log y$ gives the
natural comparison for
$d\log x\wedge d\log y$ and
$d\log(x+z)\wedge d\log y$.  This uses the local chain Gysin
normalization of A27.  By contrast, the ordinary contraction
by $D=\pi\partial_x$ is exactly $\eta_{\mathrm{lin}}$.
Both homotopies have the same differential, since their
difference
\[
                    \xi=\pi x^{-p}y^{-1}dy
\tag{A31a.6}
\]
is closed: the $x$ derivative contributes $-p\pi=0$ in $A$,
and the $y$ derivative has a repeated $dy$.

For nonexactness, identify the module of generalized fractions
with
\[
 H^2_{(x,y)}(A[x,y])
   =A[x^{\pm1},y^{\pm1}]
          /(A[x^{\pm1},y]+A[x,y^{\pm1}]).
\tag{A31a.7}
\]
Its $I$-part has $k$-basis $\pi x^{-a}y^{-b}$ for $a,b\ge1$.
Thus $\xi$ is nonzero.  A $dy$ coefficient of an exact
degree-one form comes from differentiating a degree-zero
fraction in $y$.  Such differentiation cannot produce
$x^{-p}y^{-1}dy$: it would require a $y^0$ term, whose
derivative is zero and whose fraction is already zero in
(A31a.7).  No $x$ derivative supplies a $dy$ coefficient.
Therefore $\xi$ is not exact.  The local cohomology of the
locally free form modules is concentrated in support degree
two, giving the asserted total degree.  Adjoining further
coordinates preserves this class, for example by evaluation
of those coordinates at zero and pullback of forms.
\end{proof}

\paragraph{What this calculation does and does not establish.}
The example is a local comparison calculation, not a global
axis or a wrong-sign cycle.  It proves that the canonical
first ramified comparison cannot be replaced by ordinary
Cartan contraction merely because the top ideal has square
zero or the immersion has codimension two.  The extra term
is a Frobenius-type closed class in the same form degree as
the semiregularity term.  It is exactly the sort of correction
excluded by the strict inequality $m>e/(p-1)$ in A31.
No implication from a rational cycle lift to a lift of its
chosen support is made.

\paragraph{Primary normalization actually checked.}
Berthelot--Ogus, \emph{$F$-isocrystals and de Rham cohomology. I},
\S3.1, printed p.~175, gives the PD logarithm
$\sum(-1)^{j-1}(j-1)!\gamma_j(x)$ and its finite-level meaning;
Lemma~3.9, p.~179, gives the exact PD existence and nilpotence
threshold.  Both original pages were read, and p.~175 was
visually reinspected for this calculation.
\url{https://gdz.sub.uni-goettingen.de/download/pdf/PPN356556735_0072/LOG_0013.pdf}.
The product with the second Cartier Gysin is the explicit
new calculation above.  This scope lemma is not a separate
staged progress-clock event.

% END work/axis/axis_ramified_translation_correction.tex

% BEGIN work/axis/axis_ramified_support_counterexample.tex
% C44_POS /root/axis_independent. Scope example; author and root whole-proof readings passed.
\section{A ramified rational class lift need not lift a semiregular support}
\label{sec:c44-a31b}

\begin{proposition}[C44-A31b: a support-lifting shortcut is false]
Let $p\equiv3\pmod4$ and $k=\overline{\mathbf F}_p$.
There are a complete mixed-characteristic DVR $R$ with residue
field $k$, a smooth projective abelian fourfold $\mathcal A/R$,
and a smooth surface $S\subset A=\mathcal A_k$ such that:
\begin{enumerate}
\item the rational algebraic cohomology class $[S]$ is the
specialization of an explicit rational algebraic cycle on
$\mathcal A_{\operatorname{Frac}R}$;
\item the ordinary Bloch semiregularity map
\[
 H^1(S,N_{S/A})\longrightarrow H^3(A,\Omega_A^1)
\tag{A31b.1}
\]
is injective;
\item $S$ has no flat closed lift inside $\mathcal A$, even
after a finite extension of complete DVRs.
\end{enumerate}
One may take $R$ of absolute ramification $p+1$. For this
choice there is already no flat closed lift of $S$ inside
$\mathcal A\otimes_RR/(\varpi^2)$.

This is a counterexample to a support-lifting shortcut, not a
nonliftable rational class. The fourfold $A$ here is abelian,
not a Lefschetz axis, and $S$ is not asserted to represent a
nonzero class in an axis kernel $A_B$.
\end{proposition}

\begin{proof}
Put $W=W(k)$ and let
\[
 \mathcal E_0/W:\quad y^2=x^3-x.
\]
Its discriminant is a unit for our odd $p$. Fix a square root
$\eta$ of $-1$ in $W$, and let
$\iota(x,y)=(-x,\eta y)$. The special elliptic curve $E/k$
is supersingular, as checked in A14b. Since $p\equiv3\pmod4$,
$\mathbf Z_p[\iota]$ is the integer ring $\mathcal O$
of the unramified quadratic extension of $\mathbf Q_p$.

Choose a cyclic subgroup $H$ of order $p$ in
$\mathcal E_{0,\overline{\operatorname{Frac}W}}[p]$.
No such subgroup is $\iota$-stable: on the two-dimensional
$\mathbf F_p$-space of $p$-torsion, $\iota$ satisfies the
irreducible polynomial $T^2+1$, so it has no invariant line.
After a finite extension of the complete DVR, $H$ is defined
over the fraction field. Its schematic closure $\mathcal H$
in $\mathcal E_0[p]$ is finite flat of rank $p$: the closure
is finite and its coordinate ring is torsion-free over the
DVR. It remains a subgroup by schematic closure of the
generic subgroup identities.

Form the quotient elliptic scheme
\[
 q:\mathcal E_0\longrightarrow
             \mathcal E'=\mathcal E_0/\mathcal H.
\tag{A31b.2}
\]
The quotient exists as a scheme and $q$ is a finite locally
free torsor of degree $p$, by Stacks, Tag~07S7. The quotient
is proper and flat. Its generic fibre is an elliptic curve,
and its special fibre is $E/\ker F\simeq E^{(p)}=E$:
a supersingular elliptic curve has a unique subgroup scheme
of rank $p$, namely $\ker F\simeq\alpha_p$.
The smooth-fibre criterion therefore makes the quotient
smooth over $R$, and its zero section makes it an elliptic
scheme; three times the zero section supplies a relative
projective embedding. We use the displayed identification of its special
fibre, under which $q_k=F$.

Let $\widehat q$ be the dual isogeny, and define the actual
endomorphism
\[
 f=q\circ(-\iota)\circ\widehat q\in
                    \operatorname{End}_R(\mathcal E').
\tag{A31b.3}
\]
The relative Frobenius satisfies $F\iota=-\iota F$, since
$\eta^p=-\eta$. As $F\widehat F=[p]$, we obtain
\[
                          f_k=p\iota.
\tag{A31b.4}
\]
On the generic fibre the rational endomorphism
\[
 v=f/p=q\circ(-\iota)\circ q^{-1}
\tag{A31b.5}
\]
is not an integral endomorphism even over the algebraic
closure of the fraction field. Indeed, if it were, the
identity $vq=q(-\iota)$ would imply
$\iota(H)\subset H$, contrary to the invariant-line
calculation. This argument survives every field extension.
Equivalently, the multiplier ring of the lattice
$\mathcal O+p^{-1}\mathbf Z_p t$, for a generator $t$ of
the chosen torsion direction, is
$\mathbf Z_p+p\mathcal O$. This is the conductor-$p$
order of the quasi-canonical lift.

There is consequently no endomorphism of $\mathcal E'$,
even after finite extension of $R$, whose special fibre is
$\iota$. If $g$ were such an endomorphism, then
$pg-f$ would reduce to zero. Milne's rigidity
Proposition~20.1 makes this morphism factor through the
base; because it preserves the zero section, it is zero.
Thus $pg=f$ and $g=v$ on the generic
fibre, which has just been excluded. No surjectivity of a
specialization map on homomorphisms has been used.

We now construct the surface and the rational cycle lift.
Set
\[
 \mathcal A=(\mathcal E')^4,\qquad
 S=\{(x,\iota x,z,0):x,z\in E\}\subset E^4=A.
\tag{A31b.6}
\]
Thus $S\simeq E^2$ is a smooth abelian surface. On the first
two factors write $D_1=\{0\}\times E$ and
$D_2=E\times\{0\}$. For an elliptic endomorphism $u$ define
its cross class by
\[
 R(u)=[\Gamma_u]-\deg(u)D_1-D_2.
\tag{A31b.7}
\]
It lies in $H^1(E)\otimes H^1(E)$ and is additive in $u$.
For example, pull back the point class by
$(x,y)\mapsto y-u(x)$; the two pure Kunneth terms are
$\deg(u)D_1$ and $D_2$, while the mixed term is linear in
the induced map on $H^1$. This proof gives the same identity
in etale and crystalline cohomology. Hence
\[
 [\Gamma_\iota]=D_1+D_2+\frac1pR(p\iota).
\tag{A31b.8}
\]
Use the actual relative endomorphism $f$ of (A31b.3) to
define the rational relative divisor cycle
\[
 \mathcal Z=D_1+D_2+\frac1p
       \bigl([\Gamma_f]-p^2D_1-D_2\bigr)
   =\frac1p[\Gamma_f]-(p-1)D_1+
                       \frac{p-1}{p}D_2.
\tag{A31b.9}
\]
The notation $D_1,D_2$ on the right denotes their relative
versions on $(\mathcal E')^2$. By (A31b.4)--(A31b.8), the
specialization of its cohomology class is $[\Gamma_\iota]$.
Taking its exterior product with the whole third elliptic
factor and the zero section of the fourth gives an explicit
rational relative codimension-two cycle whose special
cohomology class is $[S]$. Its generic class is algebraic;
in particular the rational crystalline class of $S$ lies
in $\operatorname{Fil}^2$ for this model.

For semiregularity, both $N_{S/A}$ and $K_S$ are trivial.
The conormal sequence is the constant split sequence of
vector spaces attached to the inclusion of the abelian
subvariety $S$. The dual of (A31b.1), with the local
generalized-fraction normalization used in A22/A27, is
\[
 H^1(A,\Omega_A^3)\longrightarrow
                  H^1(S,N_{S/A}^{\vee}\otimes K_S).
\tag{A31b.10}
\]
It is the tensor product of restriction
$H^1(A,\mathcal O_A)\twoheadrightarrow
H^1(S,\mathcal O_S)$ and the surjection of constant
coefficient spaces
\[
 \bigwedge^3 H^0(A,\Omega_A^1)
       \twoheadrightarrow
 H^0(S,N_{S/A}^{\vee})\otimes H^0(S,K_S).
\tag{A31b.11}
\]
Restriction is split by projection to factors one and three.
For (A31b.11), split the cotangent space into the two tangent
and two conormal directions and project its third exterior
power onto the summand with two tangent directions and one
conormal direction. Thus (A31b.10) is surjective. Serre
duality proves the required injectivity of (A31b.1).

Suppose now that a flat closed lift $\mathcal S\subset
\mathcal A$ exists, after any finite extension of complete
DVRs if necessary. It is smooth, since its special fibre is
smooth and properness excludes a nonsmooth locus supported
only on the generic fibre. Project it to factors one and
three:
\[
 t:\mathcal S\longrightarrow(\mathcal E')^2.
\tag{A31b.12}
\]
The special map is an isomorphism. The non-quasi-finite
locus of the proper map $t$ is closed and misses the special
fibre, so properness over the DVR makes it empty. Hence $t$
is finite. The map
$\mathcal O_{(\mathcal E')^2}\to t_*\mathcal O_{\mathcal S}$
is an isomorphism modulo the uniformizer. Its coherent
cokernel vanishes by Nakayama and properness. Flatness over
the DVR then makes its kernel vanish modulo the uniformizer,
and the same argument kills that kernel. Therefore $t$ is
an isomorphism.

The second coordinate of this graph, restricted along the
zero section of the third factor, is a morphism
$g:\mathcal E'\to\mathcal E'$ reducing to $\iota$.
Subtracting $g(0)$ gives a zero-preserving morphism with the
same reduction. By rigidity it is a homomorphism of elliptic
schemes (Milne, Corollary~20.2(a)), contrary to the
nonintegrality established above. This proves the asserted
failure of a support lift, including after finite extension.

Finally one can make the ramification and the first failure
explicit. The formal group of $\mathcal E_0$ is the canonical
height-two formal $\mathbf Z_p$-module with its unramified
quadratic action. Quotienting by the cyclic direction $H$
is precisely a level-one quasi-canonical lift. Gross,
Proposition~5.3(1)--(3), gives its field of definition over
the maximal unramified base with ramification degree $p+1$,
and says that its formal modulus is a uniformizer. In
particular it is not isomorphic to the canonical lift
modulo $\varpi^2$.

If $\iota$ lifted as an endomorphism modulo $\varpi^2$,
rigidity would preserve its relation $\iota^2=-1$ and so
give the full $\mathcal O=\mathbf Z_p[\iota]$ action on
the formal group. The height-one $\mathcal O$-module has
trivial deformation space, as used in the proof of Gross,
Proposition~2.1. It would therefore be the canonical lift
modulo $\varpi^2$, a contradiction. The projection and
Nakayama argument (A31b.12) works equally over the Artin
ring $R/(\varpi^2)$: its flatness supplies the required
vanishing of the Tor term. Thus a flat surface lift to that
ring is already impossible.
\end{proof}

\paragraph{Sources actually read and exact scope.}
Gross, \emph{On canonical and quasi-canonical liftings},
Invent.~Math.~84 (1986), 321--326,
\url{https://doi.org/10.1007/BF01388810}.
The author preprint is the final six scanned sheets in
\url{https://archive.mpim-bonn.mpg.de/id/eprint/3578/1/preprint_1985_30.pdf}.
Proposition~2.1 and its proof are on preprint p.~2, and
the complete lattice construction, definition and proof of
Proposition~5.3 are on preprint pp.~8--11 (PDF sheets
76--78, using one-based sheet numbers). These statements
and proofs were read; all three sheets containing the complete
level-one construction were also visually inspected.
Milne, \emph{Abelian varieties}, in \emph{Arithmetic Geometry}
(1986), Proposition~20.1 and Corollary~20.2(a),
pp.~145--146, were read from the original page images:
\url{https://www.jmilne.org/math/articles/1986b.pdf}.
The arbitrary-scheme base in those rigidity statements is
needed for the Artin-ring assertion. The printed proof of
Proposition~20.1 is a reference to Mumford's rigidity lemma;
that reference is used as a published input, and its full
proof was not independently reread here. Corollary~20.2(a)'s
short printed proof was read. The injectivity consequence
needed above is deduced explicitly from Proposition~20.1;
Corollary~20.2(c) instead states commutativity and is not the
injectivity locator. Stacks, Tag~07S7,
Lemma~66.14.2, was read with its complete proof for the
finite-flat projective quotient construction:
\url{https://stacks.math.columbia.edu/tag/07S7}.

\paragraph{Relation to the unresolved axis question.}
The example shows that even an actual rational algebraic
class lift, free ambient Hodge cohomology and an injective
ordinary semiregularity map do not force the chosen smooth
support to lift through the first ramified range. For
$p\ge7$ the ramification $e=p+1$ gives
$\lfloor e/(p-1)\rfloor+1=2$, exactly the initial order
required by A31; the example fails at that order. At $p=3$
that required order is three, and failure already occurs
at order two. This does not contradict A31, whose initial
smooth embedded lift is an explicit hypothesis. It does
not exclude a stronger argument using the particular
primitive axis class of A14/A18, and it is not counted as
a separate staged progress-clock event.

% END work/axis/axis_ramified_support_counterexample.tex

% BEGIN work/axis/axis_uniform_abelian_obstruction.tex
% C44_POS /root. Author and independent entire-proof readings passed.
\section{A limit on uniform transfer of the whole variable cohomology}
\label{sec:c44-a32}

\paragraph{Scope.}
This is a restriction on a proposed uniform correspondence construction,
not a counterexample to the axis sign. It concerns the whole variable
local system in a full-monodromy family. It does not rule out a
correspondence defined only on an individual special axis or on a
smaller geometric locus.

\begin{proposition}[C44-A32: the Hodge Lie algebra obstruction]
Let $A/\mathbf C$ be an abelian sixfold, let $L$ be ample, and let
$B\subset A$ be a smooth intersection of two members of $|L|$.
Write
\[
 V_B=\ker\bigl(\kappa_*:H^4(B,\mathbf Q)
                       \longrightarrow H^8(A,\mathbf Q)(2)\bigr),
 \qquad c=\chi(A,L)>0.
\]
Then
\[
                         h^{4,0}(V_B)=62c-5\ge57.
                         \tag{A32.1}
\]
Suppose these axes vary in a connected smooth complex algebraic
family whose connected monodromy on $V$ is $\mathrm{SO}(V,q)$.
For a very general point, $V_B$ is not an object of the rational
Hodge category generated by the first cohomology of abelian
varieties.

Consequently, on this family, or on a dominant generically finite
base change, every morphism of rational variations of Hodge structure
from $V$ to an object generated by first cohomology of abelian
schemes is zero. The same holds for a morphism in the reverse
direction. In particular a relative algebraic correspondence cannot
inject the whole variable local system into such an abelian
construction.
\end{proposition}

\begin{proof}
The Koszul resolution and Riemann--Roch on the abelian sixfold give
\[
 \chi(B,\mathcal O_B)=\chi(A,\mathcal O_A)
           -2\chi(A,L^{-1})+\chi(A,L^{-2})=62c.
\]
For $q<4$, weak Lefschetz and Hodge decomposition give
$h^{0,q}(B)=\binom6q$. Thus
$\chi(B,\mathcal O_B)=h^{0,4}(B)-10$, and
$h^{4,0}(B)=62c+10$. The ambient image has
$h^{4,0}(A)=\binom64=15$. Its pairing is nondegenerate by
hard Lefschetz on $A$, since $[B]=c_1(L)^2$, so it is a direct
Hodge summand. Subtracting it gives (A32.1). Here $c$ is a
positive integer; for an ample bundle on an abelian variety it is
also $h^0(A,L)$.

The variable space is primitive, by C44-A01, and carries its
polarizable rational Hodge structure. At a very general point,
the connected monodromy group is contained in the derived
Mumford--Tate group. We use precisely the inclusion in Milne,
\emph{Shimura varieties and moduli}, Theorem~6.19(a), pp.~27--28;
the full proof of that inclusion was read in this run.
Therefore the Mumford--Tate Lie algebra contains
\[
        \mathfrak{so}(V_B,q)\simeq\bigwedge\nolimits^2 V_B(4)
                         \tag{A32.2}
\]
as a Hodge substructure. The identification sends
$u\wedge v$ to the skew endomorphism
$x\mapsto q(v,x)u-q(u,x)v$, with the indicated Tate twist.
In particular
\[
             \bigwedge\nolimits^2 V_B^{4,0}
                \subset\mathfrak{so}(V_B,q)^{4,-4}
                         \tag{A32.3}
\]
is nonzero by (A32.1).

If $V_B$ were generated by abelian first cohomology, a finite
product $C$ of those abelian varieties would suffice. The
Tannakian inclusion would induce a surjection
$\operatorname{MT}(H^1(C))\twoheadrightarrow\operatorname{MT}(V_B)$.
Its differential is a surjection of Hodge structures.
But $\operatorname{Lie}\operatorname{MT}(H^1(C))$ is a Hodge
substructure of $\operatorname{End}(H^1(C))$, whose only types
are $(-1,1),(0,0),(1,-1)$. A quotient cannot have the nonzero
type in (A32.3). This is exactly the necessary criterion recorded
by Brommer-Wierig--Lam, \S1.3, equation~(1.3.1); the explicit
axis Hodge number and its application here are the additional
calculation.

The standard representation of $\mathrm{SO}(V_B,q)$ is
irreducible in the present rank. Hence $V_B$ is a simple Hodge
structure at a very general point. A nonzero map from it into
an abelian-generated Hodge structure would be injective, and a
nonzero map to it would be surjective. Either would put $V_B$
in that full tensor subcategory, contrary to what was proved.
Thus both maps vanish on a very general fibre. A map of local
systems has constant rank on a connected base, so it vanishes
on every fibre. After a dominant generically finite base change,
one may remove the branch locus and obtain a finite etale cover.
The connected monodromy is unchanged on such a cover, so the
same argument applies there.
\end{proof}

\begin{lemma}[C44-A32a: nonempty axis families in this scope]
There are actual mixed-characteristic families of genuine
Lefschetz axes whose complex generic family satisfies the
monodromy hypothesis of C44-A32. They can be chosen with
ambient variety a sixth power of an elliptic curve and with
arbitrarily high product polarization.
\end{lemma}

\begin{proof}
Start with an elliptic curve over a finite field of odd
characteristic which has an explicit characteristic-zero lift,
for example $y^2=x^3-x$ in characteristic three. Apply C44-A18
with the empty collection of prescribed centres, a sufficiently
high product polarization, and a chosen first divisor $D_0$.
Its variable local system has full orthogonal geometric
monodromy. All these choices descend to a finite field.

Lift the polarized elliptic product and $D_0$ over the complete
unramified trait, and then base change to a complete strictly
henselian trait $R$ with algebraically closed residue field.
Ample cohomology vanishing makes the section module free with
base change, so the second sections give an affine parameter
space over $R$. Let $U_R$ be its nonempty smooth open on which
the intersections are smooth and the original pencils are
genuine Lefschetz pencils. The openness and nonemptiness are
the jet conditions of C44-A06; the special fibre contains the
open used in C44-A18. Both geometric fibres of $U_R$ are
connected, being nonempty opens of affine space.
The primitive summand of $R^4f_*\mathbf Q_\ell$ is a lisse
orthogonal sheaf on $U_R$: restriction, Gysin and the inverse
of the ambient lisse hard-Lefschetz isomorphism give its
cohomological projector. No relative algebraicity of that
inverse is needed for this monodromy argument.

We need only a one-way monodromy containment, not a
specialization isomorphism for a nonproper base. Here is the
finite-cover argument. A connected finite etale cover
$U'\to U_R$ is regular and integral. If its generic fibre
were not geometrically connected, the algebraic closure $K'$
of $K=\operatorname{Frac}R$ in its function field would be a
nontrivial finite extension of $K$. At the generic point of
any component of the nonempty special fibre of $U'$, smoothness
over $R$ gives a discrete valuation for which a uniformizer
of $R$ still has valuation one. Restriction to $K'$ therefore
has ramification index one. Its residue field is a finite
extension of the algebraically closed residue field of $R$,
and hence is that field. A finite extension of a complete
discrete valuation field with perfect residue field has
degree equal to ramification index times residue degree.
Thus $K'=K$, a contradiction.
Every connected finite etale cover therefore stays connected
on the geometric generic fibre. This proves the surjectivity
\[
                  \pi_1(U_{\overline K})\longrightarrow\pi_1(U_R).
                         \tag{A32.4}
\]
The special geometric monodromy is contained in the total-space
image; (A32.4) puts that image inside the geometric generic
image. Since the special closure is already the full
orthogonal group, the generic closure is the same group.

All chosen defining sections may be lifted using algebraic
coefficients. For the displayed elliptic curve, the basis of
sections with pole order at most $N$ at the zero section is
$1,x^i,yx^j$, with $i\ge1$, $j\ge0$, $2i\le N$ and
$2j+3\le N$; it is defined
over $\mathbf Z_{(3)}$. Tensor products give the product
section basis. The finitely many residue coefficients have
Teichmuller representatives in a finite cyclotomic extension.
Thus the generic family can be embedded into $\mathbf C$.
Comparison with Betti
cohomology identifies its connected monodromy as
$\mathrm{SO}(V,q)$, as required.
\end{proof}

\paragraph{Relation to the proposed transfer and Newton approaches.}
This does not disprove C44-A08's relative-form transfer criterion.
It proves that a uniform realization of the \emph{whole} variable
space by abelian tensors is unavailable on the displayed
full-monodromy families. A special-fibre algebraic correspondence
need not extend to that family, and the algebraic subspace at a
special point need not be a nonzero relative sub-local-system.
Those possibilities remain open.

There is also a published positive-characteristic obstruction of
the same kind. Brommer-Wierig--Lam, Proposition~3.1.2 and its
complete proof, state that if a lisse sheaf over a finite-field
base is generated by an abelian scheme, all slopes of its
monodromy Lie algebra lie in $[-1,1]$. For our orthogonal
weight-four variable sheaf the geometric Lie algebra is
$\bigwedge^2V(4)$. A slope pair with sum outside $[3,5]$
would therefore obstruct such a uniform abelian realization.
No such slope pair for the present axes is claimed here.
This criterion is different from a witness which makes the
algebraic axis kernel vanish, and does not settle H9-T5 or
the universal sign.

\paragraph{Sources and programme location.}
P.~Brommer-Wierig and Y.~H.~J.~Lam,
\emph{Local systems which do not come from abelian varieties},
J.~Ecole polytechnique Math.~12 (2025), 341--350,
\S1.3, equation~(1.3.1), and Proposition~3.1.2;
\url{https://arxiv.org/pdf/2408.02475}, pp.~2--3 and~5.
J.~S.~Milne, \emph{Shimura varieties and moduli}, version~1.00,
Theorem~6.19(a) and its proof, pp.~27--28;
\url{https://www.jmilne.org/math/xnotes/svh100.pdf}.
The remaining inputs are the checked C44-A01/A06/A18 and
standard smooth-proper comparison. In Grothendieck's programme
this restricts route~(b) to JAC1964 QH9-axis, cote~44
Th.~5, scan~004, with the primitive qualification of
scans~008--009. It adds no proof of the sign and is not
counted as a separate staged movement event.

% END work/axis/axis_uniform_abelian_obstruction.tex

% BEGIN work/scan_w/orthogonal_crystal_hasse.tex
% C44_POS /root/scan_w. Local descent input; no new closed-axis sign claimed.
\section{The rational quadratic form of the central crystalline slope space}
\label{sec:c44-a33}

Throughout this section $p$ is odd,
$W=W(\overline{\mathbf F}_p)$, $K=W[1/p]$, and $\sigma$ is the
absolute Witt Frobenius. For a quadratic space over $\mathbf Q_p$
we write $\epsilon_p$ for the ordinary Hasse invariant
$\prod_{i<j}(a_i,a_j)_p$ of a diagonalization.
In particular it is the invariant of the quadratic form on a
rational vector space, not the second Stiefel--Whitney class of
a Galois representation.

\begin{proposition}[C44-A33: central descent from an integral orthogonal crystal]
Let $(D,q,\Phi)$ be a finite-dimensional quadratic $F$-isocrystal
over $K$, with
\[
 q(\Phi x,\Phi y)=\sigma(q(x,y)).
\]
Assume that $D$ contains a self-dual $W$-lattice $M$.
The normalized Frobenius $\Phi$ need not preserve $M$.
Let $h_a$ denote the multiplicity of $a$ in the elementary-divisor
exponents of $\Phi M$ relative to $M$. These integers occur
symmetrically: $h_a=h_{-a}$.
Write $D_\lambda$ for the Newton slope summands and suppose
$D_0\ne0$. Set
\[
 H(M,\Phi)=\sum_{a>0}a h_a,\qquad
 S(D,\Phi)=\sum_{\lambda>0}\lambda\dim_KD_\lambda.
 \tag{A33.1}
\]
Both sums are integers. The fixed space
\[
 T_0=D_0^{\Phi=1}
 \tag{A33.2}
\]
is a nondegenerate $\mathbf Q_p$-quadratic space of dimension
$\dim_KD_0$, and its ordinary determinant has even $p$-adic
valuation. Its Hasse invariant is
\[
 \boxed{\quad
 \epsilon_p(T_0)=(-1)^{\,H(M,\Phi)-S(D,\Phi)}.
 \quad}
 \tag{A33.3}
\]
No Hodge filtration or algebraicity of $T_0$ is assumed here.
In particular, $T_0$ need not be the span of geometric cycle classes.
\end{proposition}

\begin{proof}
First, the slope decomposition and compatibility of the form
show that $D_\lambda$ pairs only with $D_{-\lambda}$.
Thus $D_0$ is nondegenerate and every nonzero slope pair is
hyperbolic. Dieudonne--Manin descent gives
$T_0\otimes_{\mathbf Q_p}K=D_0$.
Compatibility with $\Phi$ makes the descended form
$\mathbf Q_p$-valued. The integrality of $S$ follows by taking
the determinant of each isoclinic summand: the sum of its
slopes is the integer valuation of the determinant of a
matrix representing its semilinear Frobenius.

We make the orientation and integral reference form explicit.
The determinant line $\det D$ is a rank-one orthogonal
isocrystal of slope zero. It has a Frobenius-fixed generator,
whose norm lies in $\mathbf Q_p^\times$. Since $\det M$
is a unit quadratic lattice, the valuation of that norm is
even. Rescaling the generator by a power of $p$ makes its
norm a unit. Choose a self-dual quadratic lattice $M_\ast$
over $\mathbf Z_p$ of rank $\dim D$ with that determinant
squareclass. Such a lattice exists, for example by a diagonal
form with unit coefficients. Over $W$ it is isometric to $M$:
every unit of $W$ is a square, and $p\ne2$.
The determinant of an isometry can be aligned with the chosen
fixed orientation, changing its sign by an integral reflection
if necessary. In this marking
\[
                 \Phi=b\sigma,\qquad
                 b\in\operatorname{SO}(M_{\ast,K}).
 \tag{A33.4}
\]
This choice is necessary: discarding the orientation can
discard precisely the unramified quadratic descent datum.

For $g\in\operatorname{SO}(M_{\ast,K})$, let
\[
       \kappa_2(g)=v_p(\operatorname{sn}(g))\pmod2
                    \in\mathbf Z/2\mathbf Z,
 \tag{A33.5}
\]
where $\operatorname{sn}$ is the spinor norm.
Equivalently, lift $g$ to $\operatorname{GSpin}$ and take the
valuation of its spinor similitude modulo two. Different lifts
change this valuation by an even integer. The value is invariant
under $\sigma$-conjugacy, because $\sigma$ preserves valuations.
It vanishes on the integral orthogonal group.
On a hyperbolic plane the isometry
$\operatorname{diag}(t,t^{-1})$ has spinor norm $t$ modulo
squares: it is the product of the reflections in
$e+f$ and $e+tf$, after a harmless interchange of their order.
Consequently the Cartan decomposition gives
\[
                  \kappa_2(b)=H(M,\Phi)\pmod2.
 \tag{A33.6}
\]
One may also obtain (A33.6) from Rapoport--Richartz,
Theorem~4.2(i): the Kottwitz invariant is the image of the
Hodge relative position. For an orthogonal group this image
is the parity of the sum of its positive standard weights.

The self-dual reference space is unramified and its special
orthogonal group is quasi-split. Kottwitz's slope-Levi
decomposition, Proposition~6.2 of his 1985 paper, therefore
allows a $\sigma$-conjugation of $b$ into the rational Levi
associated with the Newton cocharacter. In the standard
representation this has the form
\[
 \mathop{\perp}\limits_{\lambda>0}
       (U_\lambda\oplus U_\lambda^\vee)
          \ \perp\ V_\ast,
 \qquad
 b=
 \bigoplus_{\lambda>0}
       \operatorname{diag}(A_\lambda,A_\lambda^{-{\mathsf t}})
          \ \oplus\ b_0 .
 \tag{A33.7}
\]
The quadratic form on the first summands is the hyperbolic
one; the central reference space $V_\ast$ has a self-dual
$\mathbf Z_p$-lattice. This last assertion also follows by
splitting off unit hyperbolic planes in the self-dual reference
lattice. The central isocrystal $(V_{\ast,K},b_0\sigma)$
has slope zero and is isometric to $D_0$.

The spinor norm of
$\operatorname{diag}(A,A^{-{\mathsf t}})$ is $\det A$
modulo squares. This follows from the calculation on a
hyperbolic plane and elementary generators of $\operatorname{GL}$.
Therefore
\[
 \kappa_2(b_0)=
        H(M,\Phi)-\sum_{\lambda>0}v_p(\det A_\lambda)
       =H(M,\Phi)-S(D,\Phi)\pmod2.
 \tag{A33.8}
\]
For the central basic orthogonal isocrystal,
Howard--Pappas, Proposition~4.2.5, computes the descended
quadratic form:
\[
  \det T_0=\det V_\ast
       \pmod{\mathbf Q_p^{\times2}},\qquad
  \epsilon_p(T_0)=(-1)^{\kappa_2(b_0)}.
 \tag{A33.9}
\]
Their reference lattice is self-dual, hence has Hasse
invariant $+1$, exactly as here. If the central rank is too
small for their stated rank convention, add fixed unit lines
until the rank is at least five and apply their proposition.
Adding $\langle1\rangle$ changes neither determinant squareclass
nor Hasse invariant nor spinor-valuation parity; thus this
also proves the rank-one and rank-two cases. Equations
(A33.8)--(A33.9) prove the formula and the assertion about
the determinant.
Howard--Pappas write a quadratic function $Q$ whose associated
bilinear form is $2Q$. This normalization makes no difference
here: at odd $p$, multiplication by the unit $2$ preserves the
Hasse invariant of a form whose determinant has even valuation,
by the scaling formula for Hilbert symbols.
\end{proof}

\paragraph{The semilinear datum must be retained.}
If the variety is defined over $\mathbf F_{p^f}$, replacing the
absolute semilinear operator $b\sigma$ by its linear $f$th power
multiplies $\kappa_2$ by $f$. For even $f$ this loses the parity.
Formula (A33.3) uses the absolute Frobenius and its integral
relative position, so it does not lose the descent invariant
under a finite residue-field extension.

\begin{proposition}[C44-A33a: the actual sixfold-axis lattice]
Let $B\subset A$ be a smooth intersection of two members of
one ample $|L|$, with $\dim A=6$, over
$\overline{\mathbf F}_p$. Assume
\[
                 p\ge7,\qquad p\nmid c,\qquad
                 c=\chi(A,L).
 \tag{A33a.1}
\]
Let $V$ be the actual variable middle cohomology, equipped
with $q(x,y)=\int_Bxy$, and use normalized crystalline
Frobenius $\Phi=p^{-2}F$. Its crystalline realization contains
a self-dual integral lattice to which C44-A33 applies.
If $D_\lambda$ are its normalized Newton slope spaces, then
\[
 \epsilon_p\bigl((D_0)^{\Phi=1}\bigr)
       =(-1)^{\,\sum_{\lambda>0}\lambda\dim D_\lambda}.
 \tag{A33a.2}
\]
To interpret this as the Hasse invariant of the rational
algebraic axis form, one must additionally know that the
crystalline classes of rational variable cycles span the \emph{whole}
central space $D_0$ and are nondegenerate. A sufficient
certificate is $\dim D_0$ rational cycles in the image of
the common algebraic correspondence
\[
 P=\Delta_B-\kappa^*\Gamma\kappa_*,
 \qquad \Gamma=(L_h^2)^{-1}
 \quad\hbox{on }H^4(A),
\]
with nonsingular intersection Gram matrix. Here $\Gamma$ is
the absolute Chow-motivic inverse used in C44-A15/A18b; the
same $P$ acts in every realization. Under that certificate their
rational span $T$ satisfies
$T\otimes_{\mathbf Q}\mathbf Q_p=(D_0)^{\Phi=1}$, and
(A33a.2) is the actual invariant of $q_T$.
The certificate also makes $T$ the full potential Tate part
in every realization, with its actual rational cycle form.
No such certificate follows merely from central eigenvalues
being roots of unity.
\end{proposition}

\begin{proof}
We first justify the integral hypotheses, rather than
applying a rational orthogonal projector to a lattice.
By C44-A05 choose an embedded mixed-characteristic lift
of $(A,L,B)$; its trait is allowed to be ramified.
C44-A22a gives constant Hodge numbers, degeneration of
the special Hodge--de Rham spectral sequence, and
\[
                  \dim H^i_{\mathrm{dR}}(B/k)=b_i(B)
 \tag{A33a.3}
\]
in every degree. Rational crystalline comparison gives
$\operatorname{rank}_W H^i_{\mathrm{cris}}(B/W)=b_i(B)$.
The crystalline universal coefficient sequence gives
\[
 \dim H^i_{\mathrm{dR}}(B/k)
       =b_i(B)+t_i+t_{i+1},
 \tag{A33a.4}
\]
where $t_i$ is the number of nonzero torsion invariant
factors in $H^i_{\mathrm{cris}}(B/W)$.
Thus every $t_i$ is zero. Integral crystalline Poincare
duality now makes the middle cup-product lattice self-dual.
The same facts hold for the abelian variety.

We must also split the ambient sublattice integrally.
The crystalline cohomology of $A$ is the exterior algebra
on its free $H^1$, and the prime-to-$p$ polarization condition
$p\nmid c$ makes $h=c_1(L)$ a unimodular symplectic bivector.
On $\bigwedge^4H^1$, let $L_h$ be exterior multiplication
and let $\Lambda_h$ be contraction with its integral inverse.
The usual exterior-algebra commutator identities show that
$\Lambda_h L_h$ has eigenvalues $12,6,2$ on
$L_h^2P^0,L_hP^2,P^4$, respectively.
Their differences have prime factors only $2,3,5$.
Thus the three spectral projectors are integral under
(A33a.1). On these summands
$\Lambda_h^2L_h^2$ acts by $144,36,4$, respectively.
It is consequently an integral automorphism. Hence
\[
        L_h^2:H^4_{\mathrm{cris}}(A/W)
                     \longrightarrow H^8_{\mathrm{cris}}(A/W)
 \tag{A33a.5}
\]
is an integral isomorphism.
By the projection formula the ambient pairing in $H^4(B)$
is the perfect form
$(x,y)\mapsto\int_Axyh^2$. Therefore
\[
 H^4_{\mathrm{cris}}(B/W)=
      \kappa^*H^4_{\mathrm{cris}}(A/W)\perp M_V
 \tag{A33a.6}
\]
is an integral orthogonal direct sum and $M_V$ is self-dual.
The summands are Frobenius-stable. This proves the required
integral assertion, including the denominators that would
have been hidden by a rational Lefschetz inverse.

Mazur's theorem identifies Frobenius elementary-divisor
multiplicities with the Hodge numbers when crystalline
cohomology is torsion-free and the Hodge spectral sequence
degenerates. Apply it to $B$ and $A$ and use the integral
direct sum (A33a.6). The positive normalized exponents
on $M_V$ are $1$ with multiplicity $h^{3,1}(V)$ and $2$
with multiplicity $h^{4,0}(V)$.
For completeness, Koszul and Riemann--Roch give
\[
 \chi(\mathcal O_B)=62c,\quad
 \chi(L^{-1}|B)=602c,\quad
 \chi(\Omega_B^1)=6(62c)-2(602c)=-832c.
\]
Weak Lefschetz and duality give
$h^{1,0}=6,h^{1,1}=36,h^{1,2}=90,h^{1,4}=20$,
so $h^{3,1}(B)=832c+80$ and $h^{4,0}(B)=62c+10$.
Subtracting the corresponding ambient numbers $120,15$
gives
\[
 h^{3,1}(V)=832c-40,\qquad h^{4,0}(V)=62c-5.
 \tag{A33a.7}
\]
Consequently
\[
                H(M_V,\Phi)=956c-50\equiv0\pmod2.
 \tag{A33a.8}
\]
Gauss--Bonnet and weak Lefschetz give
$\dim V=5\int_Ah^6-165=3600c-165$, which is odd.
The symmetry of the Newton slopes therefore ensures $D_0\ne0$.
C44-A33 gives (A33a.2).

Finally, crystalline cycle classes of codimension two are
fixed by $p^{-2}F$. If the indicated Gram matrix is
nonsingular, those classes are independent in crystalline
cohomology. Their number equals $\dim D_0$, so they span
$D_0$ after extension to $K$ and span its fixed space over
$\mathbf Q_p$. All pairings are the same rational
intersection numbers. This proves the asserted geometric
interpretation without a Tate-conjecture or $D(B)$ assumption.

For the last assertion, put $d=\dim D_0$. The common Chow
correspondence $P$ and the nonsingular Gram matrix give $d$
independent variable cycle classes in every realization.
They are fixed after a finite extension of the field of
definition, so the common potential Tate multiplicity
$\tau$ is at least $d$. The common characteristic-polynomial
comparison of C44-A18b, using Katz--Messing Theorem~1,
identifies this multiplicity with its crystalline value.
Every normalized root of unity has $p$-adic slope zero,
so $\tau\le\dim D_0=d$. Equality follows. The cycle basis
therefore spans the entire potential Tate space in every
realization. For any further variable cycle its projection
onto this basis has rational coefficients, obtained by
inverting the same rational Gram matrix; its orthogonal
remainder vanishes in every realization. This also justifies
the use of the actual rational cycle form in the next result.
\end{proof}

\paragraph{Independent integral check.}
The integral Lefschetz calculation above can be checked
directly. In a symplectic basis, sort the degree-four
exterior monomials by the factors in which they use just
one of the two basis vectors. There are $s=4,2,0$ such
factors. Apart from a unimodular matrix pairing those
single choices with their complements, the ambient Gram
blocks are twice the disjointness matrices on $j$-subsets
of a $(2j+2)$-element set, for $j=(4-s)/2=0,1,2$.
These matrices are $[1]$, $J_4-I_4$, and the adjacency
matrix of $KG(6,2)$. Their eigenvalues are respectively
$1$; $3,-1,-1,-1$; and $6$ once, $-3$ five times,
$1$ nine times. Thus no prime other than $2,3$ occurs
in the principal symplectic Gram determinant. This
independently verifies the perfectness used above.

\begin{proposition}[C44-A33b: the remaining global test]
Keep the actual-axis hypotheses and the full central
algebraicity certificate of C44-A33a. Let $K_{\mathrm{alg}}$
be a known positive nondegenerate rational algebraic
span and let $R=K_{\mathrm{alg}}^\perp\cap T$.
Suppose $2\le\dim R\le3$.
Write $d_K=\det q_{K_{\mathrm{alg}}}$,
$d_R=\det q_R$, and
$S=\sum_{\lambda>0}\lambda\dim D_\lambda$.
Then
\[
 \epsilon_p(q_R)=
       (-1)^S\epsilon_p(q_{K_{\mathrm{alg}}})
                         (d_K,d_R)_p,
 \tag{A33b.1}
\]
with the evident empty-span convention.
The determinant of $q_R$ is positive by C44-A17.
The form $q_R$ is positive definite if and only if
\[
 \left[
 (-1)^S\epsilon_p(q_{K_{\mathrm{alg}}})(d_K,d_R)_p
 \right]\prod_{\ell\ne p}\epsilon_\ell(q_R)=+1.
 \tag{A33b.2}
\]
The away-from-$p$ invariants in this formula are those
of the actual orthogonal Tate subspaces. They are not
asserted to be determined by the Frobenius polynomial.
\end{proposition}

\begin{proof}
The central certificate gives $T=K_{\mathrm{alg}}\perp R$
with the actual rational cycle forms, so the orthogonal
sum formula for Hasse invariants and (A33a.2) give
(A33b.1). C44-A17 applies to the fully algebraic
nondegenerate potential Tate part after subtraction of
the known positive span and gives positive determinant.
Therefore the real negative index of $R$ is even.
In rank two or three it is either zero or two. Its
real Hasse invariant is $+1$ in the first case and
$-1$ in the second. Hilbert reciprocity gives
$\epsilon_\infty(q_R)=\prod_{\ell<\infty}\epsilon_\ell(q_R)$,
which is exactly (A33b.2). Only finitely many factors
are nontrivial.
\end{proof}

\paragraph{Why the central algebraicity condition cannot be omitted.}
Even a non-Tate unit-root isocrystal becomes trivial as
an isocrystal over $W(\overline{\mathbf F}_p)[1/p]$.
Thus its $\Phi$-fixed vectors also occur in $T_0$.
If a proper algebraic subspace $T_{\mathrm{alg}}\subset T_0$
has a central nonalgebraic orthogonal complement $U$,
formula (A33.3) determines only the invariant of
$T_{\mathrm{alg}}\perp U$. The independent factor
$\epsilon_p(U)$ remains. Roots of unity in a characteristic
polynomial and a spanning family of algebraic classes are
distinct hypotheses.

\paragraph{Why Frobenius graph pairings do not supply the missing global factor.}
On a regular reciprocal summand, let $E/E^+$ be its
etale algebra with involution and write the local form as
\[
       q_a(x,y)=\operatorname{Tr}_{E/\mathbf Q_\ell}
                                      (a x\overline y),
                  \qquad a\in(E^+)^\times .
 \tag{A33.10}
\]
The adjoint of multiplication by $u\in E$ is multiplication
by $\overline u$, independently of $a$. Consequently the
trace pairing on the actual Frobenius-correspondence algebra
is $\operatorname{Tr}_{E/\mathbf Q_\ell}(u\overline v)$,
again independently of $a$. Actual algebraicity of the
graphs does not detect the hermitian norm class of $a$.
The positive graph Gram matrices used in A28/A30 therefore
do not determine the away-from-$p$ factors of (A33b.2).

\paragraph{Checks against established supersingular examples.}
For a supersingular K3 crystal the normalized Hodge exponents
are $-1,0^{20},1$ and all Newton slopes are zero.
Formula (A33.3) gives Hasse invariant $-1$, agreeing with
Howard--Pappas's basic orthogonal calculation.
For the primitive cohomology of a supersingular Fermat cubic
fourfold at an odd prime $p\equiv2\pmod3$, the exponents
are $-1,0^{20},1$ again, and the primitive lattice is
self-dual at $p\ne3$. The formula again gives $-1$.
This agrees independently with its established positive
algebraic primitive form: its complex reference has
signature $(20,2)$, the other finite local forms agree,
and Hilbert reciprocity forces the change at $p$.
For the full middle cohomology of a supersingular abelian
fourfold, the positive Hodge sum is $16+2=18$ and the
formula gives $+1$. These are consistency checks, not
new fixed-axis examples or new positivity claims.

\paragraph{Exact primary sources and audited scope.}
Howard--Pappas,
\emph{Rapoport--Zink spaces for spinor groups},
\emph{Compositio Mathematica} 153 (2017), 1050--1118,
Proposition~4.2.5 and its entire displayed proof,
author pp.~40--41 in
\url{https://arxiv.org/pdf/1509.03914},
give the central fixed-space invariant used in (A33.9);
Section~4.1.1 states the self-dual reference normalization.
The cited statement and its proof were read.
Kottwitz,
\emph{Isocrystals with additional structure},
\emph{Compositio Mathematica} 56 (1985), 201--220,
Propositions~5.4 and 6.2 with their displayed proofs,
printed pp.~216 and 219 in
\url{https://www.numdam.org/item/CM_1985__56_2_201_0.pdf},
provide the basic classification input and the slope Levi.
Rapoport--Richartz,
\emph{On the classification and specialization of
$F$-isocrystals with additional structure},
\emph{Compositio Mathematica} 103 (1996), 153--181,
Theorem~4.2(i) and its proof, printed pp.~179--180 in
\url{https://www.numdam.org/item/CM_1996__103_2_153_0.pdf},
give the invariant/relative-position identity, also
proved directly by spinor norms above.

For the geometric Frobenius-Hodge identification, Ogus,
\emph{Frobenius and the Hodge spectral sequence},
\emph{Advances in Mathematics} 162 (2001),
equation~(0.0.1) and its exact hypotheses on author p.~1
in \url{https://math.berkeley.edu/~ogus/preprints/1975a.pdf},
recall Mazur's theorem; the original cited locator is
Berthelot--Ogus, \emph{Notes on crystalline cohomology},
Theorem~8.26. The statement in Ogus was read; no claim
is made to have rechecked that book's full proof.
The crystalline derived base-change and perfectness
inputs are the standard ones recorded in Stacks,
Remarks~60.24.9 and 60.24.12--13, Tags~07MU and 07MX,
in \url{https://stacks.math.columbia.edu/tag/07MI}.
They give the universal coefficient sequence used in
(A33a.4); the hypotheses are proper smooth and the
structure crystal.
Integral crystalline duality is also recovered by
Tang, \emph{Syntomic cycle classes and prismatic
Poincare duality}, \emph{Compositio Mathematica} 160
(2024), 2322--2365, Theorem~6.4 and Remark~6.6,
\url{https://doi.org/10.1112/S0010437X24007334}.
Those passages and the displayed proof of Theorem~6.4
were read; Remark~6.6 identifies the classical source
as Berthelot, Theorem~VII.2.1.3.
The Fermat cubic sign check uses Ancona,
Proposition~A.9 and its proof, author pp.~53--54
in \url{https://arxiv.org/pdf/1806.03216v3}.

\paragraph{Programme status.}
C44-A33a computes a previously isolated local datum under
explicit actual-geometric conditions, without lifting the
central projector. It does not prove that an arbitrary
axis has a fully algebraic central slope space, does not
compute its complementary away-from-$p$ norm classes,
and does not certify a new closed axis of residual rank
two or three satisfying (A33b.2). The universal sign row
is therefore unchanged, and no progress-clock event is
claimed for a formal criterion alone.

% END work/scan_w/orthogonal_crystal_hasse.tex

% BEGIN work/axis/axis_ramified_transfer_diagnostics.tex
% C44_POS /root/axis_independent. Scoped diagnostic; no new main-gate movement.
\section{Further limits of the ramified class-transfer attempts}
\label{sec:c44-a31c}

The following calculations concern possible extensions of C44-A31b to the
existing C44-A14/A18 primitive class. They supply neither a ramified lift of
that class nor an all-ramification nonliftability theorem. In particular,
no additional auxiliary question or progress event is attached to them.

\begin{proposition}[C44-A31c: three diagnostics for rational class transfer]
The following statements hold.

\begin{enumerate}
\item The ambient tensor obstruction of C44-A14 is preserved under
isogeny. For a fixed smooth axis $\kappa:B\hookrightarrow X$ and
$w=P[S]$, where
\[
 P=1-\kappa^*(L_h^2)^{-1}\kappa_*,
\]
a different rational cycle class $z$ satisfies $Pz=nw$, for
$n\in\mathbf Q^\times$, exactly when
\[
 z-n[S]=\kappa^*(a_z-na_S),\qquad
 a_z=(L_h^2)^{-1}\kappa_*z,\quad
 a_S=(L_h^2)^{-1}[S]_X.
 \tag{A31c.1}
\]
Equality of ambient pushforwards alone does not imply this identity.

\item Let $X=M\times N$, where $M,N$ are abelian threefolds.
Let $D_1,\ldots,D_s$ be distinct smooth integral divisors of $M$,
and suppose a smooth intersection $B\in|H|^2$, with $H$ ample, contains
$Z_i=D_i\times\{0\}$ for every $i$. Write
$h=c_1(H)|_M$ and $d_i=[D_i]$. Then
\[
 \int_B[Z_i][Z_j]
   =3\delta_{ij}\int_Mh^2d_i-2\int_Mhd_id_j.
 \tag{A31c.2}
\]
Pairwise intersections of the $D_i$ are allowed to be nonreduced.
If $\sum_i x_id_i=0$ in rational cohomology of $M$, then
$v=\sum_i x_i[Z_i]_B$ belongs to $A_B$ and
\[
       \int_Bv^2=3\sum_i x_i^2\int_Mh^2d_i>0
       \qquad\text{if }(x_1,\ldots,x_s)\ne0.
 \tag{A31c.3}
\]
Thus an ambient divisor relation produces a nonzero primitive class in
this construction, not an equality of the surface classes on $B$.
For fixed supports and $H=H_0^{\otimes t}$, their primitive projections
are linearly independent for all sufficiently large $t$, whenever a
smooth such $B$ exists.

\item Fix $p\equiv3\pmod4$. There are elliptic schemes $E_n$ over finite
ramified extensions of $W(\overline{\mathbf F}_p)$ and actual rational
divisor classes $\rho_n$ on their generic squares, for every $n\ge1$,
with the following properties. Their specializations are the same
class $R(i)$ on $E^2$, where $E:y^2=x^3-x$ and $i^2=-1$.
Each $\rho_n$ is primitive for the product principal polarization and
has polarized norm
\[
                         -\int_{E_n^2}\rho_n^2=2.
 \tag{A31c.4}
\]
Its minimal integral Betti denominator at $p$ is exactly $p^n$.
Gross's corresponding level-$n$ quasi-canonical formal lift has
minimal ramification degree $p^{n-1}(p+1)$.
This is an ambient divisor example. It makes no assertion about a
primitive axis class in C44-A14/A18.
\end{enumerate}
\end{proposition}

\begin{proof}
For the first statement, begin with the fourfold construction of A14 and
write $\tau=[S]\in H^4(A_4)$. Let $\phi:A_4\to A'_4$ be an isogeny,
let $T$ be the reduced image of $S$, and let $d$ be the generic degree of
$S\to T$. Proper pushforward and isogeny pullback give
\[
 \phi_*[S]=d[T],\qquad
 \phi^*[T]=\frac{\deg\phi}{d}[S].
 \tag{A31c.5}
\]
The map on $H^1(-,\mathbf Q_5)$ is invertible. Consequently the Lie
stabilizer of the transported four-tensor is conjugate to the zero Lie
stabilizer checked in A10/A14; a nonzero scalar does not change it.
The same Hodge-circle contradiction excludes an algebraic lift of this
ambient class on every abelian-scheme model, including ramified ones.
For the sixfold version, the intrinsic tensor-support formulation in A10
is likewise invariant under an invertible change of the $H^1$ coordinates.
This does not require $T$ to be smooth or a lift of its chosen support.

Equation (A31c.1) is exactly the expansion of
$Pz=nP[S]$. In contrast, the condition
$\kappa_*z=n[S]_X$ only says that $z-n[S]$ is primitive.
If $z$ is the specialization of an actual generic cycle on a lift of
$B$, the relative Albanese construction of A29 makes its ambient
pushforward a liftable class. Since $[S]_X$ in A14 is not liftable,
the ambient correction in (A31c.1) must cancel that full ambient
obstruction. Neither the existence of the projector nor a relation of
ambient pushforwards constructs such a $z$.

We next prove the exact intersection formula. The normal sequence for
$Z_i\subset B\subset X$ and the product embedding of $M$ give
\[
 N_{Z_i/X}=\mathcal O_{D_i}(D_i)\oplus\mathcal O_{D_i}^{\oplus3},
 \qquad N_{B/X}|_{Z_i}=H|_{Z_i}^{\oplus2}.
\]
Thus
\[
             c_2(N_{Z_i/B})=3h^2-2hd_i,
\]
and the self-intersection formula gives the diagonal terms of
(A31c.2). If $i\ne j$, the scheme
$C_{ij}=Z_i\cap Z_j$ is a Cartier divisor in $Z_i$:
the distinct integral divisors $D_i,D_j$ have no common component.
Its normal line is $\mathcal O_{C_{ij}}(D_j)$. The Cartesian square
of the regular immersion $Z_j\hookrightarrow B$ and $Z_i\to B$
therefore has excess bundle of rank one, with first Chern class
\[
 c_1(N_{Z_j/B}|_{C_{ij}})-c_1(N_{C_{ij}/Z_i})
                   =(d_j-2h)-d_j=-2h.
\]
The excess-intersection formula gives
\[
                  \int_B[Z_i][Z_j]
                     =-2\int_{C_{ij}}h=-2\int_Mhd_id_j.
\]
Only the regularity of the embedding $C_{ij}\hookrightarrow Z_i$
is used here. A Cartier intersection may be nonreduced, so no
transversality or smoothness of $C_{ij}$ has been silently assumed.

If $\sum_i x_id_i=0$, its pushforward from $M$ to $X$ is zero.
Hence $v\in A_B$. Summing (A31c.2) gives
\[
 \int_Bv^2
   =3\sum_i x_i^2\int_Mh^2d_i
                      -2\int_Mh\Bigl(\sum_i x_id_i\Bigr)^2.
\]
The last term vanishes, while each $\int_Mh^2d_i$ is positive
because $H$ is ample and $D_i$ is a nonzero effective surface.
This proves (A31c.3).

For the last assertion in part (2), put
\[
 d'_i=\int_Mh_0^2d_i>0,\qquad
 C_{ij}=\int_Mh_0d_id_j,\qquad
 a_i=(L_{h_0}^2)^{-1}[Z_i]_X,
 \quad G_{ij}=\int_Xa_ia_jh_0^2.
\]
The projection formula gives the exact projected Gram matrix
\[
 \left(\int_B P[Z_i]\,P[Z_j]\right)_{ij}
       =3t^2\operatorname{diag}(d'_i)-2tC-t^{-2}G.
 \tag{A31c.6}
\]
The first term is positive definite and dominates the two fixed
matrices as $t\to\infty$. Thus this matrix is positive definite
for all sufficiently large $t$. This is a conditional calculation
for any smooth $B$ containing the indicated supports; no existence
assertion about an extra family is needed.

For part (3), use the canonical CM elliptic curve
$E_0:y^2=x^3-x$, with $\mathcal O=\mathbf Z[i]$.
Its $p$-adic CM algebra is the unramified quadratic extension of
$\mathbf Q_p$. Choose a primitive cyclic subgroup of order $p^n$
in the generic fibre, and form the quotient
$q_n:E_0\to E_n$. After a finite extension its closure is finite
flat, and the quotient is an actual elliptic scheme, exactly as in A31b.
The choice can be made so that, in the rank-one
$\mathcal O_p$ Tate lattice, the quotient lattice is
\[
       \Lambda_n=\mathcal O_p+p^{-n}\mathbf Z_p
                =p^{-n}(\mathbf Z_p+p^n\mathcal O_p).
\]
Its stabilizer order is
$\mathcal O_{p,n}=\mathbf Z_p+p^n\mathcal O_p$:
after multiplication by $p^n$ it is the order itself, whose
endomorphisms by multiplication are exactly its own elements.
The geometric generic CM endomorphism order therefore has conductor
$p^n$ at $p$. This assertion concerns the geometric endomorphism ring,
so it remains valid after every further field extension.

The special isogeny is purely inseparable of degree $p^n$ and hence
is the $n$th Frobenius followed by an isomorphism. Use that isomorphism
to identify the special fibre with $E$. Since
$F^ni=(-1)^niF^n$, the actual generic endomorphism
\[
       f_n=q_n\bigl((-1)^ni\bigr)\widehat q_n
\]
has special fibre $p^ni$. The rational endomorphism
\[
        v_n=p^{-n}f_n
              =q_n\bigl((-1)^ni\bigr)q_n^{-1}
\]
has degree one. The stabilizer-order calculation says that
$p^av_n$ is an integral geometric generic endomorphism exactly when
$a\ge n$.

For an actual endomorphism $u$ of an elliptic curve, set
\[
 D_1=\{0\}\times E_n,\quad D_2=E_n\times\{0\},\qquad
 R(u)=[\Gamma_u]-(\deg u)D_1-D_2.
\]
This is the $H^1\otimes H^1$ component of the graph class and is
additive in $u$ in cohomology. Consequently
\[
                  \rho_n=\frac{1}{p^n}R(f_n)
 \tag{A31c.7}
\]
is an actual rational divisor class specializing to $R(i)$.
It is orthogonal to both $D_1$ and $D_2$, hence primitive for
$D_1+D_2$. The graph intersection numbers are
\[
 \Gamma_u^2=D_1^2=D_2^2=0,\quad
 D_1D_2=1,\quad \Gamma_uD_1=1,
 \quad\Gamma_uD_2=\deg u.
\]
They give $R(u)^2=-2\deg u$. Since $\deg f_n=p^{2n}$,
equation (A31c.4) follows.

Finally choose a complex embedding of the generic CM data. The product
principal polarization identifies the integral cross-term lattice in
$H^2(E_n^2,\mathbf Z)$ with integral endomorphisms of $H^1(E_n,\mathbf Z)$.
The class $\rho_n$ corresponds to the rational Hodge endomorphism $v_n$.
If $p^a\rho_n$ were integral, then $p^av_n$ would preserve that lattice.
A complex-linear map preserving the lattice descends to an endomorphism
of the complex elliptic curve, contradicting the order calculation
unless $a\ge n$. Conversely $p^n\rho_n=R(f_n)$ is integral.
Thus its least denominator is exactly $p^n$, and integral multiples
obtained by clearing it have polarized norm $2p^{2n}$.

Gross, Proposition~5.3, constructs the corresponding quasi-canonical
formal lift at every level $n$ and gives its minimal ramification
$p^{n-1}(p+1)$. The quotient construction above may be carried out after
a further finite extension if needed to define the elliptic curve and
the chosen isogeny together. The denominator and norm conclusions do not
depend on choosing the minimal field.
\end{proof}

\paragraph{The direct use of the A31b graph relation.}
Take $M=E^2\times E$ and the four divisors given by
$\Gamma_i$, $\Gamma_{pi}$, $\{0\}\times E$, and
$E\times\{0\}$ in the first two factors, each times the third $E$.
The A31b ambient divisor relation has coefficients
\[
                  \left(1,-\frac1p,p-1,-\frac{p-1}{p}\right).
\]
If a smooth axis contains these four supports, part (2) makes the
corresponding primitive class $v$ nonzero. A lift of the rational
combination of the last three graph supports would therefore specialize,
after primitive projection, to $P[Z_1]-v$, not to $P[Z_1]$.
For a sufficiently high fixed power of the polarization,
(A31c.6) even makes the four projected classes independent.
Thus that construction cannot give a nonzero rational multiple of the
first primitive class by substituting the ambient graph identity.

\paragraph{The extendable liaison attempt.}
If a residual construction gives $[S]+[T]=\kappa^*a$ on $B$, then
$P[T]=-w$. However, if the full ambient class
$\gamma=\kappa_*\kappa^*a$ lifts on the chosen abelian model,
an algebraic lift of $[T]_X=\gamma-[S]_X$ would give an algebraic lift of
$[S]_X$ by subtraction. This contradicts A14.
In particular, liaison using only extendable ambient divisor data does
not supply the desired different representative. Nonliftable ambient
correction data are not ruled out by this observation, but they do not
by themselves construct a generic cycle satisfying (A31c.1).

\paragraph{Exact sources and check scope.}
The excess-intersection input is Fulton, \emph{Intersection Theory},
second edition, Theorem~6.3 and its complete displayed proof,
printed p.~102, with Corollary~6.3 on p.~103. The primary book text was
read at \url{https://djvu.online/file/Pl4avuZzLcI6I}; its hypothesis is a
Cartesian square of regular immersions, and does not require the
intersection scheme to be smooth or reduced.
Gross, \emph{On canonical and quasi-canonical liftings},
Proposition~5.3(1)--(3), with its entire proof on preprint pp.~9--11,
was reread in
\url{https://archive.mpim-bonn.mpg.de/id/eprint/3578/1/preprint_1985_30.pdf}.
The stabilizer-order construction and orbit count are explicit there.
The actual elliptic quotient and additive graph cross term use the
already checked A31b construction. No assertion of special-Hom to
generic-Hom surjectivity is used.

The existing A14/A18 primitive class remains unresolved on arbitrary
ramified models. The diagnostics above preserve the distinction between
a chosen support, an ambient class, a primitive projection, and a
different actual cycle representative.

% END work/axis/axis_ramified_transfer_diagnostics.tex

% BEGIN work/scan_w/quaternion_frobenius_sign.tex
% C44_POS /root/scan_w. Candidate A34/A34a; independent full review required.
\section{An imaginary quadratic action reversed by finite-field Frobenius}
\label{sec:c44-a34}

The two Frobenius operations in this section are different.
On crystalline cohomology over
$K=W(\overline{\mathbf F}_p)[1/p]$, normalized absolute Frobenius
$\Phi$ is semilinear and commutes with every geometric algebraic
correspondence of degree zero. The normalized geometric Frobenius endomorphism
$F$ of a model over $\mathbf F_{p^f}$ is $K$-linear. It can conjugate
a geometric endomorphism defined only over a quadratic extension.
Both operators act on the same crystalline realization, and $F$
commutes with $\Phi$. The slopes of $\Phi$ are the valuations of
the eigenvalues of $F$, divided by $f$.

We use the ordinary Hasse invariant and the signed geometric
primitive pairing throughout. An imaginary quadratic action will
be written
\[
 E=\mathbf Q(u),\qquad u^2=-d,\quad d>0,\qquad u^\dagger=-u.
 \tag{A34.1}
\]
We choose $d$ to be a positive squarefree integer.
Here $\dagger$ is the adjoint for the pairing. The endomorphism
$u$ is required to be algebraic, and the action on the Betti
realization below is required to come from an actual extension
in mixed characteristic.

\begin{theorem}[C44-A34: quaternion control of the algebraic Tate form]
Let $p$ be odd, unramified and inert in $E$, and let $f$ be odd.
Suppose a geometric cohomology block has the following comparison
data.
\begin{enumerate}
\item Its rational Betti realization $(V_{\rm B},q_{\rm B})$
is a primitive polarized Hodge substructure of weight $2r$,
with the signed geometric pairing. Its $E$-action is by Hodge
endomorphisms satisfying (A34.1). For every $\ell\ne p$ its
special-fibre realization $(V_\ell,q_\ell)$ is identified with
$(V_{\rm B},q_{\rm B})\otimes\mathbf Q_\ell$, including this action.
The rational space $V_{\rm B}\otimes\mathbf Q_p$ admits a
self-dual $\mathbf Z_p$-lattice.
\item Its crystalline realization $(D,q_D,\Phi)$ contains an
self-dual $W(\overline{\mathbf F}_p)$-lattice $M$ stable under $u$.
Stability refers to this integral generator, not to multiplication
by every element of the field $E$.
Write $H=\sum_{a>0}a h_a$ for its positive elementary-divisor
sum as in C44-A33. Its Hodge comparison identifies
$h_a=h^{r+a,r-a}(V_{\rm B})$. In particular
\[
                  \epsilon_\infty(q_{\rm B})=(-1)^H.
 \tag{A34.2}
\]
\item Normalized geometric Frobenius $F$ is an isometry on
every special-fibre realization, satisfies $Fu=-uF$, and has one common
rational characteristic polynomial, including crystalline
cohomology. All complex conjugates of its eigenvalues have
absolute value one. On $D$, both $u$ and $F$ commute with the
absolute semilinear operator $\Phi$, and the slope normalization
is the one stated above.
\item There is a nondegenerate rational space $T$ of actual
algebraic cycle classes, with its rational intersection form,
whose realization is the entire generalized potential Tate
part of $F$ in every $V_\ell$ and in $D$. Thus the eigenvalues
on this part are precisely the roots of unity, with their
full multiplicities. The same classes identify
$T\otimes\mathbf Q_p$ with a subspace of $D_0^{\Phi=1}$.
\end{enumerate}
Then the negative index of the actual rational form $q_T$ is
divisible by four. In particular, if $T$ contains a known
positive subspace of codimension at most three, then $q_T$
is positive definite.

The whole slope-zero space $D_0$ is not required to be algebraic.
No semisimplicity conjecture for Frobenius and no extension of
the projector onto $T$ are assumed.
\end{theorem}

\begin{proof}
The assertion is immediate if $T=0$, so suppose $T\ne0$.
The full potential-Tate property makes $T$ stable under $u,F$:
their images are algebraic and in the same realized subspace,
and pairing with a rational Gram basis reconstructs rational
coordinates for those images.
The potential Tate and non-Tate generalized primary spaces
are orthogonal and nondegenerate: the pairing pairs an
eigenvalue only with its reciprocal. They are stable under
$u$, because $Fu=-uF$. Denote the non-Tate part by $U$.

\emph{Step 1: common semisimple operators.}
The semisimple part $F_s$ is a rational polynomial in $F$,
chosen using the common characteristic polynomial. Indeed,
the usual lifting of the roots modulo powers of each separable
irreducible factor constructs one polynomial that gives the
Jordan semisimple part for every operator annihilated by that
characteristic polynomial. This is the ordinary additive
Jordan decomposition in characteristic zero. Its uniqueness
gives
\[
       F_s^\dagger=F_s^{-1},\qquad F_su=-uF_s.
 \tag{A34.3}
\]
For example the first equality follows by taking semisimple
parts in $F^\dagger=F^{-1}$; the second follows by conjugating
the decomposition by $u$. On crystalline cohomology this
rational polynomial still commutes with $\Phi$.

On $U$ put
\[
                    z=F_s-F_s^{-1},\qquad c=z^2.
 \tag{A34.4}
\]
Then $u,z,uz$ are skew-adjoint, $uz=-zu$, and $c$ is
self-adjoint and commutes with both $u$ and $z$.
Every eigenvalue $\lambda$ of $F_s$ in $U$ belongs to a quartet
\[
          \{\lambda,-\lambda,\lambda^{-1},-\lambda^{-1}\}.
 \tag{A34.5}
\]
The four members are distinct, since $\pm1,\pm i$ are roots
of unity and have been removed. Their common $c$-value is
$(\lambda-\lambda^{-1})^2$. It is strictly negative at every
complex embedding, by purity. Consequently the common
squarefree algebra generated by $c$ is a product
\[
                      C=\prod_i C_i
 \tag{A34.6}
\]
of totally real number fields. Its distinguished elements
$c_i$ are totally negative. All these fields and packet
multiplicities are determined by the common rational
Frobenius characteristic polynomial.

\emph{Step 2: the packet calculation, including repetitions.}
On a factor over a characteristic-zero field $k$, write
$A=C_i\otimes k$. Trace duality recovers from $q$ a unique
$A$-bilinear symmetric form $b$ with $q=\operatorname{Tr}_{A/k}b$.
Its adjoint identities for $u,z$ are the same as for $q$.
It suffices to work on a field factor $L$ of $A$.
Choose $v$ with $a=b(v,v)\ne0$. The four vectors
\[
                         v,\quad uv,\quad zv,\quad uzv
 \tag{A34.7}
\]
are pairwise orthogonal, since $u,z,uz$ are skew-adjoint.
Their respective norms are
\[
                         a,\quad da,\quad -c_i a,\quad -dc_i a.
 \tag{A34.8}
\]
Each is nonzero. The four vectors are therefore independent,
and their span is stable under $u,z$. Its orthogonal
complement is stable as well. Induction proves that $b$ is
an orthogonal sum of scalar multiples of
\[
     n_i=\langle1,d,-c_i,-dc_i\rangle,
 \tag{A34.9}
\]
the norm form of the quaternion algebra $(-d,c_i)$.
This argument works whether that algebra splits or is a
division algebra, and for arbitrary packet multiplicity.
Write $m_i$ for the number of norm-form summands. In particular
the packet rank over $C_i$ is $4m_i$.

Here is the precise trace-form invariant used below. If
$L/k$ is finite separable of degree $e$, $a,c\in L^\times$
and $d\in k^\times$, then
\[
 \begin{split}
 \det\operatorname{Tr}_{L/k}
       \langle a,da,-ac,-dac\rangle&\in k^{\times2},\\
 \epsilon_k\left(\operatorname{Tr}_{L/k}
       \langle a,da,-ac,-dac\rangle\right)
      &=(-1,-1)_k^e\,(-d,N_{L/k}c)_k.
 \end{split}
 \tag{A34.10}
\]
For the Hasse identity, put
$\theta=\operatorname{Tr}_{L/k}\langle a,-ac\rangle$.
It has rank $2e$ and determinant
$\delta=(-1)^eN_{L/k}c$ modulo squares: the two field
trace discriminants cancel, as does $N(a)^2$.
The form in (A34.10) is $\theta\perp d\theta$.
The diagonal definition of the Hilbert symbol gives
\[
 \epsilon(\theta\perp d\theta)
  =(d,\delta)_k(d,-1)_k^e(\delta,-1)_k
  =(-1,-1)_k^e(-d,Nc)_k.
\]
Its determinant is visibly a square. This proves (A34.10)
without a transfer theorem. It proves, in particular, that
the unknown scalar $a$ changes neither invariant.
The same calculation applies to products of fields and to
orthogonal sums; all cross determinant symbols are trivial.

Define the positive rational comparison space
\[
          N^+=\mathop{\perp}\limits_i
                \left(\operatorname{Tr}_{C_i/\mathbf Q}n_i\right)^{m_i}.
 \tag{A34.11}
\]
It is positive definite because $d>0$ and all real
embeddings of $c_i$ are negative. Its determinant is a square.
For every $\ell\ne p$, the actual non-Tate space $U_\ell$
has the same dimension, determinant and Hasse invariant as
$N^+\otimes\mathbf Q_\ell$. No rational isometry between a
nonexistent Betti Frobenius space and $N^+$ is asserted.

\emph{Step 3: retain the central non-Tate summands at $p$.}
Over $\mathbf Q_p$, split the factors of $C$ into those for
which $v_p(c_i)\ge0$ and those for which $v_p(c_i)<0$.
Call the corresponding parts of $N^+$ respectively $N_0^+$
and $N_{\ne0}^+$. The first factors have only slope-zero
eigenvalues: if $v_p(\lambda)\ne0$, then
\[
             v_p\big((\lambda-\lambda^{-1})^2\big)
                          =-2|v_p(\lambda)|<0.
 \tag{A34.12}
\]
Conversely valuation zero for $\lambda$ makes $c_i$ integral.

The commuting actions of $u,F_s$ on the absolute fixed space
$D_0^{\Phi=1}$ have the same adjoint identities for its actual
$\mathbf Q_p$-valued form.
The cycle space $T\otimes\mathbf Q_p$ is a nondegenerate
summand, and its orthogonal complement $U_0$ has exactly the
central non-Tate packets. Apply the calculation of Step 2
directly to this descended space. It gives
\[
             \det U_0\in\mathbf Q_p^{\times2},\qquad
             \epsilon_p(U_0)=\epsilon_p(N_0^+).
 \tag{A34.13}
\]
Thus no assertion that a slope-zero non-Tate vector is
algebraic has entered the argument.

Set $S=\sum_{\lambda>0}\lambda\dim D_\lambda$, where these
$\lambda$ are normalized slopes of the absolute $\Phi$.
For a noncentral quartet with $c$-valuation $v<0$, the two
positive slopes are both $-v/(2f)$, with the packet's
multiplicity. Hence
\[
       \sum_{i,\,v_p(c_i)<0}m_i\,
                    v_p\bigl(N_{C_{i,p}/\mathbf Q_p}c_i\bigr)
                              =-fS.
 \tag{A34.14}
\]
The notation here sums over the individual local field
factors of $C_i\otimes\mathbf Q_p$.
Since $p$ is odd and inert in $E$, $-d$ represents an
unramified nonsquare unit and
$(-d,x)_p=(-1)^{v_p(x)}$. Also $(-1,-1)_p=1$.
Equations (A34.10) and (A34.14), and the oddness of $f$, imply
\[
                       \epsilon_p(N_{\ne0}^+)=(-1)^S.
 \tag{A34.15}
\]
C44-A33 applies to the actual integral crystal $M$ and gives
$\epsilon_p(D_0^{\Phi=1})=(-1)^{H-S}$.
Using (A34.13), and then (A34.15), yields
\[
 \epsilon_p(T)
   =(-1)^{H-S}\epsilon_p(N_0^+)
   =(-1)^H\epsilon_p(N^+).
 \tag{A34.16}
\]
This is the missing rational $p$-adic comparison. It uses the
absolute semilinear descent; passing first to an even power
of finite-field Frobenius would lose the required parity.

\emph{Step 4: global reciprocity.}
For every $\ell\ne p$, the Betti comparison and the square
determinant of the non-Tate norm packets give
\[
             \epsilon_\ell(T)
                =\epsilon_\ell(q_{\rm B})\epsilon_\ell(N^+).
 \tag{A34.17}
\]
A self-dual quadratic $\mathbf Z_p$-lattice, for odd $p$,
diagonalizes with unit coefficients, so
$\epsilon_p(q_{\rm B})=1$.
Hilbert reciprocity for the rational forms $q_{\rm B}$ and
$N^+$, together with (A34.2) and positivity of $N^+$, gives
\[
            \prod_{\ell\ne p}\epsilon_\ell(T)
                         =(-1)^H\epsilon_p(N^+).
 \tag{A34.18}
\]
Multiplying by (A34.16) and applying reciprocity to the
actual rational form $T$ proves $\epsilon_\infty(T)=1$.
Finally its negative index is even: the $E$-action gives
the orthogonal complex structure $u/\sqrt d$ on $T_{\mathbf R}$,
so both positive and negative indices are even. For an even
negative index $b$, the real Hasse invariant is $(-1)^{b/2}$.
It follows that $b$ is divisible by four. A positive subspace
of codimension at most three bounds $b$ by three, forcing
$b=0$ as claimed.
\end{proof}

\begin{corollary}[C44-A34a: an actual axis criterion with two residual Tate dimensions]
Let $B$ be a smooth intersection of two members of one ample
$|L|$ in an abelian sixfold $X/\mathbf F_{p^f}$, with
$p\ge7$, $p\nmid\chi(X,L)$ and $f$ odd. Suppose that the
geometric pair $(X,L,B)$ has an automorphism $g$ of order three
preserving $L$, and that
\[
          FgF^{-1}=g^{-1},\qquad p\equiv2\pmod3.
 \tag{A34a.1}
\]
Assume the actual pair and this automorphism admit an embedded
mixed-characteristic lift, with the given polarization.
On the variable middle cohomology put
\[
               e=1-\frac{1+g+g^2}{3},\qquad u=g-g^2.
 \tag{A34a.2}
\]
Let $K$ be a known positive, nondegenerate rational span of
actual variable cycle classes in $eV$, stable under $u$ and $F$.
Allow $K=0$. If the generalized potential Tate multiplicity
of $F$ on $eV\cap K^\perp$ is at most two, then the actual
algebraic intersection form on $eV$ is positive definite.

If, in addition, the complementary invariant block
$(1-e)V$ satisfies the one-direction criterion C44-A19
after removing a known positive span, the whole fixed-axis
form on $A_B$ is positive definite. Thus the two-dimensional
imaginary-quadratic residual block is allowed in addition
to one unresolved rational Tate direction in the other block.
No algebraicity of those possible residual directions is
part of the conclusion.
\end{corollary}

\begin{proof}
We first verify the integral and comparison conditions, and
then justify the cycle-spanning step; the latter is necessary
because no Tate conjecture is assumed.
The automorphism preserves the geometric pairing and the
variable space. The idempotent $e$ is self-adjoint; on its
image $u^2=-3$ and $u^\dagger=-u$. The assumed lift makes this
an actual imaginary-quadratic Hodge action on the Betti block.
C44-A33a supplies a self-dual variable crystalline lattice.
Because $3$ is a unit and $g$ is integral, its orthogonal
decomposition by $e,1-e$ is integral, and $eM$ is self-dual.
The lattice $eM$ is stable under $u$, and its rational span
is stable under absolute Frobenius. The normalized $\Phi$
need not preserve the lattice itself.

The Betti $p$-lattice is self-dual for the same concrete
reason. Integral weak Lefschetz identifies $H_3(B,\mathbf Z)$
with the free group $H_3(X,\mathbf Z)$; the universal coefficient
sequence and Poincare duality make the free middle lattice
unimodular. The integral exterior-algebra proof of C44-A33a
applies to $H^1(X,\mathbf Z_p)$: since $p\nmid\chi(L)$ and
$p\ge7$, $L_h^2:H^4(X,\mathbf Z_p)\to H^8(X,\mathbf Z_p)$
is an isomorphism with perfect induced pairing. Its orthogonal
variable complement, and then its $e$-summand, are self-dual.

The Hodge modules in the lift are free and constant by
C44-A22a. Taking the direct summand of a tame idempotent
preserves that fact and the degeneration of the Hodge
spectral sequence. Mazur's elementary-divisor theorem,
in the precise scope used in A33a, identifies the Hodge
numbers of this summand with those of the integral crystal.
Hodge--Riemann gives
\[
  b_-(eV_{\rm B})=2\sum_{j>0,\,j\text{ odd}}h^{2+j,2-j}(eV_{\rm B}),
\]
so (A34.2) follows with this block's own $H$. It is not
necessary for that $H$ to have the parity of the whole variable
space. Normalized Frobenius is an isometry, (A34a.1) reverses
$u$, and $\mathbf Q(\sqrt{-3})$ is unramified inert at $p$.
Absolute crystalline Frobenius commutes with $u$ by naturality.

The characteristic polynomials of this block are common
rational polynomials in all the stated realizations. Here one
can use actual algebraic Kunneth projectors, rather than assume
that individual $g$ commute with finite-field Frobenius.
Katz--Messing, Theorem~2(1), supplies a common rational
polynomial in a finite-field Frobenius that projects onto each
cohomological degree. The traces of an algebraic correspondence
composed with that projector are rational intersection numbers
and independent of the cohomology theory. Apply this to
$ePF^n$, for every $n$, where
\[
 P=\Delta_B-\kappa^*\Gamma\kappa_*,\qquad
                  \Gamma=(L_h^2)^{-1}
 \tag{A34a.3}
\]
is the common Chow-motivic inverse used in C44-A15/A18b.
Newton identities give the common polynomial. The crystalline
statement uses the same trace formula and cycle correspondences;
Katz--Messing Theorems~1 and~2(1) provide their common projectors.
Purity follows because this is a subspace of middle cohomology.
Removing $K$ by its rational Gram projector preserves all
these comparisons and the $u,F$ actions.

Let $R$ be the potential Tate part in this orthogonal complement.
It has even dimension: over an algebraically closed coefficient
field, $u$ has two opposite eigenspaces paired perfectly by the
form, so their dimensions are equal. Hence $\dim R=0$ or $2$.
If its algebraic cycle image is zero, rational Gram projection
identifies the entire algebraic block with the known $K$, and
there is nothing further to prove.

Otherwise fix one $\ell\ne p$ at which that image is nonzero.
When $\dim R=2$, its two $u$-eigenspaces are lines. The
semisimple part $F_s$ exchanges them. In a basis in which their
pairing is $\bigl(\begin{smallmatrix}0&1\\1&0\end{smallmatrix}\bigr)$,
it has matrix
\[
              \begin{pmatrix}0&a\\a^{-1}&0\end{pmatrix}.
 \tag{A34a.4}
\]
In particular $F_s^2=1$. The operators $u,F_s$ generate the
split quaternion algebra $(-3,1)=M_2(\mathbf Q_\ell)$ on $R$,
so no nonzero proper subspace is stable under both. The algebraic
cycle image is stable under both operators: they are rational
polynomials in actual algebraic correspondences. It therefore
spans $R$. Choose a rational linear combination $v$ of these
actual cycles with $a_0=q(v,v)\ne0$; one exists because the
form on $R$ is nondegenerate and the coefficient field has
characteristic zero. Then $v,uv$ have rational Gram matrix
\[
                         \operatorname{diag}(a_0,3a_0).
 \tag{A34a.5}
\]

Apply the common correspondence $eP$ and the rational Gram
projection off $K$ to all representatives in this argument.
Thus their membership in the indicated subspaces holds in
every realization. The nonzero rational Gram determinant in
(A34a.5) proves independence at every $\ell$ and in crystalline
cohomology. Dimension comparison with the common Frobenius
polynomial then shows that these cycles together with $K$
span the entire potential Tate space everywhere. Their
crystalline classes are fixed by normalized absolute Frobenius,
as all geometric algebraic cycle classes are. They consequently
supply exactly the space $T$ required in C44-A34, even if the
central slope space also has non-Tate summands.

C44-A34 makes the negative index of this actual $T$ divisible
by four. Its positive subspace $K$ has codimension two, so
the index is zero. For any further algebraic class, its
coefficients in this basis are rational: pair with the basis
and invert its rational nonsingular Gram matrix. This proves
positivity on the whole rational algebraic block, without
assuming homological and numerical equivalence coincide.
Finally the $e,1-e$ summands are geometrically orthogonal;
C44-A19 supplies the asserted complementary sign.
\end{proof}

\paragraph{Source scope and present status.}
The quaternion decomposition and the trace-Hasse formula above
are proved directly, including repetitions and split local
quaternion algebras. The quadratic reciprocity input is the
usual product formula for rational Hilbert symbols, in the
same ordinary-Hasse convention as C44-A17/A33. The crystalline
descent input is exactly C44-A33, with Howard--Pappas,
Proposition~4.2.5, and Kottwitz's rational Newton-Levi
decomposition in its checked scope. The elementary-divisor
identification has the torsion-freeness and Hodge degeneration
hypotheses explicitly verified above; it is not inferred from
weak admissibility alone.
For the additional primary comparison locator, Katz--Messing,
\emph{Inventiones mathematicae}~23 (1974), Theorem~1 and its
proof are on pp.~74--75; Theorem~2(1) and its full polynomial
projector proof are on pp.~76--77. These actual pages were
read for this argument. Their Theorem~1 has no subpart~(i).
No arbitrary numerical-motive polarization transfer is used.

The corollary is an actual geometric sufficient criterion.
An eligible closed axis with the stated two-dimensional
residual multiplicity must still be exhibited before this
section is counted as a newly instantiated fixed-axis case.
In particular the possible two residual Tate classes have
not been asserted to be algebraic, and no standard conjecture
for every fixed axis follows from this section alone.

% END work/scan_w/quaternion_frobenius_sign.tex

% BEGIN work/axis/c3_axis_sign.tex
% C44_POS /root/axis_independent; GPT-6 Astra Pro; effort not exposed.
% Complete candidate C44-A35. Independent whole-proof audit pending.
\section{Actual cyclic axes with an imaginary-quadratic residual Tate line}
\label{sec:c44-a35}

This construction instantiates C44-A34a. The two possible residual
algebraic directions are kept together in one imaginary-quadratic
block. No relative projector onto either of those individual
directions is assumed.

\begin{theorem}[C44-A35: genuine Lefschetz axes with two residual directions]
Let $p\ge11$ be a prime with $p\equiv2\pmod3$. There are smooth
axes $B\subset X=E^6$ over finite fields of odd degree over
$\mathbf F_p$, where $E:y^2=x^3+1$, with the following properties.
\begin{enumerate}
\item The axis is the smooth intersection of two members of the
same ample system, and their original pencil on $X$ is a genuine
Lefschetz pencil: every singular member has exactly one ordinary
double point.
\item The variable algebraic space contains an explicit positive
Frobenius-fixed line $K=\mathbf Qw$. On its orthogonal complement
the total generalized root-of-unity multiplicity of normalized
Frobenius is exactly two. Both dimensions belong to the same
$\mathbf Q(\sqrt{-3})$ block, on which finite-field Frobenius
conjugates the quadratic field.
\item The entire rational algebraic axis space $A_B$ is positive
definite. Its dimension is either one or three; algebraicity of
the two additional potential Tate directions is not asserted.
\end{enumerate}
More precisely these conclusions hold for a Zariski-dense set
of closed parameters of odd degree in the nonempty family
constructed below, for every sufficiently large integer $t$
prime to $p$. They use neither a Tate conjecture nor 032.
\end{theorem}

\begin{proof}
Write $\zeta$ for a primitive cube root in $\mathbf F_{p^2}$,
and put $\chi(g)=\zeta$. On $X=E^6$ define
\[
 g=(1,\zeta,\zeta,\zeta^{-1},\zeta^{-1},\zeta^{-1}).
 \tag{A35.1}
\]
Here $\zeta$ denotes $(x,y)\mapsto(\zeta x,y)$ on the elliptic
curve. Thus $g^3=1$ and arithmetic Frobenius over $\mathbf F_p$
conjugates $g$ to $g^{-1}$. The tangent characters, up to
interchanging the names $\chi,\bar\chi$, are
\[
                         1\oplus2\chi\oplus3\bar\chi.
 \tag{A35.2}
\]
The coordinate formula defines an action of the diagonalizable
group scheme $\mu_3$ over $\mathbf F_p$. Its two nonzero
weights on sections are defined over $\mathbf F_p$ itself.
Explicitly, if $g^*s=\zeta s$, conjugation of both $g$ and
$\zeta$ shows $g^*\sigma(s)=\zeta\sigma(s)$; thus each
section weight space descends separately. This distinction
will matter below: finite-field Frobenius on cohomology,
whose coefficient field is fixed, interchanges the two
geometric $g$-eigenspaces.

Let $\eta_i$ be the pullback of the point class on the $i$th
elliptic factor. Start with
$M=\boxtimes_{i=1}^6\mathcal O_E(3[0])$, and give
\[
 H=M\otimes g^*M\otimes(g^2)^*M,
 \qquad L=H^t,\qquad a=9t,
 \qquad h=c_1(L)=a\sum_{i=1}^6\eta_i
 \tag{A35.3}
\]
the norm linearization. This linearization has trivial fibre
character at every fixed point: the action cyclically
permutes three identical one-dimensional factors. The bundle
and its linearization have their natural quadratic descent.
The underlying bundle is the product of
$\mathcal O_E(9t[0])$, and
\[
                 c=\chi(X,L)=a^6=3^{12}t^6,
                 \qquad p\nmid c.
 \tag{A35.4}
\]

Prescribe the abelian surface
\[
              S=E_1\times\{0\}_2\times\{0\}_3
                    \times\{0\}_4\times\{0\}_5\times E_6.
 \tag{A35.5}
\]
The choice of a $\bar\chi$ factor in (A35.5) is material.
Over $\mathbf F_{p^2}$ take section pairs
\[
 \mathcal P=H^0(X,I_SL)_\chi\oplus H^0(X,I_SL)_{\bar\chi}.
 \tag{A35.6}
\]
Each summand is preserved by coefficient descent, as just
checked; their direct sum is the base change of an affine
space over $\mathbf F_p$, denoted by the same symbol. A pair
$(s_\chi,s_{\bar\chi})$ defines its common zero locus and its
two-dimensional pencil. At a rational point of the descended
parameter space the two sections, their eigenmembers, the
pencil, and its base are all rational over that field.

\paragraph{Uniform jets and the fixed curves.}
We take $t$ sufficiently large once and for all. The following
jet generation can be imposed simultaneously. Sections of
$I_SL$ separate jets of order at most two at any bounded
number of points off $S$, and their normal first jets along
$S$ have the corresponding fibrewise generation and bounded
pointwise independence. One direct uniform justification is
relative Serre vanishing on the proper schemes of bounded
length subschemes and on the normal-jet targets: after
twisting by a fixed power, the relevant evaluation maps are
surjective in every fibre. Equivalently, first make the
coherent ideal and normal-jet sheaves globally generated and
multiply by products of hyperplanes separating the bounded
sets of points. In particular the bound can cover twelve
geometric points, second jets at a free orbit, two distinct
free orbits, and four quotient-point evaluations used below.
At a fixed point take the corresponding equivariant jet
target. Averaging over $C_3$, with $3$ invertible, preserves
these surjections onto each character component.

The fixed locus of $g$ is the disjoint union
\[
       X^g=E_1\times\ker(\zeta-1)^5
                         =\coprod_{243}E_1.
 \tag{A35.7}
\]
Every such curve lies in the base of every pair, because the
fibre of $L$ is trivial as a representation while the sections
have nontrivial characters. On a fixed curve disjoint from
$S$, the available first derivatives of the two sections lie
in the two normal character spaces of ranks $2$ and $3$.
On each of the three fixed curves contained in $S$, requiring
vanishing on $S$ leaves normal ranks $2$ and $2$.
For either section the locus where its allowed derivative
vanishes has codimension at least two in the parameter
space, while a fixed curve has dimension one. A general pair
therefore has both derivatives nonzero at every fixed point.
They occupy different character spaces, so are independent.
This proves smoothness of the common base along all fixed
curves and smoothness there of both eigenmembers.

Away from the fixed curves the action is free. On $S$ the
normal derivatives give an arbitrary $2\times4$ matrix;
rank at most one has codimension three, larger than
$\dim S=2$. Thus a general base is smooth along all of $S$.
Off $S\cup X^g$, the failure of smoothness of a common zero
requires two zero values and a rank drop of a $2\times6$
derivative matrix. These have total codimension
$2+(6-1)=7>6$. The same equivariant jet calculation at a
free orbit therefore excludes every other singular point
for a general pair. Each eigenmember is also smooth away
from the prescribed loci: vanishing of a value and all six
derivatives again gives seven conditions in dimension six.
All these assertions are open conditions on the parameter
space and hold simultaneously on a nonempty open.

\paragraph{The original pencil has one node per singular member.}
For clarity this is the pencil on $X$, not a pencil in the
parameter space of axes. Its two eigenmembers are smooth by
the preceding paragraph. At every point of the base the two
gradients are independent, so every pencil member is smooth
there, including along $S$ and $X^g$.

For a nonexceptional member and a point outside the base,
the conditions that its value and its six derivatives vanish
impose seven independent jet conditions. The incidence
space has the six point variables and one member variable;
therefore its general parameter fibre is finite. A
degenerate quadratic part imposes one further condition,
which can be avoided by the second-jet surjectivity. Since
$p\ne2$, the remaining critical points are ordinary double
points. Two critical points in distinct free $C_3$-orbits
for the same member impose fourteen conditions on the
thirteen variables consisting of the two points and the
member. They can therefore be avoided. If two points in the
same orbit were critical for the same nonexceptional member,
then at one point two distinct pencil members would have
zero first jet. Their two generators would both have zero
first jet, contradicting the already established smooth
base. Nonexceptional parameters have free $C_3$-orbits;
this is why an orbit of points causes no repeated critical
point in a single original member.

There is at least one critical point. Indeed the first-jet
bundle has rank seven and
$c(J^1L)=(1+h)^7$, since $\Omega_X^1$ is trivial. The rank
drop of the two pencil sections in $J^1L$ is, for the just
constructed transverse pencil, a reduced zero-dimensional
scheme of degree
\[
                            \int_Xc_6(J^1L)=7\int_Xh^6>0.
 \tag{A35.8}
\]
There is no contribution from the smooth base, where the
two first jets are independent. This also verifies the
usual nonempty singular-fibre condition. Denote by
$\Omega\subset\mathcal P$ the nonempty descended open of
pairs with these smooth-base and genuine-pencil properties.

\paragraph{An explicit positive invariant cycle and exact ranks.}
Let $Q_0=\{1,6\}$ and $R_0=\{2,3,4,5\}$. In the squarefree
algebra generated by the $\eta_i$ put
\[
 \alpha=\frac1{12a^2}\left(
    3\eta_1\eta_6
    -\sum_{i\in Q_0,\ j\in R_0}\eta_i\eta_j
    +\sum_{i<j,\ i,j\in R_0}\eta_i\eta_j\right).
 \tag{A35.9}
\]
Multiplication using $\eta_i^2=0$ gives
$h^2\alpha=[S]_X$. For any smooth base in this family,
\[
       w=[S]_B-\kappa^*\alpha\in V_B,
       \qquad q(w,w)=6a^2-\frac1{4a^2}>0.
 \tag{A35.10}
\]
For completeness the first term is the self-intersection
$\int_Sc_2(N_{S/B})=3\int_Sh^2=6a^2$:
$N_{S/X}$ is trivial and its quotient is $(L|S)^{\oplus2}$.
The correction is $\int_X\alpha[S]_X=1/(4a^2)$.
The line $K=\mathbf Qw$ is constant, $g$-invariant and
Frobenius-fixed after twist $(2)$. All its representatives
and its rational Gram projector are actual algebraic
correspondences.

Write $V$ for the variable middle sheaf, with pairing
$q(x,y)=\int_Bxy$ and twist $(2)$, and put
$W=V\cap K^\perp$. Over split coefficients and the quadratic
constant-field extension,
\[
              W=W_0\oplus W_\chi\oplus W_{\bar\chi}.
 \tag{A35.11}
\]
Here $W_0$ is orthogonal and $W_{\bar\chi}=W_\chi^\vee$ by
the geometric pairing; the other character pairings vanish.
The latter two sheaves together form the rational
$E_0=\mathbf Q(\sqrt{-3})$ block. We use $E_0$ for this
coefficient field to distinguish it from the elliptic curve.

The exact rank calculation is as follows. On $H^1(X)$ the
character multiplicities are $(2,5,5)$ for
$(1,\chi,\bar\chi)$. Taking exterior powers gives
\[
\begin{array}{c|rrr}
 j&1&\chi&\bar\chi\\\hline
 0&1&0&0\\
 1&2&5&5\\
 2&26&20&20\\
 3&70&75&75\\
 4&165&165&165
\end{array}
 \tag{A35.12}
\]
These are the coefficients of
$(1+z)^2(1+\chi z)^5(1+\bar\chi z)^5$ in
$\mathbf Z[\chi]/(\chi^3-1)$.
The tangent-normal sequence for $B$ gives
$\chi(B)=5\int_Xh^6=3600c$. The tame fixed-point Euler
formula gives $\chi(B;g)=\chi(B^g)=0$, because
$B^g=X^g$ is the union (A35.7) of elliptic curves. Hence
the three Euler-character multiplicities are all $1200c$.
Weak Lefschetz and its dual identify all degrees other
than the middle with the ambient exterior powers. Their
lower alternating multiplicities are $(-45,-60,-60)$.
Subtracting twice these numbers and then the ambient
degree-four multiplicities gives
\[
 \operatorname{rank}V_0=1200c-75,\qquad
 \operatorname{rank}V_\chi=
 \operatorname{rank}V_{\bar\chi}=1200c-45.
 \tag{A35.13}
\]
Consequently
\[
           m=\operatorname{rank}W_0=1200c-76\ \text{is even},
 \qquad n=\operatorname{rank}W_\chi=1200c-45\ \text{is odd}.
 \tag{A35.14}
\]
Both ranks exceed eight.

\paragraph{The free quotient and its character moments.}
For this calculation extend constants to $\mathbf F_{p^2}$,
so that $g$ and its characters are defined. Remove
$S\cup X^g$ and put
\[
      U_X=X\setminus(S\cup X^g),\qquad Y=U_X/C_3.
 \tag{A35.15}
\]
The scheme $Y$ is smooth and geometrically connected of
dimension six, and $U_X\to Y$ is a connected finite etale
$C_3$-torsor. The two-character section bundle descends to
a rank-two bundle on $Y$. For a finite extension with
$Q$ elements and any four distinct rational points of $Y$,
evaluation of $\mathcal P(\mathbf F_Q)$ onto their four
two-dimensional fibres is surjective. Indeed their inverse
images consist of at most twelve geometric points; use the
uniform evaluation above and then take the prescribed
character components. Surjectivity descends from the
algebraic closure. This also handles nonsplit torsor fibres.

For a uniformly random parameter in the entire affine
space $\mathcal P(\mathbf F_Q)$, let $I_y$ be the indicator
that the descended section vanishes at $y$. Then the
$I_y$ are four-wise independent Bernoulli variables of
parameter $v=Q^{-2}$. Put $X_y=I_y-v$ and
\[
       T_0=\sum_{y\in Y(\mathbf F_Q)}X_y,
       \qquad T_\chi=\sum_{y\in Y(\mathbf F_Q)}w_yX_y,
       \qquad w_y=\chi(h_y),
 \tag{A35.16}
\]
where $h_y$ is the torsor Frobenius. Replacing $\chi$ by
its inverse merely conjugates the second expression; we
choose the convention matching the trace on $W_\chi$.
Write $M_Q=\#Y(\mathbf F_Q)$ and $A_Q=\sum_yw_y^2$.
The compact-support trace formula and weights for the
connected torsor imply
\[
        M_Q=Q^6+O(Q^{11/2}),\qquad A_Q=O(Q^{11/2}).
 \tag{A35.17}
\]
The second assertion follows because $\chi^2$ is nontrivial
on the geometric torsor, so the degree-twelve compact-support
invariants vanish. All remaining weights are at most
eleven. Also $|w_y|=1$ in the chosen complex embedding.

Let $\sigma^2=v(1-v)$ and
$\mu_4=v(1-v)^4+(1-v)v^4$. Expanding the two-pair and
four-equal index patterns gives the exact identities
\[
\begin{split}
 \mathbf E(T_0^4)
    &=3M_Q(M_Q-1)\sigma^4+M_Q\mu_4,\\
 \mathbf E(|T_\chi|^4)
    &=2M_Q^2\sigma^4+|A_Q|^2\sigma^4
                         +M_Q(\mu_4-3\sigma^4).
\end{split}
 \tag{A35.18}
\]
Every pattern with a singleton index has zero expectation.
In the second formula two of the pairings contribute
$M_Q^2-M_Q$, while the remaining one contributes
$|A_Q|^2-M_Q$. Thus, without assuming independence of more
than four points,
\[
        Q^{-8}\mathbf E(T_0^4)\longrightarrow3,
        \qquad Q^{-8}\mathbf E(|T_\chi|^4)\longrightarrow2.
 \tag{A35.19}
\]

We next identify the actual trace functions, including
the prescribed cycle. Temporarily omit the Tate twist and
write $W_\rho^{\rm un}$ for weight four. On the smooth open,
uniformly in the parameter and in $Q$,
\[
      T_\rho=\operatorname{Tr}(F_Q\mid W_\rho^{\rm un})
                                      +O(Q^{3/2}),
                   \qquad \rho=0,\chi.
 \tag{A35.20}
\]
Here is a direct verification. For a $C_3$-stable scheme
$T$ let $N_j(T)=\#\operatorname{Fix}(g^jF_Q;T)$. The
character average $\frac13\sum_j\chi^{-j}N_j(T)$ on a
free torsor is its weighted point count on the quotient.
Apply this to the zero locus and subtract $Q^{-2}$ times
the count on $U_X$. For the full proper varieties, weak
Lefschetz, its dual Gysin map and hard Lefschetz on $X$
give
\[
       N_j(B)-Q^{-2}N_j(X)
        =\operatorname{Tr}(g^jF_Q\mid V^{\rm un})
                                                +O(Q^{3/2}).
 \tag{A35.21}
\]
In fact $H^i(B)$ for $i>4$ cancels with
$H^{i+4}(X)(2)$, and the ambient part in degree four
cancels with $H^8(X)(2)$ by $L_h^2$. All the remaining
terms have weight at most three after the displayed
normalization. Finite-field twists by $g^j$ justify the
same trace formula for each $j$.

The removed fixed elliptic curves contribute $O(Q)$ to
$B$ and $O(Q^{-1})$ after the $Q^{-2}$ ambient factor;
these are within the error. For the removed surface $S$,
its invariant top cohomology contributes exactly
$Q^2+O(Q^{3/2})$ to the invariant average of
$(1-Q^{-2})N_j(S)$. This removes the line $K$. Its
nontrivial character average has no degree-four term,
because $g$ acts trivially on the top cohomology of the
geometrically connected surface $S$; it is
$O(Q^{3/2})$. Intersections of $S$ with the fixed curves
are also of order $O(Q)$. This proves (A35.20).

Purity makes geometric monodromy reductive. On the smooth
open both normalized traces and the corresponding
$T_\rho/Q^2$ are uniformly bounded, so (A35.20) changes
their fourth absolute moments by $o(1)$. Katz's
Lemmas 2.1.4--2.1.6 identify the limiting upper fourth
moment with the dimension of the geometric invariant
tensors. Omitting the other parameters can only decrease
the sums in (A35.18), since their summands are nonnegative.
The density of the smooth open tends to one. Consequently
\[
            M_4(G_0,W_0)\le3,
            \qquad M_4(G_\chi,W_\chi)\le2.
 \tag{A35.22}
\]
There is no estimate here for the cohomology of a singular
fibre. The three orthogonal contraction tensors in the
first space and the two general linear contraction tensors
in the second are independent, since their ranks exceed
eight. They give the opposite inequalities. Thus both
inequalities in (A35.22) are equalities.

Here the use of a limiting upper moment does not suppose
that all geometric invariant tensors are arithmetically
fixed. More explicitly, for $Q=p^{2j}$ the trace formula
on the smooth parameter open has leading term the trace
of the $j$th power of arithmetic Frobenius on the
geometric invariant tensors of
$W_\rho^{\otimes2}\otimes(W_\rho^\vee)^{\otimes2}$.
Its eigenvalues have complex absolute value one by
purity. Along a subsequence their finitely many phases
simultaneously tend to one, so the limiting upper
average is the dimension of that invariant space.
All lower compact-support terms tend to zero. This is
the limsup statement of the cited Katz lemmas and
justifies the inequalities even when those phases
are nontrivial.

\paragraph{Reflections in the variable character sheaves.}
At one free point outside $S\cup X^g$, use the independent
equivariant second jets to impose a rank-one derivative
matrix for the two equations and a nondegenerate quadratic
part after eliminating one equation. The base then has an
ordinary double point of dimension four. Equivariance
imposes the same singularity at its free orbit of three
points. The remaining jet and dimension arguments above,
now with this finite orbit prescribed, make the base smooth
everywhere else and along $S\cup X^g$. Vary the transverse
constant term at the chosen point; the resulting local
one-parameter smoothing is regular at all three nodes.
Its generic parameter can be chosen in the dense open
$\Omega$. This degeneration of axes is allowed to have
three nodes; it is distinct from the original pencil on
$X$, whose one-node condition has already been proved.

SGA 7 II, Expose XV, Theorem 3.4(i),(iii), gives three
mutually orthogonal vanishing cycles of square $\pm2$,
permuted regularly by $C_3$. They are nonzero and independent
because these squares are nonzero. One tame inertia element
acts as the product of their three reflections: all three
smoothing parameters are nonzero unit multiples of the
same local parameter. It acts as $-1$ on their span and
as $1$ on the orthogonal complement. The cycles belong to
$W$: ambient classes and the globally prescribed cycle
$w$ are fixed, and the Picard--Lefschetz formula forces
their pairings with all three independent node cycles to
vanish. The regular three-dimensional representation has
one line in each character. Hence this inertia element is
a reflection on $W_0$ and an order-two pseudoreflection
on $W_\chi$ (and on $W_{\bar\chi}$).

Katz's Theorem 1.4.6 now gives
$G_0=O(W_0)$; his Theorem 1.4.2(5) gives
\[
                         G_\chi\supset SL(W_\chi).
 \tag{A35.23}
\]
The hypotheses in the latter theorem are exactly
reductivity, fourth moment two, a reflection and rank
greater than eight. In particular no unproved determinant
finiteness or additional irreducibility hypothesis is
needed. Passing to the open with the original-pencil
condition does not change these geometric images:
restriction to a dense open of a normal connected base
surjects on its etale fundamental group.

\paragraph{A nonempty good open in an odd arithmetic coset.}
Return to the descended family over $\mathbf F_p$, and use
normalized weight-zero Frobenius. Let $\mathcal H$ be the
joint geometric monodromy closure on $W_0$ and the
imaginary-quadratic block. Over an algebraically closed
coefficient field its identity component projects onto
$SO(W_0)$ and its projection to $GL(W_\chi)$ contains
$SL(W_\chi)$, by (A35.23). We do not require these
projections to be independent.

Choose an element of the actual arithmetic image whose
constant-field degree is odd. It interchanges the two
nontrivial character spaces. In bases making their pairing
\(\bigl(\begin{smallmatrix}0&I\\I&0\end{smallmatrix}\bigr)\),
every such isometry has the form
\[
       F_E=\begin{pmatrix}0&A\\A^{-t}&0\end{pmatrix},
       \qquad F_E^2|W_\chi=AA^{-t}.
 \tag{A35.24}
\]
Multiplying the chosen element by the geometric inertia
reflection above, if necessary, makes its determinant on
$W_0$ equal to $+1$. This multiplication preserves its odd
constant-field degree and the anti-quadratic form (A35.24).
Denote the resulting element by $F_*$. Its connected
geometric coset $F_*\mathcal H^0$ is irreducible.

For any prescribed finite set of roots of unity there
is a nonempty open in this coset where $F|W_0$ avoids
that set. Indeed $m$ is even,
the determinant is $+1$, and the projection of this coset
contains a full translate of $SO_m$. On an even-dimensional
special orthogonal group choose a torus element with
generic reciprocal eigenvalue pairs; this avoids every
prescribed finite set of roots of unity.

On the imaginary-quadratic part there are necessarily at
least two root-of-unity eigenvalues. The matrix
$C=AA^{-t}$ is similar to $C^{-1}$ (for example
$A^{-1}CA=A^{-t}A$, which is the transpose of $C^{-1}$),
and has determinant one. Since $n$ is odd, the reciprocal
pairing of its eigenvalues forces an eigenvalue $1$:
the multiplicity of $-1$ is even by the determinant.
The characteristic polynomial of $F_E$ is
$\det(T^2I-C)$, so it has both $+1$ and $-1$.

This is the only forced root-of-unity pair; for a finite
prescribed set the additional exclusions form an open.
In every
coset with prescribed determinant of $A$, choose
\[
       A=\operatorname{diag}\left(
         b,\begin{pmatrix}0&t_1\\1&0\end{pmatrix},\ldots,
                  \begin{pmatrix}0&t_{(n-1)/2}\\1&0\end{pmatrix}
                                      \right).
 \tag{A35.25}
\]
Then $C$ has eigenvalues
$1,t_1,t_1^{-1},\ldots,t_{(n-1)/2},t_{(n-1)/2}^{-1}$.
The scalar $b\ne0$ adjusts the determinant of $A$ without
affecting these eigenvalues. Choosing the $t_i$ to avoid
the prescribed finite root-of-unity conditions makes
$+1,-1$ the only such eigenvalues of $F_E$, both simple.
Multiplication by the geometric $SL_n$ allows every $A$
with a fixed determinant. Thus this nonempty open occurs
on the projection of every odd coset
$F_*\mathcal H^0$.

To make the exclusions uniform at all closed points,
let $D=\operatorname{rank}V=3600c-165$ and take
\[
                      N=\operatorname{lcm}(1,\ldots,2D^2).
 \tag{A35.26}
\]
The full normalized Frobenius polynomial is rational.
A root of unity in it has degree at most $D$ and hence
order at most $2D^2$, as in C44-A12. Thus it is an
$N$th root of unity. The invariant exclusion is detected
by $\det(F^N-I\mid W_0)\ne0$. The minimal imaginary
exclusion is detected by nonvanishing of the coefficient
of $(T-1)^2$ in
\[
                         \det(TI-F^N\mid W_\chi\oplus W_{\bar\chi}).
 \tag{A35.27}
\]
On the odd coset there are always at least two eigenvalues
one for $F^N$. That coefficient is nonzero precisely when
their total algebraic multiplicity is two. These are
regular conjugacy-invariant polynomial conditions.
Each pulls back to a nonzero function on the irreducible
coset $F_*\mathcal H^0$ by the preceding constructions.
Their product is therefore nonzero. This proves that the
two required opens meet simultaneously, without a product
monodromy assertion.

The actual geometric image is Zariski dense in
$\mathcal H$; its intersection with the identity component
is dense there. It follows that an element of the actual
odd arithmetic coset meets both conditions. The
polynomials are continuous in the $\ell$-adic topology, so
an open congruence neighbourhood of this element continues
to satisfy them. Include both the quadratic constant-field
cover and the invariant determinant character in the
finite quotient defining this neighbourhood. The finite
quotient Chebotarev argument of C44-A12 then gives a
closed point of odd degree in the desired conjugacy
classes. Repeating the argument after removing any
proper closed subset proves Zariski density. Thus for
these closed axes
\[
           \tau(W_0)=0,\qquad
           \tau(W_\chi\oplus W_{\bar\chi})=2,
           \qquad \tau(V)=3.
 \tag{A35.28}
\]
Here $\tau$ denotes the full generalized root-of-unity
multiplicity, with no semisimplicity assumption.

\paragraph{Actual lifting data and the sign conclusion.}
The product elliptic curve, its order-three action, the
bundle (A35.3), and the surface (A35.5) lift together over
the Witt vectors of the quadratic residue-field extension.
Use the same equation $y^2=x^3+1$, the Teichmuller lift of
$\zeta$, and the same product point bundles and surface.
The norm linearization and its descent lift as well.
For large $t$, relative Serre vanishing makes sections of
$I_SL$ lift modulo the maximal ideal. Taking each
character component is exact because $3$ is invertible.
Equivalently use the two $\mu_3$-weight direct summands of
the free section module over the Witt vectors of the
odd-degree field; their reduction is (A35.6). Hence the
actual selected pair lifts with all
these data. Smoothness of its proper special base and
of its two eigenmembers makes the corresponding relative
schemes smooth over the trait. The variable projector,
the tame self-adjoint idempotent
\[
          e=1-\frac{1+g+g^2}{3},\qquad u=g-g^2,
                   \qquad u^2=-3\quad\text{on }eV,
 \tag{A35.29}
\]
therefore have exactly the extending geometric and Hodge
meaning required in C44-A34a. The base field has odd
degree over $\mathbf F_p$, so normalized finite-field
Frobenius reverses $u$. The prime $p\equiv2\pmod3$ is
unramified inert in $E_0$. Finally (A35.4) and $p\ge11$
verify the self-dual integral lattice conditions of that
corollary.

Apply C44-A34a to $eV$ with its known positive subspace
equal to zero. Its potential Tate multiplicity is exactly
two by (A35.28); hence its entire actual algebraic space
is positive. The invariant block has potential Tate
space exactly $K$, whose Gram value is (A35.10).
Every algebraic invariant class is a rational multiple
of $w$: its coefficient is its rational intersection
with $w$ divided by $q(w,w)$. Thus it is positive too.
The two blocks are orthogonal and sum to $V$, proving
positivity on all of $A_B$.

If the algebraic image in $eV$ is nonzero, the actual
$u,F_s$ argument in C44-A34a spans both residual
dimensions and gives a rational nondegenerate Gram
basis. Otherwise that image is zero. Together with the
line $K$, these alternatives prove the asserted
dimensions one or three. They do not decide whether
either of the two additional potential Tate eigenvectors
is algebraic.
\end{proof}

\paragraph{Exact sources and scope.}
The tame alternating trace formula used in (A35.13) is
Deligne--Lusztig, \emph{Representations of reductive groups
over finite fields}, \emph{Annals of Mathematics} 103
(1976), Theorem 3.2, printed p.~119. The character
moment estimate uses the compact-support weight bound
in Deligne, \emph{La conjecture de Weil II},
Corollary 3.3.4, printed p.~206. The exact weighted
fourth moments, all rank counts and the odd arithmetic
coset calculation are proved above.

The monodromy inputs are Katz,
\emph{Larsen's alternative, moments, and the monodromy of
Lefschetz pencils}, Theorem 1.4.2(5) and its proof
pp.~531--535, Theorem 1.4.6 p.~535, and
Lemmas 2.1.4--2.1.6 pp.~540--541,
\url{https://web.math.princeton.edu/~nmk/LAMMLF.pdf}.
The nonzero local vanishing cycles and their simultaneous
quadratic inertia are SGA 7 II, Expose XV,
Theorem 3.4(i),(iii), pp.~195--196, in odd characteristic.
The geometric jet arguments and the original-pencil
condition were verified explicitly here, rather than
inferred merely from the smoothness of the base.

These actual closed axes have two residual potential
Tate dimensions after their explicit positive line is
removed. Their sign follows from the imaginary-quadratic
Frobenius reversal and quaternion packet calculation
of C44-A34, not from the one-direction criterion C44-A15.
In Grothendieck's terms, this settles a nonempty additional
family of the JAC1964 QH9-axis term in H9-T7, corresponding
to cote~44, Theorem~5, scan~004, with the primitive
restriction in scans~008--009. It does not prove positivity for an
arbitrary fixed axis, and it does not exhibit a wrong-sign
cycle or an algebraic representative of every Tate
eigenvector. No scan~049 product-form assertion follows.
The separate ambient abelian term in a full H9-T7 height
statement remains conditional on the unrefereed 032/Milne
2026 preprints; the axis sign proved here is unconditional.



% END work/axis/c3_axis_sign.tex

% BEGIN work/arithmetic/quaternion_rank_four_diagnostic.tex
% C44_POS /root/fermat_aj. Formal sharpness diagnostic, not a geometric counterexample.
\section{The remaining rank-four sign is invisible to the comparison data}
\label{sec:c44-a34b}

\paragraph{C44-A34b: scope.}
The following is a formal rational-form and orthogonal-crystal diagnostic.
It is not an actual motive, family of varieties, or space of algebraic
cycle classes, and hence is not a counterexample to a geometric standard
conjecture. It tests the linear, integral, Hodge, and finite-field-descent
data used in C44-A34. Those data allow an all-negative residual space of
rank four, both when the normalized Frobenius eigenvalues are $\pm1$
and when they are $\pm i$. Thus an exclusion of that case must use an
additional input beyond the comparisons proved there.

Put
\[
 E=\mathbf Q(I),\quad I^2=-3,\qquad
 B=(-3,-1)=E\oplus EJ,\quad J^2=-1,\quad Ja=\overline aJ.
\]
Let $\mathcal O=\mathbf Z[(-1+I)/2]$ and
$\mathcal R=\mathcal O\oplus\mathcal OJ$. This is an order in $B$.
On its rational vector space define the two symmetric forms
\[
 q_+(x,y)=\operatorname{Trd}(x\overline y),\qquad
 q_-(x,y)=-\operatorname{Trd}(x\overline y).
 \tag{A34b.1}
\]
Thus $q_\pm(x,x)=\pm2\operatorname{Nrd}(x)$.
In the basis $1,(-1+I)/2$ on each of the two $E$-summands,
the Gram matrix of $q_+$ is the orthogonal sum of two copies of
$\bigl(\begin{smallmatrix}2&-1\\-1&2\end{smallmatrix}\bigr)$.
Its determinant is $9$. The form $q_+$ is positive definite and
$q_-$ is negative definite of rank four.

Write $u=L_I$ for left multiplication by $I$. Consider either of
the following integral operators on $\mathcal R=E\oplus EJ$:
\[
 F_+(x+yJ)=\overline x+\overline yJ,
 \qquad
 F_-(x+yJ)=J(x+yJ)=-\overline y+\overline xJ.
 \tag{A34b.2}
\]
Both preserve both forms, satisfy $Fu=-uF$, and have finite order.
The two common rational characteristic polynomials are respectively
\[
             (t-1)^2(t+1)^2,\qquad (t^2+1)^2.
 \tag{A34b.3}
\]
Here $F_+^2=1$ and $F_-^2=-1$. Also $u^\dagger=-u$ for either
choice of sign of the form.

\paragraph{Finite-place comparisons, including the integral lattices.}
At every finite prime $\ell$ there is an element
$b_\ell\in\mathcal R\otimes\mathbf Z_\ell$ with reduced norm $-1$.
For odd $\ell$, choose $x,y\in\mathbf Z_\ell$ with
$x^2+y^2=-1$ and take $b_\ell=x+yJ$.
Such a pair exists: over $\mathbf F_\ell$ the sets of squares and
$-1$ minus the squares each have $(\ell+1)/2$ elements, so they
intersect. At least one coordinate of a resulting solution is
nonzero. Fixing the other coordinate, the nonzero partial
derivative is a unit; successive correction of that coordinate
modulo $\ell^n$ lifts the solution. At $\ell=2$, take
\[
                  b_2=2+tI,\qquad t^2=-5/3.
 \tag{A34b.4}
\]
The indicated $2$-adic unit has a square root since $-5/3\equiv1
\pmod8$: given an odd solution modulo $2^n$, $n\ge3$, adding
$0$ or $2^{n-1}$ changes its square by $0$ or $2^n$ modulo
$2^{n+1}$, giving a compatible convergent sequence.
Thus $\operatorname{Nrd}(b_2)=4+3t^2=-1$.
In every case $b_\ell^{-1}=-\overline{b_\ell}$ belongs to the
same order. Right multiplication by $b_\ell$ therefore gives an
integral isometry
\[
 (\mathcal R\otimes\mathbf Z_\ell,q_+)
             \xrightarrow{\ \sim\ }
 (\mathcal R\otimes\mathbf Z_\ell,q_-).
 \tag{A34b.5}
\]
It commutes with the left $E$-action. For $F_-=L_J$ it also
commutes with $F_-$ at every prime. Thus in the $\pm i$ example
even the two indicated operators are preserved by all finite-place
comparisons. For $F_+$ this extra compatibility at $2$ is not
needed: C44-A34 requires the Betti comparisons to intertwine $E$,
and does not posit a Betti Frobenius operator. The special-fibre
operators $F_+$ themselves are the same rational operator at all
coefficient primes.

\paragraph{Absolute crystal and genuine finite-field descent data.}
Fix $p\ge7$ with $p\equiv2\pmod3$, and fix an odd $f$.
Set $W=W(\overline{\mathbf F}_p)$, $K=W[1/p]$ and
\[
 M=\mathcal R\otimes W,\quad D=M[1/p],\quad
 q_D=q_-,\quad \Phi=1\otimes\sigma.
 \tag{A34b.6}
\]
The lattice $M$ is self-dual because $p\nmid9$. It is stable under
$u$ and $\Phi$. The latter is a semilinear isometry commuting with
both choices of $F$. All normalized slopes and elementary-divisor
exponents are zero, so $H=S=0$ and
\[
     D^{\Phi=1}=(\mathcal R\otimes\mathbf Q_p,q_-).
 \tag{A34b.7}
\]
In particular its Hasse invariant is $+1$, as predicted by C44-A33.

This is also compatible with the full finite-field descent relation,
not just with the valuations of the Frobenius eigenvalues. Let
$m=2$ for $F_+$ and $m=4$ for $F_-$, and on
$\mathcal R\otimes W(\mathbf F_{p^{fm}})$ put
\[
                       \delta=F^{-1}\sigma^f.
 \tag{A34b.8}
\]
Its $m$th power is the identity. It is the semilinear descent datum
for the finite unramified extension of degree $m$ over
$W(\mathbf F_{p^f})$. Finite etale descent therefore gives a free
self-dual lattice $M_f$ over $W(\mathbf F_{p^f})$, whose extension
to $W$ is $M$. Since $F$ has rational matrices, $\Phi$ commutes
with $\delta$ and preserves $M_f$. On this descended module,
\[
                            \Phi^f=F.
 \tag{A34b.9}
\]
Indeed, $\delta v=v$ means $\sigma^fv=Fv$.
Furthermore $\delta u\delta^{-1}=-u$, so the geometric $E$-action
is defined after the quadratic extension. Its integral order-three
generator $g=(-1+u)/2$ satisfies
$\delta g\delta^{-1}=g^{-1}$, exactly as in the inversion pattern
of C44-A34a. No claim that this formal descent comes from a variety
is made.

\paragraph{Hodge data and the missing real sign.}
Give $(\mathcal R\otimes\mathbf Q,q_+)$ the pure weight-four Hodge
structure of type $(2,2)^4$. The signed primitive form is a
polarization, and the $E$-action is by Hodge endomorphisms.
Its $p$-lattice is self-dual. The normalized Hodge multiplicities
are $h_0=4$ and $h_a=0$ for $a\ne0$, agreeing with (A34b.6).
Equivalently the unnormalized crystal is $p^2\Phi$ and has all
Hodge and Newton slopes equal to two. The filtration with
$\operatorname{Fil}^2=D$ and $\operatorname{Fil}^3=0$ is weakly
admissible: every subisocrystal has equal Hodge and Newton
numbers. The unnormalized finite-field Frobenius has eigenvalues
$p^{2f}$ times those in (A34b.3), so it has the required purity.
Take $T=(\mathcal R\otimes\mathbf Q,q_-)$ as a \emph{formal}
rational space in the role of the full potential Tate part, and
take its scalar extensions as all special-fibre realizations.
Equations (A34b.5)--(A34b.9) give their common rational pairings,
operators and absolute fixed-space identification. The designation
of $T$ does not manufacture algebraic cycle representatives.

All these forms and lattices meet the stated comparison constraints,
but $T$ has negative index four. In the $F_-$ example, the adjoint
on the quaternion endomorphism algebra is the canonical involution
for both $q_+$ and $q_-$; its positive adjoint trace form consequently
does not distinguish the two signs either. The missing invariant
is the ordered real sign of the scalar multiplying the definite
quaternion norm form. Its local reduced-norm classes at every
finite prime are trivial, as proved explicitly above. A global
rational change of basis cannot remove the sign, since the two
real forms have opposite definite signatures.

This gives no additional theorem about actual cycles. It shows
why the present finite-place, integral, Hodge and descent tests
alone do not exclude the all-negative rank-four case. Actual
geometric or motivic positivity would be an additional input;
neither a counterexample to that positivity nor a new closed-axis
instance is supplied. No progress-clock movement is attached to
this diagnostic.

For the descent step the exact background references are the
Stacks Project, Proposition~35.3.9, Tag~023N (effective module
descent), and Lemma~35.7.6, Tag~05B2 (finite local freeness).
The perfect pairing and its inverse descend as morphisms.
In the programme this limits a proof route for JAC1964
QH9-axis and H9-T7, hence cote~44, Theorem~5, scan~004,
and its primitive restriction on scans~008--009. It proves
no failure of that geometric sign or of cote~16 scan~049.

% END work/arithmetic/quaternion_rank_four_diagnostic.tex

% BEGIN work/scan_w/non_tate_quaternion_criterion.tex
% C44_POS /root/scan_w. A36/A36a scope relaxation; not a new instantiated family.
\section{The quaternion correspondence need act only on the non-Tate space}
\label{sec:c44-a36}

The extending imaginary-quadratic action in C44-A34 is stronger
than the arithmetic proof needs. What must have an actual Betti
and integral crystalline comparison is the whole quadratic block.
The skew correspondence itself may exist only on the special
fibre, and may vanish on the entire potential Tate space.
The distinction is useful because it avoids both an extending
Tate projector and any assertion that a numerical polarization
is the geometric pairing.

\begin{theorem}[C44-A36: a special-fibre non-Tate correspondence suffices]
Let $p$ be odd, $f$ odd, and $d$ a positive squarefree integer
such that $p$ is unramified and inert in
$E=\mathbf Q(\sqrt{-d})$.
Let a signed geometric primitive cohomology block of weight
$2r$ over $\mathbf F_{p^f}$ have the following data.
\begin{enumerate}
\item A nondegenerate rational Betti form $q_{\rm B}$ from an
actual mixed-characteristic lift, with its Hodge--Riemann sign,
and isometries
$q_{\rm B}\otimes\mathbf Q_\ell\simeq q_\ell$ for every $\ell\ne p$.
The rational $p$-adic form $q_{\rm B}\otimes\mathbf Q_p$
has a self-dual $\mathbf Z_p$-lattice.
\item An absolute normalized orthogonal crystal
$(D,q_D,\Phi)$ with a self-dual
$W(\overline{\mathbf F}_p)$-lattice $M$.
Writing $H=\sum_{a>0}a h_a$ for its positive elementary-divisor
sum, assume the actual Hodge comparison gives
\[
           \epsilon_\infty(q_{\rm B})=(-1)^H.
 \tag{A36.1}
\]
In particular this holds when the elementary-divisor
multiplicities are the Hodge numbers of the primitive block.
\item Normalized geometric Frobenius $F$ is an isometry and
has one common rational characteristic polynomial in all
$\ell$-adic and crystalline realizations, with all complex
eigenvalue absolute values equal to one. Its relation with
absolute crystalline Frobenius is the one in C44-A34:
$F$ is linear, commutes with $\Phi$, and its eigenvalue
valuations divided by $f$ are the slopes of $\Phi$.
\item There is a nondegenerate rational space $T$ of actual
cycle classes which spans the full generalized potential
Tate space in every realization. Its crystalline classes
identify $T\otimes\mathbf Q_p$ with a subspace of
$D_0^{\Phi=1}$.
\item Let $U$ be the non-Tate generalized primary complement.
There is one actual special-fibre rational algebraic
correspondence inducing on $U$, in every realization, an
operator $u$ with
\[
               u^\dagger=-u,\qquad u^2=-d,
               \qquad Fu=-uF.
 \tag{A36.2}
\]
No action of $u$ on $T$ is required. No extension of this
correspondence or of the projector onto $U$ is required,
and $u$ need not preserve any integral lattice.
\end{enumerate}
Then the negative index of the actual rational intersection
form $q_T$ is divisible by four. A positive subspace of $T$
of codimension at most three consequently makes $q_T$
positive definite.

The identities in (A36.2) are simultaneous realization
conditions. Identities of the actual correspondences in
the Chow category suffice. Without an additional theorem,
their validity in just one cohomology theory is not being
promoted to all realizations by numerical equivalence.
\end{theorem}

\begin{proof}
If $T=0$ there is nothing to prove. The root-of-unity and
non-root-of-unity generalized primary spaces are orthogonal,
since the pairing pairs only reciprocal eigenvalues.
Their common cohomological projectors are rational
polynomials in $F$, obtained by the Chinese remainder
theorem from the common characteristic polynomial.
Thus, if desired, replace the given correspondence by
$P_UuP_U$. It is then zero on $T$ and still satisfies
(A36.2) on $U$ in every realization. These are actual
polynomial correspondences with common cohomological
idempotent realizations; no Chow-idempotence assertion is
needed. Absolute crystalline Frobenius commutes with this
degree-zero algebraic correspondence by naturality.

\emph{The non-Tate calculation uses no Betti action.}
Replace $F$ by its rational Jordan semisimple part $F_s$,
using the common polynomial as in C44-A34, and put
$z=F_s-F_s^{-1}$ on $U$. Then $u^2=-d$, $z^2=c$,
$uz=-zu$, and
\[
                 u^\dagger=-u,\quad z^\dagger=-z,
                       \quad (uz)^\dagger=-uz.
 \tag{A36.3}
\]
The eigenvalues of $F_s$ occur in distinct quartets
$\{\lambda,-\lambda,\lambda^{-1},-\lambda^{-1}\}$.
Consequently the algebra generated by $c$ is a product of
totally real fields $C_i$, with totally negative distinguished
elements $c_i$ and packet ranks $4m_i$.

For clarity, the local form calculation can be made directly.
Trace duality gives a $C_i$-bilinear symmetric form $b$.
For any $v$ of nonzero norm $a$, the four vectors
$v,uv,zv,uzv$ are pairwise orthogonal with respective norms
\[
                         a,\quad da,\quad -c_i a,\quad -dc_i a.
 \tag{A36.4}
\]
Their span and its orthogonal complement are stable under
$u,z$. Induction gives scalar multiples of the quaternion
norm form $n_i=\langle1,d,-c_i,-dc_i\rangle$, including
split quaternion algebras and repeated packets. For every
finite local extension $L/k$ of degree $e$,
\[
 \begin{split}
 \det\operatorname{Tr}_{L/k}(a n_i)&\in k^{\times2},\\
 \epsilon_k\bigl(\operatorname{Tr}_{L/k}(a n_i)\bigr)
       &=(-1,-1)_k^e(-d,N_{L/k}c_i)_k.
 \end{split}
 \tag{A36.5}
\]
Indeed write the trace form as $\theta\perp d\theta$,
where $\theta=\operatorname{Tr}_{L/k}\langle a,-ac_i\rangle$
has rank $2e$ and determinant $(-1)^eN(c_i)$ modulo squares,
and expand the diagonal Hilbert symbols. In particular the
invariant is independent of every unknown scalar $a$.

Define the positive rational comparison form
\[
        N^+=\mathop{\perp}\limits_i
                \bigl(\operatorname{Tr}_{C_i/\mathbf Q}n_i\bigr)^{m_i}.
 \tag{A36.6}
\]
It has square determinant. For each $\ell\ne p$, (A36.5)
shows that the actual $U_\ell$ has the same determinant
and Hasse invariant as $N^+\otimes\mathbf Q_\ell$.
This is only a special-fibre calculation; $u$ was not
applied to $V_{\rm B}$ at any point.

\emph{The central non-Tate part descends by absolute Frobenius.}
Both $u$ and $F_s$ commute with $\Phi$, so they act on
the actual fixed space $D_0^{\Phi=1}$, with (A36.2)--(A36.3)
on its non-Tate orthogonal complement to $T$.
Apply the same norm calculation there. If $N_0^+$ denotes
the factors in (A36.6) with $p$-adic $c_i$-valuation at
least zero, this gives
\[
    D_0^{\Phi=1}=T_{\mathbf Q_p}\perp U_0,\qquad
    \det U_0\in\mathbf Q_p^{\times2},\qquad
    \epsilon_p(U_0)=\epsilon_p(N_0^+).
 \tag{A36.7}
\]
The remaining factors have $v_p(c_i)<0$. Each corresponding
quartet has two positive absolute slopes
$-v_p(c_i)/(2f)$. For
$S=\sum_{\lambda>0}\lambda\dim D_\lambda$ it follows that
\[
       \sum_{i,\,v_p(c_i)<0}m_i v_p(Nc_i)=-fS.
 \tag{A36.8}
\]
At an odd inert prime $(-d,x)_p=(-1)^{v_p(x)}$ and
$(-1,-1)_p=1$. Since $f$ is odd, the noncentral factors
of $N^+$ have Hasse invariant $(-1)^S$.
C44-A33 applies to the whole integral crystal $M$, not
to an unproved integral projection onto $U$ or $T$.
It gives
\[
              \epsilon_p(T)=(-1)^H\epsilon_p(N^+).
 \tag{A36.9}
\]
This uses (A36.7), the invariant $(-1)^{H-S}$ of the
whole central fixed space, and (A36.8). It makes no
algebraicity assertion about $U_0$.

\emph{Real Hasse invariant and determinant.}
At every $\ell\ne p$, orthogonality and the square norm
determinants give
\[
 \epsilon_\ell(T)=\epsilon_\ell(q_{\rm B})\epsilon_\ell(N^+),
 \qquad
 \det T=\det q_{\rm B}\pmod{\mathbf Q_\ell^{\times2}}.
 \tag{A36.10}
\]
The self-dual Betti $p$-lattice has Hasse invariant one.
Using (A36.1), positivity of $N^+$, and rational Hilbert
reciprocity in (A36.9)--(A36.10), exactly as in C44-A34,
therefore gives
\[
                            \epsilon_\infty(T)=1.
 \tag{A36.11}
\]

It remains to prove that the negative index is even,
without putting a nonexistent $E$-action on $T$.
Both determinants in (A36.10) are rational. Their ratio
is a square in every $\mathbf Q_\ell$ except possibly
$\mathbf Q_p$. The deleted-one-place squareclass
injection used in C44-A15 therefore makes this ratio
a square in $\mathbf Q$. To recall why, a nonsquare
rational ratio defines a nontrivial quadratic number
field; Chebotarev supplies an inert finite prime outside
the one excluded prime, where the ratio is not a square.
Hodge--Riemann makes the negative index of $q_{\rm B}$
even, so its ordinary determinant is positive. Hence
\[
                                \det T>0.
 \tag{A36.12}
\]
The negative index $b$ of $T$ is consequently even.
Its real Hasse invariant is then $(-1)^{b/2}$; (A36.11)
forces $4\mid b$. A positive subspace of codimension
at most three makes $b\le3$, proving the last assertion.
\end{proof}

\begin{corollary}[C44-A36a: whole-variable actual-axis form]
Let $B\subset X$ be a smooth $(L,L)$ axis in an abelian
sixfold over $\mathbf F_{p^f}$, with $p\ge7$, $f$ odd,
and $p\nmid\chi(X,L)$. On its whole variable middle
cohomology suppose that the full potential Tate space has
an actual nondegenerate rational cycle Gram certificate.
Suppose also that its non-Tate complement has a common
special-fibre algebraic correspondence satisfying (A36.2)
for a positive squarefree $d$ with $p$ inert in
$\mathbf Q(\sqrt{-d})$.

Then the actual algebraic variable form has negative index
divisible by four. In particular it is positive definite
if the certificate contains a known positive subspace of
codimension at most three. Neither the correspondence
$u$ nor the certificate is required to lift.
\end{corollary}

\begin{proof}
C44-A05 supplies an embedded lift of the underlying axis,
without $u$ or the cycles. C44-A22a and C44-A33a give free
crystalline cohomology, the actual self-dual variable lattice,
the integral ambient orthogonal splitting, and its Hodge
elementary divisors. The same exterior-algebra argument
on the integral Betti lattice gives a self-dual
$\mathbf Z_p$-lattice, as in C44-A34a. For the whole variable
space the positive Hodge sum is $956\chi(L)-50$, an even
integer. Thus all comparison and integral hypotheses of
C44-A36 hold independently of the special correspondence.

Use the common variable correspondence
$P=\Delta_B-\kappa^*(L_h^2)^{-1}\kappa_*$ on all certificate
representatives and on $u$. The common Frobenius primary
projectors then isolate the stated spaces in every
realization. A nonzero rational Gram determinant proves
independence in every realization; the certificate's full
potential-Tate dimension proves that it spans there.
The common polynomial and purity are the geometric
comparison already verified in C44-A15/A18b and the
Katz--Messing trace argument in C44-A34a. Algebraic
crystalline classes are fixed by normalized absolute
Frobenius. C44-A36 now applies.
As before, pairing with the rational nonsingular Gram
basis reconstructs rational coordinates for every further
algebraic class; no assumption $D(B)$ is introduced.
\end{proof}

\paragraph{What is and is not supplied by this relaxation.}
The proof requires only the non-Tate quaternion relation on
the special fibre. It removes an extending-field-action
hypothesis and does not replace it by an extending Tate
projector. Its determinant argument uses the exact
deleted-place squareclass statement checked in the proof
of Ancona--Marmora, Corollary~4.4, printed p.~629; the
crystalline and local norm inputs are those proved in
C44-A33/A34.

It does not produce $u$ from a reciprocal or even Frobenius
polynomial alone. Such a polynomial specifies stable
conjugacy data, whereas (A36.2) constrains the actual local
coefficient forms. It also does not make an arbitrary
subspace of the potential Tate block algebraic or
nondegenerate. The full cycle certificate in A36a must
be supplied, or a separate argument must prove the
all-or-nothing algebraicity used in A34a.

No additional closed-axis family with a genuinely
special-only correspondence has been certified here.
The tame order-three family C44-A35 already has a
geometric extending action; it does not become a new
instance merely by forgetting that extension. This is
a precise scope relaxation, not a separate progress-clock
event or a new universal fixed-axis sign theorem.

% END work/scan_w/non_tate_quaternion_criterion.tex

% BEGIN work/scan_w/universal_quaternion_boundary.tex
% C44_POS /root/scan_w. A36b reserved by root; diagnostic, not a new sign case.
\section{A geometric boundary of the quaternion method}
\label{sec:c44-a36b}

\begin{proposition}[C44-A36b: no automatic quaternion action]
Take any of the actual closed axes selected in C44-A30, and retain its
nonzero residual primitive block $W$, its common algebraic projector
$P_W$, and normalized Frobenius $F$. Thus
\[
 \dim W=d=3600t^6-166,\qquad
 \operatorname{Spec}(F)=
 \{\lambda_1,\lambda_1^{-1},\ldots,
                         \lambda_{d/2},\lambda_{d/2}^{-1}\},
\]
and a product $\prod_i\lambda_i^{n_i}$ is a root of unity only when
every $n_i$ is zero. These are the maximal-connected-Frobenius-torus
axes of A30, with their known positive variable class; no additional
assumption is made here.

For every positive integer $a$, the following statements hold.
\begin{enumerate}
\item $F^a$ has distinct eigenvalues and no eigenvalue which is a root
of unity. Its spectrum contains no pair $\{\mu,-\mu\}$.
\item The image on $W$ of all rational algebraic correspondences over
the algebraic closure is precisely the commutative algebra
\[
                  \operatorname{End}_{\rm alg}(W)=\mathbf Q[F].
 \tag{A36b.1}
\]
Here the left side means the image in the chosen Weil realization,
with the same equality in every such realization. It is not an
assertion that homologically trivial Chow correspondences vanish.
\item No nonzero linear operator on $W$ can satisfy
$uF^a=-F^a u$. In particular there is no faithful action on $W$ of a
quaternion algebra by geometric algebraic correspondences.
\end{enumerate}
Thus an anticommuting quaternion structure cannot be deduced merely
from the fact that a variety is a smooth fixed axis in an abelian
sixfold. This is a method boundary on actual axes, not a negative
intersection example; the axes in question already have the positive
form proved in A30.
\end{proposition}

\begin{proof}
Suppose two eigenvalues of $F^a$ agree, or are negatives of one
another. In either case their quotient is a root of unity. In the
notation of the statement this is a relation between two of the
weights $e_i,-e_i$ of the maximal torus. Distinct weights give a
nonzero relation, contrary to A30. Equal weights cannot give a
negative pair in characteristic zero. Similarly an eigenvalue which
is a root of unity would give a nonzero multiple of one weight as a
relation. This proves the first assertion. In particular $F^a$ is
regular semisimple and
\[
 \mathbf Q[F^a]=\mathbf Q[F],\qquad
 Z_{\operatorname{End}(W_\ell)}(F^a)
                       =\mathbf Q_\ell[F].
 \tag{A36b.2}
\]
The first equality follows because both rational algebras have
dimension $d$ and one is contained in the other. The second is the
centralizer of an operator with $d$ distinct eigenvalues. The same
argument works with the crystalline coefficient field.

Let $A$ be the action of a rational algebraic correspondence
$P_W\Gamma P_W$ over the algebraic closure. The finitely many
components of $\Gamma$ and the projector are defined over a finite
extension of the base field. Accordingly $A$ commutes with $F^b$ for
some positive integer $b$. By (A36b.2), over any coefficient field
it has a unique expression
\[
                         A=\sum_{i=0}^{d-1}c_i F^i.
 \tag{A36b.3}
\]
The coefficients are rational, as follows. Put
\[
 G_{ij}=\operatorname{Tr}_W(F^{i+j}),\qquad
 b_j=\operatorname{Tr}_W(AF^j),\qquad 0\leq i,j<d.
 \tag{A36b.4}
\]
All these numbers are rational intersection numbers of algebraic
correspondences. The projector vanishes on the other cohomological
degrees, and degree four is even, so the usual alternating
correspondence trace is the ordinary trace displayed here.
The matrix $G$ is the trace matrix of the finite etale algebra
$\mathbf Q[F]$. Equivalently it is the product of a Vandermonde
matrix on the distinct roots with its transpose. It is nonsingular.
The system $Gc=b$ therefore gives $c\in\mathbf Q^d$.
Conversely each rational polynomial in $F$, with $P_W$ on both
sides, is induced by an actual rational algebraic correspondence.
This proves (A36b.1), without a new appeal to a Tate conjecture or
to numerical--homological equivalence.

Finally $uF^a=-F^a u$ sends the $\mu$-eigenline into the
$-\mu$-eigenspace. That space is zero by the first assertion, so
$u=0$. Quite independently, a nonzero unital action of a quaternion
algebra would embed a noncommutative simple algebra into the
commutative algebra (A36b.1), which is impossible. This includes
split quaternion algebras. All arguments persist after any finite
extension of the field of definition.
\end{proof}

\paragraph{Why the canonical positive involution does not repair this.}
Purity gives the positive involution $F\mapsto F^{-1}$ on the
commutative algebra $\mathbf Q[F]$. Its positivity is a statement
about $\operatorname{Tr}(aa^\dagger)$ in that algebra. It does not
recover the coefficient of the geometric form on a simple
Frobenius component. Explicitly, on a one-dimensional space over
a CM field $C$ with maximal real subfield $C^+$, all forms
\[
          q_h(x,y)=\operatorname{Tr}_{C/\mathbf Q}(h x\bar y),
                       \qquad h\in(C^+)^\times,
 \tag{A36b.5}
\]
give the same adjoint $a^\dagger=\bar a$ to multiplication by $a$.
Choosing $h$ positive or negative at a real embedding changes the
real signature without changing that involution. Over a local
field the corresponding hermitian norm class likewise remains
extra information. Formula (A36b.5) is a statement about the method's
input, not a construction of geometrically negative algebraic
classes. The quaternion packet argument succeeds because its
four-dimensional norm form removes precisely this scalar dependence;
the commutative Frobenius algebra alone does not do so.

\paragraph{An automatic Artin doubling, and its exact limitation.}
There is an actual way to manufacture the anticommuting operator
without an automorphism of the original axis. Tensor a quadratic
cohomology block $(V,q,F)$ with the two-point Artin motive
$h^0(\operatorname{Spec}\mathbf F_{q^2})$ over $\mathbf F_q$.
In the geometric point basis its form and Frobenius are
\[
 q'=q\perp q,\qquad
 F'=\begin{pmatrix}0&F\\F&0\end{pmatrix},\qquad
 u=\begin{pmatrix}0&-1\\1&0\end{pmatrix}.
 \tag{A36b.6}
\]
The matrix units are actual correspondences on the two geometric
points, defined over the quadratic constant extension. Direct
multiplication gives
\[
              u^2=-1,\qquad u^\dagger=-u,\qquad uF'=-F'u.
 \tag{A36b.7}
\]
This supplies $\mathbf Q(i)$, not an arbitrary imaginary field.
For an application of the existing comparison the residue prime
must be odd and the other hypotheses of A34/A36, or the split-prime
version A37, must still be checked. In particular this construction
does not meet the unramified quadratic-field hypothesis of A41 at
two. It does not by itself supply a full nondegenerate cycle
certificate or repair a degenerate algebraic span.

If the original full potential Tate space is a nondegenerate actual
cycle space $T$, the new one is exactly $T\oplus T$: multiplying a
Frobenius eigenvalue by $-1$ does not change whether it is a root of
unity. Its actual form is $q_T\perp q_T$, and hence
\[
                       \nu(T\oplus T)=2\nu(T).
 \tag{A36b.8}
\]
Even when every remaining comparison hypothesis is satisfied, the
quaternion conclusion $4\mid\nu(T\oplus T)$ gives only
$2\mid\nu(T)$. This is the determinant parity already supplied by
A15/A17. It does not exclude a negative plane in the original
algebraic space. Thus the automatic Artin construction supplies the
algebra, but loses exactly the factor of two needed to improve the
rank-two or rank-three sign gate.

\paragraph{Source and continuation scope.}
The actual selected-axis input is exactly A30. Its published
maximal-torus theorem is Chin, \emph{Independence of $\ell$ of
monodromy groups}, Theorem~5.7 and Corollary~5.8, with the arithmetic
hypotheses checked in the complete A30 proof. For the distinction
between a positive involution on the Frobenius algebra and a
geometric intersection form, see Milne, \emph{On the Tate and
Standard Conjectures over Finite Fields}, version~1.1, July~23,
2019, p.~1 and Section~2, pp.~6--8; the trace construction is
recalled on p.~5. The displayed countervariation of $h$ supplies
the needed elementary verification here.

These diagnostics leave the universal fixed-axis sign open. A
different construction would have to control the actual local
hermitian coefficient forms, or provide some geometric input of
equivalent strength, without merely doubling the unknown sign.
No new eligible axis, standard-conjecture closure, or progress-clock
movement is claimed in this subsection.

% END work/scan_w/universal_quaternion_boundary.tex

% BEGIN work/scan_w/split_prime_quaternion_sign.tex
% C44_POS /root/scan_w. A37--A37b; one continuing quaternion extension.
\section{Split coefficient primes and the characteristic-five lattice}
\label{sec:c44-a37}

\begin{theorem}[C44-A37: the split-prime quaternion comparison]
In C44-A34 or C44-A36, replace the assumptions that $p$ is
unramified inert in $E=\mathbf Q(\sqrt{-d})$ and $f$ is odd
by the assumptions that $p$ is odd and split in $E$, with
arbitrary positive residue degree $f$.
Retain all the other hypotheses, including the actual
self-dual Betti and crystalline lattices and the full
potential-Tate cycle space $T$. Then the conclusions of
those theorems still hold: the negative index of $q_T$ is
divisible by four, and a positive subspace of codimension
at most three forces positivity.

This applies also in the exact scope of C44-A36: the
special-fibre skew correspondence need exist only on
the non-Tate space $U$, and need not lift or preserve a
lattice. No action on the Tate space is required.
\end{theorem}

\begin{proof}
If $T=0$, the conclusion is immediate; otherwise $D_0\ne0$.
Only the rational $p$-adic comparison changes. All
non-Tate quaternion packets and the positive rational
comparison form $N^+$ are constructed exactly as in
C44-A34, or in the non-Tate-only proof C44-A36.

Choose $\alpha\in\mathbf Q_p^\times$ with $\alpha^2=-d$.
On $U$ the two $u$-eigenspace projectors have coefficients
in $\mathbf Q_p$:
\[
                     U=U_+\oplus U_-,
                    \qquad u|U_\pm=\pm\alpha.
 \tag{A37.1}
\]
Because $\alpha$ is fixed by the absolute Witt Frobenius,
and $u\Phi=\Phi u$, each eigenspace is an absolute
$F$-isocrystal. The linear geometric Frobenius $F$
anticommutes with $u$ and commutes with $\Phi$.
It therefore induces isomorphisms of absolute
$F$-isocrystals
\[
                          F:U_+\xrightarrow{\sim}U_-.
 \tag{A37.2}
\]
These isomorphisms preserve every normalized Newton
slope, not just the total dimensions.

Every nonzero slope of $D$ lies in $U$, because roots
of unity have valuation zero. If
\[
       S_+=\sum_{\lambda>0}\lambda
                    \dim_K(U_+\cap D_\lambda),
\]
then (A37.2) gives
\[
       S=\sum_{\lambda>0}\lambda\dim_KD_\lambda
                            =2S_+\in2\mathbf Z.
 \tag{A37.3}
\]
Here $S_+$ is an integer. One may take the determinant
of the entire positive-slope isocrystal inside $U_+$:
its slope is $S_+$ and is the valuation of the determinant
of a semilinear matrix over $K=W(\overline{\mathbf F}_p)[1/p]$.
The valuation group of $K$ is $\mathbf Z$. This argument
requires neither an integral $u$-stable lattice nor an
assumption about a filtered subobject.

At this split prime $-d$ is a square. The elementary
trace-Hasse formula proved in (A34.10) gives, for every
local packet, every scalar, and every multiplicity,
\[
    \epsilon_p\bigl(\operatorname{Tr}_{L/\mathbf Q_p}
           \langle a,da,-ac,-dac\rangle\bigr)
       =(-1,-1)_p^{[L:\mathbf Q_p]}(-d,Nc)_p=1.
 \tag{A37.4}
\]
All these packet determinants are squares. In particular
\[
         \epsilon_p(N^+)=\epsilon_p(N_0^+)
                           =\epsilon_p(N_{\ne0}^+)=1.
 \tag{A37.5}
\]
The same calculation applies directly on the actual
absolute fixed non-Tate space $U_0\subset D_0^{\Phi=1}$.
No inference from an unramified base extension back to
$\mathbf Q_p$ is used.

C44-A33 gives
$\epsilon_p(D_0^{\Phi=1})=(-1)^{H-S}$.
Using (A37.3) and the orthogonal decomposition
$D_0^{\Phi=1}=T_{\mathbf Q_p}\perp U_0$ therefore yields
\[
                \epsilon_p(T)=(-1)^H
                              =(-1)^H\epsilon_p(N^+).
 \tag{A37.6}
\]
This is the same rational $p$-adic identity used in
(A34.16) and (A36.9), now with no condition on $f$.
The other finite-place comparisons, the self-dual
Betti $p$-lattice, and Hilbert reciprocity consequently
give $\epsilon_\infty(T)=1$ exactly as before.
For the A34 formulation the rational quadratic action
on $T$ makes its negative index even. For the A36
formulation, comparison of its determinant with the
positive Betti determinant at every finite prime other
than $p$, followed by the deleted-place squareclass
injection, gives that evenness. Thus in either formulation
the negative index is divisible by four, proving the
theorem and the positive-codimension consequence.
\end{proof}

\begin{corollary}[C44-A37a: split-prime actual-axis criterion]
Let the actual axis, polarization and order-three
automorphism satisfy the hypotheses of C44-A34a,
except that
\[
                         p\equiv1\pmod3
 \tag{A37a.1}
\]
and the residue degree $f$ is unrestricted.
In particular retain the actual relation
$FgF^{-1}=g^{-1}$, the embedded lift of the given data,
$p\ge7$, and $p\nmid\chi(L)$.
Then the same conclusion holds: a known positive
$u,F$-stable algebraic span in the imaginary-quadratic
block, with at most two residual potential Tate
dimensions, makes its full algebraic form positive.
The complementary invariant block may be handled by
C44-A19, exactly as in C44-A34a.

The whole-variable criterion C44-A36a likewise holds
when its coefficient prime is split, with arbitrary $f$,
under its explicitly retained full cycle Gram certificate.
\end{corollary}

\begin{proof}
The actual integral, Hodge and common-realization
checks in C44-A34a do not use inertness or odd residue
degree. The field $\mathbf Q(\sqrt{-3})$ is now split
at $p$. Its special linear Frobenius must still reverse
$u=g-g^2$; this is an explicit hypothesis, not a
consequence of the residue congruence. Apply C44-A37
in place of C44-A34. The residual two-dimensional
space is still irreducible under $u,F_s$, which generate
$(-3,1)=M_2(\mathbf Q_\ell)$ there. Thus the same
all-or-nothing algebraic spanning and rational Gram
repair supply the full $T$ exactly when needed. A
positive subspace of codimension two and the divisibility
by four finish the proof. For C44-A36a, retain its
full Gram certificate and apply the non-Tate-only
version of C44-A37 directly.
\end{proof}

\begin{proposition}[C44-A37b: the integral axis splitting at $p=5$]
The actual-axis integral and Hodge assertion C44-A33a
holds with $p\ge5$ in place of $p\ge7$, under the same
condition $p\nmid\chi(X,L)$. Consequently the actual-axis
versions C44-A34a, C44-A36a and C44-A37a have the same
range improvement whenever their other hypotheses hold.

In particular the actual family and every conclusion
of C44-A35 also hold for $p=5$. It gives genuine
Lefschetz axes with a known positive line and two
remaining potential Tate dimensions in the imaginary
quadratic block, and proves positivity without 032.
No assertion for $p=2$ or $p=3$ is made by this extension.
\end{proposition}

\begin{proof}
The earlier exclusion of $5$ came from separately
forming spectral projectors for three Lefschetz
eigenvalues. The inverse itself does not require those
projectors.

Let $R$ be $\mathbf Z_p$ or
$W(\overline{\mathbf F}_p)$, and let $H^1$ be the free
rank-twelve module with the unimodular polarization
bivector $h$. This unimodularity follows from
$p\nmid\chi(L)$; the polarization degree is $\chi(L)^2$.
Write $L=L_h$ for multiplication by $h$ and
$\Lambda=\Lambda_h$ for contraction by its integral
inverse. On $\bigwedge^4H^1$ set
\[
                              A=\Lambda L.
 \tag{A37b.1}
\]
Over the fraction field its eigenvalues on
$L^2P^0,LP^2,P^4$ are $12,6,2$. Hence
\[
              A^3-20A^2+108A-144I=0.
 \tag{A37b.2}
\]
The integral module is free, so this rational
exterior-algebra identity is an identity over $R$.
For $p\ge5$, $144$ is a unit and
\[
                A^{-1}=\frac{A^2-20A+108I}{144}.
 \tag{A37b.3}
\]
This only inverts $2$ and $3$.

There is also an exact integral commutator calculation.
Let $H=[\Lambda,L]$, which acts by $6-k$ on degree
$k$. Then on degree four
\[
    A^2=\Lambda L\Lambda L
       =\Lambda(\Lambda L-H)L
       =\Lambda^2L^2-\Lambda HL
       =\Lambda^2L^2,
 \tag{A37b.4}
\]
because $L$ carries degree four into degree six, where
$H=0$. Thus
\[
                  (L^2)^{-1}=A^{-2}\Lambda^2
       :\bigwedge^8H^1\longrightarrow\bigwedge^4H^1.
 \tag{A37b.5}
\]
The displayed map is a left inverse by (A37b.4).
Its source and target are free of the same rank $495$;
equivalently their determinants multiply to a unit,
so it is the full inverse. No factor $5$ is inverted.

All subsequent arguments of C44-A33a now apply
unchanged. The cohomology and Hodge freeness of
C44-A22a, integral Poincare duality, and the above
integral ambient inverse give the self-dual actual
variable lattice. The same proof on Betti $p$-adic
cohomology gives its self-dual integral complement.
Mazur's identification uses the verified torsion-free
crystalline cohomology and Hodge spectral-sequence
degeneration; it has no further exclusion of $5$ in
this scope. In addition the fibre dimension is four,
strictly less than $5$. The orthogonal-crystal sources
in C44-A33 require $p$ odd, which $5$ satisfies.
The Hodge sum $956\chi(L)-50$ and every rational
correspondence calculation are unaffected.

For the promised actual characteristic-five family,
repeat C44-A35 with $p=5$ and $5\nmid t$.
The elliptic equation is smooth, $3$ remains invertible,
the order-three fixed points are etale, and the norm
polarization has $\chi(L)=(9t)^6$ prime to $5$.
The fixed-curve derivative ranks, the surface normal
rank-drop counts, and the uniform equivariant jet
surjections are valid in this characteristic. Ordinary
quadratic nodes and their simultaneous tame reflections
require $p\ne2$, not $p\ge7$.
The positive critical-point degree $7\int h^6$ can
be divisible by $5$ without making the transverse
critical scheme nonreduced; its reducedness comes
from the independently imposed jet conditions.
The exact character ranks, the fourth moments,
Katz's characteristic-zero monodromy criteria on
$\ell$-adic coefficients, and the finite-quotient
Chebotarev selection therefore remain valid.

The actual $j=0$ product, the surface, the norm bundle,
and the two $\mu_3$-weight sections still lift over
the same unramified Witt model. The just-proved
lattice calculation supplies the sole missing
characteristic-five input to C44-A34a. Since $5$
is inert in $\mathbf Q(\sqrt{-3})$ and the selected
degree is odd, that criterion proves precisely the
sign and dimension conclusions in C44-A35.
\end{proof}

\paragraph{Scope of the extensions.}
The split-prime argument retains the actual semilinear
descent datum. It does not obtain a rational Hasse
invariant merely by extending scalars to split a
quadratic field. Its evenness assertion (A37.3) is
an isocrystal determinant calculation on two
$\mathbf Q_p$-defined eigenspaces. The integral
improvement (A37b.3)--(A37b.5) is an exterior-algebra
identity over $\mathbf Z[1/6]$, not an appeal to an
unproved integral hard-Lefschetz theorem.

These are parts of the same developing quaternion
comparison and its actual-axis scope. They are not
separately staged progress-clock events. Earlier
checked files and their hashes have been retained.

% END work/scan_w/split_prime_quaternion_sign.tex

% BEGIN work/axis/c3_twisted_split_axes.tex
% C44_POS /root/axis_independent; GPT-6 Astra Pro; effort not exposed.
% C44-A38 complete candidate. Depends on checked A35 and candidate A37.
\section{Twisted cyclic axes at split primes}
\label{sec:c44-a38}

The split-prime version C44-A37 of the quaternion sign criterion
has actual geometric instances. Here the Frobenius reversal is
provided by a quadratic twist of the ambient abelian variety.
The descent of the defining sections differs from that in
C44-A35, and the distinction is included explicitly in the proof.

\begin{theorem}[C44-A38: split-prime axes with a positive three-dimensional span]
Let $p\ge7$ satisfy $p\equiv1\pmod3$. There are smooth
$(L,L)$ axes in abelian sixfolds over finite fields of odd
degree over $\mathbf F_p$ which are bases of genuine
Lefschetz pencils and for which the following hold.
\begin{enumerate}
\item Three explicit surface classes have primitive
projections spanning a positive definite space $K$ of
dimension three.
\item On $K^\perp$ inside the middle Gysin kernel the
generalized potential Tate multiplicity is exactly two.
These two dimensions form one imaginary-quadratic line,
with finite-field Frobenius reversing the field action.
The full variable potential Tate multiplicity is five.
\item The entire algebraic axis space is positive
definite, and has rational dimension either three or
five. No algebraicity assertion for the two possible
additional Tate directions is made.
\end{enumerate}
For every sufficiently large chosen polarization of the
form below, a Zariski-dense set of closed parameters of
odd degree has these properties. The sign is unconditional.
\end{theorem}

\begin{proof}
Let $E:y^2=x^3+1$ over $\mathbf F_p$, and fix
$\zeta\in\mathbf F_p$ of order three. Over an algebraic
closure put $\widetilde X=E^6$ and define
\[
      g=(\zeta,\zeta,\zeta,\zeta^{-1},\zeta^{-1},\zeta^{-1}),
      \qquad
      \tau(x_1,x_2,x_3,x_4,x_5,x_6)
                  =(x_4,x_5,x_6,x_1,x_2,x_3).
 \tag{A38.1}
\]
Then $\tau^2=1$ and $\tau g\tau=g^{-1}$. Let $X$ be
the quadratic twist of $E^6$ by $\tau$. It is an abelian
variety; equivalently it is
$\operatorname{Res}_{\mathbf F_{p^2}/\mathbf F_p}
 (E_{\mathbf F_{p^2}}^3)$. Under the geometric identification
with $E^6$, its finite-field Frobenius is the product of
$\tau$ with the ordinary coordinate Frobenius. Thus at
every odd-degree field it reverses $g$.

Choose a finite field $k/\mathbf F_p$ of odd degree and
$P\in E(k)$ with $\zeta P\ne P$. Such a choice exists:
the fixed subgroup has three geometric points, whereas
the numbers of points over odd-degree extensions tend
to infinity. We may enlarge $k$ by an odd-degree extension
once and fix this datum. The twist remains nontrivial as
a quadratic descent over $k$.

Let $S_0$ be the following translate of an abelian surface
in the geometric coordinates:
\[
 S_0=\{(x_1,P,0,x_4,P,0):x_1,x_4\in E\},
 \qquad S_i=g^iS_0\quad(i=0,1,2),
 \qquad Z=S_0\amalg S_1\amalg S_2.
 \tag{A38.2}
\]
These three surfaces are disjoint: their second-coordinate
offsets are the three distinct points $P,\zeta P,\zeta^2P$.
They avoid every point fixed by $g$, for the same reason.
The automorphism $\tau$ preserves $S_0$ and takes $S_i$
to $S_{-i}$. Coordinate Frobenius over $k$ fixes the
offsets. Hence $Z$ descends to $X/k$, and the Frobenius
permutation of its three geometric components is
\[
                             i\longmapsto-i.
 \tag{A38.3}
\]

Take $M=\boxtimes_{i=1}^6\mathcal O_E(3[0])$, its
$C_3$-norm $H$, and $L=H^t$ for a sufficiently large
integer $t$ prime to $p$. The product bundle and its
norm linearization are $\tau$-equivariant and descend
to the twist. On every $g$-fixed point the fibre character
of $H$ is trivial. Write
\[
          a=9t,\qquad h=a\sum_{i=1}^6\eta_i,
          \qquad c=\chi(X,L)=a^6,
          \qquad p\nmid c.
 \tag{A38.4}
\]
Over the quadratic extension of $k$ consider
\[
        \mathcal P=H^0(\widetilde X,I_ZL)_\chi
                         \oplus H^0(\widetilde X,I_ZL)_{\bar\chi}.
 \tag{A38.5}
\]
The twisted descent is $s\mapsto\tau^*\sigma(s)$, where
$\sigma$ is the coefficient automorphism of that quadratic
extension. Since $\sigma(\zeta)=\zeta$ and
$\tau g\tau=g^{-1}$, it interchanges the two section
weight spaces. Their dimensions agree by $\tau$.
Consequently (A38.5) descends to an affine parameter
space over $k$. An odd-degree rational parameter has a
pair of conjugate eigenmembers; their two-dimensional
span, pencil and common base are defined over its field.
After geometric base change the parameter space is still
the full direct sum in (A38.5). In particular no
$\tau$-invariance is imposed on an individual geometric
axis. Such an extra invariance would be a different
construction and is not used here.

\paragraph{Smoothness and the original pencil.}
The fixed locus now consists of $3^6=729$ points,
and the normal tangent character ranks there are $3,3$.
All these points are forced into the base, since the
bundle fibre character is trivial and both section
characters are nontrivial. Uniform equivariant jets
of $I_ZL$, obtained by the bounded-length Serre argument
of A35, give nonzero derivatives in both character
spaces at every one of these finitely many points.
They are independent. Along the free smooth surface
union $Z$, the two normal derivatives give an arbitrary
$2\times4$ matrix. Its rank-drop locus has codimension
three, greater than the dimension two of $Z$. Away
from $Z$ and the fixed points, the action is free and
the usual two-value plus derivative-rank-drop condition
has codimension seven in dimension six. Thus the
generic base is smooth. Each of the two eigenmembers
is smooth as well: at a free point singularity imposes
seven first-jet conditions, and at a fixed point its
available three-dimensional derivative is nonzero.

Every member is smooth along the base. Off the base,
the seven first-jet conditions for a singular point
of a specified member have a seven-dimensional
incidence parameter consisting of the point and the
member. Second jets exclude degenerate quadratic
parts. Two points in different free $g$-orbits in one
member would impose fourteen conditions on thirteen
variables, so are absent generically. Two points in
the same orbit would force two distinct members to
have zero first jet at one point, and hence both
generators to have zero first jet at a base point,
contradicting smoothness of the base. The character
ratio $\chi/\bar\chi$ has order three, so the
nonexceptional member parameters have free orbits.
The exceptional eigenmembers are smooth, as checked.
Finally $c_6(J^1L)=7h^6$ gives a nonempty finite
critical scheme of positive degree, exactly as in
A35. We have therefore obtained a nonempty descended
open $\Omega\subset\mathcal P$ of genuine original
Lefschetz pencils. Twisting introduced no simultaneous
node symmetry on any individual geometric pencil.

\paragraph{The positive orbit of surfaces.}
The three surfaces are parallel translates with
varying coordinates $1,4$. Let $Q_0=\{1,4\}$ and
$R_0=\{2,3,5,6\}$. The class
\[
 \alpha=\frac1{12a^2}\left(
    3\eta_1\eta_4
    -\sum_{i\in Q_0,\ j\in R_0}\eta_i\eta_j
    +\sum_{i<j,\ i,j\in R_0}\eta_i\eta_j\right)
 \tag{A38.6}
\]
satisfies $h^2\alpha=[S_i]_X$ for each $i$.
Put $w_i=[S_i]_B-\kappa^*\alpha$ and
$K=\operatorname{span}_{\mathbf Q}\{w_0,w_1,w_2\}$.
The same squarefree calculation as A35, together with
the disjointness of the supports, gives its exact Gram
matrix
\[
                     6a^2 I_3-\frac1{4a^2}J_3.
 \tag{A38.7}
\]
Its two eigenvalues on the augmentation space are
$6a^2$, and its invariant eigenvalue is
$6a^2-3/(4a^2)>0$. In particular all three classes
are independent. The group $C_3$ acts regularly on
$K$, and normalized Frobenius acts by (A38.3), with
eigenvalues $1,1,-1$. Both the pairing and its rational
Gram projector commute with these actions. The invariant
part $K_0$ has dimension one; the imaginary part $K_E$
has dimension two and is positive, $g$-stable and
Frobenius-stable.

\paragraph{Exact variable ranks.}
The character multiplicities of $H^1(\widetilde X)$
are $(0,6,6)$. The traces of $g$ on exterior degrees
zero through four are the first five coefficients of
$(1-z+z^2)^6$, namely
\[
                            1,-6,21,-50,90.
 \tag{A38.8}
\]
Since the fixed locus of the base is all 729 fixed
points, the tame fixed-point formula gives its
alternating trace at $g$ equal to 729. The ordinary
Euler characteristic is $3600c$. Subtracting lower
and upper ambient cohomology and the ambient middle
part as in A35 gives
\[
      \operatorname{Tr}(g\mid V)=729-90-2(1+6+21+50)=483,
      \qquad \dim V=3600c-165.
 \tag{A38.9}
\]
Thus
\[
       \dim V_0=1200c+267,
       \qquad \dim V_\chi=\dim V_{\bar\chi}=1200c-216.
 \tag{A38.10}
\]
On $W=V\cap K^\perp$ the invariant rank is
$1200c+266$, which is even, and each character rank
is $1200c-217$, which is odd. All exceed eight.

\paragraph{Moments, reflections and geometric monodromy.}
Over the quadratic constant-field extension the action
and all three centres are individually defined.
Remove $Z$ and the 729 fixed points, and divide by
$C_3$. The resulting smooth connected sixfold $Y$
has a connected finite etale $C_3$-torsor above it.
Four-point evaluation of (A38.5) onto the descended
rank-two bundle is onto, by the same twelve-point
equivariant evaluation used in A35. Thus the centred
zero indicators have Bernoulli parameter $Q^{-2}$
and are four-wise independent.

For the invariant character their exact normalized
fourth moment tends to three. For either nontrivial
character, with weights $w_y\in\mu_3$, the exact
formula is
\[
 \mathbf E\left|\sum_yw_y(I_y-Q^{-2})\right|^4
     =2M_Q^2\sigma^4+\left|\sum_yw_y^2\right|^2\sigma^4
                         +M_Q(\mu_4-3\sigma^4),
 \tag{A38.11}
\]
with the notation of (A35.18). The connected torsor
gives $M_Q=Q^6+O(Q^{11/2})$ and
$\sum_yw_y^2=O(Q^{11/2})$, so division by $Q^8$ gives limit two.

The trace comparison is the same cancellation
$N_j(B)-Q^{-2}N_j(X)$ as in (A35.21). The removed
fixed points contribute only $O(1)$. There is one
changed surface contribution which must be retained:
over the quadratic split field the torsor
$Z\to Z/C_3=S_0$ is split, because its three
components are individually rational. Its character
weighted count is $\#S_0(\mathbf F_Q)$ for
\emph{each} character, and its removal subtracts
\[
                 (1-Q^{-2})\#S_0(\mathbf F_Q)
                                      =Q^2+O(Q^{3/2}).
 \tag{A38.12}
\]
This removes precisely the one weight-four constant
line in each character of the regular representation
$K$. Hence the mock traces approximate those of
$W_0,W_\chi$ with uniform error $O(Q^{3/2})$.
The nonnegative moment argument of A35, including
the simultaneous-phase limsup for invariant tensors,
therefore gives exact fourth monodromy moments
three and two respectively.

Prescribing one free orbit of three ordinary double
points away from $Z$ and the fixed locus gives a
regular one-parameter smoothing. Independent second
jets and the same rank counts give no other
singularities. Its three nonzero, pairwise orthogonal
vanishing cycles are the regular representation of
$C_3$, and are orthogonal to the ambient image and
to $K$. One inertia element is a reflection on
$W_0$ and an order-two pseudoreflection on
$W_\chi$. The checked Katz criteria therefore give
\[
             G_0=O(W_0),\qquad G_\chi\supset SL(W_\chi).
 \tag{A38.13}
\]
All moment, purity and reflection inputs are exactly
those checked in A35; the three-centre subtraction
(A38.12) is the only new trace term.

\paragraph{Odd arithmetic selection.}
An odd-degree Frobenius on the descended family
reverses the two character spaces. Its matrix on
their sum is
$\bigl(\begin{smallmatrix}0&A\\A^{-t}&0\end{smallmatrix}\bigr)$.
Their common rank is odd by (A38.10) after removing
$K$, so $+1,-1$ are forced. The block construction
(A35.25) shows that every anti-$SL$ coset contains
an open avoiding any other prescribed finite set
of roots of unity, with those two eigenvalues
simple. The even-dimensional invariant block has
an open avoiding the same finite set of roots of
unity in its determinant-one component. Multiplying an odd arithmetic element
by the geometric node reflection if needed selects
that component without changing oddness.

The projections of the connected joint geometric
group dominate $SO(W_0)$ and contain $SL(W_\chi)$.
The two polynomial exclusions therefore meet in
the irreducible connected coset; no product
monodromy assertion is required. Bound orders by
$2(3600c-165)^2$, take one common exponent, and use
the determinant test on $W_0$ and the coefficient
of $(T-1)^2$ for the $N$th power on the two-character
block, exactly as in (A35.26)--(A35.27). A finite
congruence quotient retaining the quadratic
constant-field cover and determinant character
then gives, by Chebotarev, a Zariski-dense set of
odd-degree closed points with
\[
                   \tau(W_0)=0,
             \qquad\tau(W_\chi\oplus W_{\bar\chi})=2.
 \tag{A38.14}
\]
Adding the Frobenius-finite regular space $K$
gives $\tau(V)=5$.

\paragraph{Integral lifting and positivity.}
Lift the same elliptic equation, its cube-root
automorphism, and the triple-swap involution over
Witt vectors. The twist is defined by finite etale
quadratic descent, so is an actual abelian scheme.
Smoothness of $E$ lifts $P$ to a point over the
Witt vectors of $k$. The product point bundles,
norm linearization and all three surface centres
then lift with their $g,\tau$ actions. For large
$t$, relative Serre vanishing and tame averaging
lift the two character section modules. Their
exchange descent gives a free descended module,
whose reduction is (A38.5); thus the actual
selected pair also lifts. The resulting proper
base and eigenmembers are smooth over the trait
because their special fibres are smooth.

The actual block
$eV=(1-(1+g+g^2)/3)V$ has its extending skew
operator $u=g-g^2$, with $u^2=-3$. Its known
positive subspace is $K_E$. Its residual
potential Tate multiplicity is two by (A38.14).
Here $p$ splits in $\mathbf Q(\sqrt{-3})$,
but the split-prime criterion C44-A37 applies:
absolute Frobenius preserves the two quadratic
eigenspaces and the linear reversing Frobenius
identifies them, making the positive Newton
sum even. The ordinary Hilbert-symbol packet
calculation then supplies the same sign conclusion.
The integral hypotheses hold by $p\ge7$ and
$p\nmid c$ in (A38.4). Consequently the whole
algebraic space in $eV$ is positive.

The invariant algebraic space is exactly the
positive line $K_0$, since its residual potential
Tate part is zero. Rational Gram inversion
identifies any other invariant algebraic class
with a rational multiple of that line. Orthogonal
summation proves positivity on $A_B$. The residual
two-dimensional algebraic image is either zero
or the entire imaginary-quadratic line by the
actual $u,F_s$ argument of A34a, unchanged in
A37. This gives dimensions three or five as
claimed, without a Tate conjecture.
\end{proof}

\paragraph{Programme and source scope.}
This is the split-prime geometric extension of the
same quaternion advance as A34--A35, not a separate
staged-clock event. All general external inputs
have the exact primary locators recorded in A35:
the tame Deligne--Lusztig Euler formula, Weil-II
compact-support weights, Katz's fourth-moment
criteria and SGA 7's nonzero node cycles. The
section descent, exact centre Gram matrix,
rank changes and all new boundary terms are
proved here. Its sign conclusion uses A37 only
in the explicit split-prime scope just stated.

For JAC1964 this supplies further actual QH9-axis
instances of H9-T7, cote~44 Theorem~5 scan~004,
with the primitive restriction in scans~008--009.
There is no assertion for every fixed axis and
no scan~049 product-form conclusion. The separate
ambient abelian term in a full H9-T7 height
statement remains conditional on the unrefereed
032/Milne 2026 preprints; the present axis sign
is unconditional.

% END work/axis/c3_twisted_split_axes.tex

% BEGIN work/scan_w/characteristic_three_eigenblock.tex
% C44_POS /root/scan_w. A39: the integral order-four block in characteristic three.
\section{The order-four variable block in characteristic three}
\label{sec:c44-a39}

\begin{proposition}[C44-A39: an integral eigenspace without the ambient inverse]
Let $X/\mathbf F_{3^f}$ be an abelian sixfold, with $f$ odd,
let $L$ be ample, and let
$\kappa:B\hookrightarrow X$ be a smooth intersection of two
members of $|L|$. Let $g$ be an order-four geometric
automorphism of $(X,L,B)$ such that
\[
             g^2=[-1]_X,\qquad FgF^{-1}=g^{-1}
                 \quad\hbox{on cohomology},
 \tag{A39.1}
\]
where $F$ is normalized geometric Frobenius.
Suppose the given data admit an actual smooth embedded
mixed-characteristic lift, including $g$, the ample bundle,
and the relation $g^2=[-1]$. Set
\[
 V=\ker\bigl(\kappa_*:H^4(B)(2)\longrightarrow H^8(X)(4)\bigr),
 \qquad e=\frac{1-g^2}{2},\qquad u=g|eH^4(B).
 \tag{A39.2}
\]
The harmless Tate twists here record normalized weights;
the pairing on $V$ is the actual intersection form
$q(x,y)=\int_Bxy$.

Then
\[
                         eH^4(B)=eV
 \tag{A39.3}
\]
in every realization. This block has a self-dual integral
crystalline lattice at $3$, a self-dual Betti
$\mathbf Z_3$-lattice, and the Hodge--Frobenius
elementary-divisor comparison required in C44-A34.
The action $u$ preserves the crystalline lattice and
satisfies
\[
                    u^2=-1,\qquad u^\dagger=-u,\qquad Fu=-uF.
 \tag{A39.4}
\]
All characteristic polynomials of $F$ on this block
are one common rational polynomial, including the
crystalline realization. Its eigenvalues are pure
of normalized weight zero.

No condition $3\nmid\chi(X,L)$ and no integral inverse
of the ambient map $L_h^2:H^4(X)\to H^8(X)$ is needed
for these assertions.
\end{proposition}

\begin{proof}
We first identify the actual integral subspace.
The automorphism $[-1]$ acts by $(-1)^j$ on $H^j(X)$,
because the cohomology of an abelian variety is the
exterior algebra on $H^1$. In particular, $e$ kills
both $H^4(X)$ and $H^8(X)$. Equivariance of Gysin gives
\[
                        \kappa_*e=e\kappa_*=0.
 \tag{A39.5}
\]
This proves (A39.3) directly. The same calculation is
valid for the actual generic lift and its Betti
realization.

This is also a primitive block for $h_B=\kappa^*c_1(L)$.
Indeed weak Lefschetz identifies $H^2(B)$ with
$\kappa^*H^2(X)$. If $x\in eH^4(B)$ and $y\in H^2(B)$,
the class $h_By$ is ambient, so
$\int_B xh_By=0$ by (A39.5) and the projection formula.
Poincare duality implies $h_Bx=0$.
Thus Hodge--Riemann applies to the actual primitive
pairing on the Betti block, with positive sign in
codimension two.

The Hodge modules of the embedded lift are free,
commute with base change, and have a degenerating
Hodge--de Rham spectral sequence by C44-A22a.
That proposition has no restriction on the residue
characteristic or on the ramification of the trait:
its exterior resolution uses divided powers.
In particular, in every degree
\[
                  \dim H^j_{\rm dR}(B)=b_j(B).
 \tag{A39.6}
\]
Rational crystalline comparison identifies
$\operatorname{rank}_W H^j_{\rm cris}(B/W)$ with
$b_j(B)$, where $W=W(\overline{\mathbf F}_3)$.
If $t_j$ is the number of nonzero torsion invariant
factors in that crystalline cohomology, derived
base change over the DVR gives
\[
 \dim H^j_{\rm dR}(B)=b_j(B)+t_j+t_{j+1}.
 \tag{A39.7}
\]
Consequently all these crystalline modules are
torsion free. Integral Poincare duality makes
\[
                         M=H^4_{\rm cris}(B/W)
 \tag{A39.8}
\]
a self-dual cup-product lattice.

Since $2$ is a unit in $W$, $e$ is an integral
idempotent on $M$. The automorphism $g$ is an isometry,
and $g^2$ is its own inverse, so $e^\dagger=e$.
It follows that
\[
                         M=eM\perp(1-e)M
 \tag{A39.9}
\]
is an integral orthogonal decomposition into free
self-dual lattices. This uses the full middle lattice
and bypasses the ambient Lefschetz map altogether.
The lattice $eM$ is stable under $g$. On it $g^2=-1$
and $g^\dagger=g^{-1}=-g$, proving the first two
identities of (A39.4).

For the Betti lattice, integral weak Lefschetz
identifies $H_3(B_{\mathbf C},\mathbf Z)$ with
the free group $H_3(X_{\mathbf C},\mathbf Z)$.
The universal coefficient theorem therefore makes
$H^4(B_{\mathbf C},\mathbf Z)$ torsion free, and
integral Poincare duality makes its middle
intersection form unimodular. The same integral
self-adjoint projector $e$ over $\mathbf Z_3$
splits off a self-dual lattice in $eH^4_{\rm B}$.
Smooth proper comparison identifies the special
$\ell$-adic block with this actual Betti block,
including its form and $g$-action, for every
$\ell\ne3$.

We spell out the small-characteristic Hodge input.
Mazur's elementary-divisor theorem assumes
torsion-free crystalline cohomology and degeneration
of the Hodge spectral sequence. Under those hypotheses
it identifies the Frobenius Hodge filtration with the
geometric Hodge filtration. There is no additional
assumption $\dim B<p$. That inequality is a sufficient
condition for a separate degeneration theorem and
is unnecessary after (A39.6)--(A39.7) and C44-A22a.
All the comparisons are natural for $g$.
Taking the integral direct summand $e$ in the free
Hodge modules and in (A39.9) therefore gives the
same elementary-divisor statement on $eM$ itself.

For normalized absolute Frobenius $\Phi=3^{-2}\Phi_{\rm cris}$,
its elementary exponents on $eM$ are $j-2$, with
multiplicity $h^{j,4-j}(eH^4_{\rm B})$. If
\[
 H_e=h^{3,1}(eH^4_{\rm B})+
                         2h^{4,0}(eH^4_{\rm B}),
 \tag{A39.10}
\]
then Hodge--Riemann gives
\[
 b_-(q_{e,\rm B})=2h^{3,1}(eH^4_{\rm B}),
 \qquad
                 \epsilon_\infty(q_{e,\rm B})=(-1)^{H_e}.
 \tag{A39.11}
\]
Here the Hodge numbers are those of this block.
No parity calculation for the entire variable space
is substituted for (A39.10). The rational span of
$eM$ is $\Phi$-stable; the normalized operator need
not preserve the lattice.

Naturality of absolute crystalline Frobenius gives
$\Phi u=u\Phi$. The linear normalized finite-field
endomorphism $F$ commutes with $\Phi$, while (A39.1)
gives $Fu=-uF$. These are distinct Frobenius
operations. The quadratic field $\mathbf Q(i)$
is unramified inert at $3$, and $f$ is odd, exactly
as required in C44-A34.

Finally the characteristic polynomials are common
and rational without supposing that $F$ commutes
with $g$. It does commute with $e$. Use a common
algebraic Kunneth projector $\pi_4$ for $B$, provided
by Katz--Messing, Theorem~2(1). The trace of
$eF^n\pi_4$ is the rational intersection number of
the associated algebraic correspondence with the
diagonal, independent of the chosen realization.
It is precisely $\operatorname{Tr}(F^n|eH^4(B)(2))$,
with the common rational normalization included.
Newton identities yield a single rational
characteristic polynomial for all $n$, in all
$\ell$-adic theories and crystalline cohomology.
Purity follows from middle cohomology. This argument
needs no ambient inverse even as a projector on
cycle representatives: applying $e$ already places
them in the common variable subspace by (A39.5).
\end{proof}

\begin{corollary}[C44-A39a: the characteristic-three sign criterion]
In C44-A39, let $K\subset eV$ be a positive
nondegenerate rational span of actual algebraic
cycle classes, stable under $u$ and $F$; allow $K=0$.
If the generalized potential Tate multiplicity of
$F$ on $eV\cap K^\perp$ is at most two, then the
intersection form on the entire rational algebraic
subspace of $eV$ is positive definite, unconditionally.

For the complementary block define the rational
self-adjoint projectors
\[
 e_+=\frac{(1+g^2)(1+g)}4,\qquad
 e_-=\frac{(1+g^2)(1-g)}4.
 \tag{A39a.1}
\]
They are geometrically orthogonal, extend in the
given lift, and $e_++e_-+e=1$.
If each $e_\pm V$ satisfies C44-A19 after removing
a known positive span, then the whole fixed-axis
form on $A_B$ is positive definite.
The possible remaining Tate directions are not
asserted to be algebraic.
\end{corollary}

\begin{proof}
All comparison and lattice hypotheses of the general
theorem C44-A34 have just been proved, directly at
$p=3$. We use that theorem, not the earlier
order-three corollary C44-A34a.
We must still supply its full cycle space $T$ whenever
the residual algebraic image is nonzero.

Let $R$ denote the potential Tate part on $eV\cap K^\perp$.
Its pairing is nondegenerate: a root of unity pairs
only with its reciprocal, so the generalized
root-of-unity part is orthogonal to its complement.
The two $u$-eigenspaces are paired perfectly with
one another. Consequently $\dim R$ is even and
equals zero or two. If its algebraic cycle image
vanishes, Gram projection onto $K$ identifies the
whole algebraic block with $K$.

Otherwise work first in one $\ell$-adic realization.
The semisimple part $F_s$, a rational polynomial
in the actual algebraic endomorphism $F$, interchanges
the two isotropic $u$-eigenlines of $R$ and is an
isometry. Its matrix in a paired basis is
\[
                       \begin{pmatrix}0&a\\a^{-1}&0\end{pmatrix}.
 \tag{A39a.2}
\]
Hence $F_s^2=1$, and $u,F_s$ generate
$(-1,1)=M_2(\mathbf Q_\ell)$ on $R$.
The algebraic image is stable under these
correspondences, so a nonzero image is all of $R$.
Choose a rational linear combination $v$ of
actual cycles with $a_0=q(v,v)\ne0$.
Then $v,uv$ have rational Gram matrix
\[
                             \operatorname{diag}(a_0,a_0).
 \tag{A39a.3}
\]

Apply $e$ and the rational Gram projector off $K$
to these cycle representatives. They lie in the
same indicated subspace in every realization.
Their nonzero rational Gram determinant proves
independence at every prime and in crystalline
cohomology. The common Frobenius polynomial then
shows that these cycles together with $K$ span
the entire potential Tate space in every realization.
Their normalized absolute crystalline classes are
fixed by $\Phi$. Thus they provide exactly the
actual rational space $T$ required in C44-A34;
no algebraicity of a larger central Newton summand
has been assumed.

That theorem makes the negative index of $q_T$
divisible by four. Its positive subspace $K$ has
codimension two, so the index is zero. Every other
algebraic class has rational coefficients in this
basis by pairing and rational Gram inversion.
This proves positivity without a Tate conjecture
or an identification of numerical and homological
equivalence.

The three displayed projectors are rational
polynomials in $g$, are self-adjoint and mutually
orthogonal, and commute with $F$ by (A39.1).
The given lift supplies their actual Hodge actions.
The rational one-direction result C44-A19 therefore
applies on $e_\pm V$ under its stated hypotheses.
It requires no integral ambient splitting at $3$.
The orthogonal direct sum proves the final claim.
\end{proof}

\paragraph{Primary scope and geometric application.}
The all-characteristic Hodge calculation is exactly
C44-A22a. The crystalline torsion and duality argument
is the full-middle part of C44-A33a, before its
separate ambient splitting. Its primary base-change
and integral-duality locators are Stacks,
Section~60.24, Remarks~60.24.9 and 60.24.12--13,
and Tang, \emph{Syntomic cycle classes and prismatic
Poincare duality}, \emph{Compositio Mathematica}~160
(2024), Theorem~6.4 and Remark~6.6.
The exact Mazur hypothesis is restated by Ogus,
\emph{Frobenius and the Hodge spectral sequence},
\emph{Advances in Mathematics}~162 (2001), author
p.~1, equation~(0.0.1) and the preceding paragraph:
\url{https://math.berkeley.edu/~ogus/preprints/1975a.pdf}.
The following paragraph uses $\dim X<p$ only to
obtain degeneration automatically. These passages
were read directly. The original cited book locator
is Berthelot--Ogus, Theorem~8.26; its entire
foundational proof is not claimed as a new reading.
The common-projector comparison uses Katz--Messing,
\emph{Inventiones}~23 (1974), Theorem~1 and its proof,
pp.~74--75, and Theorem~2(1) and its proof,
pp.~76--77, in the precise previously checked scope.

The proposition isolates the actual small-characteristic
input for an order-four axis family. It does not
assert that every axis has such an automorphism.
An actual family, its genuine Lefschetz pencils,
and the required residual multiplicities must be
checked separately. This is a continuation of the
same quaternion comparison, not a separately staged
progress-clock event.

% END work/scan_w/characteristic_three_eigenblock.tex

% BEGIN work/arithmetic/c4_characteristic_three_axes.tex
% C44_POS /root/fermat_aj. Candidate C44-A40; whole-proof audit required.
\section{Actual characteristic-three axes with a quaternion-controlled block}
\label{sec:c44-a40}

The integral input here is a summand of the full middle cohomology.
It is not an integral splitting off of the ambient middle image.
This distinction permits characteristic three, where an ambient
Lefschetz inverse used in earlier constructions need not be integral.

\begin{theorem}[C44-A40: free cyclic axes in characteristic three]
Let $E:y^2=x^3-x$ over $\mathbf F_3$ and $X=E^6$.
For every sufficiently large integer $t$ prime to three, there
are smooth $(L,L)$ axes $B\subset X$ over finite fields of odd
degree over $\mathbf F_3$, with $c_1(L)=16t\sum_{j=1}^6\eta_j$,
having the following properties.
\begin{enumerate}
\item Their original pencils are genuine Lefschetz pencils.
The geometric order-four automorphism $g=(i,i,i,i,i,i)$
acts freely on $B$ and is reversed by finite-field Frobenius.
\item Four explicit parallel abelian surfaces have primitive
projections spanning a positive definite space $K$ of dimension four.
On $W=V\cap K^\perp$, where $V$ is the middle Gysin kernel,
the generalized potential Tate multiplicities are zero in the
trivial character, one in the sign character, and two in the
imaginary-quadratic character pair. The total potential Tate
multiplicity of $V$ is seven.
\item The actual rational algebraic axis space is positive
definite and has dimension between four and seven. Its possible
additional imaginary-quadratic contribution has dimension zero
or two. No algebraicity of the remaining potential Tate classes
is asserted.
\end{enumerate}
These conclusions hold on a Zariski-dense set of closed
parameters of odd degree in the family constructed below.
The axis positivity is unconditional.
\end{theorem}

\begin{proof}
\leavevmode\par
\paragraph{The action, the centres, and their descent.}
The action of $\mu_4$ on $E$ is given over $\mathbf F_3$ by
\[
             z\cdot(x,y)=(z^2x,zy),\qquad z^4=1.
 \tag{A40.1}
\]
It preserves the equation. Its diagonal action on $X$ has
geometric generator $g=i$, with $g^2=[-1]_X$.
Write $\chi(g)=i$. The tangent character is $\chi$ on all
six coordinates. Frobenius over every odd-degree field
conjugates $g$ to $g^{-1}$.

Take $k=\mathbf F_{27}=\mathbf F_3(\theta)$, where
$\theta^3-\theta-1=0$. The polynomial has no root in
$\mathbf F_3$, hence is irreducible. The point
\[
                         P=(\theta,1)\in E(k)
 \tag{A40.2}
\]
lies on $E$ and is not in $E[2]$, since its $y$-coordinate
is nonzero. Put
\[
 S_j=E_1\times E_2\times\{i^jP\}_3\times\{0\}_4
                  \times\{0\}_5\times\{0\}_6,
 \qquad Z=\coprod_{j=0}^3 S_j.
 \tag{A40.3}
\]
These surfaces are disjoint. Their third coordinates are
$(\theta,1),(-\theta,i),(\theta,-1),(-\theta,-i)$.
They avoid the entire nonfree locus
\[
                        F=X^{g^2}=E[2]^6,
 \tag{A40.4}
\]
which is finite and reduced. The union $Z$ descends to $k$.
Indeed the $k$-Frobenius fixes $P$, sends $i$ to $-i$,
and therefore permutes its components by $j\mapsto-j$.
After every odd-degree extension of $k$, that permutation
is unchanged.

Let $M=\boxtimes_{j=1}^6\mathcal O_E(4[0])$ and give
\[
 H=\bigotimes_{j=0}^3(g^j)^*M,\qquad L=H^t,
 \qquad a=16t,\qquad h=a\sum_j\eta_j,
 \qquad c=\chi(X,L)=a^6
 \tag{A40.5}
\]
the norm linearization. Reindexing the four factors gives
its descent under $g\mapsto g^{-1}$. At every point with
nontrivial stabilizer the stabilizer acts trivially on
the fibre of $H$: its action only permutes the identical
one-dimensional factors in each stabilizer orbit. This
also applies to points fixed by $g^2$ but not by $g$.

Over $k_2=\mathbf F_{27^2}$ form
\[
       \mathcal P=H^0(X,I_ZL)_1\oplus H^0(X,I_ZL)_\chi.
 \tag{A40.6}
\]
Both summands descend separately to $k$. In explicit
semilinear terms, if $g^*s=is$, then
$g^*\sigma(s)=\sigma((g^{-1})^*s)=\sigma(-is)=i\sigma(s)$.
The invariant summand is preserved as well. Equivalently,
these are the weights zero and one of the $\mu_4$-action,
which are defined over $k$. Thus $\mathcal P$ is an affine
space over $k$, with no constraint relating its two geometric
section coordinates. A parameter gives two actual sections
$s_0,s_1$, their pencil, and their common base.

\paragraph{Smooth bases and genuine original pencils.}
Choose $t$ sufficiently large. The uniform Serre-vanishing
jet argument used in C44-A35 applies to $I_ZH^t$: it supplies
second jets at bounded sets of points away from $Z$, normal
first jets along $Z$, and simultaneous values on at least
sixteen geometric points. The latter number covers four
points in the free quotient and all their inverse images.
Taking the prescribed character summands is exact because
four is invertible. Hence equivariant jets with these
characters have precisely the allowed surjectivity.

At every point of $F$, the section $s_1$ vanishes, while
$s_0$ can be chosen nonzero simultaneously at all these
finitely many points. Thus the common base avoids $F$.
At a $g$-fixed point all six tangent directions have
character $\chi$, so the available differential of $s_1$
has rank-six target. At a point fixed only by $g^2$, its
stabilizer acts by $-1$ on all six tangent directions
and on the section character, giving the same available
target; its derivative at the other point of that orbit
is determined equivariantly. A general $s_1$ therefore
has nonzero derivative at every point of $F$.

Along $Z$ the group action is free. The two normal
derivatives form an arbitrary $2\times4$ matrix; rank
at most one has codimension three, exceeding $\dim Z=2$.
Away from $Z\cup F$, singularity of the common zero locus
requires two zero values and rank drop of a $2\times6$
matrix, of total codimension seven in dimension six.
These incidences are avoided by a general parameter.
Each eigenmember is also smooth: the normal derivative
failure along $Z$ has codimension four, the failure off
$Z\cup F$ imposes seven first-jet conditions, and the
points of $F$ were already handled. This proves a
nonempty open of smooth bases with smooth eigenmembers.

Every pencil member is smooth along the base because
its two defining gradients are independent there.
A noneigenmember avoids $F$, since $s_0$ is nonzero there.
Off the base a singular point of a specified member
imposes seven independent first-jet conditions. The six
point variables and the member variable give a finite
critical locus for a general pair. Second jets exclude
degenerate quadratic parts; characteristic three is odd,
so the remaining points are ordinary double points.
Two points in distinct free orbits in the same member
impose fourteen conditions on thirteen variables and
are absent generically. Two points of the same orbit
in one noneigenmember would give zero first jets for
two distinct members at one point. Their generators
would then both have zero first jet, contradicting
smoothness of the base. Here the character ratio
$\chi/1$ has order four, so every noneigenmember has
a free orbit in the pencil. The two eigenmembers
are smooth as already proved.

There is a singular member: the reduced rank-drop
scheme of the two sections in $J^1L$ has degree
\[
              \int_Xc_6(J^1L)=7\int_Xh^6>0,
 \tag{A40.7}
\]
since $c(J^1L)=(1+h)^7$ on the abelian sixfold.
It has no contribution along the smooth base. We have
therefore obtained a nonempty descended open
$\Omega\subset\mathcal P$ of genuine original
Lefschetz pencils. The bases are free $C_4$-varieties.

\paragraph{The exact positive centre space.}
Let $Q_0=\{1,2\}$ and $R_0=\{3,4,5,6\}$. The class
\[
 \alpha=\frac1{12a^2}\left(
  3\eta_1\eta_2-\sum_{u\in Q_0,\ v\in R_0}\eta_u\eta_v
                 +\sum_{u<v,\ u,v\in R_0}\eta_u\eta_v\right)
 \tag{A40.8}
\]
satisfies $h^2\alpha=[S_j]_X$ in every Weil cohomology
for every $j$, by multiplication in the squarefree
algebra $\eta_j^2=0$ and translation invariance of point classes.
For $\kappa:B\hookrightarrow X$ put
\[
       w_j=[S_j]_B-\kappa^*\alpha,\qquad
       K=\operatorname{span}_{\mathbf Q}\{w_0,w_1,w_2,w_3\}.
 \tag{A40.9}
\]
These are actual rational cycles in the Gysin kernel.
The normal sequence gives
$[S_j]_B^2=\int_{S_j}3h^2=6a^2$.
Distinct surfaces are disjoint, while
$\int_X\alpha[S_j]_X=1/(4a^2)$. Hence their Gram matrix is
\[
                       G=6a^2I_4-\frac1{4a^2}J_4.
 \tag{A40.10}
\]
Its eigenvalues are $6a^2$ on the augmentation space
and $6a^2-a^{-2}$ on the invariant line, all positive.
Thus $K$ is four-dimensional and carries the regular
$C_4$-representation. Its rational Gram projector is
an algebraic correspondence. It commutes with $g$ and
normalized Frobenius, which acts by $w_j\mapsto w_{-j}$.
In particular $K$ has potential Tate multiplicity four.

\paragraph{The character ranks.}
Work over split coefficient fields and the quadratic
constant-field extension, and index characters by
$0,1,2,3$ modulo four. The character multiplicities
in $H^j(X)$, for $0\le j\le4$, are
\[
\begin{array}{c|rrrr}
 j&0&1&2&3\\\hline
 0&1&0&0&0\\
 1&0&6&0&6\\
 2&36&0&30&0\\
 3&0&110&0&110\\
 4&255&0&240&0
\end{array}
 \tag{A40.11}
\]
These are the coefficients of
$(1+\chi z)^6(1+\chi^{-1}z)^6$ modulo $\chi^4-1$.
The ordinary Euler characteristic of $B$ is
$5\int_Xh^6=3600c$. The tame fixed-point Euler
formula gives zero at $g,g^2,g^3$, since all those
fixed loci on $B$ are empty. Thus every Euler-character
multiplicity is $900c$. The lower alternating
multiplicities are $(37,-116,30,-116)$.
Weak Lefschetz and duality supply the same upper
multiplicities. Subtracting these twice and then the
ambient degree-four row yields
\[
 \begin{aligned}
 \dim V_0&=900c-329,& \dim V_2&=900c-300,\\
 \dim V_1&=900c+232,& \dim V_3&=900c+232.
 \end{aligned}
 \tag{A40.12}
\]
After subtracting the regular centre space, the ranks are
\[
 m_0=900c-330\ \hbox{even},\qquad
 m_2=900c-301\ \hbox{odd},\qquad
 n=900c+231\ \hbox{odd}
 \tag{A40.13}
\]
for $W_0,W_2,W_1$ respectively; $W_3=W_1^\vee$.
All these ranks exceed eight. The full rank is
$D=\dim V=3600c-165$.

\paragraph{Fourth moments and geometric monodromy.}
Over $k_2$ put $U=X\setminus(Z\cup F)$ and $Y=U/C_4$.
The map $U\to Y$ is a geometrically connected finite
etale torsor, and $Y$ is smooth of dimension six.
The pair of section characters descends to a rank-two
bundle. Evaluation at any four distinct rational
points of $Y$ is onto: after base change its inverse
image has at most sixteen points, covered by the
uniform evaluation above, and equivariant averaging
gives the prescribed targets. This includes nonsplit
torsor fibres.

For a field with $Q$ elements containing $k_2$, a
uniformly random parameter in $\mathcal P(\mathbf F_Q)$
therefore gives four-wise independent zero indicators
$I_y$ of probability $v=Q^{-2}$. Set $X_y=I_y-v$ and
\[
 T_j=\sum_{y\in Y(\mathbf F_Q)}\chi^j(h_y)X_y,
       \qquad j=0,1,2,
 \tag{A40.14}
\]
where $h_y$ is torsor Frobenius, with the convention
matching the chosen character trace. Let
$M_Q=\#Y(\mathbf F_Q)$ and $A_Q=\sum_y\chi^2(h_y)$.
The compact-support weight estimate gives
\[
            M_Q=Q^6+O(Q^{11/2}),\qquad A_Q=O(Q^{11/2}).
 \tag{A40.15}
\]
For the second bound the nontrivial quadratic character
of the connected geometric torsor has no top
compact-support invariants.

Put $\sigma^2=v(1-v)$ and
$\mu_4=v(1-v)^4+(1-v)v^4$. The real weights in
$T_0,T_2$ square to one. Four-wise independence and
the two-pair/four-equal index expansion give
\[
 \begin{split}
 \mathbf E(T_j^4)&=3M_Q(M_Q-1)\sigma^4+M_Q\mu_4
                                      \qquad(j=0,2),\\
 \mathbf E(|T_1|^4)&=2M_Q^2\sigma^4+|A_Q|^2\sigma^4
                             +M_Q(\mu_4-3\sigma^4).
 \end{split}
 \tag{A40.16}
\]
Consequently their fourth moments divided by $Q^8$
tend to $3,3,2$ respectively.

These are moments of the residual cohomology. Indeed,
the equivariant trace formula and weak Lefschetz give
\[
 \#\operatorname{Fix}(g^bF_Q;B)
  -Q^{-2}\#\operatorname{Fix}(g^bF_Q;X)
    =\operatorname{Tr}(g^bF_Q\mid V^{\rm un})+O(Q^{3/2}).
 \tag{A40.17}
\]
The upper ambient cohomology cancels by Gysin, and its
middle part cancels by rational hard Lefschetz on $X$;
the remaining terms have weight at most three.
Removing $Z$ removes $(1-Q^{-2})$ times its point
counts. Its top cohomology is the regular representation,
so in every character it removes exactly the leading
$Q^2$ contribution of $K_j$, with error $O(Q^{3/2})$.
The finite locus $F$ contributes only $O(Q^{-2})$,
since $B$ avoids it. Thus on the smooth parameter open
\[
                    T_j=\operatorname{Tr}(F_Q\mid W_j^{\rm un})
                                      +O(Q^{3/2}).
 \tag{A40.18}
\]
The error is uniform. Purity bounds the normalized
traces. Katz's limiting-upper-moment lemmas therefore
give $M_4(W_0)\le3$, $M_4(W_2)\le3$, and
$M_4(W_1)\le2$. Removing nonsmooth parameters can
only decrease the nonnegative fourth-moment sums.
The independent orthogonal contraction tensors give
the reverse inequalities three for both real blocks;
the two general-linear contractions give the reverse
inequality two for $W_1$. These are equalities.

To supply reflections, prescribe an ordinary quadratic
singularity of the base at one free point off $Z\cup F$.
The independent second jets impose the same singularity
at its four-point orbit. The remaining jet incidences
avoid all other singularities. Varying its transverse
constant term gives a regular simultaneous smoothing.
This is a degeneration of bases, separate from the
one-node condition on each original pencil member.
Its four nonzero, mutually orthogonal vanishing cycles
have squares $\pm2$ and form the regular representation.
The local inertia acts as the product of their four
reflections. The cycles are perpendicular to ambient
classes and to $K$, which are fixed by this inertia.
Thus inertia is a reflection in each real block and
an order-two pseudoreflection on $W_1$.
The moment criteria now yield
\[
           G_0=O(W_0),\qquad G_2=O(W_2),\qquad
                         G_1\supset SL(W_1).
 \tag{A40.19}
\]
Restricting to $\Omega$ does not change these geometric
images, since a dense open of a normal connected base
surjects on its etale fundamental group. This uses
precisely the geometric argument and primary inputs
already isolated in C44-A35, now with four-node orbits
and the additional real character.

\paragraph{Simultaneous selection in an odd arithmetic coset.}
Use normalized Frobenius and let $\mathcal H$ be the
joint geometric monodromy closure of the three blocks.
Its connected component projects onto $SO(W_0)$ and
$SO(W_2)$, and its $W_1$-projection contains $SL(W_1)$.
Choose an element of the actual arithmetic image with
odd constant-field degree. Multiplying by the inertia
element just constructed, if necessary, makes its
$W_0$ determinant $+1$. This may also change its
$W_2$ determinant, which is harmless. Denote the result
by $F_*$; the coset $F_*\mathcal H^0$ is irreducible.

On $W_0$ its projection contains an even-dimensional
special orthogonal coset with determinant $+1$, which
has elements avoiding any prescribed finite set of
roots of unity. On the odd-dimensional $W_2$, every
orthogonal element has the forced eigenvalue equal
to its determinant, either $1$ or $-1$. In each
determinant component a reciprocal-pair torus makes
this the only root of unity, and makes it simple.

On $W_1\oplus W_3$, odd Frobenius reverses the character
spaces and, in dual bases, has the form
\[
 F_E=\begin{pmatrix}0&A\\A^{-t}&0\end{pmatrix},\qquad
             F_E^2|W_1=AA^{-t}.
 \tag{A40.20}
\]
The latter matrix has reciprocal spectrum and
determinant one. Since $n$ is odd, it has eigenvalue
one; hence $F_E$ has both eigenvalues $1$ and $-1$.
There need be no others from a prescribed finite root
set. At fixed determinant of $A$, choose
\[
 A=\operatorname{diag}\left(
 b,\begin{pmatrix}0&t_1\\1&0\end{pmatrix},\ldots,
   \begin{pmatrix}0&t_{(n-1)/2}\\1&0\end{pmatrix}\right).
 \tag{A40.21}
\]
The eigenvalues of $AA^{-t}$ are $1,t_1,t_1^{-1},\ldots$.
Choose the $t_j$ generically and use $b$ to adjust the
determinant. The $SL(W_1)$ projection supplies these
choices in every odd arithmetic coset.

Take $N=\operatorname{lcm}(1,\ldots,2D^2)$. The full
normalized Frobenius polynomial on $V$ is rational and
has degree $D$. Every root of unity in it has order
at most $2D^2$ by C44-A12, hence is an $N$th root.
The three desired conditions are the nonvanishing of
\[
 \det(F^N-I\mid W_0),\qquad
 [(T-1)^1]\det(TI-F^N\mid W_2),\qquad
 [(T-1)^2]\det(TI-F^N\mid W_1\oplus W_3).
 \tag{A40.22}
\]
Here the bracket means the coefficient in the
expansion at $T=1$. The forced multiplicities are
respectively zero, at least one, and at least two.
Each displayed condition is a nonempty open on the
projection of the irreducible coset $F_*\mathcal H^0$,
by the preceding explicit matrices. Their pullbacks
therefore meet. No independence of the three
monodromy projections is needed.

Zariski density of the actual geometric image gives
an actual arithmetic element in this open. A finite
congruence neighbourhood preserves these polynomial
inequalities. Include the quadratic constant-field
cover in the finite quotient. The finite-quotient
Chebotarev argument of C44-A12 then gives a closed
parameter of odd degree meeting all three conditions.
It applies after deleting any proper closed subset,
so these parameters are Zariski dense. We obtain
\[
            \tau(W_0)=0,\qquad \tau(W_2)=1,
       \qquad \tau(W_1\oplus W_3)=2,\qquad \tau(V)=7.
 \tag{A40.23}
\]
These are full generalized multiplicities. In particular
the two residual imaginary eigenvalues are simple.

\paragraph{The actual lift and the integral middle summand.}
Use the same elliptic equation over the Witt vectors.
The point $P$ lifts explicitly to $(\widetilde\theta,1)$,
where $\widetilde\theta^3-\widetilde\theta-1=0$ is its
Hensel lift; the derivative $3\widetilde\theta^2-1$ is a
unit. The action, norm bundle, four centres and their
quadratic descent all lift. Relative Serre vanishing
for $I_ZH^t$ and tame averaging lift each of the two
section weight modules, so the actual selected sections
lift with their prescribed reductions. Their proper
common base and their eigenmembers are smooth over the
trait by the Jacobian criterion and properness. A
complex generic fibre can be obtained by the finite-type
descent used in C44-A05.

The rational orthogonal idempotents are
\[
 e_0=\tfrac14(1+g+g^2+g^3),\qquad
 e_2=\tfrac14(1-g+g^2-g^3),\qquad
                       e_E=\tfrac12(1-g^2).
 \tag{A40.24}
\]
They extend over the lift and are fixed by Frobenius
conjugation, so their special blocks have common rational
Frobenius polynomials. More explicitly, their compositions
with Frobenius powers are rational algebraic correspondences;
the degree-four projector and trace argument of C44-A34a
give rational, coefficient-independent power traces,
including crystalline traces. Newton identities give
the asserted block polynomials. The cycle Gram projector
onto $K$ has the same property.

On ambient cohomology $g^2=[-1]_X$ acts by $+1$ in
every even degree. Consequently
\[
           e_EH^4(X)=e_EH^8(X)=0,
          \qquad e_EH^4(B)=e_EV.
 \tag{A40.25}
\]
The latter space is primitive; its product with the
ample class lies in the $e_E$ part of $H^6(B)$, which
is zero by dual weak Lefschetz.

The crucial integral fact is that $e_E$ is a self-adjoint
idempotent over $\mathbf Z_3$. On the full middle Betti
lattice it therefore cuts out a self-dual direct
summand. The full lattice is torsion free and unimodular:
integral weak Lefschetz gives the free degree-three
homology of the ambient abelian variety, universal
coefficients give torsion-free degree-four cohomology,
and Poincare duality gives the middle perfect pairing.
On the crystalline side, C44-A22a gives all Hodge
freeness and first-page degeneration for this actual
embedded lift in characteristic three. Crystalline
universal coefficients and the de Rham comparison
give torsion-free crystalline cohomology and its full
middle perfect pairing. Applying the integral
self-adjoint $e_E$ again gives a self-dual summand.
The elementary exponents of Frobenius on it are the
actual block Hodge numbers by Mazur's theorem and
functoriality of its filtrations. There is no condition
$\dim B<3$ in that theorem once freeness and degeneration
have been supplied. These are also the integral
hypotheses isolated in C44-A39.

In particular no integral ambient Lefschetz inverse
is invoked. Nor is the generally nonunimodular known
cycle space $K_E=e_EK$ removed before applying the
quaternion theorem. On the full block $e_EH^4(B)$ put
$u=g$. Then $u^2=-1$, $u^\dagger=-u$, and the actual
odd-degree finite-field Frobenius satisfies $Fu=-uF$.
The prime three is unramified and inert in $\mathbf Q(i)$.
All comparison and lattice hypotheses of C44-A34
therefore hold on this full block.

\paragraph{The algebraic sign, without a Tate assertion.}
The real blocks are covered by C44-A19: their known
positive subspaces $K_0,K_2$ each have dimension one,
and their residual multiplicities are zero and one.
Thus their entire actual algebraic spaces are positive.

The imaginary block contains the positive space
$K_E$ of dimension two. Its residual potential Tate
space $R$ has rank two and simple eigenvalues $1,-1$.
If its actual algebraic image is zero, only $K_E$
contributes and positivity is immediate. Otherwise
the actual operators $u$ and the semisimple part
$F_s$, projected off $K_E$ using its rational Gram
matrix, act on $R$ with
\[
           u^2=-1,\qquad F_s^2=1,\qquad F_su=-uF_s.
 \tag{A40.26}
\]
They generate the full matrix algebra on this plane.
Their action on any nonzero cycle therefore makes
the span of actual cycle classes the entire realized
plane. Its nondegenerate pairing permits a rational
combination $v$ of these cycles with $b=q(v,v)\ne0$.
The actual cycles $v,uv$ have Gram matrix
$\operatorname{diag}(b,b)$, so form a common nondegenerate
rational basis in every realization. Their Gram
determinant ensures independence at every coefficient
prime and in crystalline cohomology; the common
potential Tate rank two ensures exhaustiveness.
Any further class has rational coordinates recovered
from its rational pairings with that basis.

Thus, in the nonzero case, the full potential Tate
space in $e_EH^4(B)$ is an actual rational cycle space
of dimension four containing $K_E$ positively.
C44-A34 applies to this full block. Its negative
index is divisible by four and is at most two,
because it contains the positive two-plane $K_E$.
It must therefore be zero. This proves positivity
in both cases, and the residual algebraic dimension
is either zero or two. Orthogonal summation with
the two real blocks proves the theorem.
\end{proof}

\paragraph{Primary inputs and programme scope.}
The exact general sources retained from C44-A35 are
Deligne--Lusztig, Theorem 3.2, printed p.~119, for the
tame fixed-point Euler trace; Deligne, Weil II,
Corollary 3.3.4, p.~206, for compact-support weights;
SGA 7 II, Expose XV, Theorem 3.4(i),(iii), pp.~195--196,
for the nonzero quadratic vanishing cycles; and Katz,
\emph{Larsen's alternative, moments, and the monodromy of
Lefschetz pencils}, Theorems 1.4.2(5) and 1.4.6,
pp.~531--535, and Lemmas 2.1.4--2.1.6, pp.~540--541,
\url{https://web.math.princeton.edu/~nmk/LAMMLF.pdf}.
The additional small-characteristic check is stated
explicitly in Ogus, \emph{Frobenius and the Hodge Spectral
Sequence}, first author page, equation (0.0.1) and the
following two paragraphs,
\url{https://math.berkeley.edu/~ogus/preprints/1975a.pdf}:
torsion freeness and first-page degeneration are the
Mazur hypotheses; a dimension bound is a separate
sufficient condition for degeneration, not an additional
requirement after C44-A22a has proved it.

This constructs actual positive JAC1964 QH9-axis forms
in characteristic three, within H9-T7 / cote~44
Theorem~5, scan~004 and the primitive qualification in
scans~008--009. It supplies no assertion for an arbitrary
fixed axis, no algebraic representative for every
remaining Tate vector, and no scan~049 product-form
statement. The separate ambient abelian term in the
full height statement remains explicitly conditional on
the unrefereed 032/Milne2026 preprints. The arithmetic Fermat AJ gate
C44-F07 is unchanged.

% END work/arithmetic/c4_characteristic_three_axes.tex

% BEGIN work/scan_w/dyadic_quaternion_sign.tex
% C44_POS /root/scan_w. A41--A41b: dyadic even-lattice comparison.
\section{The dyadic quaternion comparison with an even lattice}
\label{sec:c44-a41}

At the prime two we retain the ordinary Hasse invariant of a
symmetric bilinear form $q$. An even unimodular lattice means
that $q$ is perfect on the lattice and $q(x,x)$ belongs to twice
the coefficient ring for every lattice vector. The associated
integral quadratic form is
\[
                              Q(x)=q(x,x)/2.
 \tag{A41.1}
\]
Its polar form is $q$ and is nondegenerate modulo two.
The reductive integral special orthogonal group is the
Dickson-kernel group of this quadratic form. It is not the
naive kernel of the determinant on the special fibre.

\begin{proposition}[C44-A41: the even-lattice dyadic descent calculation]
Let $W=W(\overline{\mathbf F}_2)$, $K=W[1/2]$, and let
$(D,q,\Phi)$ be an even-dimensional quadratic isocrystal.
Assume that $D$ contains an even unimodular $W$-lattice $M$.
Write
\[
 H=\sum_{a>0}a h_a,\qquad
 S=\sum_{\lambda>0}\lambda\dim_K D_\lambda,\qquad
 s=\dim_K\bigoplus_{\lambda>0}D_\lambda,
 \tag{A41.2}
\]
where the $a$ are the elementary exponents of $\Phi$ on $M$
and the $\lambda$ are the Newton slopes. Both $H,S$ are integers.

There is an orientation-compatible unramified even unimodular
rational reference form $q_*$ of dimension $\dim D$.
It has a rational orthogonal decomposition
\[
                         q_*=q_{*,0}\perp\mathbb H^{s},
 \tag{A41.3}
\]
where $\mathbb H$ denotes the bilinear hyperbolic plane.
For the actual rational fixed form
$T_0=D_0^{\Phi=1}$ one has
\[
 \det T_0=\det q_{*,0}\quad\bmod\mathbf Q_2^{\times2},
 \qquad
 \epsilon_2(T_0)=\epsilon_2(q_{*,0})(-1)^{H-S}.
 \tag{A41.4}
\]
In particular, if $s=2m$ is even, then
\[
 \det T_0=\det q_*\quad\bmod\mathbf Q_2^{\times2},
 \qquad
 \epsilon_2(T_0)=\epsilon_2(q_*)(-1)^{m+H-S}.
 \tag{A41.5}
\]
The reference is specified in the proof; it is not
automatically the rational Betti form of a lift.
\end{proposition}

\begin{proof}
We give the dyadic basic calculation explicitly, so no
odd-prime restriction in Howard--Pappas is being removed.

\paragraph{The integral reference and its orientation.}
An even unimodular quadratic lattice over $\mathbf Z_2$
has, rationally, one of the two forms
\[
 \mathbb H^m,\qquad
 \mathbb H^{m-1}\perp A,\qquad
        A=\begin{pmatrix}2&1\\1&2\end{pmatrix}.
 \tag{A41.6}
\]
Here and below matrices represent bilinear forms.
For clarity, this elementary assertion can be seen by
reducing $Q$ modulo two. A primitive isotropic vector
lifts by Hensel's lemma because its polar differential
is nonzero. If $v$ is such a lift and $q(v,w)=1$, replacing
$w$ by $w-Q(w)v$ splits an integral hyperbolic plane.
Iterate until at most an anisotropic binary form remains.
In that binary case a Gram matrix is
$\bigl(\begin{smallmatrix}2a&b\\b&2c\end{smallmatrix}\bigr)$
with $a,b,c$ odd. Its signed discriminant is the unramified
quadratic squareclass, and its quadratic form is a unit
multiple of the norm of the unramified quadratic extension.
Every unit is a norm in that extension, so it is rationally
isometric to $A$. This proves (A41.6).
The two possibilities have distinct determinant squareclasses,
whose ratio is $-3$. Thus even unimodular rational forms
of a fixed even dimension and determinant are isometric.
Over $W$ the residue field is algebraically closed, so the
same isotropic-vector argument splits the entire lattice.

Choose an integral isometry over $W$ with a split reference.
Relative to its coefficient Frobenius, $\Phi$ is an
orthogonal matrix followed by Frobenius, and the determinant
of that matrix is $+1$ or $-1$. In the first case use the
split reference over $\mathbf Q_2$. In the second case
use the second reference in (A41.6): over the unramified
quadratic extension its last plane is hyperbolic, and
coefficient Frobenius interchanges its two isotropic lines.
This changes the orientation sign. In both cases we obtain
\[
                  \Phi=b\sigma_*,\qquad b\in SO(q_*)(K),
 \tag{A41.7}
\]
with $M$ identified with the scalar extension of the even
unimodular reference lattice and $\sigma_*$ preserving it.
This is the asserted orientation-compatible reference.
Its special orthogonal group is unramified and quasi-split.
The smooth integral group is valid also at two: for an
even-rank nonsingular quadratic form it is the kernel
of the Dickson invariant.

\paragraph{Hodge parity and the rational slope Levi.}
The hyperspecial Cartan decomposition has diagonal
cocharacters with exponents
$(a_1,-a_1,\ldots,a_m,-a_m)$ on a hyperbolic basis.
The coroot lattice of type $D_m$ consists of the integral
vectors with even coordinate sum. Consequently the
Kottwitz invariant is
\[
                           \kappa(b)=H\pmod2.
 \tag{A41.8}
\]
Changing the sign of one coordinate by the unramified
outer action preserves this parity. This calculation is
also the hyperspecial Hodge-point identity of
Rapoport--Richartz, Theorem~4.2(i); that theorem includes
residue characteristic two. In rank two one can stabilize
by an additional fixed hyperbolic plane.

Kottwitz's rational Newton-Levi theorem applies because
$SO(q_*)$ is quasi-split. It realizes the nonzero slope
spaces as hyperbolically paired rational weight spaces
and gives (A41.3). The remaining rational reference
$q_{*,0}$ is again an unramified even unimodular form:
removing hyperbolic planes from either form in (A41.6)
has that property. The general-linear factors contribute
the valuations of their determinants to the Kottwitz
invariant. The sum of those positive-slope determinants is
$S$, an integer because each determinant is an isocrystal
of integral valuation over $K$. The central zero-slope
orthogonal factor therefore has invariant
\[
                            \kappa_0=H-S\pmod2.
 \tag{A41.9}
\]
This also follows directly from the map on cocharacter
lattices: each general-linear determinant maps to its
parity in the type-$D$ quotient.

\paragraph{The basic comparison at two.}
We verify that changing $\kappa_0$ changes the ordinary
Hasse invariant by $-1$, while keeping the determinant.
For an unramified even reference of rank at least four,
split off an integral quadratic hyperbolic plane with
basis $e,f$ and $Q(xe+yf)=xy$. In its orthogonal
complement choose an integral vector $x_3$ with
$Q(x_3)=a\in\mathbf Z_2^\times$; such a vector exists
in either remaining hyperbolic or unramified norm plane.
Set
\[
             v=2^{-1}e+f,\qquad b_1=r_{x_3}r_v,
 \tag{A41.10}
\]
where $r_z$ is the reflection in the vector $z$.
The two vectors are orthogonal. Thus $b_1^2=1$,
$b_1$ is rational of determinant one, and $b_1\sigma$
has Newton slope zero. Reflection in $x_3$ is integral,
since its quadratic norm is a unit. Reflection in $v$
interchanges $e,f$ with coefficients $-2,-2^{-1}$.
Its relative elementary exponents are therefore
$1,-1,0,\ldots,0$, and (A41.8) gives $\kappa(b_1)=1$.

Take a unit $\Delta$ defining the unramified quadratic
extension of $\mathbf Q_2$, for example $\Delta=5$,
and choose $\eta^2=\Delta$ in $K$ with $\sigma(\eta)=-\eta$.
On the rational three-space spanned by $e,f,x_3$,
the actual fixed space of $b_1\sigma$ has quadratic
diagonal form
\[
           \langle-2,\ 2\Delta,\ a\Delta\rangle;
       \quad\hbox{the reference is }
           \mathbb H\perp\langle a\rangle.
 \tag{A41.11}
\]
Indeed $w=-2^{-1}e+f$ is fixed by $b_1$ whereas
$v,x_3$ change sign. The fixed basis $2w,2\eta v,\eta x_3$
gives the displayed diagonal. The determinants agree
modulo squares. Direct multiplication of Hilbert symbols,
using the reference diagonal $\langle-2,2,a\rangle$, gives
\[
       \frac{\epsilon_2(\langle-2,2\Delta,a\Delta\rangle)}
            {\epsilon_2(\langle-2,2,a\rangle)}
                       =(-2a,\Delta)_2=-1.
 \tag{A41.12}
\]
The last equality is the norm criterion for an unramified
quadratic extension: units are norms and an element is
a norm exactly when its valuation is even.
These are diagonals for $Q$; multiplying both forms by
two returns to the bilinear convention $q$. Their equal
dimension and determinant make the Hasse \emph{ratio}
unchanged. Orthogonally adding the unchanged complement
also leaves that ratio unchanged.

There are exactly two basic classes for this $SO$.
This is Kottwitz, Proposition~5.6, or equivalently its
proof via Proposition~5.4 for $GSpin$, whose derived
group is simply connected and whose spinor similitude
has valuation in $\mathbf Z$, with scalar multiplication
changing it by two. Both statements are over arbitrary
$p$-adic fields. Since the reference class and (A41.10)
have invariants zero and one, (A41.12) proves
\[
      \epsilon_2(\text{basic fixed form})
           =\epsilon_2(\text{reference})(-1)^{\kappa_0}.
 \tag{A41.13}
\]
The determinant is unchanged because the descent takes
place through $SO$, equivalently because its top exterior
form is fixed. A zero-slope factor of rank two is covered
by adding a fixed hyperbolic plane to both sides and then
cancelling it; equality of the determinants makes the
cross Hilbert symbols identical. The zero-dimensional
case has both invariant and Hasse ratio zero and one,
respectively. This proves (A41.4) in every central rank.

Finally
$\det\mathbb H^{2m}=1$ modulo squares and
\[
                   \epsilon_2(\mathbb H^{2m})
                     =(-1,-1)_2^{m(2m-1)}=(-1)^m.
 \tag{A41.14}
\]
The determinant cross symbol with $q_{*,0}$ is then one.
Combining (A41.3), (A41.4), and (A41.14) proves (A41.5).
\end{proof}

\begin{theorem}[C44-A41a: the dyadic quaternion signature congruence]
Retain the common-realization, algebraic correspondence,
purity, actual full potential-Tate cycle space $T$, and
two-Frobenius hypotheses of C44-A34, with the following
dyadic replacements.
The prime is $2$, the imaginary quadratic field
$E=\mathbf Q(\sqrt{-d})$ is unramified inert there, and
the finite-field degree $f$ is odd. The crystalline lattice
and the rational Betti $\mathbf Z_2$-lattice are
\emph{even} unimodular. The crystalline lattice is stable
under the integral skew generator $u$.

The actual Betti form need not be primitive; in particular,
the equality (A34.2) is replaced by its actual real Hasse
invariant below. Denote its negative index by $\nu_{\rm B}$,
and retain the verified
Hodge comparison for the positive normalized elementary
sum $H$. Its extending $E$-action is skew for this actual
form. Then
\[
                            \nu(T)\equiv\nu_{\rm B}+2H\pmod4.
 \tag{A41a.1}
\]
In the primitive Hodge--Riemann situation
$\nu_{\rm B}\equiv2H\pmod4$, so the negative index is
divisible by four, as in C44-A34.
More generally, if a known nondegenerate algebraic
summand $A\subset T$ has negative index $\nu(A)$ and
\[
 \nu_{\rm B}+2H\equiv\nu(A)\pmod4,\qquad
                        \dim(T\cap A^\perp)\le3,
 \tag{A41a.2}
\]
then the actual form on $T\cap A^\perp$ is positive
definite.
\end{theorem}

\begin{proof}
The rational semisimple Frobenius, repeated quaternion
packets, and the positive rational comparison form
$N^+$ on the non-Tate part $U$ are exactly the algebraic
constructions in C44-A34. They did not exclude two.
Every packet has square determinant at every finite
place and obeys the already proved trace formula
\[
 \epsilon_2\!\left(
       \operatorname{Tr}_{L/\mathbf Q_2}
              \langle a,da,-ac,-dac\rangle\right)
             =(-1,-1)_2^{[L:\mathbf Q_2]}(-d,Nc)_2.
 \tag{A41a.3}
\]
Here $c=(F_s-F_s^{-1})^2$ on its coefficient factor.
It is totally negative at the real embeddings, which
makes the norm model $N^+$ positive.

\paragraph{Matching the actual reference determinant.}
At every finite prime $\ell\ne2$ the common comparison
and the square determinant of $U$ show that
$\det T/\det q_{\rm B}$ is a square in $\mathbf Q_\ell$.
Deleted-place squareclass injectivity therefore makes
it a square in $\mathbf Q$: a nonsquare rational number
would define a quadratic field with infinitely many
nonsplit primes away from two. Thus
\[
                  \det T=\det q_{\rm B}
                          \quad\bmod\mathbf Q_2^{\times2}.
 \tag{A41a.4}
\]

At two split the coefficient factors of the non-Tate
packets according as $v_2(c)\ge0$ or $v_2(c)<0$.
The former are the actual central non-Tate fixed
space $U_0$; the latter have total dimension $4m$
for some integer $m$. Therefore $s=2m$ in C44-A41.
The actual operators $u,F_s$ act on $U_0$ and commute
with absolute Frobenius, so the same packet argument
gives $\det U_0$ a square over $\mathbf Q_2$.
Since
\[
                         T_0=T_{\mathbf Q_2}\perp U_0,
 \tag{A41a.5}
\]
(A41.5) gives $\det q_*=\det T$. By (A41a.4) the
orientation-compatible reference and the Betti form
have the same determinant. Both admit even unimodular
$\mathbf Z_2$-lattices and have the same dimension.
The classification (A41.6) consequently identifies them:
\[
                              q_*\simeq q_{\rm B,\mathbf Q_2}.
 \tag{A41a.6}
\]
This proves the reference comparison; it was not
assumed from an isometry over $W[1/2]$.

\paragraph{The extra dyadic Hasse factors.}
Write $N_0^+,N_{\ne0}^+$ for the local positive models
corresponding to the central and noncentral factors.
On $U_0$ the scalar-independent packet formula gives
$\epsilon_2(U_0)=\epsilon_2(N_0^+)$.
For the noncentral factors the exact slope accounting
of A34 gives
\[
                \sum_i m_i v_2(N_{L_i/\mathbf Q_2}c_i)=-fS,
                 \qquad \sum_i m_i[L_i:\mathbf Q_2]=m.
 \tag{A41a.7}
\]
Because $E/\mathbf Q_2$ is unramified,
$(-d,x)_2=(-1)^{v_2(x)}$. Also $(-1,-1)_2=-1$.
Thus (A41a.3), the oddness of $f$, and the square
packet determinants give
\[
                       \epsilon_2(N_{\ne0}^+)=(-1)^{m+S}.
 \tag{A41a.8}
\]
Apply (A41.5) and (A41a.6) to (A41a.5).
All cross determinants with a packet form are squares.
The result is
\[
 \begin{aligned}
 \epsilon_2(T)
   &=\epsilon_2(q_{\rm B})(-1)^{m+H-S}
                                  \epsilon_2(N_0^+)\\
   &=\epsilon_2(q_{\rm B})(-1)^H\epsilon_2(N^+).
 \end{aligned}
 \tag{A41a.9}
\]
The extra $(-1)^m$ from the hyperbolic slope pairs
is essential: it cancels the extra $(-1,-1)_2^m$
in the quaternion packets. Neither factor can be
discarded by substituting $p=2$ into the odd-prime proof.

At every other finite place the actual comparison
still gives
$\epsilon_\ell(T)=\epsilon_\ell(q_{\rm B})\epsilon_\ell(N^+)$.
Hilbert reciprocity and positivity of $N^+$ therefore give
\[
                     \epsilon_\infty(T)
                           =(-1)^H\epsilon_\infty(q_{\rm B}).
 \tag{A41a.10}
\]
The skew imaginary quadratic action makes both
$\nu(T)$ and $\nu_{\rm B}$ even; rational stability
of $T$ follows by the same cycle Gram inversion as
in A34. For an even negative index $\nu$ the real
ordinary Hasse invariant is $(-1)^{\nu/2}$.
Equation (A41a.10) is precisely (A41a.1).
Orthogonal subtraction of $A$ proves (A41a.2):
the residual negative index is divisible by four
and lies between zero and three, hence vanishes.
\end{proof}

\begin{corollary}[C44-A41b: a full-middle cyclic block on an actual axis]
Let $X/\mathbf F_{2^f}$, with $f$ odd, be an abelian
sixfold and $B\subset X$ a smooth $(L,L)$ axis.
Suppose an order-three geometric automorphism $g$
preserves these data, satisfies $FgF^{-1}=g^{-1}$,
and extends with the actual bundle and axis to a
smooth embedded mixed-characteristic lift. Put
\[
       e=1-\frac{1+g+g^2}{3},\qquad u=g-g^2,\qquad
       A=\kappa^* eH^4(X)\subset eH^4(B).
 \tag{A41b.1}
\]
Assume $A$ has a common nondegenerate rational
cycle basis in all realizations, with known
negative index $\nu(A)$. If
\[
 \nu(A)\equiv h^{1,1}(eH^2(X_{\mathbf C}))
                 -2h^{2,0}(eH^2(X_{\mathbf C}))\pmod4
 \tag{A41b.2}
\]
and the potential Tate multiplicity on $eV$ is
at most two, then the actual algebraic form on
$eV$ is positive definite.
The invariant variable block may independently
be handled by C44-A19 under its one-direction
hypothesis.
\end{corollary}

\begin{proof}
We work on the full block $eH^4(B)$ until the final
rational orthogonal subtraction. C44-A22a, the
crystalline universal coefficient argument, and
integral Poincare duality give the full self-dual
crystalline lattice. Integral weak Lefschetz and
the universal coefficient theorem give the full
self-dual Betti middle lattice. The projector $e$
is integral at two and self-adjoint, so it splits
both lattices into self-dual summands.
For any vector $x$ in either $e$-summand,
$x+gx+g^2x=0$ and the isometry of $g$ give
\[
                         q(x,x)=-2q(x,gx).
 \tag{A41b.3}
\]
Thus these summands are even unimodular, exactly
as required in C44-A41a. No integral ambient
Lefschetz inverse has been used.

The skew action has $u^2=-3$. Its field is the
unramified quadratic extension at two, it preserves
the crystalline lattice, and the two Frobenius
relations are those of A34. Free Hodge modules
and first-page degeneration from A22a give the
natural Mazur comparison on this direct summand.
The common rational Frobenius polynomials, including
the crystalline one, follow from the algebraic
Kunneth projector and trace argument already checked
in A34/A39. A rational ambient projector is available
from the abelian Lefschetz inverse; no integrality
of that inverse is asserted.

We calculate the nonprimitive Hodge correction.
The $e$-part of $H^0$ is zero, and $H^2(B)=H^2(X)$.
The Betti Lefschetz decomposition on the full block is
\[
                eH^4(B)=P^4_e\perp h\,eH^2(X).
 \tag{A41b.4}
\]
The middle intersection form is positive on the
$(2,2)$ primitive part and has negative index
$2h^{3,1}(P^4_e)$ on the primitive summand.
On the second summand it is negative on its
$(1,1)$ primitive source and positive on its
$(2,0)+(0,2)$ source. Consequently
\[
 \begin{aligned}
 \nu_{\rm B}&=2h^{3,1}(P^4_e)+h^{1,1}(eH^2(X)),\\
 H&=h^{3,1}(P^4_e)+h^{2,0}(eH^2(X))
                                      +2h^{4,0}(P^4_e).
 \end{aligned}
 \tag{A41b.5}
\]
This proves that $\nu_{\rm B}+2H$ is congruent modulo
four to the right side of (A41b.2). It does not
incorrectly apply the primitive sign to the full
middle block.

Rational hard Lefschetz gives the orthogonal decomposition
$eH^4(B)=A\perp eV$ with $A$ nondegenerate.
Its assumed cycle basis is entirely potential Tate.
If the residual algebraic image in $eV$ is zero,
the conclusion is immediate. Otherwise the at-most-two
potential Tate dimensions form an imaginary-quadratic
plane. Exactly as in A34a, $u$ and $F_s$ generate
$(-3,1)=M_2(\mathbf Q_\ell)$ there. Their actual
correspondences make any nonzero algebraic image
span the plane. After applying the common variable
projector, a rational nonisotropic vector $v$
and $uv$ have Gram matrix
$\operatorname{diag}(a,3a)$, $a\ne0$.
Together with the assumed basis of $A$, these are
common actual cycles spanning the full potential
Tate space of $eH^4(B)$ in every realization.

Apply A41a to this full space with its even
unimodular lattice still intact. Equations
(A41b.2) and (A41b.5) make the negative index after
orthogonally subtracting $A$ divisible by four.
The remaining plane has dimension two, so it is
positive. Rational Gram inversion accounts for
every further algebraic class. No Tate or numerical
versus homological equivalence conjecture is used.
Finally the invariant block is orthogonal and is
covered by A19 when its separately stated hypotheses
hold.
\end{proof}

\paragraph{Exact primary scope.}
Kottwitz, \emph{Isocrystals with additional structure},
\emph{Compositio Mathematica}~56 (1985), uses $p>0$
without an oddness restriction on pp.~201,203.
Section~4.5 on p.~214 identifies zero Newton slope
with rational descent; Propositions~5.4 and 5.6,
with their displayed proofs on pp.~216--217, give
the basic classification; Proposition~6.2 and its
proof on p.~219 give the rational Newton Levi.
These passages were read, including the actual
page image of Proposition~5.6:
\url{https://www.numdam.org/item/CM_1985__56_2_201_0.pdf}.
The explicit calculation (A41.10)--(A41.12) supplies
the dyadic Hasse normalization independently.

Rapoport--Richartz, \emph{Compositio Mathematica}~103
(1996), Section~4.1 and Theorem~4.2(i), pp.~178--180,
provide the hyperspecial relative-position identity
in the unramified quasi-split scope used here.
The statement and the whole displayed proof of
part~(i) were read:
\url{https://www.numdam.org/item/CM_1996__103_2_153_0.pdf}.
For the integral group at two, Conrad,
\emph{Reductive group schemes}, Definition~C.2.10,
Theorem~C.2.11 and its explanation, pp.~317--318,
and Proposition~C.3.10 with its displayed proof,
pp.~325--326, give the correct smooth special
orthogonal group and its type-$D$ root datum:
\url{https://math.stanford.edu/~conrad/papers/luminysga3smf.pdf}.
Those passages were read; their references and
root calculations left to the reader are not
represented as a new audit of the entire monograph.

The full-middle criterion needs an actual algebraic
ambient form and the stated residual multiplicity.
The local formula alone is not an assertion about
an arbitrary fixed axis. A concrete characteristic-two
family and its wild vanishing-cycle monodromy are
separate geometric work. This is part of the same
quaternion comparison and not a separate staged
progress-clock event.

% END work/scan_w/dyadic_quaternion_sign.tex

% BEGIN work/axis/c3_characteristic_two_axes.tex
% C44_POS /root. C44-A42, whole-proof author and independent audits complete.
\section{A characteristic-two cyclic family and its wild reflection}
\label{sec:c44-a42}

\begin{theorem}[C44-A42: a characteristic-two family]
Let $E/\mathbf F_2$ have equation $y^2+y=x^3$, and put $X=E^6$.
For every sufficiently large odd integer $t$, there is a
Zariski-dense set of closed parameters of odd degree giving
smooth complete intersections $B=(L,L)\subset X$, with
$c_1(L)=9t\sum_{j=1}^6\eta_j$, such that:
\begin{enumerate}
\item the original pencil is a genuine Lefschetz pencil;
\item the diagonal geometric order-three action is free on $B$;
\item the potential Tate multiplicities of the variable space
are one in its invariant block and two in its imaginary
quadratic block;
\item the actual rational algebraic axis form is positive definite.
\end{enumerate}
The dimension of that algebraic space is one of $0,1,2,3$;
its imaginary contribution has dimension zero or two.
No algebraicity of a remaining potential Tate direction is asserted.
The arithmetic sign step uses precisely the characteristic-two
comparison of C44-A41. Both the comparison and this geometric proof
have completed whole-proof author and independent audits.
\end{theorem}

\begin{proof}
\leavevmode\par
\paragraph{The family and its original pencils.}
Let $\zeta$ be a geometric primitive cube root of unity and set
$g(x_j,y_j)=(\zeta x_j,y_j)$ in every factor. This is the
geometric generator of a $\mu_3$-action defined over $\mathbf F_2$.
It acts by $\zeta$ on all six tangent directions at a fixed point.
Its fixed locus is
\[
 F=(E^g)^6,\qquad \#E^g=3,\qquad \#F=729.
 \tag{A42.1}
\]
Indeed $E^g$ consists of infinity and the two points with $x=0$;
the equation for $y$ there is separable. All these points are
already $\mathbf F_2$-rational. Since the group has prime order,
$F$ is its entire nonfree locus.

Take $M=\boxtimes_{j=1}^6\mathcal O_E(3[0])$ and
\[
 H=\bigotimes_{j=0}^2(g^j)^*M,\qquad
 L=H^t,\qquad a=9t,\qquad h=a\sum_j\eta_j,\qquad
 c=\chi(X,L)=a^6.
 \tag{A42.2}
\]
The norm linearization has trivial fibre character at every
point of $F$. It descends under Frobenius, which reverses $g$.
Let $\chi(g)=\zeta$ and use the affine parameter space
\[
 \mathcal P=H^0(X,L)_1\oplus H^0(X,L)_\chi.
 \tag{A42.3}
\]
Each weight module descends separately to $\mathbf F_2$:
if $g^*s=\zeta s$, then
$g^*\sigma(s)=\sigma((g^{-1})^*s)=
\sigma(\zeta^{-1}s)=\zeta\sigma(s)$.
Thus a parameter is an actual pair $(s_0,s_1)$ of sections,
without a relation between the two geometric coordinates.

For $t$ sufficiently large, uniform Serre vanishing gives
simultaneous jets at every bounded set of geometric points
used below. Taking a weight summand is exact because three
is invertible. At a free orbit its jets are unrestricted
at one chosen representative, the others being determined.
In particular values at four points of the free quotient
are independent, since their inverse image has at most
twelve points. The same argument supplies second jets
and the infinitesimal deformation directions used below.

At a fixed point $s_1$ vanishes, while $s_0$ can be chosen
nonzero. Each such evaluation is onto, so the finitely many
nonvanishing conditions have nonempty intersection over
the algebraic closure. The allowed differential of $s_1$
at that point is an arbitrary covector in a space of
dimension six. Its vanishing is avoided simultaneously
at all fixed points. Consequently the common base avoids
$F$, and both eigenmembers are smooth there.

On the free locus singularity of the common base imposes
two zero values and rank at most one on a $2\times6$
matrix of differentials. This has codimension seven,
greater than the dimension six of the point variable.
Each eigenmember has the usual seven independent singular
first-jet conditions. The corresponding incidences show
that a general base and both eigenmembers are smooth.

Every member is smooth along a smooth base. A noneigenmember
also avoids $F$. Away from the base, its singularity at a
specified point imposes seven first-jet conditions.
Allowing the point and pencil parameter leaves a finite
critical scheme for a general pair. The quadratic term
in its six local variables may have nondegenerate polar
form even in characteristic two: $x_1x_2+x_3x_4+x_5x_6$
is an example. Degenerate polar forms form a proper
closed locus among the second jets. Thus all these
critical points are ordinary nondegenerate quadratic
points, and the critical scheme is reduced.

Two points from distinct free orbits lying in one singular
member impose fourteen conditions on thirteen variables.
They are absent generically. Two points from the same
orbit in a noneigenmember would give zero first jets of
two distinct members at one point, since the character
ratio $\chi/1$ has order three. Both generator first
jets would then vanish at a base point, contradicting
smoothness of the base. The eigenmembers are smooth.
Finally the critical scheme is not empty: its length is
\[
 \int_X c_6(J^1L)=7\int_Xh^6=7\cdot720c>0,
 \tag{A42.4}
\]
using $c(J^1L)=(1+h)^7$ and the rank-drop formula for
two sections of this rank-seven bundle. There is no
contribution at the smooth base. A positive integer
length remains positive when it is divisible by two;
reducedness was established by the polar Hessian, not
by reducing this integer modulo two.
We have obtained a nonempty descended open
$\Omega\subset\mathcal P$ of genuine original pencils.

\paragraph{Exact character ranks.}
The characters of $H^1(X)$ are six copies each of
$\chi$ and $\chi^{-1}$. In degrees zero through four
the multiplicities and the trace at $g$ are
\[
\begin{array}{c|rrr|r}
j&1&\chi&\chi^{-1}&\operatorname{Tr}(g)\\ \hline
0&1&0&0&1\\
1&0&6&6&-6\\
2&36&15&15&21\\
3&40&90&90&-50\\
4&225&135&135&90
\end{array}
\tag{A42.5}
\]
They follow by expanding $(1+\chi z)^6(1+\chi^{-1}z)^6$
modulo $\chi^3=1$. The Euler characteristic of $B$ is
$3600c$. The fixed-point Euler traces of $g,g^2$ on $B$
are zero because their fixed loci are empty. The lower
alternating multiplicities are $(-3,-81,-81)$.
Weak Lefschetz and duality give the same upper
multiplicities. Subtracting these twice and then
the ambient degree-four row yields
\[
 \dim V_0=1200c-219,\qquad
 \dim V_\chi=\dim V_{\chi^{-1}}=1200c+27.
 \tag{A42.6}
\]
Both displayed ranks are odd and exceed eight.
The total rank is $D=3600c-165$. There are no imposed
surface centres and no centre subtractions in this family.

\paragraph{Fourth moments.}
Over the quadratic constant-field extension put
$U=X\setminus F$, $Y=U/C_3$. This is a smooth
six-dimensional quotient of a geometrically connected
finite etale torsor. A uniformly random parameter gives
four-wise independent zero indicators $I_y$, for
$y\in Y(\mathbf F_Q)$, of probability $v=Q^{-2}$.
For $j=0,1$ put
\[
 T_j=\sum_y\chi^j(h_y)(I_y-v),\qquad
 M_Q=\#Y(\mathbf F_Q),\qquad A_Q=\sum_y\chi^2(h_y),
 \tag{A42.7}
\]
where $h_y$ is torsor Frobenius with the matching trace
convention. Geometric connectedness and compact-support
weights give $M_Q=Q^6+O(Q^{11/2})$ and
$A_Q=O(Q^{11/2})$. The nontrivial character has no
top compact-support invariant.

With $\sigma^2=v(1-v)$ and
$\mu_4=v(1-v)^4+(1-v)v^4$, the two-pair/four-equal
expansion, justified by four-wise independence, gives
\[
\begin{split}
\mathbf E(T_0^4)&=3M_Q(M_Q-1)\sigma^4+M_Q\mu_4,\\
\mathbf E(|T_1|^4)&=2M_Q^2\sigma^4+|A_Q|^2\sigma^4
                         +M_Q(\mu_4-3\sigma^4).
\end{split}
\tag{A42.8}
\]
After division by $Q^8$ their limits are three and two.
These are the moments of the variable cohomology.
The equivariant trace formula subtracts $Q^{-2}$
times the ambient point count: upper ambient terms
cancel by Gysin, and its middle part by rational
hard Lefschetz. The remaining terms have weight at
most three. The removed finite locus contributes
only $O(Q^{-2})$, because $B$ avoids it. Hence,
uniformly on the smooth open,
\[
 T_j=\operatorname{Tr}(F_Q\mid V_j^{\rm un})
                                      +O(Q^{3/2}).
 \tag{A42.9}
\]
Purity bounds the normalized traces. Removing
nonsmooth parameters only decreases the nonnegative
fourth-moment sums. Katz's limiting-upper-moment
argument therefore bounds the two geometric invariant
dimensions by three and two. The three orthogonal
and two general-linear contraction tensors supply
the reverse bounds. Their moments are exactly three
and two. All coefficient representations here are
of characteristic zero; the residue characteristic
two does not change these group-theoretic criteria.
Purity also makes the geometric monodromy groups
reductive, by the exact input recalled in Katz,
Section 2.1.1, printed p.~539 (Deligne, Weil II,
1.3.8 and 3.4.3(iii)).

\paragraph{A wild reflection, with its smoothing specified.}
We need a degeneration of the four-dimensional bases,
not a singular member of the original five-dimensional
pencil. At one free point prescribe an ordinary
quadratic singularity of the base, and hence at its
three-point orbit, with completed local hypersurface
equation
\[
             x_1x_2+x_3x_4+z^2=0.
 \tag{A42.10}
\]
All other singularities are excluded by the same
independent jet incidences. The two infinitesimal
smoothing directions are represented by $1,z$:
the local first-order deformation space is
$k[z]/(z^2)$. The parameter space maps onto them,
because both values and first jets at this free
orbit can be prescribed.

The characteristic-two ordinary-quadratic normal
form, SGA 7 II, XV, Corollary 1.3.2(ii) and
Remark 1.3.3, therefore gives a trait in the
parameter space whose local equations, after
henselian coordinates, are
\[
 x_1x_2+x_3x_4+z^2+b(t)z+c(t)=0,\qquad
 b(t)=t+O(t^2),\quad c(t)=t+O(t^2).
 \tag{A42.11}
\]
One may first use the formally smooth versal
two-parameter base and then take the curve with
these prescribed tangent directions; the supplied
independent global jets ensure the required local
versality. The total space is regular at the special
node, since $c'(0)=1$, and the generic fibre is
smooth, since $b(t)\ne0$. The three local equations
are transported into one another by $g$.

The associated projective quadric is the sum of
two hyperbolic planes and
$z^2+bzw+cw^2$. Its even Clifford centre is the
quadratic algebra
\[
                  r^2+r=c/b^2.
 \tag{A42.12}
\]
For clarity, rescaling $w$ makes the binary form
$z^2+zv+(c/b^2)v^2$; in its Clifford algebra the
product of the two basis vectors has square plus
itself equal to $c/b^2$. Adding hyperbolic planes
does not change that quadratic centre, equivalently
the two families of maximal isotropic subspaces.
Here $c/b^2=t^{-1}+O(1)$. It is not an
Artin--Schreier coboundary in $\bar k((t))$:
the leading negative pole of $r^2+r$ has even
order whenever $r$ has a pole. The quadratic
character of inertia is therefore nontrivial.

SGA 7 II, XV, Proposition 3.2.1 and Complement 3.2.2
now give a vanishing cycle of square two with this
quadratic inertia character. This is a wild reflection;
no assertion of tame monodromy in characteristic two
is made. The three cycles are nonzero and mutually
orthogonal, because they have disjoint local supports.
They are perpendicular to the ambient classes, which
are fixed by the degeneration. Their span is the
regular $C_3$-representation. The same inertia element
acts as the product of the three reflections, and
therefore as a reflection on $V_0$ and an order-two
pseudoreflection on $V_\chi$.

The moment criteria of Katz, Theorems 1.4.6 and
1.4.2(5), consequently give
\[
                 G_0=O(V_0),\qquad G_\chi\supset SL(V_\chi).
 \tag{A42.13}
\]
Restriction to the dense normal open $\Omega$
does not decrease the geometric image. The group
criterion is over a characteristic-zero coefficient
field and permits the just-produced wild reflection.

\paragraph{Odd-degree closed parameters.}
Odd arithmetic Frobenius reverses the two nontrivial
character spaces. In paired bases its imaginary block
has the form
\[
 F_E=\begin{pmatrix}0&A\\A^{-t}&0\end{pmatrix},\qquad
 F_E^2|V_\chi=AA^{-t}.
 \tag{A42.14}
\]
Its rank $n=1200c+27$ is odd. The reciprocal spectrum
of $AA^{-t}$ and its determinant one force an
eigenvalue one. Thus $F_E$ has eigenvalues $1,-1$.
At fixed determinant of $A$, the matrices
\[
 A=\operatorname{diag}\left(
 b,\begin{pmatrix}0&t_1\\1&0\end{pmatrix},\ldots,
 \begin{pmatrix}0&t_{(n-1)/2}\\1&0\end{pmatrix}\right)
 \tag{A42.15}
\]
make those two eigenvalues simple and avoid every
other member of any prescribed finite root set,
by choosing the $t_j$ generically and adjusting $b$.
The $SL$ projection supplies this open in every
odd arithmetic coset.

The invariant orthogonal block is also odd-dimensional.
In either determinant component its forced eigenvalue
is the determinant, $1$ or $-1$, and a reciprocal-pair
torus makes this the only root in a prescribed finite
set and makes it simple. Work in one translate of
the connected component of the joint geometric
monodromy group. It is irreducible and its projections
contain the connected groups just used. The two
nonempty opens therefore intersect; independence
of the projections is not needed.

The full normalized characteristic polynomial on
$V$ is rational. Its roots of unity have order at
most $2D^2$ by C44-A12. With
$N=\operatorname{lcm}(1,\ldots,2D^2)$ the two
conditions are the nonvanishing of the coefficients
of $(T-1)^1$ and $(T-1)^2$ in the respective
polynomials $\det(T-F^N)$. They are nonempty
polynomial opens on the indicated odd coset.
An actual arithmetic element meets them by Zariski
density. A finite congruence neighbourhood preserves
both inequalities. Including the quadratic
constant-field cover in that quotient and applying
finite-quotient Chebotarev gives odd-degree closed
points, densely after deletion of any proper closed
subset. At them
\[
                    \tau(V_0)=1,\qquad \tau(eV)=2,
 \qquad e=1-\tfrac13(1+g+g^2).
 \tag{A42.16}
\]
These are generalized multiplicities, so the two
imaginary eigenvalues are simple.

\paragraph{The full integral block and its known ambient part.}
The elliptic equation, the $\mu_3$ action, the norm
bundle and the two weight modules lift over Witt
vectors. Relative Serre vanishing and tame averaging
lift the actual selected pair of sections. Properness
and the Jacobian criterion give a smooth embedded lift.
The rational projectors extend in this actual lift.
Their traces against Frobenius powers are rational
algebraic correspondence intersections, so the
Katz--Messing projector argument gives common
rational block polynomials, including crystalline
cohomology. Purity and the comparison hypotheses
are the actual geometric ones.

Use the full block $eH^4(B)(2)$, not an integral
primitive complement. By C44-A22a, crystalline
universal coefficients, integral Poincare duality
and integral weak Lefschetz, its Betti and crystalline
lattices are self-dual summands of the full middle
lattices. The projector $e$ is self-adjoint and
integral at two. Their bilinear form is even:
for $x\in eM$,
$x+gx+g^2x=0$ and $g$ is an isometry, so
\[
                         q(x,x)=-2q(x,gx).
 \tag{A42.17}
\]
Thus $q(x,x)/2$ defines the required integral
quadratic lattice. Mazur's Hodge comparison applies
by the torsion freeness and first-page degeneration
already provided by C44-A22a. Put $u=g-g^2$;
then $u^2=-3$, $u^\dagger=-u$ and $Fu=-uF$.
The quadratic field $\mathbf Q(\sqrt{-3})$ is
unramified inert at two.

The ambient even cohomology of $X$ is entirely
spanned by actual rational divisor products.
Indeed $\#E(\mathbf F_2)=3$ gives Frobenius polynomial
$T^2+2$; its square is scalar after normalization.
Tate's endomorphism theorem and coordinate/graph
divisors span $H^2(E^6)$, whose products span all
even cohomology. Rational intersection-matrix
inversion supplies a common rational cycle basis.
The unconditional positivity for these Lefschetz
classes is the published result isolated with its
exact scope in C44-A28; it uses no 032 input.

The ambient $eH^4(X)$ has dimension $270$ and embeds
by restriction into $eH^4(B)$ nondegenerately, with
orthogonal complement $eV$. Its Lefschetz
decomposition has primitive degree-four part of
dimension $240$ and the $L H^2(X)_e$ part of
dimension $30$, with no degree-zero contribution.
The preceding unconditional Lefschetz-class sign
therefore gives its actual negative index exactly
thirty. All of it is potential Tate. No integral
splitting of this ambient summand is required.

\paragraph{The nonprimitive Hodge correction and the sign.}
For the Betti block let $\nu_{\rm B}$ be its negative
index and set
$H=h^{3,1}(eH^4)+2h^{4,0}(eH^4)$.
Primitive Hodge--Riemann and Lefschetz decomposition give
\[
 \nu_{\rm B}-2H
 \equiv h^{1,1}(eH^2)-2h^{2,0}(eH^2)
 \equiv -30\equiv2\pmod4.
 \tag{A42.18}
\]
Here the scalar tangent action has
$h^{1,1}(eH^2)=0$, $h^{2,0}(eH^2)=15$.
This explicitly retains the nonprimitive correction;
the primitive formula is not applied to the full block.

If the algebraic image in $eV$ is zero, its sign is
settled. Otherwise $u$ and $F_s$ generate the split
matrix algebra on its two-dimensional potential
Tate space, as in C44-A34a. A nonzero cycle image
spans it. A rational vector $v$ of nonzero norm and
$uv$ have Gram matrix $\operatorname{diag}(b,3b)$.
Together with the full ambient cycle basis they
give a common actual nondegenerate rational basis
for the entire potential Tate part $T$ of
$eH^4(B)$. The Gram determinants prove independence
in every realization; the common Frobenius
polynomial proves exhaustiveness. Their absolute
normalized crystalline classes are fixed.

The characteristic-two comparison of C44-A41 applies
to this full even unimodular block. It yields
\[
                         \nu(T)\equiv\nu_{\rm B}-2H
                                      \equiv2\pmod4.
 \tag{A42.19}
\]
Subtracting the ambient negative index thirty
gives $\nu(eA_B)\equiv0\pmod4$. This index is
at most two; it is zero. Rational Gram inversion
controls any further algebraic class, as in
C44-A34a. The imaginary algebraic contribution
is consequently either zero or a positive plane.
The invariant block has potential multiplicity
one and is positive by C44-A19. The orthogonal
sum is precisely the required positive axis form.
\end{proof}

\paragraph{Primary sources and programme scope.}
The new characteristic-two geometric input is SGA 7 II,
Expose XV, Proposition 1.3.1(ii), Corollary 1.3.2(ii)
and Remark 1.3.3, printed pp.~174--176; 2.2.5(D),
pp.~184--185; Proposition 3.1.2, p.~188; and
Proposition 3.2.1 with Complement 3.2.2,
pp.~188--189. The Artin--Schreier computation above
identifies the nontrivial character required by
that source; ordinary quadratic singularity alone
would not suffice to infer it. The moment and
group inputs are exactly Katz, \emph{Larsen's
alternative, moments, and the monodromy of Lefschetz
pencils}, Theorems 1.4.2(5), 1.4.6 and
Lemmas 2.1.4--2.1.6, printed pp.~531--535,
540--541. The tame fixed-point Euler and weight
inputs retain the exact C44-A35 locators. The unconditional
Lefschetz-class sign is Milne, \emph{Polarizations and Grothendieck's
standard conjectures} (2002), Remarks 3.7--3.8, author p.~9,
in precisely the scope recorded in C44-A28.

This is a characteristic-two instance of
the same quaternion advance for JAC1964 H9-T7 /
QH9-axis, cote 44 Theorem 5, scan 004, and its
primitive restriction on scans 008--009. It does
not claim the universal sign or a product form
for cote 16 scan 049. The separate ambient
abelian contribution to full height positivity
remains conditional on the unrefereed
032/Milne2026 preprints. No part of the present
axis argument uses them.

% END work/axis/c3_characteristic_two_axes.tex

% BEGIN work/axis/c3_characteristic_two_centres.tex
% C44_POS /root/axis_independent; GPT-6 Astra Pro; effort not exposed.
% C44-A42a scoped refinement; author proof checked, independent audit pending.
\section{Six known positive classes in the characteristic-two family}
\label{sec:c44-a42-centres}

\begin{proposition}[C44-A42a: six prescribed classes in characteristic two]
The construction of C44-A42 admits the following refinement over
$k=\mathbf F_8$. For every sufficiently large odd $t$, a
Zariski-dense set of closed parameters of odd degree gives genuine
Lefschetz-pencil axes with a known positive six-dimensional
algebraic variable space $K$. The residual potential Tate
multiplicities are one in the invariant block and two in the
imaginary block. The full variable potential Tate multiplicity
is nine, and the actual algebraic axis space is positive
definite of dimension one of $6,7,8,9$.
This is a refinement of the same dyadic quaternion argument,
not a separate progress-clock event.
\end{proposition}

\begin{proof}
Retain $E:y^2+y=x^3$, $X=E^6$, the scalar geometric
order-three action $g$, the norm polarization and the
section weights $1,\chi$ from C44-A42. Thus
$h=a\sum_i\eta_i$, $a=9t$ and $c=a^6$.
Let $\theta\in\mathbf F_8$ satisfy
$\theta^3+\theta+1=0$, and put
\[
                 P=(\theta+1,\theta^2+\theta)\in E(\mathbf F_8).
 \tag{A42a.1}
\]
Indeed both $x(P)^3$ and $y(P)^2+y(P)$ equal $\theta^2$.
Since $x(P)\ne0$ and $P\ne0_E$, the point is not fixed by $g$.

Take the following two translates of the same abelian surface:
\[
 \begin{split}
 S_{0,0}&=\{(x_1,x_2,P,0_E,0_E,0_E):x_1,x_2\in E\},\\
 S_{1,0}&=\{(x_1,x_2,0_E,P,0_E,0_E):x_1,x_2\in E\},\\
 S_{\epsilon,j}&=g^j S_{\epsilon,0}
       \quad(\epsilon\in\{0,1\},\ j\in\mathbf Z/3).
 \end{split}
 \tag{A42a.2}
\]
These six surfaces are pairwise disjoint. Within one orbit
the nonzero offset has three distinct values; between the
orbits the nonzero offset is in a different coordinate.
They avoid the fixed locus $F=(E^g)^6$. Let $Z$ be their union.
Frobenius of $\mathbf F_8$ fixes $P$ and takes $g$ to $g^{-1}$,
so it permutes $S_{\epsilon,j}$ by $j\mapsto-j$.
Thus the smooth disjoint union $Z$ descends to $k$.

Use the affine parameter space
\[
             \mathcal P_Z=H^0(X,I_ZL)_1
                              \oplus H^0(X,I_ZL)_\chi.
 \tag{A42a.3}
\]
Each weight space descends separately, exactly as in A42;
no equality of their dimensions is asserted. For $t$ large,
Serre vanishing supplies all the bounded simultaneous jets
used below, including jets away from $Z$ and first normal
jets along $Z$. Tame projection to a weight space preserves
surjectivity at any chosen representative of a free orbit.

Along $Z$ the two conormal values form a $2\times4$ matrix.
Its rank-drop locus has codimension three, exceeding the
dimension two of the point variable. A general pair therefore
has a smooth common base along $Z$. The singularity condition
for either one of its divisors has codimension four there,
again exceeding two. At the finite fixed locus the ideal
$I_Z$ is the unit ideal. Thus the invariant section can still
be nonzero at every fixed point, and the weight-one section
can still have nonzero differential there. On the complement
of $Z\cup F$ the original jet incidences of A42 are unchanged.

Consequently the genuine original-pencil argument applies:
six-variable polar Hessians are nondegenerate at critical
points, two distinct orbit conditions have the same
codimension fourteen, and two same-orbit singular points
would force two zero generator first jets at a base point.
There is no critical point on the smooth base. The positive
Porteous number remains $7\int_Xh^6$. We obtain a nonempty
descended open of genuine axes containing all six surfaces.

Write $Q_0=\{1,2\}$ and $R=\{3,4,5,6\}$ and set
\[
 \alpha=\frac1{12a^2}\left(
 3\eta_1\eta_2-
 \sum_{i\in Q_0,\,j\in R}\eta_i\eta_j+
 \sum_{i<j,\ i,j\in R}\eta_i\eta_j\right).
 \tag{A42a.4}
\]
The squarefree product-ring calculation gives
$h^2\alpha=[S_{\epsilon,j}]_X$ and
$\int_X\alpha[S_{\epsilon,j}]_X=1/(4a^2)$.
For each contained surface put
$w_{\epsilon,j}=[S_{\epsilon,j}]_B-\kappa^*\alpha$.
These are rational algebraic variable classes. Their normal
bundle in $X$ is trivial, so
$c_2(N_{S/B})=3h^2|S$ and $[S]_B^2=6a^2$.
Different surfaces have zero intersection on $B$.
Projection off the common ambient part therefore gives
\[
       \bigl(q(w_{\epsilon,j},w_{\epsilon',j'})\bigr)
                  =6a^2 I_6-\frac{1}{4a^2}J_6.
 \tag{A42a.5}
\]
Its eigenvalues are $6a^2$ with multiplicity five and
$6a^2-6/(4a^2)$ once. They are positive. In particular
these classes are independent, and
$K=\langle w_{\epsilon,j}\rangle\cong\mathbf Q[C_3]^2$.
It is stable under both $g$ and normalized Frobenius;
all its Frobenius eigenvalues are $\pm1$.

The bases still avoid $F$, so their character ranks are
exactly those in A42. Rational orthogonal subtraction of
$K$ gives a lisse residual system $W=V\cap K^\perp$ with
\[
        \dim W_0=1200c-221,\qquad
        \dim W_\chi=\dim W_{\chi^{-1}}=1200c+25.
 \tag{A42a.6}
\]
Both ranks are odd and greater than eight.

We verify the only change to the moment calculation.
Over even extensions of $k$, all six components are
defined and each three-component orbit maps to its
quotient as a split $C_3$-torsor. In particular its
character weight is one on every rational quotient
point. Put $Y=(X\setminus(F\cup Z))/C_3$ and take the
same four-wise independent zero indicators of probability
$Q^{-2}$. Their fourth moments still have limits three
and two, since $Y$ has dimension six and its nontrivial
torsor character has no top compact-support invariant.
Removing each surface orbit removes
\[
              (1-Q^{-2})\#S_{\epsilon,0}(\mathbf F_Q)
                            =Q^2+O(Q^{3/2})
 \tag{A42a.7}
\]
from every character trace. With the two orbits this
subtracts precisely the $2Q^2$ trace of $K$ in each
character. Thus the mock traces approximate the residual
traces on $W$ to the same $O(Q^{3/2})$ error as in A42.
Nonnegativity on discarded parameters and the invariant
phase-limsup argument are unchanged.

At a free point outside $Z$, the mixed jet generation
also supplies the degeneration
$x_1x_2+x_3x_4+z^2+b(t)z+c(t)$ with
$b(t)=t+O(t^2)$ and $c(t)=t+O(t^2)$.
The imposed centres remain smooth throughout its
punctured trait. Its three nonzero local vanishing
cycles are orthogonal to the fixed classes in $K$,
as well as to the ambient classes. The precise wild
Clifford-character proof of A42 therefore gives the
same reflection and pseudoreflection on $W_0,W_\chi$.
The moment criteria again give full orthogonal and
at least special-linear geometric projections.

Apply the odd arithmetic-coset and finite-root-set
selection of A42 to the two odd ranks in (A42a.6).
It gives densely many odd-degree closed parameters with
$\tau(W_0)=1$ and $\tau(eW)=2$. The six known classes
are potential Tate, so $\tau(V)=6+1+2=9$.

All the data lift over Witt vectors. In particular
smoothness of the elliptic model lifts $P$ to a point
over $W(k)$; it is not necessary to prescribe separate
Teichmuller coordinates for that point. The six
translated surfaces then lift as a disjoint union,
with the same descent permutation. The actual pair
of sections lifts by relative Serre vanishing and the
weight decomposition. Hence the extending imaginary
action and the actual full even unimodular block of
A41 are available, as are the relative classes in $K$.

For the sign, work on $eH^4(B)$, with the rational
known algebraic summand
$A'=\kappa^*eH^4(X)\perp eK$.
Its negative index is thirty, since $eK$ is a positive
four-dimensional space. The ambient Betti Hodge correction
is still $\nu_{\rm B}-2H\equiv-30\pmod4$.
If the residual algebraic image in $eW$ is zero,
its sign is settled. Otherwise the actual $u,F_s$
actions and rational Gram construction of A34a give
a common nondegenerate cycle basis for that entire
two-dimensional potential Tate plane. Together with
the known basis of $A'$ it supplies the full certificate
for A41a. That theorem makes the negative index of
the residual plane divisible by four, hence zero.
The invariant residual direction is positive by
A15/A19 after the positive two-dimensional $K_0$.
The full axis space is positive. Its invariant
dimension is two or three and its imaginary dimension
is four or six, giving exactly the stated alternatives.
\end{proof}

The source inputs are the already audited exact A42
wild-monodromy passages and A41 dyadic comparison,
with the elementary centre Gram calculation above.
This supplies a nonzero positive axis space in the
same characteristic-two family of JAC1964 QH9-axis /
H9-T7, cote~44 Theorem~5 scan~004 and primitive
scans~008--009. It does not settle the arbitrary fixed
axis or the product-form problem on cote~16 scan~049.
The separate ambient implications for full height
positivity remain conditional on the unrefereed
032/Milne~2026 preprints.

% END work/axis/c3_characteristic_two_centres.tex

% BEGIN work/axis/elliptic_power_height_corollary.tex
% C44_POS /root; GPT-6 Astra Pro; effort not exposed.
% C44-A42b. Direct height consequence of the checked axis families.
\section{Unconditional full heights for the elliptic-power families}
\label{sec:c44-a42b}

\begin{corollary}[C44-A42b: the full original-pencil height]
Let $X/\overline{\mathbf F}_p$ be an abelian $2m$-fold
geometrically isomorphic to a power of one elliptic curve.
Let $h$ define a genuine Lefschetz pencil with smooth axis $B$,
and suppose its algebraic axis form
$(-1)^{m-1}\int_B ww'$ is positive definite.
Then the signed height of JAC1964 H9-T7 on the entire group
$\mathcal H=\CH^m(Y)_{\mathrm{hom},\mathbf Q}$ is positive
semidefinite, its radical is exactly
\[
 \mathcal H_0=\operatorname{im}\bigl(
 \CH^m(\operatorname{Bl}_B X)_{\mathrm{hom},\mathbf Q}
       \longrightarrow\CH^m(Y)_{\mathbf Q}\bigr),
 \tag{A42b.1}
\]
and the induced form on $\mathcal H/\mathcal H_0$ is
positive definite. This conclusion is unconditional.

In particular it applies to all the checked cyclic families
C44-A35 (including $p=5$ by A37b), A38, A40, A42 and A42a.
It is a direct consequence of their axis signs and the
published Lefschetz-class theorem for their explicit
ambient varieties, not an additional progress-clock advance.
\end{corollary}

\begin{proof}
\leavevmode\par
\paragraph{The ambient algebraic classes.}
We record why powers of a single elliptic curve have no
algebraic classes outside their divisor-generated algebra.
Descend $E$ to a finite field $\mathbf F_q$ and fix
$\ell\ne p$. All classes and all cycles under consideration
are allowed to descend to a further finite extension.
Use the exterior algebra description
\[
 H^*(E^n,\mathbf Q_\ell)
       =\bigwedge^* H^1(E^n,\mathbf Q_\ell).
 \tag{A42b.2}
\]

First suppose $E$ is ordinary. Over a coefficient field
splitting its Frobenius polynomial, write
$H^1(E^n)=U\oplus U'$, where the eigenvalues are
$\alpha,\beta$, with $\alpha\beta=q$ and
$\dim U=\dim U'=n$. The ratio $\alpha/\beta$ is not a
root of unity: its two $p$-adic valuations are nonzero
and opposite, since the slopes are zero and one.
A degree-$2r$ exterior monomial with $a$ factors from
$U$ and $b$ from $U'$, $a+b=2r$, has normalized
Frobenius eigenvalue
\[
 q^{-r}\alpha^a\beta^b=(\alpha/\beta)^{a-r}.
 \tag{A42b.3}
\]
It is a root of unity if and only if $a=b=r$.
Every such balanced monomial is a product of $r$
balanced degree-two monomials. Thus all potential
Tate classes are generated by the potential Tate
classes in degree two. This equality descends to
$\mathbf Q_\ell$, because it holds after a faithful
coefficient-field extension.

For supersingular $E$, pass to a finite extension
on which Frobenius is scalar on $H^1(E)$.
Then every degree-two class is Tate, and products
of degree-two wedges span every even exterior power.
For the particular supersingular curves in the stated
applications, the scalar-power assertion is already
explicit in A35, A40 and A42. In either case Tate's
endomorphism theorem realizes the degree-two Tate
space by coordinate and graph divisors. Equivalently,
its divisor form follows by identifying rational
Neron--Severi classes with symmetric homomorphisms
to the dual abelian variety. Consequently all potential
Tate spaces are spanned over $\mathbf Q_\ell$ by
actual products of rational divisor classes.

This scalar-extension assertion also gives the required
\emph{rational} statement about algebraic classes.
The Poincare pairing between complementary potential
Tate spaces is perfect. Indeed Frobenius is semisimple,
and every eigenvalue pairs only with its inverse
on normalized complementary degree. Choose bases
of the two Tate spaces consisting of actual divisor
products $d_i$ and $e_j$. Their intersection matrix
$M=(\int d_i e_j)$ is rational and invertible.
For any rational algebraic class $z$, the numbers
$\int z e_j$ are rational. The coordinates of $z$
in the basis $d_i$ are therefore rational, by
multiplication by $M^{-1}$. Hence the entire rational
algebraic class space equals the divisor-generated
class space in every codimension. This applies over
$\overline{\mathbf F}_p$; the field extension used in
the proof imposes no rationality requirement on $z$.
A form of $E^n$, such as the quadratic twist in A38,
has the same geometric conclusion.

\paragraph{Primitive sign and the height quotient.}
Milne's published Lefschetz-class theorem gives
unconditionally the signed Hodge standard form on
primitive divisor-generated classes, for every
ample rational polarization and in arbitrary
characteristic. By the preceding paragraph, it
applies to the whole algebraic primitive space
\[
 A_X=\ker\bigl(L_h:\Alg^m(X)\longrightarrow
                    H^{2m+2}(X)(m+1)\bigr).
 \tag{A42b.4}
\]
In middle degree its form is exactly
$(-1)^m\int_Xzz'$, so it is positive definite.

Now apply the already checked JAC1964 H9-T7,
(H9.17)--(H9.18), in its original-pencil scope.
The specialization injectivity required there is
automatic for genuine abelian Lefschetz pencils
by the inherited JAC-P07; no finite cover is taken.
The canonical isometry is
\[
 \mathcal H/\mathcal H_0\simeq A_X\perp A_B,
 \qquad Q=(-1)^m\int_Xzz'+(-1)^{m-1}\int_Bww'.
 \tag{A42b.5}
\]
Both summands are positive definite. Thus the quotient
form is positive definite, and its pullback has
radical exactly $\mathcal H_0$, proving the claim.
The named axis theorems verify the second summand;
their ambients are geometric sixth powers of their
stated elliptic curves, including the form in A38.
\end{proof}

\paragraph{Exact inputs and programme scope.}
The divisor realization input is Tate,
\emph{Endomorphisms of abelian varieties over finite fields},
Inventiones Math.~2 (1966), Main Theorem and equation~(1),
printed p.~134. The unconditional sign input is Milne,
\emph{Polarizations and Grothendieck's standard conjectures}
(2002), Remarks~3.7--3.8, author p.~9; it is precisely
the Lefschetz-class result, rather than the conditional
Theorem~3.3 of that paper. The short exterior-algebra
and rational-Gram arguments above spell out the
specialization to the ambients used here.

For cote~44 Theorem~5, scan~004, this proves the full
signed original-pencil height conclusion on these
families, with the primitive normalization required
by scans~008--009. It uses the Chow blowup formula,
Weil cohomology and the primitive standard form in
H9-T7. No assertion about arbitrary finite covers,
all special axes, or the product form needed on
cote~16 scan~049 follows. For a general abelian
ambient, the separate use of 032 Corollary~8.4 and
Milne's 2026 report remains explicitly conditional
on those unrefereed preprints. Neither is used in
this corollary. For these explicit elliptic-power ambients,
the conclusion sharpens the conservative conditional full-height
remarks appended to A35, A38, A40, A42 and A42a.

% END work/axis/elliptic_power_height_corollary.tex

% BEGIN work/axis/formal_k_effectivity_gate.tex
% C44_POS /root/axis_independent; GPT-6 Astra Pro; effort not exposed.
% Scoped route appendix; diagnostic A31d. No clock movement.
\section{The remaining effectivity step for continuous cycle classes}
\label{sec:c44-formal-k-effectivity}

The arbitrary-ramification continuous $K$-theory theorem
does not supply the actual cycle required in route (a).
Here is an exact, elementary effectivity condition, and
an obstruction to obtaining it from ranks and numerical
positivity alone. This is a diagnostic of the remaining
step, not a new primitive lifting theorem.

\begin{lemma}[C44-A31d: finite special-fibre representatives]
Let $R$ be a complete discrete valuation ring with
uniformizer $\pi$ and finite residue field $k$, and let
$\mathcal X/R$ be smooth projective. Put
$X_n=\mathcal X\otimes_R R/(\pi^n)$.
Let $(\xi_n)\in\varprojlim_n K_0(X_n)$ be a compatible
integral system. The following conditions are equivalent.
\begin{enumerate}
\item There are vector bundles $E,F$ on $\mathcal X$
such that $[E|X_n]-[F|X_n]=\xi_n$ for every $n$.
\item There is a finite list of pairs of vector bundles
$(E^{(j)}_1,F^{(j)}_1)$ on $X_1$ such that, for
arbitrarily large $n$, there are vector bundles
$E_n,F_n$ on $X_n$ with
\[
 [E_n]-[F_n]=\xi_n,
 \qquad
 (E_n|X_1,F_n|X_1)\cong
            (E^{(j)}_1,F^{(j)}_1)
\]
for some member of that fixed list.
\end{enumerate}
The same assertion applies after multiplying the entire
system by one fixed integer. It concerns a specified
compatible system; it does not assert that this system
is the only possible lift of its special cohomology class.
\end{lemma}

\begin{proof}
The first condition gives the second with a one-element
list. For the converse, consider the tree whose vertices
at depth $n$ are isomorphism classes of pairs of vector
bundles on $X_n$ with the indicated $K_0$ difference and
special fibre in the fixed list. Connect a pair to its
restriction at depth $n-1$, and add one root above the
finite first level. Compatibility of the $\xi_n$ makes
this a tree of admissible restrictions. The hypothesis
gives vertices at arbitrarily large depths.

Every vertex has only finitely many children. Indeed,
$X_n\hookrightarrow X_{n+1}$ has square-zero ideal
$\pi^n\mathcal O/\pi^{n+1}\mathcal O\cong\mathcal O_{X_1}$.
For a vector bundle $G_n$, if one lift exists, its lifts
with identified reduction form a torsor under
\[
                  H^1(X_1,\mathcal E nd(G_1)).
\]
For completeness, trivialize $G_n$ on a finite affine
cover. Two lifts have transition matrices differing by
$1+\pi^n a_{ij}$; the transition identity is the
additive Cech cocycle condition for $\mathcal E nd(G_1)$,
and changes of local trivialization are its coboundaries.
The cover may be chosen affine with affine intersections,
since $X_1$ is separated, so it computes this coherent
cohomology. Its dimension over $k$ is finite by properness,
and its underlying set is finite because $k$ is finite.
Forgetting the identified reduction cannot increase the
number of isomorphism classes. Apply this to both bundles.

The elementary finite-branching tree lemma now gives an
infinite branch. Equivalently, among finitely many children
of a vertex with descendants at arbitrary depths, one has
descendants at arbitrary depths; choose it successively.
Choose restriction isomorphisms along the branch. They
give compatible systems of actual bundles, not merely
equalities in $K_0$.

Grothendieck existence algebraizes the two coherent
systems to coherent sheaves $E,F$ on $\mathcal X$.
They are locally free. At a point of the special fibre,
their completed stalks are free: compatible bases can be
lifted successively through the nilpotent thickenings.
Faithful flatness of completion of a noetherian local
ring then gives local freeness at that point. The
non-locally-free loci are closed and proper over $R$;
they avoid its closed fibre, so their closed images in
$\operatorname{Spec}R$ are empty. Thus $E,F$ are bundles
everywhere. Their restrictions have the required classes.
\end{proof}

\begin{lemma}[Fixed rank and class do not bound representatives]
Let $B$ be a smooth projective variety admitting a
surjective morphism $f:B\to E$ to an elliptic curve
with a rational point $o$. There are vector bundles
$G_n$ on $B$ such that
\[
       \operatorname{rk}G_n=2,
       \qquad [G_n]=2[\mathcal O_B]\quad\hbox{in }K_0(B),
       \qquad h^0(B,G_n)\ge n.
\]
In particular they have the same Chern classes and
Hilbert polynomial, but are not contained in a bounded
flat family of coherent sheaves. This holds on the
positive axes constructed in C44-A35 as well.
\end{lemma}

\begin{proof}
Set $L=\mathcal O_E(o)$ and
$G_n=f^*(L^n\oplus L^{-n})$.
The successive divisor exact sequences on $E$ give
\[
     [L^n]=[\mathcal O_E]+n[\mathcal O_o],
     \qquad
     [L^{-n}]=[\mathcal O_E]-n[\mathcal O_o].
\]
Here $\mathcal O_o$ is perfect, so the identities hold
in $K_0$ and pull back along $f$. This proves the exact
class identity. Riemann--Roch on the elliptic curve
gives $h^0(E,L^n)=n$ for $n>0$, and surjectivity of
$f$ makes the pullback on sections injective.

In any bounded flat family on the fixed projective
variety, upper semicontinuity and quasi-compactness of
the parameter scheme bound $h^0$. The displayed growth
therefore rules out boundedness, despite the fixed
rank, exact $K_0$ class and Hilbert polynomial.

For an axis $B\subset E^6$ of A35, a coordinate
projection restricts to a surjection $f:B\to E$.
Indeed, with $h=a\sum\eta_i$,
\[
          \int_B f^*[o]h^3
                    =\int_{E^6}\eta_1h^5=5!a^5>0.
\]
The projection could not be constant, and a nonconstant
proper map from the connected variety to $E$ is
surjective. Hence the example is compatible with the
proved positivity of the algebraic primitive form.
\end{proof}

\paragraph{What this does and does not obstruct.}
Adding $G_n$ to both members of any bundle-pair
representative preserves its virtual class and every
primitive intersection number. Even the separate
ranks and Hilbert polynomials remain independent of
$n$. Thus these numerical data alone do not bound
all representatives. The example does not rule out
a bounded choice: the class $2[\mathcal O_B]$ itself
has an evident bounded representative. For the
particular primitive continuous class, the unresolved
issue is precisely the existence of a finite list as
in the first lemma, or another effectivity argument.
No such list has been obtained from positivity here.

The finite residue-field hypothesis in that lemma is
substantive. A finite-dimensional deformation space
over $\overline{\mathbf F}_p$ need not have finitely
many points. Passing to that field without a uniform
finite field of definition does not preserve the
finite-branching proof. Neither lemma asserts
non-liftability, or liftability, of the A14/A18
primitive class on an arbitrary ramified model.

\paragraph{Primary source boundary.}
Grothendieck existence is used in its projective form,
Stacks Lemma~30.24.3, Tag~0885; its full displayed proof
was read, as was the effectivity and flatness formulation
in Lemma~99.5.10, Tag~08W9:
\url{https://stacks.math.columbia.edu/tag/0885},
\url{https://stacks.math.columbia.edu/tag/08W9}.
The two finite-branching and unboundedness arguments
above are proved here.

Antieau--Mathew--Morrow--Nikolaus, \emph{On the Beilinson
fiber square}, Theorem~D and Question~1.4, author p.~5,
and Definition~4.1, Examples~4.3--4.4, pp.~19--20,
\url{https://arxiv.org/pdf/2003.12541}, provide the
continuous criterion and distinguish it from actual
proper-scheme algebraization. Rationalization after
the inverse limit supplies a fixed common denominator
for one continuous class; it does not supply compatible
bundle representatives. Multiplication by a denominator
divisible by $p$ can annihilate a nonzero obstruction
in a first-order $k$-vector space, so it does not
repair the initial-support condition of A31.

The newer sources retain this distinction. Hu,
\emph{On the algebraic $K$-theory of smooth schemes
over truncated Witt vectors}, version~2 of
23~February~2026, author p.~3 and Lemma~11.3/
Proposition~11.4, p.~87,
\url{https://arxiv.org/pdf/2507.12458v2}, concerns
infinitesimal and continuous groups. Efimov,
\emph{Localizing invariants of inverse limits},
version~2 of 10~February~2025, Theorems~0.1--0.2
and equation~(0.2), pp.~3--4,
\url{https://arxiv.org/pdf/2502.04123v2}, computes
continuous $K$-theory of a larger nuclear category;
its compact objects are identified separately with
perfect complexes. Neither displayed theorem gives
the required identification with actual proper
$K_0$. These exact passages were read; no full-proof
audit of Efimov's 121-page paper is claimed.

For JAC1964 QH9-axis, cote~44 Theorem~5 scan~004
and primitive scans~008--009, the two lemmas locate
the remaining class-effectivity problem in route (a).
They neither settle H9-T7 for arbitrary axes nor
add a staged-clock movement. The independent
positive-family results are unaffected; their ambient
implications remain conditional on the unrefereed
032/Milne~2026 preprints.

% END work/axis/formal_k_effectivity_gate.tex

% BEGIN work/axis/first_ramified_support_jets.tex
% C44_POS /root/axis_independent. C44-A31e diagnostic continuation; author proof read, independent review pending.
% No assertion of primitive nonliftability at e=p-1 is made.
\section{Finite support jets at the first ramified obstruction}
\label{sec:c44-first-ramified-jets}

\paragraph{Scope.}
The coherent restriction theorem below is stronger than A22, and in
characteristic three it has a two-term de Rham version.  The final
paragraph states the remaining obstruction explicitly.  In particular,
none of these statements identifies the ordinary normal obstruction
with the first ramified crystalline period.

\begin{proposition}[C44-A31e: primitive restriction to fixed finite thickenings]
Let $S$ be a fixed smooth surface in an abelian sixfold $X$, let $H$ be
ample, and fix $m\geq1$. For all sufficiently large $t$, let
$B\subset X$ be any smooth intersection of two divisors in
$|L|$, where $L=H^t$, containing $S$. Write $J=I_{S/X}$,
$I=I_{S/B}$, $S_m=\operatorname{Spec}_B(\mathcal O_B/I^m)$,
and $F=\Omega_B^3$. Then restriction induces a surjection
\[
 \ker\bigl(H^1(B,F)\xrightarrow{\kappa_*}
                         H^3(X,\Omega_X^5)\bigr)
       \longrightarrow H^1(S_m,F|S_m).
 \tag{A31e.1}
\]
The bound is uniform in $B$. The dual local-to-global map
\[
 H^1\bigl(B,\mathcal Ext^2_B(\mathcal O_{S_m},\Omega_B^1)\bigr)
       \longrightarrow H^3(B,\Omega_B^1)
 \tag{A31e.2}
\]
is injective, and its pairing with the coherent primitive test space
in (A31e.1) is already nondegenerate on its source.
\end{proposition}

\begin{proof}
The case $m=1$ is the restriction part of A22. Suppose $m\geq2$.
For every integer $a$ there is an exact sequence on $X$
\[
 0\longrightarrow J^{m-2}L^{a-2}
   \longrightarrow (J^{m-1}L^{a-1})^{\oplus2}
   \longrightarrow J^mL^a
   \longrightarrow \kappa_*(I^mL^a|B)\longrightarrow0.
 \tag{A31e.3}
\]
Here $J^0=\mathcal O_X$. To check exactness, work first at a point
of $S$. The two equations $f_0,f_1$ of $B$ have independent images
in $J/J^2$, since $B$ is smooth. Complete them to four regular
parameters generating $J$. The associated graded ring for $J$ is
a polynomial ring over $\mathcal O_S$. Thus
\[
 (f_0,f_1)\cap J^m=(f_0,f_1)J^{m-1},
 \qquad (J^{m-1}:f_i)=J^{m-2}.
 \tag{A31e.4}
\]
The ordinary Koszul relation between $f_0,f_1$ therefore gives
(A31e.3). Outside $S$, it is their usual Koszul resolution. These
local descriptions use the actual sections of $L$, so glue to
the displayed global sequence.

Serre vanishing for the finitely many fixed coherent sheaves
$J^{m-2},J^{m-1},J^m$ gives, for $t$ sufficiently large,
\[
 \begin{split}
 H^j(B,I^mL^3)&=0 \qquad(j>0),\\
 H^j(B,I^mL^2)&\simeq H^{j+2}(X,J^{m-2}) \qquad(j>0).
 \end{split}
 \tag{A31e.5}
\]
These identifications are natural for $J^{m-2}\hookrightarrow
\mathcal O_X$ and the corresponding untwisted Koszul resolution
of $L^2|B$. The quotient $\mathcal O_X/J^{m-2}$, when nonzero,
is supported on a scheme of dimension two. Consequently
\[
 \begin{split}
 H^2(B,I^mL^2)&\xrightarrow{\sim}H^2(B,L^2|B),\\
 H^1(B,I^mL^2)&\twoheadrightarrow H^1(B,L^2|B).
 \end{split}
 \tag{A31e.6}
\]
For $m=2$, both statements are isomorphisms, since $J^0=
\mathcal O_X$.

Tensor the exact tangent--adjunction sequence
\[
 0\longrightarrow F\longrightarrow(L^2|B)^{\oplus6}
                      \longrightarrow(L^3|B)^{\oplus2}
                      \longrightarrow0
 \tag{A31e.7}
\]
by $I^m$. This remains exact because its quotient is locally
free. Equations (A31e.5)--(A31e.6) show that
$H^2(B,I^mF)\to H^2(B,F)$ is an isomorphism. Hence
$H^1(B,F)\to H^1(S_m,F|S_m)$ is surjective. Moreover the
image of $H^1(B,I^mF)$ surjects onto
$H^1(B,L^2|B)^{\oplus6}$. The latter is exactly the coherent
Gysin target after the double Koszul boundary, by A22.
Subtracting such a class kills the ambient Gysin without
changing restriction to $S_m$. This proves (A31e.1).

The power of a regular two-generator ideal is perfect of grade
two: its usual Hilbert--Burch resolution has length two for its
quotient. Thus $S_m$ is a Cohen--Macaulay surface with dualizing
sheaf
$\omega_{S_m}=\mathcal Ext^2_B(\mathcal O_{S_m},\omega_B)$.
Since $F^*=\Omega_B^1\otimes\omega_B^{-1}$, one has
\[
 \mathcal Ext^2_B(\mathcal O_{S_m},\Omega_B^1)
             =\omega_{S_m}\otimes(F|S_m)^*.
 \tag{A31e.8}
\]
Serre duality on $S_m$ and on $B$ identifies the dual of the
restriction map with (A31e.2). This is also its local-cohomology
description: the Ext sheaf maps to $\mathcal H^2_S(\Omega_B^1)$,
and the ensuing cohomology map is the local-to-global map. The
trace adjunction proves compatibility with restriction. Dualizing
(A31e.1) now gives both assertions about (A31e.2).
\end{proof}

\paragraph{Why this coherent injection does not identify the ramified cocycle.}
Take $p>2$, $R=W(k)[\pi]/(\pi^{p-1}-p)$ and $R_2=R/\pi^2$,
with the inherited canonical divided powers. Locally let $S$ be
defined by $(x,y)$ in a smooth fourfold. For displacement
$\pi(a\partial_x+b\partial_y)$ the first ramified PD homotopy
has the ordinary term
\[
 \pi\left(\frac{a}{x}\,d\log y-
                       \frac{b}{y}\,d\log x\right)
 \tag{A31e.9}
\]
and the additional closed term $\pi\xi$, where
\[
 \xi=\frac{a^p}{x^p}\,d\log y
                  -\frac{b^p}{y^p}\,d\log x.
 \tag{A31e.10}
\]
This follows by applying the convergent PD logarithm separately
in the two coordinates, as in A31a. Products of two of its
surviving coefficients vanish in $R_2$. Only powers $1$ and $p$
survive, since $v_\pi(\pi^j/j)=j-(p-1)v_p(j)$ equals one
exactly at these two values. Over the more general boundary
$p=u\pi^{p-1}$, the extra term is multiplied by
$\bar u^{-1}$.

The closed term (A31e.10) is not coordinate invariant as an actual
section of $\mathcal H^2_S(\Omega_B^1)$. For an explicit check,
take $a=1,b=0$ and change coordinates to $x'=x$, $y'=y+x$.
The normal components in the new coordinates are $a'=b'=1$.
In generalized fractions, expanding in $x/y$ and discarding
terms without both poles gives
\[
 \xi'-\xi=d\vartheta,\qquad
 \vartheta=\sum_{j=0}^{p-2}
          \frac{(-1)^j}{j+1}\,x^{j-p+1}y^{-j-1}.
 \tag{A31e.11}
\]
Indeed the $dx$ coefficient of $d\vartheta$ is
$\sum_{j=0}^{p-2}(-1)^j x^{j-p}y^{-j-1}$, while the
$dy$ coefficient is
$\sum_{j=1}^{p-1}(-1)^j x^{j-p}y^{-j-1}$. These are exactly
the two coefficients in $\xi'-\xi$. Every denominator $j+1$
is a unit in $k$. The expression and its differential are
nonzero in local cohomology. The function $\vartheta$ is killed
by $I^{p-1}$, whereas $\xi$ is killed by $I^p$.

Thus the descent of the correction can involve an actual
degree-zero-form Cech term. Its natural finite-pole complex is
\[
 \left[
 \mathcal Ext^2_B(\mathcal O_{S_{p-1}},\mathcal O_B)
       \xrightarrow{d}
 \mathcal Ext^2_B(\mathcal O_{S_p},\Omega_B^1)
 \right],
 \tag{A31e.12}
\]
with terms in form degrees zero and one. These Ext sheaves are
viewed as the corresponding pole-bounded submodules of local
cohomology. Differentiation raises total pole order by at most
one, so defines the displayed map. The computation proves
that replacing the total descent class by a raw Cech cocycle
in its second term requires justification; (A31e.1) alone does
not supply it.

\begin{proposition}[A two-term primitive jet restriction in characteristic three]
In the preceding geometric setting, suppose $k=\overline{\mathbf F}_3$
and $B$ is a genuine Lefschetz axis. For $t$ sufficiently large,
let
\[
 \mathcal J_3=
 \left[\Omega_B^3/I^3\Omega_B^3
               \xrightarrow{d}\Omega_B^4/I^2\Omega_B^4\right],
 \tag{A31e.13}
\]
in degrees three and four. This differential is well-defined
because $d(I^3)\subset I^2\Omega_B^1$. Then
\[
 \ker\left(F^3H^4_{\mathrm{dR}}(B/k)
                 \xrightarrow{\kappa_*}F^5H^8_{\mathrm{dR}}(X/k)\right)
          \twoheadrightarrow\mathbb H^4(B,\mathcal J_3).
 \tag{A31e.14}
\]
This is a statement about finite jet complexes on the special
fibre, not a comparison theorem for the canonical ramified period.
\end{proposition}

\begin{proof}
An embedded lift exists by A05. A22a applied to that lift gives
$E_1$ degeneration for $B$, as well as the usual degeneration
for the abelian variety $X$. In particular
$\mathbb H^4(B,\Omega_B^{\geq3})=F^3H^4_{\mathrm{dR}}(B/k)$
and every differential induced by $d$ on coherent Hodge
cohomology is zero.

The kernel of restriction to (A31e.13) is
\[
 \mathcal K_3=[I^3\Omega_B^3\xrightarrow{d}I^2\Omega_B^4]
 \tag{A31e.15}
\]
in the same two degrees. From (A31e.5)--(A31e.7) one has
\[
 \begin{split}
 H^2(B,I^3\Omega_B^3)&\xrightarrow{\sim}H^2(B,\Omega_B^3),\\
 H^j(B,I^2\Omega_B^4)&\xrightarrow{\sim}H^j(B,\Omega_B^4)
                                                   \quad(j=1,2).
 \end{split}
 \tag{A31e.16}
\]
The second line uses $\Omega_B^4=L^2|B$ and $m=2$, whose
leftmost term in (A31e.3) is $\mathcal O_X$. The maps
$d:H^j(I^3\Omega_B^3)\to H^j(I^2\Omega_B^4)$ for $j=1,2$
are zero: their composites with the injective maps in the
second line of (A31e.16) are the zero global Hodge differentials.
The two-term hypercohomology spectral sequence consequently
gives a commutative diagram of short exact sequences
\[
 \begin{array}{ccccccccc}
 0&\to&H^1(I^2\Omega_B^4)&\to&\mathbb H^5(\mathcal K_3)
       &\to&H^2(I^3\Omega_B^3)&\to&0\\
 &&\downarrow\wr&&\downarrow&&\downarrow\wr\\
 0&\to&H^1(\Omega_B^4)&\to&\mathbb H^5(\Omega_B^{\geq3})
       &\to&H^2(\Omega_B^3)&\to&0.
 \end{array}
 \tag{A31e.17}
\]
Thus the middle vertical map is an isomorphism. The exact
triangle of restriction proves surjectivity from the full
$F^3H^4_{\mathrm{dR}}(B/k)$ to $\mathbb H^4(\mathcal J_3)$.

We can kill its ambient Gysin while keeping this restriction
unchanged. First, $\mathbb H^4(\mathcal K_3)$ surjects onto
$H^1(I^3\Omega_B^3)$ because its next differential was just
shown to be zero. By (A31e.5)--(A31e.7), the image of this last
group in the lower graded Gysin target
$H^3(X,\Omega_X^5)$ is surjective. Second, the double-boundary
map
\[
 H^0(B,I^2\Omega_B^4)\twoheadrightarrow H^2(X,\Omega_X^6)
 \tag{A31e.18}
\]
is surjective. Indeed use (A31e.3) with $m=a=2$; its two positive
terms have no higher cohomology, and its left term is
$\mathcal O_X$. The map to the ordinary Koszul resolution has
identity on this left term, so its double boundary is precisely
the restriction of the top graded Gysin. Passing through the
cokernel of $d:H^0(I^3\Omega_B^3)\to H^0(I^2\Omega_B^4)$
does not change (A31e.18), since an exact de Rham class has zero
Gysin. These two graded surjections show
\[
 \mathbb H^4(\mathcal K_3)\twoheadrightarrow
                              F^5H^8_{\mathrm{dR}}(X/k).
 \tag{A31e.19}
\]
Subtract such a class from any lift of a requested element of
$\mathbb H^4(\mathcal J_3)$. Its restriction is unchanged and
its Gysin is zero. This proves (A31e.14).
\end{proof}

\paragraph{The precise map still not controlled.}
For the actual A14 surface in an A18 axis, take a candidate
mixed-characteristic lift with $e=p-1$ and $p=3$. Its reduction
modulo $\pi^2$ has an ordinary embedded obstruction
$o\in H^1(S,N_{S/B})\otimes(\pi)/(\pi^2)$. The ordinary
component is the injective semiregularity map of A22. The
canonical PD descent also has the Frobenius-shaped component
(A31e.10) and the coordinate-change homotopies (A31e.11).
One must identify its full supported two-term hypercohomology
class, and prove that its vanishing forces $o=0$.
Neither (A31e.2) nor (A31e.14) proves that the ordinary and
Cartier-shaped contributions cannot cancel in that group.
In particular the absence of an ordinary support lift cannot
be inferred solely from the nonzero coherent semiregularity
image, and supersingularity of $H^1(S)$ does not compute this
supported Cartier operation. No such noncancellation assertion
for the specified A14/A18 immersion has been established here.

There is also a tempting automatic initial support which does
not repair this step. Since $B_2$ is in characteristic $p$,
absolute Frobenius factors through its special fibre. The
inverse image of $S$ is an initial local complete intersection
with equations $(x^p,y^p)$ and multiplicity $p^2$. But its
coherent codimension-two semiregularity expression is
\[
 \frac{a\,d(y^p)-b\,d(x^p)}{x^py^p}=0.
 \tag{A31e.20}
\]
Thus the ordinary injectivity result for the smooth surface
does not extend to this Frobenius-thickened substitute. This
is the explicit obstruction to using the initial Frobenius
construction behind Berthelot--Ogus Lemma 3.10 as a direct
replacement for the initial smooth support in A31.

\paragraph{Inputs and programme scope.}
The proofs use the already checked A22 tangent/Koszul/Gysin
calculation and A22a degeneration, and give the new ideal-power
and coordinate-change calculations in full. The PD logarithm
is Berthelot--Ogus, \emph{$F$-isocrystals and de Rham cohomology. I},
\S3.1, printed p.~175; the boundary and initial Frobenius
construction are Lemmas 3.9--3.10, p.~179, reread here.
\url{https://gdz.sub.uni-goettingen.de/download/pdf/PPN356556735_0072/LOG_0013.pdf}.
For cote~44, Th.~5, scan~004 and the primitive warning at
scans~008--009, this is a diagnostic refinement of route (a)
in QH9-axis. It provides neither a primitive class lift nor an
all-ramification nonliftability theorem, and makes no new
staged progress-clock claim.

% END work/axis/first_ramified_support_jets.tex

% BEGIN work/axis/axis_first_ramified_noncancellation.tex
% C44_POS /root/axis_independent. Author whole proof read; independent final review pending.
\section{The first ramified step for sufficiently general prescribed-centre axes}
\label{sec:c44-a31f}

\begin{theorem}[C44-A31f: exclusion of ramification two in characteristic three]
Let $k=\overline{\mathbf F}_3$, let $S\subset X$ be the smooth
nonliftable surface in the abelian sixfold of A14, and fix an ample
line bundle $H$ on $X$. For all sufficiently large $t$ there is a
nonempty open set of pairs in $H^0(X,I_{S/X}H^t)^2$ whose axes
$B$ are smooth genuine Lefschetz-pencil axes and have the following
property. For every complete mixed-characteristic DVR $R$ with residue
field $k$ and absolute ramification index $e\leq2$, and every smooth
projective abstract lift $\mathcal B/R$ of $B$, the fixed primitive class
\[
 w=[S]_B-\kappa^*(L_h^2)^{-1}[S]_X,
                 \qquad h=c_1(H^t),
 \tag{A31f.1}
\]
does not lie in $F^2H^4_{\mathrm{dR}}(\mathcal B_K/K)$ under the
canonical crystalline comparison.

The closed A18 axes may be selected inside this open. On those axes
the nonzero primitive etale class also cannot lift algebraically on
the geometric generic fibre of any such model. A finite field
extension used to define a generic cycle does not weaken this
assertion. No conclusion is made for models with $e\geq3$.
\end{theorem}

\paragraph{Notation for the local calculation.}
For the moment let $S$ be any smooth surface in a smooth fourfold $B$
over $k$, with ideal $I$ and normal bundle $N$. Put
\[
 M_m^a=\mathcal Ext^2_B(\mathcal O_B/I^m,\Omega_B^a)
           \subset\mathcal H^2_S(\Omega_B^a),\qquad
 C=[M_2^0\xrightarrow{d}M_3^1].
 \tag{A31f.2}
\]
The two terms of $C$ have form degrees zero and one. Local cohomology
with supports has the additional shift by two. The inclusion in
(A31f.2) identifies $M_m^a$ with the annihilator of $I^m$: this follows
from $R\Gamma_I(\Omega_B^a)=\mathcal H^2_I(\Omega_B^a)[-2]$
and applying $R\mathcal Hom(\mathcal O_B/I^m,-)$.
In normal coordinates $x,y$, its fractions have total normal
denominator at most $m+1$.

Write $F_S:S\to S^{(1)}$ for relative Frobenius, and let $N^{(1)}$
denote base change of $N$ along Frobenius of $k$. This is distinguished
from applying a Frobenius operation to $H^1(S,\mathcal O_S)$.
After Frobenius pushforward, $C$ is a complex of coherent
$\mathcal O_{S^{(1)}}$-modules: cubes of elements of $I$ annihilate
both terms, and $d$ is linear for cubed functions.

\begin{lemma}[The two obstruction summands in the finite-pole quotient]
The differential in (A31f.2) is injective. Let $Q$ be its cokernel,
so that $\mathbb H^2(C)=H^1(S^{(1)},Q)$. There are natural injections
\[
 j:F_{S*}N\hookrightarrow Q,\qquad
 c:N^{(1)}\hookrightarrow Q,
 \tag{A31f.3}
\]
with disjoint images. The first is ordinary contraction with the
fundamental fraction. The second is inverse Cartier on that
contraction. The quotient by their direct sum has a finite
filtration with successive factors among
\[
 \begin{gathered}
 F_{S*}\bigl(\det N\otimes\Gamma^{m-1}N\otimes\Omega_S^1\bigr)
                                       \quad(m=1,2,3),\qquad
 F_{S*}\operatorname{Sym}^2N,\\
 F_{S*}(N\otimes\det N),\qquad
 N^{(1)}\otimes B_S^1,\qquad
 B_S^1=F_{S*}\mathcal O_S/\mathcal O_{S^{(1)}}.
 \end{gathered}
 \tag{A31f.4}
\]
Here the displayed factors may be joined by extensions; a splitting
is neither asserted nor needed.
\end{lemma}

\begin{proof}
Work locally in etale normal and tangent coordinates. The highest
normal pole of a nonzero function in $M_2^0$ has exponents $a,b\geq1$
with $a+b\leq3$. At least one, in fact both, is nonzero modulo three.
Its normal derivative is nonzero and has one higher total pole;
different fraction monomials cannot cancel that coefficient. Hence
$d$ is injective, and is strict for the pole filtration, shifted
by one.

Ordinary contraction identifies $N$ with the conormal part of
$M_1^1$ by
\[
 (a,b)\longmapsto\eta(a,b)
          =\frac{a}{x}\,d\log y-\frac{b}{y}\,d\log x.
 \tag{A31f.5}
\]
This is the natural coherent semiregularity map before taking
cohomology. No image of $dM_2^0$ can have pole one unless it is zero,
so it gives the first injection in (A31f.3).

The inverse-Cartier representative of (A31f.5) on the Frobenius
twist is the closed form
\[
 \xi(a,b)=\frac{a^3}{x^3}\,d\log y
                         -\frac{b^3}{y^3}\,d\log x.
 \tag{A31f.6}
\]
Changing a lift of $a$ or $b$ by an element of $I$ changes this
fraction by zero. Its coordinate independence in $Q$ can be checked
without assuming a split formal neighbourhood. Units multiplying
one normal coordinate do not change its local cohomology class.
For $x'=x+t y$, $y'=y$, the normal components are $a'=a+t b$, $b'=b$,
and direct expansion gives
\[
 \xi'-\xi=d\left(b^3\left(-\frac{t}{xy^2}
                                  +\frac{t^2}{2x^2y}\right)\right).
 \tag{A31f.7}
\]
The formula includes the terms in $dt$. For $x'=x+c_0y^2$, $y'=y$
one obtains
\[
                 \xi'-\xi=d\left(-\frac{b^3c_0}{xy}\right),
 \tag{A31f.8}
\]
including $dc_0$. The opposite shears give the corresponding
formulas with the coordinates interchanged. An invertible linear
normal change is generated locally by elementary shears and units.
After its removal, the quadratic terms are units times a coordinate
or the pure quadratic shears just considered. Higher terms are
invisible in (A31f.6): a pure $y^r$ correction with $r\geq3$ either
loses the $y$ pole after division by $y^3$, or has derivative
coefficient $3=0$ at $r=3$; mixed terms are absorbed into a unit.
The same holds in the other coordinate. Changes of tangent
coordinates only change coefficient lifts modulo $I$.
This proves the asserted coordinate independence. Equivalently it
is the usual natural Cartier inverse on local cohomology, now with
its exact homotopies explicitly bounded by $M_2^0$.

For clarity, the pole filtration also gives its complete module
description. Its degree-$m$ normal part before quotienting by $d$ is
$\det N\otimes\Gamma^{m-1}N\otimes N^*$, and the leading
normal differential comes from $\det N\otimes\Gamma^{m-2}N$.
For $m=1,2,3$ its cokernel is $\operatorname{Sym}^mN$.
This follows by writing the two normal derivatives; at $m=3$ the
identification $\Gamma^2N=\operatorname{Sym}^2N$ uses only
$2^{-1}$, which is available. The tangential quotient is
$\det N\otimes\Gamma^{m-1}N\otimes\Omega_S^1$. Thus
\[
 0\to F_{S*}\operatorname{Sym}^mN\to\operatorname{gr}_mQ
   \to F_{S*}(\det N\otimes\Gamma^{m-1}N\otimes\Omega_S^1)\to0.
 \tag{A31f.9}
\]
The image of $j$ is the normal term at pole one. The leading
term of $c$ is the pure-cube map
$N^{(1)}\to F_{S*}F_S^*N^{(1)}\to F_{S*}\operatorname{Sym}^3N$.
It is injective, and has zero intersection with lower poles.
For a rank-two bundle in characteristic three,
\[
 0\to F_S^*N^{(1)}\to\operatorname{Sym}^3N
                         \to N\otimes\det N\to0.
 \tag{A31f.10}
\]
In a basis $u,v$, the quotient has basis $u^2v,-uv^2$, whose
linear-change matrices are those of $N\otimes\det N$; this verifies
the last identification, including unipotent changes. Finally
projection formula identifies
$F_{S*}F_S^*N^{(1)}/N^{(1)}$ with $N^{(1)}\otimes B_S^1$.
Equations (A31f.9)--(A31f.10) prove (A31f.4).
\end{proof}

\begin{lemma}[Generality gives noncancellation on cohomology]
For the fixed smooth $S\subset X$, sufficiently high $L=H^t$, and a
nonempty open set of smooth axes $B$ through $S$, one has
\[
 H^0\bigl(S^{(1)},Q/(jF_{S*}N+cN^{(1)})\bigr)=0.
 \tag{A31f.11}
\]
For such axes the map
\[
 H^1(S,N)\longrightarrow H^1(S^{(1)},Q),
          o\longmapsto j(o)+\lambda c(o^{(1)})
 \tag{A31f.12}
\]
is injective for every $\lambda\in k$. Here $o^{(1)}$ means base
change of the class to the Frobenius twist.
\end{lemma}

\begin{proof}
Let $E=N_{S/X}$, a fixed rank-four bundle. The normal sequence and
adjunction give
\[
 0\to N\to E\to (L|S)^{\oplus2}\to0,
                   \qquad \det N=K_S\otimes L^{-2}|S.
 \tag{A31f.13}
\]
Every determinant-twisted factor in (A31f.4) embeds into a fixed
bundle tensored with $L^{-2}|S$: use $N\subset E$ and its divided
or symmetric powers. These have no global sections when $t$ is
sufficiently large. This argument is uniform in the smooth axis.

It remains to arrange
\[
 H^0(S,\operatorname{Sym}^2N)=0,\qquad
 H^0(S^{(1)},N^{(1)}\otimes B_S^1)=0.
 \tag{A31f.14}
\]
The bundle $B_S^1$ is locally free, since relative Frobenius is
finite flat on the smooth surface and the inclusion of constants
is locally a direct summand of the Frobenius monomial basis.
The section spaces in (A31f.14) lie respectively in the fixed
finite-dimensional spaces
$H^0(S,\operatorname{Sym}^2E)$ and
$H^0(S^{(1)},E^{(1)}\otimes B_S^1)$.

Finitely many normal two-planes can exclude every nonzero section
of these spaces. Explicitly, start with either fixed vector space
or a nonzero subspace left after earlier conditions. Choose a
nonzero section and a point where its value is nonzero. There is
a two-plane in the four-dimensional normal fibre for which this
value does not lie in its symmetric square, respectively in its
tensor product with the fibre of $B_S^1$. Imposing that kernel
plane cuts the remaining section subspace strictly. Repeat,
choosing distinct points, until it is zero. For the second space
the planes are Frobenius twists of normal planes; every such
geometric plane is available because $k$ is perfect.

For large $t$, sections of $I_{S/X}L$ prescribe arbitrary first
normal jets at this fixed finite set of points. This follows from
Serre vanishing for the appropriate fixed ideal sheaves. Choose
the two derivatives so that their common kernels are the chosen
planes. General sections subject to these finite conditions are
still smooth along $S$: normal rank drop for a $2$-by-$4$ matrix
has codimension three, greater than $\dim S=2$. They are smooth
off $S$ by the same imposed-centre jet argument as A06. The
additional genuine-pencil conditions of A06 are open and can be
met with these finite normal jets, since their higher-jet tests
are away from this fixed finite set and all members are already
smooth on $S$. This constructs the required axes. The vanishings
(A31f.14) are open by upper semicontinuity in the normal-bundle
family.

All factors in (A31f.4) now have zero $H^0$, proving (A31f.11).
The exact quotient sequence consequently injects
\[
 H^1(S^{(1)},F_{S*}N\oplus N^{(1)})\hookrightarrow H^1(S^{(1)},Q).
 \tag{A31f.15}
\]
The first component of the class in (A31f.12) is $o$ itself,
so (A31f.15) proves its injectivity. No injectivity of Frobenius
on a coherent cohomology group is used.
\end{proof}

\paragraph{Finite-jet duality.}
There is a perfect pairing
\[
 \mathbb H^2(B,C)\times
 \mathbb H^4\!\left(B,
       [\Omega_B^3/I^3\Omega_B^3\to\Omega_B^4/I^2\Omega_B^4]\right)
                  \longrightarrow k.
 \tag{A31f.16}
\]
Termwise this is Serre duality: $H^2(M_2^0)$ pairs with
$H^0(\Omega_B^4|S_2)$, and $H^1(M_3^1)$ with
$H^1(\Omega_B^3|S_3)$. The differentials are adjoint up to the
total-complex sign. This must be checked since $d$ is initially
$k$-linear, not $\mathcal O_B$-linear. The local residue of
$d\gamma$ is zero: differentiating to a coefficient with exponent
$-1$ would have to differentiate exponent zero. Applying this
identity to a product gives integration by parts. It is compatible
with the localization and Cech trace defining Serre duality, so
identifies the dual differentials. The two finite spectral
sequences of this pairing therefore give the perfect pairing
(A31f.16). Equivalently, one can first push by Frobenius to make
the differentials coherent-linear and use the same residue trace
to identify their adjoints.

By A31e, the primitive subspace of $F^3H^4_{\mathrm{dR}}(B/k)$
surjects onto the second factor in (A31f.16). Thus vanishing of
all primitive periods detects the entire first factor, including
its degree-zero-form descent component.

\begin{lemma}[Canonical PD descent at the boundary]
Let $R$ have residue field $k$, uniformizer $\pi$, and $3=u\pi^2$
with $u\in R^\times$. Let $\mathcal B/R$ be a smooth projective
lift of $B$, and let $o\in H^1(S,N)\otimes(\pi)/(\pi^2)$ be the
obstruction to a smooth embedded lift $S_2\subset\mathcal B/\pi^2$.
The integral crystalline class of $S$, under canonical comparison,
has first primitive period functional represented in (A31f.16) by
\[
                    j(o/\pi)+\bar u^{-1}c((o/\pi)^{(1)}).
 \tag{A31f.17}
\]
In particular this assertion uses the canonical class and its
actual PD descent, not a choice of transport with the desired
bounded homotopies.
\end{lemma}

\begin{proof}
The ideal $(\pi)$ has canonical divided powers at $e=2=p-1$.
Berthelot--Ogus Remark 2.6(i), printed p.~171, explicitly includes
the equality boundary and identifies crystalline cohomology of
the special fibre over this PD base with the de Rham cohomology
of its smooth lift. Its first comparison triangle, followed by
the usual comparison at $(3)$, identifies the rationalization
with the fixed comparison used in (A31f.1). Thus the class
$\beta=[S]_{\mathrm{cris}}$ has a specified integral image.

Choose local smooth lifts of $S$ modulo $\pi^2$. Their normal
differences are the actual Cech obstruction $o$, exactly as in
A27. Compute the canonical comparison by its normalized PD
Poincare resolution of the special-fibre diagonals. The local
PD-polynomial contraction formulas are valid over $R/\pi^2$ with
its specified base divided powers; their validity here follows
directly from $z^{[n]}dz\mapsto z^{[n+1]}$. The local Gysin is
the product of two Cartier Gysins, natural under this PD base
change, so its normalization agrees with the canonical crystalline
cycle class. In adapted coordinates its chain map sends its source unit to
$\nu=dx\wedge dy/(xy)$. The source comparison contraction kills
$1$; the target contraction is therefore the descent homotopy
of this specified fundamental class.

For a single displacement $\pi(a,b)$, its degree-one-form term is
\[
 \pi\left(\eta(a,b)+\bar u^{-1}\xi(a,b)\right).
 \tag{A31f.18}
\]
Indeed the PD logarithm is $\sum_{n\geq1}(-1)^{n-1}z^n/n$.
At $z=\pi a$ only $n=1,3$ survive modulo $\pi^2$, and
$\pi^3/3=u^{-1}\pi$. The product with the other Cartier Gysin
gives exactly (A31f.18). Tangential displacements have zero
normal residue. This computes the actual one-step contraction,
not merely another primitive for its differential.

We verify a bound on the \emph{actual} second comparison
homotopies; otherwise a closed function with larger Frobenius
poles could be missed. Write $z$ for PD difference coordinates
on a diagonal and $H$ for the normalized ordered-coordinate
PD contraction. For another adapted coordinate system put
$z'=f(x+z)-f(x)$ and write $H'$ for its contraction. Both are
defined on the full de Rham complex over $R$, with the base
forms held fixed in their respective coordinate splittings.
They satisfy
\[
 dH+Hd=dH'+H'd=1-P,\qquad PH'=HP=0,\qquad
 d(HH')-(HH')d=H'-H,
 \tag{A31f.19}
\]
where $P$ is the same zero-diagonal augmentation. Although
the horizontal splittings differ, the full differential and
the augmentation agree: $z'=0$ is $z=0$, and both exact
difference differentials restrict to zero there. These
identities prove that $HH'\nu$ is an actual comparison of
the canonical contractions on the closed Gysin form.

Here is a coefficient bound for this actual homotopy. Give
each $z_i,dz_i$ weight one and base differentials weight zero.
Coordinate changes and their inverse power series have integral
coefficients and do not lower this weight. For a term of weight
$n$ in a two-contraction output, its ordinary Taylor coefficient
is a sum of integral coefficients divided by at most two
integers $j,k$, with
\[
                         1\leq j\leq n-1,\qquad1\leq k\leq n.
 \tag{A31f.20}
\]
To see this, ordered PD integration, after inverting $3$, is
ordinary monomial integration in one variable. At the first
integration one relative differential remains, so its integration
exponent is at most the weight minus one. The second exponent
is at most the weight. Intermediate coordinate substitutions
have no denominators other than invertible Jacobian coefficients
and cannot lower weight. This proves (A31f.20) coefficient by
coefficient, including nonhomogeneous substitutions by considering
their homogeneous pieces.

The contractions themselves and the coordinate changes are
integral on the PD algebra. At $z_i=\pi a_i$, a product of two
positive-degree divided powers is zero modulo $\pi^2$. Also
$v_\pi(\gamma_n(\pi a_i))\geq s_3(n)$, where $s_3(n)$ is its
base-three digit sum. Consequently only a pure monomial
$\gamma_{3^r}(z_i)$ could contribute at order $\pi$. Its
coefficient, by (A31f.20), is a sum of integral multiples of
$(3^r)!/(jk)$, with $j,k$ possibly varying from summand to summand.
For $r\geq2$ its $3$-adic valuation is at least
\[
 \frac{3^r-1}{2}-(r-1)-r
                  =\frac{3^r+1}{2}-2r\geq1.
 \tag{A31f.21}
\]
It therefore vanishes after evaluation modulo $\pi^2$.
Weight one is impossible for a two-contraction output.
Only weight three remains. Such an input has at most one
normal Taylor derivative of the simple fundamental fraction,
so its total normal denominator is at most three. Hence the
actual degree-zero-form comparison lies in $\pi M_2^0$.
The same argument for a single contraction leaves weights one
and three and gives the bound $\pi M_3^1$, in agreement with
(A31f.18). Invertible changes of normal frames only multiply
these pole bounds by regular coefficients; the bounds are
intrinsic annihilator bounds.

Composition of two infinitesimal comparisons is a separate
normalization check. On the double diagonal use the normalized
integral PD comparison homotopy, which is zero when either
comparison is the identity. Every monomial in its function
component consequently has positive divided-power degree in
each of the two displacement sets. Evaluation at the two
sets of $\pi$-multiples makes their product zero modulo
$\pi^2$. In common coordinates this is also visible directly:
(A31f.18) is additive in the normal displacement, since cubing
is additive in characteristic three. Thus composition contributes
no hidden function component. Changes of adapted coordinates
are the distinct $HH'$ homotopies already bounded above. Their
triple-overlap coherence follows either from these normalized
comparison maps or from injectivity of $d$ on $M_2^0$ after
applying $d$ to the coherence identity.

Use these actual normalized contractions in the Cech nerve
of the PD diagonal, as in the crystalline Poincare comparison
of A27. Its local degree-two term is the natural Gysin $\nu$;
its one-step terms are (A31f.18); its second comparison terms
have just been proved to lie in $\pi M_2^0$. Higher contractions
annihilate a two-form. Thus the reduction of its degree below
two is an actual total cocycle in $\pi C$. In particular no
closed degree-zero fraction outside $M_2^0$ has been discarded.
Since $d:M_2^0\to M_3^1$ is injective, the class of this cocycle
is determined in $H^1(Q)$ by its one-step components. Coordinate
independence is the canonical one in (A31f.3), so it is exactly
(A31f.17).

Finally pair with a lifted class in $F^3H^4_{\mathrm{dR}}(\mathcal B/R)$.
The local form-degree-two term gives zero by dimension; the
remaining terms pair with its degree-three and degree-four
pieces, exactly by (A31f.16). The reduction of the period
divided by $\pi$ is therefore the claimed functional. This
also shows that no chosen splitting of the Hodge filtration
has entered the calculation.
\end{proof}

\begin{proof}[Proof of the theorem]
Choose the nonempty open of axes supplied by (A31f.11), taking
$t$ large enough for A22, A31e, A06 and A18 as well. The case
$e=1$ is already A29. Suppose $e=2$ and (A31f.1) were filtered.
For the arbitrary abstract lift, A29 recovers its relative
Albanese embedding and the free primitive test lattice
$\mathscr T_R$. Its reduction is the entire primitive
$F^3$ test space on $B$.

The integral class $\beta$ of the last lemma pairs to zero
with every element of $\mathscr T_R$: its rational primitive
projection is filtered, ambient classes annihilate these tests
by projection formula, and the scalar period lies in the
torsion-free ring $R$. Reducing this equality to order $\pi$
and using (A31f.16)--(A31f.17) and the surjection of A31e
gives
\[
                 j(o/\pi)+\bar u^{-1}c((o/\pi)^{(1)})=0.
 \tag{A31f.22}
\]
The injectivity in (A31f.12) forces $o=0$. A smooth embedded
surface $S_2$ therefore exists. The finite initial gate in A31
now applies, with $n_0=2$, and extends this surface through all
orders. Its algebraization contradicts the geometric
nonliftability of the A14 surface. This proves the filtration
assertion.

All extra conditions on the pair were nonempty open conditions
defined after a finite extension of the finite field of definition.
They can be imposed when choosing the fixed first member in
A18 and when shrinking its subsequent parameter open. Restricting
a lisse system to a nonempty open preserves its geometric monodromy
image, and A18's closed-point selection applies on this smaller
open. Its common-realization-kernel lemma and the generic algebraic
primitive projector from A29 therefore give the class-level
assertion exactly as in A31. If a generic cycle requires a finite
field extension, filtration membership of the fixed vector still
descends to the original $K$; its model's ramification remains
the one excluded here.
\end{proof}

\paragraph{Source normalization and exact scope.}
The crystalline comparison is Berthelot--Ogus, \emph{$F$-isocrystals
and de Rham cohomology. I}, Remark 2.6(i), p.~171, at the equality
boundary; the logarithm is \S3.1, p.~175, and the existence threshold
is Lemma 3.9, p.~179. The normalized PD polynomial contractions
and the natural Gysin sending $1$ to the fundamental fraction are
the precise primary constructions already read in A27. Their
use over the present ramified PD base is justified by the local
formulas in this proof; no unramified-base theorem is being
quoted as a general ramified statement. The new factorial estimate
(A31f.20)--(A31f.21) is proved above.
\url{https://gdz.sub.uni-goettingen.de/download/pdf/PPN356556735_0072/LOG_0013.pdf}.

For cote~44, Th.~5, scan~004 and the primitive caution on
scans~008--009, this excludes an additional actual ramification
range in route (a) of QH9-axis for the specified nonzero classes.
It does not prove a negative sign, a universal lifting obstruction,
or the universal H9-axis conjecture. Any separate appeal to the
full ambient Hodge standard conjecture outside the elliptic-power
cases of A42b remains conditional on the unrefereed
032/Milne2026 preprints.

% END work/axis/axis_first_ramified_noncancellation.tex

% BEGIN work/axis/axis_even_ramification_boundary.tex
% C44_POS /root/axis_independent. Complete author proof; independent review pending.
\section{The equality-boundary order for even ramification}
\label{sec:c44-a31g}

\paragraph{Scope.}
This corollary belongs to the same developing argument as C44-A31f.
It makes no separate progress-clock claim.  It reduces the explicitly
required initial support order for every even ramification index in
characteristic three; it does not assert that this initial support exists.

\begin{corollary}[C44-A31g: one fewer prescribed order at the PD boundary]
Let $S\subset B\subset X$ be one of the sufficiently general prescribed-centre
axes of C44-A31f, over $k=\overline{\mathbf F}_3$.  Let $R$ be a complete
mixed-characteristic discrete valuation ring with residue field $k$,
uniformizer $\pi$, and absolute ramification index $e=2m$, $m\geq1$.
Let $\mathcal B/R$ be any smooth projective lift of the abstract fourfold
$B$.  Suppose there is a smooth closed lift
\[
                S_m\subset\mathcal B\otimes_RR/(\pi^m)
\tag{A31g.1}
\]
of the specified immersion $S\subset B$.  Then its fixed rational
crystalline primitive class
\[
          w=[S]_B-\kappa^*(L_h^2)^{-1}[S]_X
\tag{A31g.2}
\]
does not belong to $F^2H^4_{\mathrm{dR}}(\mathcal B_K/K)$ under the
canonical Berthelot--Ogus comparison.

For the closed A18 axes selected inside the A31f open, the corresponding
nonzero primitive etale class cannot specialize from an algebraic class
on the geometric generic fibre of such a model.  When $m=1$ the initial
condition is the original surface and this is exactly the new case of
A31f.  For $m>1$ the condition (A31g.1) remains an additional hypothesis.
\end{corollary}

\begin{proof}
Write $3=u\pi^{2m}$ with $u\in R^\times$.  The ideal $J=(\pi^m)$ has
the canonical divided powers at the equality boundary
$m=e/(3-1)$.  They need not be PD-nilpotent.  Berthelot--Ogus,
Remark~2.6(i), printed p.~171, gives the integral de Rham comparison
at this equality boundary.  Since $m\leq e$, the quotient
$A_0=R/(\pi^m)$ is a local Artin $k$-algebra; its coefficient structure
comes from $W(k)\to R$ followed by reduction, which kills $3$.
The smooth regular immersion (A31g.1) therefore has its natural integral
crystalline Gysin class over the PD base $(R,J)$.  Write
\[
                    \beta\in H^4_{\mathrm{dR}}(\mathcal B/R)
\tag{A31g.3}
\]
for its specified image.  The rational identity
$\beta\otimes K=[S]_{\mathrm{cris}}\otimes K$ under the fixed comparison
is proved exactly as in A31: the rational comparisons of
Berthelot--Ogus Theorem~2.1 and Corollary~2.2 respect pullback and cup
products, trace is normalized by ample top Chern classes, and the
projection formula identifies the Gysin classes.  The first triangle
of Remark~2.6(i), comparing $m$ with $e$, followed by its usual
$(3)$-comparison diagram at $e$, identifies this with the canonical
comparison.  Thus no new transport of $[S]$ has been selected.

Suppose that (A31g.2) were filtered.  A29 recovers the relative Albanese
embedding and the free primitive test lattice $\mathscr T_R$ for this
possibly ramified abstract model.  Projection formula, the preceding
rational identity, and torsion-freeness give
\[
                      \int_{\mathcal B/R}\beta\omega=0
                          \quad(\omega\in\mathscr T_R).
\tag{A31g.4}
\]
Let $o\in H^1(S,N_{S/B})\otimes(\pi^m)/(\pi^{m+1})$ be the actual
obstruction to extending the given $S_m$ to a smooth embedded $S_{m+1}$.
Local lifts exist, and their normal differences are multiples of
$\pi^m$.  Compute the canonical Gysin descent by the same normalized
PD contractions used in A31f, now for the base PD ideal $J$ and the
source $S_m$.  Its source unit is again killed by the comparison
homotopy.  All coefficients relevant to the top ideal reduce to $k$.

Here are the valuation checks which make that computation identical.
For every $n\geq1$,
\[
 v_\pi\!\left(\gamma_n(\pi^m a)\right)
       \geq mn-2m v_3(n!)=m s_3(n).
\tag{A31g.5}
\]
A product of two positive divided-power monomials consequently has
valuation at least $2m\geq m+1$ and vanishes in the top step.
Only a pure $\gamma_{3^r}$ could have valuation $m$.  In an actual
two-contraction coordinate-comparison homotopy, its coefficient is a
sum of integral multiples of $(3^r)!/(jk)$, with
$1\leq j\leq3^r-1$ and $1\leq k\leq3^r$, as proved in A31f.
For $r\geq2$ each summand has $3$-adic valuation at least
\[
                 \frac{3^r+1}{2}-2r\geq1.
\tag{A31g.6}
\]
It therefore acquires an additional $\pi$-valuation $e=2m$ and
vanishes modulo $\pi^{m+1}$.  Weight one is impossible for two
contractions.  The remaining weight three has total normal denominator
at most three, so the actual function homotopy lies in
$\pi^m M_2^0$.  Normalized composition homotopies have positive
divided-power degree in both displacement sets and vanish by (A31g.5).
Thus the canonical descent retains exactly the same bounded complex
$[M_2^0\to M_3^1]$ as A31f, including its actual degree-zero component.

For a single contraction the logarithm has only its terms of degrees
one and three at valuation $m$:
\[
       \frac{\pi^{3m}}3=u^{-1}\pi^m,
       \qquad
       \pi^m\bigl(\eta(a,b)+\bar u^{-1}\xi(a,b)\bigr)
\tag{A31g.7}
\]
is the resulting normal overlap form.  Consequently the first period
of (A31g.3), divided by $\pi^m$ and reduced modulo $\pi$, is represented
in A31f's supported complex by
\[
                     j(o/\pi^m)+\bar u^{-1}c((o/\pi^m)^{(1)}).
\tag{A31g.8}
\]
The finite-jet Serre--Stokes duality and primitive test surjection of
A31e--A31f apply on the same special pair $S\subset B$.  Equations
(A31g.4) force (A31g.8) to vanish.  The injectivity established in
A31f for every scalar multiple of its ordinary and Cartier summands
therefore gives $o=0$.

A smooth embedded $S_{m+1}$ exists.  Since
\[
                  m+1=\left\lfloor\frac e{3-1}\right\rfloor+1,
\tag{A31g.9}
\]
this is precisely the initial lift required in A31.  That theorem
extends it through every further order and algebraizes it, contradicting
the geometric nonliftability of the A14 surface.  This proves the
filtration statement.  The class-level statement for the selected
A18 axes follows from the common-realization-kernel and generic
primitive-projector argument of A18 and A29, exactly as in A31f.
\end{proof}

\paragraph{Grothendieck's question and limitation.}
For JAC1964 QH9-axis and H9-T7, this sharpens the route-(a) initial-support
gate attached to cote~44, Theorem~5, scan~004 and its primitive margin,
scans~008--009.  It does not prove positivity for an arbitrary fixed axis,
does not construct a primitive class lift at larger ramification, and
does not imply the product-form assertion of scan~049.  The unrefereed
032/Milne2026 preprints are not used in this corollary; a separate general
ambient height assertion based on those papers remains conditional on
them, outside the unconditional elliptic-power cases of A42b.

% END work/axis/axis_even_ramification_boundary.tex

% BEGIN work/scan_w/ambient_support_jet_restriction.tex
% C44_POS /root/scan_w. Author whole proof read; final independent reviews pending.
\section{Primitive finite jets before the first ambient Cartier pole}
\label{sec:c44-a31h}

\paragraph{Scope.}
This section supplies the finite-jet step needed to extend the equality-boundary
calculation of A31f to odd characteristic.  It does not construct a new
geometrically nonliftable surface in every characteristic.  The conclusion
about a support lift is stated for an actual specified smooth surface; a
nonliftability conclusion requires its geometric nonliftability as a separate
hypothesis.  The argument is part of the same developing lifting calculation,
not a separate staged progress-clock claim.

\begin{proposition}[C44-A31h: uniform primitive jet restriction]
Let $k=\overline{\mathbf F}_p$ with $p>2$, let $S$ be a fixed
smooth projective surface in an abelian sixfold $X$, and fix an ample bundle
$H$ on $X$.  Fix an integer
\[
                         3\leq M<4p.
\tag{A31h.1}
\]
For all sufficiently large $t$, every smooth genuine axis
$\kappa:B\hookrightarrow X$ cut out by two members of $|L|$, $L=H^t$,
and containing $S$, satisfies
\[
 \ker\!\left(F^3H^4_{\mathrm{dR}}(B/k)
                 \xrightarrow{\kappa_*}F^5H^8_{\mathrm{dR}}(X/k)\right)
 \twoheadrightarrow
 \mathbb H^4\!\left(B,
       [\Omega_B^3/I^M\Omega_B^3
                  \xrightarrow{d}\Omega_B^4/I^{M-1}\Omega_B^4]\right),
\tag{A31h.2}
\]
where $I=I_{S/B}$ and the two terms have degrees three and four.
The lower bound on $t$ is uniform in the smooth axis.
\end{proposition}

\begin{proof}
Write $J=I_{S/X}$ and $s=M-3$.  We first spell out the compatibility
between the ideal-power Koszul maps and exterior differentiation.

For $m\geq2$ the exact ideal-power Koszul sequence of A31e, tensored
by a locally free sheaf $E$, is
\[
 0\to J^{m-2}E\to(J^{m-1}EL)^{\oplus2}
       \to J^mEL^2\to\kappa_*(I^mEL^2|B)\to0.
\tag{A31h.3}
\]
Its double boundary, preceded by the natural adjunction map
$\Omega_B^a\to\Omega_X^{a+2}L^2|B$, defines the weighted coherent
Gysin maps
\[
 H^j(B,I^m\Omega_B^a)
       \longrightarrow H^{j+2}(X,J^{m-2}\Omega_X^{a+2}),
                         \qquad a=3,4.
\tag{A31h.4}
\]
For $a=4$ adjunction is an isomorphism.  These maps, together with
their Cech representatives, commute with the differentials which
lower the ideal exponent by one.

Here is an explicit verification of that last assertion, including
the bundle frames.  On a trivializing affine open let $f_0,f_1$ be the
two equations and let $\widetilde\eta$ lift an $I^m$-valued form on $B$.
The actual chain Gysin is represented by the generalized fraction
\[
       \left[\frac{df_0\wedge df_1\wedge\widetilde\eta}{f_0f_1}\right].
\tag{A31h.5}
\]
The two divisions lower the $J$-order by two, exactly as in (A31h.3).
The numerator may be chosen in $J^m$.  Changing its lift by an element
of $(f_0,f_1)$ uses coefficients in $J^{m-1}$, because
$(f_0,f_1)\cap J^m=(f_0,f_1)J^{m-1}$.  The cancelled-denominator
terms are precisely the middle Cech/Koszul boundaries in (A31h.3).
Differentiation of (A31h.5) gives
\[
 d\left(\frac{df_0\wedge df_1\wedge\widetilde\eta}{f_0f_1}\right)
       =\frac{df_0\wedge df_1\wedge d\widetilde\eta}{f_0f_1};
\tag{A31h.6}
\]
the differentiated denominators repeat $df_0$ or $df_1$ and vanish.
One can check the double boundary itself without choosing a connection.
On a common trivializing intersection, a Cech lift of a cocycle satisfies
\[
 \delta\widetilde\eta=f_0a_0+f_1a_1
                         +df_0\wedge b_0+df_1\wedge b_1,
\tag{A31h.6a}
\]
where $a_i$ have $J$-order at least $m-1$ and the normal-form terms
can be taken of order at least $m$.  To see these orders, use
$f_0,f_1$ as two of the four regular normal parameters for $J$ and
separate tangent and normal differentials.  After wedging by
$df_0\wedge df_1$, the last two terms disappear.  Applying $\delta$
once more and using the Koszul syzygy gives its leftmost cocycle
of order $m-2$.  Differentiate these same two identities: terms
obtained by differentiating $f_i$ or a normal-form term disappear
after the same wedge, while the other terms are precisely the
derivatives of $a_i$ and of that leftmost cocycle.  The leftmost
cocycle for $d\eta$ is therefore the differential of the old one,
up to its Cech coboundary, in $J^{m-3}\Omega_X^{a+3}$.
This is (A31h.6) on the actual iterated boundary, with all three
ideal orders accounted for.  If $f_i$ is replaced by
$u_if_i$, the additional $d\log u_i$ terms have at most one of the
two denominators.  They are the same Cech boundaries, with regular
coefficients, and respect these ideal orders.  Thus computing the two
successive Cech boundaries before or after (A31h.6) gives the same
map on cohomology.  This constructs the compatible map of the two-term
total Cech complexes, including its usual frame homotopies.  In particular
the leftmost, untwisted term has exactly the ordinary differential
$d:J^{m-2}\Omega_X^{a+2}\to J^{m-3}\Omega_X^{a+3}$; no connection
on $L$ or scalar multiple of $d$ has been chosen.

Put
\[
 \begin{split}
 K_M&=[I^M\Omega_B^3\longrightarrow I^{M-1}\Omega_B^4],\\
 A_s&=[J^{s+1}\Omega_X^5\longrightarrow J^s\Omega_X^6],\\
 P_s&=\Omega_X^{\geq5}/A_s.
 \end{split}
\tag{A31h.7}
\]
The degrees are $3,4$ for $K_M$ and $5,6$ for the other complexes.
The preceding construction gives Gysin on hypercohomology, raising
total degree by four, compatible with the ordinary Gysin and both
inclusions into the untruncated complexes.

Serre vanishing in (A31h.3) is uniform in $B$.  Using
\[
 0\to\Omega_B^3\to(\Omega_X^5L^2)|B
                      \to(L^3|B)^{\oplus2}\to0
\tag{A31h.8}
\]
after fixing the trivialization of $\Omega_X^6$, it gives the following
maps in the rows needed below:
\[
 \begin{array}{rcll}
 H^1(I^M\Omega_B^3)&\twoheadrightarrow&H^3(J^{s+1}\Omega_X^5),&\\
 H^2(I^M\Omega_B^3)&\xrightarrow{\sim}&H^4(J^{s+1}\Omega_X^5),&\\
 H^j(I^{M-1}\Omega_B^4)&\xrightarrow{\sim}&H^{j+2}(J^s\Omega_X^6),
                                                   &j=1,2,\\
 H^0(I^{M-1}\Omega_B^4)&\twoheadrightarrow&H^2(J^s\Omega_X^6).&
 \end{array}
\tag{A31h.9}
\]
For the first two rows tensor (A31h.8) by $I^M$; the quotient remains
locally free before tensoring, and its positive twists have no higher
cohomology.  The last row is the double boundary of (A31h.3).
All these maps are the maps just checked in (A31h.4)--(A31h.6).
The two-term hypercohomology sequences therefore give
\[
 \mathbb H^5(K_M)\xrightarrow{\sim}\mathbb H^9(A_s),\qquad
 \mathbb H^4(K_M)\twoheadrightarrow\mathbb H^8(A_s).
\tag{A31h.10}
\]
For precision, in degree five the right graded terms are the kernels
of the degree-two coherent differentials, and both relevant maps in
(A31h.9) are isomorphisms.  The left graded terms are cokernels in
degree one; surjectivity of the first row and bijectivity of the third
row identify these cokernels.  In degree four, a closed degree-one
class on $X$ has a closed preimage, since its differential maps to
zero through the injective third row.  The last row surjects on the
remaining cokernel.  This proves both assertions in (A31h.10) without
assuming the differentials on the ideal-valued cohomology are zero.

The decisive calculation is
\[
                      \mathbb H^8(P_s)=\mathbb H^9(P_s)=0.
\tag{A31h.11}
\]
For $s=0$, the degree-six quotient is zero and the assertion follows
from the two-dimensional support of the other quotient.  Suppose
$s>0$.  Write $S_{a,X}=\operatorname{Spec}_X(\mathcal O_X/J^a)$.
Both terms of $P_s$ are supported on a surface, so its degree-eight
hypercohomology is precisely
\[
 \operatorname{coker}\left(
 H^2(\Omega_X^5|S_{s+1,X})
         \xrightarrow{d}H^2(\Omega_X^6|S_{s,X})\right),
\tag{A31h.12}
\]
and its degree-nine hypercohomology is zero.  Powers of the regular
four-generator ideal $J$ have Cohen--Macaulay quotients of codimension
four: their associated graded pieces are the locally free symmetric
powers of $J/J^2$ on the smooth surface, and the depth lemma applies
to the finite filtration.  Set
\[
 M_{a,X}^b=\mathcal Ext_X^4(\mathcal O_X/J^a,\Omega_X^b)
             =\operatorname{ann}_{J^a}\mathcal H_S^4(\Omega_X^b).
\tag{A31h.13}
\]
Serre duality identifies the dual of the arrow in (A31h.12) with
\[
                d:H^0(M_{s,X}^0)\longrightarrow H^0(M_{s+1,X}^1).
\tag{A31h.14}
\]
This assertion for the $k$-linear differential follows by residue
integration by parts: the residue of a differentiated generalized
fraction is zero, since a coefficient of exponent $-1$ would have
to come from differentiating exponent zero.  Applying this to a
product gives the adjoint identity.  It commutes with the Cech trace
defining Serre duality.  Alternatively push by Frobenius first, so
that all differentials are coherent-linear, and use the same trace.

The sheaf arrow underlying (A31h.14) is injective.  Indeed, in four
normal coordinates a fraction in $M_{s,X}^0$ has four positive
denominator exponents with sum at most $s+3=M$.  Take the highest
total-pole part of a nonzero fraction.  A zero normal derivative
would require every one of its four exponents to be divisible by
$p$: for a fixed differential component the differentiated fraction
has a unique monomial preimage, so its coefficient cannot cancel
with another monomial.  Four positive multiples of $p$ have sum
at least $4p$, excluded by (A31h.1).  Thus a highest normal derivative
is nonzero.  Tangential derivatives do not raise that pole order
and cannot cancel it.  This proves injectivity, hence (A31h.11).

It follows that
\[
 \mathbb H^8(A_s)\twoheadrightarrow
                         \mathbb H^8(\Omega_X^{\geq5}),\qquad
 \mathbb H^9(A_s)\xrightarrow{\sim}
                         \mathbb H^9(\Omega_X^{\geq5}).
\tag{A31h.15}
\]
The ordinary, unweighted versions of (A31h.8)--(A31h.9) also give
\[
 \mathbb H^5(\Omega_B^{\geq3})
                      \xrightarrow{\sim}\mathbb H^9(\Omega_X^{\geq5}).
\tag{A31h.16}
\]
Consequently the map $\mathbb H^5(K_M)\to
\mathbb H^5(\Omega_B^{\geq3})$ is an isomorphism.  The restriction
triangle on $B$ proves surjectivity onto the jet hypercohomology in
(A31h.2).  Moreover (A31h.10) and (A31h.15) show that the image of
$\mathbb H^4(K_M)$ in $\mathbb H^4(\Omega_B^{\geq3})$ surjects,
under Gysin, onto $\mathbb H^8(\Omega_X^{\geq5})$.  This image is
the restriction kernel.  Subtracting such a class kills
the ambient Gysin without changing the prescribed jets.  Finally
the degeneration established in A22a, using an embedded lift of this
smooth axis, identifies these truncated hypercohomology groups with
$F^3H^4(B)$ and $F^5H^8(X)$.  This proves (A31h.2).
\end{proof}

\begin{proposition}[The ordinary and Cartier summands for odd $p$]
Let $S\subset X$ be as above.  For sufficiently high $L$ there is a
nonempty open set of smooth genuine axes containing $S$ with the following
property.  Put $N=N_{S/B}$ and, on $B$,
\[
 M_m^a=\mathcal Ext_B^2(\mathcal O_B/I^m,\Omega_B^a),\qquad
 C_p=[M_{p-1}^0\xrightarrow{d}M_p^1],\qquad Q_p=\operatorname{coker}d.
\tag{A31h.17}
\]
The terms have form degrees zero and one; local cohomology contributes
the additional support shift two.  After Frobenius pushforward these
are coherent modules on $S^{(1)}$: $p$-th powers from $I$ annihilate
both terms, and differentiation is linear for $p$-th powers.
There are canonical disjoint injections
\[
 j:F_{S*}N\hookrightarrow Q_p,\qquad c:N^{(1)}\hookrightarrow Q_p,
\tag{A31h.18}
\]
and their sum induces an injection on $H^1$.  In particular
$o\mapsto j(o)+\lambda c(o^{(1)})$ is injective for every $\lambda\in k$.
The primitive $F^3$ period tests detect all of $\mathbb H^2(C_p)=H^1(Q_p)$.
\end{proposition}

\begin{proof}
The normal fractions in $M_{p-1}^0$ have two positive denominator
exponents with sum at most $p$.  The leading normal derivative is
injective, since a zero derivative would require sum at least $2p$.
It is strict for the pole filtration, shifted by one.  The ordinary
and inverse-Cartier representatives, for normal coordinates $x,y$, are
\[
 \eta(a,b)=\frac{a\,dy-b\,dx}{xy},\qquad
 \xi(a,b)=\frac{a^p\,dy}{x^py}-\frac{b^p\,dx}{xy^p}.
\tag{A31h.19}
\]
Here inverse Cartier is applied to the ordinary contraction on the
Frobenius twist, so it does not assume that $\eta$ is closed.
The ordinary Cartier isomorphism is natural under localization.
To pass to local cohomology, use the support Cech complex for the two
regular generators of $I$.  Its local cohomology on each locally free
$\Omega_B^a$ is concentrated in degree two; the same holds for the
locally free Cartier cohomology sheaves on the Frobenius twist.
The two spectral sequences of this supported de Rham complex
therefore both collapse to that row.  They give the natural Cartier
isomorphism on the cohomology of $\mathcal H_S^2(\Omega_B^\bullet)$.
It gives coordinate independence of $\xi$ modulo exact local-cohomology
forms.  Its exact primitives can be chosen in $M_{p-1}^0$.
To verify the bound, take a primitive of an exact form of total pole
at most $p+1$.  If its highest function pole has total order $n>p$,
its normal derivative must vanish.  All its positive exponents are
then divisible by $p$, so $n\geq2p>p+1$.  Its tangential derivative
must vanish as well.  In etale tangent coordinates its coefficients
are therefore $p$-th powers.  Subtract these closed $p$-th-power fractions
and repeat.  This lowers the highest pole until the primitive lies
in $M_{p-1}^0$.  The argument can be made after completion in normal
coordinates; faithful flatness gives the same annihilator bound.
Since $d$ is injective on that bounded module, its exact homotopies
are unique and satisfy the coordinate-transition coherence identities.
This proves the second map in (A31h.18); the first is the natural
coherent contraction.  Their leading poles, one and $p$, are distinct.

The normal part of the degree-$m$ pole quotient of $Q_p$, for
$1\leq m\leq p$, is $\operatorname{Sym}^mN$, and the tangential
part is $\det N\otimes\Gamma^{m-1}N\otimes\Omega_S^1$.
Indeed the leading normal differential is
$\det N\otimes\Gamma^{m-2}N\to
\det N\otimes\Gamma^{m-1}N\otimes N^*$.
All factorials through $(p-1)!$ are units.  In symmetric-power bases
this is the ordinary two-variable Koszul map, whose cokernel is
$\operatorname{Sym}^mN$.  At the top pole one has the exact sequence
\[
 0\to F_S^*N^{(1)}\to\operatorname{Sym}^pN
          \to\operatorname{Sym}^{p-2}N\otimes\det N\to0.
\tag{A31h.20}
\]
In a basis $u,v$, the final arrow sends $u^{p-i}v^i$ to
$i\,u^{p-i-1}v^{i-1}(u\wedge v)$, for $1\leq i\leq p-1$;
it is equivalently the differential of a homogeneous degree-$p$
polynomial, divided by $u\,dv-v\,du$.  This verifies its equivariance,
kernel and surjectivity in every odd characteristic.

It follows that $Q_p/(jF_{S*}N+cN^{(1)})$ has a finite filtration
with factors among
\[
 \begin{gathered}
 F_{S*}(\det N\otimes\Gamma^{m-1}N\otimes\Omega_S^1)
                                                \quad(1\leq m\leq p),\\
 F_{S*}\operatorname{Sym}^mN\quad(2\leq m\leq p-1),\qquad
 F_{S*}(\operatorname{Sym}^{p-2}N\otimes\det N),\qquad
 N^{(1)}\otimes B_S^1,
 \quad B_S^1=F_{S*}\mathcal O_S/\mathcal O_{S^{(1)}}.
 \end{gathered}
\tag{A31h.21}
\]
The Frobenius sheaves are on $S^{(1)}$; $N^{(1)}$ denotes base change,
not the absolute Frobenius endomorphism of a coherent cohomology group.

Write $E=N_{S/X}$.  The normal sequence gives $N\subset E$ and
$\det N=K_S L^{-2}|S$.  Every determinant-twisted factor in
(A31h.21) embeds into a fixed bundle tensored with $L^{-2}|S$, and
therefore has no global sections for large $t$, uniformly in $B$.
The other factors embed into the finitely many fixed section spaces
$H^0(\operatorname{Sym}^mE)$, $2\leq m\leq p-1$, and
$H^0(E^{(1)}\otimes B_S^1)$.  Finitely many prescribed normal kernel
two-planes exclude every nonzero section of these spaces.  At a point
where a section is nonzero, some two-plane excludes its value: the
intersection of the corresponding symmetric-power subspaces over all
two-planes is zero, as is the intersection of their tensor products
with a fixed vector space.  Iterating cuts the remaining finite-dimensional
section space to zero.  Perfectness of $k$ makes every twisted plane
available for the last factor.

For large $t$, two sections of $JL$ prescribe these finitely many normal
first jets.  General sections subject to them remain smooth along $S$:
rank drop for their $2$-by-$4$ normal matrix has codimension three,
greater than $\dim S=2$.  The off-centre and genuine-pencil jet conditions
of A06 remain nonempty with these fixed finite conditions.  Upper
semicontinuity makes the vanishing of all these section spaces an open
condition.  Thus the quotient in (A31h.21) has zero $H^0$ on a nonempty
open of genuine axes.  Its exact sequence injects the $H^1$ of the
direct sum (A31h.18).  The first coordinate of the stated obstruction
pair is $o$, proving noncancellation for every scalar.

Serre duality and residue integration by parts give the perfect pairing
\[
 \mathbb H^2(C_p)\times
 \mathbb H^4[\Omega_B^3/I^p\Omega_B^3
                    \to\Omega_B^4/I^{p-1}\Omega_B^4]\longrightarrow k.
\tag{A31h.22}
\]
This is the same termwise Ext/Serre pairing as A31f; the residue of an
exact fraction is zero in every characteristic.  The first proposition,
with $M=p$, says that primitive $F^3$ tests surject onto its second
factor.  This proves the last assertion.
\end{proof}

\begin{corollary}[The equality boundary in odd characteristic]
Keep the nonempty open of axes in the preceding proposition.  Let $R$
be a complete mixed-characteristic DVR with residue field $k$ and
absolute ramification index $e=p-1$.  For every smooth projective
abstract lift $\mathcal B/R$ of such an axis, if the fixed primitive
class
\[
             [S]_B-\kappa^*(L_h^2)^{-1}[S]_X
\tag{A31h.23}
\]
lies in $F^2$ under canonical crystalline comparison, the specified
immersion $S\subset B$ has a smooth embedded lift over $R/(\pi^2)$.
Consequently, when $S$ is geometrically nonliftable, this primitive
class is not filtered on any such model with $e\leq p-1$.

More generally, let $e=(p-1)m$ with $m\geq1$, and suppose that the
specified immersion has a smooth embedded lift $S_m$ over $R/(\pi^m)$.
Filtration membership of (A31h.23) forces an extension of this given
lift to $S_{m+1}$.  If $S$ is geometrically nonliftable, filtration
membership is therefore impossible under this explicit initial-support
hypothesis.  For $m>1$ no existence of $S_m$ is asserted.

In particular, the actual centres of A14b give genuine closed axes in
$E^6$ for every prime $p\equiv3\pmod4$ whose nonzero primitive etale
class cannot lift algebraically on the geometric generic fibre of any
abstract smooth projective model with $e\leq p-1$.  These axes can be
chosen with $A_B=\mathbf Qw$ and $w^2>0$, as in A18.  The same class-level
exclusion holds at $e=(p-1)m$ under the displayed supplied-support
hypothesis.  Thus the first new equality cases beyond characteristic
three include $(p,e)=(7,6)$ and $(11,10)$; no new centre construction is
needed for these instances.
\end{corollary}

\begin{proof}
Write $p=u\pi^{p-1}$.  The equality-boundary comparison of
Berthelot--Ogus Remark 2.6(i) gives the specified integral crystalline
Gysin class $\beta$ of $S$ in the de Rham cohomology of $\mathcal B$;
its rationalization is the fixed comparison used in (A31h.23).
We check the general-$p$ bounds in the actual normalized PD calculation
of A31f, rather than discarding its degree-zero descent component.

Modulo $\pi^2$ the logarithm for a displacement $\pi(a,b)$ retains
precisely its terms of degrees one and $p$, since
\[
 v_\pi(\pi^n/n)=n-(p-1)v_p(n)=1
                            \quad\Longleftrightarrow\quad n=1,p.
\tag{A31h.24}
\]
Thus its one-step Gysin contraction is
$\pi(\eta(a,b)+\bar u^{-1}\xi(a,b))$.
For two coordinate splittings the normalized contractions satisfy
$d(HH')-(HH')d=H'-H$, with the same zero-diagonal augmentation.
In weight $n$ the ordinary coefficient of this actual two-contraction
homotopy is a sum of integral multiples of $1/(jk)$, with
$1\leq j\leq n-1$ and $1\leq k\leq n$, as in A31f's full-differential
calculation.  Coordinate substitutions have integral coefficients and
do not lower weight.  A product of two positive divided powers of
$\pi$ vanishes modulo $\pi^2$.  A pure divided power can have valuation
one only in degree $p^a$, since its valuation is the base-$p$ digit sum.
For $a\geq2$ its coefficient has $p$-adic valuation at least
\[
        v_p((p^a)!)-(a-1)-a
           =\frac{p^a-1}{p-1}-2a+1\geq1.
\tag{A31h.25}
\]
At $a=2$ the final expression is $p-2\geq1$, and its increment when
$a$ is increased by one is $p^a-2>0$.
These terms vanish after evaluation.  Weight one is impossible for
two contractions; only weight $p$ remains.  Its input has at most
$p-2$ normal Taylor derivatives of the simple fundamental fraction,
so the resulting function has total denominator at most $p$ and
lies in $\pi M_{p-1}^0$.  The corresponding one-contraction calculation
leaves weights one and $p$, with total denominator at most $p+1$,
and lies in $\pi M_p^1$.  Normalized composition homotopies have
positive divided-power degree in each displacement set and vanish
modulo $\pi^2$.  Higher contractions annihilate the two-form Gysin.

The actual canonical Cech descent therefore lies in $\pi C_p$.
The uniqueness of bounded exact primitives proved above gives its
coordinate-transition coherence.  Its image in $H^1(Q_p)$ is exactly
\[
                   j(o/\pi)+\bar u^{-1}c((o/\pi)^{(1)}),
\tag{A31h.26}
\]
where $o$ is the ordinary obstruction to the specified smooth $S_2$.
The chain Gysin sends its source unit to the fundamental fraction;
the source contraction kills that unit.  Thus this is the canonical
class of the immersion, with no independently chosen transport or
higher-pole closed function omitted.

For the arbitrary abstract model, A29 supplies its relative Albanese
embedding and the free full primitive test lattice.  If (A31h.23) is
filtered, projection formula and $F^2F^3=0$ make its integral Gysin
periods vanish on that lattice.  Their order-$\pi$ reductions are
(A31h.26) paired with (A31h.22).  The jet surjection and noncancellation
force $o=0$, giving $S_2$.  This is the initial order required by A31,
since $\lfloor e/(p-1)\rfloor+1=2$.  If $S$ is geometrically
nonliftable, its subsequent all-order induction is a contradiction.
For $e<p-1$ A31 already starts from the given special surface.

For the last assertion write $p=u\pi^{(p-1)m}$ and use the equality
PD ideal $(\pi^m)$.  The integral class is now the natural Gysin of
the given $S_m$ over that PD base.  Its rationalization is the fixed
class of $S$, by A31's pullback, trace and comparison-triangle argument.
For a displacement $\pi^m a$ one has
\[
 v_\pi(\gamma_n(\pi^m a))\geq m s_p(n),\qquad
                  \frac{\pi^{pm}}p=u^{-1}\pi^m.
\tag{A31h.27}
\]
Mixed positive divided powers have valuation at least $2m\geq m+1$.
The factorial estimate (A31h.25) kills all pure weights $p^a$ with
$a\geq2$ at this top step as well.  Thus the same actual bounded
descent in $\pi^m C_p$, divided by $\pi^m$, is (A31h.26) with
$o/\pi$ replaced by $o/\pi^m$.  All its coefficients reduce to $k$.
The same primitive tests and noncancellation give the obstruction
zero and extend the prescribed $S_m$ to $S_{m+1}$.  This is exactly
$\lfloor e/(p-1)\rfloor+1$, so A31 applies.  The argument uses the
given smooth support on the equality boundary and does not replace
the canonical divided powers by trivial ones.

Finally, A14b constructs the required geometrically nonliftable surface
in $E^4$ for every $p\equiv3\pmod4$, and embeds it in $E^6$ by two zero
coordinates.  The finite-plane and smoothness conditions above define
a nonempty open in the same prescribed-centre parameter space used in
A18.  They may be imposed on its fixed first divisor and its subsequent
hyperplane open.  Restricting a lisse system to a nonempty open preserves
its geometric monodromy image, so A18 selects closed axes there with
$A_B=\mathbf Qw$, $w^2>0$, and its common-realization kernel property.
The algebraic generic primitive projector of A29 and that common-kernel
property identify the crystalline specialization of any putative generic
algebraic primitive class with (A31h.23).  It would be filtered, contrary
to the result just proved.  If the generic cycle is defined only after
a finite extension, membership of this fixed vector in the filtration
descends to the original fraction field, so the original model's
ramification is still the one being excluded.
\end{proof}

\paragraph{Exact inputs and remaining scope.}
The ideal-power resolution, tangent adjunction and arbitrary-model
primitive lattice are the checked A31e/A22 and A29 constructions.
The ordinary Cartier isomorphism is Deligne--Illusie,
\emph{Relevements modulo $p^2$ et decomposition du complexe de de Rham},
Theorem 1.2, printed pp.~249--250; the two supported spectral sequences
needed here were spelled out above.  Only this Cartier theorem is used
from that source, not a dimension-restricted degeneration theorem.
\url{https://publications.ias.edu/sites/default/files/Number57.pdf}.
The natural local Gysin and its de Rham normalization are Gros,
\emph{Classes de Chern et classes de cycles en cohomologie de Hodge--Witt
logarithmique}, II, \S3.4, printed pp.~43--44, and II, \S4, p.~51,
as read in A27.  The filtered compatibility needed here is proved
explicitly in (A31h.3)--(A31h.6), rather than attributed to an unstated
ideal-power theorem.  The equality-boundary comparison and PD logarithm
are Berthelot--Ogus, \emph{$F$-isocrystals and de Rham cohomology. I},
Remark 2.6(i), p.~171; \S3.1, p.~175; and Lemma 3.9, p.~179.
The normalized PD contraction and its composition argument are exactly
those whose complete canonical-descent proof was checked in A31f.
The new numerical bound (A31h.25) is written out above.

For cote~44, Th.~5, scan~004 and the primitive caution at scans~008--009,
this is a support-lifting criterion at the first equality boundary.
The A14 characteristic-three centre recovers A31f, while the already
checked A14b centres give the stated actual instances in every
$p\equiv3\pmod4$.  No new nonliftable centre is constructed here.
No ramification $e>p-1$ is excluded without the explicit additional
support hypothesis, and no universal fixed-axis sign assertion is made.

% END work/scan_w/ambient_support_jet_restriction.tex

% BEGIN work/axis/axis_second_ramified_noncancellation.tex
% C44_POS /root/axis_independent. Author, root and independent whole-proof checks passed.
\section{Radial descent immediately beyond the PD boundary}
\label{sec:c44-a31i}

\paragraph{Review status.}
This section contains the rational canonical-descent calculation
for $p=3,e=3$ and the independently checked finite-pole quotient proof.
The author, root and independent readers have checked the whole integrated proof.  This is part of the
same developing lifting argument as A31f--A31h, with no separate staged
clock claim.

\subsection{The canonical rational class and its local resolution}

Let $S\subset B$ be a smooth surface in a smooth projective fourfold over
$k=\overline{\mathbf F}_3$.  Let $\mathcal B/R$ be a smooth projective
lift, where $R$ is a complete mixed-characteristic DVR and
\[
                         3=u\pi^3,\qquad u\in R^\times.
\tag{A31i.1}
\]
Write $\beta$ for the fixed rational crystalline class of $S$ transported
to $H^4_{\mathrm{dR}}(\mathcal B_K/K)$ by the usual comparison.
No preservation of an integral lattice by this comparison is assumed.
Put $I=I_{S/B}$ and
\[
 M_a^b=\mathcal Ext_B^2(\mathcal O_B/I^a,\Omega_B^b),\qquad
 C_9=[M_8^0\xrightarrow d M_9^1],\qquad Q_9=\operatorname{coker}d.
\tag{A31i.2}
\]
The two form degrees are zero and one; the support degree is two.

The checked rational comparison C44-A31-SA1 identifies $\beta$ with the ordinary
canonical rigid fundamental class.  The local supported divisor class is
the actual logarithmic cocycle of Petrequin, Propositions 3.9--3.10 and
Lemma 3.12; Proposition 6.6 identifies it with Gysin of the source unit.
Its transverse product, by Lemma 7.1, has local generalized-fraction
representative
\[
                            \nu=\frac{dx\wedge dy}{xy}.
\tag{A31i.3}
\]
Here $x,y$ are regular equations for a local smooth lift of the surface.
This representative is independent of a regular-sequence frame in local
cohomology.  The source unit is annihilated by every comparison contraction.
The following construction specifies the complete comparison of these local
representatives, including the function in Cech degree two.

We make the comparison frames and their evaluations explicit.  Choose small
etale adapted charts on $B$, with tangent coordinates $t_1,t_2$ on $S$ and
normal parameters $x,y$.  Lift them to the smooth formal ambient charts.
Local smooth formal support lifts exist by the relative Jacobian criterion;
no compatibility between these local supports is imposed.  After an etale
refinement of an intersection, formal etaleness identifies each local
support with the same formal tangent chart.  It is then a graph
\[
             x=u_i(t),\qquad y=v_i(t),\qquad u_i,v_i\in\pi\mathcal O.
\tag{A31i.3a}
\]
Use graph zero as reference, and put $x_0=x-u_0(t)$,
$y_0=y-v_0(t)$.  Before evaluation, introduce independent displacement
vertices $(z_i,w_i)$ and the one common universal fraction
\[
 \nu_i=d\log(x_0+z_i)\wedge d\log(y_0+w_i),\qquad
        (z_i,w_i)\longmapsto(u_0-u_i,\ v_0-v_i).
\tag{A31i.3b}
\]
The evaluation is an actual morphism of full de Rham complexes and includes
the differentials of the graph functions.  Every displacement is a
$\pi$-multiple.  Collapsing universal vertices is performed \emph{before}
this evaluation, so $P_i\nu_j=\nu_i$ is literal in this common comparison
frame.  It is not an assertion that the point-diagonal collapse identifies
different actual supports.  The graph differences modulo $\pi^2$, with
their normal-frame transition matrices, are exactly the ordinary Cech
normal obstruction.  Thus no smooth global $S_2$ has been assumed.

These frames are compatible with refinements and face maps by restricting
the graph functions and forgetting vertices.  We also address changes of
the common tangent chart, where regraphing matters.  Apply the same ambient
coordinate change at all universal vertices.  To rewrite each transformed
support as a graph, formal etaleness reparametrizes its tangent coordinates.
The resulting additional adapted coordinate changes $f_i$ preserve the zero
normal ideal and have the same reduction modulo $\pi$.  Choose one of them
as $f_0$ and factor $f_i=f_0a_i$.  Each $a_i$ preserves $(x,y)$ and is the
identity modulo $\pi$.  In coordinates the formal interpolation
\[
                           a_{i,s}=\mathrm{id}+s(a_i-\mathrm{id})
\tag{A31i.3c}
\]
is an isomorphism of pairs over the unit parameter disc: its Jacobian is
the identity modulo $\pi$, and its two normal coordinates generate the
same ideal throughout.  By the intrinsic local fundamental-class identity,
the pullback of $dx\wedge dy/(xy)$ along this pair isomorphism is exactly
that fraction in the \emph{sheaf} $\mathcal H^2_{(x,y)}(\Omega^2)$,
including parameter differentials.  This is equality of local-cohomology
sections, not merely equality of total de Rham classes.  Equivalently,
the local Gysin of the source unit is natural for this isomorphism of
pairs; terms containing a parameter differential have a cancelled normal
denominator.  All contractions in these extra pair-preserving directions
are consequently actual zero sections.  Their pullbacks through a support
displacement remain zero by functoriality of the local Gysin and local
cohomology maps.  Regraphing therefore introduces no uncomputed closed
function.  What remains is the simultaneous change of common coordinates,
whose explicit frame and composition identities are proved below.

All computations take place first on the tubes of the appropriate formal
diagonals and their localized de Rham Cech complexes.  On a sufficiently
small strict neighborhood, the relevant displacement ratios have norm less
than one.  The geometric power series and logarithms converge there.
Division of their weight-$n$ coefficients by $n$, or by $nj$ with
$1\leq j\leq n$, preserves convergence: the valuations of these integers
grow only logarithmically with $n$.  Shrinking a strict neighborhood is
allowed in the overconvergent-support complex.  Thus the homotopies below
are actual rational analytic homotopies in the resolution used for the
canonical rigid class, not merely formal expressions assigned a class
afterwards.  Evaluation of displacements at multiples of $\pi$ will have
the stronger integral bounds proved below.

\subsection{Normalized radial contractions and exact Cech coherence}

On a formal diagonal use difference coordinates $z$ and base coordinates
$x$.  Give $z_i,dz_i$ weight one and base functions and differentials
weight zero.  Let $E=\sum z_i\partial_{z_i}$, and define on positive
homogeneous weight $n$
\[
                           h=\frac{\iota_E}{n}.
\tag{A31i.4}
\]
Set $h=0$ on weight zero.  For the full de Rham differential over $R$,
including differentiation of base coefficients, one has
\[
                 dh+hd=1-P,\qquad Ph=hP=0,
\tag{A31i.5}
\]
where $P$ is zero-diagonal augmentation.  These are the ordinary Euler
identities; no divided powers on $(\pi)$ have been introduced.

Here is a compatible simplicial construction.  On an intersection with
vertices $0,\ldots,r$, write $\nu_i$ for the pullback of (A31i.3) from
vertex $i$.  Let $P_i$ collapse the vertices after $i$ to vertex $i$,
leaving the earlier vertices fixed.  Let $h_i$ be radial contraction in
these tail differences, with all head coordinates held fixed.  In addition
to (A31i.5), the normalizations give
\[
             P_i h_j=0\quad(i<j),\qquad
             P_j(\nu_k-\nu_j)=0\quad(j<k).
\tag{A31i.6}
\]
The first identity follows because the output of $h_j$ has positive
degree in a difference between a later vertex and vertex $j$.

Define the one-step and two-step homotopies by
\[
 \eta_{ij}=h_i(\nu_j-\nu_i),\qquad
 \kappa_{ijk}=h_i h_j(\nu_k-\nu_j).
\tag{A31i.7}
\]
Restriction to a face gives the same operators on the forms involved;
unused tail variables make no contribution.  The first identity in
(A31i.5) gives $d\eta_{ij}=\nu_j-\nu_i$.  On a triple put
$\omega=\nu_k-\nu_j$.  Then $P_j\omega=0$ and $P_i h_j\omega=0$,
so
\[
 d(h_i h_j\omega)=(h_j-h_i)\omega
                    =\eta_{jk}-\eta_{ik}+\eta_{ij}.
\tag{A31i.8}
\]
On four vertices the alternating sum of the two-step homotopies is
\[
 (h_jh_k-h_ih_k+h_ih_j)(\nu_l-\nu_k).
\tag{A31i.9}
\]
The differential of $h_i h_j h_k$ is the displayed operator, except
for augmentation terms.  These terms vanish by (A31i.6) and
$P_k(\nu_l-\nu_k)=0$.  A product of three contractions annihilates
a two-form.  Since $d(\nu_l-\nu_k)=0$, (A31i.9) is exactly zero.
Thus $d\kappa=\delta\eta$ and $\delta\kappa=0$.  With total
differential $\delta+(-1)^{\mathrm{Cech\ degree}}d$, the signed
degree-zero term is $-\kappa$.  This proves the full cocycle identity,
not just that a putative composition function is closed.

We also verify frame coherence explicitly.  Let $h_i,g_i$ be the
radial contractions in two adapted coordinate systems; their geometric
augmentations $P_i$ are identical.  On a pair set
\[
                       K_{ij}=h_i g_i(\nu_j-\nu_i).
\tag{A31i.10}
\]
Then $\eta^g_{ij}-\eta^h_{ij}=dK_{ij}$.  On a triple, writing
$\omega=\nu_k-\nu_j$, the expression
$\kappa^g-\kappa^h-\delta K$ is
\[
                  (g_ig_j-h_ih_j-h_jg_j+h_ig_i)\omega.
\tag{A31i.11}
\]
It is the endomorphism differential of
$h_i g_i g_j-h_i h_j g_j$ applied to $\omega$.  Every extra
augmentation term is zero: $P_i h_j=0$, $P_j g_j=0$, and
$P_j\omega=0$.  That three-contraction expression itself vanishes
on a two-form, so (A31i.11) is zero, rather than an unspecified
closed function.  The same three-contraction identity, with one common
augmentation, proves coherence for three frame changes.  These identities
include changes with coefficients depending on tangent coordinates and
nonlinear normal changes, since the full differential and the geometric
augmentations were used throughout.

Consequently (A31i.7), together with these specified pair coboundaries,
descends on arbitrary intersections of adapted charts.  It is the
normalized diagonal-resolution descent of the local source unit under
the actual chain Gysin.  In a common ordered Cartesian frame this agrees
with the familiar two-Cartier formula
\[
 \eta_{ij}^{\mathrm{ord}}=f_{ij}\beta_j-\alpha_i g_{ij},\qquad
 \kappa_{ijk}^{\mathrm{ord}}=-f_{ij}g_{jk},
\tag{A31i.12}
\]
where $\alpha_i=d\log x_i$, $\beta_i=d\log y_i$,
$f_{ij}=\log(x_j/x_i)$ and $g_{ij}=\log(y_j/y_i)$.
Its two identities follow directly from additivity of the logarithm,
as in C44-A31-SA1.  Radial and ordered contractions with the same augmentation
differ by the actual homotopy $hH_{\mathrm{ord}}$, since
$d(hH_{\mathrm{ord}})-(hH_{\mathrm{ord}})d=H_{\mathrm{ord}}-h$.
Thus the normalization agrees with the canonical local supported divisor
classes and their product.  The preceding exact frame and composition
identities construct the complete supported cocycle.  Only after this
construction may purity, or equivalently its source-unit Gysin
normalization, identify its supported class with the fundamental class.
No closed function was deleted in making that identification.

\subsection{The sharp two-contraction denominator bound}

For a change of difference coordinates
$z'=f(x+z)-f(x)$, the coordinate substitutions and their inverses have
integral coefficients and do not lower weight.  Their weight-preserving
part commutes with $\iota_E$.  This last assertion also holds when its
linear matrix depends on $x$: the term $dA\,z$ in
$d(Az)=A\,dz+dA\,z$ has the same weight, and contraction gives $Az$.

In a weight-$n$ term of $hg$, write $j$ for the weight on which the
inner radial contraction integrates.  Its coefficient has denominators
at most $nj$.  If $j=n$, the final coordinate pullback contributes only
its weight-preserving part, and the two Euler contractions compose to
zero.  Therefore only
\[
                           1\leq j<n
\tag{A31i.13}
\]
occurs.  This argument applies termwise to a nonlinear substitution by
its homogeneous expansion.

The same bound holds for the actual tail-composition operator $h_i h_j$.
Expand in head weight $r$ relative to vertex $i$ and tail weight $j$
relative to vertex $j$.  The outer total weight is $n=r+j$.
If $r=0$, no positive-weight head coefficient or differential remains;
the outer Euler contraction is the same as the inner tail contraction,
whose square is zero.  Hence $r>0$, giving (A31i.13).  Nonlinear
coordinate substitutions only add nonnegative weight, and terms with
inner weight equal to final weight are again killed by the common
linear Euler operator.  For the ordered-to-radial comparison
$hH_{\mathrm{ord}}$, the first ordered integration exponent is also
at most $n-1$: a second relative differential must remain for $h$ to
contract.  This proves the bound without falsely asserting that radial
and ordered contractions have zero composite in equal homogeneous degree.

After evaluation of all differences at multiples of $\pi$, a weight-$n$
two-contraction coefficient has valuation at least
\[
       n-3\left(v_3(n)+\lfloor\log_3(n-1)\rfloor\right).
\tag{A31i.14}
\]
For $n\geq2$ this is nonnegative, and it is zero precisely at
$n=3,6,9$.  For completeness, direct inspection gives the assertion
through $n=12$.  For $13\leq n\leq26$ the subtracted quantity is at
most twelve.  If $3^a\leq n<3^{a+1}$ with $a\geq3$, it is at most
$6a$, and $3^a>6a$.  These inequalities also show convergence to
infinite valuation.  A single contraction has denominator $n$ and
valuation at least $n-3v_3(n)$; only $n=3$ has valuation zero.

Two contractions can act only on the two-relative-differential part
of the Gysin form.  A weight-$n$ input therefore has at most $n-2$
normal Taylor derivatives of its simple two-denominator fraction.
Its resulting function has total normal denominator at most $n$.
The corresponding one-contraction bound is $n+1$.
Regular coefficients and invertible changes of normal frame preserve
these intrinsic ideal-annihilator bounds.  Thus, modulo $\pi$, every
actual function term in the construction lies in $M_8^0$, and every
one-form term lies in $M_3^1\subset M_9^1$.  This includes the nonzero
closed function of total pole six in (A31i.12); it has not been
discarded.  It also includes any surviving closed weight-nine frame term.

\subsection{The leading Cartier class}

Let $o\in H^1(S,N_{S/B})\otimes(\pi)/(\pi^2)$ be the ordinary
obstruction to a smooth embedded lift of the specified surface.
For a local displacement $\pi(a,b)$ the weight-three part of the
radial one-form, after evaluating $3=u\pi^3$, is $\bar u^{-1}A(a,b)$,
where, in characteristic three,
\[
\begin{split}
 A(a,b)={}&(a\,dy-b\,dx)
  \left(\frac{a^2}{x^3y}+\frac{ab}{x^2y^2}+\frac{b^2}{xy^3}\right)\\
 &-(a\,db-b\,da)
  \left(\frac{a}{x^2y}+\frac{b}{xy^2}\right).
\end{split}
\tag{A31i.15}
\]
Direct differentiation, including the $da,db$ terms, gives
\[
 A(a,b)-\left(\frac{a^3dy}{x^3y}-\frac{b^3dx}{xy^3}\right)
   =d\left(-\frac{a^2b}{x^2y}+\frac{ab^2}{xy^2}\right).
\tag{A31i.16}
\]
The primitive on the right lies in $M_2^0$.  Subtracting this specified
pair coboundary changes the function terms by bounded Cech coboundaries
and leaves them inside $M_8^0$.

The complete leading descent is therefore a supported total class
$\gamma\in\mathbb H^2(C_9)$.  Under the natural edge map to $H^1(Q_9)$
its image is
\[
                    \bar u^{-1}c((o/\pi)^{(1)}).
\tag{A31i.17}
\]
Here $c$ is the canonical inverse-Cartier map represented by the two
fractions in parentheses in (A31i.16), followed by the map to $Q_9$.
The superscript $(1)$ is base change by Frobenius, not the absolute
Frobenius endomorphism on a coherent cohomology group.

Only the degree-below-two descent and its pairings have been proved
integral.  This does not assert that $\beta$ lies in the full integral
de Rham lattice.  In fact, the valuation bounds give convergent integral
one-form and function cochains, and their residues against any integral
$F^3$ test converge in $R$.  Their reduction is exactly the finite
supported pairing of $\gamma$ with the test's jet class.  The local
degree-two Gysin form pairs to zero with form degrees at least three
on a fourfold.  Thus a rationally zero primitive period has zero reduced
pairing with $\gamma$, without any hypothesis on the full comparison
lattice.

\subsection{Finite-pole noncancellation input}

For sufficiently general high-degree prescribed-centre axes, the map
\[
                     H^1(S^{(1)},N^{(1)})
                           \xrightarrow{c}H^1(Q_9)
\tag{A31i.18}
\]
is injective by the complete finite-pole quotient proof inserted below.
This is stronger than the bounded $Q_3$ assertion of A31f: it controls
the possible kernel through global sections of $Q_9/Q_3$ and retains
the closed-function terms in the larger complex.

% Auxiliary fragment for C44-A31i. No separate finding ID.
% This fragment proves only the finite-pole quotient statement.
% It does not assert a canonical ramified comparison or an e=3 theorem.
\subsection{Enlarging the supported pole complex without cancellation}
\label{subsec:c44-a31i-q9}

\begin{lemma}[Uniform finite-pole quotient lemma]
Let $k=\overline{\mathbf F}_p$, $p>2$, let $S$ be a fixed smooth projective
surface in an abelian sixfold $X$, and fix an ample bundle $H$ on $X$.
Fix an integer $M\geq p$.  For all sufficiently large $t$, there is a
nonempty open set of smooth genuine axes $B\subset X$, cut out by two
members of $|L|$, $L=H^t$, and containing $S$, with the following
property.  Write
\[
 I=I_{S/B},\quad N=N_{S/B},\quad E=N_{S/X},\qquad
 M_m^a=\mathcal Ext_B^2(\mathcal O_B/I^m,\Omega_B^a),
\]
and put
\[
 C_m=[M_{m-1}^0\xrightarrow d M_m^1],\qquad
 Q_m=\operatorname{coker}(M_{m-1}^0\xrightarrow d M_m^1).
\tag{Q9.1}
\]
The complex degrees are zero and one.  Then $Q_p\to Q_M$ is an
injection of sheaves and
\[
                  H^0(Q_M/Q_p)=0.
\tag{Q9.2}
\]
Consequently $H^1(Q_p)\to H^1(Q_M)$ is injective.  This open set can be
intersected with the open set in the ordinary/Cartier proposition of
A31h.  In particular the map
\[
                  H^1(N^{(1)})\xrightarrow c H^1(Q_M)
\tag{Q9.3}
\]
is injective there.

For an additional statement which retains the closed function terms, set
\[
 D=\{\theta\in M_{M-1}^0:d\theta\in M_p^1\},
                \qquad C'=[D\xrightarrow d M_p^1].
\tag{Q9.4}
\]
There is a canonical decomposition
$C'=C_p\oplus\ker(d|M_{M-1}^0)[0]$, and
\[
                 \mathbb H^2(C')\longrightarrow\mathbb H^2(C_M)
                       \quad\hbox{is injective}.
\tag{Q9.5}
\]
All sheaf and complex assertions involving $d$ can equivalently be
made after relative Frobenius pushforward.  The resulting sheaves are
coherent on a finite thickening of $S^{(1)}$ in $B^{(1)}$.
For $M>p$ they are not asserted to be annihilated by the reduced
ideal of $S^{(1)}$.
\end{lemma}

\begin{proof}
\leavevmode\par
The case $M=p$ is immediate, so assume $M>p$.
\paragraph{A polynomial representation embedding.}
If $V$ has dimension two and $R$ is a homogeneous polynomial
$\operatorname{GL}(V)$-representation of degree $d$, there is an
equivariant injection
\[
             R\hookrightarrow
                 \bigl(\operatorname{Sym}^d(V\oplus V)\bigr)^{\oplus r},
                        \qquad r=\dim R.
\tag{Q9.6}
\]
Indeed its polynomial action extends to $\operatorname{End}(V)$.
Send $v$ to the polynomial function $A\mapsto R(A)v$, with values in
the underlying vector space of $R$, now given the trivial action.
Evaluation at the identity proves injectivity.  Right translation
$A\mapsto Ag$ intertwines this map with $v\mapsto R(g)v$.
The degree-$d$ coordinate polynomials of the two rows of $A$, under
this action, form $\operatorname{Sym}^d(V\oplus V)$.  This proves
(Q9.6), in arbitrary characteristic and without semisimplicity.
A subrepresentation and a quotient of homogeneous polynomial
representations are again homogeneous polynomial representations:
stability under all matrices follows by polynomial continuation
from the dense open invertible matrices.

Applied to the subbundle $N\subset E$, (Q9.6) embeds $R(N)$ into
the fixed bundle
$\bigl(\operatorname{Sym}^d(E\oplus E)\bigr)^{\oplus r}$.
If $d>0$, finitely many choices of normal kernel two-planes at
distinct points of $S$ force $H^0(R(N))=0$.  Here is the exact
finite-dimensional argument.  At a point let $P,Q\subset E_x$ be
complementary two-planes.  The spaces
$\operatorname{Sym}^d(P\oplus P)^{\oplus r}$ and
$\operatorname{Sym}^d(Q\oplus Q)^{\oplus r}$ have zero intersection,
by the direct-sum grading of the symmetric algebra.  Thus some
two-plane excludes any given nonzero value in the fixed ambient
fiber.  A nonzero section of the fixed ambient bundle is nonzero
at some point not previously used.  Requiring its value to lie in
the chosen two-plane subspace cuts the remaining section space
properly.  Repetition terminates because that section space is
finite dimensional.

\paragraph{An intrinsic weight filtration.}
In two regular normal parameters $x,y$, a function fraction
$a/(x^\alpha y^\beta)$ has total normal pole
$n=\alpha+\beta$, where $\alpha,\beta\geq1$.  Its leading coefficient
bundle at total pole $n$ is
\[
                  D_n=\det N\otimes\Gamma^{n-2}N.
\tag{Q9.7}
\]
Here $\Gamma$ denotes divided powers.  We filter functions by
$M_{n-1}^0$, so that $n$ is their weight.
On one-forms use the intrinsic filtration
\[
 W_n=M_{n-1}^1+
       \operatorname{image}\bigl(M_n^0\otimes_{\mathcal O_B}
                    \langle dI\rangle_{\mathcal O_B}
                    \longrightarrow\mathcal H_S^2(\Omega_B^1)\bigr).
\tag{Q9.8}
\]
Locally its normal one-forms have total pole at most $n+1$,
and its tangent one-forms have total pole at most $n$.
Multiplication by $I$ lowers this filtration by one.
It is independent of the normal and tangent coordinates: under
$x'=A(z)x+O(I^2)$ the term $dA\cdot x$ has the same weight as
$A\,dx$, while replacing a tangent coordinate by a function with
a normal linear term only lowers the weight of that new normal
one-form term.  Formula (Q9.8) is also a coordinate-free description
of these assertions.

At weight $n$ there is therefore an exact sequence of bundles
\[
 0\longrightarrow T_n=D_n\otimes\Omega_S^1
 \longrightarrow \operatorname{gr}_n W
 \longrightarrow V_n=
       \det N\otimes\Gamma^{n-1}N\otimes N^*
                  \simeq N\otimes\Gamma^{n-1}N
 \longrightarrow0.
\tag{Q9.9}
\]
No splitting of this sequence is chosen.  The leading normal part
of differentiation is the equivariant polynomial map
\[
 \delta_n:D_n\longrightarrow V_n,\qquad
 \frac{a}{x^\alpha y^\beta}\longmapsto
       -\frac{\alpha a\,dx}{x^{\alpha+1}y^\beta}
       -\frac{\beta a\,dy}{x^\alpha y^{\beta+1}}.
\tag{Q9.10}
\]
In particular
\[
                         G_n=\operatorname{coker}\delta_n
\tag{Q9.11}
\]
is a homogeneous polynomial $\operatorname{GL}_2$-bundle of
degree $n$.  Divided powers in (Q9.7)--(Q9.10) are essential;
there is no use of a characteristic-zero symmetric-power
identification in degrees at least $p$.

\paragraph{The exceptional kernels and their tangent differential.}
The kernel $K_n$ of (Q9.10) is zero unless $n=pr$ with $r\geq2$.
Indeed the differentiated monomial in each normal component has
a unique monomial preimage, so both positive denominator exponents
must be divisible by $p$.  For $n=pr$ there is canonically
\[
 K_n=F_S^*K_n',\qquad
              K_n'=\det N^{(1)}\otimes\Gamma^{r-2}N^{(1)}.
\tag{Q9.12}
\]
This is the leading part of Frobenius on local cohomology:
the fraction with denominator exponents $\alpha',\beta'$ on the
twist goes to the fraction with exponents $p\alpha',p\beta'$.
Thus the transition matrices for these local frames are $p$th
powers.  The tangent part of $d$ on $K_n$ is exactly the canonical
Cartier connection
\[
                K_n\longrightarrow K_n\otimes\Omega_S^1,
\tag{Q9.13}
\]
which differentiates only its local coefficient in these frames.
The Frobenius pushforward of its image is
$K_n'\otimes B_S^1$, where
$B_S^1=\operatorname{image}(F_{S*}\mathcal O_S
                              \xrightarrow d F_{S*}\Omega_S^1)$.
The sheaf $F_{S*}\Omega_S^1/B_S^1$ is fixed by $S$.
It is locally free, as follows from etale coordinates and the
usual Cartier complex.

The bundle $D_n/K_n$, when $K_n\ne0$, still has a determinant factor.
More explicitly, after factoring $\det N$ out of (Q9.12), its kernel
inside $\Gamma^{n-2}N$ is
\[
              (\det N)^{p-1}\otimes
                       F_S^*\Gamma^{r-2}N^{(1)}.
\tag{Q9.14}
\]
Both this representation and its quotient are polynomial, of
degree $n-2$.  Consequently $D_n/K_n$ is $\det N$ tensored with
a homogeneous polynomial representation of degree $n-2$.
This statement is also valid when $K_n=0$.

\paragraph{Strictness and the actual quotient.}
First $Q_p\to Q_M$ is injective.  If a function of total pole at
most $M$ has derivative in $M_p^1$, take its highest function pole
$n>p$.  For $n=p+1$ a zero leading normal derivative is impossible,
since $p+1<2p$.  For $n>p+1$ its leading normal derivative and its
tangent derivative both vanish.  Thus its normal exponents are
multiples of $p$ and its leading coefficients are Cartier-horizontal.
They are locally $p$th powers.  The leading part therefore lifts
to an actually closed $p$th-power local-cohomology fraction of
the same pole bound.  Subtract it and decrease the pole.
After finitely many steps the primitive lies in $M_{p-1}^0$.
Since $d$ is injective on that latter module, this also proves the
canonical decomposition asserted in (Q9.4).

For completeness the same argument proves strictness of the image
filtration used to compute
\[
          \mathcal Q=Q_M/Q_p
             =M_M^1/(M_p^1+dM_{M-1}^0).
\tag{Q9.15}
\]
Use (Q9.8) on the target modulo $M_p^1$ and the total-pole
filtration on the source modulo $M_{p-1}^0$.
If a highest graded source term maps to zero, at weight $p+1$
the normal map has zero kernel.  At every higher weight its
normal kernel is (Q9.12), and the kernel of the remaining
Cartier connection consists of horizontal coefficients.
As just proved, these lift locally to actually closed fractions,
not merely to elements with zero graded derivative.  Subtracting
such a lift lowers the source weight without changing its
derivative.  This proves strictness sheaf-locally and rules out
any later differential from a normal kernel to a lower pole.
The Cartier assertion used here is the degree-zero Cartier
isomorphism for a smooth scheme over the perfect field $k$;
the bounded local lifts may be obtained by lifting their
coefficients from $S^{(1)}$ to $B^{(1)}$ and taking Frobenius.

The quotient consequently has the following actual finite
filtration, after Frobenius pushforward.  The weight $p+1$ factor
is $F_{S*}G_{p+1}$, because $M_p^1$ already removes the tangent
part at that weight.  For $p+2\leq n\leq M$, the weight-$n$
factor has a filtration with factors
\[
 \begin{split}
 &F_{S*}G_n,\\
 &F_{S*}\bigl((D_n/K_n)\otimes\Omega_S^1\bigr),\\
 &K_n'\otimes(F_{S*}\Omega_S^1/B_S^1)
                           \quad\hbox{if }K_n\ne0.
 \end{split}
\tag{Q9.16}
\]
Indeed projection in (Q9.9) gives the normal quotient $G_n$;
its kernel is $T_n/dK_n$.  Filtering that kernel by
$(K_n\otimes\Omega_S^1)/dK_n$ gives exactly the last two
terms of (Q9.16), by (Q9.13) and the projection formula.
There is also the final tangent factor
\[
                         F_{S*}(D_{M+1}\otimes\Omega_S^1).
\tag{Q9.17}
\]
If $M=p$, the quotient is zero and these lists are empty.
For $M>p$, all the displayed factors are on the reduced
$S^{(1)}$, although their extensions can live on its finite
thickening.  Cohomology of this finite filtration is unaffected
by that distinction.

\paragraph{Vanishing of all factors.}
The normal sequence and adjunction give
\[
                         \det N=K_S\otimes L^{-2}|S.
\tag{Q9.18}
\]
By (Q9.6) and (Q9.14), each factor in the second line of
(Q9.16), and the factor (Q9.17), embeds before Frobenius
pushforward into a fixed bundle tensored with $L^{-2}|S$.
It has no $H^0$ for all sufficiently large $t$, uniformly in
the axis.  The third line has the same property on $S^{(1)}$:
use (Q9.12), the inclusion $N^{(1)}\subset E^{(1)}$, (Q9.6),
and $\det N^{(1)}=K_{S^{(1)}}\otimes(L^{(1)})^{-2}|S^{(1)}$.
The auxiliary factor $F_{S*}\Omega_S^1/B_S^1$ is fixed.
These vanishing assertions also follow for a fixed coherent
auxiliary factor: all the representation injections are
locally split as injections of bundles, and a fixed coherent
sheaf tensored by a sufficiently negative ample power has
zero $H^0$ when it is torsion free.  Here the auxiliary
Cartier factor is locally free, so no torsion issue occurs.

There remain only the finitely many $F_{S*}G_n$,
$p+1\leq n\leq M$.  Each $G_n$ is polynomial of positive
degree $n$.  The finite normal-plane construction following
(Q9.6) makes every one of their section spaces zero
simultaneously.  The required finitely many normal first
jets are realized by two global sections of $I_{S/X}L$
for large $t$.  General sections with these jets have
surjective normal matrix along $S$: rank drop of a
$2$-by-$4$ matrix has codimension three, greater than
$\dim S=2$.  The off-centre smoothness and genuine-pencil
conditions remain the nonempty open conditions of A06.
Upper semicontinuity of the finitely many bundle section
spaces makes the resulting vanishing an open condition.
This also permits simultaneous imposition of the finite
normal-plane conditions already used in A31h.
Every factor in (Q9.16)--(Q9.17) now has zero $H^0$.
The filtration proves (Q9.2), hence the asserted $H^1$
injections.

Finally the quotient complex $C_M/C'$ has injective
differential by the definition of $D$, and is
quasi-isomorphic to $\mathcal Q[-1]$.  Its degree-one
hypercohomology is $H^0(\mathcal Q)=0$.  The long exact
sequence of this quotient proves (Q9.5).  In particular a
closed degree-zero term in $C'$ cannot cancel a nonzero
$H^1(Q_p)$ component after passing to $C_M$.
\end{proof}

\paragraph{The characteristic-three instance.}
For $p=3$ and $M=9$, the normal factors are $F_{S*}G_n$
for $4\leq n\leq9$; their ranks before Frobenius pushforward
are respectively
\[
                              5,\ 6,\ 8,\ 8,\ 9,\ 12.
\]
Only $K_6=F_S^*(\det N^{(1)})$ and
$K_9=F_S^*(\det N^{(1)}\otimes N^{(1)})$ occur.
Weights $5,\ldots,9$ have the tangent corrections in
(Q9.16), and weight ten is (Q9.17).
Thus $H^1(Q_3)\to H^1(Q_9)$ is injective on the same
nonempty open, and (Q9.5) applies to
$D=\{\theta\in M_8^0:d\theta\in M_3^1\}$.
This is only the algebraic noncancellation statement.
An actual ramified application must separately prove that
its canonical descent lies in this bounded complex.

\paragraph{Other fixed bounds.}
The proof is valid for every fixed $M\geq p$; it has no
ambient restriction $M<4p$.  That restriction enters only
when A31h is used to detect $\mathbb H^2(C_M)$ by primitive
period tests.  Thus the algebraic lemma also applies to
$M=2p$ for $p\geq5$, while making no assertion here about
the canonical ramified descent or its bound.


\begin{theorem}[C44-A31i: the next ramification index in characteristic three]
For all sufficiently high degrees there is a nonempty open of genuine
axes containing the A14 surface such that their nonzero primitive
class $w$ is not filtered on any abstract smooth
projective lift over a DVR of residue characteristic three and absolute
ramification index at most three.  The A18 closed axes chosen in this open have
no algebraic lift of this primitive class on the geometric generic
fibre of any such model.
\end{theorem}

\begin{proof}
Choose the intersection of the nonempty opens of A31f, A31h and the
finite-pole quotient lemma at $M=9$.  The indices one and two are A31f,
so it remains to treat index three.  For a putative filtered $w$,
A29 supplies the relative Albanese embedding
and the full free primitive $F^3$ test lattice.  The rational periods of
$\beta$ on this lattice vanish by the projection formula: its ambient
correction pairs to zero and the remaining primitive part is filtered.
The integral lower-degree calculation above therefore makes the pairing
of $\gamma$ with every primitive test zero.

Termwise Ext/Serre duality and residue integration by parts pair
$\mathbb H^2(C_9)$ perfectly with
\[
 \mathbb H^4[\Omega_B^3/I^9\Omega_B^3
                         \longrightarrow\Omega_B^4/I^8\Omega_B^4].
\tag{A31i.19}
\]
This pairing does not require $d:M_8^0\to M_9^1$ to be injective.
A31h with $M=9<4\cdot3$ supplies the full primitive test surjection.
Consequently $\gamma=0$, so (A31i.17) vanishes.  The proved injection
(A31i.18) gives $o=0$.  A smooth embedded $S_2$ exists, and this is the
initial order $\lfloor3/(3-1)\rfloor+1=2$ required by A31.  Its
all-order induction and algebraization contradict the geometric
nonliftability of the A14 surface.

The A18 class-level assertion follows from the common-realization
kernel and the generic algebraic primitive projector of A29, as in
A31f.  A cycle defined after a finite extension still gives filtration
membership of the fixed vector over the original fraction field, so
that extension does not change the model excluded in the statement.
\end{proof}

\paragraph{Source and programme scope.}
C44-A31-SA1 supplies the canonical arbitrary-ramification rational comparison,
not an integral lattice assertion.  The local normalization used here is
Petrequin, BSMF 131 (2003), Propositions 3.9--3.10, pp.~77--79,
Lemma 3.12, pp.~79--80, Proposition 6.6 and its complete displayed
proof, pp.~109--111, and Lemma 7.1, p.~113.  These particular proofs
were read in full for this candidate; their cited foundational inputs
remain as cited there.
\url{https://www.numdam.org/article/BSMF_2003__131_1_59_0.pdf}.
The full mixed frame/composition and denominator arguments above are
new calculations, not statements attributed to those sources.

For JAC1964 QH9-axis, this targets the next route-(a) initial-support
gate at cote~44, Theorem~5, scan~004 and the primitive caution at
scans~008--009.  The next ramification range has passed the complete
integrated proof checks.  It is not a universal
axis-sign proof, an all-ramification class theorem, or a scan~049
product-form assertion.

% END work/axis/axis_second_ramified_noncancellation.tex

% BEGIN work/axis/axis_next_boundary_odd_characteristic.tex
% C44_POS /root/axis_independent. Author, root and independent whole-proof checks passed.
\section{One further ramification index in odd characteristic}
\label{sec:c44-a31j}

\paragraph{Scope and dependency.}
This is a general-characteristic addendum to the same canonical radial
argument as C44-A31i.  It uses that argument's complete common-graph,
frame-prism and source-unit normalization, and its checked finite-pole
quotient lemma.  It makes no separate staged progress-clock claim and
does not assert existence of a nonliftable centre in every characteristic.

\begin{theorem}[C44-A31j: exclusion through absolute ramification $p$]
Let $p>2$, let $S$ be a smooth geometrically nonliftable projective surface
in an abelian sixfold $X$ over $k=\overline{\mathbf F}_p$, and fix an ample
bundle $H$ on $X$.  For all sufficiently large $t$, there is a nonempty
open of smooth genuine axes $B\subset X$, cut out by two members of
$|H^t|$ and containing $S$, such that the fixed primitive class
\[
                    w=[S]_B-\kappa^*(L_h^2)^{-1}[S]_X,
                    \qquad h=c_1(H^t),
\tag{A31j.1}
\]
is not filtered on any abstract smooth projective lift over a complete
mixed-characteristic DVR with residue field $k$ and absolute ramification
index $e\leq p$.

The A14b surfaces give actual such centres for every prime
$p\equiv3\pmod4$.  A18 can select closed axes in this open with
$A_B=\mathbf Qw$ and $w^2>0$; their nonzero primitive etale class cannot
specialize from an algebraic class on the geometric generic fibre of any
such model.  The new equality cases beyond A31h include $(p,e)=(7,7)$
and $(11,11)$.  No index $e>p$ is excluded here.
\end{theorem}

\begin{proof}
The case $p=3$ follows from the proof of A31i, which uses only the
smoothness and geometric nonliftability of its prescribed centre;
the coherent and genericity lemmas there hold for every fixed smooth
$S$.  All cases $e\leq p-1$ are A31h.  Assume
$p\geq5$ and $e=p$, and write $p=u\pi^p$ with $u\in R^\times$.
Choose the nonempty open where A31h's primitive jet restriction at
$M=2p$ and the finite-pole quotient lemma in A31i at the same $M$
hold.  These conditions can be imposed together with the smoothness,
genuine-pencil and ordinary/Cartier conditions of A31h.

The canonical graph and radial comparison construction in A31i is over
characteristic zero and uses no special property of residue characteristic
three until its valuations and leading-form reduction.  Its exact
nested-augmentation and mixed-frame identities therefore apply unchanged.
In particular it constructs the actual supported rational cycle class,
not a newly selected transport.  Its two-contraction coefficients have
denominators at most $nj$, with $1\leq j<n$.  Evaluation at the
$\pi$-multiple displacements gives the bound
\[
                   n-p\bigl(v_p(n)+\lfloor\log_p(n-1)\rfloor\bigr).
\tag{A31j.2}
\]
For $n\geq2$ this is nonnegative, with equality exactly at $n=p,2p$.
Here is a complete verification.  For $n<p$ the subtracted quantity
is zero.  At $n=p$ the expression is zero.  For $p<n<p^2$, if
$p\nmid n$ it is $n-p>0$; if $n=ap$, $2\leq a\leq p-1$, it is
$p(a-2)$, zero only at $a=2$.  At $n=p^2$ it equals $p(p-3)>0$.
For $p^a<n<p^{a+1}$ with $a\geq2$, it is at least
$p^a-2pa>0$: this holds at $a=2$ because $p\geq5$, and the
difference increases thereafter.  The same bound proves positivity
at every remaining power $p^a$, and shows convergence to infinite
valuation.  A single contraction has valuation at least
$n-pv_p(n)$, whose unique zero for $n\geq1$ is $n=p$.

The normal-pole count of A31i now places the actual leading function
terms in $M_{2p-1}^0$, of total denominator at most $2p$.  The leading
one-form terms lie in $M_p^1$, of total denominator at most $p+1$.
All closed function terms, including the product of two leading
logarithms, are retained.  The leading total descent is consequently
a class
\[
                 \gamma\in\mathbb H^2(C_{2p}),\qquad
                 C_{2p}=[M_{2p-1}^0\to M_{2p}^1].
\tag{A31j.3}
\]
Only its degree-below-two cochains and their primitive periods have
been proved integral; no full integral crystalline--de Rham lattice
comparison is asserted.

We compute the leading one-form, including variable coefficients.
For a normal displacement $\pi(a,b)$, the weight-$p$ part of radial
contraction, reduced in characteristic $p$, is $\bar u^{-1}A_p$, with
\[
\begin{split}
 A_p={}&(a\,dy-b\,dx)
       \sum_{i+j=p-1}\frac{a^i b^j}{x^{i+1}y^{j+1}}\\
 &-(a\,db-b\,da)
       \sum_{i+j=p-2}\frac{a^i b^j}{x^{i+1}y^{j+1}}.
\end{split}
\tag{A31j.4}
\]
The signs follow because $p-1$ is even and $p-2$ is odd.  There is
the exact identity
\[
 A_p-\left(\frac{a^pdy}{x^py}-\frac{b^pdx}{xy^p}\right)
    =d\left(\sum_{r=1}^{p-1}
                   \frac{a^r b^{p-r}}{r x^r y^{p-r}}\right).
\tag{A31j.5}
\]
Indeed differentiation of the $r$th summand gives its $da$ coefficient
$a^{r-1}b^{p-r}/(x^ry^{p-r})$ and its $db$ coefficient
$-a^rb^{p-r-1}/(x^ry^{p-r})$, since $(p-r)/r=-1$ in $k$.
Its $dx$ coefficient is $-a^rb^{p-r}/(x^{r+1}y^{p-r})$, and its
$dy$ coefficient is $a^rb^{p-r}/(x^ry^{p-r+1})$.
Summing these four coefficient formulas gives precisely the mixed
terms on the left of (A31j.5); the two pure terms have been subtracted.
Thus the identity includes all $da,db$ terms.  Its primitive has total
normal denominator $p$ and belongs to $M_{p-1}^0$.

Let $o$ be the actual obstruction to the specified smooth embedded
lift $S_2$.  After the indicated bounded pair coboundary is subtracted,
the image of (A31j.3) in $H^1(Q_{2p})$ is
\[
                         \bar u^{-1}c((o/\pi)^{(1)}).
\tag{A31j.6}
\]
The map $H^1(N^{(1)})\to H^1(Q_{2p})$ is injective by A31h and
the complete finite-pole quotient lemma of A31i.  This conclusion
does not assume Frobenius injectivity on $H^1(\mathcal O_S)$.

Suppose (A31j.1) were filtered.  A29 gives the relative Albanese
embedding and its free full primitive test lattice for the arbitrary
abstract model.  The rational Gysin periods vanish on that lattice;
the actual convergent lower-degree descent makes these periods
integral, with reduction equal to the finite supported pairing with
$\gamma$, exactly as in A31i.  Residue/Serre duality pairs
$\mathbb H^2(C_{2p})$ perfectly with
\[
 \mathbb H^4[\Omega_B^3/I^{2p}\Omega_B^3
                         \to\Omega_B^4/I^{2p-1}\Omega_B^4].
\tag{A31j.7}
\]
Since $2p<4p$, A31h says that primitive tests surject onto this
factor.  Hence $\gamma=0$, and the injection applied to (A31j.6)
forces $o=0$.  A smooth $S_2$ exists.  The required initial order
in A31 is
\[
                         \left\lfloor\frac p{p-1}\right\rfloor+1=2.
\tag{A31j.8}
\]
That theorem extends the support through all orders and algebraizes
it, contradicting the geometric nonliftability of $S$.

Finally A14b supplies the stated actual centres.  The additional
conditions form a nonempty open of the same prescribed-centre system,
so the A18 closed-point construction applies there.  Its common
realization kernel and the generic primitive projector of A29
convert a generic algebraic lift of the primitive etale class into
filtration membership of (A31j.1).  As in A31i, passage to a finite
field of definition of that cycle does not change filtration
membership of the fixed vector over the original fraction field.
This proves the class-level assertion.
\end{proof}

\paragraph{Exact dependency and programme statement.}
The canonical rational Gysin comparison is the checked C44-A31-SA1;
the complete local graph, actual Cech composition, frame normalization
and radial denominator construction are A31i.  The only new computations
in this addendum are (A31j.2)--(A31j.5), proved above, and their use
with the already proved finite-pole and primitive-jet results.
No additional external comparison theorem is invoked.

For JAC1964 QH9-axis and the H9-T7 primitive margin of cote~44,
Theorem~5, scan~004 and scans~008--009, this further restricts
route~(a) for actual nonzero primitive classes.  It is not a
universal sign theorem or an all-ramification nonliftability theorem,
and it gives no scan~049 product-form assertion.  The argument uses
neither of the unrefereed 032/Milne2026 preprints.

% END work/axis/axis_next_boundary_odd_characteristic.tex

% BEGIN work/axis/axis_closed_layer_peeling.tex
% C44_POS /root/axis_independent. Author, root and independent whole-proof checks passed; see dated registers.
\section{Removing closed function layers at the next ramification index}
\label{sec:c44-a31k}

\paragraph{Scope.}
This is a continuation of the same lifting argument as A31f--A31j.
The new step removes actual closed function layers before reading the
Cartier obstruction.  It neither assumes an initial global support lift
nor assumes flatness of supported de Rham hypercohomology.  It makes no
separate staged progress-clock claim.

\begin{theorem}[C44-A31k: exclusion through ramification $p+1$]
Let $p>2$, let $S$ be a smooth geometrically nonliftable projective surface
in an abelian sixfold $X$ over $k=\overline{\mathbf F}_p$, and let $H$ be
an ample bundle on $X$.  For all sufficiently large $t$, there is a
nonempty open of smooth genuine axes $B\subset X$ cut out by two members
of $|H^t|$ and containing $S$, such that
\[
              w=[S]_B-\kappa^*(L_h^2)^{-1}[S]_X,
              \qquad h=c_1(H^t),
\tag{A31k.1}
\]
is not filtered on any abstract smooth projective lift over a complete
mixed-characteristic DVR with residue field $k$ and absolute ramification
index $e\leq p+1$.

The A14b centres give actual instances for every $p\equiv3\pmod4$.
The A18 construction can select closed axes in this open with
$A_B=\mathbf Qw$ and $w^2>0$.  On these axes the nonzero primitive
etale class of $w$ cannot specialize from an algebraic class on the
geometric generic fibre of any such model.  In particular the new
indices include $(p,e)=(3,4),(7,8),(11,12)$.  No assertion for
$e>p+1$ is made here.
\end{theorem}

\subsection{The exact finite-pole inputs}

For a smooth axis containing $S$, set
\[
 I=I_{S/B},\qquad
 M_m^a=\mathcal Ext_B^2(\mathcal O_B/I^m,\Omega_B^a),\qquad
 C_M=[M_{M-1}^0\xrightarrow d M_M^1],\qquad
 Q_M=\operatorname{coker}d.
\tag{A31k.2}
\]
The complex degrees are zero and one; its support shift is two.
Use
\[
                    M=2p\quad(p\geq5),\qquad M=9\quad(p=3).
\tag{A31k.3}
\]
In either case $M<4p$.  Choose the nonempty open where the primitive
jet theorem A31h and the complete finite-pole quotient lemma in A31i
hold at this $M$, together with A31h's ordinary/Cartier conditions.
There are only finitely many such open conditions.  Put
\[
       \mathcal Z_M=\ker(d:M_{M-1}^0\longrightarrow M_M^1).
\tag{A31k.4}
\]
After Frobenius pushforward these are coherent sheaves on a finite
thickening of $S^{(1)}$.  Two consequences of the checked inputs are
\[
\begin{split}
 &H^2(\mathcal Z_M)\longrightarrow\mathbb H^2(C_M)
                                      \quad\hbox{injective},\\
 &H^1(N_{S/B}^{(1)})\xrightarrow c H^1(Q_M)
                                      \quad\hbox{injective}.
\end{split}
\tag{A31k.5}
\]
The first is the restriction of A31i's injection
$\mathbb H^2(C_p\oplus\mathcal Z_M[0])\to\mathbb H^2(C_M)$.
It is stronger than an assertion about the edge map to $H^1(Q_M)$.
Also the primitive test lattice, on any abstract model as in A29,
surjects after reduction onto the perfect residue/Serre dual of
$\mathbb H^2(C_M)$ by A31h.  Thus vanishing of all reduced primitive
periods forces a supported class in $\mathbb H^2(C_M)$ to be zero.

\subsection{Closed fractions over a nilpotent characteristic-$p$ base}

\begin{lemma}[Canonical closed lifts of the finite-pole function sheaf]
Let $R$ be a complete mixed-characteristic DVR with perfect residue field
$k$, uniformizer $\pi$, and $e=v_\pi(p)$.  Let $\mathcal B/R$ be smooth
with special fibre $B$, and let $S\subset B$ be smooth of codimension two.
For $1\leq n\leq\min(p,e)$, write $A_n=R/\pi^n$ and
$\mathcal B_n=\mathcal B\times_R A_n$.
There is a canonical closed-function copy
\[
                         \mathcal Z_M\otimes_k A_n
\tag{A31k.6}
\]
in the actual supported relative de Rham complex of $\mathcal B_n$,
for the values of $M$ in (A31k.3).  It is flat over $A_n$.
It lies in the function pole bound $M$ measured using any local smooth
graph lift of $S$.  It requires no global smooth $S_n$.
Every Cech cochain in $\mathcal Z_M$ therefore lifts to a closed cochain
modulo $\pi^n$ with this same pole bound.
\end{lemma}

\begin{proof}
Since $n\leq e$, the ring $A_n$ has characteristic $p$; since $n\leq p$,
its maximal ideal has $p$th power zero.  It has a canonical coefficient
section $s:k\to A_n$:
\[
                   s(a)=\widetilde{a^{1/p}}^{\,p}.
\tag{A31k.7}
\]
Changing the lift changes the right side by an element of $(\pi^p)=0$.
The characteristic-$p$ binomial formula proves additivity and
multiplicativity, and reduction gives $a$.  Moreover
$F_{A_n}=s\circ F_k\circ\operatorname{red}$.
Consequently the relative Frobenius twist is canonically
\[
             \mathcal B_n^{(1)}
                    =B^{(1)}\times_{k,s}A_n.
\tag{A31k.8}
\]
The subscheme $S^{(1)}\times_k A_n$ thus exists in the twist independently
of the existence of a smooth support lift inside $\mathcal B_n$.

For a smooth scheme over any characteristic-$p$ base, relative
Frobenius is finite flat, and degree-zero Cartier identifies the
structure sheaf of the twist with the closed functions in the relative
de Rham complex.  These assertions hold over the nonreduced base $A_n$.
Here one can check them in etale coordinates: the polynomial algebra
is free over its relative $p$th-power subalgebra on the monomials with
exponents between zero and $p-1$, and differentiation kills exactly
the constant exponent class in each variable.  The relative Frobenius
square for an etale map is cartesian, so this calculation descends.

The same identification holds with local support cohomology.  Indeed
the support in the twist is defined by a regular sequence of length
two.  Its flat Frobenius pullback is also such a sequence.  Local
cohomology of each locally free term, and of each Cartier cohomology
sheaf, is concentrated in support degree two.  The two spectral
sequences for supported de Rham cohomology therefore identify
\[
 \mathcal H^2_{S^{(1)}\times A_n}
                    (\mathcal O_{\mathcal B_n^{(1)}})
 \;\simeq\;
 \ker\!\left(\mathcal H^2_S(\mathcal O_{\mathcal B_n})
                    \xrightarrow d
             \mathcal H^2_S(\Omega^1_{\mathcal B_n/A_n})\right).
\tag{A31k.9}
\]
Here $S$ on the right denotes the underlying closed support, not a
chosen smooth lifting.

Set $s_M=M/p$, so $s_M=2$ for $M=2p$ and $s_M=3$ for $(p,M)=(3,9)$.
The submodule of the left side of (A31k.9) with total normal pole at
most $s_M$ is the constant base change of
\[
 \mathcal Ext^2_{B^{(1)}}\!
       \left(\mathcal O_{B^{(1)}}/(I^{(1)})^{s_M-1},
                                  \mathcal O_{B^{(1)}}\right).
\tag{A31k.10}
\]
Its Frobenius image is precisely (A31k.6).  In normal coordinates it
sends $a/(x'^\alpha y'^\beta)$ to the relative Frobenius fraction with
denominator $x^{p\alpha}y^{p\beta}$.  Conversely a closed fraction has
normal denominator exponents divisible by $p$ and Cartier-horizontal
tangent coefficients.  This proves both the stated pole bound and the
identification with $\mathcal Z_M$ on the special fibre.

Every local graph lift has normal generators congruent modulo $\pi$ to
the same special generators.  Their $p$th powers agree modulo $\pi^n$:
$(f+\pi a)^p=f^p$ in $A_n$.  The Frobenius support ideal and all the
fractions from (A31k.10) therefore have the indicated bound in each
local graph chart.  This is also immediate from the canonical relative
Frobenius map; it does not choose a splitting of the formal neighbourhood.
The source (A31k.10) is a constant base change, hence is $A_n$-flat.
This proves the lemma, including compatibility with restrictions,
coordinate changes and Cech coboundaries.
\end{proof}

\paragraph{Local division with a controlled pole bound.}
We also need a cochain assertion, rather than a flatness assertion about
hypercohomology.  In a completed smooth graph chart with tangent
coordinates $t_1,t_2$ and normal coordinates $x,y$, the bounded
function module has the negative monomial basis
\[
                    x^{-\alpha}y^{-\beta},
       \qquad \alpha,\beta\geq1,\quad \alpha+\beta\leq M
\tag{A31k.11}
\]
over the smooth tangent coefficient ring.  The bounded form modules
have the corresponding bases with $dt_i,dx,dy$.  These are flat over
$A_n$, and their inclusions into the full supported fraction modules
have flat successive pole quotients.  Thus a bounded cochain reducing
to zero modulo $\pi$ is divisible by $\pi$ within that bounded module;
after division its reduction modulo $\pi^{n-1}$ retains the same bound.
The etale version follows by flat descent from these completed
coordinate calculations.  Only such local cochain modules are being
used; supported hypercohomology is not asserted to be flat.

A cochain supplied by (A31k.6) can also be lifted to an actual convergent
cochain on the characteristic-zero tubular comparison used in A31i.
Locally lift its fractions as finite sums of $R$-scalar multiples of
$p$th powers of fractions.  Their relative differential is divisible
by $p$, hence by $\pi^n$.  On overlaps their reductions are precisely
the canonical closed cochain (A31k.6).  The local expansions on the
common graph charts converge on the same strict neighbourhoods as
in A31i.  Their terms beyond the required truncation vanish modulo
$\pi^n$.  The actual total coboundary of such a lift is consequently
integral, and modulo $\pi^n$ it has only the bounded function component
given by its Cech coboundary.

\subsection{The removal of a closed layer}

Here is the exact operation used below.  Suppose an actual rational
supported descent has been multiplied by a power of $\pi$ so that its
degree-one and degree-zero cochains $(\eta,\kappa)$ are integral.
Suppose, modulo $\pi^n$ for some $2\leq n\leq\min(p,e)$, their poles
are bounded by $C_M$, all their primitive periods vanish, and
$\eta\equiv0\pmod\pi$.  Its special total class is then represented
by the Cech two-cocycle
\[
                    \kappa_0\in\check Z^2(\mathcal Z_M).
\tag{A31k.12}
\]
The equality $d\kappa=\delta\eta$ and the exact composition coherence
of A31i give its closedness and Cech cocycle condition.  Vanishing
primitive periods and A31h give zero in $\mathbb H^2(C_M)$.
By the first injection (A31k.5), its class is zero in
$H^2(\mathcal Z_M)$ itself.  On an affine adapted Cech refinement,
choose
\[
                  b_0\in\check C^1(\mathcal Z_M),
                  \qquad \delta b_0=\kappa_0.
\tag{A31k.13}
\]
This is a sheaf-cohomology coboundary, not a choice of a smooth $S_2$.

Lift $b_0$ to the closed cochain (A31k.6) modulo $\pi^n$, and then
to an actual convergent cochain as above.  Subtract its actual total
coboundary, with the sign fixed by the total differential.  The
function component becomes zero modulo $\pi$, while the one-form
component is unchanged modulo $\pi^n$.  Both components are now
divisible by $\pi$.  Divide them, and the represented rational class,
by $\pi$.  The local saturation calculation (A31k.11) preserves the
bound $C_M$ modulo $\pi^{n-1}$.  The canonical closed module ensures
this simultaneously in every common graph chart, so the reduction
is a well-defined supported total class, not unrelated local quotients.

Subtracting a genuine total coboundary does not change any rational
period; dividing multiplies it by $\pi^{-1}$.  The new lower cochains
are integral and still have zero primitive periods.  Their reduction
pairing with the special jet complex is the actual residue pairing.
The original degree-two local Gysin cochain still contributes zero
against form degrees at least three on a fourfold, regardless of its
denominators.  This is exactly the lower-degree period argument of
A31i, now applied after a cochainwise operation.  No integral comparison
of the full crystalline and de Rham lattices has been used.

The operation can be repeated if the resulting one-form component is
still zero modulo $\pi$.  At the next step one uses the canonical
closed module over $A_{n-1}$.  All higher corrections discarded at a
stage are divisible by the full current power of $\pi$; after division
they remain zero modulo the new current power.

\subsection{Exact coefficients through the target order}

\begin{proof}[Proof of the theorem]
A31j covers $e\leq p$, so consider $e=p+1$ and write $p=u\pi^{p+1}$.
Use the actual common-graph, frame-prism and source-unit comparison of
A31i, with its rational identification C44-A31-SA1.  This construction
uses only local smooth support lifts.  Its leading normal Cech
displacements are $\pi(a,b)$, and their class is the actual first
support obstruction $o\in H^1(N_{S/B})$ after identifying
$(\pi)/(\pi^2)$ with $k$.

For a one-contraction term of displacement weight $n$, the valuation
is at least $n-e v_p(n)$.  For a two-contraction function term it is
at least
\[
             v_e(n)=n-e\bigl(v_p(n)+\lfloor\log_p(n-1)\rfloor\bigr),
                         \qquad n\geq2.
\tag{A31k.14}
\]
The bound uses the actual denominators $nj$ with $j<n$ proved in A31i,
for composition and frame changes, including ordered-to-radial
normalization.  It applies to all the function cochains, not just to
the product of two particular logarithms.  Terms of weight $n$ have
total normal function pole at most $n$ and normal one-form pole at
most $n+1$.

\paragraph{The case $p\geq5$.}
For $e=p+1$, the only negative values in (A31k.14) are
\[
                          v_e(p)=-1,\qquad v_e(2p)=-2.
\tag{A31k.15}
\]
For $n<p$ the subtracted term is zero.  If $p<n<p^2$ and $p\nmid n$,
the value is $n-(p+1)\geq0$.  If $n=ap$, $2\leq a<p$, it is
$p(a-2)-2$, negative only at $a=2$.  At $n=p^2$ it equals
$p^2-3(p+1)>0$.  On every interval $p^a<n<p^{a+1}$ with $a\geq2$,
it is greater than or equal to $p^a-2(p+1)a>0$; the latter is
positive already at $(p,a)=(5,2)$ and increases thereafter.
The same estimate covers all remaining power endpoints.
It also tends to infinity.  The one-contraction bound is negative
only at $n=p$, where it is $-1$.

Multiply the entire rational descent by $\pi^2$ and work through
order one, that is, modulo $\pi^2$.  The function cochains can
involve only weights $p,2p$ at these orders; every other weight
has scaled valuation at least two.  The one-form cochains start
at order one, with only the weight-$p$ term:
\[
        \eta'=\pi^2\eta
             \equiv\pi\,\bar u^{-1} A_p(a,b)\pmod{\pi^2}.
\tag{A31k.16}
\]
Here the variable-coefficient expression $A_p$ is (A31j.4), and its
class modulo exact forms is the inverse-Cartier expression
$a^pdy/(x^py)-b^pdx/(xy^p)$ by the fully differentiated identity
(A31j.5).  Its exact primitive has function pole $p$.
Thus the entire lower descent through the target order lies in
the local bound $C_{2p}$, including every closed function term.

Apply the closed-layer operation once with $n=2$ and $M=2p$.
After subtraction of an actual boundary and division by $\pi$,
the reduced supported class $\gamma$ has edge
\[
                \gamma\longmapsto
                   \bar u^{-1}c(o^{(1)})\in H^1(Q_{2p}).
\tag{A31k.17}
\]
All its primitive periods are zero if (A31k.1) is filtered.
The primitive jet theorem therefore gives $\gamma=0$, and the
second injection (A31k.5) forces $o=0$.  A smooth embedded $S_2$
exists.  Since
\[
                 \left\lfloor\frac{p+1}{p-1}\right\rfloor+1=2
                       \qquad(p\geq5),
\tag{A31k.18}
\]
this is precisely the initial order required by A31.

\paragraph{The case $p=3,e=4$.}
Now the negative values of (A31k.14) occur exactly at
\[
                    n=3,6,9,\qquad
                    v_4(n)=-1,-2,-3,\ \hbox{respectively}.
\tag{A31k.19}
\]
The values through $12$ follow directly from (A31k.14).
For $13\leq n\leq26$, splitting according to $v_3(n)=0,1,2$
gives lower bounds $5,3,2$, respectively.  For $n\geq27$, the
interval estimate $3^a-8a>0$ for $a\geq3$ covers all remaining
terms, including the power endpoints.  The one-contraction bound
$n-4v_3(n)$ is negative only at $n=3$; at $n=9$ it is already $1$.

Multiply the actual descent by $\pi^3$ and retain terms modulo
$\pi^3$.  The function terms through order two have weights only
$3,6,9$, hence have total pole at most nine.  In particular weight
$12$ starts at order three and disappears; weight $4$ also starts
at order three.  The one-form term begins at order two:
\[
                    \eta'=\pi^3\eta
                    \equiv\pi^2\bar u^{-1}A_3(a,b)
                                                    \pmod{\pi^3}.
\tag{A31k.20}
\]
These statements control every coefficient through the target order,
not merely the smallest valuation of the series.

Use the closed-layer operation with $M=9$ first over $A_3$.
The closed module here is the Frobenius image of the source
function fractions of total pole at most three; it includes both
the closed pole-six and pole-nine terms.  After one boundary
subtraction and division, the one-form cochain is unchanged
through order one and equals
$\pi\bar u^{-1}A_3$ modulo $\pi^2$.
Use the operation a second time over $A_2$.
After the second division the reduced edge is
\[
                         \bar u^{-1}c(o^{(1)})\in H^1(Q_9).
\tag{A31k.21}
\]
At each step the complete lower descent stays within $C_9$ at the
current precision.  The higher terms discarded modulo $\pi^3$
cannot reappear at either later reduction.  The first primitive
jet test and the $H^2(\mathcal Z_9)$ injection remove the first
closed layer; the same two results remove the second.  The final
primitive jet test and Cartier injection force $o=0$.
This proves existence of a smooth embedded $S_2$.

For $e=4$ the A31 initial order is three.  Apply the already checked
scaled equality-boundary corollary of A31h (or A31g), with
$e=(p-1)m$ and $m=2$, to this actual $S_2$.
Under the same filtration hypothesis it extends to a smooth $S_3$.
This last step explicitly uses the supplied $S_2$ just proved to
exist.  It does not assume it at the start of the argument.

In both cases A31 extends the supplied initial support through all
orders and algebraizes it, contradicting the geometric nonliftability
of $S$.  A29 supplies the relative Albanese embedding and the free
full primitive test lattice for every abstract model in the theorem,
so the argument did not presuppose an embedded or complete-intersection
presentation of that model.

Finally intersect the chosen nonempty open with the A18 closed-point
selection open.  A14b provides the stated actual centres and A06
provides the positive norm.  The common realization kernel of A18
and the algebraic generic primitive projector of A29 convert any
generic algebraic lift of the primitive etale class into filtration
membership of the fixed class (A31k.1).  If the cycle is defined only
after a finite extension, membership of this fixed vector in the
base-changed filtration descends to the original fraction field.
Thus the class-level conclusion follows from the model-level
filtration contradiction.
\end{proof}

\paragraph{Primary source locators and exact dependency.}
Relative Frobenius over a general smooth characteristic-$p$ base is
finite syntomic by Stacks Project, Lemma~50.19.5, tag~0FW2
(\url{https://stacks.math.columbia.edu/tag/0FW2}); the etale Frobenius
square is Lemma~41.14.3, tag~0EBS
(\url{https://stacks.math.columbia.edu/tag/0EBS}).
The degree-zero Cartier calculation over an arbitrary characteristic-$p$
base is also contained in Esnault--Viehweg, \emph{Lectures on Vanishing
Theorems}, 1992, Section~9.13, Theorem~9.14 and Claim~9.15, author
pages~101--103, with the base setup on page~93
(\url{https://page.mi.fu-berlin.de/esnault/books/esvibuch.pdf}).
The nilpotent-base factorization (A31k.7)--(A31k.8), supported
spectral-sequence passage, bounded closed module and the actual
cochain operations are proved here.  The canonical rational cycle
comparison and all its graph/frame coherence remain the checked
C44-A31-SA1 and A31i; no new integral comparison theorem is invoked.

For JAC1964 QH9-axis and the H9-T7 primitive margin of cote~44,
Theorem~5, scan~004 and scans~008--009, this further restricts route~(a)
for actual nonzero primitive classes.  It is not a universal sign
theorem, an all-ramification theorem, or a scan~049 product-form
assertion.  It uses neither of the unrefereed 032/Milne2026 preprints.

% END work/axis/axis_closed_layer_peeling.tex

% BEGIN work/axis/axis_supplied_support_threshold.tex
% C44_POS /root/axis_independent. Author, root and independent whole-proof checks passed; see dated registers.
\section{A smaller supplied-support threshold in every ramification degree}
\label{sec:c44-a31l}

\paragraph{Scope.}
This addendum improves the finite-order hypothesis in A31.
It does not assert that the required initial support exists.
It belongs to the same developing lifting argument as A31f--A31k
and makes no separate staged progress-clock claim.

\begin{theorem}[C44-A31l: supplied support at order $\lceil e/p\rceil$]
Let $p>2$, let $S$ be a fixed smooth surface in an abelian sixfold
$X$ over $k=\overline{\mathbf F}_p$, and fix an ample bundle $H$.
For all sufficiently large $t$, there is a nonempty open of smooth
genuine axes $B\subset X$ containing $S$ and cut out by two members
of $|H^t|$ with the following property.

Let $\mathcal B/R$ be any abstract smooth projective lift of $B$,
where $R$ is a complete mixed-characteristic DVR with residue field
$k$, uniformizer $\pi$ and $e=v_\pi(p)$.
Suppose the fixed primitive class
\[
                w=[S]_B-\kappa^*(L_h^2)^{-1}[S]_X
                \qquad(h=c_1(H^t))
\tag{A31l.1}
\]
is filtered.  If a smooth embedded lift
$S_m\subset\mathcal B_{R/\pi^m}$ of $S$ is supplied and
\[
                         (p-1)m\leq e\leq pm,
\tag{A31l.2}
\]
then this particular $S_m$ extends to a smooth embedded $S_{m+1}$.

Consequently a supplied smooth $S_{\lceil e/p\rceil}$ extends through
all orders and algebraizes to a smooth projective lift of $S$.
If $S$ is geometrically nonliftable, no such model can satisfy both
the filtration condition and this supplied-support hypothesis.
The conclusion applies to the actual A14b centres and their A18
closed axes with $A_B=\mathbf Qw$; a generic algebraic lift of the
primitive class on a candidate model must therefore fail this
explicit initial-support condition.
\end{theorem}

\begin{proof}
Choose the fixed generic conditions of A31h and the finite-pole
quotient lemma in A31i at
\[
                     M=2p\ (p\geq5),\qquad M=9\ (p=3),
\tag{A31l.3}
\]
together with the ordinary/Cartier conditions at $M=p$.
They are independent of $e$ and $m$.  A29 supplies the relative
Albanese embedding and the full free primitive test lattice for
the abstract model.  Write
\[
                      p=u\pi^e,\qquad r=pm-e,\qquad 0\leq r\leq m.
\tag{A31l.4}
\]

\paragraph{The actual source class and the specified support.}
On an affine adapted cover, extend the given smooth $S_m$ to local
smooth formal supports.  Smoothness of the ambient formal scheme
and of the local support permits these local extensions; compatibility
between them is not asserted.  In the common tangent graph charts
of A31i their normal displacements agree modulo $\pi^m$, because
they extend the same specified $S_m$.  Hence the comparison
displacements are $\pi^m(a,b)$, and their first coefficients represent
the actual obstruction to extending this $S_m$, in
\[
                    H^1(N_{S/B})\otimes(\pi^m)/(\pi^{m+1}).
\tag{A31l.5}
\]
The common-graph, actual Cech composition and frame-prism construction
of A31i therefore applies with every displacement divisible by
$\pi^m$.

Its rational cycle is still the fixed class of the reduced $S$.
Indeed each local formal smooth support has special fibre $S$ with
multiplicity one, and the local source unit is sent to its regular
sequence Gysin fraction $df\wedge dg/(fg)$.
The common graph evaluations and pair-preserving frame changes
preserve that sheaf-level fundamental section.  The actual rational
comparison C44-A31-SA1 identifies its reduction with $[S]$.
Agreement to order $m$ changes only the size of the displacements,
not this source unit or its normalization.
We are not taking the fundamental cycle of the Artin scheme $S_m$
as a cycle on the special fibre; doing so would introduce the
irrelevant length $m$.  There is no factor $m[S]$ in this comparison.

\paragraph{Uniform all-weight bounds.}
Put
\[
 A(n)=v_p(n)+\lfloor\log_p(n-1)\rfloor\quad(n\geq2),\qquad
 b_p(n)=n-pA(n).
\tag{A31l.6}
\]
The two-contraction terms in the actual comparison have valuation at
least
\[
       mn-eA(n)=m\,b_p(n)+rA(n).
\tag{A31l.7}
\]
The exact denominator restriction $nj$ with $j<n$, for both
composition and frame terms, is the one proved in A31i.
For every odd $p$ one has $b_p(n)\geq0$.
Its zeros are exactly $p,2p$ when $p\geq5$, and $3,6,9$ when $p=3$.
Here is a full elementary check.  Below $p$ the value is $n$;
at $p$ it is zero.  Between $p$ and $p^2$, a nonmultiple of $p$
has value $n-p>0$, while $n=ap$, $2\leq a<p$, has value
$p(a-2)$.  At $p^2$ the value is $p(p-3)$.
For $p\geq5$ and $p^a<n<p^{a+1}$ with $a\geq2$, it is at least
$p^a-2pa>0$; the inequality holds at $a=2$ and increases thereafter.
The same bound covers all remaining power endpoints.
For $p=3$, the values through $12$ can be evaluated directly, with
the additional zero at $9$.
For $13\leq n\leq26$, the cases $v_3(n)=0,1,2$ give bounds
$7,6,6$, respectively.  For $n\geq27$, the interval bound
$3^a-6a>0$ for $a\geq3$ finishes the verification.
These bounds also tend to infinity, so the actual series converge.

For one-contraction terms the corresponding identity is
\[
       mn-e v_p(n)=m\bigl(n-pv_p(n)\bigr)+r v_p(n).
\tag{A31l.8}
\]
The integer in parentheses is nonnegative and vanishes exactly at
$n=p$: if $v_p(n)=a\geq1$, then $n\geq p^a$, and
$p^a-pa>0$ for $a\geq2$; for $a=1$ equality occurs only at $p$.
If $p\nmid n$, it is the positive integer $n$.

\paragraph{The interior range $0<r<m$.}
Divide the actual lower descent by $\pi^r$.
Equations (A31l.7)--(A31l.8) show that its one-form reduction has only
weight $p$.  The function reduction also has only weight $p$:
at weight $2p$ its valuation is $2r>r$, and at the additional
weight $9$ for $p=3$ it is $3r>r$.
Every weight with $b_p(n)>0$ has valuation at least $m>r$.
Thus the complete reduced lower descent belongs to
\[
                   \mathbb H^2(C_p),\qquad
                   C_p=[M_{p-1}^0\xrightarrow d M_p^1].
\tag{A31l.9}
\]
The general variable-coefficient identity (A31j.5), including
$da,db$, identifies its edge with
$\bar u^{-1}c(o_m^{(1)})$, where $o_m$ is (A31l.5) after
identifying $(\pi^m)/(\pi^{m+1})$ with $k$.
No ordinary weight-one contribution survives, because its valuation
is $m>r$.

The lower cochains and their primitive pairings are integral after
this division; the degree-two Gysin cochain again pairs to zero
against $F^3$ on the fourfold.  Since (A31l.1) is filtered, all
these primitive periods vanish.  A31h's perfect finite-jet testing
at $M=p$ forces (A31l.9) to be zero.  Its Cartier injection forces
$o_m=0$.  Therefore the specified $S_m$ extends to a smooth
$S_{m+1}$.

\paragraph{The two endpoints.}
If $r=m$, then $e=(p-1)m$, and this is exactly the supplied-support
equality-boundary corollary of A31h.  That result retains both the
ordinary and Cartier contributions and uses their joint injection.
It applies to the actual $S_m$ supplied here.

If $r=0$, then $e=pm$.  Every valuation in (A31l.7)--(A31l.8)
is $m$ times the corresponding valuation at $e=p$.
The complete leading comparison is therefore precisely the one
proved in A31i for $p=3$ and A31j for $p\geq5$, with the same
pole bounds (A31l.3).  All the closed function terms are retained.
The sole leading one-form is the Cartier term of $o_m$.
Primitive jet testing at $M<4p$ and the finite-pole injection force
$o_m=0$ just as there.  The actual reduced cycle normalization
was checked above, so this invocation is a comparison for the
given support, not a substitution of a new transport.
This proves the one-step assertion in all cases.

\paragraph{Iteration to the algebraization threshold.}
Set
\[
          m_0=\lceil e/p\rceil,\qquad
          n_0=\left\lfloor e/(p-1)\right\rfloor+1.
\tag{A31l.10}
\]
If $m_0\geq n_0$, restrict the supplied support to order $n_0$
and apply A31 directly.
Otherwise for every integer $m$ with $m_0\leq m<n_0$ one has
$e\leq pm$ and $(p-1)m\leq e$.
The one-step result extends the specified support successively
from $S_{m_0}$ to $S_{n_0}$.
A31 then extends it through all orders and algebraizes it to a
smooth projective lift of $S$.
This contradicts geometric nonliftability when that hypothesis
is imposed.

The actual A14b and A18 assertion follows as in A31k.
A generic algebraic lift of the primitive etale class gives
filtration membership of the fixed vector by the common realization
kernel and the algebraic generic projector.
The argument does not supply $S_{m_0}$ on an arbitrary model;
it proves that no candidate class-lifting model can have one.
\end{proof}

\paragraph{Programme statement and sources.}
All comparison normalization and graph/frame coherence are the
checked C44-A31-SA1 and A31i; the exact local source identity for
the specified higher-order support is proved above.
The new all-weight computation is (A31l.7), and the primitive
testing and finite-pole inputs are A31h--A31j.
No additional external theorem is invoked.
For JAC1964 QH9-axis and the H9-T7 primitive margin of cote~44,
Theorem~5, scan~004 and scans~008--009, the remaining initial
support requirement is now $\lceil e/p\rceil$, rather than
$\lfloor e/(p-1)\rfloor+1$.
It is an explicit conditional gate, not an all-ramification
nonliftability theorem or a universal sign theorem.
There is no scan~049 product-form assertion and no use of the
unrefereed 032/Milne2026 preprints.

% END work/axis/axis_supplied_support_threshold.tex

% BEGIN work/axis/axis_first_frobenius_range.tex
% C44_POS /root/axis_independent. Author, root and two independent whole-proof checks passed; see dated registers.
\section{The full first Frobenius range of ramification}
\label{sec:c44-a31m}

\paragraph{Scope.}
This is the iteration of the actual closed-layer operation in A31k,
combined with the supplied-support theorem A31l.
It is part of the same developing lifting argument and makes no
separate staged progress-clock claim.

\begin{theorem}[C44-A31m: exclusion through $e=2p-1$ for $p\geq5$]
Let $p\geq5$, let $S$ be a smooth geometrically nonliftable projective
surface in an abelian sixfold $X$ over $k=\overline{\mathbf F}_p$,
and fix an ample bundle $H$.  For all sufficiently large $t$, there
is a nonempty open of smooth genuine axes $B\subset X$, containing
$S$ and cut out by two members of $|H^t|$, such that the fixed
primitive class
\[
                  w=[S]_B-\kappa^*(L_h^2)^{-1}[S]_X
                  \qquad(h=c_1(H^t))
\tag{A31m.1}
\]
is not filtered on any abstract smooth projective lift over a complete
mixed-characteristic DVR with residue field $k$ and
\[
                              e\leq2p-1.
\tag{A31m.2}
\]
The actual A14b centres, and the A18 closed axes with
$A_B=\mathbf Qw$ and $w^2>0$, give corresponding nonliftable
primitive etale classes for every $p\equiv3\pmod4$, $p\geq7$.
Examples of the excluded ranges are $e\leq13$ at $p=7$ and
$e\leq21$ at $p=11$.  No assertion at $e\geq2p$ follows from
this theorem.
\end{theorem}

\begin{proof}
Choose the fixed nonempty open used in A31k at $M=2p$ and in A31l.
The case $e\leq p$ is A31j.  Write
\[
                    e=p+r,\qquad 1\leq r\leq p-1,
                    \qquad p=u\pi^{p+r}.
\tag{A31m.3}
\]
Use the actual canonical lower descent of the fixed cycle $[S]$
from A31i, on common local graph charts, without a global $S_2$.
Its one-contraction and two-contraction valuation bounds are
\[
\begin{split}
 s_e(n)&=n-e v_p(n),\\
 v_e(n)&=n-e\bigl(v_p(n)+\lfloor\log_p(n-1)\rfloor\bigr),
                                              \qquad n\geq2.
\end{split}
\tag{A31m.4}
\]
The latter applies to every actual composition and frame-comparison
function cochain, by the exact $nj$, $j<n$, denominator argument
in A31i.

\paragraph{All weights through the required order.}
We claim
\[
 \begin{split}
 &s_e(p)=-r,\qquad s_e(n)>-r\quad(n\ne p),\\
 &v_e(p)=-r,\qquad v_e(2p)=-2r,\qquad
                v_e(n)>-r\quad(n\ne p,2p).
 \end{split}
\tag{A31m.5}
\]
For the first line, a nonmultiple of $p$ has positive value.
If $v_p(n)=1$ and $n\ne p$, its value is at least $p-r>0$.
If $v_p(n)=a\geq2$, its value is at least
$p^a-a(2p-1)>0$.
This last inequality holds at $(p,a)=(5,2)$, increases with $p$,
and increases with $a$, since the successive difference is
$(p-1)p^a-(2p-1)>0$.

For the second line, weights below $p$ have value $n$.
For $p<n<p^2$ a nonmultiple of $p$ has value
$n-(p+r)\geq1-r>-r$.
At $n=ap$, $2\leq a<p$, the value is
$(a-2)p-2r$; this equals $-2r$ for $a=2$ and exceeds $-r$
for $a\geq3$, because $p-r>0$.

It remains to check all $n\geq p^2$.
For $a\geq2$ put
\[
                         F_a(p)=p^a-(4a-3)p+2a-2.
\tag{A31m.6}
\]
One has $F_2(p)=p^2-5p+2\geq2$ for $p\geq5$, and
\[
                 F_{a+1}(p)-F_a(p)=(p-1)p^a-4p+2>0.
\tag{A31m.7}
\]
At $n=p^a$,
\[
          v_e(n)+r
             =p^a-(2a-1)p-(2a-2)r
             \geq F_a(p)>0.
\tag{A31m.8}
\]
For $p^a<n<p^{a+1}$, if $v_p(n)\leq a-1$, then
$v_e(n)+r\geq F_a(p)+1>0$.
If $v_p(n)=a$, then $n\geq2p^a$ and
\[
          v_e(n)+r
             \geq2p^a-(4a-1)p+2a-1
             =F_a(p)+(p^a-2p+1)>0.
\tag{A31m.9}
\]
This proves every case of (A31m.5), and the same estimates tend
to infinity with $n$, preserving convergence of the comparison.

Multiply the whole rational descent by $\pi^{2r}$.
All its lower cochains are integral by (A31m.5).
We retain only the precision
\[
                            A_{r+1}=R/\pi^{r+1},
                            \qquad r+1\leq p,
\tag{A31m.10}
\]
not precision $2r+1$.  Through order $r$, the function terms have
weights only $p,2p$ and therefore total normal pole at most $2p$.
Every other weight has scaled valuation at least $r+1$.
The one-form terms begin exactly at order $r$, with
\[
                    \eta'=\pi^{2r}\eta
                    \equiv\pi^r\bar u^{-1}A_p(a,b)
                                                    \pmod{\pi^{r+1}}.
\tag{A31m.11}
\]
Here $a,b$ are the first coefficients of the actual local graph
displacements, representing the obstruction $o$ to $S_2$.
The variable-coefficient identity (A31j.5) identifies the class
of $A_p$ with the inverse-Cartier image of this obstruction.
Every degree-zero term at the retained orders is kept; none is
discarded on the ground that it is closed.

\paragraph{Exactly $r$ removals of closed layers.}
Suppose (A31m.1) were filtered.  Its primitive periods vanish in
the full test lattice of A29.  Start with the scaled lower
descent at precision $r+1$.
For $j=0,\ldots,r-1$, after $j$ removals the current precision is
\[
                           n_j=r+1-j\geq2,
\tag{A31m.12}
\]
and the one-form cochain starts at order $r-j>0$.
The special total class is therefore a closed function layer in
$H^2(\mathcal Z_{2p})$.
Primitive jet testing makes its image in
$\mathbb H^2(C_{2p})$ zero, and the strong finite-pole injection
of A31i makes it a genuine Cech coboundary in
$\mathcal Z_{2p}$.

Since $n_j\leq r+1\leq p$ and $n_j\leq e$, the canonical
relative-Frobenius copy
$\mathcal Z_{2p}\otimes_k A_{n_j}$ of A31k lifts the Cech
one-cochain as an actual closed cochain at this precision.
Subtract its genuine total coboundary and divide by $\pi$.
The local flat pole quotients proved in A31k preserve the bound
$C_{2p}$ at precision $n_j-1$, in every common graph chart.
The one-form cochain is unchanged before division modulo
$\pi^{n_j}$; after division its first order decreases by one.
Actual rational periods remain zero.
This proves the induction step without asserting flatness of
supported hypercohomology or existence of a smooth global
support at any intermediate stage.

After precisely $r$ steps the reduced total class belongs to
$\mathbb H^2(C_{2p})$ and has edge
\[
                      \bar u^{-1}c(o^{(1)})
                                             \in H^1(Q_{2p}).
\tag{A31m.13}
\]
It has zero primitive periods.  A31h applies because $2p<4p$,
so it is zero.  The Cartier injection gives $o=0$.
Thus a smooth embedded $S_2$ exists in the original model.
There has been no passage to a new ramified base.

For $p<e\leq2p-1$, one has $\lceil e/p\rceil=2$.
The supplied-support theorem A31l therefore extends this actual
$S_2$ through all orders and algebraizes it.
That contradicts geometric nonliftability of $S$.
The relative Albanese embedding and primitive tests in A29
make the argument valid for every abstract model in the theorem.

The A14b/A18 class-level assertion follows by the same common
realization-kernel and algebraic generic-projector argument as
in A31k.  An algebraic generic lift after a finite field
extension would put the fixed vector in the base-changed Hodge
filtration, hence in the filtration over the original fraction
field, contradicting what was just proved.
\end{proof}

\paragraph{Exact dependency and programme statement.}
The proof uses the actual canonical comparison and all graph/frame
coherence of C44-A31-SA1 and A31i, the proved closed-cochain operation
of A31k, and the supplied-support iteration of A31l.
Its new estimates (A31m.5)--(A31m.9) are proved for all weights,
and its precision requirement is exactly (A31m.10).
It invokes no additional external comparison theorem and makes
no full integral comparison-lattice assertion.

For JAC1964 QH9-axis and the H9-T7 primitive margin of cote~44,
Theorem~5, scan~004 and scans~008--009, this is a further restriction
of route~(a) for actual nonzero primitive classes.  It remains
distinct from the universal sign question and gives no
scan~049 product-form assertion.  It uses neither of the
unrefereed 032/Milne2026 preprints.

% END work/axis/axis_first_frobenius_range.tex

% BEGIN work/scan_w/arbitrary_pole_primitive_testing.tex
% C44_POS /root/scan_w. Author, root and independent complete proof reads passed; see dated registers.
\section{Primitive period tests at arbitrary finite pole order}
\label{sec:c44-a31n}

\paragraph{Scope.}
This is a finite-pole statement on the special fibre.  It removes the
restriction $M<4p$ from detection of a nonzero one-form edge of a supported
class.  It does not assert surjectivity onto all finite jets when $M\geq4p$,
and it does not supply a pole bound or an integral canonical comparison in
an additional ramification range.  Those inputs must still be proved for
each application.

\begin{proposition}[C44-A31n: arbitrary-pole primitive testing modulo closed functions]
Let $k=\overline{\mathbf F}_p$, $p>2$, let $S$ be a fixed smooth projective
surface in an abelian sixfold $X$, and fix an ample bundle $H$ and an
integer $M\geq3$.  For all sufficiently large $t$, let
$\kappa:B\hookrightarrow X$ be a smooth genuine axis containing $S$ and
cut out by two sections of $L=H^t$.  Put $I=I_{S/B}$, $J=I_{S/X}$ and
\[
 \begin{split}
 \mathcal J_M&=[\Omega_B^3/I^M\Omega_B^3
                    \xrightarrow{d}\Omega_B^4/I^{M-1}\Omega_B^4],\\
 C_M&=[M_{M-1,B}^0\xrightarrow{d}M_{M,B}^1],\\
 M_{a,B}^b&=\mathcal Ext_B^2(\mathcal O_B/I^a,\Omega_B^b),\\
 Z_M&=\ker(d:C_M^0\to C_M^1),\qquad Q_M=\operatorname{coker}(d:C_M^0\to C_M^1),\\
 P_B&=\ker\bigl(F^3H^4_{\mathrm{dR}}(B)
                         \xrightarrow{\kappa_*}F^5H^8_{\mathrm{dR}}(X)\bigr).
 \end{split}
\tag{A31n.1}
\]
The degrees of $\mathcal J_M$ are $3,4$ and those of $C_M$ are $0,1$.
Under the finite-jet Serre/residue pairing
\[
             \mathbb H^4(\mathcal J_M)^*
                          \simeq\mathbb H^2(C_M),
\tag{A31n.2}
\]
the annihilator of the restriction of $P_B$ is contained in
\[
       \operatorname{im}\bigl(H^2(Z_M)\longrightarrow\mathbb H^2(C_M)\bigr).
\tag{A31n.3}
\]
In particular, if $\gamma\in\mathbb H^2(C_M)$ has nonzero image under the
edge map $\mathbb H^2(C_M)\to H^1(Q_M)$, then some primitive $F^3$ class
has nonzero period against $\gamma$.

More precisely, write $s=M-3$ and $T=\lfloor M/p\rfloor$.  If $T<4$ the
annihilator is zero.  If $T\geq4$, put $q=T-3$.  Its elements are obtained
by the following composite:
\[
 H^0(M_{q,X^{(1)}}^0)\xrightarrow{\rho_q}
 H^2(M_{q+2,B^{(1)}}^0)\xrightarrow{F_B}
 H^2(Z_M)\longrightarrow\mathbb H^2(C_M).
\tag{A31n.4}
\]
Here $M_{a,X}^b=\mathcal Ext_X^4(\mathcal O_X/J^a,\Omega_X^b)$,
$\rho_q$ is the double connecting map of the dual ideal-power Koszul
sequence on the Frobenius twist, and $F_B$ means the relative Frobenius
map on local-cohomology fractions.  Frobenius pushforward makes all
maps linear on the twist; (A31n.4) is understood in that convention.
\end{proposition}

\begin{proof}
\textit{1. The exact ambient defect for arbitrary $M$.}
Use the complexes of A31h,
\[
 K_M=[I^M\Omega_B^3\to I^{M-1}\Omega_B^4],\quad
 A_s=[J^{s+1}\Omega_X^5\to J^s\Omega_X^6],\quad
 P_s=\Omega_X^{\geq5}/A_s.
\tag{A31n.5}
\]
The weighted Koszul/Gysin proof of A31h.3--A31h.10 does not use
$M<4p$.  For any fixed $M$, its uniform Serre vanishings give
\[
 \mathbb H^4(K_M)\twoheadrightarrow\mathbb H^8(A_s),\qquad
 \mathbb H^5(K_M)\xrightarrow{\sim}\mathbb H^9(A_s).
\tag{A31n.6}
\]
The ordinary versions give a surjection in degree $4\to8$ and an
isomorphism in degree $5\to9$.  Since $P_s$ is supported on a surface,
$\mathbb H^9(P_s)=0$, without a bound on $M$.  The compatible restriction
triangles therefore give the exact sequence
\[
          P_B\longrightarrow\mathbb H^4(\mathcal J_M)
                    \longrightarrow\mathbb H^8(P_s)\longrightarrow0.
\tag{A31n.7}
\]
Here is the diagram chase, including the primitive adjustment.  Lift
the connecting image of a class of $\mathbb H^8(P_s)$ through the
degree-five isomorphism in (A31n.6).  Its image in ordinary degree-five
cohomology vanishes, so it lifts to a jet class.  The remaining difference
comes from ordinary degree eight on $X$, and the ordinary degree-four
Gysin surjection supplies a correction.  Conversely, a jet class mapping
to zero has zero connecting image by the same isomorphism, and therefore
comes from an ordinary degree-four class on $B$.  Its Gysin comes from
$\mathbb H^8(A_s)$.  Subtract a preimage in $\mathbb H^4(K_M)$ using
(A31n.6).  This changes no jets and makes the class primitive.  A22a
identifies the ordinary truncated groups with the displayed Hodge
filtrations.  This proves (A31n.7).

As in A31h.12--A31h.14, residue integration by parts and Serre duality give
\[
 \mathbb H^8(P_s)^*
   =H^0\ker\bigl(d:M_{s,X}^0\longrightarrow M_{s+1,X}^1\bigr).
\tag{A31n.8}
\]
For $s=0$ this group is zero.  Dualizing (A31n.7), its image in
$\mathbb H^2(C_M)$ is exactly the annihilator in the statement.  The issue
is the actual two-term transpose, not just its function-valued coherent
double boundary.  We now give a bounded complex that computes it.
For $M<4p$, A31h already proves that (A31n.8) is zero.  We may therefore
assume $M\geq4p$ in the remaining construction; in particular $s\geq1$.

\textit{2. The finite Koszul--de Rham model of that transpose.}
Let $E=L^{-1}\oplus L^{-1}$.  Locally trivialize $L$ and write its two
equations as $f_0,f_1$.  They are two of the four regular normal
parameters for $S\subset X$.  For $a\geq1$ set
\[
 \mathcal R_a^{-j}=M_{a+j,X}^0\otimes\Lambda^jE,
                  \qquad 0\leq j\leq2.
\tag{A31n.9}
\]
The Koszul differential is multiplication by $f_0,f_1$.  This is the exact
dual ideal-power Koszul sequence except for its leftmost kernel:
\[
 0\to\kappa_*M_{a+2,B}^0\to M_{a+2,X}^0L^{-2}
   \to(M_{a+1,X}^0L^{-1})^{\oplus2}\to M_{a,X}^0\to0.
\tag{A31n.10}
\]
Thus $\mathcal R_a\simeq\kappa_*M_{a+2,B}^0[2]$.  Exactness can be checked
in four normal coordinates: it is the dual of the truncated Koszul
sequence for the quotients by $J^a,J^{a+1},J^{a+2}$ and by $I^{a+2}$.
Their quotients are Cohen--Macaulay of codimension four, so no other Ext
row intervenes.  With generators $\epsilon_i$ of degree $-1$, write
$\partial\epsilon_i=f_i$.  The kernel inclusion is the usual residue
inclusion
\[
       w\longmapsto \epsilon_0\epsilon_1\frac{w}{f_0f_1}.
\tag{A31n.11}
\]

To include one-forms, retain the formal generators $d\epsilon_i$ and set
$\partial(d\epsilon_i)=-df_i$.  Their Koszul degree is $-1$ and their
form degree is $1$.  Ordinary $\Omega_X^1$ terms in Koszul degree $-j$
have coefficient module $M_{s+1+j,X}^1$.  A term with one $d\epsilon_i$
and $j$ ordinary $\epsilon$'s has coefficient module
$M_{s+1+j,X}^0$ and Koszul degree $-j-1$; it has the indicated $L^{-1}$
twists for all its generators.  Call this one-form resolution
$\mathcal T_s$.  Locally it is the cone of the conormal map from the
Koszul resolution with coefficients $E$ to the resolution with
coefficients $\Omega_X^1$, both with index $s+1$.

The only cohomology of $\mathcal T_s$ is in Koszul degree $-2$ and equals
$\kappa_*M_{s+3,B}^1$.  Indeed its two possible cohomology groups are the
kernel and cokernel of
\[
 M_{s+3,B}^0\otimes E|B\xrightarrow{df_0,df_1}
             M_{s+3,B}^0\otimes\Omega_X^1|B.
\tag{A31n.12}
\]
The conormal sequence of the smooth embedding is locally split and has
cokernel $\Omega_B^1$; hence this map is injective and has the asserted
cokernel.  Exterior differentiation, including differentiation of the
$\epsilon_i$, defines a map $\mathcal R_s\to\mathcal T_s$ which
anticommutes with $\partial$.  It respects the displayed bounds: an
ordinary derivative increases the permitted denominator order by one;
replacing an $\epsilon_i$ by $d\epsilon_i$ uses exactly the same
coefficient module in the conormal part.  Its map on the degree-$-2$
cohomology is ordinary $d$ on $B$.  Explicitly, in the supported fraction
module, where a cancelled $f_i$ denominator is zero, one has
\[
\begin{split}
 d\left(\epsilon_0\epsilon_1\frac{w}{f_0f_1}\right)
       -\epsilon_0\epsilon_1\frac{dw}{f_0f_1}
   =\partial\left(
       \epsilon_0\epsilon_1d\epsilon_0\frac{w}{f_0^2f_1}
       +\epsilon_0\epsilon_1d\epsilon_1\frac{w}{f_0f_1^2}\right).
\end{split}
\tag{A31n.12a}
\]
Both terms on the right have precisely the allowed extra one-form pole.
The residue of $dw$ is its restriction to $\Omega_B^1$.
Consequently the total complex
\[
       \mathcal E_s=\operatorname{Tot}
                [\mathcal R_s\xrightarrow{d}\mathcal T_s]
                      \simeq\kappa_*C_M[2]
\tag{A31n.13}
\]
is a finite model; form degrees are $0,1$, and total differential is
$D=\partial+d$ with the indicated signs.

This construction chooses no connection on $L$.  Under a local frame
change the equations and generators both transform by a unit $u$;
the rule $d(u\epsilon)=du\,\epsilon+u\,d\epsilon$ is part of the
one-form resolution.  Its extra term respects the same bound.  Thus
the local resolutions and their differentials glue.  Equivalently,
they are the form-degree-at-most-one part of the Koszul differential
resolution, with the explicit ideal-power bounds above.

There is a strict inclusion of complexes
\[
 [M_{s,X}^0\to M_{s+1,X}^1]\longrightarrow\mathcal E_s
\tag{A31n.14}
\]
into Koszul degree zero.  It is the transpose used in (A31n.8).
For completeness this identifies the actual differential as well as
the two coherent boundary maps.  Dualizing the two ideal-power
Koszul sequences gives (A31n.10).  Dualizing the adjunction sequence
$0\to\Omega_B^3\to\Omega_X^5L^2|B\to(L^3|B)^2\to0$
gives the conormal sequence (A31n.12), including its normal relations
$df_i$.  The transpose of exterior differentiation is $d$, by the
residue integration-by-parts identity.  On the Koszul relations its
Leibniz rule is exactly $\partial(d\epsilon_i)=-df_i$.
These are all terms in (A31n.13).  In a local supported Cech resolution
the primal map is the actual fraction
$df_0\wedge df_1\wedge\widetilde\eta/(f_0f_1)$; its derivative is
$df_0\wedge df_1\wedge d\widetilde\eta/(f_0f_1)$.
Pairing with these fractions gives (A31n.14) and the same signs.
Under frame change the $du$ terms are the displayed formal-generator
terms, or, on the primal side, the cancelled-denominator boundaries.
Thus (A31n.14) computes the full finite derived transpose of (A31n.7),
including its frame homotopies.  It is not an arbitrarily chosen lift
of its function-only double boundary.

\textit{3. A bounded inverse-Cartier map through the entire resolution.}
Every closed section in (A31n.8) is uniquely the Frobenius of a section
$\psi$ of local cohomology on $X^{(1)}$.  This is degree-zero Cartier
with support.  One can check it directly in etale tangent and normal
coordinates: a closed fraction has normal exponents divisible by $p$
and coefficients in the Frobenius image.  The assertion is intrinsic
because ordinary Cartier commutes with localization.  More formally,
the supported Cech complex for the four-generator regular ideal has
local cohomology only in degree four, both on the locally free form
sheaves and on their Cartier cohomology sheaves.  The two spectral
sequences therefore give this degree-zero supported Cartier identity.
The same coordinate check shows its bound: a function in $M_{s,X}^0$
has total normal pole at most $s+3=M$, so its Frobenius root has total
pole at most $T=\lfloor M/p\rfloor$.  If $T<4$ it is zero.  Otherwise
$\psi\in H^0(M_{q,X^{(1)}}^0)$, $q=T-3$.

It remains essential to extend Frobenius through the \emph{whole}
finite resolution.  Mapping only its function terms does not suffice,
since the factors $f_i^{p-1}$ have nonzero derivatives.  In the formal
Koszul--de Rham algebra acting on the coefficient modules of
$\mathcal E_s$, modulo form degree two, put
\[
       t_i=f_i^{p-1}\epsilon_i
                         -f_i^{p-2}\epsilon_i\,d\epsilon_i.
\tag{A31n.15}
\]
The total degree of each summand is $-1$.  Direct calculation gives
\[
 \begin{split}
 D(f_i^{p-1}\epsilon_i)
   &=f_i^p+f_i^{p-1}d\epsilon_i
                         -f_i^{p-2}df_i\,\epsilon_i,\\
 D(f_i^{p-2}\epsilon_i d\epsilon_i)
   &=f_i^{p-1}d\epsilon_i
                         -f_i^{p-2}df_i\,\epsilon_i
                                      \pmod{\text{form degree }2}.
 \end{split}
\tag{A31n.16}
\]
Hence $Dt_i=f_i^p$.  The odd elements $t_0,t_1$ anticommute and have
square zero.  Sending coefficients to their $p$th powers and
$\epsilon_i^{(1)}$ to $t_i$ consequently defines a chain map from
the \emph{function-only} Koszul resolution on the twist:
\[
                  \widetilde F:\mathcal R_q^{(1)}
                                  \longrightarrow F_*\mathcal E_s.
\tag{A31n.17}
\]
This is a map of additive complexes linear on the Frobenius twist;
no differentiation is imposed on its source coefficients.

We verify every bound in (A31n.17).  A coefficient of Koszul degree
$-j$ in the source has total pole at most $T+j$.  Its function-valued
image multiplies its $p$th power by $j$ factors $f_i^{p-1}$ and thus
has total pole at most
\[
                   p(T+j)-j(p-1)=pT+j\leq M+j.
\tag{A31n.18}
\]
This is the target module $M_{s+j,X}^0$.  A one-form term uses exactly
one correction in (A31n.15), replacing one exponent $p-1$ by $p-2$.
Its pole bound is $pT+j+1\leq M+j+1$.  It has $j$ ordinary $\epsilon$'s
and one $d\epsilon$, so its target coefficient module is precisely
$M_{s+1+j,X}^0$.  There are no other terms in form degrees at most one.
All multiplications use regular coefficients; there is no division by
$f_i$ on a local-cohomology module in this argument.

The map also glues exactly in line-bundle frames.  If
$(f_i,\epsilon_i)$ is replaced by $(u f_i,u\epsilon_i)$, substitute
$d(u\epsilon_i)=du\,\epsilon_i+u\,d\epsilon_i$ in (A31n.15).
The additional term contains $\epsilon_i\,du\,\epsilon_i=0$, leaving
exactly $u^p t_i$.  This is the transition on the Frobenius-twisted
source.  Coordinate changes for the other normal parameters are
already respected by Frobenius on local cohomology.  The construction
therefore provides a global bounded chain map, including the normal
correction missing from the function-only computation.

On Koszul degree zero it sends $\psi$ to $\psi^p$ with no one-form
term.  It thus factors the actual map (A31n.14) on every closed
ambient input through $\mathcal R_q^{(1)}$.  After the quasi-isomorphisms
(A31n.10) and (A31n.13), this is a derived map
\[
          M_{q+2,B^{(1)}}^0[2]\longrightarrow F_*C_M[2].
\tag{A31n.19}
\]
After cancelling the shifts, the source is a sheaf in degree zero
and the target is a complex in degrees zero and one.  The standard
truncation triangle shows that every such morphism factors uniquely
through $H^0(F_*C_M)=F_*Z_M$: the other cohomology sheaf is in degree
one, and the required negative Ext groups vanish.  Its map to $Z_M$
is ordinary Frobenius.  Indeed the function component of the residue
kernel (A31n.11) maps to
\[
 \epsilon_0\epsilon_1(f_0f_1)^{p-1}
                 \left(\frac{w}{f_0f_1}\right)^p
            =\epsilon_0\epsilon_1\frac{w^p}{f_0f_1}.
\tag{A31n.20}
\]
The one-form corrections in (A31n.15) make this a chain map and cannot
add a separate derived one-form component to (A31n.19), by the same
truncation argument.  This is the point that a coherent double-boundary
calculation alone would not establish.

Taking global hypercohomology in degree zero now proves exactly the
factorization (A31n.4): the inclusion of the degree-zero source term
into $\mathcal R_q^{(1)}$ induces $\rho_q$, and (A31n.19) induces
Frobenius into $H^2(Z_M)$.  By (A31n.8) every element of the annihilator
is obtained this way.  This proves (A31n.3).

Finally the truncation triangle
$Z_M\to C_M\to Q_M[-1]$ identifies the kernel of the edge map to
$H^1(Q_M)$ with the image of $H^2(Z_M)$.  Its next term is $H^3(Z_M)=0$
because the support has dimension two.  In particular a class with
nonzero edge cannot belong to the annihilator, proving the assertion
about primitive periods.
\end{proof}

\begin{corollary}[Finite-pole Cartier detection]
Fix $M\geq p$.  On the sufficiently general high-degree axes of the
finite-pole enlargement lemma in A31i, any class
$\gamma\in\mathbb H^2(C_M)$ whose edge is a nonzero image of
$c(o^{(1)})\in H^1(Q_p)$ has a nonzero primitive $F^3$ period.
The same conclusion holds for any nonzero ordinary-plus-Cartier edge
covered by that lemma.  Thus an actual integral canonical comparison
with such an edge cannot vanish against all primitive periods,
whatever its proved finite pole bound $M$.
\end{corollary}

\begin{proof}
For fixed $M$, take the intersection of the sufficiently high-degree
open conditions used here and in the enlargement lemma.  That lemma
injects the specified $H^1(Q_p)$ summand into $H^1(Q_M)$.
Proposition A31n detects the resulting nonzero edge.  No inference is
made about a canonical comparison whose finite bound, integrality,
or specified edge has not been established.
\end{proof}

\paragraph{Source and proof scope.}
The cohomological vanishings and fraction Gysin are the already checked
A31h.3--A31h.10 and A22a.  The precise external Cartier input is
Deligne--Illusie (1987), Theorem~1.2, printed pp.~249--250, with the
localization and support passage proved above.  The new finite transpose
argument is the explicit Koszul--de Rham construction
(A31n.9)--(A31n.20); it does not cite a general filtered-transfer or
semiregularity assertion.  In particular no implication from an
unbounded supported class to a chosen bounded representative is used.

% END work/scan_w/arbitrary_pole_primitive_testing.tex

% BEGIN work/axis/axis_frobenius_supplied_support.tex
% C44_POS /root/axis_independent. Author, root and independent whole-proof checks passed; see dated registers.
\section{Frobenius support over a supplied finite-order lifting}
\label{sec:c44-a31o}

\paragraph{Scope.}
This addendum scales A31m using an actual supplied support.
The Frobenius twist is no longer asserted to be constant; the given
support supplies its smooth centre.  The conclusion improves a
conditional initial-order requirement and makes no separate staged
progress-clock claim.

\begin{theorem}[C44-A31o: supplied support at order $\lceil e/(2p-1)\rceil$]
Let $p\geq5$, let $S$ be a fixed smooth surface in an abelian sixfold
$X$ over $k=\overline{\mathbf F}_p$, and fix an ample bundle $H$.
For all sufficiently large $t$, there is a nonempty open of smooth
genuine axes $B\subset X$ containing $S$ and cut out by two members
of $|H^t|$ with the following property.

Let $\mathcal B/R$ be an abstract smooth projective lift of $B$,
where $R$ is a complete mixed-characteristic DVR with residue field
$k$, uniformizer $\pi$, and $e=v_\pi(p)$.
Suppose the fixed primitive class
\[
                      w=[S]_B-\kappa^*(L_h^2)^{-1}[S]_X
                      \qquad(h=c_1(H^t))
\tag{A31o.1}
\]
is filtered.  A supplied smooth embedded lift
$S_m\subset\mathcal B_{R/\pi^m}$ extends to a smooth $S_{m+1}$ if
\[
                              pm<e\leq(2p-1)m.
\tag{A31o.2}
\]
Together with A31l and A31, this implies that a supplied
$S_{\lceil e/(2p-1)\rceil}$ extends through all orders and algebraizes.

For a geometrically nonliftable $S$, no model can satisfy both
filtration membership and that supplied-support hypothesis.
This applies to the actual A14b centres and their A18 axes.
The characteristic-three supplied threshold remains
$\lceil e/3\rceil$ from A31l in this addendum.
No initial support is asserted to exist automatically beyond the
unconditional range in A31m.
\end{theorem}

\begin{proof}
Choose the same fixed open as A31m at $M=2p$, also satisfying A31l.
These geometric conditions do not depend on $m,e$.
The new point is the replacement for the constant closed sheaf
used in A31k.

\paragraph{The actual Frobenius base factorization.}
Let $A_j=R/\pi^j$.  Assume a smooth $S_m$ is supplied, and let
\[
                             1\leq n\leq\min(pm,e).
\tag{A31o.3}
\]
The ring $A_n$ has characteristic $p$.  The $p$th-power map
$R\to A_n$ factors through the ring homomorphism
\[
              \phi_{m,n}:A_m\longrightarrow A_n,\qquad
                              a\bmod\pi^m\longmapsto a^p\bmod\pi^n.
\tag{A31o.4}
\]
It is well-defined because changing a lift by $\pi^m b$ changes
its $p$th power by zero in $A_n$: the cross coefficients are
divisible by $p$, and $\pi^{pm}=0$ in $A_n$.
The characteristic-$p$ binomial formula gives additivity and
multiplicativity.  This formulation does not require a map
$A_n\to A_m$ when $n<m$.

It follows directly from the definition of the relative twist that
\[
       \mathcal B_n^{(1)}
                    =\mathcal B_m\times_{A_m,\phi_{m,n}}A_n.
\tag{A31o.5}
\]
The supplied support therefore gives the actual smooth closed subscheme
\[
             S'_n=S_m\times_{A_m,\phi_{m,n}}A_n
                                     \subset\mathcal B_n^{(1)}.
\tag{A31o.6}
\]
Its special fibre is $S^{(1)}$, and its embedding has codimension two.
Regularity can be checked by the relative normal coordinates:
the equations for $S_m$ extend to relative etale coordinates,
and their base changes do the same.  No flatness of
$\phi_{m,n}$ is required for this smooth pair base change.

\paragraph{The finite-pole closed module and cochain lifting.}
Degree-zero Cartier and the two supported spectral sequences of
A31k identify the local support functions on the twist with the
closed relative support functions on $\mathcal B_n$.
Inside them use the source simple-pole module
\[
       \mathcal P_n=
       \mathcal Ext^2_{\mathcal B_n^{(1)}}
              (\mathcal O_{S'_n},\mathcal O_{\mathcal B_n^{(1)}}).
\tag{A31o.7}
\]
Its relative Frobenius image is a closed module
$\mathcal Z_{2p,n}$ whose special fibre is exactly
$\mathcal Z_{2p}=\ker(d:M_{2p-1}^0\to M_{2p}^1)$.
In local normal coordinates it consists of the relative Frobenius
images of $a/(x'y')$, with function pole $2p$.
The source module is an invertible sheaf on the smooth $S'_n$,
so it is flat over $A_n$.  Its image is isomorphic to it as a
sheaf on the twist.

Take an affine adapted Cech cover and its affine intersections.
The coherent module $\mathcal P_n$ surjects to its special fibre
on sections there: its kernel is coherent, so the obstruction
would lie in a vanishing affine $H^1$.
Thus every special Cech one-cochain in $\mathcal Z_{2p}$ lifts
to a cochain in the actual closed module $\mathcal Z_{2p,n}$.
There is no assertion of a canonical section or of flat
hypercohomology.  What is needed and supplied is a globally
defined sheaf in which these Cech cochains and their coboundaries
can be formed.

The local formal extensions of the specified $S_m$ have generators
congruent modulo $\pi^m$ on common graph charts.  Their $p$th
powers agree modulo $\pi^{pm}$, hence modulo $\pi^n$.
Consequently $\mathcal Z_{2p,n}$ has the same pole bound in
every one of these charts.  Changes of generators for the
same smooth support are handled by the natural Ext description
(A31o.7); no frame or formal-neighbourhood splitting is selected.

Lift any chosen cochain to an actual convergent characteristic-zero
cochain by finite sums of $R$-scalar multiples of $p$th powers
of fractions.  This is possible from the defining tensor
description of relative Frobenius; the source equations map
to the $p$th powers of equations of $S_m$.
Its relative differential is divisible by $p$, and hence by
$\pi^n$.  Its restrictions and Cech coboundary reduce to the
ones in the single global sheaf $\mathcal Z_{2p,n}$.
The normal negative-monomial description and flat pole quotients
of A31k prove the same local saturation after division by $\pi$.
Thus the complete actual closed-layer operation of A31k works
at every precision (A31o.3), with this possibly nonconstant
Frobenius support.  Higher errors vanish at the current precision
and remain outside the next precision after division.

\paragraph{All-weight estimates for the scaled displacements.}
Write
\[
                 e=pm+r,\qquad 1\leq r\leq(p-1)m,
                 \qquad p=u\pi^e.
\tag{A31o.8}
\]
Extend the actual $S_m$ locally to smooth formal supports.
By the checked source-unit argument in A31l, the comparison is
the fixed reduced cycle $[S]$, with no factor $m$.
Its displacements are $\pi^m(a,b)$, and their first coefficients
represent the obstruction $o_m$ to extending this particular $S_m$.
Set $A(n)=v_p(n)+\lfloor\log_p(n-1)\rfloor$ for $n\geq2$.
The one- and two-contraction bounds are
\[
             mn-e v_p(n),\qquad mn-eA(n).
\tag{A31o.9}
\]
The full proof of the inequalities in A31m applies with the real
parameter $\rho=r/m\in(0,p-1]$ in place of its integer parameter.
For clarity, the resulting statements are
\[
 \begin{split}
 &mp-e=-r,\qquad mn-e v_p(n)>-r\quad(n\ne p),\\
 &mp-e=-r,\qquad 2mp-2e=-2r,\qquad
                        mn-eA(n)>-r\quad(n\ne p,2p).
 \end{split}
\tag{A31o.10}
\]
Here are the relevant all-weight bounds explicitly.
Between $p$ and $p^2$, a nonmultiple of $p$ has the second
quantity plus $r$ equal to $m(n-p)\geq m$.
For $n=ap$, $a\geq3$, it is at least $mp-r\geq m$.
Below $p$ it is positive, and $p,2p$ give the displayed equalities.
For $n=p^a$, $a\geq2$, the quantity plus $r$ is at least
\[
                     m F_a(p),\qquad
                     F_a(p)=p^a-(4a-3)p+2a-2>0.
\tag{A31o.11}
\]
For $p^a<n<p^{a+1}$, if $v_p(n)\leq a-1$, its bound is
$m(F_a(p)+1)$; if $v_p(n)=a$, its bound is
$m(F_a(p)+p^a-2p+1)$.
The proof $F_2(p)\geq2$ and
$F_{a+1}-F_a=(p-1)p^a-4p+2>0$ covers every $a$.
For the first quantity in (A31o.9), the cases $v_p(n)=0,1$
are immediate and $v_p(n)=a\geq2$ gives the positive lower
bound $m[p^a-a(2p-1)]$.
This verifies (A31o.10) for all weights and proves convergence.

Scale the actual lower descent by $\pi^{2r}$.
It is integral, and through order $r$ its function terms have
weights only $p,2p$.  Its one-form term starts at order $r$,
with coefficient $\bar u^{-1}A_p(a,b)$.
Use precision $r+1$, for which
\[
                       r+1\leq(p-1)m+1\leq pm<e.
\tag{A31o.12}
\]
Perform exactly $r$ closed-layer removals as in A31m.
At each current precision $n=r+1-j$, the actual module
$\mathcal Z_{2p,n}$ just constructed supplies the necessary
closed Cech lift.  The primitive jet theorem and the injection
$H^2(\mathcal Z_{2p})\to\mathbb H^2(C_{2p})$ make each
special closed layer a Cech coboundary.
Subtracting its actual total boundary and dividing preserves
the pole bound and all zero rational primitive periods.

The final special total class has edge
$\bar u^{-1}c(o_m^{(1)})\in H^1(Q_{2p})$.
Primitive testing at $2p<4p$ and the Cartier injection force
$o_m=0$.  This proves the one-step assertion.

\paragraph{The initial order for iteration.}
Put
\[
          m_0=\left\lceil\frac e{2p-1}\right\rceil,\qquad
          n_0=\left\lfloor\frac e{p-1}\right\rfloor+1.
\tag{A31o.13}
\]
One has $m_0\leq n_0$.  For $m_0\leq m<n_0$,
$(p-1)m\leq e\leq(2p-1)m$.
If $e\leq pm$, use A31l; if $e>pm$, use the just proved
one-step result.  Each extends the specified support by one
order.  At $n_0$, A31 supplies all further extensions and
algebraization.  The class-level consequence for the A14b/A18
examples follows by the common realization-kernel and generic
primitive-projector argument, as in A31k.
The contradiction for geometrically nonliftable $S$ requires
the supplied $S_{m_0}$ and does not manufacture it on an
arbitrary model.
\end{proof}

\paragraph{Sources and programme statement.}
The general-base Cartier and relative Frobenius inputs are the
exact primary statements recorded and proved locally in A31k.
The factorization (A31o.4), the nonconstant smooth support
(A31o.6), affine coherent cochain lifting and all scaled
inequalities are proved here.  Canonical rational normalization
remains C44-A31-SA1 and A31i, with the specified-support
source identity from A31l.
No new comparison theorem or flatness of supported
hypercohomology is invoked.

For JAC1964 QH9-axis and the H9-T7 primitive margin of cote~44,
Theorem~5, scan~004 and scans~008--009, this reduces the
remaining initial support order to $\lceil e/(2p-1)\rceil$
for $p\geq5$.  It is a conditional lifting gate, not universal
sign positivity or all-ramification primitive nonliftability.
It gives no scan~049 product-form assertion and uses neither
of the unrefereed 032/Milne2026 preprints.

% END work/axis/axis_frobenius_supplied_support.tex

% BEGIN work/arithmetic/axis_uniform_ample_conormal.tex
% C44_POS /root/fermat_aj. Author, root and independent entire-proof reads PASS; see dated registers.
\section{Ample conormal bundles and one open for every finite pole bound}
\label{sec:c44-a31p}

\paragraph{Scope.}
The finite-pole quotient lemma inside A31i permitted its open set to
depend on the pole bound.  The result below gives one nonempty open
which works for every finite bound.  It concerns the algebraic
supported complexes.  Primitive period detection and the pole bound
of an actual ramified comparison remain separate inputs.

\begin{lemma}[Positive polynomial representations in rank two]
Let $Y$ be an integral projective variety of positive dimension over
an algebraically closed field of characteristic $p>0$, and let $P$ be
an ample vector bundle of rank two.
For every nonzero finite-dimensional homogeneous polynomial
$\operatorname{GL}_2$-representation $R$ of positive degree,
the associated bundle $R(P)$ is ample.
There is no semisimplicity assumption on $R$.
\end{lemma}
\begin{proof}
\leavevmode\par
A simple constituent of a homogeneous polynomial representation of
degree $d$ has highest weight $(a,b)$ with
\[
                  a\geq b\geq0,\qquad a+b=d.
\tag{A31p.1}
\]
Restriction to the diagonal matrix monoid has weights with
nonnegative exponents and total degree $d$; a highest weight is one
of these weights.  Highest-weight classification gives the indicated
simple constituents.  Write
\[
       a-b=\sum_{i\geq0}m_i p^i,\qquad 0\leq m_i<p.
\]
The corresponding simple module is
\[
       (\det)^b\otimes
           \bigotimes_{i\geq0}(\operatorname{Sym}^{m_i} k^2)^{[i]},
\tag{A31p.2}
\]
where $[i]$ means that the matrix entries in the action are raised
to their $p^i$-th powers.

The rank-two formula has a direct proof.
For $m=\sum m_i p^i$, multiplication embeds the tensor product in
(A31p.2), without its determinant factor, into
$\operatorname{Sym}^m k^2$ as the span of
\[
 x^{m-j}y^j,\qquad
       j=\sum j_i p^i,\quad 0\leq j_i\leq m_i.
\tag{A31p.3}
\]
The digit monomials are distinct, so the map is injective.
It is equivariant by the characteristic-$p$ binomial formula.
The only vectors fixed by all substitutions $y\mapsto y+tx$ in
$\operatorname{Sym}^m k^2$ are the multiples of $x^m$:
on $x\ne0$ a fixed vector gives a polynomial in $y/x$ invariant
under every translation, hence a constant.
Every nonzero subrepresentation of (A31p.3) has a nonzero fixed
vector for this unipotent group, and therefore contains $x^m$.
The opposite-unipotent orbit of $x^m$ spans (A31p.3).
Indeed
$(1+T)^m=\prod_i(1+T^{p^i})^{m_i}$ shows that its nonzero
coefficients are exactly the digit monomials in (A31p.3).
Evaluation at sufficiently many distinct $t\in k$ extracts every
coefficient.  Thus (A31p.3) is simple, of highest weight $(m,0)$.
Highest-weight uniqueness proves (A31p.2).

For bundles, the factor with digit $m_i$ is
$\operatorname{Sym}^{m_i}(F_Y^{i*}P)$, with the corresponding
field twists understood.  Finite Frobenius pullback preserves
ampleness.  Positive symmetric powers and the determinant of an
ample bundle are ample: they are quotients of positive tensor
powers, and tensor products of ample bundles on projective
varieties are ample in arbitrary characteristic.
The latter is Barton's published positive-characteristic theorem.
At least one of $b,m_0,m_1,\ldots$ is positive because $d>0$.
Consequently each simple constituent in (A31p.2) gives an ample
bundle.  A composition series of $R$ gives an exact filtration
by these bundles.  Extensions of ample bundles on a projective
variety are ample, so $R(P)$ is ample.

Divided powers will be used only as polynomial representations.
No identification of $\Gamma^dP$ with $\operatorname{Sym}^dP$ is
made, and no claim that every polynomial representation is a
quotient of a tensor power is needed.
\end{proof}

\begin{lemma}[Uniform twisted vanishing]
Fix a very ample line bundle $A$ on a smooth integral projective
variety $Y$ of positive dimension and finitely many vector bundles
$G_1,\ldots,G_v$.
Choose $c\geq1$ so that each $G_j$ embeds in a finite direct sum of
copies of $A^c$.
If $N$ has rank two and
\[
                           N^*\otimes A^{-c}
                      \quad\hbox{is ample},
\tag{A31p.4}
\]
then, simultaneously for every homogeneous polynomial
$\operatorname{GL}_2$-representation $R$ of every degree $d\geq1$,
\[
                       H^0(Y,R(N)\otimes G_j)=0
                         \quad(1\leq j\leq v).
\tag{A31p.5}
\]
\end{lemma}
\begin{proof}
Such a $c$ exists because $G_j^*\otimes A^c$ is globally generated
for large $c$; dualizing its evaluation gives a locally split
injection.
Put $P=N^*\otimes A^{-c}$, so $N=P^*\otimes A^{-c}$.
Homogeneity gives
\[
 R(N)\otimes G_j\hookrightarrow
       \bigl(R(P^*)\otimes A^{-c(d-1)}\bigr)^{\oplus r_j}.
\tag{A31p.6}
\]
The dual of $R(P^*)$ is a positive-degree polynomial construction
on $P$.  If the polynomial action of $R$ is $\rho$, its dual
polynomial action is
\[
                       \rho^\sharp(g)=\rho(g^{\mathsf t})^{\mathsf t}.
\tag{A31p.7}
\]
This is a representation, is polynomial of the same degree, and
on invertible matrices equals $\rho(g^{-\mathsf t})^{-\mathsf t}$.
Thus $R(P^*)^*=R^\sharp(P)$, without asserting that the ordinary
contragredient representation is polynomial.
The preceding lemma makes this dual bundle ample.
The dual of an ample bundle has no global sections on a
positive-dimensional integral projective variety.
To check this, restrict a hypothetical nonzero map from the ample
bundle to $\mathcal O$ to the normalization of a curve on which
it is nonzero; its image is a quotient line bundle of nonpositive
degree, contrary to ampleness.
Multiplication by a nonzero section of $A^{c(d-1)}$ injects
$R(P^*)\otimes A^{-c(d-1)}$ into $R(P^*)$.
Therefore (A31p.6) proves (A31p.5).
The one integer $c$ is independent of both $d$ and $R$.
The conclusion cannot be extended to arbitrary fixed sheaves
supported at finitely many points; all auxiliary sheaves used
below are vector bundles.
\end{proof}

\begin{lemma}[Generic ample cokernels on a surface in characteristic $p$]
Let $Y$ be an integral projective surface over an algebraically
closed field, let $A$ be very ample, and let $F_0$ have rank four.
For all sufficiently large integers $d$, a nonempty open of maps
\[
                         (A^{-d})^{\oplus2}\longrightarrow F_0
\tag{A31p.8}
\]
has a locally free ample cokernel of rank two.
This holds in arbitrary characteristic.
\end{lemma}
\begin{proof}
\leavevmode\par
We give the algebraic proof, including its amplification step.
First consider maps on $\mathbf P^2$ from
$\mathcal O(-d_0)^{\oplus2}$ to $\mathcal O^{\oplus v}$,
where $v\geq6$ and $d_0\geq v$.
A general map has rank two everywhere: matrix rank drop has
codimension $v-1>2$.
For its cokernel $E_0$, the projective bundle $\mathbf P(E_0)$ is
the intersection of two general divisors of bidegree $(d_0,1)$ in
$\mathbf P^2\times\mathbf P^{v-1}$.
Its projection to $\mathbf P^{v-1}$ is finite.
Cut the two-dimensional fibres successively.
For each positive-dimensional fibre component, the condition
that a degree-$d_0$ divisor contain it imposes at least $d_0+1$
independent linear conditions: $d_0+1$ distinct points can be
separated by products of linear forms.
The parameter space of fibres has dimension $v-1$.
At either step the bad incidence therefore has positive
codimension, since $d_0+1>v-1$.
This is the finiteness argument of Lazarsfeld Lemma~6.3.43;
upper semicontinuity makes the good condition open.
It works for nonreduced fibres by using their
positive-dimensional irreducible components.
The tautological line bundle of $\mathbf P(E_0)$ is the pullback
of $\mathcal O_{\mathbf P^{v-1}}(1)$ by the resulting finite map.
Hence $E_0$ is ample.

Let $\nu_k:\mathbf P^2\to\mathbf P^2$ be the finite map defined
by the $k$-th powers of the homogeneous coordinates.
It satisfies $\nu_k^*\mathcal O(1)=\mathcal O(k)$ even when
inseparable.
For $\xi=c_1(\mathcal O_{\mathbf P(E_0)}(1))$ and base hyperplane
class $H$, choose rational $\epsilon>0$ such that
$\xi-\epsilon H$ is ample.
Finite pullback shows that
$\xi_{\nu_k^*E_0}-k\epsilon H$ is ample.
For any fixed integer $m$ and sufficiently large $k$,
\[
        \xi_{\nu_k^*E_0}-mH
       =(\xi_{\nu_k^*E_0}-k\epsilon H)+(k\epsilon-m)H
\tag{A31p.9}
\]
is an ample class plus a nef class, hence is ample.
Thus $\nu_k^*E_0\otimes\mathcal O(-m)$ is ample.
This proves the required amplification using only ample line
bundles and finite pullback.

Choose $m_0$ so that $F_0\otimes A^m$ is globally generated for
every $m\geq m_0$.
Six general sections generate this rank-four bundle:
failure has codimension three, greater than $\dim Y$.
Fix $v=6$ and, for each such $m$, a surjection
\[
                         (A^{-m})^{\oplus v}\twoheadrightarrow F_0.
\tag{A31p.10}
\]
Three general sections of $A$ define a finite surjective
$f:Y\to\mathbf P^2$ with $f^*\mathcal O(1)=A$.
They have no common zero, and no curve can be contracted
because $A$ is ample; properness then gives finiteness.

Write every sufficiently large integer as
\[
                 d=kd_0+m_0+i,\qquad 0\leq i<d_0.
\]
By (A31p.9), simultaneously for these finitely many $i$,
the bundle $\nu_k^*E_0\otimes\mathcal O(-m_0-i)$ is ample
for large $k$, and has presentation
\[
 0\longrightarrow\mathcal O(-d)^{\oplus2}
 \longrightarrow\mathcal O(-m_0-i)^{\oplus v}
 \longrightarrow
      \nu_k^*E_0\otimes\mathcal O(-m_0-i)\longrightarrow0.
\tag{A31p.11}
\]
Pull back to $Y$.  The ample-cokernel condition is a nonempty
open in the space of maps between the two left bundles.
Compose a map in this space with (A31p.10), for $m=m_0+i$.
The composite has rank two everywhere on another nonempty open:
at each point its values range over all $4$-by-$2$ matrices,
whose rank-drop locus has codimension three.
The two opens meet in the irreducible affine parameter space.
The composite cokernel is locally free and is a quotient of
the ample cokernel of the first map, hence is ample.
We have produced one map (A31p.8) with ample locally free cokernel.
Openness of ampleness in proper families proves the assertion
for the full space of maps (A31p.8).
No separability or characteristic-zero tensor decomposition is used.
\end{proof}

\begin{theorem}[C44-A31p: one open for every finite pole bound]
Let $k=\overline{\mathbf F}_p$, $p>2$, let $S$ be a smooth integral
projective surface in an abelian sixfold $X$, and let $H$ be ample
on $X$.  Fix a very ample bundle $A$ on $S$.
For every fixed $c\geq1$ and every sufficiently large $t$, there
is a nonempty open of smooth genuine axes $B\subset X$, containing
$S$ and cut out by two members of $|H^t|$, on which
\[
                         N_{S/B}^*\otimes A^{-c}
                         \quad\hbox{is ample}.
\tag{A31p.12}
\]
After choosing $c$ once depending on $S$ and $p$, there are
$t_0$ and, for each $t\geq t_0$, one nonempty open $\mathcal U_t$
of such axes with the following properties simultaneously for
every finite integer $M\geq p$.

Use the notation of the A31i finite-pole quotient lemma:
\[
 \begin{split}
 I&=I_{S/B},\qquad N=N_{S/B},\\
 M_m^a&=\mathcal Ext_B^2(\mathcal O_B/I^m,\Omega_B^a),\\
 C_m&=[M_{m-1}^0\xrightarrow d M_m^1],\qquad
 Q_m=\operatorname{coker}(d:M_{m-1}^0\to M_m^1).
 \end{split}
\tag{A31p.13}
\]
Then $Q_p\to Q_M$ is injective and
\[
                          H^0(Q_M/Q_p)=0.
\tag{A31p.14}
\]
The canonical ordinary and Cartier maps give an injection
\[
                  H^1(F_{S*}N)\oplus H^1(N^{(1)})
                     \longrightarrow H^1(Q_M).
\tag{A31p.15}
\]
Put
\[
 \begin{split}
 Z_M&=\ker(d:M_{M-1}^0\to M_M^1),\\
 D_M&=\{\theta\in M_{M-1}^0:d\theta\in M_p^1\},\\
 C'_M&=[D_M\xrightarrow d M_p^1].
 \end{split}
\]
There is a canonical decomposition and injection
\[
       C'_M=C_p\oplus Z_M[0],\qquad
       \mathbb H^2(C'_M)\hookrightarrow\mathbb H^2(C_M).
\tag{A31p.16}
\]
All maps involving $d$ are maps of additive sheaves, or of coherent
sheaves after relative Frobenius pushforward.
The latter can live on a finite thickening of $S^{(1)}$; they
are not asserted to be bundles on the reduced surface.
The open $\mathcal U_t$ may be intersected with any previously
specified nonempty open in the same axis parameter space.
No intersection of infinitely many open sets is used.
\end{theorem}
\begin{proof}
\leavevmode\par
\paragraph{Realizing the ample cokernel as an axis conormal.}
Write $E=N_{S/X}$ and $L=H^t$.
The normal matrix of the two defining sections gives
\[
        0\longrightarrow (L^{-1}|S)^{\oplus2}
          \longrightarrow E^*
          \longrightarrow N_{S/B}^*\longrightarrow0.
\tag{A31p.17}
\]
Every map on the left occurs as a normal matrix for large $t$:
\[
 H^0(X,I_{S/X}L)\longrightarrow H^0(S,E^*\otimes L|S)
\tag{A31p.18}
\]
is surjective by $H^1(X,I_{S/X}^2L)=0$.
Hence the axis parameter space surjects onto the affine space
of normal matrices.

There is no divisibility restriction on $t$.
Choose $a>0$ with $T=H^a|S$ very ample.
For each residue class of $t$ modulo $a$, choose
$0\leq s<a$ with $t+s$ divisible by $a$.
Choose $b$ once so that
\[
                      H^s|S\otimes T^b\otimes A^{-c}
\]
is globally generated for every one of these finitely many $s$.
Apply the preceding lemma to
\[
 (T^{-(t+s)/a-b})^{\oplus2}
        \longrightarrow E^*\otimes H^{-s}|S\otimes T^{-b}.
\tag{A31p.19}
\]
This is (A31p.17) tensored with the fixed bundle
$H^{-s}|S\otimes T^{-b}$, so its map parameter space is unchanged.
It gives an ample cokernel
$N_{S/B}^*\otimes H^{-s}|S\otimes T^{-b}$.
Tensoring with the displayed globally generated line bundle gives
(A31p.12).  This tensor product is ample because it is a quotient
of a finite direct sum of copies of the original ample bundle.
There are only finitely many $s$, so one threshold works for all
sufficiently large $t$.

Pull back the nonempty open of normal matrices through
(A31p.18) and intersect it with the smoothness and genuine-pencil
open of A06.  Both are nonempty opens in the same irreducible
parameter space.
Rank drop along $S$ has codimension three; off $S$ the already
established axis jet conditions apply.
This proves the geometric assertion and the claimed compatibility
with other nonempty open conditions.

\paragraph{The fixed auxiliary bundles.}
On $S$ take the finite collection
\[
                              \mathcal O_S,\quad\Omega_S^1.
\]
On $S^{(1)}$ take
\[
          \mathcal O_{S^{(1)}},\quad B_S^1,\quad
               F_{S*}\Omega_S^1/B_S^1,
\tag{A31p.20}
\]
where $B_S^1=\operatorname{im}(d:F_{S*}\mathcal O_S\to
F_{S*}\Omega_S^1)$.
These are vector bundles: in etale coordinates the Cartier
complex is a tensor product of the one-variable complexes with
basis $1,z,\ldots,z^{p-1}$, whose kernels, images and cokernels
split over the Frobenius-twisted coordinate ring.
Choose $c$ so that these bundles embed in sums of $A^c$,
respectively $(A^{(1)})^c$.
The uniform twisted-vanishing lemma and its field-twisted version
then apply on the single open (A31p.12).
They kill every positive homogeneous polynomial bundle in $N$,
or in $N^{(1)}$, tensored with the indicated fixed bundle.
The integer $c$ is independent of a representation or pole bound.

\paragraph{Every finite-pole quotient on this one open.}
The intrinsic filtration and strictness (Q9.7)--(Q9.17) of A31i
are sheaf calculations valid for every fixed finite $M$;
they do not use genericity.
Recall
\[
 \begin{split}
 D_n&=\det N\otimes\Gamma^{n-2}N,\\
 G_n&=\operatorname{coker}\bigl(D_n\to N\otimes\Gamma^{n-1}N\bigr),\\
 K_n&=\ker\bigl(D_n\to N\otimes\Gamma^{n-1}N\bigr).
 \end{split}
\tag{A31p.21}
\]
The bundles $G_n$ and $D_n/K_n$ are homogeneous polynomial
constructions in $N$ of degree $n$.
For the exceptional kernels, if $n=pr$ with $r\geq2$, then
\[
       K_n=F_S^*K_n',\qquad
       K_n'=\det N^{(1)}\otimes\Gamma^{r-2}N^{(1)},
\tag{A31p.22}
\]
and (Q9.14) proves that the quotient remains polynomial.
Otherwise $K_n=0$.
The bundle $K_n'$ is polynomial of positive degree $r$ on
$S^{(1)}$.

After Frobenius pushforward the finite filtration of $Q_M/Q_p$
has only the following types of factors:
\[
 \begin{gathered}
 F_{S*}G_n,\qquad
 F_{S*}\bigl((D_n/K_n)\otimes\Omega_S^1\bigr),\\
 K_n'\otimes(F_{S*}\Omega_S^1/B_S^1),\qquad
 F_{S*}(D_{M+1}\otimes\Omega_S^1).
 \end{gathered}
\tag{A31p.23}
\]
Every one has zero $H^0$ by (A31p.20).
This holds for every relevant $n$, however large.
Therefore (A31p.14) follows for every finite $M$ on the same open.
The sheaf injection $Q_p\to Q_M$ is the already proved bounded
strictness statement, independent of these section vanishings.

All factors in (A31h.21), the quotient of $Q_p$ by its disjoint
ordinary and Cartier subbundles, are also positive polynomial
bundles tensored with members of the same finite collection.
Its $H^0$ is zero.
This proves (A31p.15) for $M=p$ and then for every $M$
by (A31p.14).

The same strictness calculation gives the decomposition in
(A31p.16).  The quotient complex $C_M/C'_M$ is quasi-isomorphic
to $(Q_M/Q_p)[-1]$; its degree-one hypercohomology is zero.
The long exact sequence gives the degree-two injection.
Thus the closed function summands are retained in full.
The one ample-cokernel open implies all the vanishings; no
infinite intersection of generic conditions is being taken.
\end{proof}

\paragraph{Exact source and proof-read scope.}
The generic-cokernel argument is Lazarsfeld,
\emph{Positivity in Algebraic Geometry II}, Theorem~6.3.65,
printed pp.~55--56 of the author's Chapter~6 PDF, with
Proposition~6.3.64, pp.~54--55, Lemma~6.3.43, pp.~43--44,
and the required line-twist case of Proposition~6.3.62, p.~54.
All these displayed proofs have been read in full.
The amplification step was also proved directly in (A31p.9).
Remarks~6.1.17, p.~15, and 6.2.16, p.~26, explicitly certify the
positive-characteristic scope of the formal ample and nef facts:
\[
 \text{\url{https://www.math.stonybrook.edu/~roblaz/Reprints/Laz.PAG.Chapt6.pdf}}.
\]
The tensor-product input is C.~M.~Barton,
\emph{Tensor Products of Ample Vector Bundles in Characteristic $p$},
\emph{Amer.\ J.\ Math.} \textbf{93} (1971), 429--438,
DOI~10.2307/2373385, as explicitly stated in the checked
Lazarsfeld Remark~6.1.17.
The original Barton article was not accessible in this source
check; no claim to have reread its entire proof is made.

The representation input is highest-weight classification and
the rank-two form of Steinberg's tensor-product theorem.
Independently opened locators are S.~Riche's author notes,
Chapter~1, Theorem~1.16 and its complete proof, pp.~22--23,
Theorem~2.12 and Example~2.13, pp.~32--33, and Appendix~B,
Exercise~1.1, p.~283:
\[
                 \text{\url{https://riche.perso.math.cnrs.fr/curso-riche-total.pdf}}.
\]
Theorem~2.12 refers to Jantzen,
\emph{Representations of Algebraic Groups}, second edition,
Proposition~II.3.16.
The digit formula, including its irreducibility argument, was
proved above; no unproved exercise is being treated as a proof.

\paragraph{Programme status.}
A31p removes the dependence of the algebraic noncancellation open
on the finite pole bound.
It does not prove primitive period detection for every larger
complex, identify an arbitrary ramified canonical descent with a
bounded complex, or construct a smooth initial support.
It does not alter Fermat arithmetic Abel--Jacobi gate F07.
It uses neither of the unrefereed 032/Milne2026 preprints.


% END work/arithmetic/axis_uniform_ample_conormal.tex

% BEGIN work/axis/axis_second_frobenius_range.tex
% C44_POS /root/axis_independent. Author, root and independent whole-proof checks passed; see dated registers.
\section{One controlled change of pole order in the Frobenius argument}
\label{sec:c44-a31q}

\paragraph{Scope.}
This continues A31k--A31o.  Above the constant-twist precision,
we use actual local cochains with a bounded extra pole term.
We do not posit a globally flat closed sheaf there.
The argument belongs to the same developing lifting advance and
makes no separate staged progress-clock claim.

\begin{theorem}[C44-A31q: the range $e\leq3p-1$ and its supplied-support form]
Let $p\geq7$, let $S$ be a fixed smooth surface in an abelian sixfold
$X$ over $k=\overline{\mathbf F}_p$, and fix an ample bundle $H$.
For all sufficiently large $t$, there is a nonempty open of smooth
genuine axes $B\subset X$, containing $S$ and cut out by two
members of $|H^t|$, with the following property.
On any abstract smooth projective lift $\mathcal B/R$ of $B$,
where $R$ is a complete mixed-characteristic DVR with residue
field $k$ and uniformizer $\pi$, write $e=v_\pi(p)$ and
\[
                     w=[S]_B-\kappa^*(L_h^2)^{-1}[S]_X
                     \qquad(h=c_1(H^t)).
\tag{A31q.1}
\]
If $w$ is filtered, a supplied smooth embedded $S_m$ extends
to a smooth $S_{m+1}$ whenever
\[
                            pm<e\leq(3p-1)m.
\tag{A31q.2}
\]
Consequently a supplied $S_{\lceil e/(3p-1)\rceil}$ extends through
all orders and algebraizes.

If $S$ is geometrically nonliftable, $w$ is therefore not filtered
on any such model with $e\leq3p-1$, since the required initial
support is then just $S_1=S$.
The A14b centres and A18 closed axes with $A_B=\mathbf Qw$ give
actual nonliftable primitive etale classes throughout this range for
$p\equiv3\pmod4$.
For example the excluded indices are $e\leq20$ at $p=7$
and $e\leq32$ at $p=11$.
At arbitrary $e$ the asserted condition is the supplied
$S_{\lceil e/(3p-1)\rceil}$; it is not asserted to exist.
\end{theorem}

\subsection{A local lift with one explicitly bounded new pole term}

Write
\[
 M_a^b=\mathcal Ext_B^2(\mathcal O_B/I_{S/B}^a,\Omega_B^b),
 \qquad C_M=[M_{M-1}^0\xrightarrow d M_M^1],
 \qquad \mathcal Z_M=\ker(d|M_{M-1}^0).
\tag{A31q.3}
\]
The following cochain statement replaces a constant closed sheaf
at precisions greater than $pm$.

\begin{lemma}[One change of normal pole order]
Suppose a smooth $S_m\subset\mathcal B_{R/\pi^m}$ is supplied.
Let
\[
                         1\leq n\leq\min(2pm,e).
\tag{A31q.4}
\]
Every special Cech one-cochain $b$ in $\mathcal Z_{2p}$
has an actual convergent characteristic-zero cochain lift
$\widehat b$ with differential divisible by $p$.
In every common smooth graph chart for local extensions of $S_m$,
its reduction modulo $\pi^n$ has the bound
\[
        \widehat b\in
             M_{2p-1}^0+\pi^{pm}M_{3p-1}^0
                                              \pmod{\pi^n}.
\tag{A31q.5}
\]
Here the modules on the right are the local relative pole modules
in that chart.  The same bound holds for its restrictions and
Cech coboundary.  No global smooth $S_{m+1}$ or global flat closed
module at this precision is assumed.
\end{lemma}

\begin{proof}
Over the perfect field $k$, absolute Frobenius gives an additive
isomorphism from the simple supported function module $M_1^0$
onto $\mathcal Z_{2p}$, with its Frobenius scalar convention.
Indeed the degree-zero Cartier description says that a closed
fraction of total normal pole at most $2p$ has the form
$a^p/(x^py^p)$; its unique root in the simple module is
$a/(xy)$.  Equivalently one can use relative Cartier and the
isomorphism on coefficients induced by the automorphism $F_k$.
Take the corresponding root Cech cochain $\rho$ of $b$.

The actual given support supplies the coherent invertible module
\[
       \mathcal E_m=
       \mathcal Ext^2_{\mathcal B_m}
                    (\mathcal O_{S_m},\mathcal O_{\mathcal B_m})
\tag{A31q.6}
\]
on $S_m$.  On each affine adapted Cech intersection,
sections of $\mathcal E_m$ surject onto its special fibre
$M_1^0$: the coherent kernel has no affine $H^1$.
Lift $\rho$ as a Cech cochain in this single sheaf.
This choice need not be a cocycle.
On local formal extensions of $S_m$, choose actual meromorphic
representatives $\widehat\rho$ of its fractions, with one
normal denominator in each direction, and set
$\widehat b=\widehat\rho^{\,p}$ locally.
Finite sums can be treated term by term.
They are actual cochains on a common tubular Cech refinement,
and $d\widehat b=p\widehat\rho^{p-1}d\widehat\rho$.
The elementary binomial expansion converges on the strict
neighbourhoods used in A31i.

We verify (A31q.5), including changes of graphs and frames.
In a common tangent chart, the normal graph generators for two
extensions of $S_m$ differ by $\pi^m(a,b)$.
To first order in this displacement,
\[
 \frac1{(x+\pi^m a)(y+\pi^m b)}
   =\frac1{xy}
     -\pi^m\!\left(\frac a{x^2y}+\frac b{xy^2}\right)
     +O(\pi^{2m}).
\tag{A31q.7}
\]
Changes in the coefficient or tangent chart add a first-order
term of pole at most three as well.
The lifted root cochains are sections of the same $\mathcal E_m$,
so their discrepancy, after these graph changes, is zero modulo
$\pi^m$ as a local cohomology section.
Consequently the full first discrepancy is $\pi^m$ times a
fraction of total pole at most three.
Changes of generators preserving a local smooth pair preserve
its simple Ext module exactly.  Thus they do not add a new pole
term; changing their lifts only contributes to the just described
$\pi^m$ discrepancy.
Terms with only one support denominator are zero as local
cohomology sections, before any contraction is taken.

In $R/\pi^n$ one has $p=0$.  Raising this calculation to the
$p$th power retains only its $p$th powers.
The first discrepancy has factor $\pi^{pm}$ and pole at most
$3p$; every second-order discrepancy has factor
$\pi^{2pm}=0$.  This proves (A31q.5).
In explicit common normal coordinates the only new denominator
terms are
\[
       -\pi^{pm}\left(
                \frac{a^p}{x^{2p}y^p}
                +\frac{b^p}{x^py^{2p}}\right)
\tag{A31q.8}
\]
times the lifted coefficient, together with coefficient terms
of pole $2p$.
Mixed binomial terms are divisible by $p$ and vanish at the
precision (A31q.4).  The same observation shows that changing
representatives by a one-denominator section does not introduce
a mixed supported term at this precision.
Restrictions and Cech coboundaries are formed from the actual
chosen cochains, so the estimate persists under them.
This is a cochain estimate in each common chart, not a claim that
the displayed weighted sum glues as a flat closed sheaf.
\end{proof}

\subsection{The all-weight bounds and the two stages of removal}

\begin{proof}[Proof of the theorem]
Choose the finite generic conditions of A31h and A31i simultaneously
at $M=2p$ and $M=3p$, as well as the conditions in A31l.
Both pole bounds are below $4p$.
These conditions are independent of $m,e$.
Let $S_m$ be supplied and put
\[
            e=pm+r,\qquad 1\leq r\leq(2p-1)m,\qquad p=u\pi^e.
\tag{A31q.9}
\]
The source-unit normalization of A31l gives the canonical
comparison of the fixed reduced $[S]$, with displacements
$\pi^m(a,b)$ and actual first extension obstruction $o_m$.
There is no multiplicity $m$ in the cycle class.

For $A(n)=v_p(n)+\lfloor\log_p(n-1)\rfloor$, the function
valuation bound is $mn-eA(n)$ and the one-form bound is
$mn-e v_p(n)$.
For functions, the only weights whose bound can be at most $-r$
are $p,2p,3p$, with respective bounds
\[
                           -r,\quad -2r,\quad pm-2r.
\tag{A31q.10}
\]
Here is the full check.  Below $p$ the bounds are positive.
Between $p$ and $p^2$, a nonmultiple has bound plus $r$
equal to $m(n-p)>0$; for $n=ap$ with $a\geq4$ it is
$(a-2)pm-r\geq2pm-r\geq m$.
For all remaining weights set
\[
                     G_a(p)=p^a-(6a-5)p+2a-2,\qquad a\geq2.
\tag{A31q.11}
\]
One has $G_2(p)=p^2-7p+2\geq2$ for $p\geq7$ and
$G_{a+1}-G_a=(p-1)p^a-6p+2>0$.
At $n=p^a$, the function bound plus $r$ is at least $mG_a(p)$.
For $p^a<n<p^{a+1}$, it is at least $m(G_a(p)+1)$ if
$v_p(n)\leq a-1$, and at least
$m(G_a(p)+p^a-3p+1)$ if $v_p(n)=a$.
These are all strictly positive.

For one-forms, the unique least bound is $-r$ at weight $p$.
For valuation one, every other weight has bound plus $r$
at least $mp$; for valuation zero it is positive.
For valuation $a\geq2$, the bound plus $r$ is at least
\[
                         m[p^a-(3a-2)p+a-1]>0.
\tag{A31q.12}
\]
At $a=2$ the bracket is $p^2-4p+1>0$ and its successive
difference is $(p-1)p^a-3p+1>0$.
This checks all weights and also gives convergence to
infinite valuation.

Scale the actual lower descent by $\pi^{2r}$ and retain precision
$r+1$.  Its lower cochains are integral.  Through order $r$ its
functions have the weighted pole bound
\[
          \kappa\in M_{2p-1}^0+\pi^{pm}M_{3p-1}^0
                                           \pmod{\pi^{r+1}},
\tag{A31q.13}
\]
and its one-form begins with
$\pi^r\bar u^{-1}A_p(a,b)$.
The weight-$3p$ function term starts at order $pm$, independently
of $r$; all other weights beyond $p,2p,3p$ vanish at the retained
precision.  Notice
\[
                       r+1\leq(2p-1)m+1\leq2pm,
                       \qquad r+1\leq e.
\tag{A31q.14}
\]

\paragraph{First stage: simple closed fractions.}
For $0\leq j<\min(r,pm)$, maintain the following invariant after
$j$ removals:
\[
 \begin{split}
  &\hbox{precision }r+1-j,\qquad
      \eta\hbox{ starts at order }r-j,\\
  &\kappa\in M_{2p-1}^0+
                    \pi^{pm-j}M_{3p-1}^0
                                  \pmod{\pi^{r+1-j}}.
 \end{split}
\tag{A31q.15}
\]
It holds initially by (A31q.13).
Because $j<pm$, the special closed layer has pole at most $2p$.
Primitive testing at $2p$ and the strong finite-pole injection
$H^2(\mathcal Z_{2p})\to\mathbb H^2(C_{2p})$ show it is a
Cech coboundary in $\mathcal Z_{2p}$ itself.
Choose its Cech one-cochain and lift it by the preceding lemma,
using the actual specified $S_m$.
The precision satisfies (A31q.4).
Subtract its genuine total coboundary and divide by $\pi$.
The lemma contributes no pole beyond $3p$, and its new
high-pole part has coefficient $\pi^{pm}$ before division.
Together with the existing $\pi^{pm-j}$ term this gives exactly
the invariant (A31q.15) at $j+1$.
The differential of the correcting cochain is divisible by
$p$, hence zero at the current precision, so the asserted
one-form order is unchanged before division.

All these operations use the actual tubular Cech cochains.
Their boundaries preserve rational primitive periods, and local
flat pole quotients justify each division as in A31k.
The lower cochains remain integral.
Their residue pairings reduce to the finite-pole pairings, while
the degree-two Gysin term pairs to zero against $F^3$.
Thus every reduced primitive period used here is indeed zero
under the filtration hypothesis.

\paragraph{Second stage: the full closed module of pole $3p$.}
If $r\leq pm$, the first stage already reaches the target order.
If $r>pm$, at $j=pm$ the remaining precision is
\[
                         r+1-pm\leq(p-1)m+1\leq pm.
\tag{A31q.16}
\]
The actual Frobenius support $S'_n\subset\mathcal B_n^{(1)}$
of A31o therefore exists at every remaining precision $n$.
Replace its simple source Ext line by
\[
       \mathcal Ext^2_{\mathcal B_n^{(1)}}
           (\mathcal O_{\mathcal B_n^{(1)}}/(I_{S'_n})^2,
                                      \mathcal O_{\mathcal B_n^{(1)}}).
\tag{A31q.17}
\]
This is the module of source function fractions of total pole
at most three.  Its relative Frobenius image is the full
closed module of pole at most $3p$.
Its special fibre is $\mathcal Z_{3p}$.
It is flat over $A_n$, by the negative-monomial basis in a
smooth pair chart, and coherent; on affine intersections its
sections surject onto the special sections.
The proof in A31o thus supplies actual closed Cech lifts,
now for $\mathcal Z_{3p}$.
It asserts neither constancy of this module nor flat
hypercohomology.

Perform the remaining $r-pm$ closed-layer removals using
primitive testing and the strong injection at $M=3p$.
Every remaining precision is at most $pm$, so (A31q.17)
fits every local extension of the specified $S_m$ exactly.
The genuine boundaries, local saturation and period argument
are the ones already proved in A31k and A31o.

After all $r$ removals, the special total class lies in
$\mathbb H^2(C_{3p})$ and its edge is
\[
                             \bar u^{-1}c(o_m^{(1)})
                                                   \in H^1(Q_{3p}).
\tag{A31q.18}
\]
The variable-coefficient identity (A31j.5) is used here,
including its $da,db$ terms.
Since $3p<4p$, primitive jet testing makes this class zero.
The Cartier injection gives $o_m=0$, so the particular $S_m$
extends to $S_{m+1}$.

\paragraph{Iteration and the unconditional first support.}
Set $m_0=\lceil e/(3p-1)\rceil$ and
$n_0=\lfloor e/(p-1)\rfloor+1$.
For every $m_0\leq m<n_0$ one has
$(p-1)m\leq e\leq(3p-1)m$.
Use A31l when $e\leq pm$ and the new one-step result when
$e>pm$.  This extends the specified initial support to $S_{n_0}$.
A31 then gives all further orders and algebraization.
If $e\leq3p-1$, $m_0=1$ and the initial support is automatic.
Geometric nonliftability of $S$ gives the contradiction.
A29 makes the proof valid for every abstract model, and the
common realization-kernel and generic-projector argument in A31k
gives the actual A14b/A18 primitive class conclusion.
\end{proof}

\paragraph{The next boundary is not covered.}
At $e=3p$ with first displacements $\pi(a,b)$, the weight-$4p$
function term reaches the same target order as the leading
Cartier one-form.  A second normal change in the lifted
$p$th-power fraction can also have pole $4p$.
The full primitive jet-surjectivity used here requires $M<4p$.
At $p=7$, the weight-$p^2$ denominator bound also reaches the
target at that boundary, so merely adding the pole-$4p$ terms
would not establish a complete estimate.
The argument therefore does not extrapolate past $3p-1$.
For $p=5$, extra weight-$p^2$ terms enter before this range
is complete, and the present proof makes no improvement over A31m.

\paragraph{Programme and exact dependency.}
The new input is the actual moving-graph cochain estimate
(A31q.5)--(A31q.8), with its two-stage precision bookkeeping.
The rational cycle normalization and graph/frame comparison
remain C44-A31-SA1, A31i and A31l; finite-pole injectivity and
primitive testing are A31h and the complete A31i quotient lemma.
No new integral comparison theorem or flat supported
hypercohomology is invoked.
For JAC1964 QH9-axis and the H9-T7 primitive margin of cote~44,
Theorem~5, scan~004 and scans~008--009, this restricts route~(a)
for actual nonzero primitive classes.
It is not universal sign positivity or all-ramification
nonliftability, gives no scan~049 product-form assertion,
and uses neither of the unrefereed 032/Milne2026 preprints.

% END work/axis/axis_second_frobenius_range.tex

% BEGIN work/scan_w/ambient_closed_layer_removal.tex
% C44_POS. Author, root and independent entire-proof readings PASS; scopes recorded in CHECK_EVENTS.
\section{Removing an ambient closed layer at finite precision}
\label{sec:c44-a31r}

\paragraph{Scope.}
The argument keeps an actual ordinary rational cohomology class and a
bounded supported representative only at the current finite precision.
An ambient class need not have an actual rational supported lift.
The pole bound $M$ is fixed, the special axis must satisfy A31n at that
bound, and the precision below is at most $p$.  No uniform-in-$M$ Serre
threshold or arbitrary-precision constant Frobenius twist is asserted.

\begin{proposition}[C44-A31r: ambient closed-layer removal]
Let $p>2$, $k=\overline{\mathbf F}_p$, and let
$S\subset B\xrightarrow{\kappa}X$ be a smooth surface in a smooth
genuine axis in an abelian sixfold, as in A31n.  Fix $M\geq p$.
Assume the primitive-testing conclusion of A31n and the injection
\[
       H^2(Z_M)\hookrightarrow\mathbb H^2(C_M)
\tag{A31r.1}
\]
of A31i, or its uniform version A31p.  These are hypotheses on the
special axis at the fixed bound $M$.

Let $\mathcal B/R$ be an abstract smooth projective lift of $B$ over a
complete mixed-characteristic DVR with uniformizer $\pi$, fraction
field $K$, residue field $k$, and $e=v_\pi(p)$.  Use the relative
Albanese closed immersion
$\mathcal\kappa:\mathcal B\hookrightarrow\mathcal X$ of A29a.
It is not assumed to be a complete intersection over $R$.
Put
\[
 \begin{split}
 A_n&=R/\pi^n,\qquad 2\leq n\leq\min(p,e),\\
 \Lambda_R&=\mathbb H^4(\mathcal B,\Omega_{\mathcal B/R}^{<2})
             =H^4_{\mathrm{dR}}(\mathcal B/R)/F^2,\\
 \mathscr T_R&=\ker\bigl(F^3H^4_{\mathrm{dR}}(\mathcal B/R)
                     \xrightarrow{\mathcal\kappa_*}
                                F^5H^8_{\mathrm{dR}}(\mathcal X/R)\bigr).
 \end{split}
\tag{A31r.2}
\]
The ordinary quotient $\Lambda_R$ is finite free and commutes with base
change.  Here $\Omega^{<2}$ is the quotient of the de Rham complex by
form degrees at least two.

Suppose $a\in H^4_{\mathrm{dR}}(\mathcal B_K/K)$ has zero pairing with
$\mathscr T_R$, and its image $\beta$ in $\Lambda_R\otimes_R K$ lies
in $\Lambda_R$.  Suppose the reduction of $\beta$ modulo $\pi^n$ is
the forgotten class of a cocycle $\xi_n$ in the \emph{actual} supported
relative complex
\[
           R\Gamma_S(\mathcal B_n,
                         \Omega_{\mathcal B_n/A_n}^{<2}).
\tag{A31r.3}
\]
In the local-cohomology Cech model write its lower components as
$(\eta_n,\kappa_n)$: one-forms in Cech degree one and functions in
Cech degree two, with the support shift two.  Assume that in every
common local smooth graph chart these components have function pole
at most $M$ and one-form pole at most $M+1$, and that
$\eta_n\equiv0\pmod\pi$.

Then there is an actual class
$b\in H^4_{\mathrm{dR}}(\mathcal X/R)$ such that
\[
                         a'=\frac{a-\mathcal\kappa^*b}{\pi}
\tag{A31r.4}
\]
again has zero primitive periods and has integral image in
$\Lambda_R$.  This image modulo $\pi^{n-1}$ has a supported cocycle
$\xi'_{n-1}$ with the same chartwise pole bounds.  Its one-form
component is the divided old component $\eta_n/\pi$ at that precision.
The change in the function component is a bounded closed ambient
representative followed by the Cech boundary of a bounded closed
one-cochain, before division.

The coset of the actual full class modulo ambient cohomology is
divided by $\pi$.  No smooth global $S_n$ or actual rational supported
lift of $b$ is required.
\end{proposition}

\begin{proof}
\textit{1. The ordinary lattice and the finite supported model.}
Lemma A29a proves, for every abstract model in this statement, that
all Hodge cohomology groups are finite free with arbitrary base change
and that the relative Hodge--de Rham spectral sequence degenerates at
$E_1$.  Its proof also states that the Hodge filtration has free
quotients with base change.  The truncation triangle therefore gives
the equality and the assertions about $\Lambda_R$ in (A31r.2), not just
the corresponding equality over $K$.  The same facts hold for the
abelian scheme $\mathcal X$.

The cup-product pairing with $F^3$ factors through $\Omega^{<2}$:
on the fourfold, forms of degree at least two multiply forms of degree
at least three to zero.  Hence it is an integral pairing on the
ordinary free lattice $\Lambda_R$.  Moreover
\[
       \langle\mathcal\kappa^*b,\tau\rangle
                       =\langle b,\mathcal\kappa_*\tau\rangle=0
                   \qquad(\tau\in\mathscr T_R).
\tag{A31r.5}
\]
This is an equality for actual classes, independent of a support
representative.

There is also an exact criterion for the integrality hypothesis on
$\beta$.  Serre duality on the Hodge graded pieces and A29a's free
degeneration give a perfect integral pairing
$\Lambda_R\times F^3H^4_{\mathrm{dR}}(\mathcal B/R)\to R$.
On the graded pieces it pairs form degrees $0,1$ with $4,3$,
respectively; both filtrations have free quotients.  Thus $\beta$
is integral if its periods against \emph{all} integral $F^3$ tests
are integral.  A canonical lower residue calculation may verify this
criterion without proving integrality of the full degree-four class.
Vanishing only against the primitive sublattice would not by itself
be this integrality criterion.

The support in (A31r.3) is the underlying closed set $S$ of the
special fibre.  Locally on $\mathcal B_n$ it can be defined by two
normal generators of any smooth graph lift, because the base ideal
is nilpotent.  Local cohomology of each locally free relative form
sheaf is concentrated in degree two.  Thus the usual Cech complex
of these local-cohomology sheaves computes (A31r.3).  We use it only
at finite precision.  The local graph bounds are not asserted to
define a global finite-pole sheaf on $\mathcal B_n$ when a global
smooth $S_n$ does not exist.

In a completed smooth graph chart, the bounded function modules have
the negative-monomial basis
\[
       x^{-a}y^{-b},\qquad a,b\geq1,\quad a+b\leq M,
\tag{A31r.6}
\]
over the smooth tangent coefficient ring; the bounded one-form
modules have the corresponding bases with relative differentials.
They and their successive pole quotients are flat over $A_n$.
Consequently a bounded cochain whose reduction is zero is uniquely
divisible by $\pi$ as a cochain over $A_{n-1}$ within the same bound.
The calculation descends from completed etale coordinates.  Applied
in every common graph chart, uniqueness makes the divided cochains
agree; the actual supported differential also commutes with division.
This is local flatness and saturation, not flatness of supported
hypercohomology.

\textit{2. The closed layer is ambient modulo a genuine closed boundary.}
The special fibre of $\xi_n$ has only the function component
$\kappa_0$, which is a Cech two-cocycle in the global special-fibre
sheaf $Z_M$.  Its image in $\mathbb H^2(C_M)$ has zero primitive
periods, since its forgotten ordinary class is the reduction of
$\beta$ and A29 supplies all the special primitive tests.
By A31n, if $T=\lfloor M/p\rfloor<4$ this image is zero.  In that case
put $\psi=0$ and use no ambient correction.  Otherwise, with $q=T-3$,
there exists
\[
                         \psi\in H^0(M_{q,X^{(1)}}^0)
\tag{A31r.7}
\]
whose composite (A31n.4) equals this image.  Injection (A31r.1)
upgrades the equality to equality in $H^2(Z_M)$ itself.  In either
case choose a Cech representative $\zeta_0$ of the corresponding
closed ambient class $F_B\rho_q(\psi)$, with $\zeta_0=0$ in the
zero case.  On an affine refinement there is
\[
         v_0\in\check C^1(Z_M),\qquad
                         \delta v_0=\kappa_0-\zeta_0.
\tag{A31r.8}
\]
This is a boundary in the actual closed-function sheaf.  It does not
assert that $\kappa_0$ itself is a boundary when $M\geq4p$.

\textit{3. Canonical closed representatives at the current precision.}
Since $n\leq\min(p,e)$, $A_n$ has characteristic $p$ and its maximal
ideal has $p$th power zero.  The coefficient section
\[
                 s:k\longrightarrow A_n,\qquad
                    c\longmapsto\widetilde{c^{1/p}}^{\,p}
\tag{A31r.9}
\]
is well-defined, and $F_{A_n}=sF_k\operatorname{red}$.
Therefore the relative Frobenius twist of the \emph{entire pair}
$\mathcal B_n\hookrightarrow\mathcal X_n$ is canonically the constant
pair $B^{(1)}\hookrightarrow X^{(1)}$ base-changed to $A_n$.
In particular the supports $S^{(1)}\times A_n$ exist on both twists.
This statement uses only the special complete-intersection equations;
it does not lift them to a complete-intersection presentation of the
given model.

The bounded source local-cohomology modules on these twists are
constant base changes.  For example on $B^{(1)}$ the source functions
of total pole at most $T$ are
$M_{T-1,B^{(1)}}^0$, and their relative Frobenius image is the global
closed module
\[
                         Z_M\otimes_k A_n.
\tag{A31r.10}
\]
If $T<2$ this module is zero.  The assertion for every $T$ is the same
local fraction calculation as in A31k: coefficients are Cartier
horizontal, both positive denominator exponents are multiples of
$p$, and the source total pole is at most $\lfloor M/p\rfloor$.
Powers of the source regular ideal are Cohen--Macaulay, so the Ext
description and its base change agree with these fractions.
The supported Cartier identity follows from the two local-cohomology
spectral sequences, exactly as in A31n.

In every local graph chart, replacing a normal generator by a
congruent generator modulo $\pi$ leaves its $p$th power unchanged
modulo $\pi^n$.  Thus (A31r.10) has the stated bound in every one of
the charts in the hypothesis.  Its Cech cochains and their boundaries
are formed in one global sheaf even though the graph ideals need not
glue.

Lift $v_0$ by (A31r.10) to a closed cochain $v_n$.
If $\psi\ne0$, its constant extension on $X^{(1)}\times A_n$ gives a
global closed ambient support class $\phi_n$ by relative Frobenius.
Its ordinary de Rham degree is four (local support degree four and
form degree zero).  Apply the function-only dual Koszul double boundary
to $\psi\otimes1$ on the constant pair and then relative Frobenius.
It gives a closed function two-cocycle $\zeta_n$ on $\mathcal B_n$,
reducing to the chosen $\zeta_0$, with function pole at most $M$ and
zero one-form component.  The degrees and bounds are those in
A31n.10 and A31n.18.  If necessary change its Cech representative by
the constant extension of a closed boundary to match $\zeta_0$.

Its forgotten ordinary class is
\[
       [\zeta_n]_{\mathrm{ord}}
                  =\mathcal\kappa_n^*[\phi_n]_{\mathrm{ord}}.
\tag{A31r.11}
\]
This equality is natural supported restriction on the constant
function complexes followed by the natural degree-zero relative
Cartier map.  It can equally be checked by the two ordinary Koszul
boundaries on the twist: the residue inclusion is
$w\mapsto w/(f_0f_1)$, and Frobenius gives the same function boundary
on the target.  Forgetting supports commutes with these maps.
The full finite transpose and its normal corrections were verified
in A31n; (A31r.11) does not select a different lift of its one-form
component.  Here all the source cochains on the twist are functions,
and their relative Frobenius images are actually closed.

\textit{4. Lifting the ordinary ambient class and its representatives.}
The group $H^4_{\mathrm{dR}}(\mathcal X/R)$ is finite free with base
change.  Choose an actual class $b$ whose reduction is
$[\phi_n]_{\mathrm{ord}}$.  If $\psi=0$, choose $b=0$.
This lifts an ordinary ambient cohomology class; it makes no
supported lifting claim.

The representative adjustment can be performed integrally.  Choose
a common affine Cech refinement for the ordinary de Rham complex and
the Cech localization complex defining the finite support.  Forgetting
supports is a map of these complexes.  Every ordinary form term on
an affine intersection is $R$-flat and surjects onto its reduction.
Choose a closed ordinary cocycle representing $b$.  Its reduction
and a forgotten cocycle representing $\phi_n$ differ by the total
differential of an ordinary degree-three cochain over $A_n$, since
their cohomology classes agree.  Lift that cochain term by term and
subtract its actual differential.  The result is an actual closed
ambient representative agreeing with the forgotten supported
representative modulo $\pi^n$.  The same procedure after restriction
to $\mathcal B$ realizes (A31r.11) at the cochain level, if desired.
It uses ordinary cochains; any intermediate comparison homotopy
need not have a uniform support pole bound.  The retained supported
representative is the already constructed bounded $\zeta_n$.
Its bound, and not a bound on an arbitrary ordinary homotopy, is what
will be preserved below.

\textit{5. Subtraction, division, and the exact preserved invariant.}
In the actual finite supported complex set
\[
                 \xi_n^{\mathrm{rem}}
                         =\xi_n-\zeta_n-Dv_n,
\tag{A31r.12}
\]
with the total sign for which the function component of $Dv_n$ is
$\delta v_n$.  Its one-form component is still $\eta_n$, because
$v_n$ and $\zeta_n$ are closed functions.  Its function component
reduces to zero by (A31r.8), and the one-form component reduces to
zero by hypothesis.  All retained cochains satisfy the original
bounds in every graph chart.  The flat saturation in (A31r.6)
therefore gives a bounded cocycle
\[
                          \xi'_{n-1}=\xi_n^{\mathrm{rem}}/\pi
\tag{A31r.13}
\]
over $A_{n-1}$.  The cocycle equations follow by dividing the actual
equations modulo $\pi^n$; uniqueness in the flat local modules
ensures all restrictions and changes of charts agree.

Let $b_B$ denote the image of $\mathcal\kappa^*b$ in $\Lambda_R$.
After forgetting supports, (A31r.11)--(A31r.13) give
\[
           (\beta-b_B)\bmod\pi^n
                      =\pi[\xi'_{n-1}]_{\mathrm{ord}}.
\tag{A31r.14}
\]
In particular $\beta-b_B\in\pi\Lambda_R$.  Put
$\beta'=(\beta-b_B)/\pi\in\Lambda_R$.
Because $\Lambda_R$ is free, multiplication by $\pi$ is injective
from $\Lambda_R/\pi^{n-1}$ to $\Lambda_R/\pi^n$.
Equation (A31r.14) proves
\[
                \beta'\bmod\pi^{n-1}
                              =[\xi'_{n-1}]_{\mathrm{ord}}.
\tag{A31r.15}
\]
This proves the maintained representative invariant without any
flatness assumption on supported hypercohomology.  It also excludes
an untracked ordinary closed-class correction after division.

The actual full class $a'$ is exactly (A31r.4), its lower image is
$\beta'$, and (A31r.5) proves that all its primitive periods still
vanish.  If a literal integral ordinary lower representative is
needed, choose a closed Cech cocycle representing $\beta'$; equality
(A31r.15) permits the same integral representative adjustment as in
step 4.  It represents the lower image of the actual rational full
class, not a substituted supported class.  The one-form component
of (A31r.13) is exactly $\eta_n/\pi$ at the new precision.
All assertions follow.
\end{proof}

\begin{corollary}[Initialization by the actual canonical lower cochains]
Suppose an actual canonical comparison of C44-A31-SA1 and A31i,
after a specified scalar multiplication, represents a rational class
$a$ whose lower one-form and function cochains are integral and
whose coefficient valuations tend to infinity with displacement
weight.  Suppose their reduction at the fixed precision $n$ is the
bounded actual supported cocycle required in Proposition A31r.
Then the image of $a$ in $\Lambda_R\otimes K$ is integral and reduces
to the forgotten class of that cocycle.  In particular the proposition
applies if its primitive periods vanish; integrality does not follow
merely from that vanishing assumption.
\end{corollary}

\begin{proof}
Use any integral test in the \emph{full} lattice
$F^3H^4_{\mathrm{dR}}(\mathcal B/R)$, with integral Cech/de Rham
representatives.  The local degree-two Gysin component contributes
zero: its product with a form of degree at least three has degree
at least five on a fourfold.  The remaining period is the actual
residue expression in the lower cochains, by the source comparison
and the chain-level trace identity checked in A31i.  Each term is
integral.  Its test jets are ordinary integral Taylor coefficients,
and residue extraction and the Cech trace introduce no inverse
factorial.  The coefficient valuations tend to infinity, so these
residue sums converge in $R$.  This establishes integral periods
against every integral $F^3$ test without imposing primitivity.
The perfect pairing in step 1 of A31r implies integrality of the
lower ordinary image.  At any fixed precision only the finitely
many retained terms contribute.  Forgetting their supported
representative is the reduction of the same actual ordinary class,
because the integral lower trace calculation gives equal pairings
with every reduced $F^3$ test.  Perfectness after base change
identifies the two classes in $\Lambda_R/\pi^n$.
This proves the asserted initialization.  No integrality of the
full crystalline or degree-four de Rham class has been deduced.
\end{proof}

\paragraph{Exact dependencies and remaining scope.}
Ordinary Hodge freeness, base change, degeneration, the Albanese
embedding, and the full primitive test lattice are respectively
C44-A29a and the primitive-lattice proposition in A29, whose proofs
use the previously checked A22a.  The annihilator and its actual
bounded transpose are C44-A31n.  The closed-function injection and
local flat pole calculation are A31i/A31k, or A31p for the uniform
algebraic injection.  The nilpotent-base coefficient section and
relative Cartier calculation are proved again in their needed scope
above.  No smooth support over $A_n$ is constructed, and no new
integral crystalline--de Rham comparison is invoked.

For precision greater than $p$, the constant-twist construction used
in step 3 need not exist without a supplied higher-order support.
Also A31p alone does not make the Serre threshold in A31n uniform in
$M$.  This proposition removes a closed ambient layer only under
its stated finite-bound and finite-precision hypotheses; it is not
an all-ramification assertion or a universal sign theorem.

% END work/scan_w/ambient_closed_layer_removal.tex

% BEGIN work/axis/axis_next_exact_ramification_boundary.tex
% C44_POS. Author, root and independent entire-proof readings PASS; scopes recorded in CHECK_EVENTS.
\section{The next exact ramification boundaries}
\label{sec:c44-a31s}

\paragraph{Scope and inputs.}
This is a continuation of the A31f--q lifting argument.
It uses the proved arbitrary-pole primitive test A31n and the
single ample-conormal open A31p.  It does not assume that a closed
supported class is a coboundary once its pole bound reaches \(4p\).
All closed layers removed below have smaller special-fibre pole bound.
The larger bound occurs only in the final class, where A31n detects
its one-form edge.

\begin{theorem}[C44-A31s]
Let \(k=\overline{\mathbf F}_p\), let \(X/k\) be an abelian sixfold,
let \(S\subset X\) be a smooth integral projective surface, and let
\(H\) be ample.  Suppose either \(p=5\) or \(p\geq7\), and put
\[
 c_p=\begin{cases}2p,&p=5,\\3p,&p\geq7,\end{cases}
 \qquad
 M_p=\begin{cases}p^2,&p=5\text{ or }7,\\4p,&p\geq11.\end{cases}
\tag{A31s.1}
\]
For every sufficiently large \(t\), there is a nonempty open of
smooth genuine axes \(B\subset X\), cut out by two members of
\(|H^t|\) and containing \(S\), with the following property.
Write
\[
 w=P_B[S]_B,\qquad
 P_B=1-\kappa^*(L_h^2)^{-1}\kappa_*,
 \qquad h=c_1(H^t).
\tag{A31s.2}
\]
Let \(R\) be a complete mixed-characteristic discrete valuation
ring with residue field \(k\), uniformizer \(\pi\), and absolute
ramification index \(e=v_\pi(p)\).  Let \(\mathcal B/R\) be any
smooth projective lift of \(B\).

If the canonical rational de Rham transport of \(w\) belongs to
\(\operatorname{Fil}^2 H^4_{\mathrm{dR}}(\mathcal B_K)\), and a
particular smooth embedded lift \(S_m\subset\mathcal B_m\) is
given, where
\[
                         m=\left\lceil e/c_p\right\rceil ,
\tag{A31s.3}
\]
then that support extends to a smooth projective characteristic-zero
lift after passing through the compatible finite-order extensions
constructed below.

Consequently, if \(S\) has no such geometric lift, then \(w\) is
not in \(\operatorname{Fil}^2\) for any such model with
\(e\leq c_p\).  More generally, on a class-lifting model at arbitrary
\(e\), a smooth embedded support through the order in (A31s.3)
cannot exist.
\end{theorem}

\begin{proof}
\leavevmode\par
\paragraph{1. The fixed geometric open and the tests.}
Intersect the A31p ample-conormal open with the finitely many
A31n testing opens needed for \(2p,3p,M_p\), and with the
smooth genuine-pencil and semiregularity opens already used in
A31q.  For \(p=5\), the \(3p\) test can harmlessly be included.
All are nonempty opens of the same irreducible parameter space
for sufficiently large \(t\).
This choice is independent of \(m\) and \(e\).

A29 supplies the relative Albanese embedding for an arbitrary
abstract model \(\mathcal B\), as well as the integral primitive
test lattice.  The A31i canonical graph construction supplies the
actual lower de Rham descent of the reduced cycle \([S]\).
For a particular supplied \(S_m\), local smooth formal extensions
can be chosen on affine adapted charts; their differences are
\(\pi^m(a,b)\) to first order.  They require no global
\(S_{m+1}\).
The obstruction to extending \(S_m\) is
\[
                         o_m\in H^1(S,N_{S/B}).
\tag{A31s.4}
\]
The canonical source is the unit of the local smooth support,
and its reduced Gysin is \([S]\) with multiplicity one.
It is not the fundamental cycle of the nonreduced Artin scheme
and is not \(m[S]\).

Write \(M_j^a=\mathcal Ext_B^2(\mathcal O_B/I^j,\Omega_B^a)\),
\(C_M=[M_{M-1}^0\to M_M^1]\), and \(Q_M=\operatorname{coker}d\),
as in A31p.
The following two already proved conclusions will be used in
different places.

First, for special pole bound \(M<4p\), a supported class with
zero primitive periods is zero in \(\mathbb H^2(C_M)\).
Together with the closed-summand injection of A31p, this makes
a closed-function Cech two-cocycle an actual Cech boundary in
the closed-function sheaf \(Z_M\).

Second, at the final bound \(M_p\), A31n says that a supported
class with zero primitive periods has zero one-form edge in
\(H^1(Q_{M_p})\).  A31p injects the Cartier summand
\(H^1(N^{(1)})\) into this group.
The coefficient map \(o\mapsto o^{(1)}\) here is the base-change
twist followed by inverse Cartier, not absolute Frobenius on
\(H^1(S,\mathcal O_S)\).

\paragraph{2. Moving closed cochains through two correction orders.}
We record the exact refinement of the root-cochain construction
in A31q.  Suppose \(S_m\subset\mathcal B_m\) is supplied and
\(n\leq e\).  Thus \(A_n=R/\pi^n\) has characteristic \(p\).
A special closed Cech cochain \(b\in Z_{2p}\) has a Frobenius
root of total normal pole at most two.  Lift that root, cochain
by cochain on affine intersections, to
\[
 \mathcal Ext^2_{\mathcal B_m}
                (\mathcal O_{S_m},\mathcal O_{\mathcal B_m}).
\tag{A31s.5}
\]
This line on \(S_m\) is \(A_m\)-flat, and reduction of sections
on each affine intersection is surjective.
Choose local smooth formal extensions of \(S_m\), lift the
resulting fractions to \(R\), and take their actual \(p\)-th
powers.  Call the resulting analytic Cech cochain \(\widehat b\).
Its differential is divisible by \(p\).

In any common graph chart, its pole bound through precision
\(n\leq 2pm+1\) is
\[
 \widehat b\in
 M_{2p-1}^0+\pi^{pm}M_{3p-1}^0+\pi^{2pm}M_{4p-1}^0.
\tag{A31s.6}
\]
This is a local bound on an actual cochain, not an assertion
that the closed sheaf is a constant family.

Here is the complete correction calculation.
Two local extensions of the same \(S_m\) differ by \(\pi^m\).
Expanding an inverse normal parameter, or a coefficient in an
etale tangent chart, gives an integral Taylor expansion: its
term of order \(q\) has coefficient divisible by \(\pi^{qm}\)
and raises total normal pole by at most \(q\).
This can be seen directly by the geometric-series expansion of
\((x+\pi^m a)^{-u}(y+\pi^m b)^{-v}\); no factorial is inverted.
The same assertion holds for an arbitrary etale change by its
integral power-series expansion.
For a simple normal fraction the first three pole bounds are
therefore \(2,3,4\).
Modulo \(p\), taking a \(p\)-th power is additive and multiplies
these orders and poles by \(p\).
All subsequent terms have order at least \(3pm\), which is
at least \(2pm+1\).
This proves (A31s.6), including its behaviour under restriction
and the Cech differential.

Changes of generators for the same regular immersion preserve
the Ext section and its intrinsic pole filtration.
Fractions with only one support denominator are zero in local
cohomology before taking the Frobenius power.  Changing a
representative by such a fraction has the same effect modulo
\(p\); the mixed binomial terms are divisible by \(p\).
This is sufficient since \(n\leq e\).
Consequently no coordinate-dependent representative has been
made canonical.

There is also the following version for \(b\in Z_{3p}\).
Its root has total pole at most three and lifts on affine
intersections to
\[
 \mathcal Ext^2_{\mathcal B_m}
       (\mathcal O_{\mathcal B_m}/I_{S_m}^2,
                                      \mathcal O_{\mathcal B_m}).
\tag{A31s.7}
\]
This module and its reduction commute with the base change:
the power of an ideal generated by a regular sequence has the
usual finite locally free resolution, and its relative pole
graded pieces are locally free on \(S_m\).
For \(n\leq pm+1\), the same calculation gives
\[
                \widehat b\in
                  M_{3p-1}^0+\pi^{pm}M_{4p-1}^0.
\tag{A31s.8}
\]
Again \(d\widehat b\) is divisible by \(p\).

Finally, once \(n\leq pm\), A31o supplies the actual Frobenius
support
\[
 S'_n=S_m\times_{\varphi_{m,n}}A_n
               \subset\mathcal B_n^{(1)},\qquad
 \varphi_{m,n}(a)=a^p.
\tag{A31s.9}
\]
Its bounded Ext modules lift special closed cochains with pole
\(3p\) without any moving correction.
The map in (A31s.9) is defined by \(R\to A_n\) and factors
through \(A_m\); it does not require a map \(A_n\to A_m\).

All divisions by \(\pi\) below are performed on cochains.
In a common graph chart the bounded local-cohomology modules
have their negative-normal-monomial basis, with flat successive
pole quotients.  The weighted sums in (A31s.6) and (A31s.8)
are therefore saturated for the indicated division, with the
displayed exponents lowered by one.  The preceding change-of-chart
calculation proves the corresponding bounds in every chart.
No flatness of supported hypercohomology is used.

\paragraph{3. All displacement weights.}
It suffices first to treat \(e=pm+r\), where
\[
 0<r\leq
 \begin{cases}pm,&p=5,\\2pm,&p\geq7.\end{cases}
\tag{A31s.10}
\]
The smaller range \(e\leq pm\) is the already proved A31l range
or the final A31 induction.
For a positive displacement weight \(N\), put
\[
 a(N)=v_p(N)+\lfloor\log_p(N-1)\rfloor\quad(N\geq2),
 \qquad a(1)=0.
\]
The canonical radial proof in A31i gives the lower bounds
\[
 v_1(N)=mN-e\,v_p(N),\qquad
 v_2(N)=mN-e\,a(N)
\tag{A31s.11}
\]
for the one-contraction and two-contraction terms respectively.
The outer integration denominator is \(N\); the inner one
has weight strictly less than \(N\).
These are bounds on every summand, including frame and
composition homotopies.

The one-form minimum is \(-r\), attained only at \(N=p\).
For two contractions the minimum is \(-2r\), at \(N=2p\).
All terms through target \(-r\) are as follows:
\[
\begin{array}{c|c|c}
 &\text{weight}&v_2(N)\\ \hline
 p\geq7&p&-r\\
       &2p&-2r\\
       &3p&pm-2r\\
       &4p&2pm-2r\\
 p=7   &p^2&4pm-3r\\ \hline
 p=5   &p&-r\\
       &2p&-2r\\
       &3p&pm-2r\\
       &p^2&2pm-3r
\end{array}
\tag{A31s.12}
\]
In the first block \(4p\), and \(p^2\) when \(p=7\), are
at the target only if \(r=2pm\); otherwise they are above it.
In the \(p=5\) block \(3p,p^2\) reach the target only if
\(r=pm\).
Rows above the target may be omitted; retaining them with
their true positive power of \(\pi\) is harmless.

For clarity, we prove that no weight has been lost.
After adding \(r\), and dividing by \(m\), the two-contraction
bound is
\[
 N-p+(1-a(N))\,e/m.
\tag{A31s.13}
\]
For \(a(N)\geq1\) it is decreasing in \(e/m\).
It is enough to examine the endpoints \(e/m=3p\) for \(p\geq7\)
and \(e/m=2p\) for \(p=5\).

At \(3p\), the weights below \(p^2\) which are not multiples of
\(p\) give a strictly positive value in (A31s.13).
The multiples \(N=bp\), \(2\leq b<p\), give \(p(b-4)\);
the listed \(2p,3p,4p\) are exactly the nonpositive ones.
At \(N=p^2\) the value is \(p(p-7)\), zero only for \(p=7\).
For \(p^2<N<p^3\), values with \(v_p(N)\leq1\) are bounded
below by \(p^2-6p>0\), and those with \(v_p(N)=2\) by
\(2p^2-10p>0\).
For a power \(N=p^a\), \(a\geq3\), the value is
\[
                        p^a-p(6a-5)>0.
\tag{A31s.14}
\]
It is positive at \(a=3,p=7\) and increases with both parameters;
the increment in \(a\) is \((p-1)p^a-6p>0\).
On the strict interval \(p^a<N<p^{a+1}\), if \(v_p(N)\leq a-1\)
the bound is larger than (A31s.14); if \(v_p(N)=a\), it is at
least \(2p^a-p(6a-2)\), also larger.
This proves the first block.

At \(p=5,e/m=2p\), a nonmultiple below \(p^2\) again gives a
positive value; the multiples \(bp\), \(2\leq b<p\), give
\(p(b-3)\).
The power \(p^2\) gives \(p(p-5)=0\).
For a power \(p^a\), \(a\geq3\), the value is
\(p^a-p(4a-3)>0\), and the same strict-interval argument proves
positivity between powers.  On \(p^2<N<p^3\), the two lower
bounds are respectively \(p^2-5p+1>0\) and \(2p^2-7p>0\).
Thus the second block is complete.
The one-contraction assertion follows in the same way from
\(N-p-(v_p(N)-1)e/m\):
for \(a=v_p(N)\geq2\), it is bounded below by
\(p^a-p(3a-2)>0\) in the first case, and
\(p^a-p(2a-1)>0\) in the second.
For \(a=1\), only \(N=p\) gives zero.

\paragraph{4. The actual lower period and its leading edge.}
Scale the actual canonical lower descent by \(\pi^{2r}\)
and retain precision \(r+1\).
This precision is at most \(e\).
Its one-form component \(\eta\) is divisible by \(\pi^r\).
After dividing by \(\pi^r\), its leading edge is a nonzero
unit times the inverse-Cartier image of \(o_m^{(1)}\).
The general variable-coefficient identity in A31j proves this
for the radial representative, including the tangent differentials;
its difference from the pure Cartier expression is the differential
of a bounded function.
All higher one-contraction weights vanish at the retained precision.

For \(p\geq7\), the function component \(\kappa\) has the bound
\[
 M_{2p-1}^0+
        \pi^{pm}M_{3p-1}^0+
        \pi^{2pm}M_{M_p-1}^0
                       \pmod{\pi^{r+1}}.
\tag{A31s.15}
\]
For \(p=5\), its bound is
\[
       M_{2p-1}^0+\pi^{pm}M_{M_p-1}^0
                       \pmod{\pi^{r+1}}.
\tag{A31s.16}
\]
These are bounds through the target coefficient, not just
bounds for the most negative coefficient.
At \(p=7\), the \(p^2\)-weight has scaled order
\(4pm-r\geq2pm\); at \(p=5\), it has order
\(2pm-r\geq pm\), as required.

The Filtration hypothesis says that all integral primitive
test periods of this actual rational class are zero.
The \(\Omega^2\) part has zero product with a primitive
degree-three test on the fourfold, so the lower descent computes
these periods integrally at the precision just specified.
This does not assert integrality of the full rational
crystalline--de Rham comparison.
Subtracting an actual total boundary and dividing by \(\pi\)
preserves the vanishing of these rational periods.

\paragraph{5. Removing the closed layers for \(p\geq7\).}
At stage \(j<r\), the current precision is \(r+1-j\) and the
one-form component starts in order \(r-j\).
For \(j<\min(r,pm)\), reduction of (A31s.15) is a closed-function
two-cocycle of pole at most \(2p\).
The small-pole primitive test and A31p show that it equals
\(\delta b_j\) for an actual cochain \(b_j\in Z_{2p}\).
Use (A31s.6) to lift \(b_j\), subtract its total boundary,
and divide the entire supported cochain by \(\pi\).
Its derivative is zero at the required precision because it
is divisible by \(p\).
Saturation gives the next bound
\[
 M_{2p-1}^0+
    \pi^{pm-j-1}M_{3p-1}^0+
    \pi^{2pm-j-1}M_{M_p-1}^0.
\tag{A31s.17}
\]
Thus this induction is valid for all the first \(\min(r,pm)\)
layers.

If \(r\leq pm\), it already reaches the target.
If \(r>pm\), the special layer at \(j=pm\) has pole at most
\(3p<4p\), so the same argument gives \(b_j\in Z_{3p}\).
When \(r<2pm\), its precision is at most \(pm\), and the global
Frobenius support (A31s.9) supplies its closed lift.
When \(r=2pm\), its precision is \(pm+1\); use (A31s.8) for
this one step.  Its extra term has order \(pm\) and pole \(4p\),
hence lies in the last summand of (A31s.15).
After subtraction and division the precision is \(pm\).

Every remaining stage has precision at most \(pm\).
Lift the \(Z_{3p}\) cochains on (A31s.9), using the bounded
Ext module of \(I_{S'_n}^2\).
There is no further moving correction.
Before the target, the \(M_p\)-pole summand still has positive
\(\pi\)-order, so the reduced closed layer always has pole
at most \(3p\).
All \(r\) removals are therefore justified by the small-pole
test, without applying it at bound \(4p\) or larger.

At the final stage we obtain a class in
\(\mathbb H^2(C_{M_p})\) with zero primitive periods and with
one-form edge equal to the nonzero unit multiple of
\(c(o_m^{(1)})\).
A31n kills that edge and A31p makes its Cartier map injective.
Thus \(o_m=0\), and the particular \(S_m\) extends to \(S_{m+1}\).

\paragraph{6. Removing the closed layers for \(p=5\).}
Here \(r\leq pm\).
For every \(j<r\), (A31s.16) has reduced closed-function pole
at most \(2p\).
Lift the boundary cochain in \(Z_{2p}\) by the root construction.
The only possible correction through precision \(r+1-j\)
has order \(pm\) and pole \(3p\), contained in \(M_p=p^2\).
It reaches the target only when \(r=pm\).
Hence all \(r\) removals use the \(2p<4p\) test.
At the target A31n with \(M_p=25\) and the A31p Cartier
injection again give \(o_m=0\).

This proves the next-support assertion throughout the full
scaled intervals \(pm<e\leq c_pm\), not just at their endpoints.

\paragraph{7. Iteration and class-level scope.}
Start with \(m=\lceil e/c_p\rceil\).
At every subsequent order \(a\geq m\), one has \(e\leq c_pa\).
If \(e>pa\), apply the interval just proved.
If \((p-1)a\leq e\leq pa\), apply A31l.
Once \(a>e/(p-1)\), A31 applies directly and constructs all
remaining compatible support lifts.
Its algebraization argument produces the claimed smooth
projective characteristic-zero support.
No initial support beyond the one explicitly supplied has
been inferred.

If \(e\leq c_p\), the required initial support is \(S_1=S\),
which is automatic.
For a geometrically nonliftable \(S\), this contradicts the
Filtration hypothesis on every abstract model in that range.

For the explicit A14b surfaces at primes \(p\equiv3\pmod4\),
the intersections with the A18 closed-point selection open
give actual nonzero primitive algebraic classes, and may be
chosen with \(A_B=\mathbf Qw\).
A18's common-realization kernel then turns any claimed
characteristic-zero rational cycle-class lift into the
Filtration hypothesis used above, regardless of its chosen
special representative.
Thus the examples include \(p=7,e\leq21\) and
\(p=11,e\leq33\).
The \(p=5\) theorem is valid under its stated nonliftable-centre
hypothesis; A14b alone does not provide that hypothesis at \(p=5\).
\end{proof}

\paragraph{Programme status and the next unresolved step.}
A31s sharpens the supplied-support threshold to
\(\lceil e/(3p)\rceil\) for \(p\geq7\), and to
\(\lceil e/10\rceil\) for \(p=5\).
It excludes the corresponding initial ranges of ramification
for the actual A14b/A18 primitive classes where those centres
exist.  It does not exclude arbitrary ramification, or assert
that the improved initial support exists.

In Grothendieck's terms this is a route-(a) obstruction for the
special axis component in QH9-axis, with the primitive correction
at cote 44, scans 008--009 and the H9-T7 reduction of Theorem 5,
scan 004.  It is not a universal H9-T7 sign proof or a scan 049
product-form assertion.
The argument uses neither of the unrefereed 032/Milne2026
preprints and does not change the independent ambient terms.
It belongs to the same developing A31 lifting advance.

At larger ramification, a closed layer of pole at least \(4p\)
can occur before the final one-form coefficient.
A31n then puts its class in a closed-function image, but does
not make it a coboundary.
Removing that ambient image at the required finite precision
is a separate step.  No conclusion is inferred from merely
continuing the pole pattern.

% END work/axis/axis_next_exact_ramification_boundary.tex

% BEGIN work/axis/axis_characteristic_three_ramification_five.tex
% C44_POS. Author, root and independent entire-proof readings PASS; scopes recorded in CHECK_EVENTS.
\section{The first ambient closed layer in characteristic three}
\label{sec:c44-a31t}

\paragraph{Scope.}
This applies the finite-precision ambient-removal proposition A31r.
The new case is \(p=3,e=5\).
The first three closed layers are genuine bounded boundaries.
The fourth may contain an ambient class; it is removed as an
actual ordinary ambient class, rather than declared to be a
supported boundary.
This is part of the same developing A31 lifting argument.

\begin{theorem}[C44-A31t]
Let \(k=\overline{\mathbf F}_3\), let \(X/k\) be an abelian
sixfold, let \(S\subset X\) be a smooth integral projective
surface with no geometric characteristic-zero lift, and let
\(H\) be ample.
For every sufficiently large \(t\), there is a nonempty open
of smooth genuine axes \(B\subset X\), containing \(S\) and
cut out by two members of \(|H^t|\), for which the nonzero class
\[
 w=P_B[S]_B,\qquad
 P_B=1-\kappa^*(L_h^2)^{-1}\kappa_*,
 \qquad h=c_1(H^t),
\tag{A31t.1}
\]
has the following property.

For every complete mixed-characteristic discrete valuation ring
\(R\) with residue field \(k\) and \(v_\pi(3)\leq5\), and every
abstract smooth projective lift \(\mathcal B/R\) of \(B\), the
canonical rational de Rham transport of \(w\) is not in
\(\operatorname{Fil}^2H^4_{\mathrm{dR}}(\mathcal B_K)\).
The A14/A18 examples can be chosen in this open with
\(A_B=\mathbf Qw\).  In those examples, the primitive rational
cohomology class cannot be the specialization of a
characteristic-zero rational algebraic cycle class on any
such model.
\end{theorem}

\begin{proof}
\leavevmode\par
\paragraph{1. The model and its maintained ordinary class.}
The cases \(e\leq4\) are A31k, so assume \(e=5\).
Choose the axis in the intersection of the A31p open,
the A31n opens for \(M=9,18\), and the previously established
smooth genuine-axis and semiregularity opens.
For high \(t\) these are nonempty opens in the same irreducible
parameter space.  A06 also makes \(w\) nonzero.
A29 gives the relative Albanese embedding for any abstract
model \(\mathcal B/R\), the free Hodge lattice, and the
integral primitive test lattice \(\mathscr T_R\).

Suppose for a contradiction that \(w\) is filtered as in the
statement.
Let \(\gamma\in H^4_{\mathrm{dR}}(\mathcal B_K)\) be the actual
canonical rational transport of \([S]\).
Its pairing with every \(\mathscr T_R\) test is zero:
the primitive part is in \(F^2\), and the remaining ambient
part is orthogonal to \(\mathscr T_R\).
We will maintain an actual rational full class \(a_j\), its
integral ordinary lower image
\[
 \beta_j\in
 \Lambda_R=\mathbb H^4(\mathcal B,\Omega_{\mathcal B/R}^{<2})
                 =H^4_{\mathrm{dR}}(\mathcal B/R)/F^2,
\tag{A31t.2}
\]
and a bounded representative of its reduction in the actual
finite supported complex.
The lattice \(\Lambda_R\) is free and has base change.
The full class need not be integral and need not have a rational
supported lift after an ambient subtraction.

Write \(M_\ell^a=\mathcal Ext_B^2(\mathcal O_B/I^\ell,\Omega_B^a)\)
and \(C_M=[M_{M-1}^0\to M_M^1]\).
A total function pole \(M\) is therefore represented in
\(M_{M-1}^0\).
The lower supported components are denoted \((\eta,\kappa)\):
one-forms in Cech degree one and functions in Cech degree two.

\paragraph{2. Every displacement weight through the target.}
The local smooth formal supports in the canonical A31i
graph construction need only lift the special surface locally.
Their first differences are \(\pi(a,b)\); no global
\(S_2\) is assumed.
For a weight \(N\), put
\[
 A(N)=v_3(N)+\lfloor\log_3(N-1)\rfloor\quad(N\geq2),
 \qquad A(1)=0.
\]
The actual one- and two-contraction terms have the valuation lower bounds
\[
                  V_1(N)=N-5v_3(N),\qquad
                  V_2(N)=N-5A(N).
\tag{A31t.3}
\]
The one-form minimum is \(-2\), only at \(N=3\).
The function minimum is \(-6\), only at \(N=9\).
The complete list of function weights at or below the
one-form target \(-2\) is
\[
\begin{array}{c|rrrrr}
 N&3&6&9&12&18\\ \hline
 V_2(N)&-2&-4&-6&-3&-2 .
\end{array}
\tag{A31t.4}
\]

Here is an all-weight check.
For \(1\leq N\leq26\), direct use of \(v_3(N)\) and the three
possible logarithmic intervals gives (A31t.4); the only other
negative value is \(V_2(4)=-1\).
At a power \(N=3^a\), \(a\geq3\),
\[
                   V_2(N)+2=3^a-10a+7>0.
\tag{A31t.5}
\]
It is positive at \(a=3\) and its increment is
\(2\cdot3^a-10>0\).
For \(3^a<N<3^{a+1}\), if \(v_3(N)\leq a-1\), the value is
larger than (A31t.5); if \(v_3(N)=a\), it is at least
\(2\cdot3^a-10a+2\), again larger.
For one contraction, at \(v_3(N)=a\geq2\) its minimum is
\(3^a-5a\geq-1\), and this increases after \(a=2\).
When \(a=1\), only \(N=3\) has value \(-2\).
This proves every assertion and also shows that the coefficient
valuations tend to infinity with \(N\).

Set \(a_0=\pi^6\gamma\) and retain its canonical lower descent
modulo \(\pi^5\).
All its lower cochains are integral, including weights beyond
the retained precision.
Its one-form component starts at order four, and its function
component has the exact through-target bound
\[
 \eta_0\in\pi^4 M_3^1\pmod{\pi^5},\qquad
 \kappa_0\in
      M_8^0+\pi^3M_{11}^0+\pi^4M_{17}^0
                              \pmod{\pi^5}.
\tag{A31t.6}
\]
The terms of weights \(3,6,9\) are in the first summand,
weight \(12\) first appears at order three, and weight \(18\)
at order four.
The leading one-form edge is
\[
              \pi^{-4}\eta_0\bmod\pi
                     =u^{-1}c(o_1^{(1)})
                            \quad\hbox{in }H^1(Q_{18}),
 \qquad 3=u\pi^5 ,
\tag{A31t.7}
\]
up to the fixed harmless sign of the homotopy convention.
Here \(o_1\in H^1(S,N_{S/B})\) is the obstruction to \(S_2\).
The variable-coefficient cubic identity of A31i includes
\(da,db\); its exact-function correction does not change this
edge.  The superscript \((1)\) is base change by the field
Frobenius, not Frobenius acting on coherent cohomology.

The initialization corollary of A31r applies.
Indeed the lower residues pair integrally with every integral
\(F^3\) test, not just the primitive tests, and converge by
(A31t.5).  The degree-two Gysin component has zero product
with such a test on a fourfold.
Full integral Serre duality therefore makes the lower image
\(\beta_0\) of \(a_0\) integral in \(\Lambda_R\), and its
reduction modulo \(\pi^5\) is the forgotten class of the
actual bounded supported cocycle just constructed.
This is not a claim that the full comparison class is integral.

\paragraph{3. Three closed boundaries with a moving pole correction.}
For \(j=0,1,2\), we maintain current precision \(5-j\),
one-form divisibility \(\pi^{4-j}\), and the function bound
\[
       M_8^0+\pi^{3-j}M_{11}^0+\pi^{4-j}M_{17}^0.
\tag{A31t.8}
\]
Modulo \(\pi\) the cocycle has only a closed function component
of total pole at most nine.
Since \(9<4p=12\), the small-pole primitive test makes its
class zero in \(\mathbb H^2(C_9)\).
The closed-summand injection of A31p then gives a genuine
Cech one-cochain \(b_j\in Z_9\) whose boundary is this layer.

A closed \(Z_9\) fraction has a Frobenius root of total normal
pole at most three.
On each affine intersection choose an integral lift of its root
in the local module
\[
             \mathcal Ext^2(\mathcal O/I_S^2,\mathcal O)
\tag{A31t.9}
\]
along a local smooth formal extension of the special support,
and take its actual cube.
This gives an actual cochain \(\widehat b_j\) with differential
divisible by \(3\).
In any common graph chart its bound through precision five is
\[
                        M_8^0+\pi^3M_{11}^0.
\tag{A31t.10}
\]
For completeness, an order-one graph or frame discrepancy
changes the root by \(\pi\) times a fraction with one higher
normal pole.  Cubing modulo \(3\) puts this term in order
\(\pi^3\) and pole at most twelve.
The second root discrepancy has order \(\pi^2\), whose cube
is \(\pi^6=0\) at the retained precision.
The expansion of inverse normal parameters has integral
coefficients and requires no factorial division.
Changing a representative by a one-denominator fraction has
no local-cohomology effect after cubing modulo \(3\), because
all mixed terms are divisible by \(3\).
This proves (A31t.10) under restriction and Cech boundary,
without a global smooth support or constant Frobenius twist
through precision five.

Subtract the actual total boundary \(D\widehat b_j\) and
divide the retained cochain by \(\pi\).
Its one-form part is unchanged before division: the derivative
of \(\widehat b_j\) is zero at current precision \(5-j\).
The local flat negative-monomial bases and their saturated
pole quotients give exactly (A31t.8) with \(j\) replaced by
\(j+1\).
The first correction in (A31t.10) is no earlier than the
displayed order \(3-j\), so the weighted bound is preserved.

At the ordinary level the rational full class is
\(a_{j+1}=a_j/\pi\).
The boundary does not change it.
The forgotten special class of \(\beta_j\) is zero, since
the finite supported special layer is a genuine boundary.
Freeness gives \(\beta_j/\pi\in\Lambda_R\).
Multiplication by \(\pi\) from
\(\Lambda_R/\pi^{4-j}\) to \(\Lambda_R/\pi^{5-j}\) is
injective, so the new finite supported cocycle forgets to
the reduction of exactly this divided ordinary class.
Thus the required ordinary invariant, not merely the primitive
periods, is maintained at all three steps.

\paragraph{4. The first nonzero ambient layer.}
After these steps the precision is two, the one-form component
starts at order one, and
\[
                  \kappa_3\in M_{11}^0+\pi M_{17}^0,
                     \qquad \eta_3\in\pi M_3^1
                                      \pmod{\pi^2}.
\tag{A31t.11}
\]
Its special closed-function layer can have pole twelve.
The small-pole test is not applicable to this layer.

Apply A31r with \(M=18,n=2\).
All its hypotheses have been maintained: \(a_3\) is an actual
rational ordinary class with zero primitive periods;
\(\beta_3\) is integral; and (A31t.11) is its actual bounded
finite supported representative.
The inequality \(2\leq\min(p,e)\) holds.
A31r gives an actual ambient class
\(b\in H^4_{\mathrm{dR}}(\mathcal X/R)\), a new full class
\[
                     a_4=(a_3-\mathcal\kappa^*b)/\pi,
\tag{A31t.12}
\]
and a supported cocycle over \(k\) with bound \(C_{18}\).
Its ordinary lower image is the integral image of exactly
(A31t.12), and all its primitive periods are zero.
The correction is a closed ambient representative plus an
actual closed boundary at precision two.
Consequently its one-form component is exactly \(\eta_3/\pi\);
the edge (A31t.7) has not been changed.

A31n at \(M=18\) now says that this edge is zero in
\(H^1(Q_{18})\).
The A31p Cartier injection implies \(o_1=0\).
This forces a smooth embedded lift \(S_2\subset\mathcal B_2\).

\paragraph{5. Support induction and actual classes.}
Return to the original Filtration hypothesis on \(w\).
With the supplied \(S_2\), the A31l interval
\[
                       (p-1)2=4<e=5<p\,2=6
\]
forces \(S_3\).
This is the initial order
\(\lfloor e/(p-1)\rfloor+1=3\) required by A31.
That theorem extends the support to all orders and algebraizes
it, contradicting the geometric nonliftability of \(S\).
The ambient subtraction in (A31t.12) has not changed the
original hypothesis used for this final support induction.

For the explicit A14 surface, all conditions imposed on \(B\)
are nonempty opens of the constrained parameter space.
Intersecting them with the A18 closed-point selection gives
actual axes with \(A_B=\mathbf Qw\), and A06 supplies the
positive nonzero line.
The common-realization kernel of A18b/c identifies a proposed
rational cycle-class lift with the same canonical primitive
crystalline class even if its special cycle representative is
different.  Such a lift would put \(w\) in \(F^2\), which was
just excluded.  Filtration membership is unchanged by testing
after a finite extension of the generic field of a fixed
model, by faithful base change of the Hodge filtration.
This proves the class-level assertion for the stated models.
\end{proof}

\paragraph{Programme status.}
This excludes the next genuine ramification case \(p=3,e=5\)
for the existing A14/A18 primitive classes, hence all \(e\leq5\)
in characteristic three.
It does not assert an arbitrary scaled initial-support
threshold at \(e=5m\), since the ambient-removal step would then
require its own precision check.
For higher \(e\), further closed layers and, eventually,
additional leading Cartier weights must be analyzed.

The result is a route-(a) obstruction in QH9-axis for the actual
primitive term of H9-T7, corresponding to cote 44,
Theorem 5 scan 004 and the primitive margin scans 008--009.
It is not a universal sign proof or a scan 049 product-form
assertion.  It uses neither of the unrefereed 032/Milne2026
preprints and makes no new claim about the separate ambient term.
It is grouped with the single A31f--s lifting advance.

% END work/axis/axis_characteristic_three_ramification_five.tex

% BEGIN work/scan_w/higher_cartier_detection.tex
% C44_POS. Author, root and independent entire-proof readings PASS; scopes recorded in CHECK_EVENTS.
\section{Higher Cartier filtrations and bounded normal-obstruction detection}
\label{sec:c44-a31u}

\paragraph{Scope.}
This is a statement about the special-fibre supported de Rham complex.
It repairs the failure of a naive pure $p^a$-power normal formula when
$a\geq2$.  It does not construct an actual ramified comparison at a new
precision, lift a global ambient supported section, or justify an
unbounded sequence of ambient removals.  Those are separate requirements.
In particular the sufficiently-high-degree hypothesis in A31n still
depends on its fixed pole bound.

\begin{proposition}[C44-A31u: higher Cartier detection]
Let $k=\overline{\mathbf F}_p$, $p>2$, let $S$ be a smooth integral
projective surface in an abelian sixfold $X$, and let $H$ be ample.
For every sufficiently large $t$, there is a nonempty open of smooth
genuine axes $B\subset X$ containing $S$, cut out by two members of
$|H^t|$, with the following properties.  Put
\[
 I=I_{S/B},\quad N=N_{S/B},\quad
 M_m^i=\mathcal Ext_B^2(\mathcal O_B/I^m,\Omega_B^i),\quad
 \mathscr M^i=\mathcal H_S^2(\Omega_B^i)=\bigcup_mM_m^i.
\tag{A31u.1}
\]
Write $\nu=[dx\wedge dy/(xy)]$ for the canonical fundamental
fraction of $S$, and let
\[
 j:N\longrightarrow M_1^1,\qquad
 (a,b)\longmapsto
       [a\,dy/(xy)-b\,dx/(xy)]
\tag{A31u.2}
\]
be the ordinary contraction.  The formulas are independent of the
choice of the two regular normal parameters and their local lifts.

\begin{enumerate}
\item For every $m\geq1$,
\[
                 H^0(B,M_m^2)=k\nu.
\tag{A31u.3}
\]
Consequently every coherent subsheaf of $\mathscr M^2$ whose sections
are locally exact supported two-forms has no global sections.

\item For every $m\geq1$, the map
$H^1(S,N)\to H^1(B,M_m^1)$ induced by $j$ is injective.
For every $m\geq2$, the map
\[
            H^1(S,N)\xrightarrow{\ d j\ }H^1(B,M_m^2)
\tag{A31u.4}
\]
is also injective.  The second map is interpreted after Frobenius
pushforward, so that it is a morphism of coherent modules on the
Frobenius twist; its assertion on cohomology is equivalently an
assertion about the underlying additive groups.

\item Use the supported Cartier isomorphism to define
$\mathscr Z_a^1\subset\mathscr M^1$ and $\mathscr B_a^1\subset
\mathscr Z_a^1$ recursively, with
$\mathscr Z_0^1=\mathscr M^1$, $\mathscr B_0^1=0$,
$\mathscr Z_1^1=\ker d$, $\mathscr B_1^1=\operatorname{im}d$, and
\[
 C^{-1}(\mathscr Z_{a-1}^1)=\mathscr Z_a^1/\mathscr B_1^1,
 \qquad
 C^{-1}(\mathscr B_{a-1}^1)=\mathscr B_a^1/\mathscr B_1^1.
\tag{A31u.5}
\]
The sheaves on the right of a Cartier map are on the corresponding
field twist.  Fix $M\geq2$ and set
\[
 \begin{split}
 Q_M&=M_M^1/dM_{M-1}^0,\\
 K_{0,M}&=Q_M,\qquad
 K_{a,M}=(\mathscr Z_a^1\cap M_M^1)/dM_{M-1}^0
                    \quad(a\geq1).
 \end{split}
\tag{A31u.6}
\]
For every $a\geq1$, the inclusion gives an injection
\[
                  H^1(B,K_{a,M})\hookrightarrow H^1(B,Q_M).
\tag{A31u.7}
\]
It is important that the denominator in (A31u.6) is the indicated
bounded exact-form sheaf; it is not silently replaced by the
intersection with all unbounded exact forms.
\end{enumerate}

These assertions hold on one open simultaneously for every finite
$m,M,a$.  Suppose, in addition, that an actual bounded one-form edge
$\eta$ is represented in $K_{a,M}$ and that its image in $H^1(Q_M)$
is zero.  If, on the $a$-th field twist, its representative satisfies
\[
                    d C^a\eta=\lambda\,d j(o),
                    \qquad \lambda\in k^*,\quad
                    o\in H^1(S^{(a)},N^{(a)}),
\tag{A31u.8}
\]
as a cohomology identity in some finite supported two-form module,
then $o=0$.  In (A31u.8) an identity of actual Cech cocycles suffices.
The hypothesis that the bounded edge really lies in $K_{a,M}$,
and the stated identity after Cartier, must be verified for an
application; they are not inferred merely from a pole estimate.
\end{proposition}

\begin{proof}
\textit{1. One ample-conormal open and the pole filtration.}
Apply the uniform twisted-vanishing and geometric construction of
A31p with the finite collection
\[
                         \mathcal O_S,\quad\Omega_S^1,
                         \quad\Omega_S^2.
\tag{A31u.9}
\]
Thus $N^*\otimes A^{-c}$ is ample for a fixed sufficiently large
$c$, and
\[
 H^0(S,R(N)\otimes\Omega_S^i)=0\quad(i=0,1,2)
\tag{A31u.10}
\]
for every positive-degree homogeneous polynomial representation
$R$ of $\operatorname{GL}_2$.  All field twists have the same
vanishing.  Only this fixed collection is used, regardless of
the subsequent pole bound or number of Cartier operations.

The regular-ideal pole filtration has exact associated quotients
\[
 M_m^i/M_{m-1}^i
   =\det N\otimes\Gamma^{m-1}N\otimes\Omega_B^i|_S,
                 \qquad M_0^i=0.
\tag{A31u.11}
\]
This follows by dualizing the homogeneous normal ideal quotients:
the dual of $\operatorname{Sym}^{m-1}N^*$ is the divided power
$\Gamma^{m-1}N$, including when $p$ divides a factorial.
The conormal filtration of $\Omega_B^2|_S$ has factors
$\Lambda^jN^*\otimes\Omega_S^{2-j}$, for $j=0,1,2$.
The corresponding factors in (A31u.11) are
\[
 \begin{array}{c|c|c}
 j&\text{factor}&\text{polynomial degree in }N\\ \hline
 0&\det N\otimes\Gamma^{m-1}N\otimes\Omega_S^2&m+1\\
 1&N\otimes\Gamma^{m-1}N\otimes\Omega_S^1&m\\
 2&\Gamma^{m-1}N&m-1.
 \end{array}
\tag{A31u.12}
\]
Here $\det N\otimes N^*=N$ is the canonical rank-two
identification.  The only degree-zero factor is $m=1,j=2$.
It is an actual submodule $\mathcal O_S\nu\subset M_1^2$,
coming from $\Lambda^2N^*\subset\Omega_B^2|_S$; it is not
merely a factor in an abstract associated graded object.
Every remaining factor has no global sections by (A31u.10).
The finite filtration therefore proves
$H^0(M_m^2)=H^0(\mathcal O_S\nu)=k\nu$ for every $m$.

The supported Cartier operator sends $\nu$ to the nonzero
fundamental fraction on the twist.  Thus no nonzero multiple of
$\nu$ is locally exact.  A global section of a coherent locally
exact two-form subsheaf lies in some finite $M_m^2$, because its
support is quasi-compact.  By (A31u.3) it is a multiple of $\nu$,
and hence is zero.  This proves the last assertion in part 1.

\textit{2. The ordinary and differentiated normal contractions.}
The conormal submodule of $M_1^1=\det N\otimes\Omega_B^1|_S$
is exactly the image $jN=\det N\otimes N^*$.  Its quotient is
$\det N\otimes\Omega_S^1$.  All higher pole factors for
one-forms have positive polynomial degree.  Consequently
\[
                         H^0(M_m^1/jN)=0
\tag{A31u.13}
\]
by (A31u.10), which proves the first injection in part 2.

For the differentiated assertion one must check the actual
leading-pole cancellation.  In normal coordinates, with local
lifts of $a,b$, differentiation of (A31u.2) gives
\[
 d j(a,b)=
 \left[\frac{da\wedge dy-db\wedge dx}{xy}
       -\frac{a\,dx\wedge dy}{x^2y}
       -\frac{b\,dx\wedge dy}{xy^2}\right].
\tag{A31u.14}
\]
Changing a lift by a member of $I$ does not change the original
one-form in local cohomology, and therefore does not change its
differential.  Modulo $M_1^2$, the last two fractions give
minus the identity onto the conormal two-form submodule
\[
 N=\det N\otimes N\otimes\Lambda^2N^*
                           \subset M_2^2/M_1^2.
\tag{A31u.15}
\]
The terms involving $da,db$ have pole at most two after normal
representatives are chosen, and (A31u.15) is intrinsic under a
change of normal frame.

Let $J$ be the inverse image of the submodule (A31u.15) in
$M_2^2$.  The exact sequence
\[
                       0\longrightarrow M_1^2
                         \longrightarrow J
                         \longrightarrow N\longrightarrow0
\tag{A31u.16}
\]
splits as coherent modules after Frobenius pushforward: the
section is $-d j$, and its composition with the last arrow is
the identity.  Frobenius pushforward makes $d j$ linear,
since $d(f^p\omega)=f^p d\omega$.
This is an actual splitting of (A31u.16), not an inference from
equal ranks in its associated graded.
It follows that $F_*M_2^2/dj(F_*N)$ has a filtration beginning
with $F_*M_1^2$, followed by the quotient of
$F_*(M_2^2/M_1^2)$ by the indicated $F_*N$ submodule.
All factors of that quotient have positive polynomial degree.
The same holds for the additional factors at every $m>2$.
Within $F_*M_1^2$ the sole degree-zero part remains $F_*\mathcal
O_S\nu$.  Therefore
\[
 H^0\bigl(F_*M_m^2/dj(F_*N)\bigr)=k\nu,
\tag{A31u.17}
\]
and every one of these sections lifts to the global fraction
$\nu$ in $F_*M_m^2$.  The long exact cohomology sequence now
gives the injection (A31u.4).

\textit{3. Supported Cartier and finite bounds.}
The ordinary Cartier isomorphism is natural under localization.
For the regular two-generator support ideal, local cohomology
of every locally free de Rham or Cartier-cohomology sheaf is
concentrated in degree two.  The two spectral sequences for
the support Cech complex and the de Rham complex therefore
give the supported Cartier isomorphism
\[
 \mathcal H^i(\mathscr M^\bullet)
                  \simeq\mathscr M^i_{B^{(1)}}.
\tag{A31u.18}
\]
Iterating its cycles-and-boundaries definition gives (A31u.5)
and $C^a:\mathscr Z_a^1/\mathscr B_a^1\simeq
\mathscr M^1_{B^{(a)}}$.  This is the supported version of
the higher Cartier filtration; it does not identify
$\mathscr Z_a^1/\mathscr B_1^1$ with the final target.

Cartier does not increase the finite normal pole bound.
For example a surviving normal term
$x^{-i}y^{-j}dx$ has $i\equiv1\pmod p$ and $j\equiv0\pmod p$,
and is sent to a term with denominator total
$(i+j-1)/p+1\leq i+j$.
A surviving tangent term has both normal exponents divisible
by $p$, and their total is divided by $p$.
These formulas are checked on the usual etale-coordinate
de Rham decomposition and commute with passage to local
cohomology.  Differentiation raises total normal pole by at
most one.  Thus every image used below is contained in a
finite supported two-form module, for example in $M_{M+1}^2$
on the appropriate twist.  No unbounded inverse is used.

For $a=1$, the quotient $K_{0,M}/K_{1,M}$ is the image of
$d:M_M^1\to\mathscr M^2$.
For $a>1$, the quotient $K_{a-1,M}/K_{a,M}$ is the image of
\[
                 d C^{a-1}:K_{a-1,M}
                          \longrightarrow\mathscr M^2_{B^{(a-1)}}.
\tag{A31u.19}
\]
This map is defined because the previous Cartier iterates are
closed, and it kills the bounded exact denominator in
(A31u.6).  Its kernel is exactly $K_{a,M}$ by definition.
Each image consists of locally exact supported two-forms,
so part 1, on the field twist, gives
\[
                         H^0(K_{a-1,M}/K_{a,M})=0.
\tag{A31u.20}
\]
The long exact sequences give injections
$H^1(K_{a,M})\hookrightarrow H^1(K_{a-1,M})$.
Their composite is (A31u.7).

Finally a class satisfying the hypotheses of (A31u.8) is zero
already in $H^1(K_{a,M})$.  Applying the defined sheaf map
$dC^a$ makes its image zero in supported two-form cohomology.
The injection (A31u.4), on the $a$-th field twist, implies
$\lambda o=0$, hence $o=0$.  This completes the proof.
\end{proof}

\paragraph{The actual radial numerators.}
The local identity relevant to an eventual application can be
proved without a divided-power or ramification assertion.
Let $a,b$ be displacement functions and, for $n\geq1$, put
\[
 \begin{split}
 B_n(a,b)={}&(a\,dy-b\,dx)
       \sum_{i+j=n-1}\frac{(-a)^i(-b)^j}{x^{i+1}y^{j+1}}\\
 &+(a\,db-b\,da)
       \sum_{i+j=n-2}\frac{(-a)^i(-b)^j}{x^{i+1}y^{j+1}},
 \end{split}
\tag{A31u.21}
\]
where the second sum is empty for $n=1$.
These are the numerators of the radial transgression
$\sum_{n\geq1}z^n B_n/n$ for the displacement $(za,zb)$.
Expanding
$d\log(x+za)\wedge d\log(y+zb)$, including the terms $da,db$,
shows $dB_n=n\nu_n$, with $\nu_n$ its degree-$n$ two-form
coefficient.  In characteristic $p$, $B_{pn}$ is therefore
closed.  Direct monomial Cartier calculation gives
\[
                             C(B_{pn})=B_n.
\tag{A31u.22}
\]
For the $dy$ terms the surviving indices satisfy
$i+1\equiv0\pmod p$ and $j\equiv0\pmod p$;
for the $dx$ terms they satisfy $i\equiv0\pmod p$ and
$j+1\equiv0\pmod p$.
For the $db$ and $da$ terms both normal denominator exponents
are divisible by $p$, and the respective coefficient exponent
of $b$ or $a$ is $-1$ modulo $p$.
Dividing those exponents by Cartier gives exactly the two
sums for $B_n$.  This computation can be made in the universal
polynomial ring in $a,b$ with $x,y$ inverted, using the standard
de Rham monomial decomposition, and then pulled back; it does
not assume that the displacement coefficients are constant.

In particular
\[
                  B_{p^a}\in\mathscr Z_a^1,
                       \qquad C^a(B_{p^a})=j(a,b).
\tag{A31u.23}
\]
If a leading edge is an actual cocycle whose local expression
is $\lambda B_{p^a}+\mu B_{p^{a+1}}$, its local expression
after $dC^a$ is $\lambda^{1/p^a}d j(a,b)$: the higher term
becomes the closed form $B_p$ and disappears under $d$.
This observation is useful also when two adjacent leading
weights tie.  To invoke (A31u.8), actual global descent and
the effect of preceding closed-layer or ambient corrections
must still be checked.

\paragraph{Why the pure-power replacement fails at the next level.}
For clarity, consider instead the shorter expression
\[
 \xi_q(a,b)=a^q\frac{dy}{x^q y}
                    -b^q\frac{dx}{x y^q},\qquad q=p^a.
\tag{A31u.24}
\]
Take $a=0,b=1$ and the constant shear $x'=x+y,y'=y$.
The transformed vector is $(1,1)$, and in local cohomology the
change is
\[
 D_q=\sum_{j=1}^{q-1}(-1)^j
       \left(-\frac{dx}{x^{j+1}y^{q-j}}
                    +\frac{dy}{x^j y^{q-j+1}}\right).
\tag{A31u.25}
\]
The terms with $p\nmid j$ are
$d((-1)^j/(j x^j y^{q-j}))$.
For $q=p$ this makes $D_p$ exact, but $D_p$ is nonzero as an
ordinary supported one-form.  For $q=p^a$, $a\geq2$, the
remaining terms with $p\mid j$ have Cartier image
\[
                              C([D_{p^a}])=D_{p^{a-1}}\ne0.
\tag{A31u.26}
\]
Thus $D_{p^a}$ is not exact.  Cartier here has target the
ordinary one-form sheaf; an exact but nonzero form in that
target is still nonzero.  The naive pure-power expressions
therefore do not define a normal-line map into $Q_M$ at
level $a\geq2$.  Formula (A31u.21) and the higher filtration
retain precisely the additional terms that this replacement
would lose.  This is a method boundary, not a geometric
counterexample to a lifting assertion.

For a first displacement of order one, the one-contraction
valuation is $n-e v_p(n)$.  Its minimum occurs at powers
$n=p^a$, with values $p^a-ae$.  The difference between
successive values is $p^a(p-1)-e$.
Hence the first tie of the $p$ and $p^2$ terms is exactly
$e=p(p-1)$, including $p=3,e=6$.
Beyond that point the problem is not merely the removal of
more closed functions: the leading one-form itself enters
the higher Cartier filtration.

\paragraph{Primary source scope.}
The higher cycles-and-boundaries definitions and the exact
Cartier iteration are Illusie, \emph{Complexe de de Rham--Witt},
Asterisque 63 (1979), pp.~83--112, section 2.3, printed p.~94,
formulas (2.3.1)--(2.3.3):
\url{https://www.numdam.org/item/AST_1979__63__83_0.pdf}.
The actual printed page and these formulas were read.
The supported passage was proved above using the same local
Cartier theorem and support spectral sequences as A31h.
All normal-pole, global-section, leading-pole splitting, and
radial numerator calculations in this section are given here;
no theorem about a ramified global comparison is being
attributed to Illusie.


\paragraph{Programme linkage.}
This is an auxiliary detection statement for route (a) of
JAC1964 QH9-axis-lift, concerning the primitive axis term of H9-T7
in cote 44, Theorem 5, scan 004, with the primitive warning on
scans 008--009.  It is a special-fibre de Rham and Cartier calculation,
not a proof of the unrestricted lifting or sign assertion.
It makes no new claim for the product-form condition I or
semisimplicity on cote 16, scan 049.  No 032 or Milne2026 input is used.

% END work/scan_w/higher_cartier_detection.tex

% BEGIN work/axis/axis_formal_cartier_saturation.tex
% C44_POS. Author, root and independent entire-proof readings PASS; scopes recorded in CHECK_EVENTS.
\section{Formal Cartier saturation and the remaining uniformity issue}
\label{sec:c44-a31v}

\paragraph{Scope.}
This diagnostic separates a finite-dimensional colimit from the
individual finite-pole groups.  It supplies a sufficient finite
test for colimit edges on one fixed axis, but does not prove that
this test holds there.  In particular it does not replace the
degree threshold depending on \(M\) in A31n by a uniform one.

\begin{proposition}[C44-A31v: a finite colimit under ample conormal]
Let \(B\) be a smooth projective fourfold over an algebraically
closed field \(k\) of characteristic \(p>0\), and let
\(S\subset B\) be a smooth surface with ample conormal bundle
\(N_{S/B}^*\).
Use the supported sheaves
\[
 M_a^b=\mathcal Ext_B^2(\mathcal O_B/I_S^a,\Omega_B^b),
 \qquad
 C_M=[M_{M-1}^0\xrightarrow d M_M^1],
 \qquad Q_M=\operatorname{coker}d .
\tag{A31v.1}
\]
Take filtered colimits over finite \(M\), and put
\[
 C_\infty=[\mathcal H_S^2(\mathcal O_B)
                 \xrightarrow d\mathcal H_S^2(\Omega_B^1)],
 \quad Z_\infty=\ker d,\quad Q_\infty=\operatorname{coker}d.
\tag{A31v.2}
\]
Then
\[
                  \mathcal E_B:=H^1(Q_\infty)
                         =\operatorname*{colim}_M H^1(Q_M)
\tag{A31v.3}
\]
is finite-dimensional.
This assertion does not say that every \(H^1(Q_M)\) maps
injectively into the colimit, or that the groups \(H^1(Q_M)\)
have stabilized.
\end{proposition}

\begin{proof}
\leavevmode\par
Write \(F:B\to B^{(1)}\) for relative Frobenius and put
\[
 B_B^1=\operatorname{im}(d:F_*\mathcal O_B\to F_*\Omega_B^1),
 \qquad E=F_*\Omega_B^1/B_B^1 .
\tag{A31v.4}
\]
Both \(B_B^1\) and \(E\) are vector bundles on \(B^{(1)}\).
In etale coordinates this follows from the tensor product
of the one-variable Cartier complexes with basis
\(1,z,\ldots,z^{p-1}\).
The two exact sequences of locally free sheaves are
\[
 0\to\mathcal O_{B^{(1)}}\to F_*\mathcal O_B\to B_B^1\to0,
 \qquad
 0\to B_B^1\to F_*\Omega_B^1\to E\to0.
\tag{A31v.5}
\]

Local cohomology with support in a smooth codimension-two
immersion is concentrated in degree two on each locally free
sheaf.  Apply it to (A31v.5).
The first sequence identifies the image of the supported
differential with \(\mathcal H^2_{S^{(1)}}(B_B^1)\);
the second identifies its cokernel.  Thus, with Frobenius
pushforward understood,
\[
 Z_\infty=\mathcal H^2_{S^{(1)}}(\mathcal O_{B^{(1)}}),
 \qquad
 Q_\infty=\mathcal H^2_{S^{(1)}}(E).
\tag{A31v.6}
\]
In particular
\[
                  H^1(Q_\infty)=H^3_{S^{(1)}}(B^{(1)},E).
\tag{A31v.7}
\]

Here is a direct finite-duality proof that this group is finite.
Let \(S_a\) be the \(a\)-th ordinary ideal-power thickening of
\(S^{(1)}\) in \(B^{(1)}\).
The definition of sections with support, applied to an
injective resolution, gives
\[
 H^3_{S^{(1)}}(B^{(1)},E)
     =\operatorname*{colim}_a
          \operatorname{Ext}^3_{B^{(1)}}(\mathcal O_{S_a},E).
\tag{A31v.8}
\]
Filtered colimits are exact, so this uses no formal-algebraization
claim.
Finite Serre duality on the smooth projective fourfold gives
compatible perfect pairings
\[
 \operatorname{Ext}^3_{B^{(1)}}(\mathcal O_{S_a},E)^*
        \simeq
 H^1\!\left(S_a,(E^*\otimes\omega_{B^{(1)}})|S_a\right).
\tag{A31v.9}
\]
For example one obtains this by applying Serre duality to
\(\mathcal O_{S_a}\otimes E^*\otimes\omega_{B^{(1)}}\);
all sheaves at this finite stage are coherent.

Set \(G=(E^*\otimes\omega_{B^{(1)}})|S^{(1)}\) and
\(P=N_{S^{(1)}/B^{(1)}}^*\).
The graded layers of the right-hand side of (A31v.9) are
\(G\otimes\operatorname{Sym}^aP\).
Since \(P\) is ample, the tautological line bundle on
\(\mathbf P(P)\) is ample.
The projective-bundle formula and Serre vanishing give
\[
 H^i(S^{(1)},G\otimes\operatorname{Sym}^aP)=0
                      \qquad(i>0,\ a\gg0).
\tag{A31v.10}
\]
No Kodaira vanishing or characteristic-zero lifting is used.
The exact thickening sequence therefore makes
\[
 H^1(S_{a+1},(E^*\omega)|S_{a+1})
       \longrightarrow H^1(S_a,(E^*\omega)|S_a)
\tag{A31v.11}
\]
an isomorphism for all sufficiently large \(a\): its adjacent
terms are precisely the \(H^1\) and \(H^2\) in (A31v.10).
Dualizing shows that the direct system in (A31v.8) is eventually
constant and finite-dimensional.

Finally, after Frobenius pushforward all finite-pole sheaves
are quasi-coherent on finite thickenings of \(S^{(1)}\).
Their union is the corresponding local-cohomology sheaf.
Filtered colimits commute with cokernels, and cohomology on
the quasi-compact separated scheme \(B^{(1)}\) commutes with
these filtered colimits.
This proves both the identity and the finiteness in (A31v.3).
It proves stabilization of the particular Ext system in
(A31v.8), not of the differently filtered groups \(H^1(Q_M)\).
\end{proof}

\begin{lemma}[The finite-pole quotients need not inject]
Even locally in smooth coordinates the maps
\(Q_M\to Q_{M+1}\) need not be injective.
\end{lemma}
\begin{proof}
Take etale tangent and normal coordinates \(t,z,x,y\), with
\(S=(x,y)\), and put \(M=2p-1\).
The function
\[
                  \theta=\frac{t}{x^py^p}
\]
has total pole \(2p\), one beyond the old function bound.
Its differential is
\[
                         d\theta=\frac{dt}{x^py^p}.
\tag{A31v.12}
\]
This lies in the old one-form bound \(M_M^1\).
It is not the derivative of a function of total pole at most
\(2p-1\): a tangent derivative preserves normal pole, and a
normal derivative has a \(dx\) or \(dy\) component.
In the negative-normal-monomial basis there can therefore be
no \(dt\,x^{-p}y^{-p}\) term from such a function.
Its class in \(Q_{2p-1}\) is nonzero.
The same class is zero in \(Q_{2p}\), where \(\theta\) is allowed.
This is a sheaf-level counterexample; it does not assert that
every such local class gives a nonzero global \(H^1\) class.
\end{proof}

\begin{proposition}[A sufficient finite test for colimit edges]
Suppose now \(S\subset B\subset X\) is an axis as in A31n,
and its conormal is ample.
Let \(P_B\) denote its finite-dimensional primitive \(F^3\)
test space.
Let
\[
 \Pi_M:\mathbb H^2(C_M)\longrightarrow P_B^*
\tag{A31v.13}
\]
be the period map, and define
\[
 W_M=\operatorname{im}\left(
        H^2(Z_M)\to\mathbb H^2(C_M)\xrightarrow{\Pi_M}P_B^*
                         \right).
\tag{A31v.14}
\]
There is a finite integer \(M_0=M_0(B,S)\) such that
\[
 W_M=W_\infty\quad(M\geq M_0),\qquad
 H^1(Q_{M_0})\twoheadrightarrow \mathcal E_B.
\tag{A31v.15}
\]
If the annihilator conclusion of A31n holds for this fixed
axis at this finite bound \(M_0\), then the induced colimit
map
\[
                    \overline\Pi_\infty:
                  \mathcal E_B\longrightarrow P_B^*/W_\infty
\tag{A31v.16}
\]
is injective.

If in addition the A31p injection of the ordinary and first
Cartier normal summands holds for every finite \(M\geq p\), these
summands inject into \(\mathcal E_B\).
Under the stated finite-testing hypothesis, any finite
supported class whose edge is a nonzero vector in one of
these summands has a nonzero primitive period.
\end{proposition}

\begin{proof}
The \(W_M\) form an increasing family of subspaces of the
fixed finite-dimensional vector space \(P_B^*\); hence they
eventually stabilize.
By (A31v.3), choose finitely many representatives of a basis
of \(\mathcal E_B\) in finite groups \(H^1(Q_M)\).
A single larger \(M_0\geq p\) contains all of them and is past the
stabilization index of \(W_M\).
This proves (A31v.15) without assuming injectivity of any
finite-pole transition map.

All \(Z_M\), and their filtered union \(Z_\infty\), have
cohomology in degrees at most two: the finite sheaves are
supported on a surface and cohomology commutes with the
filtered union.
The truncation triangle for \(C_\infty\) thus gives the
right-exact sequence
\[
 H^2(Z_\infty)\longrightarrow\mathbb H^2(C_\infty)
                       \longrightarrow\mathcal E_B\longrightarrow0.
\tag{A31v.17}
\]
Period maps commute with the filtered union.
Quotienting their target by \(W_\infty\) makes (A31v.16)
well-defined.

Let \(v\in\mathcal E_B\) have zero image in (A31v.16).
Choose \(v_0\in H^1(Q_{M_0})\) mapping to \(v\).
Since \(H^3(Z_{M_0})=0\), choose a lift
\(\gamma_0\in\mathbb H^2(C_{M_0})\).
Its period lies in \(W_\infty=W_{M_0}\).
Subtract the image of a class of \(H^2(Z_{M_0})\) with
that period.
The resulting class has zero primitive periods at the same
finite bound \(M_0\).
The A31n annihilator assertion at that bound says that its
edge in \(H^1(Q_{M_0})\) is zero.
Subtracting a closed class did not change the edge, so
\(v_0=0\), and hence \(v=0\).

For the last assertion, A31p injects the original normal
summands into every \(H^1(Q_M)\), compatibly with all transition
maps.  A nonzero vector cannot become zero in the filtered
colimit, since that would mean that it was already zero at
some finite stage.
Injectivity of (A31v.16) then makes its period nonzero modulo
\(W_\infty\), and in particular nonzero in \(P_B^*\).
This argument concerns the original summands, not every
possible nonzero edge at an arbitrary finite pole bound.
\end{proof}

\begin{lemma}[Why the original ambient Serre argument is not uniform]
The positive-twist Serre vanishings in A31n cannot in general
hold for every ideal power with a fixed axis degree, even when
the axis is chosen in the A31p ample-conormal open.
\end{lemma}
\begin{proof}
Take \(X=U\times V\), where \(U\) is an abelian surface and
\(V\) an abelian fourfold, and \(S=U\times\{0\}\).
Choose very ample \(H_U,H_V\), put \(H=H_U\boxtimes H_V\), and let
\(J=I_{S/X}\).
For fixed positive integers \(t,u\), put \(A=H^{tu}\).
The \(s\)-th ambient thickening is
\(U\times\operatorname{Spec}(\mathcal O_{V,0}/\mathfrak m^s)\).
The \(V\)-line bundle is trivial on this local Artin scheme,
whose length is
\(\binom{s+3}{4}\).
Consequently
\[
 h^0(X,A\otimes\mathcal O_X/J^s)
       =h^0(U,H_U^{tu})\binom{s+3}{4}.
\tag{A31v.18}
\]
The exact sequence of \(J^sA\) implies
\[
 h^1(X,J^sA)\geq
       h^0(U,H_U^{tu})\binom{s+3}{4}-h^0(X,A).
\tag{A31v.19}
\]
The right side tends to infinity.
This estimate does not even require vanishing of \(H^1(X,A)\).
For \(p>2\) and sufficiently large fixed \(t\), A31p supplies smooth
genuine axes containing this \(S\) with ample conormal, but
the ambient groups in (A31v.19) are independent of the
choice of axis.
Thus that open condition cannot make these particular
Serre vanishings uniform in \(s\).
This is a limitation of the proof input, not a counterexample
to uniform primitive period detection.
\end{proof}

\paragraph{Sources and exact proof-read scope.}
The filtered-colimit cohomology fact is Stacks,
Lemma 30.6.1, tag 07TB, whose complete printed proof was read:
\[
              \text{\url{https://stacks.math.columbia.edu/tag/07TB}}.
\]
Finite Serre duality used in (A31v.9) is Stacks,
Lemma 48.27.1(5),(7), Remark 48.27.2, and Lemma 48.27.5,
whose full printed statements and proofs in the section were read:
\[
              \text{\url{https://stacks.math.columbia.edu/tag/0FVU}}.
\]
This records a read of those printed proofs and their cited
dependency structure, not a fresh reconstruction of all of
Grothendieck duality.
For (A31v.10), the elementary projective-bundle proof is also
Lazarsfeld, Theorem 6.1.10, implication (i) to (ii), author
Chapter 6 PDF, printed pp.11--12; the complete theorem proof
was opened and read.
Its arbitrary-characteristic scope is Remark 6.1.17, p.15:
\[
 \text{\url{https://www.math.stonybrook.edu/~roblaz/Reprints/Laz.PAG.Chapt6.pdf}}.
\]
The supported Cartier and finite ideal-power Ext calculations
are the previously checked A31i/n/p inputs, and were spelled
out again in the needed form above.

\paragraph{Programme boundary.}
The integer \(M_0\) in (A31v.15) depends on the fixed axis.
A31n proves its hypothesis when the axis degree is sufficiently
large for a \emph{previously fixed} bound.
The argument here does not compare that degree threshold with
the unknown \(M_0(B,S)\), and therefore does not exhibit an axis
satisfying the sufficient finite test.
It also does not assert stabilization or injectivity of all
finite groups \(H^1(Q_M)\).

In QH9-axis this identifies a precise finite saturation issue
inside route (a), not an additional ramification range or a
universal H9-T7 sign conclusion.
There is no use of either unrefereed 032/Milne2026 preprint,
and no countable intersection of nonempty opens is assumed to
have an \(\overline{\mathbf F}_p\)-point.
This is a scoped diagnostic within the existing lifting work,
with no separate progress-clock claim.


\paragraph{Grothendieck and predecessor linkage.}
The remaining test belongs to JAC1964 QH9-axis-lift and the
primitive axis term H9-T7 of cote 44, Theorem 5, scan 004,
with the primitive qualification on scans 008--009.  The argument
uses coherent local cohomology, Cartier theory and finite duality.
It supplies neither the universal sign nor a new product-form
condition I for cote 16, scan 049.

% END work/axis/axis_formal_cartier_saturation.tex

% BEGIN work/arithmetic/fermat_recent_source_audit.tex
% C44_POS. Author, root and independent entire-proof readings PASS; scopes recorded in CHECK_EVENTS.
% No arithmetic-gate movement is claimed. Author/peer status is in dated registers.
\section{Further exact boundaries for the Fermat arithmetic gate}
\label{sec:c44-f07-sa3}

\paragraph{C44-F07-SA3: status and scope.}
This section tests additional primary results against the actual
groups of C44-F07, F17--F18 and F25.  It does not prove scalar
Abel--Jacobi injectivity, finite generation of the discrete
\'etale-motivic group, or a uniform bound on the actual descent
cokernel \(D=\overline A^{G_{\mathbf Q}}/A\).
The earlier source exclusions in F07-SA1 and F07-SA2 are not
reproved.  No unrefereed 032/Milne2026 input is used.

\subsection{The special CM tensor does not acquire faithfulness from a Kummer surface}

Use the actual Chow motives and pairings of C44-F18 over
\(k=\overline{\mathbf Q}\), with coefficient field
\(K=\mathbf Q(i)\) when indicated.  Thus
\[
 N_2\otimes K=U^2\oplus\overline U^2,\qquad
 U\otimes\overline U\simeq\mathbb L_K,\qquad
 N_4\otimes K=U^4\oplus\overline U^4.
 \tag{SA3.1}
\]
In particular \(U\) is an invertible Chow motive.  Let \(I\) on
\(N_2\) be induced by \([i]\) on the \emph{first} factor of
\(E^2\); it has eigenvalues \(i,-i\) on the two displayed lines.
Using the diagonal action on both factors would instead give
\(-1\) on both lines and would not give this decomposition.

Here is a direct test of the proposed endomorphism argument.
Put
\[
 \mathcal R_{\rm CM}=
 \ker\bigl[\operatorname{End}_{\rm CHM}(N_2)
               \longrightarrow\operatorname{End}_{\rm hom}(N_2)\bigr].
\]
There are exact vector-space identifications and an ideal identity
\[
 \operatorname{End}_{\rm CHM}(N_2)
       =K\oplus\mathcal R_{\rm CM},\qquad
 \mathcal R_{\rm CM}\simeq\operatorname{CH}^2(N_4),\qquad
 \mathcal R_{\rm CM}^{\,2}=0.
 \tag{SA3.2}
\]
These statements concern this actual motive; they do not exhibit
a nonzero element of the ideal.

Indeed, after extending coefficients to \(K\), each diagonal
endomorphism ring
\(\operatorname{End}(U^2)\) and
\(\operatorname{End}(\overline U^2)\) is \(K\), because tensoring
by an invertible motive identifies it with
\(\operatorname{End}(\mathbf1_K)\).
Both off-diagonal Hom groups have zero Betti realization:
the two Hodge lines have different types \((2,0)\) and \((0,2)\).
Any composite of two off-diagonal maps is a diagonal scalar with
zero realization and is therefore zero in Chow.
The diagonal pair descends to the actual algebra \(K=\mathbf Q[I]\).
It realizes all Hodge endomorphisms of the rational rank-two
structure, since the centralizer of its two complex Hodge lines
is \(\mathbf Q[I]\).
Thus the off-diagonal part is exactly the homological kernel.
Finally the pairing gives
\[
 \operatorname{End}(N_2)
  =\operatorname{CH}^2(N_2\otimes N_2),\qquad
 N_2\otimes N_2\simeq N_4\oplus\mathbb L^2\oplus\mathbb L^2,
 \tag{SA3.3}
\]
by (SA3.1).  This identifies the off-diagonal part with the
group in (SA3.2).  Equivalently \(Ir=-rI\) for
\(r\in\mathcal R_{\rm CM}\).  The CM centralizer itself is
exactly \(K\), but arbitrary Chow correspondences need not
centralize the CM action.

The Kummer construction changes neither this group nor its
vanishing question.  Kahn--Murre--Pedrini,
\emph{On the transcendental part of the motive of a surface},
the July 2005 author version
\url{https://webusers.imj-prg.fr/~bruno.kahn/preprints/surfaces4.pdf},
Remark 7.6.13, pp.~41--42, gives
\(t_2(E^2)\simeq t_2(\operatorname{Km}(E^2))\).
Here one may choose \(t_2(E^2)=N_2\): the two factor divisors,
the diagonal, and the graph of \([i]\) span its rank-four
N\'eron--Severi space.  Their intersection projector and the
canonical degree-two projector are distinguished, so their
complement is the projector of \(N_2\) already used in F18.
Theorem 7.6.12, pp.~40--41, separates
finite-dimensionality and nilpotence, assertions (1)--(3),
from faithfulness of the endomorphism realization, assertion (6).
The complete displayed proof of that theorem and the following
remark were read.  In our test the radical is already
square-zero by (SA3.2).  Reapplying a nilpotence theorem therefore
does not prove that it is zero.

\subsection{The mixed elliptic category retains the same vanishing hypothesis}

Jin Cao, \emph{Motives for an elliptic curve},
arXiv:1708.08092v2, 5 May 2018,
\url{https://arxiv.org/pdf/1708.08092v2},
uses a contravariant embedding with
\(\Phi(h^1(E))=M_1(E)[1]\) and
\(\Phi(\mathbb L)=K(1)[2]\); see the proof of Lemma 2.1,
pp.~4--5.  Write \(M_1(E)_K=M\oplus\overline M\).
The exact translation of the surviving F18 group is
\[
 \operatorname{CH}^2(N_4)\otimes K
  \simeq
 \operatorname{Hom}_{\rm DM}(M^4,K(2))
 \oplus
 \operatorname{Hom}_{\rm DM}(\overline M^4,K(2)).
 \tag{SA3.4}
\]
This follows by applying \(\Phi\) to
\(\operatorname{Hom}_{\rm CHM}(\mathbb L^2,N_4)\);
the shifts \([4]\) on source and target cancel.

The vanishing in (SA3.4) is exactly the CM zeroth-vanishing
condition of Definition 2.5 and Conjecture 2.6, p.~5,
at \(i=2,j=0\).
Example 2.8, p.~6, proves the \(i=1\) case by divisor
endomorphisms, which is the already eliminated F18 component.
The unconditional triangulated equivalence of Theorem 8.15,
p.~26, does not assert the required vanishing.
The CM Tannakian conclusion of Corollary 8.19, p.~27,
assumes the vanishing properties and the additional generalized
Beilinson--Soul\'e conjecture.
The convention passages, Definition 2.5--Example 2.8 including
their displayed proofs, and Theorem 8.15 through Corollary 8.19
were read; no entire-proof reading of the long dg-equivalence
construction is claimed.

\subsection{Uniform torsion results control a different descent group}

Charles--Pirutka, \emph{Finitude uniforme pour les cycles de
codimension 2 sur les corps de nombres},
arXiv:2207.05805v1, 12 July 2022,
\url{https://arxiv.org/pdf/2207.05805v1},
Theorem 2.10, p.~9, bounds the exponent of
\[
 \ker\bigl[\operatorname{CH}^2(X)
          \longrightarrow\operatorname{CH}^2(\overline X)^{G_k}\bigr]
 \tag{SA3.5}
\]
under the standing hypothesis \(H^2(X,\mathcal O_X)=0\).
Its complete proof, together with Lemma 2.11 and its complete
proof on p.~10, was read.
The argument uses the degree-three \'etale restriction kernel,
the N\'eron--Severi comparison, and a Picard-curve isogeny.
It does not bound the \emph{cokernel} in F25.
For the actual projected Fermat group the restriction kernel
is already zero by F25, whereas \(J\) itself does not satisfy
the displayed coherent-cohomology hypothesis.
Projecting away \(H^2\) does not change a theorem about (SA3.5)
into a theorem about its cokernel.
The family bounds in Theorems 1.2 and 1.4 also concern primary
Chow torsion or its integral-model image, not \(D\).

For the same reason, the general \(K_2\) descent theorem cannot
be applied with the function field silently replaced by its
constant number field.
Kahn, \emph{Descente galoisienne et \(K_2\) des corps de nombres},
\emph{K-Theory} 7 (1993), 55--100,
\url{https://webusers.imj-prg.fr/~bruno.kahn/preprints/descentegaloisienne.pdf},
Theorem 2.1, pp.~60--62, explicitly leaves
\(\ker[H^3(F,\mathbf Q/\mathbf Z(2))
\to H^3(E,\mathbf Q/\mathbf Z(2))]\)
in the description of \(H^2(G,K_2(E))\).
Theorem 2.2 and Example 2.1, pp.~62--64, do not supply
a general finiteness theorem for that degree-two group.
The relative Theorem 3.1 and Corollary 1, pp.~67--69,
give the degree-one vanishing.
The displayed statements and the relative theorem's complete
proof were read.  F25 needs the degree-two tame-symbol kernel
over the transcendence-degree-three function field.

\subsection{Self-products of CM elliptic curves: exact image and field scope}

Jonathan Dee, \emph{Selmer Groups of Hecke Characters and Chow
Groups of Self Products of CM Elliptic Curves}, 4 January 1999,
\url{https://arxiv.org/pdf/math/9901155},
Theorem 5.2.2, pp.~30--31, controls the cycle-map image
intersected with the second Hochschild--Serre filtration.
Its hypotheses include \(p>2i+1\).
For \(i=2,p=5\), the strict inequality fails.
Remark 5.2.1 explicitly retains the strict bound because of the
global input, despite an improved local endpoint.
More fundamentally the relevant geometric representation is
\(H^{2i-2}\), which is \(H^2\) here; its projected piece
vanishes already.
The complete displayed argument in Section 5.2, including the
theorem's proof, was read, together with the Section 5.1
decomposition and stated rational image theorem.
No conclusion concerns the scalar cycle-map kernel, F25's
descent cokernel, or actual finite generation in F07.

The recent local-to-global result has a different group as well.
Jonathan Wills, \emph{On the Local-to-Global Principle for
Zero-Cycles on Self Products of Elliptic Curves with CM},
arXiv:2509.13641v1,
\url{https://arxiv.org/html/2509.13641v1},
Theorem 5.2 establishes an exact local-to-global complex of
inverse limits of Albanese kernels modulo \(p^n\), for specified
extensions and the hypothesis
\(p\mid\#\overline E(\mathbf F_v)\).
Theorem 1.2, Proposition 5.1 with its full displayed proof,
and Theorem 5.2 with its proof were read.
This is a result for zero-cycles on the surface \(E^2\).
It does not identify the discrete group of codimension-two
cycles on the odd motive in \(E^3\), and it does not bound
\(\overline A^{G}/A\).

Finally Gazaki--Love, \emph{Torsion phenomena for zero-cycles
on a product of curves over a number field},
\emph{Research in Number Theory} 10 (2024), article 35,
\url{https://doi.org/10.1007/s40993-024-00519-4},
Proposition 4.8 concerns the \emph{componentwise} zero-cycle
subgroup when \(E(\mathbf Q)\) has rank one and \(E(K)\)
has rank two.
The proof reduces to four symbols formed from one point and
its CM multiple, using Proposition 4.5.
Both displayed proofs were read.
It supplies no statement uniform over all finite fields of
definition, and no generation statement for
\(\operatorname{CH}^2(N_4)\) or for the odd Fermat group.
Thus none of these checked results closes the actual arithmetic
or residual Chow gate.


\paragraph{Programme linkage.}
The unchanged predecessor gates are JAC1964 QFermat-AJ,
Q35-full, AR-16a and P18.  They concern the arithmetic counterpart
of the primitive height-sign problem in cote 44, Theorems 4--5,
scans 002 and 004, with the primitive warning on scans 008--009.
The argument is formulated in Chow motives and motivic descent;
it supplies no new assertion for conditions A, C or I on
cote 16, scans 049--050.

% END work/arithmetic/fermat_recent_source_audit.tex

% BEGIN work/arithmetic/ramified_rigid_source_audit.tex
% C44_POS, /root/fermat_aj. Primary-source audit and local diagnostic.
% No 032/Milne positivity input. No higher-ramification nonliftability theorem.

\section{Source audit: rigid Gysin transport beyond the first ramification range}
\label{C44-A31-rigid-source-audit}

\noindent\textbf{Audit identifier: C44-A31-SA1.}
The canonical rational comparison exists at arbitrary ramification.  The
issue remaining in the proposed extension of C44-A31f is the calculation
of its \emph{actual} supported descent representative and the detection
of its new closed terms by primitive periods.  This audit does not prove
nonliftability at $p=3,e\geq3$, does not change the Fermat gate C44-F07,
and is not a separate progress-clock advance.

\paragraph{Primary locators and the scope of their proofs.}
Petrequin, \emph{Classes de Chern et classes de cycles en cohomologie
rigide}, BSMF~131 (2003), 59--121, has the following relevant locators:
Lemma~3.4, pp.~72--74, convergent logarithms over a general coefficient
DVR; Proposition~3.9, p.~77, the supported Cech cocycle;
Proposition~3.10 and Lemma~3.12, pp.~78--80, change-of-choice and
embedding independence; Proposition~3.15, p.~80, scalar extension.
The complete displayed proofs were read.  Remark~3.16, p.~80, and
Section~5.1, p.~89, then fix a Cohen ring.  Theorem~5.23 and
Corollaries~5.29--5.30, pp.~97--101, compare Chern classes; their
displayed proofs were read.  Theorem~5.22 itself cites an unpublished
1990 letter for its underlying comparison.  Proposition~2.8, p.~69,
and the complete proof of Proposition~6.6, pp.~109--111, identify the
supported divisor class with Gysin of $1$.  Lemma~7.1, p.~113, treats
regular-immersion composition and transverse products; its displayed
argument was read, with its cited foundational results retained.
The checked published PDF is
\url{https://www.numdam.org/article/BSMF_2003__131_1_59_0.pdf}.

For the underlying comparison one can instead use Berthelot,
\emph{Finitude et purete cohomologique en cohomologie rigide},
Inventiones~128 (1997), 329--377.  In the checked author PDF,
Propositions~1.8--1.9, pp.~12--14, prove scalar extension and
rigid--crystalline comparison; both complete proofs were read.
Theorem~3.8, pp.~22--23, allows any complete mixed-characteristic DVR.
Sections~5.2--5.4, pp.~34--36, construct Gysin from algebraic de Rham
Gysin by analytification and exact overconvergent-support functors;
the entire construction and Lemma~5.3 proof were read.
Remark~3.9(iv), p.~23, explicitly leaves general independence from
the chosen lifted pair unchecked.  Thus that source alone must not
be cited for such independence.  Petrequin's comparison with the
canonical fundamental class supplies the degree-zero class needed here.
The author-version pagination is used throughout this paragraph:
\url{https://wstein.org/people/berthelo/publis/Finitude.pdf}.

\paragraph{The ordinary cycle-class comparison needs no lift of the cycle.}
Let $k$ be perfect, $W=W(k)$, $K_0=\operatorname{Frac}(W)$,
and let $R$ be a complete mixed-characteristic DVR with residue
field $k$ and fraction field $K$.  Let $X/k$ and $S/k$ be smooth
projective, and let $i:S\hookrightarrow X$ have codimension $r$.
Write $c_{\mathrm{rig}}(S)$ and $c_{\mathrm{cris}}(S)$ for their
ordinary degree-$2r$ fundamental classes.  The comparison in
Berthelot's Proposition~1.9 is induced by a differential-compatible
map of algebras, is functorial, and is multiplicative.  The Chern-class
comparison just cited fixes its top-degree trace normalization:
on every geometrically connected component, a top power of an
ample class has trace its nonzero rational degree in both theories.
Since top cohomology is one-dimensional, the two trace maps agree.
For every complementary-degree class $a$, the defining identity
\[
 \operatorname{Tr}_X\bigl(c(S)\cup a\bigr)
   =\operatorname{Tr}_S(i^*a)
 \tag{A31.SA1}
\]
is therefore preserved.  Perfect Poincare duality implies that
rigid--crystalline comparison carries $c_{\mathrm{rig}}(S)$ to
$c_{\mathrm{cris}}(S)$.  This argument is about the actual ordinary
class, not an arbitrary choice of supported chain map.

If $\mathcal X/R$ is a proper smooth lift of $X$, choose the formal
completion $\widehat{\mathcal X}$ as the frame in the rigid
definition.  Its tube of $X$ is the whole analytic generic fibre.
Properness removes the compactification boundary; proper comparison
then gives
\[
 H^*_{\mathrm{cris}}(X/W)\otimes_W K
 \simeq H^*_{\mathrm{rig}}(X/K_0)\otimes_{K_0}K
 \simeq H^*_{\mathrm{rig}}(X/K)
 \simeq H^*_{\mathrm{dR}}(\mathcal X_K/K).
 \tag{A31.SA2}
\]
No smooth lifting of $S$ is assumed.  In particular, the image of
$c_{\mathrm{cris}}(S)$ in the last term is a well-defined canonical
rational class even when $S$ has no compatible integral support.

For the identification of (A31.SA2) with the usual comparison, the
published source is Berthelot, \emph{Geometrie rigide et cohomologie
des varietes algebriques de caracteristique $p$}, MSMF~23 (1986),
pp.~11--14, Sections~1.2--1.5, followed by Section~2.4, p.~19.
Section~1.2 starts from the arbitrary-ramification rational comparison
and identifies formal and analytic de Rham cohomology by Cech;
Sections~1.3--1.5 extend this analytic interpretation through tubes
and diagonal embeddings.  Proper rigid cohomology is the naive
theory in Section~2.4.  These passages were read.  This article
is an exposition and defers full proofs.  The complete proof of the
arbitrary-ramification comparison was also visually read in
Berthelot--Ogus, \emph{$F$-isocrystals and de Rham cohomology I},
Inventiones~72 (1983), Theorem~2.4, pp.~169--170, with
Corollary~2.5 and Remark~2.6, p.~171.  It uses the PD ideal $(p)$,
crystalline base change and Frobenius comparison of deformations;
it does not put divided powers on $(\pi)$ at arbitrary ramification.
The 1986 primary PDF is
\url{https://www.numdam.org/article/MSMF_1986_2_23__R3_0.pdf}.
The original 1983 scan already present in the workspace is
\texttt{work/axis/primary/BerthelotOgus1983.pdf}; its source is
\url{https://gdz.sub.uni-goettingen.de/download/pdf/PPN356556735_0072/LOG_0013.pdf}.
No integral-lattice preservation is inferred from (A31.SA2).

\paragraph{What local logarithms do prove.}
On a strict neighborhood on which $|\pi a/x|<1$, the series
\[
 \log(1+\pi a/x)
   =\sum_{n\geq1}(-1)^{n+1}\frac{\pi^n}{n}(a/x)^n
 \tag{A31.SA3}
\]
converges, independently of the ramification index $e=v_\pi(p)$.
Indeed $v_\pi(\pi^n/n)=n-e\,v_p(n)$ tends to infinity, and on
each sufficiently small affinoid neighborhood the usual strict
norm bound gives uniform convergence.  This is a convergence
statement, not a statement that all coefficients are integral.
Petrequin's construction shows that these logarithms occur in
canonical supported divisor cocycles and in their change-of-lift
coboundaries.  Transverse products give the corresponding local
two-Cartier representative.  A global codimension-two cycle need
not be the intersection of two global Cartier divisors, so these
local facts by themselves do not give its complete descent data.

Here is the two-Cartier computation with its composition term kept.
Work where the indicated ratios admit the convergent logarithms,
and put
\[
 \alpha_i=d\log x_i,\quad \beta_i=d\log y_i,\quad
 f_{ij}=\log(x_j/x_i),\quad g_{ij}=\log(y_j/y_i),
 \qquad
 \eta_{ij}=f_{ij}\beta_j-\alpha_i g_{ij}.
\]
Then
\[
 d\eta_{ij}=\alpha_j\wedge\beta_j-
             \alpha_i\wedge\beta_i,
 \qquad
 \eta_{12}-\eta_{02}+\eta_{01}
       =-d(f_{01}g_{12}).
 \tag{A31.SA4}
\]
For completeness, the second identity follows by substituting
$f_{02}=f_{01}+f_{12}$ and $g_{02}=g_{01}+g_{12}$: its left side is
$-f_{01}(\beta_2-\beta_1)-(\alpha_1-\alpha_0)g_{12}$.
Thus the ordered composition homotopy has function
$\kappa_{012}=-f_{01}g_{12}$, with $d\kappa=\delta\eta$.
Also $\delta\kappa=0$ by the same additive identities.
With the total differential $\delta+(-1)^{\text{Cech degree}}d$,
the corresponding total cocycle uses $-\kappa$ in Cech degree two.
This sign convention distinguishes a primitive of $\delta\eta$
from the signed term inserted in the total cocycle.

\paragraph{The surviving term at $p=3,e=3$.}
Write $3=u\pi^3$ with $u\in R^\times$.  If
$x_1/x_0=1+\pi A$ and $y_2/y_1=1+\pi B$, then reduction of
(A31.SA3) coefficientwise in the $\pi$-adically completed localized
ring, modulo $\pi$, gives
\[
 \overline f_{01}=\overline u^{-1}\overline A^3,
 \qquad
 \overline g_{12}=\overline u^{-1}\overline B^3,
 \qquad
 \overline\kappa_{012}
   =-\overline u^{-2}\overline A^3\overline B^3.
 \tag{A31.SA5}
\]
Only $n=3$ survives in each logarithm: for $n\neq3$ one has
$n-3v_3(n)\geq1$.  To check the latter, write $n=3^a m$ with
$3\nmid m$.  For $a=0$ it is clear; for $a=1$ it is
$3(m-1)>0$ unless $m=1$; for $a\geq2$ it is at least
$3^a-3a\geq3$.  The displayed function is closed in characteristic
three.  If $A=a/x$ and $B=b/y$, it includes the genuine local
cohomology term
\[
 -\overline u^{-2}\frac{a^3b^3}{x^3y^3}.
 \tag{A31.SA6}
\]
For example $a=b=1$ gives a nonzero negative Laurent monomial in
$H^2_{(x,y)}(k[[x,y]])$; no generic nonvanishing assertion about
all geometric coefficients is needed.  This term has total
denominator exponent six and lies outside the degree-zero
finite-pole piece $M_2^0$ used in C44-A31f.  At $e=2$ both
logarithms have coefficients divisible by $\pi$, so their product
vanishes modulo $\pi^2$; that argument cannot be reused at $e=3$.

\paragraph{The exact remaining calculation.}
To prove a higher-ramification period obstruction by this route one
still has to construct compatible transition and composition
homotopies for arbitrary changes of the local regular sequence,
identify their total cocycle with the supported canonical class
under (A31.SA2), and place that \emph{actual} cocycle in a finite
complex that retains terms such as (A31.SA6).  Finally one must
prove a primitive restriction/duality statement detecting the
resulting class in that larger complex.  Cohomological purity can
identify a completed supported class from its residues; it does
not discard a closed Cech-degree-two term or bound its poles.
Neither the convergence of (A31.SA3), the mere existence of
(A31.SA2), nor replacement of a computed homotopy by an arbitrary
primitive accomplishes this calculation.

This is a source audit and a local obstruction to reusing the
first-ramified complex.  It supplies neither a motivic
counterexample nor a proof of lifting or nonlifting at higher
ramification.

% END work/arithmetic/ramified_rigid_source_audit.tex

% BEGIN work/scan_w/quaternion_centralizer_comparison.tex
% C44_POS /root/scan_w. A36c reserved by root.
% Structural comparison within the quaternion advance; no new instance claimed.
\section{Quaternion complements and their simply connected centralizers}
\label{sec:c44-a36c}

\begin{theorem}[C44-A36c: an actual quaternion complement]
Let $p$ be any prime, put
$K=\operatorname{Frac}W(\overline{\mathbf F}_p)$, and let $\sigma$
be its absolute Frobenius. Consider a geometric quadratic cohomology
block of rank $n\geq3$, with a common algebraic cohomological
projector and the following data.
\begin{enumerate}
\item Its actual mixed-characteristic lift supplies a rational Betti
form $q_{\rm B}$, and comparisons identify
$q_{\rm B}\otimes\mathbf Q_\ell$ with its geometric form for every
$\ell\ne p$. The form $q_{\rm B}\otimes\mathbf Q_p$ has a self-dual
$\mathbf Z_p$-lattice. The actual normalized crystalline block
$(V_K,q_K,\Phi)$ has a self-dual $W(\overline{\mathbf F}_p)$-lattice
$M$. At $p=2$ both lattices are assumed even; in particular $n$ is
even and at least four. Write
\[
                  H=\sum_{a>0}a h_a
 \tag{A36c.1}
\]
for the positive elementary-divisor sum of $\Phi$ relative to $M$.
Any identification of $h_a$ with Hodge numbers is an actual
comparison hypothesis, as in A33/A41.
\item $T$ is a nondegenerate finite rational space of actual cycle
classes, with its rational geometric Gram form $q_T$. It embeds
with that Gram matrix in every realization; its normalized
crystalline classes are fixed by $\Phi$. Let $U=T^\perp$ in the
whole block. The projection onto $T$, and hence onto $U$, is the
common algebraic Gram projector. No lift of this projector or of
the cycles in $T$ is assumed.
\item On $U$ there is a common action by actual rational algebraic
correspondences of
\[
             \mathscr B=\prod_i B_i,
 \tag{A36c.2}
\]
where $B_i$ is a totally definite quaternion algebra over a totally
real number field $C_i$. The adjoint of the action is the canonical
quaternion involution. The identities hold in every realization.
The absolute crystalline operator $\Phi$ commutes with this abstract
rational algebra action, as follows from naturality for actual
degree-zero correspondences. No extension of the action to the lift,
and no integral stability under $\mathscr B$, is assumed.
\end{enumerate}
Then, for the ordinary real negative indices,
\[
                      \nu(q_T)\equiv\nu(q_{\rm B})+2H\pmod4.
 \tag{A36c.3}
\]
For a signed primitive block with its actual Hodge comparison,
$\nu(q_{\rm B})\equiv2H\pmod4$, so $4\mid\nu(q_T)$. A known
positive subspace of $T$ of codimension at most three then makes
$q_T$ positive definite.

More generally, if a known nondegenerate summand $A\subset T$ has
$\nu(A)\equiv\nu(q_{\rm B})+2H\pmod4$ and its orthogonal
complement in $T$ has dimension at most three, that complement is
positive definite. Application to the entire algebraic axis space
still requires a separate certificate that $T$ exhausts that space;
the theorem does not assert it from the quaternion action.
\end{theorem}

\begin{proof}
The empty complement is allowed, with the empty positive norm model
below. We give the proof for the nonempty factors.

\emph{1. A positive rational norm model, with all finite local forms.}
The central idempotents are self-adjoint and split $U$ orthogonally.
Trace duality turns the form on the $i$th factor into a symmetric
$C_i$-bilinear form whose adjoint on $B_i$ is canonical conjugation.
Choose pure quaternion generators
\[
 u_i^2=-d_i,\qquad z_i^2=c_i,\qquad u_i z_i=-z_i u_i,
 \qquad d_i\text{ totally positive},\quad c_i\text{ totally negative}.
 \tag{A36c.4}
\]
Such a choice exists because the reduced norm on the trace-zero
part of a totally definite quaternion algebra is totally positive.
Choose $u_i\ne0$ and then a nonzero trace-zero $z_i$ perpendicular
to $u_i$. The associated reduced-norm quadratic form is
\[
                   n_i=\langle1,d_i,-c_i,-d_i c_i\rangle.
 \tag{A36c.5}
\]
We use the symmetric bilinear form with this diagonal matrix;
its adjoint on left multiplication is canonical conjugation.

Over any characteristic-zero local centre field, a vector $v$ of
nonzero norm $a$ gives four mutually orthogonal vectors
$v,u_i v,z_i v,u_i z_i v$, of norms
$a,d_i a,-c_i a,-d_i c_i a$. They are independent and their
orthogonal complement is again $B_i$-stable. Induction therefore
expresses the centre-valued form as a sum of scalar multiples of
$n_i$, including when $B_i$ is split. In particular its centre-rank
is $4m_i$. These integers are common to all realizations: the
characteristic polynomial of a primitive element of $C_i$ has
rational, realization-independent coefficients by correspondence
traces, and its conjugate roots must have the same multiplicity.

Over a nonarchimedean local field, a four-dimensional form of square
determinant and its nonzero scalar multiple have the same Hasse
invariant. Indeed the usual scaling formula has factor
$(a,\det n_i)$ in rank four, which is one. Local classification
then gives $a n_i\simeq n_i$. Applying the field trace preserves
an isometry, so for every $\ell\ne p$ the actual form on $U$ is
isometric to the corresponding scalar extension of the rational form
\[
       N^+=\mathop{\perp}\limits_i
                    \bigl(\operatorname{Tr}_{C_i/\mathbf Q}n_i\bigr)^{m_i}.
 \tag{A36c.6}
\]
The crystalline comparison after extension to $K$ is constructed
separately in Step 2; no rational $\mathbf Q_p$-descent of $U$ is
asserted here. All scalars in the actual form may be different in
different realizations. They do not affect the preceding finite-place
isometry calculation.

The model $N^+$ is positive definite over $\mathbf R$. Its ordinary
determinant is a rational square: transfer of a rank-four form
multiplies the determinant by $\operatorname{disc}(C_i)^4$ and
applies the norm to $\det n_i$, itself a square. Consequently
orthogonality and the comparisons away from $p$ give
\[
 \det q_T=\det q_{\rm B}\pmod{\mathbf Q_\ell^{\times2}},\qquad
 \epsilon_\ell(q_T)\epsilon_\ell(N^+)=\epsilon_\ell(q_{\rm B})
                      \quad(\ell\ne p).
 \tag{A36c.7}
\]
The deleted-one-place squareclass injection of A15 makes the first
equality an equality of rational squareclasses. Thus it also holds
at $p$ and at the real place. This step does not infer the
$p$-adic Hasse invariant from an unramified scalar extension.

\emph{2. A quaternion-linear comparison over the unramified closure.}
We now keep track of the action, rather than only its quadratic
isometry class. For each place $v\mid p$ of $C_i$, the algebra
$B_i\otimes C_{i,v}$ splits over an unramified extension of
$C_{i,v}$. Such an extension is contained in the compositum with
the maximal unramified extension of $\mathbf Q_p$. Therefore every
factor of $\mathscr B\otimes K$ is a matrix algebra of degree two
over the corresponding factor of $C_i\otimes K$.

The preceding orthogonal decomposition into scalar norm forms is
valid over these factor fields as well. In a split quaternion
algebra the reduced norm is surjective: on $M_2(L)$ it is the
determinant. If $r$ has reduced norm $a$, right multiplication by
$r$ is an isometry from $(B_i,a n_i)$ to $(B_i,n_i)$, and commutes
with left multiplication by $B_i$. The orthogonal sum of its
inverse maps on these blocks produces an actual
$\mathscr B\otimes K$-linear isometry
\[
                         j:(N^+)_K\xrightarrow{\sim}U_K.
 \tag{A36c.8}
\]
This is stronger than an unstructured isometry of quadratic spaces.
It matches the action of the abstract rational centre and hence
all its coefficient-field factors.

Let $\sigma_N$ be coefficient Frobenius on the rational model
$N^+$. Define
\[
                 b_U=j^{-1}\Phi j\,\sigma_N^{-1}.
 \tag{A36c.9}
\]
It is $K$-linear and isometric. Both $\Phi$ and $\sigma_N$ commute
with the abstract rational algebra $\mathscr B$, so $b_U$ commutes
with it. This proves the assertion even when $\sigma$ permutes the
factors of $C_i\otimes K$: the two semilinear maps induce the same
permutation, which cancels in (A36c.9). No individual centre factor
is silently declared Frobenius-stable.

\emph{3. The centralizer is simply connected.}
Let $G_N$ be the algebraic isometry group of $N^+$ which centralizes
$\mathscr B$ pointwise. It is a product, over the centres, of forms
of symplectic groups. To see this directly after splitting $B_i$,
write the module as $L^2\otimes M_i$. Canonical involution on
$M_2(L)$ is adjunction for the standard alternating form on $L^2$.
The symmetric form on the tensor product is consequently the tensor
of that alternating form with a nondegenerate alternating form on
$M_i$, whose dimension is $2m_i$. Its centralizing isometry group
is $\operatorname{Sp}(M_i)$.

Thus $G_N$ is a semisimple simply connected group, specifically a
product of restrictions of scalars of inner forms of
$\operatorname{Sp}_{2m_i}$. It is connected and lies in
$\operatorname{SO}(N^+)$. Formula (A36c.9) is an element of
$G_N(K)$, with its actual rational form and Frobenius descent;
there is no replacement by a split group over $\mathbf Q_p$.

Kottwitz's functorial homomorphism on points of connected reductive
groups kills a semisimple simply connected group. Therefore the
image of $b_U$ has zero Kottwitz invariant in every orthogonal group
into which this centralizer is embedded. This also follows by
lifting the embedding to the spin cover: the identity component of
its pullback is a central isogeny onto the simply connected source,
and hence is an isomorphism. This argument does not require $b_U$
to be basic, and it treats its nonzero Newton slopes as well as its
slope-zero part.

\emph{4. Keep the rational reference twist at $p$.}
Put
\[
                   R=(T_{\mathbf Q_p},q_T)\perp(N^+_{\mathbf Q_p},q_N).
 \tag{A36c.10}
\]
The actual crystalline cycle vectors identify $V_K$ with $R_K$,
under which its semilinear operator is
\[
                          b_R\sigma_R,
                    \qquad b_R=1_T\oplus b_U.
 \tag{A36c.11}
\]
Here $R$ is only a rational comparison space. We do not identify
it with the Betti form over $\mathbf Q_p$.

By (A36c.7), $R$ and $q_{{\rm B},p}$ have the same dimension and
determinant squareclass. They become isometric over an unramified
quadratic extension: that extension kills the difference of their
two-torsion local Brauer invariants, and local quadratic-form
classification applies. In particular choose an oriented isometry
\[
                          g:R_K\xrightarrow{\sim}(q_{\rm B})_K
 \tag{A36c.12}
\]
defined over some finite unramified extension. Orientations may be
matched because the determinant ratio is a square in $\mathbf Q_p$.
The element
\[
                           b_0=g\sigma(g)^{-1}
 \tag{A36c.13}
\]
lies in $\operatorname{SO}(q_{{\rm B},p})(K)$. It is basic of
Newton slope zero, and the quadratic form on the fixed vectors of
$b_0\sigma$ is exactly $R$. Thus $b_0$ retains the possible
nontrivial rational Hasse difference between $R$ and the Betti
reference.

The actual crystalline operator in that reference has coefficient
\[
 b=g b_R\sigma(g)^{-1}
             =(g b_Rg^{-1})\,b_0.
 \tag{A36c.14}
\]
The first factor is in the image of the simply connected group of
Step 3, transported by the actual $K$-isometry $g$. Functoriality of
the Kottwitz homomorphism over $K$ makes its invariant zero. Since
the homomorphism is additive under multiplication,
\[
                              \kappa(b)=\kappa(b_0).
 \tag{A36c.15}
\]
This is the point at which omitting the reference twist would give
an unjustified rational-descent assertion.

The actual lattice $M$ and the Betti reference lattice become
integrally isometric over $W(\overline{\mathbf F}_p)$. At odd $p$
this is the classification of unimodular forms over that strictly
henselian ring; at two both lattices are even and reduce to the
same split nonsingular quadratic lattice. Using such an isometry
to compute the Cartan invariant gives
\[
                               \kappa(b)=H\pmod2.
 \tag{A36c.16}
\]
Changing the isometry only changes $b$ by $\sigma$-conjugacy.
The relevant orthogonal fundamental group is $\mathbf Z/2$, with
trivial Frobenius action, so its invariant does not change.
At two the reductive integral orthogonal group is the Dickson-kernel
group used in A41, not the naive determinant kernel.

The basic-form calculation of A33, and its direct dyadic proof A41,
applied to (A36c.13) now gives
\[
 \frac{\epsilon_p(R)}{\epsilon_p(q_{\rm B})}
                  =(-1)^{\kappa(b_0)}=(-1)^H.
 \tag{A36c.17}
\]
Both rational reference forms and their orientation have been kept
throughout. At two the factor $\epsilon_2(q_{\rm B})$ is retained;
it has not been set to one. As $N^+$ has square determinant,
\[
 \epsilon_p(q_T)\epsilon_p(N^+)
                      =(-1)^H\epsilon_p(q_{\rm B}).
 \tag{A36c.18}
\]

\emph{5. Rational reciprocity and the real sign.}
Combine (A36c.18) with (A36c.7) at every other finite place. Hilbert
reciprocity for the three rational forms, and positivity of $N^+$,
give
\[
                  \epsilon_\infty(q_T)
                         =(-1)^H\epsilon_\infty(q_{\rm B}).
 \tag{A36c.19}
\]
Their rational determinant squareclasses agree, so their negative
indices have the same parity. Write
$\nu(q_T)=\nu(q_{\rm B})+2k$. The ratio of their real Hasse
invariants is then $(-1)^k$. Equation (A36c.19) gives
$k\equiv H\pmod2$, proving (A36c.3).

For a signed primitive Hodge structure, its negative real blocks
are exactly those with odd positive Hodge displacement. Thus
$\nu(q_{\rm B})=2\sum_{a>0,\ a\ {
m odd}}h_a\equiv2H\pmod4$.
The positive-codimension-three conclusions follow by orthogonal
subtraction and the bound on the remaining negative index.
\end{proof}

\paragraph{Relation to A34/A36 and exact scope.}
The algebra generated there by $u$ and $F_s-F_s^{-1}$ on the
non-Tate space is a product of the totally definite quaternion
algebras of this theorem. Hence this centralizer proof gives a
structural formulation of the same mechanism. It needs neither
an inert imaginary quadratic field nor odd residue-field degree;
it permits ramified quaternion centres and retains their actual
Frobenius permutation. At two it still requires the actual even
unimodular whole-block lattices. No integral primitive splitting
or lift of the quaternion correspondences has been introduced.

For the congruence itself, $T$ need not be the whole potential Tate
space: its actual nondegenerate Gram certificate and the action
on its entire orthogonal complement suffice. If one uses the
non-Tate complement from A36, its full potential-Tate cycle
certificate supplies precisely this hypothesis. Neither version
constructs the quaternion action on an arbitrary axis; A36b gives
actual axes on whose non-Tate part no such action is possible.
No genuinely different eligible axis has been verified here, and
no separate progress-clock movement is claimed.

\paragraph{Primary inputs and their exact scope.}
Kottwitz, \emph{Isocrystals with additional structure. II},
Compositio Math.~109 (1997), pp.~293--299, Section~7.1--7.5,
constructs the functorial group homomorphism used in Steps 3--4;
Section~7.4 kills semisimple simply connected groups and
Section~7.5 compares it with the invariant on $B(G)$.
Its setup permits every prime. The basic rational-descent
classification is Kottwitz (1985), Section~4.5 and Proposition~5.6,
printed pp.~214 and 217; its ordinary quadratic normalization at
two is proved explicitly in A41. The hyperspecial Hodge invariant
is Rapoport--Richartz (1996), 4.1 and 4.2(i), pp.~178--180.

For the unramified splitting used in Step 2, see Milne,
\emph{Class Field Theory}, version~4.03 (2020), Chapter~IV,
pp.~139--140: a local division algebra has an unramified maximal
subfield, which splits it. The classical-group identification is
Milne, \emph{Algebraic Groups}, corrected CUP 2022 edition,
Section~24g and Theorem~24.47, pp.~529--530. The split tensor-form
calculation is supplied above, so no assertion about a simply
connected centralizer is inferred merely from a quadratic
isometry after scalar extension.

% END work/scan_w/quaternion_centralizer_comparison.tex

% BEGIN work/arithmetic/quaternion_picard_source_audit.tex
% C44_POS /root/fermat_aj. Exact source boundary; no new sign/Tate theorem.
\section{The rank-one Picard-motive boundary of a residual quadratic plane}
\label{sec:c44-a34-sa1}

\paragraph{C44-A34-SA1: what the geometric correspondences already give.}
Work over $\overline{\mathbf F}_p$ with rational coefficients, and fix
one Weil cohomology for the category of homological motives. Normalize
the relevant middle motive by the twist $(2)$, so that its pairing
takes values in $\mathbf1$. In the rank-two residual potential-Tate
plane of C44-A34a, let $R$ denote the resulting homological summand.
The actual correspondences $u,F_s$ have, in this category, the relations
\[
 u^2=-3,\qquad F_s^2=1,\qquad F_su=-uF_s,
 \qquad u^\dagger=-u,\quad F_s^\dagger=F_s.
 \tag{A34-SA1.1}
\]
The potential-Tate projector, the orthogonal projector off the known
positive span, and the semisimple Frobenius polynomial are algebraic
as cohomological correspondences. Their relations therefore hold
in homological motives, and hence in numerical motives. This does
not assert that the chosen representatives are idempotents modulo
rational equivalence.

The self-adjoint projectors
\[
             e_\pm=(1\pm F_s)/2
 \tag{A34-SA1.2}
\]
give nondegenerate rank-one homological summands $L_\pm$. The map $u$
interchanges them and is an isomorphism, with inverse $-u/3$.
Consequently, writing $L=L_+$,
\[
                       R\simeq L\oplus L.
 \tag{A34-SA1.3}
\]
Each restricted perfect form identifies $L$ with its dual. More
precisely, evaluation $q_L:L\otimes L\to\mathbf1$ and coevaluation
formed with that duality are inverse: their composite on $\mathbf1$
is the categorical dimension $1$, and the other composite is the
identity on the one-dimensional realization of $L\otimes L$.
Faithfulness of the chosen homological realization makes the latter
an equality of morphisms. Thus
\[
                         L^{\otimes2}\simeq\mathbf1
 \tag{A34-SA1.4}
\]
in homological motives and in numerical motives. No semisimplicity
conjecture, standard conjecture $D$, or finite-dimensionality of the
Chow motive of the axis is used in this argument. It constructs a
two-torsion Picard class in those categories, not an Artin motive.

\paragraph{The exact remaining existence statement.}
For this particular $L$, the following conditions are equivalent:
\[
 \begin{gathered}
 \text{the residual plane has a nonzero algebraic cycle class},\\
 \operatorname{Hom}_{\mathrm{hom}}(\mathbf1,L)\ne0,\\
 L\simeq\mathbf1\text{ in homological motives},\\
 L_{\mathrm{num}}\simeq\mathbf1\text{ in numerical motives}.
 \end{gathered}
 \tag{A34-SA1.5}
\]
Indeed, project a nonzero cycle class to the two lines and, if needed,
apply $u$; this gives a nonzero morphism to $L$. Conversely such a
morphism is represented by an algebraic cycle. Its nonzero
one-dimensional realization has nonzero self-pairing, so the
restricted form constructs its inverse after dividing by that
nonzero rational number. This proves the first three equivalences.
Their implication to numerical triviality is immediate. For the
reverse implication lift a nonzero numerical morphism
$\mathbf1\to L_{\mathrm{num}}$ to an algebraic correspondence
representing a homological morphism to $L$. Its realization cannot
vanish, since a homologically trivial cycle is numerically trivial.
This returns to the second condition. Once one class exists, $u$
and $F_s$ give the full plane as in C44-A34a.

The one-dimensional realizations of $L_\pm$ have normalized Frobenius
$+1$ and $-1$, respectively; after a quadratic finite-field extension
they are both trivial representations. This is a realization
statement. It does not produce the morphism in (A34-SA1.5).
The $M_2$ action on $R=L\oplus L$ acts on its multiplicity space
and does not remove the common invertible factor $L$.

\paragraph{The exact published Picard conjecture.}
Kahn, \emph{Zeta Functions and Motives}, \emph{Pure and Applied
Mathematics Quarterly} 5 (2009), Conjecture~14.1, printed p.~548,
conjectures that invertible motives are generated by the Lefschetz
motive and quadratic Artin motives. Proposition~14.2 and its proof,
pp.~548--549, derive this from the relevant essential injectivity of
realization, in particular from Hodge or Tate hypotheses. The author
PDF is \url{https://webusers.imj-prg.fr/~bruno.kahn/preprints/qjpam.pdf}.
The adequate-equivalence and finite-dimensionality conventions are
stated in \S13, p.~545; an invertible rank-one even object such as
(A34-SA1.4) is already positive finite-dimensional in its own
homological or numerical category.

Over $\overline{\mathbf F}_p$, the Artin factors in that conjecture
are trivial. Since the present object has weight zero, the required
case is precisely that its particular two-torsion Picard class be
trivial. This is weaker than the full Picard conjecture but is not
proved by the cited proposition. Kimura nilpotence is used there
only after an isomorphism modulo homological equivalence has been
obtained. It does not create a nonzero morphism from the unit.
Likewise, conservation of existing isomorphisms is not essential
injectivity on isomorphism classes.

\paragraph{Connectedness and finite-field classification do not supply it.}
In a Tannakian setting, an invertible object of square one defines
a character with values in $\mu_2$. Connectedness of the relevant
motivic group, or just absence of this particular character, would
kill it. Such connectedness is additional motivic information.
Kahn, \emph{Chow--Lefschetz motives}, Remark~5.19, author p.~23,
explicitly separates effective descent from expected connectedness
over separably closed fields and notes that the restricted
Lefschetz-motive groups can be disconnected:
\url{https://webusers.imj-prg.fr/~bruno.kahn/preprints/Chow-Lefschetz-motives-Indag-corr.pdf}.
This is a warning about a restricted category of correspondences,
not a counterexample in the full category of numerical motives.

Milne, \emph{Motives over finite fields}, Proposition~2.6,
author p.~15, and Theorem~2.18, p.~22, classify simple motives by
their Frobenius numbers under the Tate assumption governing \S2.
The first paragraph of the proof of Proposition~2.6 invokes Tate
exactly to turn an equality of Frobenius data into a nonzero motivic
Hom. Thus it cannot identify the present line with the unit without
the missing input:
\url{https://jmilne.org/math/articles/1994aP.pdf}.

A related search lead, Jossen's \emph{On the relation between Galois
groups and motivic Galois groups}, Theorem~9.15, p.~50, concerns
\emph{integral Weil-finite Nori motives} for subfields of
$\mathbf C$; that base-field/category scope is stated on p.~1.
It is not a result about rational homological or numerical motives
over $\overline{\mathbf F}_p$:
\url{https://www.jossenpeter.ch/PdfDvi/MoFGandGalois.pdf}.
Only its primary statement and scope were checked for this audit;
no full-paper proof certification is asserted.

\paragraph{Chow and geometric scope.}
To assert the same splitting inside Chow motives requires actual
Chow idempotents or an additional justified idempotent-lifting
argument. Even after such a lift, finite-dimensionality alone does
not supply the unit morphism in (A34-SA1.5). The imaginary quadratic
action also does not show that the axis summand is of abelian type;
no Hodge or Tate input for abelian varieties is imported on that
basis. No supplied conditional 032 or Milne2022 statement is used
in this audit.

The useful conclusion is the exact boundary (A34-SA1.5), not its
discharge. No applicable unconditional Artin/Tate classification
for this geometric summand was located in the checked primary
sources. This is a source and gate audit, with no new closed-axis
example, no assertion that the potential plane is algebraic, and no
progress-clock movement.

% END work/arithmetic/quaternion_picard_source_audit.tex

% BEGIN work/arithmetic/fermat_regulator_source_audit.tex
% C44_POS /root/fermat_aj. Exact source scope, no new theorem.
% No 032 or Milne preprint input is used.
\section{Source audit: arithmetic models of regulator images do not supply the Chow lattice}
\label{C44-F07-source-audit-two}

\noindent\textbf{Audit identifier: C44-F07-SA2.}
The following checks test three possible ways to discharge the actual
group gate in C44-F07. None supplies finite generation of its projected
degree-four etale-motivic group $B$, vanishing of the descent obstruction
$D$ of C44-F25, or the family Chow vanishing in JAC-P18.

\paragraph{An arithmetic model for the algebraic representative.}
Achter--Casalaina-Martin--Vial,
\emph{Normal functions for algebraically trivial cycles are algebraic
for arithmetic reasons}, author PDF, Theorem~1 on p.~2,
Section~1.3 on p.~6, and Proposition~2.1 with its proof on pp.~7--9,
descend the Abel--Jacobi \emph{image} and its normal functions:
\url{https://www.math.colostate.edu/~achter/math/algnf0819.pdf}.
The kernel $P$ in Section~1.3 is the kernel of a morphism
$J(C)\to J_a^3(J)$ of complex tori. It is not identified with
the kernel of the correspondence on rational Chow groups.
The proof of Proposition~2.1 passes from torsion to all points by
density for morphisms of abelian varieties, after composing the cycle
family with Abel--Jacobi.

For our actual projector, any Hodge morphism
$H^1(C)(-1)\to N$ vanishes, since its source has types
$(2,1),(1,2)$ and $N$ has types $(3,0),(0,3)$.
Consequently the rational $N$-part of this algebraic representative
is zero. These arithmetic models therefore give no injectivity on
the projected algebraically trivial Chow group. Corollary~2.3 on
p.~9 establishes embedding independence of its Abel--Jacobi kernel,
not vanishing of that kernel. This is consistent with C44-F17 and
the exact one-Lefschetz-factor boundary in C44-F07-SA1.

\paragraph{The explicit cubic CM example at five has a different twist.}
Bars, \emph{On the Tamagawa number conjecture for Hecke characters},
22-page author PDF, Section~2, p.~5, explicitly replaces the full
motivic group by a constructed subspace in the weak conjecture:
\url{https://mat.uab.cat/~francesc/TamagHeckefi.pdf}.
Theorem~3.5 and Corollary~3.6, with their short proofs on p.~7,
and Theorem~5.12 with its proof on p.~16, concern the constructed
module $\mathcal R_\theta$. Remark~5.14 on pp.~16--17 retains this
restriction even when its regulator is injective.

Section~7, pp.~20--21, explicitly treats infinity type $(3,0)$
at $p=5$. Its constructed motivic group has coefficient twist
$w+l+1=4+l$, with $l\ge0$. Our group has cohomological degree
four and twist two; it would require $l=-2$.
The negative-$l$ extension in Section~6.2, p.~17, still requires
$0<-l\le\min(a,b)$ and $-w-2l\le-3$. At $(a,b)=(3,0)$ both
conditions fail for $l=-2$. This exclusion holds before any
comparison of the finite Hecke characters of the elliptic models.
Thus neither the twist range nor the constructed-subspace conclusion
provides the missing actual-group theorem.

\paragraph{The small-motive descent theorem has an essential geometric hypothesis.}
Colliot-Thelene, \emph{Descente galoisienne sur le second groupe de
Chow: mise au point et applications}, author version dated
28~August~2015, Theorem~4.1 and its complete proof on printed p.~122,
assume universal geometric $\operatorname{CH}_0=\mathbf Z$:
\url{https://www.imo.universite-paris-saclay.fr/~jean-louis.colliot-thelene/ChowH3_28August2015.pdf}.
This particular PDF is paginated 109--134.
The proof uses that hypothesis to replace geometric $K$-cohomology
by Picard data. It does not infer the replacement from a rational
point, cohomological dimension two, or complex multiplication.
It does not apply to $J$: its nonzero $H^{3,0}$ already prevents
the displayed geometric hypothesis. The cubic summand retains
that Hodge type, so simply inserting its projector does not
validate the same small-motive argument. The actual arithmetic
tame-symbol kernel is still the one in C44-F25.

\paragraph{Check scope.}
The cited statements and the specified displayed proofs were read.
For the first source the construction in Section~1.3 and the complete
two-step proof of Proposition~2.1 were checked; no new audit of all
its earlier cited papers is claimed. These are source exclusions,
not new vanishing results, and they do not change the review clock.

% END work/arithmetic/fermat_regulator_source_audit.tex

% BEGIN work/axis/nonliftable_surface_prime_extension.tex
% C44_POS /root/axis_independent. Side extension; no separate progress-clock event.
\section{The nonliftable surface construction in every characteristic
\texorpdfstring{$p\equiv3\pmod4$}{p congruent 3 modulo 4}}
\label{sec:c44-a14b}

\paragraph{Review status.}
This is a checked side extension of the checked C44-A10/A14
construction, reusing its exact fixed tensor certificate.  It is
not a new claim about the unresolved universal axis sign.

\begin{proposition}[C44-A14b: varying the residue characteristic]
For every prime $p\equiv3\pmod4$, over
$k=\overline{\mathbf F}_p$ the elliptic curve
$E:y^2=x^3-x$ has a smooth connected surface $S\subset E^4$
with all the properties (A14.1)--(A14.2).  In particular $S$ is
geometrically nonliftable to characteristic zero as an abstract
smooth projective variety.  It is defined over a finite subfield
of $k$.

There are genuine Lefschetz axes $B\subset E^6$ containing
$S\times\{0\}^2$ for which $A_B=\mathbf Qw$ with $w^2>0$.
The nonzero primitive class $w$ cannot lift algebraically on the
geometric generic fibre of any abstract smooth projective model
over a complete mixed-characteristic DVR of ramification index
$e<p-1$.  For a larger ramification index, the precise additional
initial-lift hypothesis of C44-A31 still gives the same exclusion.
\end{proposition}

\begin{proof}
We give the actual replacement divisors, so that this extension
does not silently reuse the characteristic-three graph formula.
Choose $\eta\in\mathbf F_{p^2}$ with $\eta^2=-1$, and put
\[
 \iota(x,y)=(-x,\eta y),\qquad
 \pi(x,y)=(x^p,y^p).
\tag{A14b.1}
\]
Then $\iota^2=-1$ and $\pi\iota=-\iota\pi$ because
$\eta^p=-\eta$.  Let $\chi$ be the quadratic character on
$\mathbf F_p$, extended by $\chi(0)=0$.  Since
\[
             (-x)^3-(-x)=-(x^3-x),\qquad \chi(-1)=-1,
\]
pairing $x$ and $-x$ gives
$\sum_x\chi(x^3-x)=0$.  Thus $\#E(\mathbf F_p)=p+1$,
the Frobenius polynomial is $T^2+p$, and
$\pi^2=-p$ by Cayley--Hamilton and Tate-module faithfulness.
Consequently $[p]=-\pi^2$ is purely inseparable, so $E$ is
supersingular.  These are the same elementary endomorphism and
point-count inputs used in A10, now with $p$ retained.

The auxiliary coefficient prime remains $5$; it differs from
every $p\equiv3\pmod4$.  On $H^1(E,\mathbf F_5)$ choose a basis
in which
\[
 I=\iota^*=\begin{pmatrix}2&0\\0&3\end{pmatrix},\qquad
 P=\pi^*=\begin{pmatrix}0&1\\-p&0\end{pmatrix}.
\tag{A14b.2}
\]
Indeed $I^2=-1$ and $\det I=1$ force both distinct eigenvalues
$2,3$ to occur.  The invertible $P$ interchanges their eigenspaces;
rescaling one eigenvector makes its upper-right entry one, and
$P^2=-p$ gives the other entry.  As in A10, write
$[0]=\lambda e_0e_1$ with $\lambda\in\mathbf F_5^\times$;
no symplectic normalization of the chosen basis is assumed.
Use the same basis and the same $\lambda$ on all four factors.

For each pair $i<j$, let $D_i=\operatorname{pr}_i^*[0]$ and
define the actual integral graph divisors
\[
 R_{ij,u}=(\operatorname{pr}_i+u\operatorname{pr}_j)^*[0]
                     -D_i-(\deg u)D_j,
 \qquad u=1,\iota,\pi,\pi\circ\iota,
\tag{A14b.3}
\]
where the four degrees are now $1,1,p,p$.  All these divisors
are defined over $\mathbf F_{p^2}$.  Put
$a=2i-2,b=a+1,c=2j-2,d=c+1$.  In the cross-coordinate order
$(ac,ad,bc,bd)$, their normalized modulo-five classes are
\[
 (0,1,4,0),\quad(0,3,3,0),\quad
 (1,0,0,p),\quad(2,0,0,3p).
\tag{A14b.4}
\]
For the last column the pullback matrix is $IP$, since pullback
reverses $\pi\circ\iota$.  Directly, the cross part for a matrix
$M$ is obtained from
\[
 (e_a+M_{11}e_c+M_{21}e_d)
 (e_b+M_{12}e_c+M_{22}e_d).
\]
The $ac,bd$ coefficient matrix of the last two columns is
$\left(\begin{smallmatrix}1&2\\p&3p\end{smallmatrix}\right)$,
whose determinant is $p\ne0$ in $\mathbf F_5$.

For $x\in\mathbf F_5$, denote its representative in
$\{0,1,2,3,4\}$ by $\langle x\rangle_5$.  The required integral
divisors are
\[
 \begin{aligned}
 F_{ac}&=3R_{ij,\pi}+4R_{ij,\pi\circ\iota},\\
 F_{ad}&=3R_{ij,1}+R_{ij,\iota},\\
 F_{bc}&=2R_{ij,1}+R_{ij,\iota},\\
 F_{bd}&=\langle-2p^{-1}\rangle_5 R_{ij,\pi}
              +\langle p^{-1}\rangle_5 R_{ij,\pi\circ\iota},
 \qquad F_{2i-2,2i-1}=D_i.
 \end{aligned}
\tag{A14b.5}
\]
Multiplying the columns in (A14b.4) verifies
$\operatorname{cl}_5(F_{ab})=\lambda e_ae_b$ for every pair
$a<b$.  At $p=3$ the last coefficients are $(1,2)$, exactly
the original formula.  Thus the only coefficient formula that
changes with the residue of $p$ modulo five is $F_{bd}$.

Use the \emph{same} seventy displayed integers $c_{abcd}$ from
A10 and set
\[
                     z_p=\sum_{a<b<c<d}c_{abcd}F_{ab}F_{cd}.
\tag{A14b.6}
\]
Its class modulo five is the same $\lambda^2u_0$.  The exact
derivative minor in the existing
\texttt{supersingular\_class\_certificate.json} is therefore
unchanged: it is
\[
 26723729698312903479574054159209851065213568
                         \equiv3\pmod5.
\tag{A14b.7}
\]
For the actual integral tensor the minor reduces to
$\lambda^{128}\cdot3=3$.  Consequently the actual
$\mathbf Q_5$ tensor has zero infinitesimal stabilizer.  This
is the fixed, independently verified certificate of A10, not
a new random search or an assertion about unknown higher
$5$-adic coordinates.

Now perform the complete geometric construction of A14 with
$z_p$ in place of $z$.  The identity
$xy=3((x+y)^2-x^2-y^2)$ in modulo-five cohomology, the direct
sums of inverse line-bundle pairs, and ample twists by fifth
powers give a split ample bundle $\mathcal E$ with
$c_2(\mathcal E)\equiv\operatorname{cl}_5(z_p)$.
The independent first-jet incidence bounds are unchanged:
the deeper matrix-rank stratum has codimension six in a
fourfold, and failure of transversality to the codimension-two
rank stratum has total codimension five.  Thus a general
maximal-minor locus is a smooth surface; nonvanishing of
$c_2(\mathcal E)$ proves it is nonempty.  The incidence
calculation gives $[S]=c_2(\mathcal E)$.

The explicit Hilbert--Burch resolution in A14 consists of
anti-ample line bundles on the abelian fourfold.  Its
cohomology gives connectedness and
$H^1(\mathcal O_S)=H^1(\mathcal O_{E^4})$ in every
characteristic.  The Kummer/Picard argument with coefficient
prime five gives $H^1(E^4,\mathbf Q_5)\simeq H^1(S,\mathbf Q_5)$.
The dual top-cup tensor and its zero-stabilizer calculation
use only this isomorphism and the same unit minor.  A
characteristic-zero lift would supply a nonzero real Hodge
complex structure in that stabilizer, contradicting it.
Every step and every invertibility used in this construction
is exactly as checked in A14; none requires the residue
characteristic to equal three.  The nonempty parameter open
and its chosen point descend to a finite field.

Embed $S$ into $E^6$ by two zero coordinates.  The prescribed
centre existence and positivity of A06, the primitive
semiregularity of A22 and the constrained moment construction
of A18 apply for every odd $p$.  Taking the polarization high
enough for all three conclusions and then a selected closed
axis gives $A_B=\mathbf Qw$, $w^2>0$.  If $e<p-1$, A31 has
$n_0=1$, so its initial smooth embedded lift is simply $S$.
It excludes algebraic lifting of $w$ for every such abstract
model.  For larger $e$ its explicit hypothesis is retained.
\end{proof}

\paragraph{Nonvacuous small ramification.}
For $p\ge7$ in this congruence class, the range
$2\le e\le p-2$ is nonempty.  Smooth models of the underlying
axis do exist: lift $y^2=x^3-x$, its zero section and the
chosen product polarization over $W(k)$, and lift the two
axis sections using ample $H^1$ vanishing and base change.
Smoothness of the proper complete intersection follows from
the smooth special fibre.  Base change then gives models
over ramified traits of every specified finite index.
The exclusion in the proposition applies to \emph{all}
abstract smooth projective models in the small-ramification
range, rather than just to these exhibited models.

\paragraph{Sources and check scope.}
The degree/Frobenius and exterior-cohomology inputs are the
Milne locators already checked in A10; the Hodge-circle,
determinantal, Hilbert--Burch and Kummer arguments are the
complete checked proofs in A10/A14.  The only new arithmetic
is (A14b.2)--(A14b.5); its four possible nonzero residues
$p\bmod5$ are verified directly by the displayed matrices.
The fixed exact tensor certificate is reused, not rerun as
an unrecorded probabilistic experiment.  This extension is
part of the A27/A29/A31 lifting reduction, with no separate
staged progress-clock event.

% END work/axis/nonliftable_surface_prime_extension.tex

% BEGIN work/scan_w/low_rank_tate_obstructions.tex
% C44_POS /root/scan_w; next gate after A15. No geometric counterexample claimed.
\section{The rank-two and rank-three potential Tate gate}
\label{sec:c44-a17}

\begin{proposition}[C44-A17: determinant parity for a fully algebraic Tate block]
Retain the fixed axis and the known positive span $K$ of C44-A15.
Let $\mathcal T_{\ell_0}$ be the full generalized
root-of-unity Frobenius subspace of the residual space
$W_{\ell_0}$, of dimension $\tau$. Suppose that this
\emph{entire cohomological subspace} is spanned over
$\mathbf Q_{\ell_0}$ by algebraic cycle classes.
Then the rational algebraic residual space is of dimension
$\tau$, its intersection form is nondegenerate, and its
ordinary Gram determinant is positive. In particular its
negative index is even. Thus:
\[
\begin{array}{c|c}
 \tau&\text{possible signatures }(s_+,s_-)\\ \hline
 1&(1,0)\\
 2&(2,0)\text{ or }(0,2)\\
 3&(3,0)\text{ or }(1,2).
\end{array}
 \tag{A17.1}
\]
Full algebraicity is an explicit hypothesis here; it is
not deduced from the Frobenius polynomial or assumed via
the Tate conjecture. No standard conjecture $D(B)$ is used.
\end{proposition}

\begin{proof}
The generalized root-of-unity part is orthogonal to the
remaining generalized eigenspaces: the inverse of a root
of unity is again a root of unity. Hence its pairing is
nondegenerate. Choose $\tau$ actual cycles spanning this
part and apply the common algebraic projector and rational
orthogonal subtraction of A15. Their Gram matrix $G_T$
is a nonsingular rational matrix. In every coefficient
realization these same cycles have matrix $G_T$, so they
are independent, have finite Frobenius orbit, and fill the
root-of-unity part, whose dimension is independent of
$\ell$. Their common finite field of definition also shows
that Frobenius on this part is semisimple of finite order.
This is a consequence of the spanning hypothesis, not an
assumption of semisimplicity on all middle cohomology.

Every other residual algebraic cycle lies in this same
subspace. Pairing with the selected basis and multiplying
by $G_T^{-1}$ gives rational coefficients, proving equality
of the rational cycle-class space with this rational span.
The nondegenerate Gram matrix also constructs an actual
special-fibre Chow projector onto it: use the selected
cycles and their rational intersection-dual cycles.
Their compositions are the identity matrix of intersection
numbers. Its image is a direct sum of $\tau$ Tate motives
and has these same independent cycle realizations, including
the crystalline realization. This observation requires
neither an extension of the projector over a trait nor
independence of an arbitrary homological equivalence relation.

Choose an even exponent killing all the roots of unity.
On the orthogonal non-Tate complement the operator
$S-S^{-1}$ is invertible and skew-adjoint. The proof of
A15 gives a positive rational number
$\delta=\psi(1)\psi(-1)$ representing its determinant
squareclass in every $\mathbf Q_\ell$, $\ell\ne p$.
The one embedded geometric lift gives the positive
rational Betti determinant $D$ of the full variable space.
Thus
\[
             \det G_T\equiv D/(\det G\,\delta)
                     \pmod{(\mathbf Q_\ell^*)^2}
                    \quad\text{for every }\ell\ne p.
\]
The missing-$p$ squareclass injection of A15 makes this
a global rational squareclass equality. Its right side
is positive. Hence $\det G_T>0$, and the negative index
of $G_T$ is even, giving (A17.1).
\end{proof}

\paragraph{The exact remaining Hilbert-symbol condition.}
For a rational nondegenerate form $q_T$ of rank two or
three with positive ordinary determinant, write
$\epsilon_v(q_T)$ for its usual Hasse invariant.
Its two possibilities in (A17.1) are distinguished by
\[
 q_T>0\quad\Longleftrightarrow\quad
 \epsilon_\infty(q_T)=1
 \quad\Longleftrightarrow\quad
 \epsilon_p(q_T)=\prod_{\ell\ne p}\epsilon_\ell(q_T).
 \tag{A17.2}
\]
The product is finite. These equivalences follow from
$\epsilon_\infty=(-1)^{s_-(s_--1)/2}$ and Hilbert reciprocity.
They isolate a $p$-adic invariant; determinant parity does
not determine it. The necessary local classification and
product formula are Ancona--Marmora, Propositions~4.1--4.2
and Theorem~4.3, printed p.~629.

\begin{proposition}[C44-A17a: sharp arithmetic ambiguity away from one prime]
For every prime $p$ there are positive definite and negative
definite rational binary forms $q_+$ and $q_-$ with the same
positive determinant such that
\[
                   q_+\otimes\mathbf Q_\ell\simeq
                   q_-\otimes\mathbf Q_\ell
                         \quad\text{for all }\ell\ne p.
 \tag{A17.3}
\]
There are likewise rational ternary forms of signatures
$(3,0)$ and $(1,2)$ satisfying (A17.3) and containing a
common positive line. Frobenius may be set to the identity
in all these arithmetic examples. Thus all coefficient
primes away from $p$, positive determinant, and a fully
Tate Frobenius polynomial do not by themselves distinguish
the two signatures in (A17.1).
These are quadratic-space examples, not algebraic varieties
or counterexamples to a standard conjecture.
\end{proposition}

\begin{proof}
If $p=2$, take $d=u=1$. If $p\equiv3\pmod4$, take
$d=p,u=1$. If $p\equiv1\pmod4$, take $d=p$ and choose
a prime $u\equiv3\pmod4$ with $(u/p)=-1$; Dirichlet's
theorem and the Chinese remainder theorem give such $u$.
Set
\[
                   q_+=\langle1,d\rangle,
 \qquad            q_-=\langle-u,-d/u\rangle.
 \tag{A17.4}
\]
Both determinants equal $d>0$, and the stated real signs
are immediate. Their finite Hasse invariants satisfy
\[
       \epsilon_v(q_+)=1,
 \qquad\epsilon_v(q_-)=(-u,-d/u)_v=(-u,-d)_v.
 \tag{A17.5}
\]
The last equality uses $(a,-a)_v=1$.

For $p=2$, the quaternion $(-1,-1)$ is ramified precisely
at $2$ and infinity. For $p\equiv3\pmod4$, the symbol
$(-1,-p)$ is negative at $p$ and infinity and positive
elsewhere: at $p$ use $(-1/p)=-1$, while at $2$ the
odd-unit formula gives $1$ since $-p\equiv1\pmod4$.
For $p\equiv1\pmod4$, the symbol $(-u,-p)$ is negative
at $p$ because $(-u/p)=(u/p)=-1$. At the auxiliary prime
$u$ it is positive, since quadratic reciprocity gives
\[
                     (-p/u)=-(p/u)=-(u/p)=1.
\]
At $2$ it is positive because $-u\equiv1\pmod4$;
at infinity it is negative; at every other finite prime
both coefficients are units and the symbol is positive.
Thus in all cases the only discrepancies in (A17.5)
are at $p$ and infinity.

The equality of dimension, determinant and Hasse invariant
over every $\mathbf Q_\ell$, $\ell\ne p$, proves (A17.3)
by local classification. Add $\langle1\rangle$ to both
forms to obtain the ternary examples and their common
positive line. The identity is an isometry with entirely
Tate characteristic polynomial in either case.

The formulas used above are independently checked in
Voight, \emph{Quaternion Algebras}, version~1.0.5,
\S12.4, especially (12.4.10) and (12.4.14), pp.~189--190.
The exact three choices of quaternion presentation above
also appear in Example~14.2.13, p.~214, after the proof
of Proposition~14.2.7 and Dirichlet Theorem~14.2.9,
pp.~213--214. Thus the local ambiguity is present at
every residue characteristic, not just at $2$.
\end{proof}

\begin{proposition}[C44-A17b: a Tate isocrystal need not be an admissible subobject]
There is, for every prime $p$, a weakly admissible filtered
$\varphi$-module over $\mathbf Q_p$ with a perfect symmetric
filtered $\varphi$-pairing and a nondegenerate rank-two
$\varphi$-fixed summand whose induced filtration is not
weakly admissible. Its Frobenius characteristic polynomial
can be rational with all complex roots of absolute value one.
Thus these properties of the full filtered module do not
make the special Tate projector a filtered projector.
\end{proposition}

\begin{proof}
Let $D$ have basis $t_1,t_2,e,f$, with symmetric pairing
\[
 q(t_i,t_j)=-\delta_{ij},\qquad q(e,f)=1,
 \qquad q(e,e)=q(f,f)=0,
 \qquad q(t_i,e)=q(t_i,f)=0.
\]
Choose the root $\lambda\in\mathbf Q_p$ of
\[
                       T^2-p^{-1}T+1
 \tag{A17.6}
\]
with valuation $-1$. Such a root exists because
$1-4p^2$ is a square in $\mathbf Q_p$ by Hensel's lemma;
choose its square root sufficiently close to $1$ and use
$(1+\sqrt{1-4p^2})/(2p)$. The other root is
$\lambda^{-1}$ of valuation $1$. Define
\[
 \varphi t_i=t_i,\qquad \varphi e=\lambda e,
 \qquad \varphi f=\lambda^{-1}f.
\]
This preserves $q$. The rational Frobenius polynomial is
$(T-1)^2(T^2-p^{-1}T+1)$; its complex roots all have
absolute value one. The last two are not roots of unity,
since their sum is the nonintegral rational number $1/p$.

The line $L=\langle t_1+e+f/2\rangle$ is isotropic.
Put
\[
 \operatorname{Fil}^1D=L,\quad
 \operatorname{Fil}^0D=L^\perp,\quad
 \operatorname{Fil}^{-1}D=D,
 \qquad \operatorname{Fil}^2D=0.
 \tag{A17.7}
\]
This filtration is self-dual for $q$. Its Hodge weights
are $1,0,0,-1$, so $t_H(D)=0=t_N(D)$.

Write $T=\langle t_1,t_2\rangle$, $E=\langle e\rangle$
and $F=\langle f\rangle$. Since the three eigenvalues
$1,\lambda,\lambda^{-1}$ are distinct, every
$\varphi$-stable subspace is $T'\oplus E'\oplus F'$
with $T'\subset T$, $E'$ either $0$ or $E$, and $F'$
either $0$ or $F$. For such a subspace $D'$ the induced
filtration has
\[
 t_H(D')=\mathbf1_{L\subset D'}-
                    \mathbf1_{D'\not\subset L^\perp},
 \qquad
 t_N(D')=-\dim E'+\dim F'.
 \tag{A17.8}
\]
If $E'=E$ and $F'=0$, the first indicator is zero and
the second is one, since $q(L,e)=1/2$; hence
$t_H(D')=-1=t_N(D')$. If $E'=0,F'=F$, one has
$t_H(D')\leq0<1=t_N(D')$. If both are present,
$t_N(D')=0$ and the second indicator is one, so
$t_H(D')\leq0$. If neither is present, the first
indicator is zero and again $t_H(D')\leq0=t_N(D')$.
This checks all subspaces and proves weak admissibility.

On the nondegenerate fixed subspace $T$, however, the
induced filtration is
\[
 \operatorname{Fil}^1T=0,\qquad
 \operatorname{Fil}^0T=\langle t_2\rangle,
 \qquad\operatorname{Fil}^{-1}T=T.
\]
Consequently $t_H(T)=-1$ whereas $t_N(T)=0$.
It is not weakly admissible. Equivalently the orthogonal
projector onto $T$ sends $L$ to $\langle t_1\rangle$
and fails to preserve the filtration. This failure
occurs despite the full module's self-dual filtered
structure and weak admissibility.
\end{proof}

\paragraph{What the rank-two theorem really supplies.}
Ancona, \emph{Standard conjectures for abelian fourfolds},
Theorem~8.1 (author arXiv:1806.03216v3, p.~25), does
\emph{not} assume a given CM action: it proves positivity
for a rank-two motive with quadratic form in mixed
characteristic, full numerical algebraic dimension two
on the special fibre, and a polarizing Betti form.
Proposition~8.3, p.~26, deduces the necessary all-realization
algebraicity and nondegeneracy under those hypotheses.
Proposition~10.8, pp.~32--33, derives a quadratic
$p$-adic field action using the orthogonal group of the
\emph{admissible rank-two object}. The isotropic-line
argument and admissibility are used explicitly in
excluding a split quadratic algebra. No such admissible
rank-two object is supplied merely by a Frobenius
projector on the special fibre; A17b shows the inference
from full weak admissibility is false.

The higher-rank signature theorem of Ancona--Marmora,
Theorem~3.2 and Remark~3.3, printed p.~626, does assume
a relative number-field action of full degree with
adjoint involution. Its remark explicitly places the
removal of those action hypotheses in Conjecture~3.1.
We do not transfer that theorem to an arbitrary
special Tate block. Neither numerical semisimplicity
nor a special-fibre identification with a sum of Tate
motives specifies the geometric form's real signature.

\paragraph{Status of the next axis gate.}
For cote~44, Theorem~5, scans~004 and~008, A17 gives a
determinant constraint when the residual block is
actually algebraic. A17a proves that the resulting
arithmetic constraints still allow both remaining
rank-two and rank-three signatures; A17b identifies
the precise missing filtered-subobject step in the
attempt to apply the rank-two $p$-adic theorem.
The residue-characteristic Hilbert invariant in
(A17.2) is not computed here for an arbitrary axis.
There is no geometric wrong-sign example and no claim
that the residual rank-two or full sixfold
multiplicity-three gate has been closed.

\paragraph{Primary sources actually read.}
Ancona's Theorem~8.1, Proposition~8.3, the complete
quadratic-form reduction in \S9, the comparison
discussion in \S10, and the complete proof of
Proposition~10.8 were read in the downloaded primary
author version
\url{https://arxiv.org/pdf/1806.03216v3}.
Ancona--Marmora's Theorem~3.2/Remark~3.3 and the full
proof of Corollary~4.4 were read at
\url{https://ems.press/content/serial-article-files/51595?nt=1}.
The explicit Hilbert symbols and the quaternion
construction of Voight cited above were read in
\url{https://jvoight.github.io/quat-book-v1.0.5.pdf}.
The filtered-module construction A17b is proved
directly above; no geometric realization theorem
for this example is invoked. In particular it does not
come with a compatible rational Betti factor whose form
polarizes its Hodge structure. It refutes only the
filtered-subobject inference, not the combined geometric
hypotheses of Ancona's theorem.

% END work/scan_w/low_rank_tate_obstructions.tex

% BEGIN work/scan_w/odd_tate_selection_axes.tex
% C44_POS /root/scan_w; odd-rank companion to A12 and application of A15.
\section{Selecting one potential Tate direction and positive spans of arbitrary rank}
\label{sec:c44-a16}

\begin{lemma}[C44-A16: odd orthogonal selection]
Let $U/\mathbf F_q$ be smooth, geometrically connected and
quasiprojective of positive dimension. Let $\mathcal W$ be
a lisse $\mathbf Q_\ell$-sheaf, $\ell\ne p$, of odd rank
$d\geq1$, endowed with a nondegenerate symmetric form
preserved by arithmetic monodromy. Suppose that its geometric
monodromy closure contains $\operatorname{SO}(\mathcal W)$
and that every closed-point Frobenius characteristic
polynomial lies in $\mathbf Q[T]$.
Then a Zariski-dense set of closed points has exactly one
root-of-unity Frobenius eigenvalue counted with algebraic
multiplicity. This eigenvalue is simple and is $1$ or $-1$.
No semisimplicity or algebraicity of this line is asserted
or required.
\end{lemma}

\begin{proof}
Put $M=\operatorname{lcm}(1,\ldots,2d^2)$, which is even.
If a root of unity of order $n$ occurs in a rational degree-$d$
polynomial, then $\varphi(n)\leq d$ and $n\leq2d^2$.
For completeness, the elementary latter bound follows from
\[
 \frac{\varphi(n)^2}{n}
       =\prod_{a^e\parallel n}a^{e-2}(a-1)^2\geq\frac12,
 \tag{A16.1}
\]
where the product is over prime powers: the possible factor
less than one is $a=2,e=1$, whose value is $1/2$.
Consequently all root-of-unity eigenvalues of any Frobenius
in the statement are killed by the exponent $M$.

For an orthogonal matrix $g$ define the conjugacy-invariant
regular function
\[
 f(g)=[T^1]\det\bigl(TI-(g^M-I)\bigr),
 \tag{A16.2}
\]
where $[T^1]$ means the coefficient of $T$. In odd dimension
every orthogonal matrix has an eigenvalue $1$ or $-1$.
Thus $g^M$ has eigenvalue one, and $f(g)\ne0$ means precisely
that its eigenvalue one has algebraic multiplicity one.
This is a statement about a characteristic polynomial, so
Jordan blocks elsewhere cause no difficulty.

The restriction of $f$ to $\operatorname{SO}_d$ is not zero.
Over an algebraic closure a maximal torus has eigenvalues
\[
             1,t_1,t_1^{-1},\ldots,t_{(d-1)/2},t_{(d-1)/2}^{-1}.
\]
Choose each $t_i^M\ne1$. Then the zero eigenvalue of
$g^M-I$ is simple, so $f(g)\ne0$. For $d=1$ the function
on $\operatorname{SO}_1$ is simply $1$.
Zariski density of the geometric monodromy image in its
closure supplies a geometric image element $\gamma$ with
$f(\gamma)\ne0$.

Let $\Gamma$ be the compact arithmetic image. By continuity
and conjugacy invariance there is an open normal subgroup
$J\subset\Gamma$ such that the inverse image of the
conjugacy class of $\gamma$ in $\Gamma/J$ is contained in
$f\ne0$. This finite quotient comes from a finite etale
Galois cover of $U$. Its selected conjugacy class lies in
the geometric subgroup, because $\gamma$ does.

Apply finite-field Chebotarev with the constant-field coset
condition, as in C44-A12. More explicitly, if the finite
quotient cover has constant-field degree $\nu$, choose
sufficiently large extension degrees divisible by $\nu$;
the chosen geometric class is then in the permitted coset.
Entin, Theorem~3 in \S4, author arXiv:1805.05454v3,
pp.~9--11, supplies points over such extensions with
Frobenius in that class. At those points $f(F)\ne0$.
By the cyclotomic bound and (A16.2), their total
root-of-unity multiplicity is exactly one.

If this point is presented as rational over an extension,
its Frobenius is a positive power of the Frobenius of its
underlying closed point. Passing to a positive power
preserves the total multiplicity of root-of-unity
eigenvalues: roots of unity remain such, and a number
whose positive power is a root of unity was already one.
Thus the underlying closed point has the same multiplicity
one. The construction works on every nonempty open subset
of $U$, whose fundamental group surjects onto that of $U$;
the monodromy image is unchanged. The selected closed
points are therefore Zariski dense. Finally, rationality
of their polynomial forces its unique root of unity to
be rational, hence $1$ or $-1$.
\end{proof}

\begin{corollary}[C44-A16a: closed supersingular axes with positive spans of arbitrary rank]
Let $N\geq0$ be any integer. With $E:y^2=x^3-x$, $A=E^6$
and $H=\sum h_i$ as in C44-A13, choose $N$ distinct points
$a_i\in E^4(K)$ over a finite number field and put
$Y_i=E^2\times\{a_i\}$. For $N=0$ take $K=\mathbf Q$
and impose no centre. For every sufficiently large $t$
there is an incidence family on a nonempty open
$U\subset\mathbf P^1_K$ of smooth genuine Lefschetz
axes $B_u=D_0\cap D_u$, $D_0,D_u\in|tH|$, containing all
the $Y_i$, with the following property.

After finitely many rational primes have been excluded,
for every remaining $p\equiv3\pmod4$ and every place of
$K$ above $p$, a Zariski-dense set of closed points of the
reduced parameter curve has the required positive axis
sign on the \emph{whole} $A_{B_u}$, unconditionally. The
ambient sixfold is supersingular. The explicit classes
\[
   w_i=[Y_i]_{B_u}-t^{-2}\kappa^*\alpha_0
\]
span a positive $N$-dimensional rational subspace $K_u$
with Gram matrix $6t^2I_N-J_N/(4t^2)$, provided
$24t^4>N$. At the selected points:
\begin{itemize}
\item if $N$ is odd, $A_{B_u}=K_u$;
\item if $N$ is even, $A_{B_u}=K_u$ or
      $A_{B_u}=K_u\perp\mathbf Qz_u$, with $q_B(z_u,z_u)>0$.
\end{itemize}
The known span $K_u$ has relative algebraic lifts. In the
even-$N$ case, no assertion is made that the possible
additional Tate line is algebraic or that a class on it
lifts. For $N=0$ the conclusion is positivity of a kernel
of dimension zero or one; nonvanishing is not claimed.
\end{corollary}

\begin{proof}
The disjoint-centre construction of A13b applies for any
$N$ up to, and including, its monodromy calculation. To
make the lack of a parity assumption at those steps
explicit, take $Z=\operatorname{Bl}_{\coprod Y_i}A$,
$L'=\pi^*\mathcal O_A(tH)(-\sum F_i)$, choose a general
$\widetilde D_0\in|L'|$, and vary a general auxiliary
hyperplane pencil on its induced embedding. For $N=0$
this just means $Z=A$ and $L'=\mathcal O_A(tH)$.
The normal-row ranks four and the rank-drop codimension
three on each two-dimensional centre give the original
smooth axes, and the nonempty open of Lefschetz lines
through $\widetilde D_0$ imposes the original-pencil
condition as in A13b. These statements do not involve
the parity of $N$.

Every centre has $H^2\alpha_0=[Y_i]_A$, normal self-intersection
$[Y_i]^2_{B_u}=6t^2$, and $\int_{Y_i}\alpha_0=1/4$.
Distinct centres are disjoint. These equalities give
the displayed Gram matrix, positive for $24t^4>N$.
The blockwise blowup formula and excess formula give
\[
 \mathcal W_N=\ker\kappa_*\cap\bigcap_iw_i^\perp
       \simeq \hbox{the auxiliary-pencil vanishing system},
 \qquad
 \operatorname{rank}\mathcal W_N=3600t^6-165-N.
 \tag{A16.3}
\]
The isomorphism is isometric, as shown in A13b; disjointness
makes all exceptional blocks orthogonal. The formula remains
valid with an empty exceptional sum when $N=0$.
For every $N$ and sufficiently large $t$ the rank is positive
and
\[
                  h^{4,0}(\mathcal W_N)=62t^6-5>0.
\]
SGA~7~II, XIX, Proposition~3.4 and \emph{Weil II},
4.4.1--4.4.2s therefore give full orthogonal geometric
monodromy in characteristic zero, independently of parity.

Spread this fixed incidence model and its finite omitted
parameter divisor over $\mathcal O_K[1/M_0]$, retaining
smooth families, disjoint centres, and finite etale
boundary. Exclude the finitely many rational primes below
all bad places or collisions, together with $2$, $\ell$
and the prime divisors of
$|\operatorname{GL}_{3600t^6-165-N}(\mathbf F_\ell)|$.
The stable-integral-image argument of A13 and A13b, using
the prime-to-$p$ specialization isomorphism of SGA~1,
XIII, 2.9--2.12, then identifies the geometric image at
every remaining place with the characteristic-zero image.
It applies equally in odd and even ranks. Thus the full
orthogonal closure persists at every remaining place.
Infinitely many rational primes $p\equiv3\pmod4$ remain,
and $E$ is supersingular at all places above them.

The Frobenius polynomial on $\mathcal W_N$ is rational:
divide the normalized middle polynomial of $B_u$ by
the ambient factor from $A$ and by $(T-1)^N$ supplied
by the known relative cycles. If $N$ is odd, the rank
in (A16.3) is even and C44-A12 supplies dense closed
points with no residual root of unity, giving
$A_{B_u}=K_u$ exactly as in A13b.
If $N$ is even, that rank is odd. C44-A16 supplies
dense closed points with residual cyclotomic
multiplicity one. C44-A15, with the known positive
span $K_u$, proves positivity of the entire algebraic
kernel and the stated dimension alternatives. This
last application does not infer an algebraic cycle
from a Frobenius root.

The selected closed parameter lifts in the fixed open
incidence model after an unramified extension of a
completed local trait. The centres and the divisor
polynomial $\alpha_0$ lift with it, giving relative
algebraic representatives of the known $w_i$. This
proves exactly the asserted liftability of $K_u$.
No construction or lifting of a possible additional
class is part of the argument.
\end{proof}

\paragraph{Programme scope and source.}
This extends the whole-axis sign family for cote~44,
Theorem~5, scan~004 and its primitive margin on scan~008
to prescribed positive spans of every rank. The axis conclusion is
unconditional. As in A15, the corresponding full-height
and $I(B)$ conclusions through C44-A01 retain the separate
unrefereed 032/Milne2026 abelian input. The general fixed-axis
case is not proved. The finite-field selection uses Entin's
complete Theorem~3 and proof in \S4, read at
\url{https://arxiv.org/pdf/1805.05454v3}, pp.~9--11, with
the constant-field congruence condition explicitly retained.
All geometric and monodromy source locators are those
verified in A13 and A13b; the sign of a possible final
algebraic line is supplied by the independently checked A15.

% END work/scan_w/odd_tate_selection_axes.tex

% BEGIN work/axis/axis_ring_obstruction.tex
% C44_POS /root/axis_independent; GPT-6 Astra Pro; effort not exposed.
\section{The cohomology-ring boundary of the tensor lifting method}
\subsection*{C44-A14a. Orthogonal freedom of the variable middle summand}

\textbf{Proposition.}
Let $X$ be an abelian variety of dimension $g=2r+2$, $r\ge1$,
over an algebraically closed field, and let
$\kappa:B\hookrightarrow X$ be a smooth complete intersection of
two ample divisors of the same class $h$. Work in rational
$\ell$-adic cohomology, with $\ell$ different from the
characteristic, and suppress Tate twists. Put
\[
 V=H^1(X),\qquad
 W=\ker\bigl(\kappa_*:H^{2r}(B)\to H^{2r+4}(X)\bigr),\qquad
 q(w,w')=(-1)^r\int_Bww'.
                                                               \tag{A14a.1}
\]
The ambient image is the graded algebra
\[
 R=\operatorname{im}\kappa^*
       \simeq\frac{\bigwedge^\bullet V}{\operatorname{Ann}(h^2)},
 \qquad
 \operatorname{Ann}(h^2)=\{\alpha:\alpha h^2=0\}.              \tag{A14a.2}
\]
Every cohomology group of $B$ is ambient except for the direct
summand $W$ in degree $2r$. If $\omega_B$ is the top class with
$\int_B\omega_B=1$, all remaining multiplication is described by
\[
   w\alpha=0\quad(w\in W,\ \alpha\in R^{>0}),\qquad
   ww'=(-1)^r q(w,w')\omega_B.                               \tag{A14a.3}
\]
Let $\operatorname{Sp}(V,h)$ denote the stabilizer of the
nondegenerate degree-two exterior tensor $h$.
There is a faithful action
\[
 \operatorname{Sp}(V,h)\times\operatorname O(W,q)
          \lhook\joinrel\longrightarrow
  \operatorname{Aut}_{\rm gr}\bigl(H^\bullet(B),h|_B,\int_B\bigr)
                                                               \tag{A14a.4}
\]
under which the first factor acts on $R$ through exterior powers
and fixes $W$ pointwise, while the second factor fixes $R$ pointwise.
In particular, the infinitesimal stabilizer of the cup ring,
polarization and \emph{any collection} of variable classes
contains the nonzero Lie algebra $\mathfrak{sp}(V,h)$.

\textbf{Proof.}
The exterior algebra description of $H^\bullet(X)$ is Milne,
\emph{Abelian Varieties}, I Theorem 12.1, printed pp.~55--56.
Self-intersection gives
\[
                         \kappa_*\kappa^*\alpha=h^2\alpha.
                                                               \tag{A14a.5}
\]
For $j\le2r$ weak Lefschetz makes $\kappa^*:H^j(X)\to H^j(B)$
injective, and it is an isomorphism for $j<2r$.
Hard Lefschetz makes $L_h^2$ injective in all these degrees.
Consequently both kernels in (A14a.2) are zero there.

For $j>2r$, the Gysin map
\[
                    \kappa_*:H^j(B)\longrightarrow H^{j+4}(X)
\]
is an isomorphism: its Poincare dual is the weak Lefschetz
isomorphism in degree $4r-j<2r$. The map $L_h^2$ onto
$H^{j+4}(X)$ is surjective. One way to see the last assertion
is to apply hard Lefschetz in the dual degree; equivalently,
it is the upper-half surjectivity in each Lefschetz string.
Equation (A14a.5) therefore makes $\kappa^*$ surjective in
these degrees and identifies its kernel with
$\ker L_h^2$. This proves (A14a.2) in every degree. The
annihilator is an ideal, as required for the displayed algebra
quotient.

In degree $2r$, $L_h^2:H^{2r}(X)\to H^{2r+4}(X)$ is an
isomorphism. Thus self-intersection gives the orthogonal
decomposition
\[
                        H^{2r}(B)=R^{2r}\perp W.
                                                               \tag{A14a.6}
\]
The form on $R^{2r}$ is
$(-1)^r\int_X\alpha\beta h^2$, which is nondegenerate by
hard Lefschetz, so $W$ is nondegenerate too.

If $w\in W$ and $\alpha$ has positive ambient degree,
projection formula gives $\kappa_*(w\kappa^*\alpha)=0$.
The product is above the middle degree; the Gysin isomorphism
just proved makes it zero. Products past the top degree are
zero as well. The product of two degree-$2r$ variable classes
is determined by their top trace, proving (A14a.3).

An element of $\operatorname{Sp}(V,h)$ acts on the exterior
algebra, preserves $h^2$ and its annihilator, and therefore
acts on $R$. It preserves the top trace because
$h^g\ne0$ generates the one-dimensional top cohomology of $X$:
fixing $h$ fixes that generator. In particular it preserves
the form on $R^{2r}$ and fixes $\omega_B$.
Extend this action by the identity on $W$.
Similarly extend an isometry of $(W,q)$ by the identity on
$R$. Formulas (A14a.3) verify every mixed or variable product;
ambient products are already preserved. The two actions commute.
Their product is faithful because the first factor is
faithful in degree one and the second is faithful on $W$.
This proves (A14a.4).

Finally, any collection of elements of $W$ is fixed by the
entire first factor. Its Lie algebra has dimension
$g(2g+1)>0$. Thus the asserted infinitesimal stabilizer is
nonzero, for every choice of those elements.
\hfill$\square$

\paragraph{Exact consequence for the lifting attack.}
C44-A10 and C44-A14 obstruct lifting by exhibiting an
intrinsic fourfold cup tensor on an eight-dimensional
$H^1$ space whose infinitesimal stabilizer is zero.
The proposition proves that this criterion cannot be
recovered from the graded cup ring of an axis, its
polarization, and variable classes alone. Those data retain
the full ambient symplectic symmetry acting in degree one.
In particular, replacing the A14 centre by several centres,
or replacing its class by a difference with zero ambient
pushforward, cannot retain that same zero-stabilizer
certificate through ordinary cup products.

This is a limitation of that specific tensor method.
It does not assert that a variable class lifts, that a
formal Hodge structure on this algebra is realized
geometrically, or that the primitive intersection form is
positive. Constraints involving the period variation,
the Hodge numbers of a lift, algebraic correspondences
outside the cup ring, or integral and crystalline data
are additional information and are not addressed by
(A14a.4).

\paragraph{Programme locator.}
The variable summand is the exact $A_B$ cohomological
container in JAC1964 H9-T7, hence the remaining term in
cote~44, Theorem~5, scan~004, with the primitive margin
on scan~008. This proposition sharpens the obstruction
to route (a), the universal lifting attack, and to
cup-product versions of route (b). It supplies no
additional instance of the desired sign and does not
close QH9-axis.

% END work/axis/axis_ring_obstruction.tex

% BEGIN work/covers/contribution.tex
\section{Special cover strata and the old axis summand}

\subsection{C44-C01: the old variable classes do not disappear}
\begin{proposition}[C44-C01]\label{prop:old-axis-cover}
Use JAC1964 H9-T18. Write $A$ for the abelian fourfold, $B$ for its
surface axis, $X_0=\operatorname{Bl}_B A$, $M_0=A\times\mathbf P^1$,
and $k_0:X_0\hookrightarrow M_0$ for the incidence embedding.
For a finite separable cover $\pi:C\to\mathbf P^1$ of degree $b$,
unramified over the critical values, put
$X_C=X_0\times_{\mathbf P^1}C$, $M_C=A\times C$, and let $g:X_C\to X_0$
and $k_C:X_C\hookrightarrow M_C$ be the maps. On algebraic middle classes,
\[
 \Alg^2(X_C)\cap\ker k_{C*}
 =g^*(\Alg^2(X_0)\cap\ker k_{0*})\ \mathbin\perp\
   (\Alg^2(X_C)\cap\ker k_{C*}\cap\ker g_*).
 \tag{C44.C01.1}
\]
The first summand is canonically the old axis space
$A_B=\ker(\kappa_*:\Alg^1(B)\to H^6(A)(1))$ and has form
\[
          (w,w')\longmapsto-b\int_Bww'.                 \tag{C44.C01.2}
\]
It is positive definite, by the surface Hodge index theorem, without 032.
Consequently the remaining H9-T18 sign is exactly the sign on the second
summand in (C44.C01.1). Statements that generic covers add no algebraic
middle classes refer to this new part (or the full trace kernel), and do
not say that the entire H9-T18 variable space vanishes.
\end{proposition}
\begin{proof}
The finite flat map $g$ has $g_*g^*=b$. Projection formula makes
$P=b^{-1}g^*g_*$ a self-adjoint algebraic idempotent. The Cartesian
incidence square gives
$k_{C*}g^*=(1_A\times\pi)^*k_{0*}$, and proper functoriality gives
$k_{0*}g_*=(1_A\times\pi)_*k_{C*}$. Thus $P$ preserves
$\ker k_{C*}$, and its image there is exactly the first summand of
(C44.C01.1); its kernel is the displayed second summand. Both are
algebraic subspaces, since $P$ and $1-P$ are algebraic. Self-adjointness
proves orthogonality.

For the old space use the codimension-two blowup decomposition from
H9-T7. A middle class has unique form $q^*z+e_*p^*w$ with
$z\in H^4(A)$ and $w\in H^2(B)$. With
$t=[\mathrm{point}]\in H^2(\mathbf P^1)$, incidence pushforward is
\[
 k_{0*}(q^*z+e_*p^*w)=(hz+\kappa_*w)+zt.
\]
The exceptional divisor is $B\times\mathbf P^1$ as a variety over
$A\times\mathbf P^1$; this proves the second term's formula, while
the first follows from $[X_0]=h+t$. Kunneth independence forces $z=0$
and $\kappa_*w=0$ in the kernel. This is the asserted axis identification.
H9-T7 equation (H9.24) gives its old form $-\int_Bww'$. Pullback under
$g$ multiplies intersections by $b$, proving (C44.C01.2). Its strict
positivity is precisely the surface input used in H9-T8. No sign input
on $X_C$ has been used.
\end{proof}

\paragraph{Archive and predecessor implications.}
This sharpens the wording of the cover gate for cote 44, Th.~5,
scan 004 and JAC1964 QH9-cover-primitivep. It uses H9-T7, H9-T8,
H9-T18 and finite-flat push-pull; it does not re-prove their height
identification. Cote 16, scan 050, language is algebraic Gysin and
restriction with their correct Lefschetz weights. In Chow motives,
$b^{-1}g^*g_*$ is the old-part projector; the unresolved piece is its
complement inside the primitive middle projector.

\subsection{C44-C02: every nonextremal fixed splitting stratum of triple covers}
\begin{theorem}[C44-C02]\label{thm:maroni}
Fix the abelian pencil of H9-T23 and integers
\[
                  1\le e\le f<2e,
                  \qquad g=e+f-2\ge2.
\]
In every characteristic, let $U_{e,f}$ be the parameter space of smooth
connected triple covers of the fixed target $\mathbf P^1$, unramified
above the pencil's critical set, with framed Tschirnhausen bundle
$\mathcal O(e)\oplus\mathcal O(f)$. It is a nonempty irreducible open
subset of the space of twisted binary cubics. At its geometric generic
point, every algebraic class in the new middle trace kernel of the
pullback pencil is zero. The exact norm-height formula and radical are
\[
 h_L(z,w)=\tfrac13h_K(\operatorname{cor}z,\operatorname{cor}w),
 \qquad \operatorname{rad}h_L=
       \operatorname{cor}^{-1}(\operatorname{rad}h_K).       \tag{C44.C02.1}
\]
For abelian-fourfold pencils the full homological height is therefore
semipositive, conditionally on 032/Milne as in H9-T8. The vanishing
assertion and norm identity themselves do not use 032.

For $d=f-e>0$ the unframed stratum has codimension $d-1$ in the full
fixed-target smooth triple-cover stack. Thus this treats in particular
the first unbalanced stratum for even $g\ge6$ (a divisor), and for odd
$g\ge9$ (codimension two). It does not assert anything for every closed
point of these strata or for the extremal splitting $f=2e$.
\end{theorem}
\begin{proof}
For $E=\mathcal O(e)\oplus\mathcal O(f)$, Wood's cubic-algebra
equivalence, in the exact form already used in H9-T24 equations
(H9.136)--(H9.139), gives the affine coefficient space
\[
 P_E=H^0\bigl(\mathcal O(2e-f)\oplus\mathcal O(e)
                \oplus\mathcal O(f)\oplus\mathcal O(2f-e)\bigr).
\]
All degrees are nonnegative, so $\dim P_E=2g+8$. The smooth,
geometrically connected, and specified unramified conditions define an
open $U_E$. Each point has the required quotient bundle by the cubic
equivalence. It remains to show nonemptiness and the monodromy property.

Set
\[
             a=2e-f-1\ge0,\qquad c=f-e\ge0.
\]
Then $2a+3c=g$, $a+c+1=e$, and $a+2c+1=f$. Apply H9-T23 with
these particular parameters (its proof allows arbitrary nonnegative
$a,c$, not only the balanced choices used later in H9-T24).
In characteristic different from three, the Kummer normalization has
the relative eigensheaves
$\mathcal O\oplus\mathcal L^{-1}\oplus\mathcal L^{-2}(D_2)$,
where $\deg\mathcal L=e$ and $\deg D_2=a+1$. Its quotient dual is
$\mathcal O(e)\oplus\mathcal O(2e-a-1)=E$.
All exponent-two points and the line bundle are fixed over the moving
branch parameter, so this is an isomorphism of relative bundles.
In characteristic three, there are $a+1$ simple poles and $c$ double
poles. The exact normalized basis calculation (H9.144) gives quotient
dual degrees $a+c+1=e$ and $a+2c+1=f$, again with constant pole
divisors. After choosing a framing, in each case this constructs an
actual morphism from the H9-T23 parameter curve into $U_E$.

The universal pullback pencil is smooth proper, and
$1-\frac13Q^*Q_*$ defines a lisse middle-cohomology direct summand,
exactly as in H9-T24 equations (H9.146)--(H9.147). Its restriction to
the constructed curve has no nonzero finite-orbit vectors by H9-T23's
middle-convolution or Fourier calculation. A finite-orbit vector for
the fundamental group of $U_E$ would remain such after restriction to
this curve. Hence none exists on $U_E$. Every algebraic class at the
geometric generic point is defined over a finite extension of the
parameter function field and consequently has a finite monodromy
orbit. Its new projection must vanish. The actual-cycle push-pull and
admissibility argument in H9-T24 then gives (C44.C02.1) and the radical.
This uses precisely the old result, with the newly specified parameter
stratum, rather than a generic conclusion specialized to a closed point.

For codimension, $\dim\operatorname{Aut}(E)=d+3$ if $d>0$ and is
$4$ if $d=0$: compute $H^0(\End E)$ from its four line summands.
Thus the $d>0$ stratum has dimension $2g+5-d$, while the balanced
stratum has dimension $2g+4$. Smooth separable covers have finite
automorphism groups, so the quotient-stack dimension calculation is
valid also in characteristics two and three. The difference is $d-1$.
The first nonbalanced value is $d=2$ for even $g$ and $d=3$ for odd
$g$; $f<2e$ gives exactly the genus bounds stated above.
\end{proof}

\paragraph{Sources and scope.}
Wood, \emph{Parametrizing quartic algebras over an arbitrary base},
Theorem 2.1 and equation (1), printed pp.~1073--1074; JAC1964 H9-T23,
(H9.127)--(H9.132), and H9-T24, (H9.136)--(H9.147). These are the
precise predecessor inputs. The new conclusion concerns geometric
generic points of special Maroni strata, over algebraic closures of
their parameter fields. It is not a proof for all special covers over
$\overline{\F}_p$. Archive target: cote 44, scan 004, Th.~5 and its
fibration margin; predecessor row QH9-cover-primitivep is narrowed in
this specified stratum scope, not closed.

% END work/covers/contribution.tex

% BEGIN work/covers/boundary_strata.tex
% C44_POS /root; candidates undergoing independent whole-proof review.
\subsection{C44-C03: characteristic-two conductor strata with a simple pole}
\begin{theorem}[C44-C03]\label{thm:AS-strata}
In characteristic two, fix an Artin--Schreier double-cover pole partition
\[
 d_i=2a_i+1\quad(1\le i\le s),\qquad
 s\ge2,\quad a_i\ge0,\quad a_1=0,
 \qquad g=\sum_i(a_i+1)-1\ge2.
\]
Fix an abelian-fourfold Lefschetz pencil as in H9-T22, and require the
branch locus of the covers to avoid its critical values. At the geometric
generic point of the indicated irreducible conductor stratum, there is
no algebraic class in the new middle trace kernel of the pulled-back
pencil. Consequently
\[
 h_L(z,w)=\tfrac12h_K(\operatorname{cor}z,\operatorname{cor}w),
 \qquad \operatorname{rad}h_L=
 \operatorname{cor}^{-1}(\operatorname{rad}h_K).
 \tag{C44.C03.1}
\]
The identity and vanishing are unconditional. Semipositivity of the
resulting height uses H9-T8 and is conditional on the unrefereed
032/Milne inputs. The assertion includes the geometric generic point
of the hyperelliptic $2$-rank $g-1$ divisor for every $g\ge2$.
It is a statement at geometric generic points, not at every closed
point of the stratum.
\end{theorem}
\begin{proof}
Order the poles and choose a coordinate for which zero and infinity
are smooth values of the fixed pencil. We construct one witness
family with its distinguished simple pole fixed at $t=0$; this
fixed-position family is not asserted to be dense in the stratum.
Put the other
poles at distinct nonzero $b_2,\ldots,b_s$, away from the critical set
$\Sigma$ of the pencil. Choose a rational function $f_0$, regular at zero
and infinity, with exactly these poles and orders $d_2,\ldots,d_s$.
The family
\[
                  w^2-w=f_0(t)+\lambda/t,
                  \qquad \lambda\in\mathbf G_m
 \tag{C44.C03.2}
\]
lies in the specified conductor stratum. Its normalization is smooth
and proper over the parameter. Here is the required local verification
also for a higher pole. If $z$ is a local parameter at a pole of order
$2a+1$, set $v=z^{a+1}w$. The integral local equation becomes
\[
 v^2-z^{a+1}v=z^{2a+2}f(z).
\]
The right side has first term $c z$, with $c$ a unit on the parameter
space. At the unique point $z=v=0$ its derivative in $z$ is $c\ne0$.
The equation is regular there, finite flat of degree two, and agrees
with the extension off the pole, so it is the normalization. Away from
the poles the derivative in $w$ is a unit. This proves the relative
smoothness and justifies smooth proper base change.

Let $V$ be the variable odd local system of the original pencil,
of rank $r>0$, with $q=\#\Sigma$. Its irreducibility and rank-one tame
transvections are the inputs used in H9-T21--T22. Change coordinates
to $x=1/t$, and put
\[
 G=\mathcal L_\psi(f_0(1/x)),\qquad
 F=(\operatorname{inv}^*V)\otimes G.
\]
Their middle extensions to $\mathbf A^1_x$ are irreducible Fourier
sheaves: $G$ has a finite wild pole, and $F$ is irreducible of rank
$r\ge2$. Both are unramified at infinity. The exact character-part
Leray calculation of H9-T22 identifies the new middle cohomology
with
\[
 H^3(A)\otimes\operatorname{NFT}_\psi(G)
       \ \oplus\ \operatorname{NFT}_\psi(F)
 \tag{C44.C03.3}
\]
over the parameter open. This is a calculation of the actual
smooth-proper cohomology, not a substitution of extension by zero at
the old transvections. The outer Leray terms are deck invariant; the
two displayed terms are its nontrivial deck character.

Set $R=\sum_{i=2}^s(d_i+1)$. The Fourier ranks are $R$ and $q+rR$,
by the finite conductor formula, Katz, \emph{Exponential Sums and
Differential Equations}, Lemma 7.3.9(2). Theorem 7.3.8(3) gives
irreducibility of the transforms. At parameter zero, Corollary
7.4.3.1 and Theorem 7.4.4(2)--(3), equivalently Theorem 7.5.4(5)--(7),
give their respective local representations
\[
 J(1,2)\oplus J(1,1)^{R-2},\qquad
 J(1,2)^{r}\oplus J(1,1)^{q+rR-2r}.
 \tag{C44.C03.4}
\]
Each is infinite because it contains a nontrivial unipotent block.
An irreducible representation with infinite image has no nonzero
finite-orbit vector: the span of the orbit would be a nonzero
subrepresentation, and the normal core of its finite-index stabilizer
would then act trivially on the whole representation. This would make
its image finite. Direct sums and tensoring with a constant vector
space preserve the absence of finite-orbit vectors. Thus (C44.C03.3)
has none.

For completeness, the space with ordered pole positions and all
reduced odd principal-part coefficients is irreducible: it is an open
in a vector bundle over the ordered configuration space, with each
leading coefficient nonzero. The above local normalization works
throughout it. It maps dominantly to the indicated conductor stratum
of covers. Its universal trace-kernel local system restricts to
(C44.C03.3). A finite-orbit vector for the universal monodromy would
remain one on this parameter curve. There are therefore no such
vectors at the geometric generic point of the stratum. Algebraic
classes, defined over finite extensions of its function field, have
finite monodromy orbit. The actual trace projector and admissibility
argument of H9-T22 and H9-T24 now give (C44.C03.1).

The genus and $2$-rank are $g$ and $s-1$. Pries--Zhu,
\emph{The $p$-rank stratification of Artin--Schreier curves},
Theorem 1.1 and Corollary 1.3, printed p.~708, identify the components
in characteristic two with partitions of $g+1$ into $s$ positive
parts $a_i+1$, all of dimension $g-1+(s-1)$. For rank $g-1$ the
partition is $(2,1,\ldots,1)$, with $g-1$ parts equal to one.
It has a simple pole and is a divisor in the hyperelliptic moduli
space of dimension $2g-1$. Fixing the target adds three dimensions
to both spaces and leaves codimension unchanged. This proves the
last assertion.
\end{proof}

\subsection{C44-C04: extremal triple-cover splitting strata}
\begin{theorem}[C44-C04]\label{thm:extremal-maroni}
Let $e\ge2$ and $g=3e-2$. In every characteristic, let $U_{e,2e}$
denote the framed space of smooth connected triple covers of the
fixed target $\mathbf P^1$ with Tschirnhausen bundle
$\mathcal O(e)\oplus\mathcal O(2e)$, unramified over the old
critical set. At its geometric generic point the new middle
trace kernel of an abelian-fourfold pencil contains no algebraic
classes. The norm-height identity and radical in (C44.C02.1) hold.
These assertions are unconditional; height semipositivity is
conditional on the unrefereed 032/Milne inputs through H9-T8.

Together with C44-C02 this settles the geometric generic point
of every admissible triple-cover splitting stratum of genus at
least two. It does not assert the same vanishing at every closed
point over $\overline{\mathbf F}_p$.
\end{theorem}
\begin{proof}
The cubic coefficient space for $E=\mathcal O(e)\oplus\mathcal O(2e)$
is $H^0(\mathcal O\oplus\mathcal O(e)\oplus\mathcal O(2e)
\oplus\mathcal O(3e))$. Its smooth connected unramified locus is
irreducible if nonempty. We construct a family in it for each
characteristic and prove that its new trace-kernel monodromy has
no finite-orbit vectors. The restriction argument of C44-C02
will then prove the assertion for this whole stratum.

\paragraph{Characteristic different from three.}
Choose a coordinate $t$ with zero and infinity outside the old critical
set $\Sigma$. Choose $3e-3$ nonzero points $b_i$ so that their
$\mu_3$-orbits are pairwise disjoint and avoid the $\mu_3$-orbits of
$\Sigma$. Put $f(t)=\prod_{i=1}^{3e-3}(t-b_i)$ and consider
\[
                   u^3=(t^3-s)f(t).
 \tag{C44.C04.1}
\]
Remove zero and the finitely many parameters for which a moving root
hits a fixed branch point or a critical value. The normalized covers
are smooth and have $3e$ simple branch points, all with exponent one;
infinity is unramified. Their genus is $3e-2$, and the relative
eigensheaf decomposition is
\[
        \pi_*\mathcal O_C=
        \mathcal O\oplus\mathcal O(-e)\oplus\mathcal O(-2e).
 \tag{C44.C04.2}
\]
The normalized local equations are the tame equations of H9-T23.
Because the degree of the branch polynomial is $3e$, this is a
relative decomposition with fixed line bundles, not only a
fibrewise numerical calculation. It supplies a framing into $U_{e,2e}$.

Fix a nontrivial cubic character $\chi$. On the $t$ line let
$G_\chi=\mathcal L_\chi(f(t))$ and
$F_\chi=V\otimes G_\chi$, always taking middle extensions.
For $\rho:\mathbf A^1_t\to\mathbf A^1_v$, $v=t^3$, set
\[
                  M_\chi=\rho_*G_\chi,
                  N_\chi=\rho_*F_\chi.
\]
Both are irreducible middle extensions. Indeed, restrict to the
$\mu_3$-stable open on which $\rho$ is a Galois cover. Each input
is irreducible. Its three conjugates are pairwise nonisomorphic:
at a fixed branch point one has nontrivial scalar inertia $\chi$,
whereas either conjugate is lisse, by the choice of the $b_i$.
The elementary induction criterion for a normal subgroup with
pairwise inequivalent conjugates therefore makes its induced
representation irreducible. Finite pushforward is its middle
extension, as is checked by inertia invariants at the omitted points.

Finite pushforward and projection formula identify the two new
character terms with
\[
          H^3(A)\otimes\operatorname{MC}_\chi(M_\chi)
                    \ \oplus\ 
                    \operatorname{MC}_\chi(N_\chi).
 \tag{C44.C04.3}
\]
This follows from the actual H9-T23 character cohomology by replacing
$\mathcal L_\chi(t-s)$ there with
$\mathcal L_\chi(t^3-s)$ and then applying $\rho_*$. The fibre
cohomology is the middle convolution, including its local invariant
conditions. Both input sheaves are tame. At infinity their
representations are respectively
\[
                 1\oplus\chi\oplus\chi^2,
                (1\oplus\chi\oplus\chi^2)^{r}.
 \tag{C44.C04.4}
\]
Here the original $V$ and $G_\chi$ are unramified at infinity,
and the ramification of $t\mapsto t^3$ supplies the regular
representation of $\mu_3$.

Katz, \emph{Rigid Local Systems}, Theorem 3.3.3 and Corollary 3.3.6
(with Corollary 4.3.11 for the characteristic-zero comparison,
as already used in H9-T23),
give irreducibility of the convolutions and their local invariant
quotient formula. In that formula write
$M_\infty/M_\infty^{I_\infty}$ for the input representation.
Each trivial one-dimensional summand in (C44.C04.4) forces a
$J(1,2)$ block in $M_\infty$. Tensoring by the nontrivial $\chi$
turns it into $J(\chi,2)$, and quotienting by inertia invariants
does not remove it. Thus both outputs have infinite local image.
For a rank check, the first input has finite drops $2$ at zero
and $1$ at each of the $3e-3$ fixed branch images; Corollary
3.3.7 gives output rank $3e-2$. The second has finite drops
$q+(3e-1)r$ in total and output rank $q+(3e-2)r$.
Collisions among old critical images do not change their total
drop; the chosen branch images are disjoint from them.
Irreducibility and infinite image exclude finite-orbit vectors,
as in C44-C03.

\paragraph{Characteristic three.}
Choose a coordinate $t$ with zero and infinity outside $\Sigma$.
Choose $e-1$ distinct nonzero points $b_j$, away from $\Sigma$, so
that the reciprocal points $b_j^{-1}$ and their negatives are
distinct and avoid the reciprocal critical set and its negatives.
For nonzero constants $c_j$, put
\[
 f_0(t)=\sum_{j=1}^{e-1}\frac{c_j}{(t-b_j)^2},
 \qquad w^3-w=f_0(t)+s/t^2,\qquad s\in\mathbf G_m.
 \tag{C44.C04.5}
\]
There are $e$ poles, all of order two, hence genus $3e-2$.
The local order-two normalization and the global basis calculation
of H9-T23--T24, (H9.131) and (H9.144), give a smooth family and
the relative splitting (C44.C04.2). Its pole divisor is fixed,
so once again it maps to the framed stratum.

Put $x=1/t$. For either nonzero additive character index
$\alpha\in\mathbf F_3^\times$, set
\[
 G_\alpha=\mathcal L_\psi(\alpha f_0(1/x)),\qquad
 F_\alpha=(\operatorname{inv}^*V)\otimes G_\alpha,
 \qquad \rho:x\longmapsto v=x^2.
\]
The two finite pushforwards $\rho_*G_\alpha$ and
$\rho_*F_\alpha$ are irreducible by the same induction argument,
now for the tame quadratic cover and its involution $x\mapsto-x$.
A finite wild pole at $b_j^{-1}$ has no pole in its negative orbit,
so the two conjugates cannot be isomorphic. Their ranks are $2$
and $2r$, so neither is an exceptional rank-one additive character.
They are Fourier sheaves. At infinity their representations are
$1\oplus\epsilon$ and $(1\oplus\epsilon)^r$, respectively,
where $\epsilon$ is the nontrivial quadratic tame character.

Projection formula converts the variable term
$\mathcal L_\psi(\alpha s x^2)$ into
$\mathcal L_\psi(\alpha s v)$. The actual new character
cohomology is consequently the sum of $H^3(A)$ tensored with
$\operatorname{NFT}_\psi(\rho_*G_\alpha)$ and
$\operatorname{NFT}_\psi(\rho_*F_\alpha)$, pulled back by
$s\mapsto\alpha s$. Katz, \emph{Exponential Sums and Differential
Equations}, Theorem 7.3.8(3), gives irreducibility. Corollary
7.4.3.1 and Theorem 7.4.4(2)--(3) show that the trivial infinity
summands produce $J(1,2)$ blocks at Fourier parameter zero.
The nontrivial quadratic summands remain nontrivial rank-one
blocks and do not cancel these unipotent blocks. Hence these
transforms also have infinite image and no finite-orbit vectors.

In both characteristics every new character has now been handled.
The direct sum has no finite-orbit vectors. The universal lisse
trace-kernel projector on $U_{e,2e}$ restricts to this sum, so
the restriction argument proves generic algebraic vanishing on
the stratum. The norm-height identity and radical follow by
the same actual-cycle push-pull as in C44-C02. All smooth connected
triple-cover splittings have $e\le f\le2e$; the inequality follows
also directly from the cubic equation, since if $f>2e$ its first
coefficient is zero and the cubic has a linear factor, contradicting
smooth connectedness. C44-C02 handles the strict inequality, while
the present result handles equality for all genus at least two.
\end{proof}

\paragraph{Programme location and residual cover question.}
C44-C03--C04 extend the monodromy method of JAC1964 H9-T21--T24 to
specified special strata. Their height meaning is cote 44, Th.~5,
scan 004, with the primitive qualification of scan 008; the
idempotent $1-b^{-1}g^*g_*$ is the new Chow-motive projector.
The argument concerns lisse sheaves and finite monodromy orbits,
using SGA~7~II, XVIII~6.7.2 for the original pencil variation as
in the predecessor proofs. It gives no positivity theorem for an
extra Tate class at a closed point. Such a class belongs to the
remaining H9-T18 primitive incidence gate, after C44-C01 removes
the positive old surface-axis summand.

% END work/covers/boundary_strata.tex

% BEGIN work/covers/as_conductor_strata.tex
\section{Every Artin--Schreier conductor partition in characteristic two}
\label{sec:c44-c05}

\begin{theorem}[C44-C05: generic variable-kernel vanishing on every conductor stratum]
Work over an algebraically closed field \(k\) of characteristic \(2\), in the
abelian-fourfold Lefschetz-pencil setup of C44-C01 and C44-C03.
Write \(f:\widetilde A\to\mathbf P^1\) for the blowup pencil, \(\Sigma\)
for its critical-value set, and \(\mathcal V\) for the geometrically
irreducible variable summand of \(R^3f_*\mathbf Q_\ell\), where
\(\ell\ne2\). At each point of \(\Sigma\) its local monodromy is a
nontrivial tame transvection; at smooth pencil values it is unramified.
If \(\mathcal V=0\), omit the variable monodromy argument below.

Fix odd positive integers \(d_1,\ldots,d_s\), with
\(g=-1+\frac12\sum_i(d_i+1)\ge2\). Let \(\mathcal H_{\mathbf d}\) be the
irreducible Artin--Schreier conductor stratum of connected degree-two
covers \(\pi:C\to\mathbf P^1\) with these pole orders, restricted to
covers branched at smooth pencil values. At its geometric generic point,
put \(X_C=\widetilde A\times_{\mathbf P^1}C\), with incidence embedding
\(k_C:X_C\hookrightarrow A\times C\). A superscript \({\rm new}\)
denotes the anti-invariant summand for the deck involution. Then
\[
 \Alg^2(X_C)\cap
 \bigl(\ker(k_{C*}:H^4(X_C)\to H^6(A\times C))\bigr)^{\rm new}=0.
\]
Consequently the intersection form on the algebraic incidence kernel
is positive definite: only the old surface-axis contribution of
C44-C01 remains. This kernel-sign assertion is unconditional.

If \(s\ge2\), the stronger assertion
\(\Alg^2(X_C)^{\rm new}=0\) holds. The degree-two norm-height and
radical identities of C44-C03 consequently apply on the same domains,
with factor \(1/2\).

If \(s=1\), no vanishing of the ambient new summand is asserted and
no degree-two norm-height identity is inferred. The height
semipositivity conclusion on the Albanese-primitive domain follows
from the preceding kernel assertion and H9-T18 only
\emph{conditionally on the unrefereed preprints 032/Milne2026}.
\end{theorem}

\begin{proof}
Choose a smooth pencil value as infinity. Put \(d=d_1\).
When \(\mathcal V\ne0\), choose a distinguished branch point
\(a\notin\Sigma\) such that
\[
 \Sigma_a=\{(\sigma-a)^{-1}:\sigma\in\Sigma\}\subset\mathbf G_m
\]
has trivial stabilizer under \(\mu_d\). Such \(a\) exist: an equality
\(\zeta\Sigma_a=\Sigma_a\), with \(\zeta\ne1\), induces the homothety
\(t\mapsto a+\zeta^{-1}(t-a)\) permuting \(\Sigma\).
For each \(\zeta\) and permutation, the equation for one point
determines \(a\), so only finitely many values are excluded.
If \(\mathcal V=0\), choose any smooth \(a\) and omit this condition.
Set \(x=(t-a)^{-1}\) and \(\rho(x)=x^d\); this map is tame since \(d\)
is odd. Choose the other branch points so their \(x\)-coordinates
have pairwise disjoint \(\mu_d\)-orbits, disjoint also from
\(\mu_d\Sigma_a\).

Choose \(f_0(t)\) with the prescribed odd pole orders
\(d_2,\ldots,d_s\) at the other branch points, regular at \(a\)
and infinity. For \(s=1\), take \(f_0=0\). Consider
\[
       w^2-w=f_0(t)+\frac{\lambda}{(t-a)^d},
       \qquad \lambda\in\mathbf G_m.                    \tag{C05.1}
\]
Every fibre has exactly the required partition. Its smooth projective
normalizations form a smooth family, with the usual restriction to
the parameter open needed for the cohomological construction.
Indeed at a pole of order \(2b+1\), write the right side as
\(c(t)t^{-(2b+1)}\), \(c(0)\ne0\), and put \(v=t^{b+1}w\).
The equation
\(v^2-t^{b+1}v=t c(t)\) is smooth at the ramified point because its
\(t\)-derivative there is \(-c(0)\ne0\). This is uniform while
the leading coefficient is a unit. Away from the poles the equation
is etale. Branching avoids \(\Sigma\), so the incidence fourfolds
obtained by base change are smooth too.

Put \(\mathcal G=\mathcal L_\psi(f_0(a+1/x))\) and
\(\mathcal F=\mathcal V_a\otimes\mathcal G\), where
\(\mathcal V_a\) is the transformed pencil local system.
Both are unramified at \(x=\infty\).
Take middle extensions and set
\(\mathcal N=\rho_*\mathcal F\) on the affine \(y\)-line.
The restriction of \(\mathcal F\) to its lisse open is irreducible.
Its different deck translates are nonisomorphic: the set where the
local monodromy has the pencil's nontrivial tame unipotent part is
exactly \(\Sigma_a\), and no nontrivial deck transformation preserves
it. The scalar wild ramification at the separate other poles cannot
remove or imitate those tame transvections.
Mackey's criterion for induction along the Galois cover
\(x\mapsto x^d\) therefore gives geometric irreducibility of
\(\mathcal N\); for \(d=1\) this is immediate.

At \(y=\infty\), the local representation of \(\mathcal N\) is the
regular representation of \(\mu_d\), repeated
\(r=\operatorname{rank}\mathcal V\) times. Its \(r\) trivial tame
blocks of length one become \(r\) unipotent blocks of length two at
Fourier parameter \(\lambda=0\). The Fourier middle extension stays
irreducible. Its input has rank \(dr\ge2\), since the nonzero variable
odd cohomology is symplectic, so it is not an exceptional punctual
transform. The displayed unipotent blocks also witness its nonzero
lisse part and infinite geometric monodromy.

The Fourier cohomology is exactly
\[
 H^1\!\left(\mathbf P^1,
   j_*(\mathcal V\otimes
      \mathcal L_\psi(f_0(t)+\lambda/(t-a)^d))\right).     \tag{C05.2}
\]
Finite pushforward by \(\rho\), followed by tensoring with
\(\mathcal L_\psi(\lambda y)\), gives the Fourier integral.
For \(\lambda\ne0\) the sheaf is wild at \(x=\infty\), with no
boundary invariants there, identifying the affine Fourier cohomology
with projective middle-extension cohomology. Finite pushforward
preserves the middle extensions at the other singular points.

An irreducible characteristic-zero representation of infinite image
has no nonzero finite-orbit vector. Otherwise the invariants of the
core of a finite-index vector stabilizer would be a nonzero invariant
subspace, hence the entire representation, forcing finite image.
Thus the variable family (C05.2) has no such vector.

If \(s\ge2\), apply the same argument to
\(\mathcal M=\rho_*\mathcal G\). Its wild pole set is nonempty
and its chosen disjoint \(\mu_d\)-orbits distinguish all deck
translates, so \(\mathcal M\) is irreducible. Its infinity inertia
is the regular representation of \(\mu_d\); Fourier therefore has
a nontrivial length-two unipotent block at zero.
For \(d=1\), the rank-one input has a finite wild pole and is
unramified at infinity, excluding both the constant sheaf and the
linear Artin--Schreier exceptional input. Consequently
\[
 H^1\!\left(\mathbf P^1,
 j_*\mathcal L_\psi(f_0(t)+\lambda/(t-a)^d)\right)
       =H^1(C_\lambda)^{\rm new}                         \tag{C05.3}
\]
has no nonzero finite-orbit vector when \(s\ge2\).
Tensoring with the geometrically constant \(H^3(A)\) preserves
this assertion.

The new degree-four Leray cohomology splits as
\[
 H^4(X_C)^{\rm new}
 =\bigl(H^3(A)\otimes H^1(C)^{\rm new}\bigr)
 \oplus H^1(\mathbf P^1,j_*(\mathcal V\otimes\mathcal L_\pi)).
                                                            \tag{C05.4}
\]
The fibrewise projector
\(\operatorname{res}\circ L_h^{-1}\circ\operatorname{Gysin}\)
canonically splits \(R^3\) into the constant \(H^3(A)\) and
\(\mathcal V\) pieces. The outer Leray terms are deck-invariant,
so only \(E_2^{1,3}\) occurs in the nontrivial deck character.
Thus (C05.4) uses no arbitrary choice of Leray splitting.
Its first summand is \(k_C^*H^4(A\times C)^{\rm new}\).
The incidence divisor has ample class
\(p_A^*h+p_C^*\pi^*c_1\mathcal O_{\mathbf P^1}(1)\).
Hence \(k_{C*}k_C^*:H^4(A\times C)\to H^6(A\times C)\)
is its hard-Lefschetz isomorphism. On the new ambient part it is
\(L_h:H^3(A)\to H^5(A)\) tensored with \(H^1(C)^{\rm new}\),
so is injective. The second summand is killed by Gysin:
the fibrewise Gysin kills variable \(H^3\), and there are no
other new Leray terms to receive its image. It is therefore
exactly \((\ker k_{C*})^{\rm new}\).

The partition stratum is irreducible. Work if necessary on its
normal-form coefficient parameter space with ordered distinct
branch points, which dominates the stratum and carries the family.
The local systems restrict along (C05.1) to those just considered.
A finite-orbit vector for the full parameter monodromy would retain
a finite orbit on this witness curve, so none exists in the variable
system, or, when \(s\ge2\), in the full new degree-four system.
A geometric-generic algebraic class is defined after a finite
extension of the parameter function field. It thus has finite orbit
under its absolute Galois group, which surjects onto the fundamental
group of the smooth coefficient space. Purely inseparable extensions
do not affect this argument. Its projection to these summands must
vanish, proving the two vanishing assertions.

C44-C01 identifies the old incidence kernel with the surface-axis
divisor contribution, whose required sign is the Hodge index theorem.
Explicitly, if \(q:X_C\to\widetilde A\) is the degree-two map,
\(z_{\rm old}=\frac12q^*q_*z\). Proper functoriality keeps \(q_*z\)
in the old incidence kernel, and pullback multiplies the form by
the positive degree \(2\).

For \(s=1\), the equation \(w^2-w=\lambda x^d\) becomes a constant
curve after adjoining a \(d\)-th root of \(\lambda\) and rescaling
\(x\). The scalar system can thus have finite monodromy and its
ambient contribution cannot be discarded. This leaves the kernel
identification unaffected, but prevents full-new vanishing and
the associated norm identity. H9-T18 and the kernel sign yield
the stated Albanese-primitive height conclusion only under the
explicitly stated unrefereed 032/Milne2026 hypotheses.
\end{proof}

\paragraph{Source and scope audit.}
The Fourier inputs are Katz, \emph{Exponential Sums and Differential
Equations}, 7.3.2, 7.3.8(3), 7.4.3.1, and 7.4.4(2)--(3),
equivalently 7.5.4(7); they impose no exclusion of characteristic two.
Conductor-stratum irreducibility, including the one-pole,
\(p\)-rank-zero case, is Pries--Zhu, Proposition 3.6, Remark 3.9
and Theorem 3.10, printed pp.~716--717.
These are geometric-generic statements on every stratum, not
assertions at every closed finite-field point.
For cote 44, scans 004 and 008, this extends the H9-T18 primitive
height criterion and the H9-T22 monodromy calculation. It narrows
QH9-cover-primitivep to further special points inside these strata,
without closing that original pointwise question.

% END work/covers/as_conductor_strata.tex

% BEGIN work/covers/separable_cover_lifts.tex
% C44_POS /root/axis_independent; GPT-6 Astra Pro; effort not exposed.
% Root's separable-cover deformation idea and incidence determinant extension.
\section{Lifting separable covers and a fixed-incidence sign criterion}
\subsection*{C44-C06. Every separable map of a smooth curve to the line lifts}

\textbf{Theorem.}
Let $k$ be a perfect field of positive characteristic, let $C/k$
be a smooth projective geometrically connected curve, and let
$f:C\to\mathbf P^1_k$ be a finite separable morphism of degree $d$.
There is a smooth projective curve $\mathcal C/W(k)$ and a finite
flat morphism
\[
                   \widetilde f:\mathcal C\longrightarrow
                                    \mathbf P^1_{W(k)}
                                                               \tag{C06.1}
\]
of degree $d$ reducing to the given map.
No Galois action or branch configuration is required to lift.

For the incidence consequence assume $k$ is finite.
Let $A/k$ be an abelian fourfold, let $L$ be ample,
and let $s_0,s_1\in H^0(A,L)$ span a pencil.
Write $M=f^*\mathcal O_{\mathbf P^1}(1)$ and let
$u,v\in H^0(C,M)$ be the pullbacks of the two homogeneous
coordinates. Suppose the incidence fourfold
\[
 Y=\{s_0\otimes u+s_1\otimes v=0\}\subset A\times C
                                                               \tag{C06.2}
\]
is smooth. After a finite residue extension,
there is an embedded smooth projective mixed-characteristic
lift of $Y\subset A\times C$, including its defining ample
divisor class. Thus every smooth incidence fourfold of this
form has a geometric lift, including those formed with wild
separable covers.

\textbf{Proof.}
Set $W_n=W(k)/p^n$. We construct compatible smooth proper
curves $C_n/W_n$ and maps $f_n:C_n\to\mathbf P^1_{W_n}$.
The initial pair is $(C,f)$. Suppose it is constructed over
$W_n$, and put $I=p^nW_{n+1}$, a one-dimensional $k$-space.
The small extension $W_{n+1}\to W_n$ has square-zero kernel $I$.

The obstruction to lifting the smooth curve $C_n$ lies in
$H^2(C,T_C)\otimes_k I$, which is zero. Choose one smooth lift
$C'_{n+1}$. It is proper because properness is unchanged under
a nilpotent thickening. The isomorphism classes of such lifts,
with their specified reduction, form a torsor under
$H^1(C,T_C)\otimes_k I$. These are SGA~1, Expos\'e III,
6.9--6.10; the underlying-curve algebraization appears in
III 7.2--7.4. We will instead algebraize the finite maps
directly after constructing them.

On affine opens of $C'_{n+1}$ the map $f_n$ lifts because the
target $\mathbf P^1_{W_{n+1}}$ is smooth. The differences of
local lifts on overlaps form a Cech cocycle, with obstruction
class
\[
                       \beta\in H^1(C,f^*T_{\mathbf P^1})\otimes I.
                                                               \tag{C06.3}
\]
This is the morphism-lifting obstruction of SGA~1 III 5.1--5.6.
Changing the lift of $C_n$ by a class
$\alpha\in H^1(C,T_C)\otimes I$ changes $\beta$ by
$df(\alpha)$, up to the harmless choice of sign for the
Cech cocycles. Indeed a gluing automorphism $1+\alpha_{ij}$
changes the difference of the two local map lifts by
composition with the differential $df$.

Since $f$ is separable, the nonzero map of line bundles
\[
       0\longrightarrow T_C\xrightarrow{df}f^*T_{\mathbf P^1}
                         \longrightarrow N_f\longrightarrow0
                                                               \tag{C06.4}
\]
is injective and its cokernel $N_f$ is supported at finitely
many ramification points. This includes wild ramification;
the torsion sheaf need not be reduced.
Hence $H^1(C,N_f)=0$, and the long exact sequence makes
\[
              H^1(C,T_C)\longrightarrow H^1(C,f^*T_{\mathbf P^1})
\]
surjective. Choose $\alpha$ to cancel (C06.3), modify the
source lift by that class, and glue the local morphisms.
This constructs $(C_{n+1},f_{n+1})$ with the required reduction.
Induction produces the compatible system for every $n$.

Each $f_n$ is finite: it is proper, and its underlying
topological map is the finite special-fibre map, so it is
quasi-finite. It is flat as well. The special map between
smooth curves is finite flat. Inductively, for the first-order
thickenings $C_n\subset C_{n+1}$ and
$\mathbf P^1_{W_n}\subset\mathbf P^1_{W_{n+1}}$, flatness
over $W_{n+1}$ identifies their ideals with their respective
special-fibre structure sheaves tensored with $I$.
The natural conormal map for $f_{n+1}$ is therefore an
isomorphism. Stacks, Lemma 37.10.1, Tag 063Y, proves its
flatness. Thus $(f_n)_*\mathcal O_{C_n}$ is locally free of
rank $d$ on $\mathbf P^1_{W_n}$. No miracle-flatness argument
over the nonregular Artin base is being used.

Stacks, Lemma 30.28.2, Tag 09ZT, algebraizes the compatible
finite maps to a finite morphism
$\widetilde f:\mathcal C\to\mathbf P^1_{W(k)}$.
Explicitly, Grothendieck existence gives the coherent sheaf
from the compatible finite locally free algebras
$(f_n)_*\mathcal O_{C_n}$; full faithfulness algebraizes
their multiplication and unit, and taking the relative
spectrum gives the map. Its completed local modules along
the special fibre are flat of rank $d$, so the algebraized
module is locally free there. The non-locally-free locus
is closed and proper over $W(k)$ and misses the special
fibre; it must therefore be empty. Thus the map is finite
flat of degree $d$ and $\mathcal C$ is flat over $W(k)$.

The special fibre of $\mathcal C$ is smooth. The nonsmooth
locus of $\mathcal C/W(k)$ is closed, and is proper over
the base because $\mathcal C$ is finite over the projective
line. As it misses the special fibre it is empty.
The source is projective: the pullback of
$\mathcal O_{\mathbf P^1}(1)$ under a finite morphism is
ample. This proves (C06.1).

For (C06.2), use the polarized abelian lift C44-A05.
After completing its discrete valuation ring $R$ and,
if necessary, enlarging the perfect residue field, the
Cohen subring $W(k)\subset R$ allows base change of
$(\mathcal C,\widetilde f)$ to $R$.
The ample line bundle $L$ extends to $\mathcal L$ on
$\mathcal A/R$. Ample line bundles on abelian varieties
have no higher cohomology, so cohomology and base change
lift both $s_0,s_1$ to sections $\widetilde s_0,\widetilde s_1$.
The cover lift supplies
$\mathcal M=\widetilde f^*\mathcal O(1)$ and sections
$\widetilde u,\widetilde v$.
Their incidence section
\[
                 \widetilde s_0\otimes\widetilde u+
                 \widetilde s_1\otimes\widetilde v
\]
defines a relative Cartier divisor $\mathcal Y$ in
$\mathcal A\times_R\mathcal C_R$. Its special section is a
nonzerodivisor, so the divisor is flat over $R$.
Its special fibre is the prescribed smooth $Y$.
The same properness argument for the nonsmooth locus
makes $\mathcal Y$ smooth over $R$. It is projective and
its inclusion in the lifted product is the required
embedded lift. \hfill$\square$

\paragraph{Scope of C06.}
This lifts the morphism as a morphism of curves.
It does not preserve a deck group, the branch divisor,
ramification multiplicities, or a chosen Hurwitz stratum.
In particular it does not assert a solution to the
Galois-cover lifting problem. Nor does it lift an
individual middle cycle on the incidence fourfold.

\subsection*{C44-C06a. Removing the separability restriction}

\textbf{Corollary.}
Every finite nonconstant morphism $f:C\to\mathbf P^1_k$
from a smooth projective curve over a perfect field of
characteristic $p$ has a finite flat lift of the same
degree over $W(k)$ with smooth projective source.
The embedded geometric lifting conclusion of C06
therefore holds for every smooth incidence (C06.2)
over a finite field.

\textbf{Proof.}
Let $t$ be the rational coordinate on the line and put
$a=f^*t\in k(C)$. Because $k$ is perfect, the kernel
of $d:k(C)\to\Omega^1_{k(C)/k}$ is $k(C)^p$.
This is Stacks, Lemma 10.158.2, Tag 031W.
Choose the largest $e\ge0$ for which
$a=b^{p^e}$ with $b\in k(C)$. Such a largest exponent
exists: a nonzero pole order of the nonconstant $a$
is divisible by $p^e$ for every such exponent and
therefore bounds it. Maximality gives $db\ne0$.
Thus $b$ defines a finite separable morphism
$h:C\to\mathbf P^1_k$, and
\[
 f=F_e\circ h,\qquad
 F_e:[x:y]\longmapsto[x^{p^e}:y^{p^e}].                      \tag{C06a.1}
\]
Here $F_e$ is the displayed $k$-morphism, acting as
the identity on constants; it is not being confused
with absolute Frobenius on a non-prime field.
The separating-function criterion follows equivalently
from $\Omega^1_{k(C)/k(b)}=0$ for this finite extension.

Lift $h$ by C06 and compose with the same homogeneous
power formula on $\mathbf P^1_{W(k)}$. That formula
is finite flat of degree $p^e$: on each standard
chart its coordinate algebra is free with basis
$1,t,\ldots,t^{p^e-1}$ over the target chart.
The composite is therefore finite flat of degree
$p^e\deg h=\deg f$, and it has the prescribed reduction.
Its source is the smooth projective lift supplied
by C06. The incidence argument of C06 then applies
unchanged. \hfill$\square$

\subsection*{C44-C07. A known positive span and one residual potential Tate direction}

\textbf{Theorem.}
Assume the data of (C06.2) are defined over a finite field,
with $Y$ smooth, and let $k:Y\hookrightarrow T=A\times C$.
All cohomology in this statement is geometric.
For every coefficient prime $\ell\ne p$ put
\[
 V_\ell=\ker\bigl(k_*:H^4(Y,\mathbf Q_\ell(2))
                         \to H^6(T,\mathbf Q_\ell(3))\bigr),
 \qquad q_Y(x,y)=\int_Yxy.                                  \tag{C07.1}
\]
Fix one coefficient prime for
$A_Y=\operatorname{Alg}^2(Y)\cap V_\ell$.
Let $K\subset A_Y$ be a finite-dimensional rational algebraic
subspace on which $q_Y$ is positive definite; $K=0$ is allowed
with the empty determinant equal to one.
After applying the common algebraic projector described
below to cycle representatives, extend the finite field
so that these representatives are defined.
Let $\tau$ be the total algebraic multiplicity of roots
of unity in the normalized Frobenius polynomial of
$K_\ell^\perp\cap V_\ell$.
If $\tau\le1$, then $q_Y$ is positive definite on all of
$A_Y$. More precisely,
\[
 A_Y=K\quad\hbox{or}\quad A_Y=K\perp\mathbf Qz,\qquad q_Y(z,z)>0,
                                                               \tag{C07.2}
\]
and only the first case occurs if $\tau=0$.
The known span $K$ need not lift.
No Tate conjecture, Frobenius semisimplicity, standard
conjecture $D(Y)$, or 032 input is used.

Writing $g=g(C)$, $d=\deg f$ and $D=\int_A h^4$,
the full variable rank is
\[
                  \dim V_\ell=D(5d+2g-2)-14-28g,
                                                               \tag{C07.3}
\]
which is even. Therefore the test with $K=0$ only detects
vanishing of $A_Y$; a residual multiplicity-one test can
be nontrivial only when $\dim K$ is odd.

\textbf{Proof.}
Choose a rational zero-cycle of degree one on $C$, after
a finite extension if desired, and let $\eta$ denote its
cohomology class. The incidence class is
\[
                          a=[Y]=h+d\eta.
                                                               \tag{C07.4}
\]
It is ample, since $L$ and $f^*\mathcal O(1)$ are ample.
Weak and hard Lefschetz give the orthogonal decomposition
\[
                    H^4(Y)=k^*H^4(T)\perp V_\ell,
 \qquad k_*k^*=L_a:H^4(T)\xrightarrow{\sim}H^6(T).
                                                               \tag{C07.5}
\]
We give one inverse correspondence valid for every
coefficient prime, rather than assuming kernel membership
is independent of the realization.

Use the algebraic Chow--Kunneth projectors of the abelian
fourfold and of the curve. On the curve the latter are
$e\times C$, $C\times e$, and their complement in the
diagonal, where $e$ is the rational degree-one zero-cycle.
Kunnemann's motivic hard Lefschetz theorem gives common
algebraic inverse operators
\[
 J=(L_h^2:H^2(A)\to H^6(A))^{-1},\qquad
 J_3=(L_h:H^3(A)\to H^5(A))^{-1}.                            \tag{C07.6}
\]
One may replace $L$ by $L\otimes[-1]^*L$ to obtain a
symmetric ample line bundle and then rescale the operators
by the appropriate rational powers of two.
The exact source is Kunnemann, \emph{Invent.\ Math.}
113 (1993), Theorems 5.1--5.2, printed p.~96.
Only his absolute Chow-motive theorem is used.

In the Kunneth decomposition, an element of $H^6(T)$
is written as $(x,y,z\eta)$ with
\[
 x\in H^6(A),\quad y\in H^5(A)\otimes H^1(C),\quad z\in H^4(A).
 \qquad b=Jx,\quad c=J(hz).
 \]
Define its inverse image in $H^4(T)$ by
\[
 \left(hb+d^{-1}(z-hc),\
                 (J_3\otimes1)y,\
                 (c-db)\eta\right).                        \tag{C07.7}
\]
Since $h(z-hc)=hz-h^2J(hz)=0$, multiplication by
$h+d\eta$ sends the first component to $x$, the odd
component to $y$, and the last total component to
\[
              d\big(hb+d^{-1}(z-hc)\big)+h(c-db)=z.
 \]
This proves the inverse formula. Every operation in it
is an algebraic correspondence, including the product
Kunneth projectors, multiplication by $h$, and the rational
factor $d^{-1}$. Thus one rational correspondence $\Gamma$
on $T$ realizes $L_a^{-1}$ simultaneously for all $\ell$.
The correspondence
\[
                       P=\Delta_Y-k^*\Gamma k_*
                                                               \tag{C07.8}
\]
is the common cohomological orthogonal projector onto
$V_\ell$. Applying it to chosen cycle representatives
preserves their initially defined classes in $A_Y$ and
their rational intersection matrix, and puts them in
$V_\ell$ in every realization. Its nonzero determinant
makes their span independent in every realization.

The embedded lift from C06 and C06a provides one rational Betti
quadratic space for $(V_\ell,q_Y)$ at all $\ell\ne p$.
The Betti kernel is primitive for $a|_Y$: the Gysin map
above the middle degree is injective by dual weak
Lefschetz, so $k_*v=0$ implies $a|_Y\,v=0$.
Hodge--Riemann makes the ordinary Gram determinant
$D_V$ of this rational Betti space positive. Indeed
type $(2,2)$ is positive, and every negative block
comes in a complex-conjugate pair of even real dimension.
This argument only lifts the geometry, not the cycles
spanning $K$.

Here the determinant proof of C44-A15 applies with
the common projector (C07.8). We spell out its sign
step to make the implication precise. Write $G$ for
the positive rational Gram matrix of $K$.
Given any algebraic class, apply $P$ and subtract its
$K$-projection using $G^{-1}$ and rational intersection
numbers. This produces one rational cycle $Z$ in the
orthogonal residual space for every $\ell$.
Algebraic cycles have finite normalized Frobenius orbit.
If $\tau=0$ all residual classes are zero.
If $\tau=1$ and a residual class is nonzero, its potential
Tate line is simple and nondegenerate; its nonzero rational
square $a_Z=q_Y(Z,Z)$ makes the same cycle nonzero in
every realization. It spans that line in each.
Here rationality makes the unique cyclotomic root equal
to $1$ or $-1$; simplicity and orthogonality of the other
generalized eigenspaces prove nondegeneracy without
semisimplicity.

After squaring normalized Frobenius $F$, put $S=F^2$.
On the orthogonal complement $U_\ell$ of that line
there are no eigenvalues $1$ or $-1$.
The operator $S-S^{-1}$ is invertible and skew-adjoint,
so its product with the Gram matrix is alternating.
Its determinant is a square. If
$\psi(T)=\det(T-S|U_\ell)$, reciprocal pairing and purity
give the common rational positive number
\[
              \delta=\det(S-S^{-1})=\psi(1)\psi(-1)>0.
 \]
Both factors are products of positive complex-conjugate
pairs; the empty product is one. Consequently
\[
                  a_Z\det(G)\delta/D_V
                           \in(\mathbf Q_\ell^*)^2
                  \quad\hbox{for every }\ell\ne p.
 \]
A rational number that is a square in all these local
fields is a rational square: otherwise its nontrivial
quadratic extension has an inert prime different from
$p$. Thus $a_Z>0$. Every other algebraic residual class
is a rational multiple of $Z$, its coefficient being
$q_Y(Z',Z)/a_Z$. This proves (C07.2).
The rationality, purity and coefficient independence
of the Frobenius polynomials are Deligne, \emph{Weil I},
Theorem 1.6, printed p.~276; quotienting the middle
polynomial by the ambient factor and $(T-1)^{\dim K}$
gives the residual polynomial. None of these steps
requires semisimplicity or realizes an algebraic class
from a Frobenius root.

For the rank formula, $c(T_{A\times C})=1+(2-2g)\eta$.
The normal sequence for the divisor (C07.4) yields
\[
 \chi_{\rm top}(Y)
     =\int_{A\times C}\bigl(a^5-(2-2g)\eta a^4\bigr)
     =D(5d+2g-2).
 \]
Weak Lefschetz and Kunneth give
\[
 b_1(Y)=8+2g,\qquad b_2(Y)=29+16g,\qquad b_3(Y)=64+56g.
 \]
Poincare duality then gives
$b_4(Y)=\chi_{\rm top}(Y)+84+84g$.
The ambient middle rank is
$b_4(A\times C)=98+112g$, proving (C07.3).
Abelian Riemann--Roch gives $D=24\chi(A,L)$, so that
rank is even. The non-root-of-unity eigenvalues of any
orthogonal transformation occur in reciprocal pairs.
Hence residual cyclotomic multiplicity has the parity
of $\dim V_\ell-\dim K$. This proves the stated parity
qualification. \hfill$\square$

\paragraph{Primary sources and programme scope.}
The deformation inputs were read directly in SGA~1,
Expos\'e III 5.1--5.6, 6.9--6.10 and 7.2--7.4:
\texttt{https://arxiv.org/pdf/math/0206203}.
These are revised pp.~59--62, 67--68 and 70,
respectively (original pp.~71--75, 82--83 and 85--86).
The finite-map algebraization and conormal flatness
criteria are
\texttt{https://stacks.math.columbia.edu/tag/09ZT}
and \texttt{https://stacks.math.columbia.edu/tag/063Y}.
The separating-function input for C06a is
\texttt{https://stacks.math.columbia.edu/tag/031W}.
For (C07.6), the primary scanned source is
\texttt{https://gdz.sub.uni-goettingen.de/download/pdf/PPN356556735\_0113/LOG\_0011.pdf}.
The local-to-global squareclass implication is also
explicit in Ancona--Marmora, \emph{Comment.\ Math.\ Helv.}
100 (2025), proof of Corollary 4.4, printed p.~629;
only that elementary implication is used.

C06 separates geometric lifting from cycle lifting
for the smooth incidences of JAC1964 H9-T18.
C07 supplies a sufficient fixed-cover sign test on
exactly its algebraic kernel, hence a new case of
cote~44, Theorem~5, scan~004 after imposing the
Albanese-primitive domain of H9-T17 and the primitive
qualification on scan~008. It does not settle all
special covers, does not assert that $\tau\le1$
always holds, and makes no claim about singular
incidences requiring a separate resolution.
The kernel sign in C07 is unconditional. Any full-height
conclusion that also uses the separate abelian primitive
terms remains conditional on the unrefereed 032 and
Milne2026 inputs specified in the prompt.

% END work/covers/separable_cover_lifts.tex

% BEGIN work/arithmetic/contribution.tex
% C44_POS, arithmetic continuation, /root/fermat_aj.
% GPT-6 Astra Pro; effort setting not exposed. 2026-11-22 (Europe/Berlin).
% New proofs below are subject to the recorded whole-proof check.

\section{The Fermat regulator: an integral projector and the exact
unramified obstruction}

The notation in this section is that of JAC1964 JAC-P13, JAC-P17,
AR-16a and AR-23.  In particular, $J/\mathbf Q$ is the Jacobian of
$x^4+y^4=z^4$, $N=h^3(J)^G$ with $G=\mu_4^2$, and
$\eta\in V:=\operatorname{CH}^2(N)_{\mathbf Q}$ is the specified
Fourier--Ceresa class.  The nonzero vector
$\alpha_5(\eta)$ spans
$S_5=H^1_f(\mathbf Q,H^3(N_{\overline{\mathbf Q}},\mathbf Q_5)(2))$
by the carried input AR-23, using AR-19 and H9-T15.  We do not reprove
that Selmer calculation, nor its nonzero cycle.  None of the deductions
in this section uses 032 or Milne's report.

The question $\ker\alpha_5=0$ remains unresolved.  The findings here
locate its integral obstruction more precisely.  They do not identify
a Selmer dimension with a Chow dimension.

\subsection{An integral model at the prime 5}

\begin{proposition}[C44-F01]\label{C44-F01}
Let $R=\mathbf Z_{(5)}$ and $S=\operatorname{Spec}\mathbf Z[1/10]$.
Let $\mathcal J/S$ be the Jacobian of the smooth Fermat curve
$\mathcal C:x^4+y^4=z^4$.  There is an exact idempotent, with
coefficients in $\mathbf Z[1/2]\subset R$, given by
\[
 e=\frac{1-[-1]^*}{2}\,\frac1{16}\sum_{g\in G}g^*.
 \tag{F01.1}
\]
The sum is formed after $S'=\operatorname{Spec}\mathbf Z[i,1/10]$
and descended to $S$.  Its rational generic-fibre motive is exactly
$N$.  Put
\[
 A=e\operatorname{CH}^2(J;R).
\]
Then restriction is an isomorphism
\[
 e\operatorname{CH}^2(\mathcal J;R)
       \xrightarrow{\ \sim\ } A,
 \qquad A\otimes_R\mathbf Q=V.                 \tag{F01.2}
\]
Moreover $e$ acts on geometric $5$-adic cohomology only in degree
three, where its image is a free rank-two direct summand.  Thus
\[
 T=eH^3_{\mathrm{et}}(J_{\overline{\mathbf Q}},\mathbf Z_5)(2)
 \tag{F01.3}
\]
is an actual $\mathbf Z_5$ lattice for the representation in the task,
and $eH^b_{\mathrm{et}}(J_{\overline{\mathbf Q}},\mathbf Z/5^n(2))=0$
for $b\ne3$ and every $n\ge1$.
\end{proposition}

\begin{proof}
The Fermat equation is smooth over $\mathbf Z[1/2]$: a singular
projective geometric point would have $x=y=z=0$, since all three
partial derivatives are units times the corresponding cubes.
The section $(1:0:1)$ gives its relative Jacobian, an abelian scheme.
The group action and inversion extend over $S$.  The two averaging
operators in (F01.1) commute and are idempotent.  They use only
powers of two as denominators.  Galois conjugation on $S'/S$ permutes
the sixteen group elements.  Restriction and norm, whose composite
is multiplication by two, descend the idempotent with the asserted
coefficients.  This descent is valid for Chow correspondences and
also for their actions on the cohomology theories used below.

The six characters of $G$ on $H^1(J)$ are
\[
 (1,1),(1,2),(2,1),(3,3),(3,2),(2,3)\quad\text{in }(\mathbf Z/4)^2;
 \tag{F01.4}
\]
these are the carried calculation in JAC-P17.  In their exterior
powers, an odd-degree invariant occurs only in degree three.  The
two invariant triples are the first three and the last three
characters.  Their types are $(3,0)$ and $(0,3)$.  The exact rational
identification with the projector of $N$ is JAC-P13/JAC-P17; in
particular (F01.1) is an integral-at-5 representative of the same
rational summand, not a new motive inferred solely from its
realization.  More explicitly, averaging is zero on $h^1(J)$
because it is induced by the zero norm endomorphism
$\sum_{g\in G}g$.  It is zero on $h^5(J)$ by the duality
identification through the $G$-invariant principal polarization.
Inversion removes the even Deninger--Murre pieces.  Thus the
rational idempotent on the full motive equals $\pi_3p_G$.
The geometric cohomology of an abelian variety is the
exterior algebra of its free degree-one cohomology.  Because the
projector has denominators prime to five, its integral image is a
direct summand and has no torsion.  The character calculation
therefore proves the assertions with both $\mathbf Z_5$ and
$\mathbf Z/5^n$ coefficients.

We next verify that $e$ annihilates the vertical terms in Chow
localization.  On any fibre of $\mathcal J$, the operator
$(1-[-1]^*)/2$ sends divisor classes into
$\operatorname{Pic}^0\otimes R$.  The endomorphism
$\sum_{g\in G}g$ of $\mathcal J_{S'}$ is zero.  Indeed its generic
fibre has zero tangent map by (F01.4), so is zero, and equality of
endomorphisms of an abelian scheme is detected on its dense generic
fibre.  Its dual is also zero, which says that
$\sum g^*$ annihilates $\operatorname{Pic}^0$.  Consequently
\[
 e\operatorname{CH}^1(\mathcal J_s;R)=0
 \quad\text{for every closed }s\in S.            \tag{F01.5}
\]
When a fibre does not contain $i$, check this after the extension
of degree at most two and descend using restriction--norm.

Chow localization gives
\[
 \bigoplus_{s\in S^{(1)}}\operatorname{CH}^1(\mathcal J_s;R)
 \longrightarrow\operatorname{CH}^2(\mathcal J;R)
 \longrightarrow\operatorname{CH}^2(J;R)\longrightarrow0.
 \tag{F01.6}
\]
The action of every automorphism in (F01.1) commutes with these
maps.  Applying the exact direct-summand functor $e$ and (F01.5)
proves (F01.2).  Notice that this argument is integral at five;
it does not invoke rational vanishing of the whole Chow group of
a special fibre.
\end{proof}

\paragraph{Locators and scope.}
The localization sequence is Fulton, \emph{Intersection theory},
Proposition~1.8 and its arithmetic-base version in \S20.2.
The Picard-scheme construction is FGA, Expos\'e~232,
Theorem~3.1 (printed pp.~148--149), applied to the smooth
projective relative curve; the canonical abelian component is
FGA, Expos\'e~236, Theorem~3.5 (printed p.~239), with the smooth
Picard schemes of the curve fibres.  The latter fibre smoothness
also follows from the vanishing of their $H^2(\mathcal O)$.
The relative-motive formalism over schemes smooth over a Dedekind
domain is recalled in the supplied Laga--Shnidman source,
\emph{Ceresa cycles of bielliptic Picard curves}, \S4.1--4.5;
the extension of Deninger--Murre to that base is explained at the
start of \S4.2, citing K\"unnemann, \S1, Remark~1.1.
The generic projector, the character list, and its rational field
of definition are JAC1964 JAC-P13 and JAC-P17.  The new assertion
here is the integral-at-5 localization isomorphism for that
projector.  The finite character check is reproduced by
\texttt{work/arithmetic/check\_fermat\_projector.py}; the script
does not prove any statement about rational equivalence.

\begin{lemma}[C44-F02]\label{C44-F02}
For the lattice $T$ in (F01.3),
\[
 H^0(\mathbf Q,T/5^nT)=0\quad(n\ge1),
 \qquad H^0(\mathbf Q,T\otimes\mathbf Q_5/\mathbf Z_5)=0.
 \tag{F02.1}
\]
Consequently
\[
 eH^3_{\mathrm{et}}(J,\mathbf Z/5^n(2))=0,
 \qquad
 eH^4_{\mathrm{et}}(J,\mathbf Z/5^n(2))
       \simeq H^1(\mathbf Q,T/5^nT).             \tag{F02.2}
\]
\end{lemma}

\begin{proof}
We make AR-23's Frobenius-at-11 check independent of a choice of
integral lattice.  By JAC-P17, the weight-three character is
$\psi_+^3$, where $\psi_+$ is the character of
$E_+:y^2=x^3+x$.  Direct counting gives
$\#E_+(\mathbf F_{11})=12$, so the degree-one Frobenius
polynomial is $X^2+11$.  Cubing its two roots gives
$X^2+11^3$ on $N$.  After the twist $(2)$, the polynomial is
\[
                 X^2+\frac1{11}.
\]
It follows that
\[
       \det(1-\operatorname{Fr}_{11}\mid T)
              =\frac{12}{11}\in\mathbf Z_5^\times.
 \tag{F02.3}
\]
For any stable $\mathbf Z_5$ lattice, this determinant condition
makes $1-\operatorname{Fr}_{11}$ invertible.  Its reductions are
invertible modulo every $5^n$, and it is invertible on the torsion
quotient as well.  Invariance under the entire Galois group implies
invariance under this Frobenius, proving (F02.1).  The conclusion
does not change if arithmetic rather than geometric Frobenius is
used, since both Frobenius and its inverse are integral units.

Apply the Hochschild--Serre spectral sequence with finite
coefficients and the integral projector $e$.  By C44-F01 the
projected geometric cohomology occurs only in row $b=3$.
The spectral sequence therefore degenerates on that row and
identifies its total degrees three and four with the two Galois
groups in (F02.1)--(F02.2).
\end{proof}

\paragraph{Check scope.}
The numerical part of (F02.3), including the point count and its
reduction $12/11\equiv2\pmod5$, is an exact integer/rational
computation in the accompanying script.  This is a lattice-uniform
form of the Frobenius test used in AR-23, not a new proof of
AR-23's nonzero Abel--Jacobi class.

\subsection{The precise obstruction after passing to finite coefficients}

Write
\[
 B=eH^4_{\mathrm{et}}(J,\mathbf Z_{(5)}(2)),
 \qquad
 C=eH^3_{\mathrm{nr}}(J,\mathbf Q_5/\mathbf Z_5(2)).
 \tag{F03.1}
\]
The cohomology defining $B$ is \emph{\'etale motivic cohomology},
not continuous $5$-adic cohomology.  The distinction is essential.
The group $C$ is unramified cohomology over the ground field
$\mathbf Q$, with residues at the codimension-one points of $J$.
The operations defining $e$ act on it by restriction, norm and
the automorphisms of $J_{\mathbf Q(i)}$.

\begin{theorem}[C44-F03]\label{C44-F03}
There is a canonical short exact sequence
\[
                 0\longrightarrow A\longrightarrow B
                   \longrightarrow C\longrightarrow0,
 \tag{F03.2}
\]
and both $A$ and $B$ have no nonzero $5$-primary torsion.
For every $n\ge1$, the kernel of the finite-coefficient cycle map
\[
 c_n:A/5^nA\longrightarrow H^1(\mathbf Q,T/5^nT)
\]
is canonically $C[5^n]$.  In particular, for
$\widehat A=\varprojlim_n A/5^nA$, the completed cycle map satisfies
\[
 \ker\!\left(\widehat c:\widehat A\longrightarrow
                    H^1_{\mathrm{cont}}(\mathbf Q,T)\right)
          = T_5C:=\varprojlim_n C[5^n],            \tag{F03.3}
\]
where the transition on $C[5^n]$ is multiplication by five.

Let
\[
 \rho:A\otimes_R\mathbf Q_5\longrightarrow
                  \widehat A\otimes_{\mathbf Z_5}\mathbf Q_5,
 \qquad K_{\mathrm{comp}}=\ker\rho.
\]
The full scalar-extended Abel--Jacobi kernel in the question has
the exact description
\[
 0\longrightarrow K_{\mathrm{comp}}
 \longrightarrow\ker\alpha_5
 \xrightarrow{\rho}
 \operatorname{im}\rho\cap
             (T_5C\otimes_{\mathbf Z_5}\mathbf Q_5)
 \longrightarrow0.                              \tag{F03.4}
\]
The intersection is taken inside
$\widehat A\otimes\mathbf Q_5$.  Neither group on the two ends
of (F03.4) has been proved to vanish.
\end{theorem}

\begin{proof}
For a smooth variety over a field, the codimension-two
comparison sequence is
\[
 0\longrightarrow\operatorname{CH}^2(J)
 \longrightarrow H^4_{\mathrm{et}}(J,\mathbf Z(2))
 \longrightarrow H^3_{\mathrm{nr}}(J,\mathbf Q/\mathbf Z(2))
 \longrightarrow0.
\]
The exact locator is Kahn, \emph{On the injectivity and
non-injectivity of the $l$-adic cycle class maps}, (1.5),
author p.~3, citing his \emph{Classes de cycles motiviques
\'etales}, Proposition~2.9.  Localizing at five and applying
the idempotent $e$ gives (F03.2).

The coefficient triangle for $\mathbf Z_{(5)}(2)$ has
$\mathbf Z/5^n(2)=\mu_{5^n}^{\otimes2}$ as its finite-coefficient
term.  Its long exact sequence and (F02.2) give
$B[5^n]=0$ and an injection
\[
 B/5^nB\hookrightarrow
       eH^4_{\mathrm{et}}(J,\mathbf Z/5^n(2)).     \tag{F03.5}
\]
The relevant coefficient comparison is the Merkurjev--Suslin
case of the norm-residue theorem; it is recorded in Kahn,
loc.~cit., (3.2) and (4.1), author pp.~7 and~10.
Since $A\hookrightarrow B$, also $A[5^n]=0$.

Apply multiplication by $5^n$ to (F03.2).  The snake lemma gives
\[
 0\longrightarrow C[5^n]\xrightarrow{\delta_n}A/5^nA
 \longrightarrow B/5^nB.
 \tag{F03.6}
\]
Explicitly, for $c\in C[5^n]$, choose $b\in B$ above $c$;
then $5^n b\in A$, and
$\delta_n(c)=5^n b\bmod5^nA$.  This is independent of the
choice of lift.  The absence of $5^n$-torsion in $B$ proves
its injectivity.  Composing with (F03.5) and the
Hochschild--Serre identification gives precisely $c_n$.
Thus (F03.6) proves the finite-coefficient assertion.

If $c\in C[5^{n+1}]$ and $b$ is a lift, then the lift $5b$
of $5c$ shows that reduction of $\delta_{n+1}(c)$ modulo $5^n$
equals $\delta_n(5c)$.  The transition is therefore multiplication
by five.  Taking inverse limits in (F03.6) is exact at the first
two terms, so the kernel of the completed finite-coefficient
cycle map is $T_5C$.  This uses left exactness only; no
unverified surjectivity of transitions is needed.

Continuous cohomology has the comparison exact sequence
\[
 0\longrightarrow\varprojlim{}^1_n
       H^0(\mathbf Q,T/5^nT)
 \longrightarrow H^1_{\mathrm{cont}}(\mathbf Q,T)
 \longrightarrow\varprojlim_n H^1(\mathbf Q,T/5^nT)
 \longrightarrow0.
\]
Its left term is zero by C44-F02.  Equivalently, in total
degree four on $J$ the potential $\varprojlim{}^1$ term is the
limit of the zero groups in (F02.2).  Thus the completed map
just constructed is the continuous cycle map, proving (F03.3).
This is why an inverse limit is legitimate in this particular
summand although it is not interchangeable with continuous
cohomology for an arbitrary variety over a number field.

After tensoring by $\mathbf Q_5$, the cycle map from $A$
factors as $\widehat c\circ\rho$.  Its rational projected
Hochschild--Serre component is the Abel--Jacobi map
$\alpha_5$ from the task.  The inclusion of its image in
$H^1_f$ is the carried AR-16a formalism and does not change
the kernel.  Flatness of $\mathbf Q_5$ over $\mathbf Z_5$
identifies the kernel of $\widehat c\otimes\mathbf Q_5$ with
$T_5C\otimes\mathbf Q_5$.  Taking the inverse image under
$\rho$ and then its image yields (F03.4).
\end{proof}

\paragraph{What this attack removes, and what it leaves.}
C44-F01 removes the ambiguity of a varying arithmetic spread for
the Chow source.  C44-F02 removes both the global degree-three
term and the continuous-cohomology $\varprojlim{}^1$ ambiguity
from the finite-coefficient calculation.  The remaining group
in (F03.3) is the actual Tate module of one specified unramified
group.  It is not the Selmer group.  A Selmer bound does not bound
this Tate module or prove the tensor-to-completion injectivity in
(F03.4).  Specialization of rational odd-motive Chow classes,
already ruled out as a detector by AR-20, supplies neither
missing assertion.

\subsection{A bounded-denominator formulation on rational planes}

For a finite free $R$-submodule $M\subset A$, let
$W_M=M\otimes_R\mathbf Q\subset V$.  Equation (F03.2) and
the absence of $5$-torsion identify $B$ with a subgroup of $V$:
if $b\in B$, choose $5^n b=a\in A$ and send $b$ to $a/5^n$.
This map is well defined and injective, since $C$ is $5$-primary
torsion.  This convention makes $B\cap W_M$ unambiguous.

\begin{proposition}[C44-F04]\label{C44-F04}
The following are equivalent for every finite free $M\subset A$:
\begin{enumerate}
\item $\alpha_5$ is injective on $M\otimes_R\mathbf Q_5$;
\item $T_5(B/M)=0$;
\item $(B\cap W_M)/M$ is finite;
\item there is an integer $b_M\ge0$ such that
      $5^{b_M}(B\cap W_M)\subset M$.
\end{enumerate}
More precisely,
\[
 \ker(\alpha_5|_{M\otimes\mathbf Q_5})
       \simeq T_5(B/M)\otimes_{\mathbf Z_5}\mathbf Q_5.
 \tag{F04.1}
\]
Choose a nonzero multiple $a_\eta\in A$ of the carried class
$\eta$.  It is enough to verify these equivalent conditions for
the modules
\[
                  M_z=R a_\eta+Rz\quad(z\in A).
 \tag{F04.2}
\]
The condition holds when $z$ lies in $\mathbf Q\eta$, by AR-23.
For arbitrary $z$ the bounded-denominator condition in (4) is
unproved.
\end{proposition}

\begin{proof}
Apply the argument of (F03.6) to
$0\to M\to B\to B/M\to0$.  Both $M$ and $B$ have no
$5$-torsion.  It gives
\[
 0\longrightarrow T_5(B/M)\longrightarrow\widehat M
                         \longrightarrow\widehat B.
\]
Since $M$ is finite free over $R$, its completion is
$M\otimes_R\mathbf Z_5$.  The injection (F03.5) implies
that the kernel of the map to $\widehat B$ is exactly the
kernel of the continuous cycle map.  Rationalizing proves
(F04.1).

All $5$-power-torsion elements of $B/M$ lie in
\[
 D_M=(B\cap W_M)/M\ \subset\ W_M/M
          \simeq(\mathbf Q_5/\mathbf Z_5)^{\operatorname{rank}M}.
\]
Conversely every element of $D_M$ is $5$-power torsion, since
$R$ already inverts the other primes.  Thus $T_5(B/M)=T_5D_M$.
A subgroup of a finite direct sum of $\mathbf Q_5/\mathbf Z_5$
is a direct sum of finitely many copies of
$\mathbf Q_5/\mathbf Z_5$ and a finite group.  One proof is to
take its Pontryagin dual: it is a quotient of the finite free
$\mathbf Z_5$ module, hence is finitely generated; the structure
theorem over the discrete valuation ring $\mathbf Z_5$ and
duality give the assertion.  It follows that $T_5D_M$ is finite
free, and vanishes exactly when $D_M$ is finite.  Finiteness of
this subgroup is equivalent to a bound on its exponent, giving
the equivalence with (4) and also with (1), by (F04.1).

Finally AR-23 makes $\alpha_5(a_\eta)\ne0$ and gives a
one-dimensional target.  If all the two-generator modules
(F04.2) have injective scalar-extended map, none can have rank
two.  Hence every $z\in A$ belongs to $\mathbf Q\eta$, so
$V=\mathbf Q\eta$ and the full map is injective.  The converse
is immediate.  On the line $\mathbf Q\eta$, AR-23 already gives
injectivity, so the equivalent bounded-denominator conclusion
is proved on that line without any hypothesis about the rest
of the Chow group.
\end{proof}

\begin{proposition}[C44-F05: an algebraic model of the missing step]
\label{C44-F05}
Even finite generation of the integral Chow-like source,
one-dimensional surjective regulator image, a nonzero distinguished
class, and a positive definite real form do not formally imply
the scalar-extended injectivity in C44-F04.  There is an explicit
model satisfying all those properties in which the completed
kernel is one-dimensional and arises from a nonzero unramified-like
Tate module.
\end{proposition}

\begin{proof}
Choose the root $a\in\mathbf Z_5$ of $a^2=6$ with $a\equiv1\pmod5$.
Hensel's lemma gives this root, and $a\notin\mathbf Q$ because
six is not a rational square.  Put
\[
 A_0=R^2,\qquad
 B_0=\{(r,s)\in\mathbf Q^2:r+as\in\mathbf Z_5\},
 \qquad f(r,s)=r+as.
\]
Then $A_0\subset B_0$, both are $5$-torsion-free, and
\[
  B_0/A_0\simeq\mathbf Q_5/\mathbf Z_5.
 \tag{F05.1}
\]
To see the last statement, take the kernel of the surjection
$(\mathbf Q_5/\mathbf Z_5)^2\to\mathbf Q_5/\mathbf Z_5$,
$(r,s)\mapsto r+as$, whose restriction on the first coordinate
is the identity.  The inverse image of this kernel in
$\mathbf Q^2$ is exactly $B_0$.

The scalar-extended regulator is
$\mathbf Q_5^2\to\mathbf Q_5$, $(r,s)\mapsto r+as$.
It is surjective, sends $\eta_0=(1,0)$ to one, and has kernel
$\mathbf Q_5(-a,1)$.  It is even injective on the underlying
rational vector space $\mathbf Q^2$.  The form
$h((r,s),(r',s'))=rr'+ss'$ is positive definite and
$h(\eta_0,\eta_0)=1$.  Neither fact removes the scalar-extended
kernel.

For a completely explicit division tower, choose integers $a_n$
with $a_n\equiv a\pmod{5^n}$ and
$a_{n+1}\equiv a_n\pmod{5^n}$, and set
\[
 b_n=5^{-n}(-a_n,1)\in B_0.
\]
Then $b_n\bmod A_0$ has order $5^n$ and
$5b_{n+1}-b_n\in A_0$.  The resulting nonzero element of
$T_5(B_0/A_0)$ maps to $(-a,1)\in\widehat A_0$.
Thus the precise obstruction in C44-F03 is present in the
model, despite all the other listed data.  This is a model of
an invalid inference, not a geometric counterexample to
Abel--Jacobi injectivity or to a height conjecture.
\end{proof}

\subsection{What finite coefficients see of the fourfold component}

\begin{proposition}[C44-F06]\label{C44-F06}
Let $F/\mathbf Q$ be finite and let $D/F$ be a smooth projective
geometrically connected curve with a point $d_0\in D(F)$.
On $X=D\times J_F$, put
\[
 q=\pi_1^D\otimes e,\qquad
 A_D=q\operatorname{CH}^2(X;R),\qquad
 B_D=qH^4_{\mathrm{et}}(X,R(2)),
 \quad R=\mathbf Z_{(5)}.
\]
Then $B_D$ is uniquely divisible by five and hence is a
$\mathbf Q$-vector space.  The motivic comparison identifies it
canonically with the rational Chow component:
\[
       B_D\simeq A_D\otimes_R\mathbf Q.
                                                        \tag{F06.1}
\]
In particular $A_D$ has no $5$-primary torsion, and the
finite-coefficient cycle map on $A_D/5^nA_D$ is zero for every $n$.
Writing
\[
 C_D=qH^3_{\mathrm{nr}}(X,\mathbf Q_5/\mathbf Z_5(2)),
\]
one has
\[
 C_D=B_D/A_D,\qquad
 C_D[5^n]\simeq A_D/5^nA_D,
 \qquad T_5C_D\simeq\widehat A_D.             \tag{F06.2}
\]
Thus finite generation of $B_D$ as an $R$-module would force
the entire rational component in JAC-P18 to vanish.  Such finite
generation is not established here; finite generation of $A_D$
alone would not have that consequence.
\end{proposition}

\begin{proof}
The point projector $\pi_1^D$ is integral.  By C44-F01 and
the K\"unneth formula, $q$ selects geometric cohomology only in
degree four, where the integral lattice after twist is
\[
 T_D=H^1_{\mathrm{et}}(D_{\overline F},\mathbf Z_5)
                  \otimes_{\mathbf Z_5}T.
\]
It is free, of rank $4g(D)$.  At any place of $F$ over five,
its rational realization has Hodge--Tate weights
$-2,-1,1,2$ (each with multiplicity $g(D)$, up to the
simultaneous opposite weight convention).  Indeed the two
factors before twist have weights $\{0,1\}$ and $\{0,3\}$.
None is zero.  The Hodge--Tate comparison for smooth proper
varieties therefore implies
\[
                 H^0(F,T_D)=0.                 \tag{F06.3}
\]
Here one may use Scholze, \emph{$p$-adic Hodge theory for
rigid-analytic varieties}, Corollary~1.8, author p.~4, applied
to the analytification, together with functoriality for the
algebraic projectors.  This use needs neither good reduction
of $D$ at five nor a Hodge or Tate conjecture.

Finite-coefficient Hochschild--Serre now gives
\[
 qH^3_{\mathrm{et}}(X,\mathbf Z/5^n(2))=0,
 \qquad
 qH^4_{\mathrm{et}}(X,\mathbf Z/5^n(2))
       =H^0(F,T_D/5^nT_D).
                                                        \tag{F06.4}
\]
As in C44-F03, the coefficient triangle implies that
$B_D[5^n]=0$ and injects $B_D/5^nB_D$ into the second group
in (F06.4).  Taking inverse limits gives
\[
 \widehat B_D\hookrightarrow
        \varprojlim_n H^0(F,T_D/5^nT_D)
        =H^0(F,T_D)=0.
\]
The maps $B_D/5^{n+1}B_D\to B_D/5^nB_D$ are surjective,
so each projection from their inverse limit is surjective.
Consequently $B_D/5^nB_D=0$ for every $n$.  Together with
$B_D[5^n]=0$ this proves unique divisibility by five.
The other primes are already inverted in $R$.

The comparison sequence (F03.2) is valid in codimension two
on $X$ as well, and is compatible with $q$.  Its quotient
is $5$-primary torsion.  Rationalizing its injection
$A_D\hookrightarrow B_D$, and using that $B_D$ is already
a rational vector space, proves (F06.1).  The finite-coefficient
cycle map factors through $B_D/5^nB_D=0$.  Applying the snake
lemma to multiplication by $5^n$ gives (F06.2), including the
transition by multiplication by five; its inverse limit is
the final displayed isomorphism.

A nonzero rational vector space is not finitely generated as
a $\mathbf Z_{(5)}$-module.  Thus the stated finite generation
condition forces $B_D=0$.  JAC-P18 has already proved that
every algebraic class in this projected component is
homologically zero, so this is the vanishing that its
componentwise Murre~(D) would require.  The elementary inclusion
$R\subset\mathbf Q$ shows why finite generation of the smaller
group $A_D$ cannot substitute for finite generation of $B_D$.
\end{proof}

\paragraph{Scope of the additional attack.}
C44-F06 does not prove Murre~(D).  It identifies the relevant
\'etale motivic group as entirely divisible, with all of its
completed Chow source accounted for by unramified cohomology.
This makes precise why torsion-freeness or finite-coefficient
cycle computations alone cannot kill the JAC-P18 obstruction.
It concerns cote~44, Th.~1--3 (scan~002) and the boundary of
Th.~5 (scan~004); it does not change the known polarization
on the incidence higher Picard factor.  In the frameworks of
SGA and Chow motives, it is the tensor summand
$h^1(D)\otimes N$ of weight four with no $(2,2)$ part,
followed by the same codimension-two \`etale-motivic comparison
as C44-F03.  Its vanishing is a rational-equivalence issue,
not the semisimplicity or signed $I$ assertion of cote~16,
scan~049.  No 032/Milne input is used.

\subsection{The remaining finiteness question on one arithmetic model}

\begin{proposition}[C44-F07]\label{C44-F07}
For the abelian scheme $\mathcal J/S$ of C44-F01, restriction
also gives an isomorphism on the projected \emph{\'etale
motivic} groups:
\[
 eH^4_{\mathrm{et}}(\mathcal J,R(2))
                 \xrightarrow{\ \sim\ } B.
                                                        \tag{F07.1}
\]
The completion $\widehat B$ is a free $\mathbf Z_5$-module
of rank exactly one.  Its rational realization is the Selmer
line already realized by $\eta$.

The following two conditions are equivalent:
\begin{enumerate}
\item the scalar-extended Fermat Abel--Jacobi map is injective;
\item $eH^4_{\mathrm{et}}(\mathcal J,R(2))$ is finitely
      generated over $R$.
\end{enumerate}
Condition (2) is not proved here.  In particular the proved
finite generation of the completion is not being substituted
for it.
\end{proposition}

\begin{proof}
Purity and localization for \'etale motivic cohomology
give the exact segment
\[
 \bigoplus_{s\in S^{(1)}}
 H^2_{\mathrm{et}}(\mathcal J_s,R(1))
 \longrightarrow H^4_{\mathrm{et}}(\mathcal J,R(2))
 \longrightarrow H^4_{\mathrm{et}}(J,R(2))
 \longrightarrow\bigoplus_{s\in S^{(1)}}
 H^3_{\mathrm{et}}(\mathcal J_s,R(1)).
                                                        \tag{F07.2}
\]
An exact source for the needed purity is Geisser,
\emph{Motivic cohomology over Dedekind rings},
Theorem~1.2(1), published p.~774: after inverting the residue
characteristic no truncation is needed.  All residue
characteristics of $S$ differ from five and are units in $R$.
The Bloch--Kato premise of that 2004 statement is supplied by
the norm-residue theorem (in this weight, Merkurjev--Suslin).
The localization of the cycle complex and passage to the
generic fibre give the displayed direct sum.  This is also
the localization used in Geisser, \emph{Hasse principles for
\'etale motivic cohomology}, \S4, diagram~(2).

The first outer group in (F07.2) is a sum of Picard groups,
because $\mathbf Z(1)_{\mathrm{et}}=\mathbf G_m[-1]$.
Its $e$-part vanishes by C44-F01.  The last is a sum of
cohomological Brauer groups localized at five.  For a fibre
$\mathcal J_s$, finite-coefficient Hochschild--Serre and the
character calculation of C44-F01 give
\[
           eH^2_{\mathrm{et}}(\mathcal J_s,\mu_{5^n})=0
           \quad(n\ge1).
\]
Indeed its projected geometric cohomology starts in degree
three.  Kummer's sequence implies
$e\operatorname{Br}(\mathcal J_s)[5^n]=0$ for every $n$.
The Brauer group is torsion, so its localization at five has
zero $e$-part.  Apply $e$ to (F07.2) to obtain (F07.1).

The lattice $T$ is a lisse $\mathbf Z_5$-sheaf on $S$.
Proper smooth base change and the same single-row calculation
give
\[
 eH^4_{\mathrm{et}}(\mathcal J,\mathbf Z/5^n(2))
          =H^1_{\mathrm{et}}(S,T/5^nT),
 \qquad eH^3_{\mathrm{et}}(\mathcal J,\mathbf Z/5^n(2))=0.
                                                        \tag{F07.3}
\]
For the second identity, Frobenius at eleven still belongs
to $\pi_1(S)$, so C44-F02 applies.  Finite-coefficient
cohomology of this arithmetic scheme is finite.  Alternatively,
in degree one this follows from the finiteness of number
field extensions of bounded degree unramified outside
$\{2,5,\infty\}$.
It follows from the coefficient comparison that each
$B/5^nB$ is finite, and that
\[
 \widehat B\hookrightarrow
           H^1_{\mathrm{cont}}(S,T).
                                                        \tag{F07.4}
\]
There is no inverse-limit correction by (F07.3).
The group on the right is a finite free $\mathbf Z_5$-module:
it is a compact inverse limit of finite groups, its quotient
modulo five injects into $H^1(S,T/5T)$, and hence compact
Nakayama gives finite generation; its $5$-torsion vanishes
because $H^0(S,T/5T)=0$.  Restriction to
$H^1_{\mathrm{cont}}(\mathbf Q,T)$ is injective by inflation.

The group $B$ is rationally $V$, by (F03.2).  Therefore the
rational image of every element of $B$ in (F07.4) belongs to
the Selmer line of AR-23.  Its intersection with the finite
free lattice $H^1_{\mathrm{cont}}(S,T)$ is a closed rank-one
lattice.  The image of $B$ is dense in $\widehat B$, whose
image in (F07.4) is closed by compactness.  Hence
$\widehat B$ lies in that rank-one lattice.  It is nonzero,
because $a_\eta\in A\subset B$ has nonzero regulator.
This proves its asserted rank, including its rational image.

If $B$ is finitely generated over $R$, torsion-freeness makes
it finite free.  Its rank equals the rank of $\widehat B$,
namely one.  As $A\subset B$ and $A\otimes_R\mathbf Q=V$,
this forces $V=\mathbf Q\eta$, and AR-23 then proves
injectivity.  Conversely, injectivity and AR-23 give
$V=\mathbf Q\eta$.  Apply C44-F04 to $M=Ra_\eta$: since
its rational span is all of $V$, the bounded-denominator
conclusion places $B$ inside $5^{-b_M}M$.  Thus $B$ is a
submodule of a finite free $R$-module and is finitely
generated.  This proves the equivalence.
\end{proof}

\paragraph{Meaning of the arithmetic finiteness gate.}
C44-F07 makes the finiteness question a statement about one
specific summand on one scheme over $\mathbf Z[1/10]$.
Neither finite generation of ordinary Chow groups nor
finite generation of continuous \(5\)-adic cohomology proves
it.  Kahn, \S4, author pp.~10--11, explicitly isolates this
distinction; his Proposition~4.1 gives counterexamples to
unqualified \`etale-motivic finiteness on arithmetic schemes.
Geisser, \emph{Hasse principles}, Conjecture~4.3, states
Lichtenbaum's finiteness expectation for regular schemes
proper over $\mathbf Z$; that conjecture is not a theorem
for the present summand, and $\mathcal J/S$ is proper over
the indicated open base, not over all of $\mathbf Z$.
No claim of direct applicability or proof of that conjecture
is made.  In cote~44 this refines the arithmetic gate of
Th.~5, scan~004; in cote~16 it specifies information about
rational equivalence that conditions $C$ and $I$ on pure
realizations do not provide.  The equivalence does not by
itself count as a strictly smaller main gate under the
progress rule of the prompt.  No 032/Milne input is used.

\subsection{The three proposed arithmetic routes: exact scope}

\paragraph{C44-FA01: Euler systems.}
The supplied Lilienfeldt--Shnidman paper,
\emph{Derivatives of Rankin--Selberg $L$-functions and heights of
generalized Heegner cycles}, arXiv:2408.04375,
Theorem~1.1/13.4, identifies a height on a Hecke-generated cycle
space modulo the real-height radical.  Its Corollary~1.4 expressly
assumes the injectivity of its scalar-extended Abel--Jacobi map.
In \S15, author pp.~71--73, equation~(99) on p.~72 is the Euler-system
conclusion that a nonzero generalized Heegner cohomology class
spans the Selmer space; the next paragraph invokes the assumed
injectivity to pass back to Chow.  The displayed implication chain
on p.~73 marks both uses of Abel--Jacobi injectivity.  That is the exact missing
step, and C44-F03--F04 describe its integral obstruction for the
Fermat summand.  No Euler-system relation at the level of
Galois cohomology bounds the denominators in (F04.2).

There are also independent obstacles to a direct substitution
in the displayed theorem: its standing assumptions in \S1.1
require an imaginary quadratic field of odd discriminant and
an unramified Hecke character, whereas the Fermat character
over $\mathbf Q(i)$ has conductor $(1+i)^4$ by JAC-P17/AR-17.
The stated range is $0<t\le k-1$, so the weight-four central
case with $k+t=2$ is not a direct instance.  These scope failures
are not the reason the existing Selmer result fails---that result
is already available as AR-19/H9-T15---but they prevent the paper
from being treated as a new Chow-generation theorem here.

The Loeffler--Zerbes constructions consulted in this run
interpolate or bound Galois-cohomology classes, not the entire
Chow group.  The exact example is Jetchev--Loeffler--Zerbes,
\emph{Heegner points in Coleman families}, Theorem~A
(published p.~125), equivalently Theorem~5.4.1, whose
specializations are \emph{Abel--Jacobi images} of generalized
Heegner cycles.  An interpolation statement about these images
does not supply the missing denominator bound.  The existing
JAC1964 audit AR-17 and its
\texttt{LATEST\_RANK\_SOURCE\_SCOPE.tex} remain in force; their
statements are not reproved here.

\paragraph{C44-FA02: Gross--Schoen heights.}
AR-24--AR-25 are carried as the exact projected-diagonal and
sixteen-intersection identities.  Zhang,
\emph{Gross--Schoen cycles and dualising sheaves},
Theorem~2.3.5, author p.~27, proves the triple-intersection
identity; Conjecture~2.4.1, author p.~28, asserts semipositivity
of the tensor-correspondence pairing.  The latter is not proved
by the former.  Even adding that conjecture would give the
semipositivity conclusion recorded in AR-25, without an
equality criterion or the injectivity in task~3.  In particular
one must not substitute a real-height calculation for
the absent faithfulness of the $5$-adic regulator.  C44-F05
shows that positivity and the stated Selmer data alone are not
a formal bridge.  The supplied Gu draft, Theorem~1.3,
Corollary~5.4 and Conjecture~1.1, gives formulas and
conditional positivity, not that missing faithfulness.  Its
statements are treated here as supplied draft inputs, not as
independently verified public results, and none is used in the
proofs of C44-F01--F05.

\paragraph{C44-FA03: the algebraically trivial radical.}
JAC-P18 isolates the possible cycle
$Z=(\pi_1^C\otimes\pi_N)_*Z$ on $C\times J$, fixed by the
degree-four product projector.  Xu--Xu,
\emph{Remarks on Murre's conjecture on Chow groups},
Corollary~3.9, author p.~9, gives (A) and (B) for the product;
the earlier definition of (D), author p.~2, is the extra
cycle-class injectivity needed to annihilate this component.
Theorem (B) leaves degree four acting on codimension two and
therefore does not remove this component.  The special
Laga--Shnidman isomorphism in their Theorem~5.2 concerns a
different invariant summand of types $(2,1),(1,2)$, identified
with $h^1(E)(-1)$.  Such a Picard-motive isomorphism cannot hold
for the Fermat $N$, whose types are $(3,0),(0,3)$.
No proof of the needed componentwise (D) was found in these
sources or in the other primary sources checked in this run.

\paragraph{Programme crosswalk and current gate.}
C44-F01--F04 concern the arithmetic extension of cote~44,
scan~004, Th.~5, and the incidence/higher-Picard boundary of
Th.~1--3 on scan~002, as made explicit in the letter on
scans~008--009.  They do not alter JAC-T4.3 or establish a new
polarization of $P^i(X)$.  The actual cycle lives in the
extreme-type motive outside that Picard component.
In cote~16, scans~049--050, the conclusion is a limitation of
the passage from $C$, $I$, and semisimplicity of cohomological
correspondences to rational Chow groups: the kernel measured
by (F03.4) is invisible to those pure-realization assertions.
The corrected $A$, $A'$, $A''$ conditions of SN-C04 are
unchanged.

In Grothendieck's later language, C44-F01 is a statement over
one scheme over $\mathbf Z$; (F03.2) is the comparison from
codimension-two cycles to \emph{\'etale motivic} cohomology;
(F02.2)--(F03.4) concern $\ell$-adic sheaves, weights, and their
arithmetic realization.  For the generic fibre, SGA~6,
Expos\'e~XIV, \S4.1--4.2 (especially (4.1) and (4.4)),
identifies the rational Chow group in codimension two with
degree two of the topological and $\lambda$ filtrations on
$K_0(J)$.  The operators in (F01.1) act there by pullback.
This $\lambda$-ring translation supplies no bound on rational
equivalences.  No new EGA/FGA theorem or standard-conjecture
implication is asserted: the modern comparison theorems are
cited under their actual sources above.

The open rows \texttt{QFermat-AJ} and \texttt{Q35-full} are
\emph{not closed}.  Their exact new test is (F03.4), or,
equivalently, bounded five-power denominators in $B$ on every
rational plane (F04.2).  The original analytic order-one result
AR-11, positive value AR-13, and Selmer-line result AR-23
remain intact.  None of the present findings proves
$\operatorname{rank}\operatorname{CH}^2(N)=1$.

% END work/arithmetic/contribution.tex

% BEGIN work/arithmetic/fermat_excess_picard.tex
% C44_POS: arithmetic algebraic-radical continuation, /root/fermat_aj.
% All statements here are unconditional.  No 032 or Milne input is used.
% This is a reduction of JAC-P18, not a proof of full arithmetic AJ injectivity.

\section{Removing ambient Picard families from the Fermat radical}

We retain JAC-P13, JAC-P17 and JAC-P18 without reproving them.
In particular $J$ is the Fermat-quartic Jacobian and $p_N$ is its
actual rational Chow projector, with realization of type
$(3,0)+(0,3)$.  Its construction from origin-preserving automorphisms
and the canonical abelian Chow--K\"unneth projectors makes it
symmetrically distinguished.  The results below are independent of
032, of the supplied Milne report, and of Milne2022 Proposition~4.2.

For the geometric argument work over an algebraically closed field
$k\subset\mathbf C$ of characteristic zero; in the arithmetic
application take $k=\overline{\mathbf Q}$.  All Chow groups and
homomorphism groups are tensored with $\mathbf Q$.  The use of this
algebraic closure loses no nonzero rational Chow class defined over
a number field: any rational equivalence over the algebraic closure
is defined over a finite extension, and restriction followed by norm
is multiplication by its degree.

\subsection{An exact correspondence identity}

\begin{lemma}[C44-F15, divisor identity]\label{C44-F15-identity}
Let $A/k$ be an abelian threefold.  Let $p$ be a symmetrically
distinguished Chow projector with $p=\pi_3^A p=p\pi_3^A$ and with
Betti realization contained in
$H^{3,0}(A)\oplus H^{0,3}(A)$.  For every
$d\in\operatorname{CH}^1(A)$, write $L_d$ for multiplication by $d$,
viewed as a correspondence of degree one.  Then
\[
                 p L_d\pi_1^A=0
       \quad\hbox{in }\operatorname{CH}^4(A\times A).
       \tag{F15.1}
\]
Also $p_*\operatorname{CH}^1(A)=0$.
\end{lemma}

\begin{proof}
The canonical abelian projectors are polynomials in the pullback
correspondence of $[n]$, for any fixed integer $n\ge2$; thus they
are symmetrically distinguished.  Graphs of homomorphisms, their
transposes, products, and compositions are distinguished by
O'Sullivan's theorem.  Here composition uses only pullback and
pushforward by projections between abelian varieties and their
intersection products.

Decompose
\[
 d=d_++d_-,\qquad d_\pm=\tfrac12(d\pm[-1]^*d).
 \tag{F15.2}
\]
The class $d_+$ is symmetric and $d_-$ lies in
$\operatorname{Pic}^0(A)_{\mathbf Q}$.  Symmetric divisor classes
are distinguished: the distinguished lift of the numerical class
is symmetric, and the difference from $d_+$ lies in
$\operatorname{Pic}^0(A)_{\mathbf Q}$, on which $[-1]^*$ acts by
$-1$, so a symmetric such difference is zero.
Consequently $p L_{d_+}\pi_1^A$ is distinguished.  Its cohomological
realization is zero.  Indeed, multiplication by the divisor class
maps $H^1(A)$ into the $(2,1)+(1,2)$ part of $H^3(A)$, and $p$
kills that part.  A cohomologically trivial cycle is numerically
trivial.  Injectivity of the distinguished subalgebra into the
numerical Chow ring now proves
$p L_{d_+}\pi_1^A=0$ in rational Chow.

For $d_-$ no distinguished-cycle assertion is needed.
The theorem of the square gives $[n]^*d_-=n d_-$.
Multiplication in the Chow motive is compatible with $[n]^*$,
and $[n]^*$ acts as $n$ on $h^1(A)$.  Hence
$L_{d_-}\pi_1^A$ has image in the exact $n^2$ eigensummand
$h^2(A)$, and is killed by $p\pi_2^A=0$.
Equivalently one may multiply on the left by each of the
polynomial projectors and compare the distinct eigenvalues
$n^j$.  This proves (F15.1).

Finally, on divisor classes the symmetric and antisymmetric
parts have $[n]^*$ eigenvalues $n^2$ and $n$ respectively.  Thus
only $\pi_2^A$ and $\pi_1^A$ act on $\operatorname{CH}^1(A)$,
which proves the last assertion.
\end{proof}

\subsection{The exact excess-Picard quotient}

For a pointed smooth projective connected curve $(C,c_0)$ let
\[
 q_C=\pi_1^C\otimes p,
 \qquad W_C(A,p)=(q_C)_*\operatorname{CH}^2(C\times A).
 \tag{F15.3}
\]
Let $i:D\longrightarrow A$ be a proper morphism from a smooth
projective connected surface, generically finite onto an integral
surface.  Put
\[
 B_i=\operatorname{im}\bigl(i^*:\operatorname{Pic}^0(A)
                  \longrightarrow\operatorname{Pic}^0(D)\bigr),
 \qquad Q_i=\operatorname{Pic}^0(D)/B_i.
 \tag{F15.4}
\]
These are an abelian subvariety and an abelian quotient, respectively;
they are not Chow groups or algebraic representatives of codimension-two
cycles.

\begin{proposition}[C44-F15, family reduction]\label{C44-F15}
With the hypotheses of the preceding lemma there is a well-defined
linear map
\[
 \overline\Psi_{C,i}:
 \operatorname{Hom}(\operatorname{Jac}(C),Q_i)_{\mathbf Q}
                  \longrightarrow W_C(A,p)
 \tag{F15.5}
\]
whose image is exactly
\[
 (q_C)_*(1_C\times i)_*
                    \operatorname{CH}^1(C\times D).
 \tag{F15.6}
\]
In particular (F15.6) is zero if $i^*:\operatorname{Pic}^0(A)
\longrightarrow\operatorname{Pic}^0(D)$ is surjective.

For each integral divisor $D_0\subset A$, choose a resolution
$i_{D_0}:\widetilde D_0\longrightarrow D_0\hookrightarrow A$.
The sum of (F15.5) over these divisors is surjective:
\[
 \bigoplus_{D_0\subset A}
 \operatorname{Hom}\bigl(\operatorname{Jac}(C),Q_{i_{D_0}}\bigr)_{\mathbf Q}
       \twoheadrightarrow W_C(A,p).
 \tag{F15.7}
\]
No injectivity of this surjection is claimed.
\end{proposition}

\begin{proof}
Choose a point $d_0\in D(k)$.  The normalized Poincar\'e
description of divisors on a product gives the exact decomposition
\[
 \operatorname{Pic}(C\times D)_{\mathbf Q}
 =\operatorname{pr}_C^*\operatorname{Pic}(C)_{\mathbf Q}
 \oplus\operatorname{pr}_D^*\operatorname{Pic}(D)_{\mathbf Q}
 \oplus\operatorname{Hom}(\operatorname{Jac}(C),
                          \operatorname{Pic}^0(D))_{\mathbf Q}.
 \tag{F15.8}
\]
For clarity, a line bundle in the last summand is trivial on
$\{c_0\}\times D$ and on $C\times\{d_0\}$.  Its family of
restrictions to $D$ defines a pointed morphism
$C\to\operatorname{Pic}^0(D)$, hence a homomorphism from
$\operatorname{Jac}(C)$.  Conversely the universal Poincar\'e
bundle constructs this line bundle.  Subtracting the two restrictions
from any line bundle yields the asserted decomposition, with
uniqueness from the seesaw principle.  This is an assertion about
\emph{divisors on the source product}, where Picard theory is exact.

The second summand of (F15.8), after pushforward, is killed by
$\pi_1^C$, which annihilates the fundamental class $[C]$.
An element of the first summand pushes forward to
$a\boxtimes i_*[D]$, with $a\in\operatorname{CH}^1(C)$.
Its projection by $p$ on $A$ is zero by the last assertion of the
lemma, since $i_*[D]\in\operatorname{CH}^1(A)$.
Notice that $\pi_1^C$ by itself need not kill a degree-zero
$a$; it is the target projector that kills this term.

It remains to treat the normalized third summand.  Suppose its
homomorphism $\phi$ takes values in $B_i$.  A surjection of abelian
varieties splits up to isogeny: choose a polarization and its
orthogonal complementary abelian subvariety, so that the addition
map from the kernel and its complement is an isogeny.  It follows
that $\phi$ lifts, with rational coefficients, to a homomorphism
$\widetilde\phi:\operatorname{Jac}(C)\to\operatorname{Pic}^0(A)$.
Let $\ell\in\operatorname{CH}^1(C\times A)$ be its normalized
Poincar\'e divisor.  Its pullback to $C\times D$ has the required
cross homomorphism.  Adjusting its normalization at $d_0$ changes
it only by the first two summands of (F15.8), already disposed of.
The projection formula gives
\[
 (1_C\times i)_*(1_C\times i)^*\ell
       =\ell\cdot\operatorname{pr}_A^*(i_*[D]).
 \tag{F15.9}
\]
The class $\ell$ is in the exact $\pi_1^A$ summand on its
$A$ factor: pullback by $1_C\times[n]$ multiplies its normalized
Poincar\'e class by $n$.  Applying $q_C$ to (F15.9) therefore
gives zero by (F15.1).  This proves that the map on cross
homomorphisms kills
$\operatorname{Hom}(\operatorname{Jac}(C),B_i)_{\mathbf Q}$.

The isogeny category of abelian varieties is semisimple, by the
same polarization-complement argument.  Hence applying
$\operatorname{Hom}(\operatorname{Jac}(C),-)_{\mathbf Q}$ to
$0\to B_i\to\operatorname{Pic}^0(D)\to Q_i\to0$ is exact,
including surjectivity onto the last term.  This proves the
well-defined factor map (F15.5), independent of a rational
splitting, and the image assertion (F15.6).

To prove (F15.7), take an integral codimension-two subvariety
$Z\subset C\times A$, so $\dim Z=2$.  Its image in $A$ has
dimension one or two.  In dimension one, writing that image as
$\Gamma$, dimension and irreducibility imply
$Z=C\times\Gamma$, so its $\pi_1^C$ projection is zero.
In dimension two its image is an integral divisor $D_0$.
Over the dense open where $\widetilde D_0\to D_0$ is an
isomorphism, take the inverse image of $Z$ and then its closure
$\widetilde Z\subset C\times\widetilde D_0$.
It is an integral divisor on a smooth variety, and proper
pushforward gives $[Z]$ with degree one.  Thus the projected
class of $Z$ belongs to (F15.6).  These classes generate the
Chow group, which proves the surjection.
\end{proof}

\subsection{Only singular ample surfaces with excess irregularity remain}

\begin{proposition}[C44-F16]\label{C44-F16}
Under the hypotheses of C44-F15, the right side of (F15.7) is
generated already by those integral divisors $D_0\subset A$ which
are ample and whose resolution has irregularity
\[
            q(\widetilde D_0)>3.
 \tag{F16.1}
\]
For such an ample divisor,
\[
       \dim Q_{i_{D_0}}=q(\widetilde D_0)-3.
 \tag{F16.2}
\]
All non-ample image surfaces contribute zero.  Every normal ample
divisor with rational singularities contributes zero, including
every smooth ample divisor.  Therefore any nonzero class in
$W_C(A,p)$ requires a singular ample image surface with nonzero
excess Picard quotient and a nonzero map from $\operatorname{Jac}(C)$
to that quotient.

Applied to $(J,p_N)$, this is a strict restriction of the family
component left open by JAC-P18.  It does not establish its vanishing
and does not settle the scalar-extended arithmetic Abel--Jacobi
injection of AR-16a/C44-F07.
\end{proposition}

\begin{proof}
First suppose $D_0$ is not ample.  The connected translation
stabilizer $B$ of $D_0$ has positive dimension, and $D_0$ is the
inverse image of a divisor in $A/B$.  Indeed, the usual
ampleness criterion for an effective divisor on an abelian variety
identifies ampleness with finiteness of its translation stabilizer;
stability under $B$ then makes the divisor saturated for the
smooth quotient map.  Choose a complementary abelian subvariety
$A'$; the sum map
\[
                    u:B\times A'\longrightarrow A
 \tag{F16.3}
\]
is an origin-preserving isogeny.  Flat pullback of $D_0$ is a union
of products $B\times\Gamma_0$, with $\Gamma_0$ an integral
divisor of $A'$.  The dimension of $B$ is one or two.

Pullback by $u$ is an isomorphism of rational Chow motives.
For example a complementary isogeny composed with $u$ is $[n]$,
and $[n]^*$ is invertible on every canonical eigensummand
$h^j$; the inverse is a polynomial in the canonical projectors
and homomorphism graphs.  Consequently the transported projector
$p'=u^*p(u^*)^{-1}$ is distinguished, remains a summand of
$h^3(B\times A')$, and has only the two extreme Hodge types.
It suffices to kill each pulled-back family, since $u^*$ is
injective and intertwines the two projectors.

If $\dim B=2$, then $A'$ is an elliptic curve and each
$\Gamma_0$ is a point.  Its product with $B$ is smooth, and
restriction $\operatorname{Pic}^0(B\times A')\to
\operatorname{Pic}^0(B\times\Gamma_0)$ is surjective.  C44-F15
kills its entire family contribution.

If $\dim B=1$, let $\Gamma\to\Gamma_0\subset A'$ be the
normalization.  The surface $B\times\Gamma$ is smooth, and
\[
 \operatorname{Pic}^0(B\times\Gamma)
       =\operatorname{Pic}^0(B)\times\operatorname{Jac}(\Gamma).
 \tag{F16.4}
\]
After the two pullback summands in (F15.8) are removed, a
divisor on $C\times(B\times\Gamma)$ is the sum of a normalized
cross divisor from $C\times B$ and one from $C\times\Gamma$.
The first homomorphism has values in $\operatorname{Pic}^0(B)$,
which is inherited from $B\times A'$, so C44-F15 kills it.
The pushforward of the second is of the form
\[
                        [B]\boxtimes\gamma,
       \qquad \gamma\in\operatorname{CH}^2(C\times A'),
 \tag{F16.5}
\]
after permuting the factors.  This statement concerns the actual
cycle: the source divisor is pulled back from $C\times\Gamma$,
and the map to $B\times A'$ is the identity on $B$.

The canonical projector $\pi_0^B$ fixes $[B]$.  On the other hand
\[
                  p'(\pi_0^B\otimes\Delta_{A'})=0
       \quad\hbox{as a rational Chow correspondence}.
 \tag{F16.6}
\]
Both factors in (F16.6) are distinguished.  In cohomology only
$H^0(B)\otimes H^3(A')$ can map into $p'$, and the third
cohomology of an abelian surface has types $(2,1)+(1,2)$.
The map is therefore zero in cohomology and hence zero in Chow
by O'Sullivan.  Equations (F16.5)--(F16.6) kill the remaining
term.  Resolving a family over $B\times\Gamma_0$ by the
degree-one strict-transform construction in F15 shows that this
argument covers every family on the original non-ample image
surface, not only a previously chosen resolution.

Now suppose $D_0$ is ample and $i:\widetilde D_0\to A$ is a
resolution.  The image of the homomorphism
$\operatorname{Alb}(\widetilde D_0)\to A$ contains a translate
of $D_0$.  It must be all of $A$: otherwise dimension forces
its image to be an abelian surface and $D_0$ to be a translate
of that surface, which is not ample.  Duality therefore shows
that $i^*:\operatorname{Pic}^0(A)\to
\operatorname{Pic}^0(\widetilde D_0)$ has finite kernel.
The image has dimension three, proving (F16.2).  F15 kills
every term with $q(\widetilde D_0)=3$.

Finally suppose that $D_0$ is normal with rational singularities.
For a resolution $r:\widetilde D_0\to D_0$, these hypotheses
give $r_*\mathcal O_{\widetilde D_0}=\mathcal O_{D_0}$ and
$R^1r_*\mathcal O_{\widetilde D_0}=0$.  Hence
$H^1(\widetilde D_0,\mathcal O)=H^1(D_0,\mathcal O)$.
Since $K_A$ is trivial, Kodaira vanishing and Serre duality give
\[
 H^1(A,\mathcal O_A(-D_0))=
 H^2(A,\mathcal O_A(-D_0))=0.
 \tag{F16.7}
\]
The exact sequence
$0\to\mathcal O_A(-D_0)\to\mathcal O_A\to
\mathcal O_{D_0}\to0$ then identifies
$H^1(A,\mathcal O_A)$ with
$H^1(\widetilde D_0,\mathcal O_{\widetilde D_0})$.
Thus $q(\widetilde D_0)=3$.  Equivalently the pullback of
Picard varieties has isomorphic tangent spaces and is an
isogeny in characteristic zero.  The asserted vanishing follows
from F15.  A smooth divisor is a special case.

Combining these eliminations with (F15.7) proves every assertion.
The last implication is only necessary: the existence of a map
to $Q_i$ does not assert that its image under
$\overline\Psi_{C,i}$ is nonzero.
\end{proof}

\paragraph{Application to the arithmetic radical, with exact scope.}
JAC-P18 says that a nonzero geometrically algebraically trivial
$N$-class would be detected by a nonzero class of
$W_C(J,p_N)$ for a curve over a finite extension.  F15--F16
remove all its families whose image surfaces are non-ample, or
whose resolutions have no Picard quotient beyond the ambient
Jacobian.  It is sufficient to kill the maps (F15.5) for
singular ample surfaces with $q(\widetilde D_0)>3$.
This is a genuinely smaller geometric family gate than demanding
Murre(D) for every $C\times J$: several complete source classes
have been annihilated in rational Chow, not merely in a realization.
It is not a proof that the remaining gate holds.  In particular
normalization of a singular divisor can increase irregularity,
and neither the Picard decomposition of source divisors nor
Poincar\'e reducibility makes their codimension-two pushforward
injective.  The full arithmetic AJ5 finite-generation problem in
C44-F07 remains unchanged.

\paragraph{Primary locators.}
O'Sullivan, \emph{Algebraic cycles on an abelian variety},
arXiv:0908.0626v1, 5 August 2009, 73 pages: introductory
Theorem (i)--(iii), printed p.~2; Theorem~5.1.6, printed
p.~55, for the multiplicative canonical eigendecomposition;
Theorem~6.2.5 and Corollary~6.2.6, printed pp.~70--71,
for distinguished cycles and the numerical section.
Primary URL: \url{https://arxiv.org/pdf/0908.0626}.
The published article is J. reine angew. Math. 654 (2011),
1--81, DOI 10.1515/crelle.2011.025; the locators just given
are explicitly those of the checked author version.

Bost--Charles, \emph{Some remarks concerning the Grothendieck
period conjecture}, author copy of the 2014 online article,
\S3.2.2, formula~(3.1), printed p.~17, and Propositions~3.5--3.6,
printed pp.~17--18, give the exact Picard-product and
Poincar\'e identifications used in (F15.8).
Primary URL:
\url{https://www.imo.universite-paris-saclay.fr/~francois.charles/crelle-periods.pdf}.

Debarre, \emph{Singularities of divisors on abelian varieties},
20 March 2006, introductory paragraph, printed p.~1, explicitly
states the quotient description of a non-ample effective divisor.
Primary author URL:
\url{https://www.math.ens.psl.eu/~debarre/DivVarAb.pdf}.
The proof above explains the subsequent isogeny, normalization,
and exact projector arguments; none of those Chow conclusions is
being attributed to that source.

% END work/arithmetic/fermat_excess_picard.tex

% BEGIN work/arithmetic/fermat_difference_surface.tex
% C44_POS /root/fermat_aj. Actual geometric diagnostic, not F07 movement.
\section{A normal difference surface with excess irregularity}
\label{sec:c44-f26}

\begin{proposition}[C44-F26]\label{C44-F26}
Let $C$ be a smooth nonhyperelliptic curve of genus three over
$k=\overline{\mathbf Q}$, let $J=\operatorname{Jac}(C)$, and let
$a:C\hookrightarrow J$ be an Abel embedding.  Give the image
\[
 D=C-C=\{a(x)-a(y):x,y\in C\}\subset J
\]
its reduced structure.  The difference map
\[
 d:C\times C\longrightarrow D
 \tag{F26.1}
\]
is a proper birational morphism, is an isomorphism away from the
diagonal and $0$, and contracts the diagonal to $0$.
The surface $D$ is an integral normal ample divisor, with
\[
                   [D]=2[\Theta]\quad\hbox{in }\operatorname{NS}(J)_{\mathbf Q},
 \tag{F26.2}
\]
where $[\Theta]$ is the canonical principal polarization.
Its only singular point is $0$.  The morphism (F26.1) is a
resolution with
\[
 q(C\times C)=6,
 \qquad
 \frac{\operatorname{Pic}^0(C\times C)}
      {d^*\operatorname{Pic}^0(J)}\simeq J
 \tag{F26.3}
\]
up to the canonical Jacobian autoduality.

In particular the normalization of $D$ is $D$ itself and its
normalization conductor has empty support, although its resolution
has three dimensions of excess Picard variety.  For the Fermat
quartic, all the projected P18 families supported over this $D$
are nevertheless zero by C44-F19.  Thus this example rules out
a conductor-only description of the general excess-Picard gate;
it is neither a nonzero radical class nor progress on C44-F07.
\end{proposition}

\begin{proof}
Nonhyperellipticity implies that an effective divisor of degree two
on $C$ has at most one section up to scalars.  Indeed, two
independent sections, after removing a possible base point, would
give a map of degree at most two to $\mathbf P^1$; degree one is
impossible for a positive-genus curve and degree two would make
$C$ hyperelliptic.  Consequently two linearly equivalent effective
degree-two divisors are equal.

Suppose $d(x,y)=d(x',y')$.  Then
\[
                       x+y'\sim x'+y,
 \tag{F26.4}
\]
so the two unordered pairs in (F26.4) are equal.  Either
$(x,y)=(x',y')$, or $x=y$ and $x'=y'$.  Also
$a(x)=a(y)$ implies $x=y$, since the Abel map is an embedding.
The fibre over $0$ therefore has support the diagonal, and every
geometric fibre away from $0$ has exactly one point.  These
arguments hold after any algebraically closed extension of $k$.
In particular the image has dimension two and the map is
generically of degree one in characteristic zero.

We check the differential before using normality of the image.
Invariant one-forms on $J$ pull back isomorphically to
$H^0(C,\omega_C)$.  After translating tangent spaces of $J$ to
the origin, the tangent line of $a(C)$ at $a(x)$ is the line
represented by the canonical point of $x$ in
$\mathbf P(H^0(C,\omega_C)^\vee)$.  The canonical morphism of a
nonhyperelliptic genus-three curve is an embedding.  Thus, when
$x\ne y$, these two tangent lines are distinct, and
\[
              (v,w)\longmapsto da_x(v)-da_y(w)
 \tag{F26.5}
\]
is injective.  Hence $d$ is unramified off the diagonal: its
relative differentials vanish at every geometric point there,
and then vanish by Nakayama's lemma.

The restriction
\[
        (C\times C)\setminus\Delta_C\longrightarrow J\setminus\{0\}
\]
is proper, as a base change of the original proper map.  It is
quasi-finite and universally injective by (F26.4), and unramified
by (F26.5).  A proper unramified universally injective morphism
is a closed immersion (equivalently, it is a proper unramified
monomorphism).  Its image is reduced and has support
$D\setminus\{0\}$, so it is precisely that reduced open
subscheme.  This proves that $D$ is smooth off $0$ and that
(F26.1) is an isomorphism there.

The image is integral because $C\times C$ is integral.  Every
integral codimension-one subvariety of the smooth variety $J$
is an effective Cartier divisor.  Thus $D$ is a hypersurface in
a regular threefold and is Cohen--Macaulay, in particular $S_2$.
It is regular in codimension one since it is smooth off a single
closed point.  Serre's $R_1+S_2$ criterion proves normality.

For (F26.2), write $m:J\times J\to J$ for addition, and let
$W_2=a(C)+a(C)$.  Riemann's theta-divisor description says that
$W_2$ represents the canonical principal polarization.  The
addition map $C\times C\to W_2$ has generic degree two, since
$C^{(2)}\to W_2$ is birational and the quotient by interchanging
the factors has degree two.  Therefore
\[
                 m_*([a(C)]\boxtimes[a(C)])=2[W_2]
                       \quad\hbox{in }\operatorname{CH}^1(J).
 \tag{F26.6}
\]
In Betti cohomology inversion acts by $+1$ on $H^4(J)$, the
degree of the curve class.  It follows that the cycle classes of
$a(C)$ and $-a(C)$ are equal, and hence
\[
 \operatorname{cl}\bigl(d_*[C\times C]\bigr)
       =\operatorname{cl}\bigl(m_*([a(C)]\boxtimes[a(C)])\bigr).
 \tag{F26.7}
\]
The left side is $\operatorname{cl}([D])$ by generic degree one.
The divisor cycle-class map is injective on
$\operatorname{NS}(J)_{\mathbf Q}$, so (F26.6)--(F26.7) give
(F26.2).  Ampleness depends only on the numerical divisor class;
therefore $D$ is ample.

The smooth surface $C\times C$ has irregularity six by the
Kunneth decomposition.  If $D$ were smooth at $0$, it would be a
smooth ample divisor in $J$ and would have irregularity three
by C44-F16.  The proper birational map between the two smooth
projective surfaces would preserve their Albanese varieties,
contradicting these irregularities.  Thus $0$ is singular and
(F26.1) is a resolution of the asserted kind.  In particular the
singularity cannot be rational, again by C44-F16.

Finally,
$\operatorname{Pic}^0(C\times C)=J\times J$ after Jacobian
autoduality.  Pullback by the difference map sends $L$ to
\[
                         (a^*L,(a^*L)^{-1}).
 \tag{F26.8}
\]
Indeed the theorem of the square on $\operatorname{Pic}^0(J)$
gives $m^*L=\operatorname{pr}_1^*L\otimes\operatorname{pr}_2^*L$,
and inversion pulls $L$ to $L^{-1}$.  The homomorphism
$a^*:\operatorname{Pic}^0(J)\to\operatorname{Pic}^0(C)$ is an
isomorphism, so the image in (F26.8) is the anti-diagonal.
The addition homomorphism $J\times J\to J$ identifies the
quotient with $J$, proving (F26.3).

For the Fermat quartic the hypotheses hold, since it is a smooth
plane quartic.  Apply C44-F19 with the two maps $a$ and $-a$.
That proposition kills every projected divisor family on the
smooth product and every integral family with image $D$, including
ones whose pushforward is not a Cartier family on $D$.  This uses
the actual rational Chow identities of F19, not merely its
realized cohomology.  Normality makes the normalization an
identity map, so there is no normalization conductor that could
account for the nonzero quotient (F26.3).  This proves the scope
assertions as well.
\end{proof}

\paragraph{Primary locators and status.}
Milne, \emph{Jacobian Varieties}, the corrected transcription dated
12 June 2021 of the published 1986 chapter, Proposition~2.2,
printed pp.~6--7, identifies invariant differentials; Theorem~5.1
and Corollary~5.5, printed pp.~16--19, give the symmetric-power
maps and their degrees; Theorem~6.6 and its proof, printed
pp.~21--23, identify $W_{g-1}$ with the canonical principal
polarization and give Jacobian autoduality:
\url{https://www.jmilne.org/math/xnotes/JVs.pdf}.
The two scheme-theoretic criteria used above are Stacks Project,
Lemma~41.7.2, Tag~04XV, and Lemma~10.157.4, Tag~031S:
\url{https://stacks.math.columbia.edu/tag/04XV},
\url{https://stacks.math.columbia.edu/tag/031S}.
Their complete displayed proofs were read.  All these are classical
unconditional inputs; no unrefereed 032/Milne-2026 statement is
used.  F26 preserves an actual geometric obstruction to a proposed
shortcut.  It supplies no new nonzero P18 class, no additional
family elimination beyond F19, and no change to F07.

% END work/arithmetic/fermat_difference_surface.tex

% BEGIN work/arithmetic/fermat_boundary_curve_tests.tex
% C44_POS, /root/fermat_aj.  A scoped P18 boundary diagnostic.
% No conditional 032/Milne input.  No vanishing or F07 movement is asserted.

\section{Boundary-curve tests for the excess Picard family}

The following refines the parameter description in C44-F17.  It does
not prove the remaining Chow maps zero.  Work over
$k=\overline{\mathbf Q}$, with rational Chow groups and the conventions
of F17.  Let $D\subset A$ be an integral ample divisor in an abelian
threefold, let $\nu:T\to D$ be its normalization, and let
$r:S\to T$ be a projective resolution.  Write $i:S\to A$ for the
composite.  Fix the extreme projector $p$ of F15 and put $N=(A,p)$;
we write $p_N=p$.  If $D$ is normal, the conductor curves below are absent.

Let $Z\subset D$ and $\widetilde Z\subset T$ be the conductor
subschemes, with their full scheme structures.  Let $C_a$ run through
all normalizations of the one-dimensional irreducible components of
$\widetilde Z_{\mathrm{red}}$, and let $E_b$ run through the
normalizations of the exceptional curves of $r$.  These are smooth
projective curves.  Zero-dimensional conductor components contribute
no Jacobian.

\begin{proposition}[C44-F27, boundary diagnostic]\label{C44-F27}
Put
\[
 Q=\operatorname{Pic}^0(S)/
       \operatorname{im}\bigl(\operatorname{Pic}^0(A)\to
                                      \operatorname{Pic}^0(S)\bigr),
 \qquad
 M=\prod_a\operatorname{Jac}(C_a)\times
                         \prod_b\operatorname{Jac}(E_b).
 \tag{F27.1}
\]
Then $Q$ is isogenous to a direct factor of $M$.  In particular, if
all these boundary curves are rational, then $Q=0$ and F15 kills
all P18 families with image $D$.

Choose a surjective homomorphism $u:M\to Q$, and let $u_F$ denote
its restriction to a factor $\operatorname{Jac}(F)$ of $M$.  Choose
basepoints, and put $h_F=u_F\circ a_F:F\to Q$.  For the actual
residual cycle $\beta_i$ of F17 define
\[
 \gamma_F=(\pi_1^F\otimes p_N)_*
                    (h_F\times1_A)^*\beta_i
                  \in\operatorname{CH}^2(F\times A).
 \tag{F27.2}
\]
Then
\[
                         \beta_i=0
                 \quad\Longleftrightarrow\quad
                         \gamma_F=0\text{ for every }F.
 \tag{F27.3}
\]
The maps $h_F$ depend on isogeny splittings; they are not claimed
to be the maps induced by the boundary inclusions.  The assertion
is an exact test for the particular F17 cycle, not vanishing of
all groups $\operatorname{Hom}(h^1(F)(-1),N)$.
\end{proposition}

\begin{proof}
\leavevmode\par
\paragraph{The normal-resolution contribution.}
First let $T$ be any normal projective surface.  Let
$B_E\subset\operatorname{Alb}(S)$ be the abelian subvariety generated
by the images of the $\operatorname{Jac}(E_b)$, using the pointed
maps of the exceptional curves into $S$.  Those images are killed
by the map $\operatorname{Alb}(S)\to\operatorname{Alb}(T)$.
Conversely, the Albanese map of $S$, followed by the quotient by
$B_E$, is constant on every fibre of $r$.  Indeed $r$ has connected
fibres: its Stein factorization has finite birational second map,
which is an isomorphism because $T$ is normal.  Each positive
dimensional fibre is a connected union of exceptional curves;
a map constant on every component is constant on the whole fibre.

Let $P=\operatorname{Alb}(S)/B_E$.  The reduced image of
$(r,a):S\to T\times P$ projects properly and quasi-finitely onto
$T$, since all its geometric fibres have a single point.  This
projection is finite and birational, hence is an isomorphism by
normality.  Thus $a$ descends to a morphism $T\to P$.  Universality
of the Albanese variety of a normal projective variety now shows
\[
 B_E=\ker\bigl(\operatorname{Alb}(S)\to
                              \operatorname{Alb}(T)\bigr)^0.
 \tag{F27.4}
\]
Here and below finite kernels do not affect an assertion up to
isogeny.  Dually,
\[
 Q_r=\operatorname{Pic}^0(S)/r^*\operatorname{Pic}^0(T)
 \quad\text{is an isogeny subfactor of}\quad
                              \prod_b\operatorname{Jac}(E_b).
 \tag{F27.5}
\]
One can also read this from restriction of line bundles to the
$E_b$: its connected kernel is $r^*\operatorname{Pic}^0(T)$.
The pullback $r^*$ itself is injective, since $r_*\mathcal O_S=
\mathcal O_T$ and the projection formula recovers a line bundle
from its pullback.

\paragraph{The normalization contribution, including nilpotents.}
Set $P_D=\operatorname{Pic}^0(D)$ and
$P_T=\operatorname{Pic}^0(T)$.  F22 proves that $P_D$ is proper
and isogenous to $\operatorname{Pic}^0(A)$; $P_T$ is an abelian
variety because $T$ is normal.  The composite
$\operatorname{Pic}^0(A)\to\operatorname{Pic}^0(S)$ has finite
kernel: its dual is the surjection $\operatorname{Alb}(S)\to A$
from F20.  Consequently $P_D\to P_T$ has finite kernel as well.

The conductor square is locally the ring identity
\[
 R=\overline R\times_{\overline R/I}R/I,
 \tag{F27.6}
\]
where $I$ is the conductor ideal.  Invertible sheaves on $D$ are
therefore equivalent to an invertible sheaf on $T$, one on $Z$,
and an identification of their restrictions to $\widetilde Z$.
This is finite-projective-module patching, applied on affine opens
and then glued; the square remains a fibre product after any
base change over the field.  In particular it applies to line
bundle families, not just to $k$-points of Picard groups.

Put $G_Z=\operatorname{Pic}^0(Z)$ and
$G_{\widetilde Z}=\operatorname{Pic}^0(\widetilde Z)$, and form
the quotient algebraic group
\[
 G=G_{\widetilde Z}/\operatorname{im}(G_Z).
 \tag{F27.7}
\]
Restriction gives $P_T\to G$.  Patching implies
\[
       \ker(P_T\to G)^0=\operatorname{im}(P_D\to P_T).
 \tag{F27.8}
\]
For clarity about identity components, form the algebraic-group
fibre product
$P_T\times_{G_{\widetilde Z}}G_Z$.  The inverse image of this fibre product in $\operatorname{Pic}(D)$
is a union of connected components, with identity component $P_D$.
This inverse image surjects onto the fibre product as an fppf sheaf,
by the preceding patching equivalence.  Its identity component therefore surjects
onto the identity component of this fibre product.  For example,
this follows by lifting its generic point fppf locally: the image
of any component of $\operatorname{Pic}(D)$ is a translate of the
closed image of the proper group $P_D$, so a component whose image
contains that generic point fills the connected component.
Projecting to $P_T$ gives (F27.8).  This argument does not replace
the conductor by its reduction in the patching step.

The image of $P_T$ in $G$ is proper, and is an abelian variety.
Thus
\[
 Q_\nu=P_T/\operatorname{im}(P_D)
 \tag{F27.9}
\]
is isogenous to a proper connected subgroup of $G$.  Its map to
the Chevalley abelian-variety quotient $G^{\mathrm{av}}$ has finite kernel, because
a subgroup that is both affine and proper is finite.

For any projective scheme $Y$ of dimension at most one, the
Chevalley abelian-variety quotient of $\operatorname{Pic}^0(Y)$ is
\[
                \prod_{F\subset Y_{\mathrm{red}}}
                       \operatorname{Jac}(\widetilde F).
 \tag{F27.10}
\]
Here the product runs through its one-dimensional irreducible
components.  To include nonreduced $Y$, successively use square-zero
nilpotent thickenings: the unit exact sequence has additive kernel
$\mathcal I$, and $H^2(Y,\mathcal I)=0$.  Restriction on Picard
identity components is consequently surjective with unipotent
kernel.  For reduced $Y$, normalization changes its Picard identity
component by a connected affine group (torus and unipotent parts),
as follows from the units sequence supported on its finite singular
locus.  Normalization gives the displayed product of Jacobians.
This proves (F27.10) and explains why nilpotent conductor structure
cannot simply be ignored in descent, although it supplies no new
abelian factor.

The quotient map $G_{\widetilde Z}\to G$ induces a surjection on
Chevalley abelian-variety quotients: any map to an abelian variety annihilates the
connected affine subgroup.  It follows that $Q_\nu$ is an isogeny
subfactor of $\prod_a\operatorname{Jac}(C_a)$.

\paragraph{Combining the two parts and detecting the Chow cycle.}
There is an exact sequence of abelian varieties, up to the already
irrelevant finite kernels,
\[
             0\longrightarrow Q_\nu\longrightarrow Q
                      \longrightarrow Q_r\longrightarrow0.
 \tag{F27.11}
\]
Poincar\'e complete reducibility splits it in the isogeny category.
Together with (F27.5) and (F27.9)--(F27.10) this proves that $Q$ is
isogenous to a direct factor of $M$.  Clearing a common integer
denominator gives the surjective actual homomorphism $u$ used in
the statement.

On degree-one rational Chow motives, $u^*:h^1(Q)\to h^1(M)$ is
split injective.  The degree-one motive of $M$ is the direct sum
of the degree-one motives of its Jacobian factors, and the Abel
map identifies each $h^1(\operatorname{Jac}(F))$ with $h^1(F)$.
Consequently the combined projected pullbacks in (F27.2) are
injective on the component $\pi_1^Q\otimes p_N$ containing
$\beta_i$.  Applying a Chow-motivic left inverse recovers $\beta_i$
from these pullbacks, exactly as in F17.  This proves (F27.3).
If every boundary curve is rational, $M=0$, so $Q=0$ and the
last assertion follows from F15.
\end{proof}

\paragraph{What the boundary argument does not kill.}
Even when $D$ is normal and $i(E_b)$ is a point, (F27.3) is not
a zero statement.  The test map $h_{E_b}:E_b\to Q$ is chosen
through Picard duality and an isogeny splitting.  Its Poincar\'e
divisor on $E_b\times S$ need not be supported on
$E_b\times E_b$.  The degree issue is visible already in the
following natural cohomological construction, for an ample
class $H$ on $S$ and an exceptional curve map $j:E_b\to S$:
\[
 H^1(E_b)(-1)\xrightarrow{j_*}H^3(S)
 \xrightarrow{L_H^{-1}}H^1(S)(-1)
 \xrightarrow{i_*}H^3(A).
 \tag{F27.12}
\]
There are also the Picard summand projections and isogeny
splittings when representing a specified test map.  The inverse
Lefschetz operator on the Picard summand is algebraic, but it is
not proved here to have exceptional support.  Thus $(i\circ j)_*=0$
on the relevant direct curve map cannot be substituted for the
composite with $L_H^{-1}$ inserted.  The negative-definite
intersection matrix of the exceptional divisor controls its
Tate classes, and supplies no such support statement for the
excess degree-one motive.

For nonnormal $D$, patching determines precisely which Picard
families descend to Cartier families on $D$; F22 kills those.
A quotient family in (F27.9) has not been shown to have a
Cartier representative after pushforward.  Its boundary origin
therefore does not allow another application of F22.
The remaining assertion is exactly the vanishing of the particular
cycles (F27.2), equivalently of F17's $\beta_i$.  F26 shows why
conductor curves alone cannot be sufficient.  F19 already kills
all product-dominated examples, including that difference surface.
No full P18 vanishing, arithmetic descent vanishing, or actual
\'{e}tale-motivic finite generation follows here.  This diagnostic
is not a new advance of the F07 gate.

\paragraph{Primary locators and scope.}
For conductor patching use D.~Ferrand, \emph{Conducteur, descente
et pincement}, Bull.~Soc.~Math.~France \textbf{131} (2003),
Th\'eor\`eme~2.2(i),(iv), printed pp.~559--560, with the complete
proof on pp.~560--563:
\url{https://perso.imj-prg.fr/olivier-debarre/wp-content/uploads/ferrand-pub/Conducteurs.03.pdf}.
The argument applies to finite projective modules, hence to line
bundles, and retains the full conductor schemes.  A second precise
locator is Stacks, Lemma~15.6.9, together with Lemmas~15.6.4,
15.6.7--15.6.8 in \url{https://stacks.math.columbia.edu/tag/0D2G}.
For the normal graph descent use Stacks, Lemma~29.55.8
\url{https://stacks.math.columbia.edu/tag/0AB1}.
For the Albanese of a normal projective variety and the reduced
curve units sequence, use J.-L.~Colliot-Th\'el\`ene and
A.~Skorobogatov, \emph{The Brauer--Grothendieck group},
author version dated 31 October 2019, \S4.1, printed p.~105,
and Lemma~7.3.1 with its displayed proof, printed pp.~160--161:
\url{https://www.imo.universite-paris-saclay.fr/~jean-louis.colliot-thelene/BGgroup_book.pdf}.
These are unconditional inputs.  The nilpotent thickening argument
and the isogeny and cycle-detection steps are supplied above.

% END work/arithmetic/fermat_boundary_curve_tests.tex

% BEGIN work/arithmetic/axis_square_trace.tex
% C44_POS /root/fermat_aj. Root-proposed route, complete author draft.
% A28 reserved by /root. No 032/Milne-2026 input in this result.
\section{Regular Frobenius and the square of a closed nonabelian axis}
\label{sec:c44-a28}

\begin{theorem}[C44-A28: a product-polarized nonabelian family]
Use the family of axes, notation, and good supersingular primes of
C44-A13. Thus $A=E^6$, $E:y^2=x^3-x$, $h=tH$, and
$B=B_u\subset A$ is a smooth intersection of two divisors of class $h$,
containing $Y=E^2\times\{o\}^4$. Put $h_B=h|_B$ and
\[
 w=[Y]_B-t^{-2}\kappa^*\alpha_0,\qquad
 \beta=\int_Bw^2=6t^2-\frac1{4t^2}>0,
 \qquad K=\mathbf Q_\ell w.
 \tag{A28.1}
\]
The expression for $\alpha_0$ is the already checked C44-A13.3.
Let $W$ be the orthogonal complement of $K$ in the variable middle
cohomology. Its rank is $d=3600t^6-166$, an even positive integer.

A Zariski-dense set of closed axes in the C44-A13 parameter curve
has all of the following properties. Write $k=\mathbf F_q$ for a
field of definition of a selected axis and put
$\Phi=q^{-2}\operatorname{Frob}_q^*$ on $W$.
\begin{enumerate}
\item The characteristic polynomial $P(T)$ of $\Phi$ is rational and
has distinct roots. None is a root of unity, and the quotient of two
distinct roots is never a root of unity. These are statements about
one common polynomial, hence apply to every $\ell\ne p$.
\item The rational algebraic classes in $W\otimes W$, regarded as
middle correspondences on $B\times B$, identify exactly with
\[
 E=\mathbf Q[\Phi]\simeq\mathbf Q[T]/P(T).
 \tag{A28.2}
\]
This is a product of CM fields. Under this identification their
actual middle intersection form is
\[
 (a,b)\longmapsto\operatorname{Tr}_{E/\mathbf Q}(a b^\dagger),
 \qquad \Phi^\dagger=\Phi^{-1}.
 \tag{A28.3}
\]
It is positive definite. Neither a Tate conjecture nor equality of
numerical and homological equivalence is an input to this assertion.
\item The Hodge standard conjecture holds for $(B,h_B)$ and for
\[
 (B\times B,\ \rho\operatorname{pr}_1^*h_B+
                      \sigma\operatorname{pr}_2^*h_B),
 \qquad \rho,\sigma\in\mathbf Q_{>0},
 \tag{A28.4}
\]
in every primitive algebraic degree and every $\ell\ne p$.
The Lefschetz standard conjecture holds for $B$, its smooth
hyperplane section, and $B\times B$. Consequently numerical and
$\ell$-adic homological equivalence coincide on $B$ and $B\times B$.
The algebra of algebraic cohomological correspondences of degree zero
on $B$ has a positive adjoint trace form and is semisimple.
\end{enumerate}
No unrefereed 032/Milne 2026 input and no Milne 2022,
Proposition 4.2, is used. The polarization assertion is (A28.4);
no assertion for an arbitrary ample class on $B\times B$ is made.
\end{theorem}

\begin{proof}
\emph{A degree-safe strengthening of the closed-point selector.}
Set
\[
 M=\operatorname{lcm}(1,2,\ldots,2d^4).
 \tag{A28.5}
\]
If $\lambda,\mu$ are roots of a rational polynomial of degree $d$,
then $[\mathbf Q(\lambda,\mu):\mathbf Q]\le d^2$.
A root of unity among $\lambda/\mu$ therefore has degree at most
$d^2$, and the elementary bound used in C44-A12 puts its order at
most $2d^4$. The same bound covers roots of unity among the roots
of the original polynomial. Thus it is enough to impose
\[
 \det(g^M-1)\,
 \operatorname{disc}\bigl(\det(T-g^M)\bigr)\ne0.
 \tag{A28.6}
\]
This is a regular nonvanishing condition on the orthogonal group.
On a maximal torus of $\operatorname{SO}_d$ its eigenvalues are
$t_1,t_1^{-1},\ldots,t_{d/2},t_{d/2}^{-1}$. Choose the $t_i$
so that their $M$-th powers and inverse $M$-th powers are pairwise
distinct and none is one. Hence (A28.6) is nonempty on the geometric
monodromy identity component. The finite-quotient and constant-field
Chebotarev argument of C44-A12 applies verbatim to this open set.
It applies on every nonempty parameter open, proving Zariski density.
Passing from a rational point over a finite extension to its underlying
closed point does not lose (1): a root or ratio which is a root of
unity remains so after every positive power.

C44-A13 identifies the variable system and its orthogonal monodromy.
It also gives the rational characteristic polynomial as the quotient
of the full middle polynomial by the ambient polynomial and $T-1$.
Those polynomials are common to the Weil realizations. All roots of
$P$ have complex absolute value one, by the Weil theorem and the
normalization $q^{-2}$. In what follows no semisimplicity theorem
for Frobenius is assumed: its semisimplicity on $W$ follows from
its simple characteristic polynomial.

\emph{Actual algebraic projectors and the Lefschetz operators.}
We spell out the operators needed below, so an abstract motive
injection is not substituted for a comparison of forms.
Let $X\hookrightarrow A$ be a smooth complete intersection of $c$
divisors of class $h$, with $n=\dim X=6-c$, and write $i$ for the
embedding. For $j<n$, weak Lefschetz identifies $H^j(X)$ with
$i^*H^j(A)$. The map
\[
 i^*\,(L_A^{6-j})^{-1}\,i_*:
 H^{2n-j}(X)\longrightarrow H^j(X)
 \tag{A28.7}
\]
is the inverse of $L_X^{n-j}$: on $i^*x$ the composite before the
inverse is $i_*L_X^{n-j}i^*x=L_A^{6-j}x$.
All operators on the right are algebraic, since the Lefschetz
standard conjecture is known for abelian varieties. The middle
hard-Lefschetz map is the identity. These inverses supply algebraic
Kunneth, primitive, and Lefschetz projectors, and the lowering operator.
In particular the lowering operator is zero on the new primitive
middle summand. The construction is relative in a smooth lifting of
the complete intersection and the polarized abelian variety.

For $X=B$, let $p_{\rm amb}$ be the identity outside the middle
and the projector
\[
 \kappa^*(L_A^2)^{-1}\kappa_*
 \quad\hbox{on }H^4(B).
 \tag{A28.8}
\]
It is self-adjoint for Poincare duality. Its complementary projector
$p_V$ is supported on the variable primitive $H^4(B)$. The projector
$p_K(x)=\beta^{-1}(\int_Bxw)w$ on the middle is represented by
$\beta^{-1}w\times w$, and put $p_W=p_V-p_K$.
These are actual algebraic cohomological projectors, mutually
orthogonal, self-adjoint, and commuting with $L_B$ and its lowering
operator. The ambient projector extends over the fixed smooth trait
lift in C44-A13. The class $w$ also lifts in that family, although
that last fact is not needed for the separate $K$-block computations.
Write $U=p_{\rm amb}H^*(B)$. Then, as graded Lefschetz modules,
\[
 H^*(B)=U\perp K\perp W,
 \tag{A28.9}
\]
where $K$ and $W$ occur only in primitive degree four.

\emph{The rational centralizer and rational reconstruction of cycles.}
On $W$, the normalized Frobenius is represented by the rational
algebraic correspondence $p_Wq^{-2}\operatorname{Frob}_q^*p_W$.
A middle algebraic class $z$ on $B\times B$ projected by
$p_W\otimes p_W$ defines an endomorphism $a_z$ of $W$, by Poincare
duality. Its defining cycle is over a finite extension of $k$, so
$a_z$ commutes with some positive power $\Phi^n$.
Condition (1) implies that $\Phi^n$ still has pairwise distinct
roots and that its centralizer equals that of $\Phi$. Hence
\[
 a_z\in\operatorname{End}(W)^{\Phi^n}
       =\mathbf Q_\ell[\Phi]
       =E\otimes_{\mathbf Q}\mathbf Q_\ell.
 \tag{A28.10}
\]

It remains essential to recover rational, rather than merely
$\ell$-adic, coefficients. Use the basis $1,\Phi,\ldots,\Phi^{d-1}$
of $E$. The matrix
\[
 G_{ij}=\operatorname{Tr}(\Phi^{i+j}),\qquad0\le i,j<d,
 \tag{A28.11}
\]
is rational and nonsingular: it is the trace pairing of the etale
algebra $E$, or equivalently a Vandermonde matrix times its transpose.
For each $j$, $\operatorname{Tr}(a_z\Phi^j)$ is rational. Indeed it
is the cohomological trace of an actual rational algebraic
correspondence supported only in degree four, and hence equals its
intersection with the diagonal by the correspondence trace formula.
There is no alternating-degree sign left, since four is even.
Writing $a_z=\sum_i c_i\Phi^i$, the equations
\[
 G(c_i)_i=(\operatorname{Tr}(a_z\Phi^j))_j
 \tag{A28.12}
\]
force every $c_i\in\mathbf Q$. Conversely each power of $\Phi$ is
an algebraic correspondence supported on $W$, so every element of
$E$ occurs. This proves (A28.2) without an assumed Tate surjectivity,
an assumed finite rational rank of cycle classes, or an assumed $D$.

Let $q_W(x,y)=\int_Bxy$. It is nondegenerate. Since the degree of
Frobenius on the fourfold is $q^4$, its normalized action on $H^4$
is an isometry of this form. Thus its adjoint is $\Phi^{-1}$.
The involution preserves each field factor of $E$: complex conjugation
sends any root $\lambda$ of its rational irreducible polynomial to
$\overline\lambda=\lambda^{-1}$, a conjugate in the same field.
For every complex embedding it is complex conjugation, so its fixed
field is totally real. No root is real, since a real number of absolute
value one is $1$ or $-1$, excluded by (1). Each factor is consequently
a CM field, and this is its usual CM involution.

Under the identification $u\otimes v\mapsto(x\mapsto q_W(v,x)u)$,
the tensor form $q_W\otimes q_W$ is
$(a,b)\mapsto\operatorname{Tr}(a b^\dagger)$.
The Kunneth cup-product sign is positive, since both factors have
even degree four. Also $W\otimes W$ is primitive in degree eight
on the eightfold $B\times B$, and its Hodge-standard sign is
$(-1)^4=+1$. Its geometric primitive form is therefore exactly
(A28.3), with no additional sign or Lefschetz factor. Finally
\[
 \operatorname{Tr}_{E/\mathbf Q}(a a^\dagger)
       =\sum_{\sigma:E\hookrightarrow\mathbf C}|\sigma(a)|^2>0
       \quad(a\ne0),
 \tag{A28.13}
\]
with the sum taken over the embeddings of the field factors.
This proves (2).

\emph{The unconditional supersingular abelian input.}
For every $N$, all algebraic classes on the power $E_{\bar k}^N$
are rational Lefschetz classes, namely classes generated by divisors.
Here is the precise argument, including its rationality point.
After a finite extension, Frobenius on $H^1(E)$ is scalar. Tate's
endomorphism theorem gives
$\operatorname{End}(E)\otimes\mathbf Q_\ell=
\operatorname{End}_{\mathbf Q_\ell}(H^1(E))$.
Coordinate divisors and graphs of these endomorphisms therefore span
$H^2(E^N,\mathbf Q_\ell(1))$. Indeed its Kunneth pieces are the
one-dimensional coordinate $H^2$'s and the four-dimensional
$H^1(E)\otimes H^1(E)$ pieces. Their exterior products span every
even cohomology group. Thus divisor products span all even
cohomology over $\mathbf Q_\ell$.

Choose complementary rational divisor-product classes giving a
nonsingular intersection matrix. Pairing any actual rational
algebraic class with these complementary classes gives rational
intersection numbers. Inverting that rational matrix expresses the
class as a rational linear combination of the divisor products.
This upgrades the $\mathbf Q_\ell$ spanning statement to the
required statement about actual rational cycle classes.
Milne, \emph{Polarizations and Grothendieck's Standard Conjectures},
Remarks 3.7--3.8, proves unconditionally the Hodge standard
conjecture on these Lefschetz classes, and records that their
homological and numerical equivalences coincide for every Weil
cohomology. Thus we may use the actual signed primitive
positivity on $E^N$ in all degrees and for every $\ell\ne p$.
This is the published Lefschetz-class theorem, not the conditional
all-abelian statement of that paper and not 032.

\emph{Ambient primitive forms, including the ambient tensor block.}
For $0\le a\le2$, restriction identifies the primitive ambient part
of $H^{2a}(B)$ with $P^{2a}(A)$. This follows either from the
Lefschetz strings, which shorten by two, or by applying $\kappa_*$
and weak Lefschetz to the primitive equation. For $x,y$ in this part,
\[
 q_{B,a}(\kappa^*x,\kappa^*y)
 :=(-1)^a\int_B\kappa^*x\,\kappa^*y\,h_B^{4-2a}
 =(-1)^a\int_Axyh^{6-2a}.
 \tag{A28.14}
\]
An ambient algebraic class on $B$ has an algebraic inverse image,
by (A28.8) and the algebraic Lefschetz inverse. Hence (A28.14) and
the preceding paragraph prove positivity on every primitive
ambient algebraic part of $B$. Its remaining algebraic primitive
part is $K$, with norm $\beta>0$; $W$ has no algebraic classes.
This proves the assertion for $B$.

For the ambient tensor block on $T=B\times B$, fix $\rho,\sigma>0$ in
(A28.4), put $h_T=\rho h_1+\sigma h_2$, and let $0\le r\le4$.
The primitive summand
\[
 V_r=P^{2r}(T)\cap(U\otimes U)
 \tag{A28.15}
\]
has an algebraic relative self-adjoint projector and inverse
pairing, by (A28.7)--(A28.9) and tensor Lefschetz theory.
The relative correspondence $(\kappa\times\kappa)_*$ gives an
injection on it into
\[
 H^{2r+8}(A\times A)(4).
 \tag{A28.16}
\]
To check injectivity, $\kappa_*$ is injective on $U$ in each degree:
in degree at most four this is $L_A^2$ on the ambient image, and
above four it is dual to the weak-Lefschetz isomorphism in the
complementary degree. The same holds on every fibre of the
smooth trait lift of C44-A13.

We apply the positive-spectrum proof of C44-A08 to this actual
relative injection, replacing its conditional abelian positivity
input by the unconditional supersingular result just proved.
For clarity, the form being recovered is precisely
\[
 q_{T,r}(x,y)=(-1)^r\int_Txyh_T^{8-2r}.
 \tag{A28.17}
\]
Choose the algebraic direct sum of signed primitive forms $Q_A$
on the target (A28.16), with its Tate normalization, put
$q_1=\alpha^*Q_A$, and $c=q_{T,r}^{-1}q_1$.
Both forms polarize the generic complex primitive Hodge structure,
so $c$ has strictly positive real spectrum there. Its characteristic
polynomial is unchanged in the special fibre, and $c^{-1}$ is a
rational polynomial in $c$. On special algebraic classes $q_1$ is
positive, by Milne's Lefschetz-class theorem on $E^{12}$, and $c$
preserves those classes and is $q_1$-self-adjoint. On the finite
rational cyclic space generated by any given class and powers of
$c$, this gives
$q_{T,r}(z,z)=q_1(c^{-1}z,z)>0$ for $z\ne0$.
This is the complete form comparison needed from C44-A08; the
positivity of an unspecified abstract motive form is not used.
It works for each $r=0,1,2,3,4$, and includes the odd-times-odd
Kunneth pieces which can contain even algebraic classes.

\emph{All remaining primitive blocks and their exact factors.}
The projectors in (A28.9), and their products, commute with $L_T$
and are self-adjoint. Thus they decompose every primitive part
orthogonally, and send rational algebraic classes to rational
algebraic classes. Every even-degree normalized ambient Frobenius
eigenvalue is a root of unity, since $A$ is supersingular.
Consequently the blocks $W\otimes U$, $U\otimes W$,
$W\otimes K$, and $K\otimes W$ have no algebraic classes:
a normalized algebraic class is fixed by a power of Frobenius,
whereas every eigenvalue on these even-degree blocks is a root
of $P$ times a root of unity. Odd total degrees contain no
algebraic cycles in the first place.

It remains to check $K\otimes U$, $U\otimes K$, and $K\otimes K$.
Because $L_Bw=0$, tensoring by $w$ leaves the Lefschetz strings of
$U$ unchanged, with their degrees shifted by four. The primitive
parts in even degrees at most eight are therefore exactly
\[
 K\otimes P^{2a}_{\rm amb}(B),\quad
 P^{2a}_{\rm amb}(B)\otimes K,
 \qquad a=0,1,2,
 \tag{A28.18}
\]
and $K\otimes K$ in degree eight. No higher point on one of these
strings is primitive: on $L_B^jx$, with $j>0$ and total degree
at most eight, the defining primitive power remains nonzero
by hard Lefschetz on that string.
For primitive ambient $x,y$ of degree $2a$, (A28.17) gives
\[
 \begin{split}
 q_{T,2+a}(w\otimes x,w\otimes y)
       &=\beta\,\sigma^{4-2a}q_{B,a}(x,y),\\
 q_{T,2+a}(x\otimes w,y\otimes w)
       &=\beta\,\rho^{4-2a}q_{B,a}(x,y),\\
 q_{T,4}(w\otimes w,w\otimes w)&=\beta^2.
 \end{split}
 \tag{A28.19}
\]
Indeed every term of $h_T^{4-2a}$ having a positive power of
$h_1$ in the first expression vanishes against $w$, so the
remaining coefficient is exactly $\sigma^{4-2a}$, not an omitted
binomial coefficient. The sign is $(-1)^{2+a}=(-1)^a$.
For the product polarization $\rho=\sigma=1$, the primitive degrees and
forms outside $U\otimes U$ are therefore:
\[
\begin{array}{c|c|c}
\text{codimension on }T&\text{block}&\text{actual signed form}\\\hline
2&K\otimes H^0(B),\ H^0(B)\otimes K
  &720t^6\beta\quad\text{on each displayed generator}\\
3&K\otimes P^2_{\rm amb}(B),\ P^2_{\rm amb}(B)\otimes K
  &\beta\bigl(-\int_Bxyh_B^2\bigr)\\
4&K\otimes P^4_{\rm amb}(B),\ P^4_{\rm amb}(B)\otimes K
  &\beta\int_Bxy\\
4&K\otimes K&\beta^2\\
4&W\otimes W&\operatorname{Tr}_{E/\mathbf Q}(a a^\dagger).
\end{array}
 \tag{A28.20}
\]
Here $\int_Bh_B^4=\int_Ah^6=720t^6$.
Extraction by the algebraic correspondence
$\beta^{-1}\langle -,w\rangle$ shows that an algebraic class in
either first type in (A28.18) has an algebraic ambient factor;
there is no unstated rational spanning assertion.
Equations (A28.14), (A28.19), and (A28.13) prove positivity on
all these blocks. Together with (A28.15)--(A28.17) and the
vanishing mixed blocks, this proves the full product-polarized
assertion in every primitive degree.

The Lefschetz standard conjecture for $B$ follows from (A28.7),
and for $T$ from the tensor Lefschetz operators. Under strong
Lefschetz and the Hodge standard conjecture, $A$ and $D$ are
equivalent; since $B$ supplies $A$, we obtain $D$ for $B$ and $T$.
Finally, Kleiman's semisimplicity theorem and its positive
adjoint trace argument apply to $B$ and $B\times B$.
Precisely, Theorem 5-6(1), printed p.19, in
\emph{The Standard Conjectures} (1994) assumes
$B(X),\operatorname{Hdg}(X\times X)$ and asserts semisimplicity.
Its proof defines the adjoint $u'=*\,{}^tu\,*$ and proves
$\operatorname{Tr}(u'u)>0$ for every nonzero algebraic
cohomological correspondence $u$. Apply it with $X=B$.
This proves the asserted correspondence-algebra conclusion.
These axes are nonabelian: weak Lefschetz gives $b_1(B)=12$,
whereas an abelian fourfold has first Betti number eight.
\end{proof}

\begin{corollary}[C44-A28a: the corrected scan-050 conditions on this family]
Put $X=B$, let $Y'\subset B$ be a smooth hyperplane section for
$h_B$, put $T=B\times B$ with $h_T=h_1+h_2$, and let
$j:Z\hookrightarrow T$ be any smooth Segre hyperplane section.
Then all three conditions on the right side of SN-C04 hold:
\[
 C_G(Y'),\qquad R_0(T,Z),\qquad G_{L,0}(T,Z).
 \tag{A28.21}
\]
The weighted condition uses $L_Tj_*$, not ordinary Gysin alone.
\end{corollary}

\begin{proof}
The variety $Y'$ is a smooth intersection of three divisors of
class $h$ in $A$, so (A28.7) gives $C_G(Y')=B(Y')$.
Write $r=j^*$ and $g=j_*$. Weak Lefschetz is an isomorphism
$r:H^6(T)\longrightarrow H^6(Z)$, and $gr=L_T$.
Let $\mathcal D_2:H^{10}(T)\to H^6(T)$ be the algebraic inverse
of $L_T^2$, supplied by $B(T)$.
For $z\in\operatorname{Alg}^3(Z)$, its cohomological inverse
under $r$ is
\[
               x=\mathcal D_2 L_Tj_*z.
 \tag{A28.22}
\]
Indeed write $z=rx_0$ in cohomology; then the right side is
$\mathcal D_2L_T^2x_0=x_0$. All operators in (A28.22) are
algebraic. Hence $x$ is algebraic and $rx=z$, proving $R_0$.
For $y\in\operatorname{Alg}^5(T)$, put
$x=\mathcal D_2y$ and $z=rx$. They are algebraic and
\[
                       L_Tj_*z=L_T^2x=y.
 \tag{A28.23}
\]
This proves $G_{L,0}$. These degrees are the exact critical
ones: $\dim T=8$, $\dim Z=7$, and the weighted map is
$\operatorname{Alg}^3(Z)\to\operatorname{Alg}^5(T)$.

For a concrete description of the hyperplane, choose a basis
$s_0,\ldots,s_N$ of the embedding by $h_B$. Then $Z$ is of the
form
\[
             \sum_{i,j} c_{ij}s_i(x)s_j(y)=0.
 \tag{A28.24}
\]
Choose $(c_{ij})$ in the nonempty smooth-hyperplane open of
this projective coefficient space, over a finite extension of
$k$ if necessary. Equations (A28.22)--(A28.23) verify the two
conditions for every such choice; no condition on $D(Z)$ or
$I(Z)$ is being assumed. The hyperplane is not $Y'\times Y'$.
\end{proof}

\paragraph{Primary source locators and scope.}
Tate, \emph{Endomorphisms of abelian varieties over finite fields},
Invent. Math.2 (1966), 134--144, Main Theorem and equation (1),
printed p.134, were read in
\url{https://pazuki.perso.math.cnrs.fr/index_fichiers/Tate66.pdf}.
Milne, \emph{Polarizations and Grothendieck's Standard Conjectures},
Ann. of Math.155 (2002), 599--610, is used only through the
unconditional Lefschetz-class Remarks 3.7--3.8; the checked
11-page author version dated 19 September 2001 has these on p.9:
\url{https://www.jmilne.org/math/articles/2002aP.pdf}.
Kleiman, \emph{The Standard Conjectures}, Proc. Sympos. Pure Math.55,
Part1 (1994), 3--20, Proposition 5-1, printed p.17, and
Theorem 5-6(1) with its positive-adjoint-trace proof, printed p.19,
were read in the primary scan
\url{https://download.uni-mainz.de/mathematik/Algebraische%20Geometrie/Lehre/WS23.Padische1994.Kleiman.Standard.Conjectures.pdf}.
The selector uses the exact primary Chebotarev statement already
checked in C44-A12. C44-A13 supplies the actual nonempty family,
its relative lift, its monodromy, and the number $\beta$; those
settled proofs are not reproduced.

This gives actual closed nonabelian varieties, positive geometric
forms on their product-polarized squares, the scan-049 positive
adjoint trace and semisimplicity conclusion, and a uniform
verification of each corrected SN-C04 condition. It does not certify
that this family's Hodge standard conjecture was previously unknown
in the literature; that historical requirement of the original
example request is kept separate. It does not prove the general
fixed-axis sign, positivity for every ample class on these squares,
or the Fermat arithmetic Abel--Jacobi injection. No assertion that
the residual motive $W$ is of abelian type is made.

% END work/arithmetic/axis_square_trace.tex

% BEGIN work/arithmetic/axis_square_tate.tex
% C44_POS /root/fermat_aj. Separate corollary; A28/A28a left unchanged.
\section{Tate classes on the regular-Frobenius axis and its square}
\label{sec:c44-a28b}

\begin{corollary}[C44-A28b]
Let $B/k$ be any selected closed axis of C44-A28, where $k$ is
finite. The Tate conjecture holds in every codimension for $B$
and $B\times B$, for every $\ell\ne\operatorname{char}(k)$ and
after every finite extension of $k$. More explicitly, the cycle maps
\[
 \mathrm{CH}^r(X_{k'})_{\mathbf Q}\otimes_{\mathbf Q}\mathbf Q_\ell
 \longrightarrow
 H^{2r}(X_{\bar k},\mathbf Q_\ell(r))^{\operatorname{Gal}(\bar k/k')}
 \tag{A28b.1}
\]
are surjective for $X=B,B\times B$ and every finite $k'/k$.
The stronger pole-order formulation also holds: numerical and
$\ell$-adic homological equivalence agree, and the rational numerical
rank in codimension $r$ is the pole order at $s=r$ of $\zeta(X_{k'},s)$.
This is independent of the positivity and $D$ conclusions of A28:
the ingredients are its actual algebraic projectors, the regular
Frobenius condition, and the elementary Tate argument for powers
of the supersingular elliptic curve. No 032/Milne-2026 input is used.
\end{corollary}

\begin{proof}
We retain the notation $A=E^6$, $\kappa:B\hookrightarrow A$,
$U,K,W$, and $\Phi$ of C44-A28. All the projectors there are
defined over $k$. We first make explicit the algebraic left inverse
to the ambient Gysin injection, so the Tate claim does not rely on
an unproved descent assertion for an abstract motive injection.
Write $U^j=U\cap H^j(B)$. For $0\le j\le4$ define
\[
 R_j=\kappa^*(L_A^{6-j})^{-1}L_A^{4-j}:
 H^{j+4}(A)\longrightarrow U^j.
 \tag{A28b.2}
\]
The final codomain means the image; $\kappa^*$ has image $U^j$
in these degrees. Weak Lefschetz and the definition of the ambient
middle summand give $U^j=\kappa^*H^j(A)$. Since
$\kappa_*\kappa^*=L_A^2$, one has
\[
 R_j\kappa_*\kappa^*x
 =\kappa^*(L_A^{6-j})^{-1}L_A^{6-j}x=\kappa^*x.
 \tag{A28b.3}
\]
For $4<j\le8$ put
\[
 R_j=L_B^{j-4}\kappa^*(L_A^{j-2})^{-1}:
 H^{j+4}(A)\longrightarrow H^j(B)=U^j.
 \tag{A28b.4}
\]
Here $(L_A^{j-2})^{-1}$ maps $H^{j+4}(A)$ to
$H^{8-j}(A)$. The identity
\[
 \kappa_*R_j=L_A^{j-4}L_A^2(L_A^{j-2})^{-1}=1
 \tag{A28b.5}
\]
shows that $R_j$ is the inverse of $\kappa_*$, because Gysin in
these degrees is an isomorphism by the dual of weak Lefschetz.
Equations (A28b.2)--(A28b.5), with the algebraic Kunneth projectors,
give an actual rational algebraic operator $R$ of degree minus four
which is a left inverse of $\kappa_*$ on $U$.
All inverses on $A$ are the algebraic abelian Lefschetz inverses
already used in A28. The displayed operators are defined over $k$;
alternatively averaging a representative over a finite field of
definition gives the same $k$-defined cohomological operator.
Consequently
\[
 U\xrightarrow{\kappa_*}H^*(A)[4](2),\qquad
 U\otimes U\xrightarrow{(\kappa\times\kappa)_*}
                  H^*(A\times A)[8](4)
 \tag{A28b.6}
\]
have algebraic left inverses $R$ and $R\otimes R$. The notation
records the degree and Tate shifts, not an assertion about complexes.

For every $N$, the proof in A28 using Tate's endomorphism theorem
shows that divisor products span $H^{2s}(E^N,\mathbf Q_\ell(s))$
over $\mathbf Q_\ell$ after a finite extension where Frobenius on
$H^1(E)$ is scalar. These products are classes of actual cycles
over finite fields. This also proves the Tate assertion over any
smaller finite field $k'$ of definition: express a fixed vector as
a $\mathbf Q_\ell$-linear combination of such cycle classes, choose
one finite extension $k''/k'$ over which all the finitely many
cycles are defined, and average the expression over
$\operatorname{Gal}(k''/k')$. The invariant vector is unchanged.
The averaged cycles are $[k'':k']^{-1}$ times their proper
pushforwards from $E^N_{k''}$ to $E^N_{k'}$, so are rational cycles
defined over $k'$. This proves (A28b.1) for every $E^N$ and every
finite $k'$. It does not require a statement that all these classes
are fixed over the initially chosen field.

The Gysin injections and their left inverses in (A28b.6) now show
that every Tate vector in an even total degree of $U$ or
$U\otimes U$ is algebraic over its field of invariance. The second
assertion includes the odd--odd Kunneth pieces: their total degree
is even, and the target is the even cohomology of $E^{12}$.
The Tate twists in (A28b.6) send degree $2r$ to degree $2r+4$
with twist $r+2$, or to degree $2r+8$ with twist $r+4$, respectively.

The line $K=\mathbf Q_\ell w$ is the explicit algebraic class
in degree four. It therefore contributes Tate classes, as do
$K\otimes K$, $K\otimes U$, and $U\otimes K$; the assertions for
the latter two follow by the same left-inverse argument or by
external products. All normalized eigenvalues on the even ambient
cohomology are roots of unity, since $E$ is supersingular. None
of the normalized eigenvalues on $W$ is a root of unity, by A28(1).
Thus $W$, $W\otimes K$, $K\otimes W$, and the even-degree
parts of $W\otimes U$ and $U\otimes W$ have no Tate vectors over
any finite extension. An odd ambient factor has odd total degree
with $W$ and contributes to none of the groups in (A28b.1).

Only the block $W\otimes W$, in middle degree eight, remains.
Via the nondegenerate middle pairing it is $\operatorname{End}(W)$,
and normalized Frobenius acts by conjugation with $\Phi$.
If $k'/k$ has degree $n$, A28(1) implies that $\Phi^n$ has distinct
eigenvalues, hence
\[
 (W\otimes W)^{\operatorname{Gal}(\bar k/k')}
  =\operatorname{End}(W)^{\Phi^n}
  =\mathbf Q_\ell[\Phi].
 \tag{A28b.7}
\]
The powers $1,\Phi,\ldots,\Phi^{d-1}$ are actual algebraic
correspondences on $B\times B$ defined over $k$, with the
$p_W$ projectors inserted as in A28. They span the right side
of (A28b.7). This proves Tate surjectivity on the last block.
The algebraic orthogonal projectors decompose every group in
(A28b.1) into the blocks just treated. Their sum proves the result
for $B$ and $B\times B$ in all codimensions and after every finite
extension. No rational-cycle reconstruction beyond A28, and no
prior numerical--homological equivalence, enters this argument.

For completeness this also gives the strong formulation without
using A28's positivity proof. Frobenius is semisimple on the ambient
supersingular summands (some power is scalar), on $K$, and on $W$
(its characteristic polynomial has distinct roots), hence on all
cohomology of $B$ and $B\times B$. Poincare duality consequently
restricts to a nondegenerate pairing between the Tate invariants in
complementary codimensions, and the proved cycle surjectivity makes
this pairing a pairing of algebraic classes. Choose rational cycles
whose classes give a $\mathbf Q_\ell$-basis in one codimension and
complementary cycles with nonsingular intersection matrix. The matrix
is rational; pairing any further rational cycle with the complementary
cycles recovers rational coordinates in that basis. Thus the numerical
rank is the dimension of the Tate invariant space, and a numerically
trivial rational class is homologically zero. By semisimplicity this
dimension is the multiplicity of the eigenvalue $|k'|^r$ in
$H^{2r}(X)$; the zeta-function trace formula and purity of weights
identify it with the pole order at $s=r$, since other cohomological
degrees have different weights and cannot cancel this factor.
\end{proof}

\paragraph{Scope.}
The corollary concerns only the closed axes satisfying the stronger
regular-Frobenius selector of A28. It asserts neither the Tate
conjecture for every axis nor any analogous statement for all
higher powers. It supplies no arithmetic Chow finite generation
or Abel--Jacobi injectivity. It makes no claim of historical priority
or of previously unknown status for the examples.

% END work/arithmetic/axis_square_tate.tex

% BEGIN work/scan_w/axis_all_powers_torus.tex
% C44_POS /root/scan_w. A30 reserved by root.
% Published maximal-torus input; explicit tensor and geometric-form argument.
\section{Maximal Frobenius tori and every power of a selected closed axis}
\label{sec:c44-a30}

\begin{theorem}[C44-A30: one dense family with all-power positivity]
Use the family, good supersingular primes, and notation of C44-A13
and C44-A28. Thus $A=E^6$, $E:y^2=x^3-x$,
$B=B_u\subset A$ is a smooth fourfold axis of a genuine Lefschetz
pencil, and
\[
 H^*(B)=U\perp K\perp W,\qquad
 K=\mathbf Q_\ell w,\qquad
 \int_Bw^2=\beta>0,\qquad d=\dim W=3600t^6-166.
 \tag{A30.1}
\]
The summands $K,W$ occur only in primitive degree four; the
projectors are the actual algebraic projectors of C44-A28.
There is a Zariski-dense set of closed parameters for which
the cyclic group generated by normalized Frobenius
\[
                 \Phi=q^{-2}\operatorname{Frob}_q^*|W
 \tag{A30.2}
\]
has Zariski closure a maximal torus of $\operatorname{SO}(W)$.
Its eigenvalues may consequently be ordered as
\[
 \lambda_1,\lambda_1^{-1},\ldots,\lambda_{d/2},\lambda_{d/2}^{-1},
 \qquad
 \prod_i\lambda_i^{n_i}\text{ a root of unity }
                   \Longrightarrow n_1=\cdots=n_{d/2}=0.
 \tag{A30.3}
\]
For every such closed axis, every $N\geq1$, every $\ell\ne p$,
and every $\rho_1,\ldots,\rho_N\in\mathbf Q_{>0}$, the Hodge
standard conjecture holds for
\[
       \left(B^N,\ h_N=\sum_{i=1}^N\rho_i
                                  \operatorname{pr}_i^*h_B\right).
 \tag{A30.4}
\]
The Lefschetz standard conjecture and equality of numerical and
$\ell$-adic homological equivalence hold on every $B^N$.
The Tate conjecture holds in every codimension on every $B^N$
over every finite extension of its field of definition,
including equality with the corresponding pole order.

The same selected axis works for all $N$; the choice does not
depend on a tensor degree. No 032/Milne-2026 assumption or
assertion that $W$ is of abelian type is used. Positivity is
asserted for the displayed product polarizations, not for
every ample class on $B^N$.
\end{theorem}

\begin{proof}
\emph{1. Exact hypotheses of the maximal-torus theorem.}
Fix an auxiliary coefficient prime $\ell\ne p$ and consider
the normalized residual lisse sheaf $\mathcal W_\ell$ on the
smooth parameter curve of C44-A13. Its geometric monodromy
contains $\operatorname{SO}_d$ in its standard representation.
Thus it is absolutely irreducible, hence semisimple, as an
arithmetic representation. Arithmetic monodromy is contained
in $\operatorname{O}_d$, since the normalized middle Poincare
form is preserved. Its identity component is exactly
$\operatorname{SO}_d$.

The sheaf is $\mathbf Q$-rational: its closed-point Frobenius
polynomial is the quotient of the common middle polynomial
by the ambient polynomial and the polynomial of $K$.
It is pure of weight zero. It satisfies all three arithmetic
boundedness hypotheses of Chin's Theorem~5.7, as follows.
If $\alpha$ is an unnormalized degree-four eigenvalue at
a point of residue cardinality $q_u$, then both $\alpha$
and $q_u^4/\alpha$ are algebraic integers. The first statement
is Weil integrality, and the second follows because
$q_u^4/\alpha$ is a complex conjugate of $\alpha$ by purity.
Consequently $\lambda=\alpha/q_u^2$ is a unit at every finite
place away from $p$, which is Chin's \emph{plainness of
characteristic $p$}, and at every place $v$ above $p$,
\[
                  -2\leq\frac{v(\lambda)}{v(q_u)}\leq2.
 \tag{A30.5}
\]
The valuation bound is uniform.

Put $b=b_4(B)=d+496$, constant in the smooth family.
Katz--Messing identifies the degree-four crystalline and
$\ell$-adic characteristic polynomials. By the classification
of isocrystals, the slopes of a rank-$b$ isocrystal have
denominators at most $b$. Thus all the ratios in (A30.5)
lie in $(1/b!)\mathbf Z$. This is exactly the smooth-proper
argument in the proof of Chin's Lemma~3.3(2), author
pp.~8--9. It applies at every $p$-adic embedding because
the rational polynomial includes all conjugate roots.
Alternatively Chin's Proposition~4.5 gives uniform bounds
for an algebraic lisse sheaf on a smooth curve.
These checks do not assume Frobenius semisimplicity.

Chin's Theorem~5.7 therefore supplies a conjugation-stable
Zariski-open dense subset of arithmetic monodromy's identity
component such that the cyclic closure of any closed-point
Frobenius in it is itself a maximal torus. The selected
Frobenius is in particular semisimple. The precise hypotheses
are semisimplicity of the lisse sheaf, number-field rationality,
purity, plainness away from $p$, and uniform valuation and
denominator bounds; all were verified above.

This controls every multiplicative relation at once.
Indeed, in a splitting field of degree at most $d!$, the
$p$-adic valuation vectors of normalized roots have entries
in the finite set $[-2,2]\cap(1/b!)\mathbf Z$.
An integer relation has zero valuation at every place if
and only if its product is a root of unity: plainness handles
the other finite places, purity handles the archimedean
places, and Kronecker's theorem handles the resulting unit.
There are therefore only finitely many possible saturated
relation lattices and diagonal Frobenius subtori. This
is the weight-zero instance of the finite-torus argument
in Chin's Theorems~5.5 and~5.7. It is not an intersection
of countably many Chebotarev conditions.

Intersect the supplied open with the compact arithmetic
monodromy image. Zariski density gives a nonempty
conjugation-stable $\ell$-adic open subset detected in a
finite quotient. C44-A12's constant-field-correct
Chebotarev argument yields closed points there.
The same reasoning works on every nonempty open of the
parameter curve, proving the asserted Zariski density.

The weights of a maximal torus in the standard representation
of $\operatorname{SO}_d$ are a basis and its negatives,
$e_1,-e_1,\ldots,e_{d/2},-e_{d/2}$.
If a nontrivial product in (A30.3) were torsion, a nonzero
multiple of its character would vanish on the cyclic
closure, contradicting maximality. This proves (A30.3).
Conversely that independence makes the cyclic closure the
whole standard maximal torus. The condition depends only
on the common rational polynomial, hence holds for every
$\ell\ne p$, and it persists under every positive power
of Frobenius. It implies the simpler selector of C44-A28.

\emph{2. Rational pair tensors and a positive linear-algebra model.}
Fix a selected closed axis, and put
\[
 E_0=\mathbf Q[\Phi]=\mathbf Q[T]/P(T),\qquad
 \tau=T\bmod P(T),\qquad \tau^\dagger=\tau^{-1}.
 \tag{A30.6}
\]
C44-A28 proves that $E_0$ is an etale product of CM fields
and that $\dagger$ is complex conjugation on each factor.
For $a\in E_0$, let $t_a\in W\otimes W$ correspond to the
endomorphism $a$ under
\[
 u\otimes v\longmapsto
       \bigl(x\longmapsto q_W(v,x)u\bigr),
       \qquad q_W(x,y)=\int_Bxy.
 \tag{A30.7}
\]
Each $t_a$ is the class of an actual rational algebraic
correspondence on $B^2$, projected by $p_W\otimes p_W$.
The tensors $t_1,t_\tau,\ldots,t_{\tau^{d-1}}$ are a
rational basis of these pair classes.

Use the auxiliary rational space
\[
 V_0=E_0,\qquad
 q_0(x,y)=\operatorname{Tr}_{E_0/\mathbf Q}(xy^\dagger),
                   \qquad F_0(x)=\tau x.
 \tag{A30.8}
\]
Its real form is positive definite, and $F_0$ is an
isometry with characteristic polynomial $P$.
This is only a model for the pairings, not an asserted
Betti realization or extending motive of $W$.
After scalar extension to $\overline{\mathbf Q}_\ell$
there is an isometry
\[
                 (V_0,q_0,F_0)\simeq(W,q_W,\Phi).
 \tag{A30.9}
\]
Both operators have the same distinct eigenvalues.
Eigenlines of eigenvalues $\lambda,\mu$ are orthogonal
unless $\lambda\mu=1$, and the pairing on each pair of
reciprocal eigenlines is nonzero by nondegeneracy.
Rescale one vector of each reciprocal pair to match
these pairings. This constructs (A30.9), intertwines
every element of $E_0$, and carries its model tensor
for $a$ to $t_a$.

For $s\geq0$, pair $2s$ ordered slots, place one of
$t_{\tau^j}$, $0\leq j<d$, on each pair, and allow
all permutations of slots. The resulting finite collection
consists of actual rational algebraic classes on $B^{2s}$;
call them pair tensors of length $2s$.
Construct the identical collection in the model (A30.8).
By (A30.9), their complete Gram matrices coincide over
$\overline{\mathbf Q}_\ell$, hence as rational matrices:
the actual entries are rational intersection numbers.
The model form $q_0^{\otimes 2s}$ is positive definite
over $\mathbf R$. Choose pair tensors giving a rational
basis of their model span. Their actual images are
linearly independent and have the identical positive
rational Gram matrix. All linear relations among the
collections agree after scalar extension by (A30.9);
no inference from an isotropic vector alone is used.

\emph{3. Every invariant tensor is generated by pairs.}
Let $T_\Phi$ be the cyclic Frobenius torus and choose
eigenvectors $v_i^+,v_i^-$ of weights $e_i,-e_i$ over
a splitting field. A tensor word in $W^{\otimes b}$
has weight zero exactly when it contains equally
many $v_i^+$ and $v_i^-$ for each $i$.
There are no such words when $b$ is odd. When $b=2s$,
each such word can be partitioned into reciprocal pairs.
In degree two the invariant tensors identify via
(A30.7) with the centralizer of $\Phi$, which is
$E_0\otimes\overline{\mathbf Q}_\ell$.
Thus both oriented pair tensors
$v_i^+\otimes v_i^-$ and $v_i^-\otimes v_i^+$
are linear combinations of the $t_{\tau^j}$.
Replacing the pairs in a balanced word gives
\[
 (W^{\otimes 2s})^{T_\Phi}
       =\operatorname{span}_{\mathbf Q_\ell}
           \{\text{pair tensors of length }2s\}.
 \tag{A30.10}
\]
Connectedness of the cyclic closure implies that the
fixed space of every positive power of Frobenius is
this same space. In particular every rational algebraic
class on this projected block lies in (A30.10), even
if its defining cycle is over a finite extension.

Choose a model rational basis $z_1,\ldots,z_m$ from
the pair tensors. It is an actual $\mathbf Q_\ell$-basis
of (A30.10), and
\[
                   G_{ij}=\int_{B^{2s}}z_i z_j
 \tag{A30.11}
\]
is rational and positive definite. The intersections
of an arbitrary actual rational algebraic class $z$
with the $z_j$ are rational. Inverting $G$ recovers
rational coefficients for $z$ in this basis.
Hence all rational algebraic classes on $W^{\otimes2s}$
are exactly the rational pair span, and their actual
tensor form $q_W^{\otimes2s}$ is positive definite.
This proves rationality and positivity without a prior
Tate or numerical-equivalence conjecture. For $s=0$
the statement is the unit class with form $1$.

\emph{4. Ambient powers and their actual primitive forms.}
The map $(\kappa^a)_*$ injects $U^{\otimes a}$,
degree by degree, into $H^*(A^a)$ with shift $4a$.
It has an algebraic left inverse on its image.
For a single factor in degree at most four, invert
$L_A^2$ on its Lefschetz strings and restrict back.
Above four the Gysin is an isomorphism onto the
corresponding ambient degree, dual to weak Lefschetz;
its inverse is given by the algebraic Lefschetz inverse
(A28.7), restriction, and the corresponding raising
operator on $B$. Tensor these degreewise maps.
All these constructions are relative over the fixed
smooth embedded trait lift.

Every even cohomology group of every supersingular
power $E^M$ is spanned by divisor-product classes,
by the unconditional argument of C44-A28.
Apply the algebraic left inverse just described to
a spanning set on $A^a$, and then the algebraic
primitive projector. Thus every even primitive
ambient space is spanned over $\mathbf Q_\ell$ by
actual rational algebraic classes.

For positive factor weights, let
$h_{U,a}=\sum_{i=1}^a\rho_i h_i$.
The actual form on $P^{2r}(B^a)\cap U^{\otimes a}$,
$0\leq r\leq2a$, is
\[
                 q_{U,a,r}(x,y)
       =(-1)^r\int_{B^a}xyh_{U,a}^{\,4a-2r}.
 \tag{A30.12}
\]
It is positive on its rational algebraic classes.
To verify the geometric form, its primitive projector
and inverse pairing are algebraic and relative by the
tensor Lefschetz operators. The relative injection
$(\kappa^a)_*$ maps it into
$H^{2r+4a}(A^a)(2a)$.
Choose on that target the direct sum of its signed
primitive polarization forms with the Tate normalization,
call it $Q_A$, and pull it back to $q_1$.
On the complex generic fibre both $q_1$ and (A30.12)
polarize the same primitive Hodge structure.
Their ratio $c=q_{U,a,r}^{-1}q_1$ is an algebraic
relative endomorphism with positive real spectrum.
Its characteristic polynomial is constant along the
trait, and $c^{-1}$ is a rational polynomial in $c$.
On special rational algebraic classes, $q_1$ is
positive by the unconditional Lefschetz-class theorem
on $E^{6a}$. For a given class, its finite rational
cyclic space under $c$ consists of algebraic classes
and is stable under $c^{-1}$. On it, $c$ is
$q_1$-self-adjoint with positive eigenvalues, so
$q_{U,a,r}(z,z)=q_1(c^{-1}z,z)>0$ for $z\ne0$.
This is precisely the checked positive-spectrum argument
of C44-A08 and C44-A28 with $a$ factors. It includes
the odd-times-odd Kunneth contributions.

We obtain a rational cycle basis of every ambient
even primitive space with positive definite rational
Gram matrix. Every other actual rational algebraic
class has rational coordinates by pairing with that
basis. For $a=0$ take the unit space with form $1$.

\emph{5. Every primitive block on every power.}
The self-adjoint algebraic projectors in (A30.1)
commute with the raising and lowering operators.
Their products decompose each primitive part of
$B^N$ orthogonally. After reordering factors a block is
\[
             W^{\otimes b}\otimes K^{\otimes c}
                                  \otimes U^{\otimes a},
                         \qquad a+b+c=N.
 \tag{A30.13}
\]
Since $W$ and $K$ have degree four and zero raising
operator, the even primitive part of (A30.13) is
\[
 W^{\otimes b}\otimes K^{\otimes c}\otimes
              \bigl(P^{2r}(B^a)\cap U^{\otimes a}\bigr),
                       \qquad 0\leq r\leq2a,
 \tag{A30.14}
\]
in total codimension $2(b+c)+r$.
The remaining Lefschetz strings have their degrees
shifted by $4(b+c)$ and their centres shifted from
$4a$ to $4N$; a positive step on a string is not
primitive. Odd total degrees have no algebraic cycles.

After a finite extension, normalized even ambient
Frobenius and the action on $K$ are one.
Before that extension their eigenvalues are roots
of unity, which gives the same potential fixed spaces.
By (A30.3), every algebraic class in (A30.14) lies in
\[
 (W^{\otimes b})^{T_\Phi}\otimes K^{\otimes c}\otimes
              \bigl(P^{2r}(B^a)\cap U^{\otimes a}\bigr).
 \tag{A30.15}
\]
For odd $b$ this is zero. For $b=2s$, it is spanned
over $\mathbf Q_\ell$ by exterior products of the
actual pair cycles, $w^{\otimes c}$, and an actual
ambient primitive cycle basis.

There is no unspecified scalar in its geometric form.
Any term of $h_N^{4a-2r}$ involving a factor occupied
by $W$ or $K$ vanishes. The surviving term is exactly
$h_{U,a}^{4a-2r}$, with coefficient one.
The degrees of $W,K$ are even, so they introduce
no additional Kunneth sign, and
$(-1)^{2(b+c)+r}=(-1)^r$.
The actual signed primitive form on (A30.15) is
\[
                  q_W^{\otimes b}\otimes
                         \beta^c\otimes q_{U,a,r}.
 \tag{A30.16}
\]
For even $b$ its Gram matrix on the displayed
exterior-product basis is rational and positive
definite. Pair an arbitrary actual rational
algebraic class with this basis and invert that
matrix to recover rational coefficients.
Thus the form is positive on the entire algebraic
block, not merely on the exhibited products.
Orthogonality proves (A30.4) for every $N$ and
every choice of positive factor weights.

\emph{6. Standard and Tate consequences.}
The algebraic Lefschetz operators of $B$ tensor
to give the Lefschetz standard conjecture on every
$B^N$ for these product polarizations.
Together with the Hodge standard conjecture,
Kleiman's Proposition~5-1 gives equality of numerical
and $\ell$-adic homological equivalence.

The same argument proves, rather than assumes,
Tate surjectivity. Frobenius on the ambient cohomology
is semisimple, since it comes from powers of the
supersingular elliptic curve. On $W$ it is semisimple
by (A30.3), and on $K$ it is scalar.
It is therefore semisimple in every cohomological
degree of every $B^N$.
After algebraic Lefschetz decomposition, its potential
even Tate spaces are exactly (A30.15); all are
spanned by actual cycle classes by Steps~3--5.

For a given finite field, express a Frobenius-fixed
class in these generators and choose one finite
extension defining the finitely many cycles involved.
Average over its finite Galois group, equivalently
push down and divide by the extension degree.
The resulting cycles are defined over the original
field and span its fixed space.
This proves Tate surjectivity over every finite
field of definition and every finite extension.
Semisimplicity identifies the invariant dimension
with the relevant pole order; numerical-homological
equality identifies it with the numerical cycle rank.
\end{proof}

\paragraph{Exact published input and prior-known scope.}
Chin, \emph{Independence of $\ell$ of monodromy groups},
\emph{J. Amer. Math. Soc.} 17 (2004), no.~3, 723--747,
is used through Theorem~5.7 and its proof, Corollary~5.8,
and the boundedness discussion above.
The checked journal-version author PDF is
\url{https://arxiv.org/pdf/math/0206147}.
Theorem~5.7 occupies author pp.~14--16,
Corollary~5.8 is on p.~16, Proposition~4.5 on p.~12,
and the smooth-proper denominator proof for
Lemma~3.3(2) on pp.~8--9.
The relevant statements and displayed proofs were read;
this does not claim an independent proof of Lafforgue's
theorem or the full isocrystal classification.
The direct geometric denominator check avoids needing
Lafforgue for that particular hypothesis.

The actual ambient positivity input is Milne's published
unconditional theorem for Lefschetz classes, with the
scope checked in C44-A28. The numerical-equivalence
consequence uses Kleiman, \emph{The Standard Conjectures}
(1994), Proposition~5-1, printed p.~17 of the checked
Mainz scan. Positive adjoint trace and semisimplicity
on the correspondence algebra of any $B^N$ also follow
from its Theorem~5-6(1) and proof, printed p.~19, using
the now-proved Hodge standard conjecture on $B^{2N}$.

Maximal Frobenius tori, reciprocal-pair generation of
invariant tensors, and CM trace forms are established
tools with published precedents. No historical novelty
claim for this family is inferred from an unsuccessful
search for an earlier explicit statement.
The contribution recorded here is the specified
application to the checked closed-axis family, with
its actual algebraic generators, rational reconstruction
and all the geometric primitive forms.

% END work/scan_w/axis_all_powers_torus.tex

% BEGIN work/scan_w/axis_all_ample_fourfold.tex
% C44_POS /root/scan_w. Scope corollary, not a separate progress-clock event.
\subsection*{C44-A30a. Every ample polarization on the selected fourfold itself}

\begin{corollary}[C44-A30a]
Let $B$ be any closed fourfold selected in C44-A30.
For every ample rational divisor class $L$ on $B$ and every
$\ell\ne p$, the Hodge standard conjecture holds for $(B,L)$.
This assertion concerns $B$ itself. For $B^N$ with $N>1$,
C44-A30 continues to assert positivity for its displayed
positive product polarizations only.
\end{corollary}

\begin{proof}
Write $N^j=N^j(B)_{\mathbf Q}$ for codimension-$j$ cycles
modulo numerical equivalence and $\rho_j=\dim_{\mathbf Q}N^j$.
C44-A30 gives numerical--homological equality in every
codimension for every $\ell\ne p$. Thus each $N^j$ is
identified with the rational algebraic cycle-class space
in that realization, compatibly with products.

For the original polarization $h$, the orthogonal
Lefschetz decomposition is
\[
                 N^2=P_h^2\perp hN^1.
\]
The first summand is positive definite by C44-A30.
On the second summand the middle intersection form is
\[
                (hD,hD')\longmapsto\int_B DD'h^2,
\]
which has signature $(1,\rho_1-1)$ by the divisor
Hodge index theorem. Consequently the total middle
intersection form on $N^2$ has signature
\[
                   (\rho_2-\rho_1+1,\ \rho_1-1).
                   \tag{A30a.1}
\]
This total form does not involve a polarization.

Now fix any ample rational $L$. The divisor Hodge
index theorem shows that the form
$\int_B DD'L^2$ on $N^1$ is nondegenerate with the
same signature $(1,\rho_1-1)$. It follows in particular
that multiplication by $L$ embeds $N^1$ into $N^2$.
The orthogonal complement of $LN^1$ is exactly
\[
             P_L^2=\ker(L:N^2\longrightarrow N^3).
                   \tag{A30a.2}
\]
Indeed, orthogonality means that $Lz$ pairs to zero
with every class in $N^1$, hence that $Lz=0$ modulo
numerical equivalence. By the already established
numerical--homological equality in codimension three,
this is also the cohomological primitive condition.
Subtracting the signature of $LN^1$ from (A30a.1)
proves that the actual middle primitive form on
$P_L^2$ is positive definite.

The codimension-one primitive form is positive by
the divisor Hodge index theorem, and the
codimension-zero form is positive because
$\int_B L^4>0$. These are all the primitive algebraic
degrees at or below the middle degree of a fourfold.
Clearing the positive denominator of $L$ changes each
form only by a positive factor. This proves the
assertion for every $\ell\ne p$.
\end{proof}

\paragraph{Exact source and scope.}
This is the fourfold argument of Ancona,
\emph{Standard conjectures for abelian fourfolds},
\emph{Invent. Math.} 223 (2021), 149--212,
Proposition~3.15 and its complete proof, author p.~12
in \url{https://arxiv.org/pdf/1806.03216v3}.
Despite the paper's title, that proposition is stated
for any fourfold under its Kunneth/Lefschetz assumptions
of Assumption~3.4. The homological conclusion here
uses the independently proved $D(B)$ of C44-A30;
no numerical--homological equivalence is silently
assumed. No continuity argument through irrational
ample classes and no assertion for arbitrary ample
classes on higher powers is used.

% END work/scan_w/axis_all_ample_fourfold.tex

% BEGIN work/arithmetic/fermat_geometric_torsion.tex
% C44_POS; /root/fermat_aj.  A geometric result, not closure of C44-F07.
\section{Geometric five-primary torsion and the surviving division problem}

Retain $J,e,R=\mathbf Z_{(5)},T$ from C44-F01 and put
$\overline J=J_{\overline{\mathbf Q}}$.  All statements below are
unconditional.  They use neither 032 nor a Milne transfer assertion.
The arithmetic Abel--Jacobi injection in the task remains open.

\begin{theorem}[C44-F23]\label{C44-F23}
For every integer $n\ge1$, restriction to the geometric function
field is injective on the indicated direct summand:
\[
 eH^3_{\mathrm{et}}(\overline J,\mathbf Z/5^n(2))
 \longrightarrow
 H^3(\overline{\mathbf Q}(J),\mathbf Z/5^n(2)).
 \tag{F23.1}
\]
Equivalently, $eN^1H^3_{\mathrm{et}}(\overline J,
\mathbf Z/5^n(2))=0$.  Consequently
\[
 \overline A:=e\operatorname{CH}^2(\overline J;R)
 \quad\hbox{has no nonzero five-primary torsion}.       \tag{F23.2}
\]

Choose any nonzero integral-at-five multiple $a_\eta\in A$ of
the carried class $\eta$.  There is an integer $v\ge0$ such that
\[
 \operatorname{res}(a_\eta)=5^r b,\quad b\in\overline A,
 \qquad\Longrightarrow\qquad r\le v.                 \tag{F23.3}
\]
Thus this one class does not become infinitely five-divisible even
after passing to the algebraic closure.  The assertion is not a
uniform bound for all geometric cycles.

Let
\[
 \overline B=eH^4_{\mathrm{et}}(\overline J,R(2)),\qquad
 \overline C=eH^3_{\mathrm{nr}}(\overline J,
                         \mathbf Q_5/\mathbf Z_5(2)),
\]
where the first group is \emph{\'etale motivic} cohomology.
There are canonical exact sequences of Galois modules
\[
 0\longrightarrow T/5^nT\longrightarrow\overline C[5^n]
       \longrightarrow\overline A/5^n\overline A
       \longrightarrow0                              \tag{F23.4}
\]
and
\[
 0\longrightarrow T\longrightarrow T_5\overline C
       \longrightarrow\widehat{\overline A}
       \longrightarrow0.                             \tag{F23.5}
\]
The completion on the right is nonzero by (F23.3).  In particular,
the geometric torsion vanishing does not identify the geometric
unramified Tate module with $T$, and does not prove
$T_5C=0$ for the arithmetic group $C$ of C44-F03.
\end{theorem}

\begin{proof}
\emph{The mod-five coniveau calculation.}
The Jacobian has good reduction over $\mathbf Z_5$ by C44-F01.
The carried isogeny of JAC-P17 is
$J\sim E_-^2\times E_+$, with $E_\pm:y^2=x^3\pm x$.
Direct counts give
\[
 \#E_-(\mathbf F_5)=8,\qquad
 \#E_+(\mathbf F_5)=4.
 \tag{F23.6}
\]
Their traces are $-2$ and $2$, so both special fibres are ordinary.
The isogeny extends over $\mathbf Z_5$; ordinarity is an isogeny
invariant.  Thus the special fibre $J_5$ is ordinary.

We use the construction in the proof of Bloch--Esnault,
\emph{The coniveau filtration and non-divisibility for algebraic
cycles}, Theorem~(1.2), not just its conclusion for the full
cohomology group.  Here the dimension is three, the coefficient
prime is five, and the absolute ramification index is one, so its
inequality $3<(5-1)/\gcd(1,5-1)$ applies.  The edge map in their
(1.2.6) is natural for automorphisms of the smooth model and gives
a nonzero map
\[
 \alpha:H^3_{\mathrm{et}}(J_{\overline{\mathbf Q}_5},\mathbf F_5)
 \longrightarrow H^0(J_{\overline{\mathbf F}_5},\Omega^3)(-3).
 \tag{F23.7}
\]
Their diagram (1.2.7), whose right vertical arrow is injective,
shows that $\alpha$ kills the kernel of restriction to the
geometric function field.

Every fourth root of unity is in $\mathbf Z_5$.  Hence the
automorphisms in $e$ act on this model and commute with (F23.7).
The three holomorphic characters from C44-F01 are
$(1,1),(1,2),(2,1)$.  Their sum is $(0,0)$ modulo four,
so $G=\mu_4^2$ acts trivially on the one-dimensional space
$H^0(J_5,\Omega^3)$.  Inversion acts there by $-1$.
These statements also follow by reducing the corresponding
free Hodge modules over $\mathbf Z_5$.  The idempotent $e$ therefore
acts as the identity on the target of (F23.7).  Since $\alpha\ne0$,
its restriction to $eH^3$ is nonzero.  It follows that
$eN^1H^3(\overline J,\mathbf F_5(2))$ is a proper subspace of
$eH^3(\overline J,\mathbf F_5(2))$.
Passage from $\overline{\mathbf Q}$ to
$\overline{\mathbf Q}_5$ causes no problem: any class supported
on a divisor remains supported after this extension; proper base
change identifies the cohomology groups.  This is also the argument
of Bloch--Esnault Corollary~(1.3).

The subspace just obtained is stable under
$G_{\mathbf Q}=\operatorname{Gal}(\overline{\mathbf Q}/\mathbf Q)$,
because Galois permutes its supports and $e$ is defined over
$\mathbf Q$.  At seven, the CM elliptic curve $E_+$ has trace zero
and Frobenius polynomial $X^2+7$.  JAC-P17 identifies the
rank-two summand with the cube of its Gaussian character.
The twisted polynomial on $T$ is therefore $X^2+1/7$.
Modulo five it is
\[
                         X^2+3,                      \tag{F23.8}
\]
which is irreducible since $2$ is not a square in $\mathbf F_5$.
Thus $T/5T$ has no nontrivial $G_{\mathbf Q}$-stable subspace.
Its proper coniveau-one subspace must be zero.  This proves
(F23.1) for $n=1$.

\emph{The higher coefficients and torsion cycles.}
For completeness, the induction in Bloch--Esnault Remark~(1.4)
works on this exact summand as well.  Put
$F=\overline{\mathbf Q}(J)$ and $H_n=eH^3(\overline J,
\mathbf Z/5^n(2))$.  The coefficient sequence gives
\[
 0\longrightarrow H_n\longrightarrow H_{n+1}
                 \longrightarrow H_1\longrightarrow0.
 \tag{F23.9}
\]
One can see this directly from C44-F01 and the freeness of $T$.
For the field $F$, Merkurjev--Suslin identifies
$H^2(F,\mu_{5^m}^{\otimes2})$ with $K^M_2(F)/5^m$.
The coefficient reduction from exponent $5^{n+1}$ to exponent
five is therefore surjective in degree two.  The coefficient
long exact sequence implies that
\[
 H^3(F,\mathbf Z/5^n(2))\longrightarrow
 H^3(F,\mathbf Z/5^{n+1}(2))
 \quad\hbox{is injective}.                           \tag{F23.10}
\]
If a class in $H_{n+1}$ restricts to zero, its reduction in $H_1$
is zero by the case already proved.  By (F23.9) it comes from
$H_n$, and (F23.10) says that this preimage also restricts to
zero.  Induction proves (F23.1) for every $n$.

The Bloch torsion-cycle map in codimension two is injective with
image $N^1H^3(\overline J,\mathbf Q_5/\mathbf Z_5(2))$.
This is the Merkurjev--Suslin theorem as used explicitly at the
start of Bloch--Esnault Proposition~(4.1).  The map and its image
are natural for automorphisms, so the same holds after applying
$e$.  Filtered direct limits are exact; (F23.1) shows that the
projected coniveau-one group with $\mathbf Q_5/\mathbf Z_5$
coefficients is zero.  This proves (F23.2).

\emph{Division over the algebraic closure.}
For every finite extension $L/\mathbf Q$, the representation
$T\otimes\mathbf Q_5$ has no $G_L$-invariants: at a good place
outside five its Frobenius eigenvalues have complex absolute
value $\operatorname{Norm}(v)^{-1/2}$, so none is one.
Thus $T^{G_L}=0$.  For any finite Galois extension $L/\mathbf Q$
one has the integral restriction isomorphism
\[
 H^1_{\mathrm{cont}}(\mathbf Q,T)
   \xrightarrow{\ \sim\ }
 H^1_{\mathrm{cont}}(L,T)^{\operatorname{Gal}(L/\mathbf Q)}.
 \tag{F23.11}
\]
Here is the coefficient-limit justification.  Apply the
inflation--restriction five-term sequence to $T/5^nT$.
The modules $(T/5^nT)^{G_L}$ are finite and form a
Mittag--Leffler system whose inverse limit is $T^{G_L}=0$.
For a fixed finite group, the inverse limits of their positive
cohomology groups are consequently zero.  Taking inverse limits
in the five-term sequence yields (F23.11); the relevant
$\varprojlim^1$ terms of these finite systems vanish.
This is precisely the argument in Bloch--Esnault
(4.1.1)--(4.1.3), printed p.~313.

Suppose $5^r b=\operatorname{res}(a_\eta)$ in $\overline A$.
By (F23.2), the equation implies $\sigma b=b$ for every
$\sigma\in G_{\mathbf Q}$.  The cycle and the finite number
of rational-equivalence relations needed below are defined
over a finite Galois extension.  Its integral Abel--Jacobi class
is invariant under that finite Galois group, after enlarging
the extension if necessary.  To justify invariance precisely,
the difference of a conjugate cycle and the original cycle
restricts to zero over $\overline{\mathbf Q}$, hence over some
finite extension.  Restriction on $H^1(-,T)$ is injective by
(F23.11), applied over the relevant ground field.  The difference
of their Abel--Jacobi classes is therefore already zero.
Equation (F23.11) descends the Abel--Jacobi class of $b$ to
$\gamma_b\in H^1_{\mathrm{cont}}(\mathbf Q,T)$.  Functoriality
and the same restriction injectivity give
\[
                    \alpha_5(a_\eta)=5^r\gamma_b.     \tag{F23.12}
\]
In this display $\alpha_5(a_\eta)$ denotes the integral class
whose rationalization is the carried nonzero regulator.
It is nonzero by AR-23.  By C44-F02,
$H^1_{\mathrm{cont}}(\mathbf Q,T)=
\varprojlim_nH^1(\mathbf Q,T/5^nT)$, so some finite coefficient
image of this class is nonzero.  If that happens at exponent
$5^{n_0}$, equation (F23.12) is impossible for $r\ge n_0$.
Taking $v=n_0-1$ proves (F23.3).

\emph{What the geometric unramified group still contains.}
The comparison sequence of C44-F03 over the algebraic closure
is $0\to\overline A\to\overline B\to\overline C\to0$.
The geometric group
$eH^3_{\mathrm{et}}(\overline J,R(2))$ is uniquely divisible
by five.  Indeed Colliot-Th\'el\`ene--Kahn Theorem~6.6(a)
identifies the obstruction to divisibility with the torsion
in $H^3(\overline J,\mathbf Z_5(2))$, which is zero;
part~(b) identifies its five-primary torsion with
$eH^2(\overline J,\mathbf Q_5/\mathbf Z_5(2))=0$.
The coefficient triangle and C44-F01 now give
\[
 \overline B/5^n\overline B=0,
 \qquad\overline B[5^n]=T/5^nT.                     \tag{F23.13}
\]
Applying the multiplication-by-$5^n$ snake lemma to the comparison
sequence, and using (F23.2), gives (F23.4).
Its transition maps are multiplication by five on the torsion
terms and reduction on the quotient $\overline A/5^n$.
The left system $T/5^nT$ has surjective transitions, so its
$\varprojlim^1$ is zero.  Taking inverse limits proves
(F23.5).  Finally (F23.3) means that the image of
$\operatorname{res}(a_\eta)$ is nonzero in some
$\overline A/5^{n_0}\overline A$, hence in its completion.
\end{proof}

\paragraph{Primary locators and the remaining arithmetic gate.}
Bloch--Esnault, Math.~Ann.~304 (1996), 303--314, was read in the
published twelve-page scan hosted by Esnault:
\url{https://page.mi.fu-berlin.de/esnault/preprints/helene/37-coniveau_filt_non_div.pdf}.
The exact inputs are Theorem~(1.2), printed p.~306; its proof,
(1.2.6)--(1.2.7) and Corollary~(1.3), p.~307;
Remark~(1.4), p.~308; and the torsion-cycle identification and
integral descent argument in Proposition~(4.1), pp.~312--313.
The proof of Theorem~(1.2), including its diagrams, and
Remark~(1.4) were visually checked in the original scan.
The needed part of Proposition~(4.1) was read directly.
Colliot-Th\'el\`ene--Kahn, \emph{Cycles de codimension 2 et
$H^3$ non ramifi\'e pour les vari\'et\'es sur les corps finis},
the 43-page author PDF, Theorem~6.6(a)--(b), printed p.~27,
is at
\url{https://www.imo.universite-paris-saclay.fr/~jean-louis.colliot-thelene/17avrilcorpsfinisCTKahn.pdf}.
It is used over the geometric field, within its stated scope;
their Theorem~6.3 with its cohomological-dimension-one hypothesis
is not used over a number field.

This theorem controls geometric torsion and the geometric division
of the one known class.  It proves neither arithmetic finite
generation nor a rank bound for the whole Chow group.  In
particular the exact quotient in (F23.5) cannot be dropped when
passing to Galois invariants.  The two arithmetic obstructions
of C44-F03 and the equivalent finite-generation gate of C44-F07
remain unresolved.  The new result is not counted as movement
of that finite-generation gate.

% END work/arithmetic/fermat_geometric_torsion.tex

% BEGIN work/arithmetic/fermat_geometric_diagnostic.tex
% C44_POS arithmetic obstruction addendum, /root/fermat_aj.
% This is a formal module construction, not an actual motive or variety.
% No progress-clock movement is claimed.

\section{A geometric-division diagnostic for the remaining arithmetic gate}

This addendum strengthens the formal diagnostic recorded after C44-F07.
It is \emph{not} a counterexample to any conjecture about algebraic cycles.
The modules below are constructed abstractly.  They are not asserted to
be Chow groups, unramified cohomology of a field, or the cohomology of a
motivic complex.  The purpose is to determine exactly which inference
cannot follow from the finite-coefficient realization and the additional
division properties established in C44-F23.

Put $R=\mathbf Z_{(5)}$, $S=\operatorname{Spec}\mathbf Z[1/10]$,
$G=\pi_1(S)$, and let $T$ be the actual rank-two lattice of C44-F01.
Write $W=T\otimes_{\mathbf Z_5}\mathbf Q_5/\mathbf Z_5$.
The carried statements C44-F02 and C44-F07 give
\[
 H^0(G,W)=0,\qquad
 L=H^1_{\rm cont}(G,T)\text{ finite free over }\mathbf Z_5.
\]
Choose a primitive generator $\xi$ of the saturated lattice in $L$
on the known rational Selmer line.  Thus
\[
 L\cap\mathbf Q_5\xi=\mathbf Z_5\xi.
 \tag{GD.1}
\]
No assertion that $\xi$ is the regulator of an integral cycle is needed
for this formal construction.

\begin{proposition}[C44-F24: formal diagnostic; no motivic counterexample]
\label{C44-F24}
For every $\lambda\in\mathbf Z_5$ there are discrete continuous
$R[G]$-modules $\overline A$, $\overline B$, $\overline C$ with the
following properties.
\begin{enumerate}
\item
$\overline A=R^2$ with trivial $G$-action,
$\overline B$ is five-divisible,
$\overline B[5^n]=T/5^n$, and
\[
 0\longrightarrow\overline A\longrightarrow\overline B
   \longrightarrow\overline C\longrightarrow0.
 \tag{GD.2}
\]
In particular, $\overline A$ has no five-primary torsion,
all five-divisions in it are unique when they exist, and
$\eta=(1,0)$ is not divisible by five.
\item
If $A=\overline A^G$, $B=\overline B^G$, and
$C=\overline C^G$, then
\[
 A=R^2,\qquad
 B=\{(r,s)\in\mathbf Q^2:r+\lambda s\in\mathbf Z_5\},
 \qquad 0\longrightarrow A\longrightarrow B
       \longrightarrow C\longrightarrow0.
 \tag{GD.3}
\]
Here $C\simeq\mathbf Q_5/\mathbf Z_5$ as an $R$-module,
$B$ is torsion-free and not finitely generated over $R$, while
$\widehat B\simeq\mathbf Z_5$.
\item
There are compatible exact sequences
\[
 0\longrightarrow T/5^n\longrightarrow\overline C[5^n]
   \longrightarrow\overline A/5^n\longrightarrow0,
 \tag{GD.4}
\]
and
\[
 0\longrightarrow T\longrightarrow T_5\overline C
   \longrightarrow\mathbf Z_5^2\longrightarrow0.
 \tag{GD.5}
\]
The connecting map on $G$-invariants of (GD.5) is
\[
 \delta:\mathbf Z_5^2\longrightarrow L,
 \qquad (a,b)\longmapsto(a+\lambda b)\xi.
 \tag{GD.6}
\]
Consequently
\[
 (T_5\overline C)^G=T_5C
   =\mathbf Z_5(-\lambda,1)\ne0.
 \tag{GD.7}
\]
\item
The formal complex $\mathcal M=\overline B[-4]$ has
\[
 \mathcal M\mathbin{\otimes}^{\mathbf L}_{R}R/5^n
       \simeq(T/5^n)[-3]
 \tag{GD.8}
\]
as complexes of discrete continuous $G$-modules.  Thus its entire
finite-coefficient derived Galois realization agrees with the
single-row realization used in C44-F01--F02, not merely with a
dimension of a Selmer group.
\end{enumerate}
The realization on the arithmetic source is
\[
 A\otimes_R\mathbf Q_5=\mathbf Q_5^2
   \longrightarrow\mathbf Q_5\xi,
 \qquad (a,b)\longmapsto(a+\lambda b)\xi.
 \tag{GD.9}
\]
If $\lambda\notin\mathbf Q$, its restriction to $A\otimes_R\mathbf Q$
is injective, whereas (GD.9) has the nonzero kernel
$\mathbf Q_5(-\lambda,1)$.  If $\lambda=0$, even its rational
restriction has the nonzero kernel $\mathbf Q(0,1)$.
\end{proposition}

\begin{proof}
Represent $\xi$ by a continuous cocycle
$c:G\longrightarrow T$.  On the free underlying module
$E=T\oplus\mathbf Z_5$ define
\[
 g(t,u)=(gt+u c(g),u).
\]
This gives the extension $0\to T\to E\to\mathbf Z_5\to0$
with class $\xi$.  Rationalizing and quotienting by $T$ gives
\[
 0\longrightarrow W\longrightarrow E_{\mathbf Q_5}/T
   \longrightarrow\mathbf Q_5\longrightarrow0.
 \tag{GD.10}
\]
Define $\phi(r,s)=r+\lambda s$ on $\mathbf Q^2$ and let
$\overline B$ be the pullback of (GD.10) along $\phi$.
We give this module the discrete topology.  For each of its elements,
the difference $g x-x$ is the reduction modulo $T$ of a continuous
map into $T\otimes\mathbf Q_5$.  The submodule $T$ is open there;
therefore the stabilizer of $x$ is open.  This verifies the required
continuity of the discrete action.  We have
\[
 0\longrightarrow W\longrightarrow\overline B
   \longrightarrow\mathbf Q^2\longrightarrow0.
 \tag{GD.11}
\]
Both outer modules are five-divisible, so $\overline B$ is
five-divisible: first divide its image in $\mathbf Q^2$, then
correct the difference by a division in $W$.  Since $\mathbf Q^2$
is torsion-free, $\overline B[5^n]=W[5^n]=T/5^n$.

For $v=(r,s)\in R^2$ one has $\phi(v)\in\mathbf Z_5$.
The element $((0,\phi(v))\bmod T,v)$ of the pullback is $G$-fixed,
because its Galois difference is $\phi(v)c(g)\in T$.
These elements define the injection $\overline A=R^2\to\overline B$.
Set $\overline C=\overline B/\overline A$.
All the assertions about torsion and division in $\overline A$,
including primitivity of $\eta$, are now literal statements about
the free $R$-module $R^2$.  In particular every existing geometric
division in this model is already fixed by $G$; no descent of such
a division requires passing to a larger field.

The connecting map from $\mathbf Q^2$ in (GD.11) sends $v$ to
the image of $\phi(v)\xi$ in $H^1(G,W)$.  By (GD.1), its kernel
is exactly the displayed subgroup $B$ in (GD.3).  Since $W^G=0$,
projection to $\mathbf Q^2$ identifies $\overline B^G$ with that
subgroup.  Furthermore
\[
 H^1(G,\overline A)=\operatorname{Hom}_{\rm cont}(G,R^2)=0:
\]
a continuous homomorphism from the compact group $G$ to the
discrete torsion-free group $R^2$ has finite image and hence is zero.
Taking invariants of (GD.2) proves the remaining exact sequence
in (GD.3).

There is a canonical isomorphism $\mathbf Q/R\simeq
\mathbf Q_5/\mathbf Z_5$.  Injectivity follows from
$\mathbf Q\cap\mathbf Z_5=R$; surjectivity follows by representing
any coset by its finite negative-five-adic expansion.  Therefore
$B/A$ is the kernel of
\[
 (\mathbf Q_5/\mathbf Z_5)^2\longrightarrow
       \mathbf Q_5/\mathbf Z_5,
 \qquad (x,y)\longmapsto x+\lambda y.
\]
This is a surjective map whose kernel is the graph
$\{(-\lambda y,y):y\in\mathbf Q_5/\mathbf Z_5\}$.
It proves $C\simeq\mathbf Q_5/\mathbf Z_5$.
In particular $B$ cannot be finitely generated over $R$, since
its quotient $C$ is not finitely generated.

The map $\phi:B\to\mathbf Z_5$ induces an isomorphism
$B/5^nB\simeq\mathbf Z/5^n$.  It is surjective because
$R(1,0)\subset B$.  If $\phi(v)\in5^n\mathbf Z_5$, then
$v/5^n\in B$, so its kernel is precisely $5^nB$.
This proves the completion assertion.  Torsion-freeness of $B$
follows from its inclusion in $\mathbf Q^2$.

Apply multiplication by $5^n$ to (GD.2).  The facts that
$\overline A[5^n]=0$ and $\overline B/5^n=0$ give (GD.4).
The transition on its left and middle terms is multiplication
by five, and on its right term is reduction.  On the left it
identifies with the surjective reduction $T/5^{n+1}\to T/5^n$.
Its first derived inverse limit therefore vanishes, so inverse
limits give (GD.5).

For $v\in R^2$, use the division
\[
 m_n=((0,\phi(v)/5^n)\bmod T, v/5^n)\in\overline B
\]
of its displayed lift in $\overline A$.  Its Galois difference is
$\phi(v)c(g)/5^n\bmod T$.  Under
$W[5^n]\simeq T/5^n$ this is the reduction of
$\phi(v)c(g)$.  Thus the connecting cocycle in (GD.5) on
$v\in R^2$ is $\phi(v)c$.  Passage through all finite reductions
gives (GD.6) on $\widehat{\overline A}=\mathbf Z_5^2$.
Since $T^G=0$ and $\xi$ is a nonzero vector in the free lattice
$L$, the kernel is exactly $\mathbf Z_5(-\lambda,1)$.
Invariants commute with inverse limits, and
$\overline C[5^n]^G=C[5^n]$, proving both equalities in (GD.7).

Finally, derived reduction of $\overline B$ is represented by
$[\overline B\xrightarrow{5^n}\overline B]$ in degrees $-1,0$.
Surjectivity of multiplication by $5^n$ and its kernel $T/5^n$
give (GD.8) after the shift $[-4]$.
Applying derived continuous invariants preserves this
identification.  The compatible connecting computations above
also give (GD.9).  If $\lambda\notin\mathbf Q$, a rational
solution of $a+\lambda b=0$ must have $a=b=0$; over
$\mathbf Q_5$ its kernel is the stated line.  For $\lambda=0$
the rational kernel is instead already visible.  This proves
every assertion.
\end{proof}

\paragraph{What this tests, and what remains missing.}
Taking $\lambda$ to be the root of $X^2-6$ in $\mathbf Z_5$
congruent to one modulo five recovers the irrational example of
C44-F05.  In it, the arithmetic Chow analogue $A$ is already
finitely generated, the geometric analogue $\overline A$ is
five-torsion-free and five-adically separated, and the known class
is primitive even geometrically.  Unique division and Galois
descent are therefore not the missing steps.  The missing control
is the saturation of $A$ in the larger invariant group $B$:
the quotient $B/A$ retains a divisible direction whose Tate module
is the irrational line in (GD.7).  Primitivity of one class bounds
division only on that specified line; it supplies no bound for
this quotient or for every rational plane appearing in C44-F04.

All rational local Bloch--Kato conditions imposed on the known
Selmer line are preserved, since every displayed extension class
is a scalar multiple of $\xi$.  This observation concerns only
rational realization.  The construction supplies no actual
algebraic cycles, purity, specialization, Gersten complex,
field-valued unramified cohomology, integral local comparison,
or proper arithmetic model.  In particular it does not assert
that the actual group in C44-F07 has the displayed form.
It only proves that the listed formal realization and division
properties do not establish its finite generation.

% END work/arithmetic/fermat_geometric_diagnostic.tex

% BEGIN work/arithmetic/fermat_actual_descent.tex
% C44_POS, /root/fermat_aj. Exact actual-group descent statement.
% This does not discharge the arithmetic finite-generation gate.

\section{Actual integral descent and the remaining $K_2$ kernel}

Retain $R=\mathbf Z_{(5)}$, $J$, $e$, $A$, $B$, $C$ and $T$
from C44-F01--F03.  Put $\overline J=J_{\overline{\mathbf Q}}$,
$G_{\mathbf Q}=\operatorname{Gal}(\overline{\mathbf Q}/\mathbf Q)$,
and
\[
 \overline A=e\operatorname{CH}^2(\overline J;R),\quad
 \overline B=eH^4_{\rm et}(\overline J,R(2)),\quad
 \overline C=eH^3_{\rm nr}(\overline J,\mathbf Q_5/\mathbf Z_5(2)).
\]
Here the middle group is still \emph{\'etale motivic} cohomology.
The superscript $G_{\mathbf Q}$ denotes invariants of the
actual discrete Galois modules, not invariants of a substituted
realization.

\begin{proposition}[C44-F25]\label{C44-F25}
Restriction is an isomorphism
\[
                 B\xrightarrow{\sim}\overline B^{G_{\mathbf Q}}.
 \tag{F25.1}
\]
The actual integral Chow descent defect
\[
       D=\overline A^{G_{\mathbf Q}}/A
 \tag{F25.2}
\]
is well defined and is a five-primary torsion group.  There is a
canonical identification
\[
             D=\ker(C\longrightarrow\overline C).
 \tag{F25.3}
\]
For every $n\ge1$ there is also a canonical identification
\[
 \ker\bigl[A/5^nA\longrightarrow
       (\overline A/5^n\overline A)^{G_{\mathbf Q}}\bigr]
              =D[5^n].
 \tag{F25.4}
\]
Consequently
\[
 \ker\bigl[\widehat A\longrightarrow
             \widehat{\overline A}^{\,G_{\mathbf Q}}\bigr]=T_5D.
 \tag{F25.5}
\]
In particular $D=0$ is equivalent to injectivity of the single
map in (F25.4) for $n=1$.

There is an explicit $K$-theoretic expression for this defect.
Set $K=\mathbf Q(i)$ and $G_K=\operatorname{Gal}(\overline{\mathbf Q}/K)$.
Let
\[
 M_0=\bigl(K_2(\overline{\mathbf Q}(J))/K_2(\overline{\mathbf Q})\bigr)
             \otimes_{\mathbf Z}R,
 \qquad
 M_1=\left(\bigoplus_{x\in\overline J^{(1)}}
                    \overline{\mathbf Q}(x)^\times\right)
             \otimes_{\mathbf Z}R.
\]
The tame symbols define $M_0\to M_1$, and
\[
 D\simeq
 \left(
 e\ker\left[H^2(G_K,M_0)\longrightarrow H^2(G_K,M_1)\right]
 \right)^{\operatorname{Gal}(K/\mathbf Q)}.
 \tag{F25.6}
\]
The projector annihilates constant $K_2$, so $eM_0$ can equivalently
be written $eK_2(\overline{\mathbf Q}(J))\otimes R$.

None of the groups in (F25.2), (F25.5), or (F25.6) is asserted to
vanish.  Even their vanishing would not by itself prove C44-F07.
\end{proposition}

\begin{proof}
Let $\mathcal E=eR\Gamma_{\rm et}(\overline J,R(2))$ with its
natural continuous Galois action.  The coefficient triangle and
C44-F01 show that multiplication by five on $H^q(\mathcal E)$
is bijective for $q\le2$: both adjacent finite-coefficient
groups vanish.  Since these are $R$-modules, they are
$\mathbf Q$-vector spaces.  The same conclusion in degree three
follows from Colliot-Th\'el\`ene--Kahn, Theorem~6.6(a)--(b),
as already used in C44-F23.  Specifically, the torsion in
$H^3(\overline J,\mathbf Z_5(2))$ is zero and
$eH^2(\overline J,\mathbf Q_5/\mathbf Z_5(2))=0$.
Thus $eH^3_{\rm et}(\overline J,R(2))$ is uniquely five-divisible.

Higher continuous Galois cohomology of a discrete rational
vector-space module is zero.  One can see this by passing to
finite quotients fixing any finite collection of elements and
averaging by the order of the finite quotient.  Consequently the
projected Hochschild--Serre spectral sequence has no term with
$p>0,q\le3$.  The degree-four edge is therefore precisely
(F25.1).  Equivalently, to avoid an unbounded-below convergence
issue, truncate $\mathcal E$ at degree three.  Its lower
truncation is quasi-isomorphic to its rationalization.  Here
$\operatorname{cd}_5(G_{\mathbf Q})\le2$: Milne,
\emph{Arithmetic Duality Theorems}, second edition (2006),
Chapter~I, Theorem~4.10(c), printed p.~57, identifies higher
cohomology of finite modules with the real-place terms.
Those terms vanish for five-primary modules; passage to filtered
colimits gives the stated bound.  The rational/torsion exact
sequences therefore bound the cohomological dimension on
discrete $R$-modules by three: higher cohomology of a rational
module vanishes, that of a five-primary torsion module vanishes
above degree two, and passage to a torsion-free $R$-module costs
at most one degree.  Thus derived invariants have bounded
amplitude, which justifies the truncation calculation even if
$\mathcal E$ is unbounded below.  On its rational lower
truncation, derived Galois invariants are exact and have
no cohomology in degrees four or five.  The upper truncation
starts in degree four, where its invariant edge is
$H^4(\mathcal E)^{G_{\mathbf Q}}$.  The truncation triangle
gives the same isomorphism.

Apply restriction to the fundamental codimension-two sequences
of C44-F03 and C44-F23.  The resulting diagram has exact top row
\[
 0\longrightarrow A\longrightarrow B\longrightarrow C\longrightarrow0
\]
and lower row exact at its first two nonzero terms
\[
 0\longrightarrow\overline A^{G_{\mathbf Q}}
  \longrightarrow\overline B^{G_{\mathbf Q}}
  \longrightarrow\overline C^{G_{\mathbf Q}}.
\]
Since the middle vertical arrow is an isomorphism, the left
vertical arrow is injective and the snake-lemma calculation gives
(F25.3).  Surjectivity of the lower row at its last term is not
needed and is not asserted.

For any class in $\overline A^{G_{\mathbf Q}}$, choose a finite
extension over which a representative is defined.  Its norm
restricts to the degree of this extension times the original
class, because that class is Galois invariant.  Thus its image
in $D$ is torsion.  All primes other than five are invertible
in $R$, so $D$ is five-primary.

By C44-F23, $\overline A$ has no five-primary torsion.  If a
$G_{\mathbf Q}$-fixed class $a$ satisfies $a=5^n b$ in
$\overline A$, then
\[
                  5^n(gb-b)=0
\]
for every $g$, so $b$ is fixed as well.  Therefore
\[
 \overline A^{G_{\mathbf Q}}/5^n
       \longrightarrow(\overline A/5^n)^{G_{\mathbf Q}}
\]
is injective.  The multiplication-by-$5^n$ sequence for
$0\to A\to\overline A^{G_{\mathbf Q}}\to D\to0$
now proves (F25.4).  Under this identification the transition on
the kernels is multiplication by five on $D[5^n]$.
Inverse limits are left exact and invariants commute with
inverse limits, giving (F25.5).  A nonzero five-primary torsion
group contains a nonzero element killed by five.  This proves
the equivalence with the single-modulus injectivity statement.

It remains to justify the field and coefficient scope of (F25.6).
The number field $K=\mathbf Q(i)$ has cohomological dimension
two.  Colliot-Th\'el\`ene--Kahn, Proposition~6.2(b), applies
to $J_K$ and identifies the kernel of restriction on unramified
$H^3(-,\mathbf Q/\mathbf Z(2))$ with the kernel of the displayed
$H^2$ map on $K_2$ and the codimension-one multiplicative groups.
Localizing at five gives exactly $M_0,M_1$ as above.  Continuous
Galois cohomology commutes with these filtered coefficient
localizations.  The theorem's identification is natural for
automorphisms, hence respects $e$ and conjugation.

Restriction and norm over $K/\mathbf Q$ identify the original
arithmetic Chow, \emph{\'etale motivic}, and unramified groups
with the conjugation invariants of their $K$-versions: the
composite is multiplication by two and two is a unit in $R$.
Taking invariants of this order-two group is exact.
Apply (F25.3) over $K$ and take conjugation invariants to obtain
(F25.6).  Finally inversion acts trivially on constant $K_2$,
so the factor $(1-[-1]^*)/2$ in $e$ annihilates it.
This proves the last assertion about $M_0$.
\end{proof}

\paragraph{Exact source and its limits.}
Colliot-Th\'el\`ene--Kahn, \emph{Cycles de codimension~2 et
$H^3$ non ramifi\'e pour les vari\'et\'es sur les corps finis},
the checked 43-page author PDF, Proposition~6.2(b), printed
p.~25, is the cohomological-dimension-two statement used in
(F25.6).  Its Theorem~6.3, also starting on p.~25, requires
cohomological dimension at most one and is \emph{not} applied
to a number field.  Theorem~6.6(a)--(b), printed p.~27, supplies
the geometric degree-three divisibility input.
Primary URL:
\url{https://www.imo.universite-paris-saclay.fr/~jean-louis.colliot-thelene/17avrilcorpsfinisCTKahn.pdf}.
The new edge calculation above verifies separately why the
projector permits (F25.1) over $\mathbf Q$.
The cohomological-dimension bound used to justify its truncation
is Milne, \emph{Arithmetic Duality Theorems}, second edition
(2006), Chapter~I, Theorem~4.10(c), printed p.~57:
\url{https://www.jmilne.org/math/Books/ADTnot.pdf}.
The exact statement was checked directly; its complete proof is
not claimed to have been reread in this run.

\paragraph{The surviving arithmetic obstruction.}
Formula (F25.5) concerns the failure of completed Chow descent.
It is only one part of the arithmetic unramified obstruction:
there is an exact sequence at its first two nonzero terms
\[
 0\longrightarrow T_5D\longrightarrow T_5C
       \longrightarrow (T_5\overline C)^{G_{\mathbf Q}}.
 \tag{F25.7}
\]
The target can already contain nonzero invariant vectors, despite
$T^{G_{\mathbf Q}}=0$, because the exact sequence in C44-F23
also contains $\widehat{\overline A}$.  C44-F24 makes the
distinction explicit at the formal-module level: its defect
$D$ is zero, whereas its arithmetic $T_5C$ is nonzero.
No claim that this model is geometric is being made.
Even vanishing of the entire $T_5C$ would still leave the
separate scalar-to-completion kernel in C44-F03 to control.
Thus (F25.6) is a precise actual-group descent gate, not a
finite-generation theorem or a closure of the full Abel--Jacobi
problem.  No 032/Milne input is used.

% END work/arithmetic/fermat_actual_descent.tex

% BEGIN work/arithmetic/fermat_source_audit.tex
% C44_POS, /root/fermat_aj. Source audit only; no new theorem.
% No 032, supplied Milne report, or Milne2022 input.

\section{Source audit: the unproved Chow step for the cubic CM motive}
\label{C44-F07-source-audit}

\noindent\textbf{Audit identifier: C44-F07-SA1.}
This section records exact primary-source scopes.  It asserts no new
vanishing theorem and does not reduce the actual finite-generation
gate C44-F07.  The geometric family gate JAC-P18 and the full
arithmetic Abel--Jacobi problem remain open.

\paragraph{Conventions and the actual family group.}
Use the contravariant rational Chow motives of C44-F17--F18, with
$\mathbb L=\mathbf Q(-1)$.  Over
$k=\overline{\mathbf Q}$ the carried Fermat motive $N$ is the cubic
CM summand $N_3$ of an elliptic cube.  After a finite number-field
extension the isogenies and projectors already exist.  The complementary
part of $h^3$ is of the form $h^1(A)(-1)$, while $N$ has types
$(3,0),(0,3)$.  Thus $N$ is a lift of the coniveau-zero summand
denoted $t_3$ in the source below.  This uses the actual Chow
identification in C44-F18; it does not infer Chow vanishing from
these Hodge types.

For a pointed smooth projective curve $C/k$, the precise group under
discussion is
\[
 W_C=(\pi_1^C\otimes p_N)_*\operatorname{CH}^2(C\times J)
 \simeq
 \operatorname{Hom}_{\rm CHM}\bigl(h^1(C)(-1),N\bigr).
 \tag{F07.SA1}
\]
The excess-quotient version is C44-F17.  Kahn instead uses effective
\emph{covariant} motives $\mathcal M$, and writes $M(1)=M\otimes L$
for the Lefschetz twist.  His Section~1.5, printed p.~596, explicitly
specifies both conventions.  Passing to the opposite category first
gives
\[
 \mathcal M\bigl(h_1(C)(1),N_{\rm cov}\bigr)
 \simeq
 \operatorname{Hom}_{\rm CHM}\bigl(N,h^1(C)(-1)\bigr).
\]
The self-dual middle projector and the curve polarization give
$N^\vee\simeq N(3)$ and $h^1(C)^\vee\simeq h^1(C)(1)$.
Dualizing the last Hom group and twisting both entries by $-3$
therefore identifies it with (F07.SA1).  This is why both the
variance and the sign of the twist must be changed; simply copying
$h_1(C)(1)$ into the contravariant notation is incorrect.

\paragraph{The proved birational statement and the conjectural step.}
Kahn, \emph{Albanese kernels and Griffiths groups}, Tunisian Journal
of Mathematics~3 (2021), 589--656, Proposition~7.6, printed
pp.~607--608, proves
\[
 \mathcal M(M(2),t_3(T))=0,\qquad
 \mathcal M_k\bigl(h(Y)(1),t_3(T)\bigr)
 \simeq
 \mathcal M_{k(Y)}\bigl(L,t_3(T)_{k(Y)}\bigr).
 \tag{F07.SA2}
\]
Here $T$ satisfies the threefold hypotheses of that section and
$Y$ is smooth projective and connected.  Proposition~7.2, printed
p.~606, verifies its Chow--Lefschetz hypothesis for abelian
threefolds.  Rationalizing the source's refined groups
``up to isogeny'' gives the ordinary rational groups used here.
The first formula has \emph{two} Lefschetz factors.  It does not
annihilate $\mathcal M(h_1(C)(1),t_3(T))$, which has only one.
The second formula is a restriction isomorphism, not a vanishing
statement.

Immediately afterward, Conjecture~7.8, printed p.~609, asks for
bijectivity of
\[
 \mathcal M(L,t_3(T))\longrightarrow
 \mathcal M_{\rm alg}(L,t_3(T)).
 \tag{F07.SA3}
\]
For the self-dual cubic $N$, these are its codimension-two cycle
groups before and after algebraic equivalence.  Proposition~7.9
on the same page derives this conjecture from the
Bloch--Beilinson--Murre conjectures.  Its proof, following the
Weil--Bloch trick in Lemma~7.10, requires the vanishing
\[
 \mathcal M\bigl(h_1(\Gamma)(1),t_3(T)\bigr)=0
 \tag{F07.SA4}
\]
over the function-field extensions that occur there.  The numerical
Hom group vanishes by the coniveau decomposition; the passage from
rational to numerical equivalence is the \emph{conjectural} input
from Jannsen's Proposition~5.8, invoked explicitly in the proof.
For $t_3(T)=N$, the Chow group in (F07.SA4) is exactly (F07.SA1)
after the translation above.  The conjecture (F07.SA3) over one
fixed ground field is not being asserted equivalent to every
family-group vanishing: the proof uses the stronger Hom statement
and base changes.  Neither Proposition~7.6 nor Kimura uniqueness
of the coniveau lift supplies that statement.

These assertions are not restricted to finite fields: Kahn's base
field is general, and the unproved step remains unproved when
$k$ is a number field or $\overline{\mathbf Q}$.  Any particular
geometric cycle and rational equivalence descend to a finite
number-field extension, and restriction--norm permits rational
descent.  This observation is not integral descent, actual motivic
finite generation, or arithmetic Abel--Jacobi injectivity.
The checked primary published PDF is
\url{https://webusers.imj-prg.fr/~bruno.kahn/preprints/tunis.pdf}.
The entire proofs of Proposition~7.6 and Proposition~7.9, and the
adjoining conventions and statements, were read in this run;
no complete-paper proof audit is claimed.

\paragraph{Two nearby arithmetic results have different scopes.}
Broe, \emph{On the Beilinson--Bloch conjecture over function
fields}, arXiv:2505.00696v2, 27~January~2026, Theorem~6.13,
pp.~32--33, proves vanishing of the Chow groups of the same type
of nontrivial CM tensor power over \emph{global function fields
of positive characteristic}.  Its finite-field model argument
does not give cycle-class injectivity on an arithmetic model over
the integers of a number field.  Theorem~A.7 identifies the cubic tensor with
$t_3(E^3)$, and Corollary~A.8 uses precisely Kahn's birational
restriction above; it does not bypass (F07.SA4).
Primary: \url{https://arxiv.org/pdf/2505.00696}.

Gonz\'alez-Avil\'es, \emph{On $K_2$ of varieties over number
fields}, Journal of K-Theory~1 (2008), 175--183,
Theorem~3.5/Corollary~3.6, p.~180, give CM vanishing for
$H^1(k,K_2(\bar k(X))/H^0(\bar X,\mathcal K_2))$.
Lemma~3.1/Remark~3.2, pp.~177--178, treat
$H^2(k,H^0(\bar X,\mathcal K_2))$.
C44-F25 instead leaves an $H^2$ kernel on the \emph{full}
function-field $K_2$, constrained by tame symbols.
The paper's Proposition~4.3, pp.~181--182, does address
$H^2(k,K_2(\bar k(X)))$, but its Hasse injection is for curves:
the proof uses $\operatorname{cd}(\bar k(X))=1$ to make that
$K_2$ divisible.  The geometric function field of $J$ has
transcendence degree three.  Thus neither the CM $H^1$ theorem
nor the curve $H^2$ theorem proves the remaining arithmetic
vanishing.  The complete published text posted by its author is
\url{https://www.researchgate.net/publication/231908479_On_K2_of_varieties_over_number_fields}.

All references in this audit are used with their stated scopes.
No 032, supplied Milne report, or Milne2022 Proposition~4.2 is
used, and none of the outstanding Fermat gates is declared closed.

% END work/arithmetic/fermat_source_audit.tex

% BEGIN work/arithmetic/fermat_cartier_families.tex
% C44_POS, /root/fermat_aj.  Unconditional family-level P18 elimination.
% No 032 or conditional Milne positivity input.  Full AJ5 remains open.

\section{Cartier families on the singular image surface}

Keep the abelian threefold $A$ and its distinguished
extreme projector $p\le\pi_3^A$ from C44-F15.
All schemes and rational Chow groups in this section
are over $\overline{\mathbf Q}$.  The Picard group
of a singular surface is used directly; it must
not be confused with that of its normalization
or resolution.

\begin{proposition}[C44-F22]\label{C44-F22}
Let $\iota:D_0\hookrightarrow A$ be an integral
ample divisor, with no normality or smoothness
assumption.  Then
\[
 \iota^*:\operatorname{Pic}^0(A)
                       \longrightarrow\operatorname{Pic}^0(D_0)
 \tag{F22.1}
\]
is an isogeny of abelian varieties.  In particular
the target is proper, although this need not hold
for the Picard variety of an arbitrary non-normal
projective variety.

For every smooth projective curve $C$ and line
bundle $\mathcal L$ on $C\times D_0$, one has
\[
 (\pi_1^C\otimes p)_*(1_C\times\iota)_*
       \bigl(c_1(\mathcal L)\cap[C\times D_0]\bigr)=0
                         \quad\hbox{in }\operatorname{CH}^2(C\times A).
 \tag{F22.2}
\]
The same holds for a rational linear combination
of Cartier divisor cycles.  Hence it holds for
an integral cycle $Z\subset C\times D_0$ if
some nonzero integer multiple of $[Z]$ is the
cycle of a Cartier divisor on $C\times D_0$.
In particular it holds if the scheme morphism
$Z\to D_0$ is finite and flat.

Thus a surviving P18 family on the residual
surfaces of F16--F21 must fail this Cartier-cycle
condition.  It cannot be represented by an
integral finite-flat family on its image surface.
No assertion that every Weil divisor family is
Cartier, or that every generically finite family
is flat, is used.
\end{proposition}

\begin{proof}
Since $A$ is smooth, $D_0$ is an effective
Cartier divisor on $A$.  The restriction exact
sequence is
\[
 0\longrightarrow\mathcal O_A(-D_0)
  \longrightarrow\mathcal O_A
  \longrightarrow\mathcal O_{D_0}\longrightarrow0.
 \tag{F22.3}
\]
Serre duality on the threefold $A$, whose
canonical bundle is trivial, and Kodaira
vanishing for the ample bundle
$\mathcal O_A(D_0)$ give
\[
 H^1(A,\mathcal O_A(-D_0))=
 H^2(A,\mathcal O_A(-D_0))=0.
 \tag{F22.4}
\]
Consequently restriction induces an isomorphism
\[
                  H^1(A,\mathcal O_A)
                    \xrightarrow{\sim}H^1(D_0,\mathcal O_{D_0}).
 \tag{F22.5}
\]
The Picard scheme exists for the projective
integral scheme $D_0$; its identity component
is of finite type and separated, and has
tangent space $H^1(D_0,\mathcal O_{D_0})$.
In characteristic zero it is smooth.
The morphism (F22.1) has isomorphism on
tangent spaces by (F22.5).  Its kernel is
therefore zero-dimensional and finite.  Its
image is closed, because the source is proper
and the target is separated, and has the full
dimension three of the smooth connected target.
Thus its image is the whole identity component.
The target is proper and the map is an isogeny,
which proves (F22.1).  No statement about
$\operatorname{Pic}^0(\widetilde D_0)$ follows
from this argument; that group may be larger.

We spell out the divisor-family decomposition
to avoid using a smoothness assumption on
$D_0$ implicitly.  Choose $c_0\in C(k)$ and
put $\mathcal L_0=\mathcal L|_{\{c_0\}\times D_0}$.
The line bundle
$\mathcal L\otimes\operatorname{pr}_{D_0}^*
\mathcal L_0^{-1}$ gives a pointed morphism
\[
                         C\longrightarrow\operatorname{Pic}^0(D_0).
 \tag{F22.6}
\]
Since the target is now known to be abelian,
it factors through a homomorphism
$\operatorname{Jac}(C)\to\operatorname{Pic}^0(D_0)$.
The isogeny (F22.1) has a rational inverse,
so after multiplying by an integer this
homomorphism lifts to
$\operatorname{Pic}^0(A)$.
The normalized Poincar\'e divisor for the
lift on $C\times A$ gives a rational divisor
class $\ell$ satisfying
\[
                         (1_C\otimes\pi_1^A)_*\ell=\ell.
 \tag{F22.7}
\]
The exact projector assertion follows either
from the Poincar\'e construction or from
target multiplication by $n$, whose pullback
on this cross divisor is multiplication by $n$.

By construction, the rational line bundle
$\mathcal L-\ell|_{C\times D_0}-
\operatorname{pr}_{D_0}^*\mathcal L_0$ has
trivial class on every fiber over $C$.
Here subtraction denotes the tensor operation
in the rational Picard group, after clearing
the same integer denominator.
The seesaw principle for the proper integral
variety $D_0$ therefore makes the difference
the pullback of a rational line bundle on $C$.
It follows that the Cartier divisor cycle of
$\mathcal L$ is a rational sum of the following
three types:
\[
  \ell|_{C\times D_0},\qquad
  \operatorname{pr}_{D_0}^*b,\qquad
  \operatorname{pr}_C^*a.
 \tag{F22.8}
\]
This argument uses line bundles and the
Picard functor of $D_0$ itself, not a
representability assertion for codimension-two
Chow groups.

The projection formula sends the first term
to
\[
                             \ell\cdot\operatorname{pr}_A^*[D_0]
                 \quad\hbox{on }C\times A.
 \tag{F22.9}
\]
It is killed by $1_C\otimes p$, using (F22.7)
and the exact identity
$pL_{[D_0]}\pi_1^A=0$ from C44-F15.
The second term in (F22.8) pushes to
$[C]\boxtimes\iota_*(c_1(b)\cap[D_0])$
and is killed by $\pi_1^C[C]=0$.
The third pushes to $a\boxtimes[D_0]$
and is killed by $p\operatorname{CH}^1(A)=0$.
This proves (F22.2), and rational linearity
gives the Cartier-cycle assertion for $Z$.

Finally suppose $Z\to D_0$ is finite flat.
The ambient morphism $C\times D_0\to D_0$
is smooth of relative dimension one, hence
flat and of finite presentation.  Each fiber
$Z_d$ is a finite closed subscheme of the
smooth curve $C_{\kappa(d)}$, and so is an
effective Cartier divisor on that curve.
The fiberwise criterion for relative effective
Cartier divisors applies: the closed immersion
is of finite presentation, and $Z$ is flat
over the base.  Therefore $Z$ is an effective
Cartier divisor on $C\times D_0$, with
$[Z]=c_1(\mathcal O(Z))\cap[C\times D_0]$.
Equation (F22.2) applies.  This criterion
requires no regularity of the base $D_0$.
\end{proof}

\paragraph{Relation with the Albanese-cycle criterion.}
For F20's surface cycle $z_{(0)}$ on $B$, let
$z_{\rm dist}$ be O'Sullivan's unique distinguished
lift of its numerical class.  In characteristic
zero numerical and homological equivalence
coincide on abelian varieties (Kleiman,
\emph{Algebraic cycles and the Weil conjectures},
Remark~3.10, printed p.~380, using Lieberman's
Theorem~2A11 on p.~377), so
\[
                 z_{(0)}-z_{\rm dist}
                    \in\operatorname{CH}^{b-2}_{(0)}(B)_{\rm hom}.
 \tag{F22.10}
\]
The distinguished lift is itself in grade zero:
multiplication by $n$ and uniqueness of the
distinguished lift give the required eigenvalue.
The operator (F20.3) kills $z_{\rm dist}$ by
the same distinguished/Hodge argument as in
F20.  Thus it depends only on the difference
in (F22.10).  Fourier identifies the latter
with a homologically trivial group
$\operatorname{CH}^2_{(0)}(B^\vee)$.
This makes the relationship to AR-15a precise;
it is a reformulation, not an additional
vanishing theorem or a discharge of AJ5.

\paragraph{Primary sources and exact scope.}
The classical characteristic-zero numerical/homological
comparison just used is checked in Kleiman's original
\emph{Algebraic cycles and the Weil conjectures},
\emph{Dix expos\'es sur la cohomologie des sch\'emas}
(1968), 359--386, Remark~3.10 and Corollary~3.9,
printed p.~380, together with Theorem~2A11
(Lieberman), printed p.~377.  The retrieved scan is
\url{https://download.uni-mainz.de/mathematik/Algebraische%20Geometrie/Lehre/WS23.Padische.1966.kleiman.weil.pdf}.
These are unconditional classical Betti statements.

For the singular Picard functor, Kleiman,
\emph{The Picard scheme}, arXiv:math/0504020v1
(1 April 2005), Theorem~5.4, printed p.~37,
gives existence and quasi-projectivity for
projective geometrically integral schemes;
Theorem~5.11, printed pp.~42--44, identifies
the tangent space; and Corollary~5.14,
printed p.~44, gives smoothness in
characteristic zero.  No smoothness or
normality of $D_0$ is imposed in these
assertions.  Primary author preprint:
\url{https://arxiv.org/pdf/math/0504020}.
The universal rigidified Picard family under
the same hypotheses is also stated explicitly
in Conrad, \emph{Math 248B: Picard functors
for curves}, Theorem~2.2, printed p.~1;
Proposition~3.1 on p.~2 computes the tangent
space before the notes specialize to curves:
\url{https://virtualmath1.stanford.edu/~conrad/248BPage/handouts/pic.pdf}.

For the finite-flat implication, Stacks
Project, Lemma~31.19.9(2) and its final
statement, Tag~062Y, gives exactly the
fiberwise criterion used in the proof:
\url{https://stacks.math.columbia.edu/tag/062Y}.
Kimura, \emph{Murre's conjectures for certain
product varieties}, arXiv:math/0509249v2
(20 September 2005), Lemma~2.2, printed p.~5,
records this implication in the same
curve-times-image-surface setting; the
stronger hypotheses of that paper's
Theorem~2.1 are not used here.
For the seesaw principle on a possibly
singular complete integral variety, Milne's
\emph{Abelian Varieties} (1986), corrected
transcription dated 2 January 2022,
Theorem~5.1 and Corollary~5.2, printed p.~7,
have the required scope:
\url{https://www.jmilne.org/math/xnotes/AVs.pdf}.
This is a classical Picard-theoretic input,
not a conditional Milne positivity statement.

\paragraph{Remaining gap.}
Finite-flat and Cartier families on the
singular ample image are now eliminated
without a hypothesis on the irregularity
of its resolution or the distinguishedness
of its Albanese surface cycle.  Families
that do not descend as Cartier cycles
remain possible.  Passing to a resolution
or flattening blowup can enlarge the Picard
variety; it does not repair the argument
for those families automatically.
The other residual P18 classes and the
full actual \'etale-motivic finite-generation
gate C44-F07 remain open.

% END work/arithmetic/fermat_cartier_families.tex

% BEGIN work/arithmetic/fermat_albanese_surfaces.tex
% C44_POS, /root/fermat_aj.  Further unconditional P18 reductions.
% No 032 or conditional Milne positivity input.  Full AJ5 remains open.

\section{Albanese surface cycles and an excess-irregularity example}

We retain the notation and hypotheses of C44-F15: $A$ is
an abelian threefold over $k=\overline{\mathbf Q}\subset
\mathbf C$, and $p\le\pi_3^A$ is a distinguished projector
with only Hodge types $(3,0),(0,3)$.  The Fermat pair
$(J,p_N)$ satisfies these hypotheses.  All Chow groups
have rational coefficients.  These results concern the
algebraically trivial obstruction of JAC-P18; they do not
prove full arithmetic Abel--Jacobi injectivity.

\subsection{A vanishing criterion on the Albanese image}

\begin{proposition}[C44-F20]\label{C44-F20}
Let $i:D\to A$ be a resolution of an integral ample
surface $D_0\subset A$.  Put $B=\operatorname{Alb}(D)$,
$b=\dim B$, and $P=B^\vee$.  The Albanese homomorphism
$\rho:B\to A$ induced by $i$ is surjective.  Choose a
translate $j:D\to B$ of the pointed Albanese map such
that $i=\rho\circ j$, and put
\[
                        z=j_*[D]\in\operatorname{CH}^{b-2}(B).
 \tag{F20.1}
\]
Suppose that its Beauville grade-zero component $z_{(0)}$
is symmetrically distinguished.  Then, for every curve $C$,
\[
 (\pi_1^C\otimes p)_*(1_C\times i)_*
                    \operatorname{CH}^1(C\times D)=0.
 \tag{F20.2}
\]
Consequently the universal excess-Picard cycle of C44-F17
for this surface is zero, and the surface contributes
nothing to the P18 obstruction.

The hypothesis holds if $z$ belongs to the subalgebra
of $\operatorname{CH}^*(B)$ generated by divisor classes.
More generally the same conclusion holds for any
factorization $i=\rho j$ through an abelian variety $B$
such that $j^*:\operatorname{Pic}^0(B)\to
\operatorname{Pic}^0(D)$ is an isogeny and $z_{(0)}$ is
distinguished; the source map $i$ may then be generically
finite onto its image rather than birational.

In particular, if $D$ is a smooth complete-intersection
surface of ample divisors in an abelian variety $B$,
and a homomorphism $\rho:B\to A$ is generically finite
on $D$ with surface image, then all projected families
supported over that image vanish.  This statement does
not assume that $D$ or its image is dominated by a
product of curves.
\end{proposition}

\begin{proof}
The surjectivity of $\rho$ for the ample image is the
argument already checked in F16: if its image were a
proper abelian subvariety, the image surface would be
a translate of an abelian surface and hence non-ample.
For a pointed Albanese map $a$, one has
$i=\rho a+x$ for a constant $x\in A(k)$.  Choose
$b_0\in B(k)$ with $\rho(b_0)=x$, and take
$j=a+b_0$.  Pullback on $\operatorname{Pic}^0$ is unchanged
by this translation and identifies $P$ with
$\operatorname{Pic}^0(D)$.

Let $\ell\in\operatorname{CH}^1(P\times B)$ be the
normalized Poincar\'e divisor, and let $r_B:P\times B\to B$
be the projection.  The all-Picard universal family
after target projection is
\[
 \Gamma=(\pi_1^P\otimes p)_*
     (1_P\times\rho)_*\bigl(\ell\cdot r_B^*z\bigr)
                       \in\operatorname{CH}^2(P\times A).
 \tag{F20.3}
\]
To verify the formula before drawing a conclusion from
it, the normalized Poincar\'e family on $P\times D$
is the pullback of $\ell$ by $1_P\times j$, up to a
divisor pulled back from $P$.  The projection formula
gives
\[
 (1_P\times j)_*(1_P\times j)^*\ell
                            =\ell\cdot r_B^*j_*[D].
 \tag{F20.4}
\]
Composing pushforwards with $1_P\times\rho$ yields
(F20.3).  The normalization term pushes to a divisor
on $P$ externally multiplied by the divisor $i_*[D]$
on $A$ and is killed by $p\operatorname{CH}^1(A)=0$.
If $j^*$ is merely an isogeny, the same construction
parametrizes every cross divisor after rational
linear combinations: homomorphisms into
$\operatorname{Pic}^0(D)$ lift up to isogeny to $P$.

We now check precisely which part of $z$ can contribute.
The Poincar\'e divisor is in Beauville grade zero on
$P\times B$, since simultaneous multiplication by $n$
pulls it back to $n^2\ell$.  Multiplication by this
divisor and proper pushforward by a homomorphism of
abelian varieties preserve the Beauville grade.
The codimension in (F20.3) is correct:
$\ell\cdot r_B^*z$ has codimension $b-1$, and
$1_P\times\rho$ has source minus target dimension
$b-3$ (and is smooth of that relative dimension in
the surjective case).
The output projector $\pi_1^P\otimes p$ has total
Chow--K\"unneth index $1+3=4$ in codimension two,
so it annihilates every Beauville grade except zero.
It follows that (F20.3) is unchanged on replacing
$z$ by $z_{(0)}$.

Under the hypothesis, every cycle and operation now
appearing in (F20.3) is distinguished: the symmetric
Poincar\'e divisor, $z_{(0)}$, the canonical projector
$\pi_1^P$, $p$, the projections, and the abelian
homomorphism $1_P\times\rho$.
O'Sullivan's functorial distinguished subalgebras
therefore contain $\Gamma$.  Its cohomology lies in
\[
                            H^1(P)\otimes\operatorname{im}p.
 \tag{F20.5}
\]
The possible Hodge types are $(4,0),(3,1),(1,3),(0,4)$,
and none is $(2,2)$.  The class of an algebraic
codimension-two cycle in (F20.5) is consequently zero.
A homologically trivial distinguished cycle is zero
in rational Chow, by injectivity of the distinguished
subalgebra into numerical cycles.  Thus
\[
                                  \Gamma=0.
 \tag{F20.6}
\]

For a homomorphism $\phi:\operatorname{Jac}(C)\to P$,
pull back (F20.3) along the pointed map
$C\to\operatorname{Jac}(C)\xrightarrow{\phi}P$ and
then apply $\pi_1^C\otimes p$.  By the Poincar\'e
identification and the projection formula this is
the projected pushforward of the cross divisor
associated with $j^*\phi$.  The observation is
linear for rational homomorphisms.  The two
single-factor terms in $\operatorname{Pic}(C\times D)$
are killed exactly as in F15.  Equation (F20.6)
therefore proves (F20.2).  The universal quotient
cycle of F17 is zero by its curve-detection statement,
or directly by restricting (F20.3) to a rational
section of the excess Picard quotient.

Every divisor on an abelian variety splits into its
grade-zero symmetric part and its grade-one
$\operatorname{Pic}^0$ part.  The grade-zero part of a
polynomial in divisors is thus the same polynomial
in their symmetric parts, in the prescribed
codimension.  Symmetric divisors are distinguished,
and distinguished cycles form a subalgebra.  This
proves the stated sufficient condition on $z$.

For the complete-intersection assertion, write
$D=H_1\cap\cdots\cap H_{b-2}$ scheme-theoretically.
Its pushforward fundamental
cycle is the product of the divisor classes $H_i$.
Moreover restriction on $\operatorname{Pic}^0$ is an
isogeny.  This last assertion can be seen directly
without a conjecture: at each smooth ample divisor
$Y\subset X$ with $\dim X\ge3$, the sequence
$0\to\mathcal O_X(-Y)\to\mathcal O_X\to
\mathcal O_Y\to0$, Serre duality and Kodaira vanishing
give
\[
                            H^1(X,\mathcal O_X)
                              \simeq H^1(Y,\mathcal O_Y).
 \tag{F20.7}
\]
Indeed $H^1(X,\mathcal O_X(-Y))$ and
$H^2(X,\mathcal O_X(-Y))$ vanish, since their duals
are positive-degree cohomology groups of
$\omega_X(Y)$.  Thus restriction on the Picard
varieties has an isomorphism on tangent spaces;
in characteristic zero it is an isogeny.
If the successive intersections are smooth, iteration
proves the assertion for $D$.  The general smooth
complete-intersection case follows from its Koszul
resolution on $B$: every nontrivial Koszul term is
a sum of negative ample line bundles, whose cohomology
vanishes in degrees below $b$ by Serre duality and
Kodaira vanishing on $B$.  The resolution has length
$b-2$, so $H^1(B,\mathcal I_D)$ and
$H^2(B,\mathcal I_D)$ vanish.  Hence restriction
$H^1(B,\mathcal O_B)\to H^1(D,\mathcal O_D)$ is an
isomorphism in this case as well.
The general factorization case of the proof applies.
Finally, every integral family cycle over its image
lifts after multiplication by the generic degree
to a divisor on $C\times D$, by the finite-flat-open
and closure argument (F19.15).  Hence the conclusion
holds for all projected image families, not just
the families already presented on $D$.
\end{proof}

\paragraph{Independence of the translated Albanese choice.}
The grade-zero condition in F20 does not depend on the
translation used in (F20.1).  For completeness this
uses only a known low-dimensional Fourier bound, not
the general Beauville conjecture.  A surface cycle on
an abelian variety of dimension $b$ has codimension
$b-2$, and Fourier sends its grade $s$ to
$\operatorname{CH}^{2+s}_{(s)}(B^\vee)$.
For $s=-1$ this is the zero negative-grade part of
the divisor group; for $s=-2$ it is a nonzero-grade
part of $\operatorname{CH}^0$, also zero; for smaller
$s$ the codimension is negative.  Thus every surface
cycle has only grades $s\ge0$.
Under Fourier, translation is multiplication by
the exponential of a $\operatorname{Pic}^0$ divisor
(up to the irrelevant sign determined by pushforward
or pullback).  Such a divisor has grade one.
Translation therefore leaves the grade-zero part
of every surface cycle unchanged.  This also shows
why no inference about arbitrary higher-dimensional
negative-grade cycles is being made here.

\subsection{An explicit class of singular ample images with excess Picard}

\begin{proposition}[C44-F21]\label{C44-F21}
Let $E$ be any elliptic curve and let
$B=A\times E$, with projection $\rho:B\to A$.
Choose very ample line bundles $H$ on $A$ and $M$ on
$E$, and put $L=H\boxtimes M$.
For a general pair of divisors $H_1,H_2\in|L|$,
the smooth connected surface
\[
                          S=H_1\cap H_2\subset B
 \tag{F21.1}
\]
has the following properties:
\begin{enumerate}
\item The projection $S\to D_0:=\rho(S)\subset A$
is finite birational, so $S$ is the normalization of
the integral image surface $D_0$.
\item The divisor $D_0$ is ample and non-normal,
and $q(S)=4$.  In particular its smooth resolution
has irregularity greater than three.
\item Every P18 family over $D_0$ is killed by
$\pi_1^C\otimes p$, by C44-F20.
\end{enumerate}
Thus the vanishing criterion actually applies inside
F16's residual singular-ample, excess-irregularity
class.  No claim is needed or made that these surfaces
are, or are not, dominated by products of curves.
\end{proposition}

\begin{proof}
The line bundle $L$ is very ample, by the product
embeddings followed by the Segre embedding.  Bertini
and connectedness of ample complete intersections
give a smooth connected surface for a general pair.
Let $N=\dim|L|$, so the pair parameter space has
dimension $2N$.  Consider the off-diagonal relative
pair variety
\[
 R=(B\times_A B)\setminus\Delta_B
                \simeq A\times(E\times E\setminus\Delta_E).
 \tag{F21.2}
\]
It has dimension five.  For any pair $(x,y)\in R$,
the two points of $B$ are distinct, so very ampleness
implies that evaluation of sections of $L$ at $x,y$
is surjective.  Requiring a divisor in $|L|$ to
contain both points imposes two independent linear
conditions.  Requiring both $H_1$ and $H_2$ to do
so imposes four.  The incidence variety of
$(x,y,H_1,H_2)$ thus has dimension
\[
                              5+2N-4=2N+1.
 \tag{F21.3}
\]
By the dimension-of-fibers theorem, for a general
pair the off-diagonal fiber product
$(S\times_A S)\setminus\Delta_S$ has dimension at
most one (or is empty).

The image of $S$ must have dimension two: if its
dimension were one or zero, the generic fibers
would have positive dimension, giving an
off-diagonal pair locus of dimension at least
three or four, respectively.  If the generically
finite map onto its surface image had degree at
least two, characteristic zero makes it generically
separable, and its off-diagonal pair locus would
have dimension two.  Both contradict the bound
just obtained.  Its generic degree is therefore one.
There are also no positive-dimensional fibers
anywhere: a positive-dimensional closed subset of
the elliptic fiber $\{a\}\times E$ must be that
whole fiber.  Such a fiber inside $S$ would yield
an off-diagonal pair locus of dimension two.
Hence the map is quasi-finite and proper, and
therefore finite.  Since $S$ is smooth and integral,
a finite birational map makes it the normalization
of $D_0$.

Write $h=\operatorname{pr}_A^*c_1(H)$ and
$m=\operatorname{pr}_E^*c_1(M)$.  In Chow one has
\[
 [S]=(h+m)^2=h^2+2hm,\qquad
 \rho_*[S]=2\deg(M)c_1(H).
 \tag{F21.4}
\]
Here $m^2=0$, and $\rho_*h^2=0$ by the projection
formula.  The generic degree is one, so
$[D_0]=\rho_*[S]$ as divisor classes.  Thus $D_0$
is linearly equivalent to a positive multiple of
the ample divisor $H$, and is ample.

The restriction argument (F20.7), twice, gives
$H^1(S,\mathcal O_S)\simeq H^1(B,\mathcal O_B)$,
so $q(S)=4$.  Alternatively the Koszul resolution
of the two divisors uses the vanishing
$H^i(B,L^{-r})=0$ for $i<4$ and $r=1,2$, and
gives the same isomorphism.
If $D_0$ were normal, its finite normalization
map would be an isomorphism, making it a smooth
ample divisor in $A$.  The same restriction
argument on $A$ would then give $q(S)=q(A)=3$,
a contradiction.  It is consequently non-normal
and singular, with smooth normalization of
irregularity four.

Finally $[S]$ is a product of divisor classes in
$B$ and restriction on $\operatorname{Pic}^0$
is an isogeny.  C44-F20 therefore annihilates
all projected families over $D_0$.
Every construction and the nonempty open choice
of the pair can be made over $\overline{\mathbf Q}$.
\end{proof}

\paragraph{Source scope and remaining obstruction.}
The distinguished-cycle inputs are exactly the
O'Sullivan theorem and functoriality already
located in C44-F15.  For the grading and the
translation observation we use Beauville (1986),
Propositions~1--2, printed p.~648, and Beauville
(1983), Proposition~4(i),(ii), printed p.~246,
with its proof on p.~247.  The latter states the
Fourier transform of a point and the translation
formula; all identities are in rational Chow.
Primary author PDFs:
\url{https://math.univ-cotedazur.fr/u/beauvill/pubs/Chow.pdf}
and
\url{https://math.univ-cotedazur.fr/u/beauvill/pubs/fourier.pdf}.
The Picard-family identification is the exact
Picard-product statement already used in F15.
The complete-intersection and incidence arguments
above are proved explicitly, using Bertini,
Kodaira vanishing, and standard proper
pushforward and finite-morphism facts.

F20's distinguished Albanese-cycle hypothesis
is a sufficient condition that has just been
verified for a specified class of image surfaces.
It is not asserted for every Albanese image.
The remaining F17 universal cycles for other
singular ample surfaces with excess irregularity
remain open, as does the full AJ5 finite-generation
gate C44-F07.  No step in this file uses 032 or a
conditional Milne positivity statement.

% END work/arithmetic/fermat_albanese_surfaces.tex

% BEGIN work/arithmetic/fermat_product_surfaces.tex
% C44_POS, /root/fermat_aj.  Unconditional further P18 reduction.
% No 032 or conditional Milne positivity input.  Full AJ5 remains open.

\section{Eliminating surfaces dominated by products of curves}

The excess Picard variety in C44-F16 need not be zero for a
surface obtained by adding two curves in an abelian threefold.
The following argument nevertheless kills all its projected
divisor families.  It uses the exact divisor correspondence
identity of C44-F15, rather than an assertion that the Picard
variety of the surface is inherited from the ambient variety.
Work over $k=\overline{\mathbf Q}\subset\mathbf C$ and with
rational Chow groups.

\begin{proposition}[C44-F19]\label{C44-F19}
Let $A$ be an abelian threefold, and let $p$ be a symmetrically
distinguished projector below $\pi_3^A$ whose Betti realization
has only Hodge types $(3,0)$ and $(0,3)$.  Let
\[
                 f:C_1\times C_2\longrightarrow A
 \tag{F19.1}
\]
be a morphism from a product of smooth projective connected
curves whose image is an integral surface $D_0$.
For every smooth projective connected curve $C$ and every
divisor class $W\in\operatorname{CH}^1(C\times C_1\times C_2)$,
one has
\[
 (\pi_1^C\otimes p)_*(1_C\times f)_*W=0
                    \quad\hbox{in }\operatorname{CH}^2(C\times A).
 \tag{F19.2}
\]
Moreover every integral codimension-two cycle
$Z\subset C\times A$ whose projection onto $A$ has image
$D_0$ satisfies
\[
                         (\pi_1^C\otimes p)_*[Z]=0.
 \tag{F19.3}
\]

In particular this applies to the sum of any two curves in
$A$ when their sum has dimension two.  For the Fermat
projector $p_N$ on $J$, all such image surfaces can be
removed from the residual family obstruction of F16--F17,
even if their smooth resolutions have irregularity greater
than three.  More generally it applies whenever there is a
generically finite dominant rational map from a product of
two smooth projective curves to the image surface: its
composition with the inclusion in $J$ extends to a morphism.
It is not asserted that every
singular ample surface with excess irregularity has this
property.
\end{proposition}

\begin{proof}
Choose points on the two curves.  The Albanese variety of
their product is the product of their Jacobians.  The
universal property of the Albanese therefore writes the
morphism (F19.1) as
\[
                       f(x,y)=f_1(x)+f_2(y),
 \tag{F19.4}
\]
where the constant translation has been absorbed in one
of the two maps.  The maps $f_j:C_j\to A$ need not send
the chosen points to zero.  Equivalently (F19.4) follows
from the rigidity lemma for maps to an abelian variety.

The Picard decomposition of a product, applied to
$C\times(C_1\times C_2)$, expresses every divisor class as
a sum of the following types:
\[
 \begin{gathered}
 \operatorname{pr}_C^*a,\qquad
 \operatorname{pr}_{C_1\times C_2}^*b,\qquad
 \operatorname{pr}_{C\times C_1}^*\mathcal L_1,\qquad
 \operatorname{pr}_{C\times C_2}^*\mathcal L_2. 
 \end{gathered}
 \tag{F19.5}
\]
Here the $\mathcal L_j$ are normalized cross divisors,
corresponding to homomorphisms
$\operatorname{Jac}(C)\to\operatorname{Jac}(C_j)$, and are
of degree zero on every fiber $\{c\}\times C_j$.
This decomposition follows from
$\operatorname{Pic}^0(C_1\times C_2)
 =\operatorname{Pic}^0(C_1)\times\operatorname{Pic}^0(C_2)$;
any cross divisor between $C_1$ and $C_2$ is already included
in $b$.  The first term of (F19.5) pushes to
$a\boxtimes f_*[C_1\times C_2]$.  Its second factor is a
divisor on $A$, and is killed by
$p\operatorname{CH}^1(A)=0$, from F15.  The second term
pushes to $[C]\boxtimes f_*b$ and is killed by
$\pi_1^C[C]=0$.  In particular we have not asserted that
$\pi_1^C$ kills all divisor classes on $C$.

It remains to treat the third term; the fourth is identical
with the two curves exchanged.  Define
\[
 \alpha=(1_C\times f_1)_*\mathcal L_1
                     \in\operatorname{CH}^3(C\times A),
 \qquad
 \gamma=(f_2)_*[C_2]\in\operatorname{CH}^2(A).
 \tag{F19.6}
\]
Let $\star_A$ denote the relative Pontryagin operation in
the $A$ factor, with the second input pulled back from $A$.
The projection formula and (F19.4) give the exact identity
\[
 (1_C\times f)_*\operatorname{pr}_{C\times C_1}^*\mathcal L_1
                          =\alpha\star_A\gamma.
 \tag{F19.7}
\]
We next isolate its target Chow--K\"unneth indices.  Put
$\alpha_a=(1_C\otimes\pi_a^A)_*\alpha$ and
$\gamma_b=(\pi_b^A)_*\gamma$.
Because $\mathcal L_1$ has degree zero on the $C_1$ fibers,
the relative degree of $\alpha$ is zero.  The projector
$\pi_6^A=A\times\{0\}$ acts on
$\operatorname{CH}^3(C\times A)$ by this relative degree,
followed by the zero section.  Thus
\[
                             \alpha_6=0.
 \tag{F19.8}
\]
Only indices $a\le5$ can occur.  This uses no injectivity
statement about degree-zero zero-cycles.

For the fixed curve class one has
\[
                         \gamma_b=0\quad(b>4).
 \tag{F19.9}
\]
Indeed Beauville's Fourier isomorphism sends
$\operatorname{CH}^2_{(s)}(A)$ to
$\operatorname{CH}^{1+s}_{(s)}(A^\vee)$.
The index $b$ corresponds to $s=4-b$.  For $b=5$ this
would be $\operatorname{CH}^0_{(-1)}(A^\vee)=0$, since
the fundamental class is in grade zero; for $b=6$ the
target has negative codimension.  The transform is
invertible, which proves (F19.9).  This is the elementary
negative-grade vanishing for curves on an abelian
threefold, not an assumption of the general Beauville
conjecture.

The exact relative Pontryagin grading is
\[
       \alpha_a\star_A\gamma_b
        \in(1_C\otimes\pi_{a+b-6}^A)_*
                         \operatorname{CH}^2(C\times A),
 \tag{F19.10}
\]
with zero understood for an index outside $[0,6]$.
For example this follows directly from multiplication by
any integer $n\ge2$: pushforward $[n]_*$ acts on the
$a$th target component by $n^{6-a}$, and it respects
Pontryagin products by the addition square.  Hence the
product has eigenvalue $n^{12-a-b}$, giving (F19.10).
All these are identities of rational Chow eigensummands;
they remain valid with the independent smooth parameter
curve $C$.

Since $p\le\pi_3^A$, only $a+b=9$ can contribute to
$(1_C\otimes p)_*(\alpha\star_A\gamma)$.
The bounds $a\le5$ and $b\le4$ leave just one pair:
\[
 (1_C\otimes p)_*(\alpha\star_A\gamma)
                =(1_C\otimes p)_*(\alpha_5\star_A\gamma_4).
 \tag{F19.11}
\]

Apply partial Fourier transform in the $A$ factor.  Its
component on target index $a$ changes that index to $6-a$
and changes total codimension by $3-a$.  Consequently
\[
 \lambda=\mathcal F_A(\alpha_5)
      \in(1_C\otimes\pi_1^{A^\vee})_*
                              \operatorname{CH}^1(C\times A^\vee),
 \qquad
 d=\mathcal F_A(\gamma_4)
                              \in\operatorname{CH}^1_{(0)}(A^\vee).
 \tag{F19.12}
\]
In particular $d$ is a symmetric divisor class.  Fourier
transform exchanges the relative Pontryagin operation
with intersection with this fixed divisor:
\[
                 \mathcal F_A(\alpha_5\star_A\gamma_4)
                                      =\lambda\cdot d.
 \tag{F19.13}
\]
These formulas are the usual Poincar\'e-correspondence
identities, with $C$ carried along; their direct proofs
use the same addition square and projection formula.

The middle Fourier component is an isomorphism of Chow
motives $h^3(A)\to h^3(A^\vee)$.  Transport $p$ along it
to the projector $p^\vee$.  It is distinguished, because
the Fourier correspondence and its inverse are built
from the distinguished Poincar\'e divisor.  Its Hodge
realization again has only types $(3,0)$ and $(0,3)$:
Fourier sends $(r,s)$ to $(3-s,3-r)$ in this degree.
Thus C44-F15 on $A^\vee$ gives the actual Chow identity
\[
                   p^\vee L_d\pi_1^{A^\vee}=0.
 \tag{F19.14}
\]
Equations (F19.12)--(F19.14) annihilate the Fourier
transform of (F19.11).  Fourier is invertible, so
$(1_C\otimes p)_*(\alpha\star_A\gamma)=0$.
Together with (F19.5)--(F19.7), and then applying the
independent source projector $\pi_1^C$, this proves
(F19.2).  In fact the two cross terms were already
killed by the target projector alone.

Finally, the morphism $C_1\times C_2\to D_0$ is proper
and generically finite, of some degree $e>0$.  There is
a dense open $U\subset D_0$ over which it is finite
flat of degree $e$.  An integral $Z$ as in the statement
has $Z_U=Z\cap(C\times U)$ dense in $Z$.
Flat pullback of $[Z_U]$ to
$C\times f^{-1}(U)$, followed by closure of its cycle
components in the smooth threefold
$C\times C_1\times C_2$, defines a divisor cycle $W$.
Every resulting component dominates $Z$, and the finite
flat degree calculation on $U$ gives
\[
                            (1_C\times f)_*W=e[Z].
 \tag{F19.15}
\]
There is no extra component over $D_0\setminus U$ in
this proper pushforward: the image of each closure is
the closure of its image on $U$, namely $Z$.
Applying (F19.2) and dividing by $e$ proves (F19.3).
This argument does not require a lift of $f$ to a
preselected resolution of $D_0$, and therefore avoids
an unjustified identification of their Picard varieties.

If only a dominant rational map $C_1\times C_2\dashrightarrow
D_0$ is given, compose it with $D_0\hookrightarrow A$.
Weil's extension theorem extends this rational map to a
morphism on the smooth product.  Its image is still $D_0$:
the image lies in the closed subvariety $D_0$ because this
holds on a dense open subset, and it contains a dense subset
of $D_0$.  Properness makes the image closed.  The extension
is therefore a generically finite dominant morphism onto
$D_0$, to which the argument already proved applies.
\end{proof}

\paragraph{Exact residual gate and limits.}
For the algebraically trivial obstruction of JAC-P18,
C44-F16 leaves only singular ample image surfaces
$D_0\subset J$ whose smooth resolutions have
$q(\widetilde D_0)>3$, through the actual excess quotient
and cycle of C44-F17.  F19 removes, in addition, every
such surface rationally dominated by a product of two
curves.  This is a genuine
annihilation of their rational Chow families, including
sums of two curves with possibly excess irregularity.
No representability-to-injectivity implication is used.
It does not eliminate the other surfaces in that gate,
and does not establish the actual \'etale-motivic
finite generation or full arithmetic AJ5 injectivity
of C44-F07.

\paragraph{Primary source locators.}
Beauville, \emph{Sur l'anneau de Chow d'une vari\'et\'e
ab\'elienne}, Mathematische Annalen~273 (1986), 647--651,
Proposition~1 and Proposition~2(a),(b), printed p.~648,
give the exact multiplication eigenspaces, Pontryagin
grading and Fourier codimension change used above.
Primary author PDF:
\url{https://math.univ-cotedazur.fr/u/beauvill/pubs/Chow.pdf}.
The page has been visually checked because the author's
PDF is an image scan.
Beauville, \emph{Quelques remarques sur la transformation
de Fourier dans l'anneau de Chow d'une vari\'et\'e
ab\'elienne}, LNM~1016 (1983), Proposition~3(i),(ii),
printed p.~243, gives the inverse and product identities;
the explicit Poincar\'e-correspondence proof in
Proposition~3', printed pp.~245--246, also shows their
compatibility with an additional parameter factor.
Primary author PDF:
\url{https://math.univ-cotedazur.fr/u/beauvill/pubs/fourier.pdf}.
The Picard-product and distinguished-divisor inputs are
the ones already checked and cited in C44-F15.
For the product splitting and rational-map extension,
Milne, \emph{Abelian Varieties}, Arithmetic Geometry
(1986), Chapter~V, Corollary~2.5 and Theorem~3.1,
are read in the author's corrected transcription dated
2 January 2022, printed p.~4 (numbering unchanged):
\url{https://www.jmilne.org/math/xnotes/AVs.pdf}.
These are the classical rigidity and Weil extension
theorems, not a use of the conditional Milne positivity
inputs flagged elsewhere in this project.

% END work/arithmetic/fermat_product_surfaces.tex

% BEGIN work/arithmetic/fermat_residual_motive.tex
% C44_POS, /root/fermat_aj.  Unconditional refinements of C44-F15--F16.
% No 032 or Milne input.  The full arithmetic AJ5 gate remains open.

\section{A universal residual cycle and the CM elliptic subquotient}

We use C44-F15--F16 and the Fourier/source-transpose conventions of
AR-15a.  We do not repeat the latter's arbitrary parameter-Jacobian
reduction.  The new point is that the parameter can be the actual
excess Picard quotient of an image surface.  Work over
$k=\overline{\mathbf Q}\subset\mathbf C$, with rational Chow groups
and the contravariant Chow-motive convention
$\mathbb L=\mathbf Q(-1)$.

\subsection{One cycle on the actual excess quotient}

\begin{proposition}[C44-F17]\label{C44-F17}
Let $i:D\to J$ be a resolution of an integral image surface, and
put
\[
 P=\operatorname{Pic}^0(D),\qquad
 B=\operatorname{im}(\operatorname{Pic}^0(J)\to P),\qquad
 Q=P/B,\qquad q=\dim Q.
 \tag{F17.1}
\]
If $q>0$, there is a canonical class
\[
 \beta_i\in
 (\pi_1^Q\otimes p_N)_*\operatorname{CH}^2(Q\times J)
 \subset\operatorname{CH}^2_{(0)}(Q\times J)_{\rm hom}.
 \tag{F17.2}
\]
Here canonical means independent of the rational splitting and
Poincar\'e normalizations used in its construction, for the given
morphism $i$.  For every pointed curve $(C,c_0)$, every residual
family map in (F15.5) is obtained by pullback of $\beta_i$,
followed by $\pi_1^C\otimes p_N$, along the pointed map
$C\to Q$ induced by the homomorphism
$\operatorname{Jac}(C)\to Q$, with rational linearity when a
denominator is present.  Thus $\beta_i=0$ kills all families with
this image surface.

Conversely, if the residual family maps vanish for all curves,
then $\beta_i=0$.  It suffices to test a single smooth curve in
$Q$ whose induced Jacobian homomorphism surjects onto $Q$.

Partial Fourier transform on $Q$ identifies the whole component
in (F17.2) with
\[
 \operatorname{Hom}_{\rm CHM}
       \bigl(h^1(Q^\vee)(-1),N\bigr).
 \tag{F17.3}
\]
A polarization can replace $Q^\vee$ by $Q$ up to isogeny.
For an isogeny decomposition $Q\sim\prod_a Q_a^{m_a}$, with
$Q_a$ simple, the component is the corresponding direct sum
of the components for $Q_a$, repeated $m_a$ times.

This describes the remaining classes exactly; it does not prove
that (F17.2) vanishes.  C44-F16 restricts the relevant surfaces
to singular ample surfaces with $q(D)>3$.
\end{proposition}

\begin{proof}
Choose a rational section $s:Q\to P$ of the quotient map
$P\to Q$, in the isogeny category.  The normalized Poincar\'e
line bundle on $P\times D$, pulled back by $s\times1_D$,
defines a rational divisor class
\[
                         \mathcal P_s\in\operatorname{CH}^1(Q\times D).
 \tag{F17.4}
\]
This is well-defined for a rational homomorphism $s$: if $ms$
is an actual homomorphism, take $m^{-1}c_1(\mathcal P_{ms})$.
The homomorphism--Poincar\'e identification is additive, so the
answer is independent of $m$.  Define
\[
 \beta_i=(\pi_1^Q\otimes p_N)_*(1_Q\times i)_*\mathcal P_s.
 \tag{F17.5}
\]
Pullback of the normalized divisor by $[n]\times1_D$ multiplies
it by $n$.  Thus the source factor in (F17.5) already belongs
to the exact $\pi_1^Q$ eigensummand.

If $s'$ is another rational section, $s'-s$ has values in $B$.
Lift this difference up to isogeny to
$\operatorname{Pic}^0(J)$.  Its Poincar\'e divisor on $Q\times D$
is pulled back from $Q\times J$, up to the two single-factor
terms.  The proof of (F15.9) applies with the abelian parameter
$Q$ in place of a curve: it uses only the projection formula
on the target and $p_N L_d\pi_1^J=0$.  Hence its projected
pushforward is zero.  A change of normalization adds a divisor
pulled back from $Q$ or from $D$.  The first is killed by
$p_N\operatorname{CH}^1(J)=0$, and the second by
$\pi_1^Q[Q]=0$.  This proves independence of all these choices.

Let $\phi:\operatorname{Jac}(C)\to Q$ be an actual homomorphism
and $a_C:C\to\operatorname{Jac}(C)$ the pointed Abel map.
Pullback by $(\phi a_C)\times1_D$ of $\mathcal P_s$ is exactly
the cross divisor for $s\phi$.  Its pushforward commutes with
this pullback on the other factor.  After applying the curve
projector it therefore gives
$\overline\Psi_{C,i}(\phi)$, by F15's independence of the
lift.  This also explains the equality for rational $\phi$:
multiply by an integer and divide, using additivity of the
Poincar\'e class.  No map of points into a Chow group is being
assumed to have an injective algebraic representative.

For the converse choose a smooth complete-intersection curve
$C\subset Q$ when $q\ge2$, and take $C=Q$ when $q=1$.
Its Jacobian maps surjectively to $Q$, since its cohomological
restriction on $H^1$ is injective by weak Lefschetz.
After choosing a basepoint and translating the map, the induced
map of Picard motives
$h^1(Q)\to h^1(C)$ is split injective in rational Chow motives:
the homomorphism of Jacobians splits up to isogeny, and the
degree-one motive functor identifies these homomorphisms with
the corresponding Chow-motive maps.  Applying the splitting on
the first tensor factor recovers any class in (F17.2) from its
projected curve pullback.  If that projected pullback is zero,
$\beta_i=0$.
Translation causes no problem: the component used is $h^1$,
and the pointed Poincar\'e description gives the same cross
homomorphism, with only the already killed normalization terms.

The total Chow--K\"unneth index of $\beta_i$ is $1+3=4$,
so its Beauville grade in codimension two is zero.  In Betti
cohomology the component is
\[
                  H^1(Q)\otimes H^3(N).
 \tag{F17.6}
\]
Its Hodge types are $(4,0),(3,1),(1,3),(0,4)$, with possible
multiplicities, and it has no $(2,2)$ part.  Consequently every
algebraic class in this component, in particular $\beta_i$,
is homologically trivial.

For completeness, the source transpose matters.  As a
correspondence, the cycle in (F17.2) has right projector
$\pi_{2q-1}^Q$, not $\pi_1^Q$.  Partial Fourier transform
interchanges the first-factor cycle indices $1$ and $2q-1$.
It increases codimension by $q-1$, giving
\[
 \widehat\beta_i\in\operatorname{CH}^{q+1}(Q^\vee\times J),
 \qquad
 \widehat\beta_i=p_N\circ\widehat\beta_i\circ\pi_1^{Q^\vee}.
 \tag{F17.7}
\]
Thus it is a map $h^1(Q^\vee)\to N(1)$, equivalently the map
in (F17.3).  These are exact Chow assertions.  The Fourier
correspondence is given by the exponential of the Poincar\'e
divisor and has inverse, up to sign, inversion followed by the
dual Fourier transform.  Its component on $h^1$ is the
$(2q-1)$st power of that divisor divided by $(2q-1)!$.
The eigensummand identities can also be checked on cohomology
and lifted to Chow by O'Sullivan, since the Poincar\'e divisor,
Fourier components, and canonical projectors are distinguished.
Finally the degree-one motive of a product is the sum of the
degree-one motives of its factors, and isogenies are isomorphisms
of rational Chow motives.  This proves the last assertion.
\end{proof}

\subsection{One CM elliptic factor: half of the tensor is killed}

Let $E/k$ have $j(E)=1728$, and let $I=[i]^*$ on $h^1(E)$.
Over the coefficient field $K=\mathbf Q(i)$ define
\[
 e_+=\tfrac12(\pi_1^E-iI),\qquad
 e_-=\tfrac12(\pi_1^E+iI),\qquad
 U=(E,e_+),\quad\overline U=(E,e_-).
 \tag{F18.1}
\]
The two lines have types $(1,0)$ and $(0,1)$ after an ordering
of the embeddings.  For $m\ge1$, the sum of the two projectors
$e_+^{\otimes m}+e_-^{\otimes m}$ has rational coefficients.
Let $N_m$ denote the rational motive it defines on $E^m$.
Thus $N_m\otimes K=U^{\otimes m}\oplus\overline U^{\otimes m}$,
and $N_m$ has rank two and types $(m,0),(0,m)$.

\begin{proposition}[C44-F18]\label{C44-F18}
Over $k=\overline{\mathbf Q}$, the Fermat motive $N$ is
isomorphic to $N_3$ by an actual distinguished Chow
correspondence.  There is a rational Chow-motive isomorphism
\[
               h^1(E)\otimes N
                 \simeq N_4\oplus (N_2\otimes\mathbb L).
 \tag{F18.2}
\]
The second summand has zero codimension-two Chow group:
\[
              \operatorname{CH}^2(N_2\otimes\mathbb L)
                   =\operatorname{CH}^1(N_2)=0.
 \tag{F18.3}
\]
Consequently
\[
 (\pi_1^E\otimes p_N)_*\operatorname{CH}^2(E\times J)
                    \simeq\operatorname{CH}^2(N_4).
 \tag{F18.4}
\]
This kills the whole rational rank-two component with Hodge
types $(3,1),(1,3)$; it leaves exactly the rank-two extreme
component with types $(4,0),(0,4)$.

If $Q$ in F17 has an isogeny factor $E^r$, the same
elimination applies to all $r$ copies of this factor.  In
particular, when $Q\sim E^r$, vanishing of
$\operatorname{CH}^2(N_4)$ is sufficient for all residual
families from this $Q$.  Its vanishing has not been proved
here.  The isogeny $J\sim E^3$, CM splitting, and Kimura
nilpotence alone do not prove it.
\end{proposition}

\begin{proof}
The carried Fermat decomposition JAC-P17 gives $J\sim E^3$
over $k$.  Isogeny graphs transport $p_N$ to a distinguished
projector on $h^3(E^3)$.  In $H^3(E^3)$ its two extreme Hodge
lines are exactly $U^{\otimes3}$ and $\overline U^{\otimes3}$;
all remaining summands have types $(2,1),(1,2)$.
The transported projector and the projector defining $N_3$
are therefore identical in cohomology.  Their difference is
distinguished and numerically trivial, so O'Sullivan makes
them identical in rational Chow.  This proves the asserted
actual isomorphism, rather than merely a realization isomorphism.

The polarized pairing on $h^1(E)$ gives
\[
                         U\otimes\overline U\simeq\mathbb L_K.
 \tag{F18.5}
\]
One can verify (F18.5) without a semisimplicity assumption for
the Chow category: take the projected cup-product pairing and
its projected dual diagonal, and rescale their nonzero
composition on $\mathbb L_K$.  They are distinguished; the
other composite is the identity in cohomology on the rank-one
summand, hence the identity in Chow by O'Sullivan, with
$K$-coefficients.  The extension of the latter theorem from
$\mathbf Q$ to $K$ here just tensors an injective map by the
field $K$.

Expanding the left side of (F18.2) over $K$ gives the four
summands
\[
 U^4,\quad\overline U^4,\quad
 U\otimes\overline U^3,\quad
 \overline U\otimes U^3.
 \tag{F18.6}
\]
Using (F18.5), the last two become
$\mathbb L_K\otimes\overline U^2$ and
$\mathbb L_K\otimes U^2$.  The pairing constructions and
tensor permutations commute with conjugation of the coefficient
field, so the two conjugate pairs descend to the rational
isomorphism (F18.2).  All maps used are actual algebraic
correspondences.  Tensor permutations may change a pairing
normalization by a sign; no positivity claim is used and the
isomorphism and vanishing are unaffected.

The motive $N_2$ is a summand of $h^2(E^2)$ with only
types $(2,0),(0,2)$.  On an abelian surface the degree-two
part of its divisor Chow group is its N\'eron--Severi group
tensored with $\mathbf Q$: the homologically trivial divisors
are in $\operatorname{Pic}^0$, on which the exact projector
is $\pi_1$.  Thus the cycle class map is injective on
$(\pi_2^{E^2})_*\operatorname{CH}^1(E^2)$.
The $N_2$ part has zero divisor cohomology by its Hodge types,
and therefore $\operatorname{CH}^1(N_2)=0$.
The Tate twist identity proves (F18.3); applying Chow groups
to (F18.2) proves (F18.4).

Finally the isogeny splitting on $h^1(Q)$ in F17 is additive.
Apply the preceding argument separately to each elliptic
$E$ summand.  For any other simple factor $Q_a$ no vanishing
follows just from the absence of $(2,2)$ classes in
$H^1(Q_a)\otimes H^3(N)$.  Even for $Q_a=E$, the remaining
$N_4$ component sits in the allowed middle, grade-zero part
of codimension two on the abelian fourfold $E^4$.
Neither the threefold Murre(D) results nor divisor
cycle-class injectivity applies to that component.
\end{proof}

\paragraph{What has and has not changed.}
F17 is the exact universal formulation of F16's residual
family obstruction; it is not counted by itself as a new
main-gate advance.  F18 does annihilate an additional specified
CM subcomponent in rational Chow.  It does not show that
every excess quotient has this CM factor, or that the
remaining extreme $N_4$ cycles vanish.  The full AJ5
finite-generation gate of C44-F07 remains open, and even
vanishing of the entire algebraic radical would leave its
possible non-algebraically-trivial kernel to be addressed.

\paragraph{Primary source locators and scope.}
The Fourier input is Beauville, \emph{Quelques remarques sur
la transformation de Fourier dans l'anneau de Chow d'une
vari\'et\'e ab\'elienne}, Lecture Notes in Mathematics~1016
(1983), 238--260: Proposition~1, printed p.~241, for the
degree interchange, and Proposition~3(i), printed p.~243,
for the inverse identity on rational Chow groups.
Primary author PDF:
\url{https://math.univ-cotedazur.fr/u/beauvill/pubs/fourier.pdf}.
The passage from cohomological projector identities to
distinguished Chow identities uses the exact O'Sullivan
locators recorded in F15.  The quotient Poincar\'e construction
uses Bost--Charles \S3.2, also checked in F15.

For the low-dimensional boundary, the carried source audit
of Vial's threefold theorem and Xu--Xu's Corollary~3.9 is
retained.  Vial, \emph{Remarks on motives of abelian type},
author version 20 July 2015, introductory Theorem~3,
printed p.~2, still makes Murre(D) for products of curves
an additional premise.  It does not prove the grade-zero
codimension-two injection on $E^4$ used in the unproved
last step above.
Primary author PDF:
\url{https://www.math.uni-bielefeld.de/~vial/abelian.pdf}.

% END work/arithmetic/fermat_residual_motive.tex

% BEGIN work/arithmetic/fermat_relative_transfer.tex
% C44_POS, /root/fermat_aj. 2026-11-22 UTC.
% This file proves the relative geometric input for the root's
% polarization-comparison argument. It does not assert that arbitrary
% abstract motivic injections preserve a chosen intersection form.

\section{Relative Fermat correspondences with good reduction}

\subsection{The exact geometric input}

\begin{proposition}[C44-F08]\label{C44-F08}
Fix an integer $d\geq2$, put
\[
 S_d=\operatorname{Spec}\mathbf Z[1/d,\zeta_{2d}],\qquad
 \epsilon=\zeta_{2d},\qquad \epsilon^d=-1,
\]
and write
\[
 F_n=\{x_0^d+\cdots+x_{n+1}^d=0\}\subset\mathbf P^{n+1}_{S_d},
 \quad C=F_1,\quad J=\operatorname{Pic}^0_{C/S_d}.
\]
There is a finite direct sum
\[
 M_n=\bigoplus_{(s,t)}h(J^{s}/S_d)(-t)^{\oplus a_{n,s,t}},
 \qquad a_{n,s,t}\in\mathbf Z_{\geq0},\quad s,t\geq0,
 \tag{F08.1}
\]
and relative rational Chow correspondences
\[
 \alpha_n:h(F_n/S_d)\longrightarrow M_n,
 \qquad
 \beta_n:M_n\longrightarrow h(F_n/S_d),\qquad
 \beta_n\alpha_n=\operatorname{id}.
 \tag{F08.2}
\]
Here $J^0=S_d$ and the convention is
$\mathbb L=\mathbf1(-1)$. All schemes in this assertion are smooth
and projective over $S_d$; $J$ is an abelian scheme. The same conclusion
holds for every finite product of the $F_n$, with the corresponding
products of Jacobians. These identities specialize to every geometric
fibre of characteristic $p\nmid d$, including characteristic $2$
when $d$ is odd.

In particular, for every such fibre and every $r$, an algebraic
correspondence extending over the same characteristic-zero model gives
a split injection
\[
 H^{2r}(F_n)\longrightarrow
 \bigoplus_{s,t}H^{2r-2t}(J^s)(-t)^{\oplus a_{n,s,t}}.
 \tag{F08.3}
\]
This assertion and its proof are unconditional. Conditional input 032
is needed only if the root's separate positive-form comparison is
used to infer the Hodge standard conjecture on the special fibre.
No assertion from Milne's unrefereed report is used here.
\end{proposition}

\begin{proof}
All the Fermat schemes are smooth over $S_d$: since $d$ is a unit,
the partial derivatives cannot vanish simultaneously at a projective
geometric point. The scheme $F_0$ consists of $d$ disjoint sections,
because the roots of $u^d=-1$ lie in the base ring and their pairwise
differences are units after $d$ is inverted. For example, this last
assertion follows from the discriminant of $u^d+1$, which is a unit
times $d^d$. Consequently
\[
 h(F_0/S_d)=\mathbf1^{\oplus d}.
 \tag{F08.4}
\]
The curve has the section $e=[1:\epsilon:0]$. Its relative Jacobian
therefore exists as an abelian scheme, and its normalized universal
Poincar\'e bundle exists on $C\times_{S_d}J$.

\paragraph{The relative Shioda--Katsura map.}
For $n\geq2$, use coordinates $[x_0:\cdots:x_n]$ on $F_{n-1}$
and $[y_0:y_1:y_2]$ on $C$. The rational map is
\[
 \phi_n:F_{n-1}\times C\dashrightarrow F_n,
 \quad (x,y)\longmapsto
 [x_0y_2:\cdots:x_{n-1}y_2:
       \epsilon x_ny_0:\epsilon x_ny_1].
 \tag{F08.5}
\]
Indeed, the sum of the $d$-th powers of the displayed coordinates is
\[
 y_2^d\sum_{i=0}^{n-1}x_i^d
      -x_n^d(y_0^d+y_1^d)
 =-y_2^dx_n^d+x_n^dy_2^d=0.
\]
Its base ideal is precisely the ideal $(x_n,y_2)$, with the equations
understood after local trivialization of the two hyperplane bundles.
To check this near the proposed centre, at least one $x_i$ with
$i<n$ is a unit in a projective chart and at least one $y_j$ with
$j<2$ is a unit. The sections defining the map then generate both
$y_2$ and $x_n$. Outside that centre at least one defining section
is nonzero. The centre is
\[
 Z_n=\{x_n=0\}\times\{y_2=0\}
       =F_{n-2}\times F_0
       \subset P_n:=F_{n-1}\times C.
 \tag{F08.6}
\]
It is smooth of relative codimension two. Hence
\[
 b_n:Y_n=\operatorname{Bl}_{Z_n}(P_n)\longrightarrow P_n
\]
is smooth and projective over $S_d$. The defining sections generate
an invertible ideal on the blowup, so (F08.5) extends to a morphism
$f_n:Y_n\longrightarrow F_n$.

This morphism has generic degree $d$ on every geometric fibre. To
see the degree explicitly, work on the nonempty target open set
\[
 A:=\sum_{i=0}^{n-1}z_i^d\ne0.
\]
On the exceptional divisor the first block of coordinates is a
scalar multiple of a point of $F_{n-2}$, so $A=0$ there. The
proper transforms of $x_n=0$ and $y_2=0$ also map into $A=0$.
Consequently every preimage on this open set has $x_ny_2\ne0$. Normalize
$x_n=y_2=1$. A preimage of a representative
$[z_0:\cdots:z_{n+1}]$ is then determined by
\[
 x_i=\lambda z_i\ (i<n),\qquad
 y_j=\lambda\epsilon^{-1}z_{n+j}\ (j=0,1),
 \qquad \lambda^d=-A^{-1}.
 \tag{F08.7}
\]
Both Fermat equations follow from this one scalar equation and the
target equation. Its $d$ solutions are distinct because $d$ is
invertible. They give distinct points, since both projective
normalizations have been fixed. Thus the map is finite \'etale of
degree $d$ on this open set, as also established in
Shioda--Katsura, Lemmas~1.1--1.3. Projection formula, or directly
composition of graph correspondences, now gives
\[
 (f_n)_*f_n^*=d\,\operatorname{id}_{h(F_n/S_d)}.
 \tag{F08.8}
\]
There is no assertion that $f_n$ is finite everywhere. The identity
is valid for a proper generically finite morphism between these
smooth schemes; contracted exceptional strata introduce no term in
the composition $(f_n)_*f_n^*$.

\paragraph{The codimension-two blowup correspondence.}
Let $j:E_n\hookrightarrow Y_n$ and $p:E_n\to Z_n$ be the exceptional
divisor and its projection. The relative blowup formula is
\[
 \Phi=(b_n^*,j_*p^*):
 h(P_n/S_d)\oplus h(Z_n/S_d)(-1)
       \xrightarrow{\ \sim\ }h(Y_n/S_d).
 \tag{F08.9}
\]
In this codimension-two case its inverse is explicitly
\[
 \Psi=((b_n)_*,-p_*j^*).
 \tag{F08.10}
\]
For completeness the four components of $\Psi\Phi$ can be checked
without a general inverse formula. They use
$(b_n)_*b_n^*=1$, $p_*p^*=0$ for the projective-line bundle $p$,
and
\[
 j^*j_*=-\xi\cdot,
 \qquad \xi=c_1(\mathcal O_{E_n}(1)),\qquad
 p_*(\xi p^*a)=a.
\]
The off-diagonal components are zero and the two diagonal
components are identities. The blowup formula for Chow groups
states that $\Phi$ is surjective as well. It holds after product
with every smooth projective test scheme over $S_d$; thus the
relative Manin identity principle gives $\Phi\Psi=1$ as Chow
correspondences. These are the usual intersection-theoretic
blowup identities over a regular base. Over a field they are
displayed, with an explicit inverse, in Vial, Proposition~5.1 and
Theorem~5.3. The same proof uses the relative intersection theory
of Fulton, \S\S20.1--20.2; no resolution of singularities is used.

Combining (F08.8)--(F08.10) gives a concrete split injection
\[
 h(F_n/S_d)\ 
 \underset{\frac1d(f_n)_*\Phi}{\stackrel{\Psi f_n^*}{\rightleftarrows}}\ 
 h(F_{n-1}/S_d)\otimes h(C/S_d)
       \oplus h(F_{n-2}/S_d)(-1)^{\oplus d}.
 \tag{F08.11}
\]
In particular, the projector onto the image in the right-hand
side is
\[
 \frac1d\,\Psi f_n^*(f_n)_*\Phi.
 \tag{F08.12}
\]
It is idempotent by (F08.8). This supplies the requested explicit
Chow projector rather than only an injection of realizations.

\paragraph{Induction to powers of the curve.}
Starting with (F08.4) and $h(F_1)=h(C)$, (F08.11) shows that
$h(F_n)$ is a retract of
\[
 \bigoplus_{t=0}^{\lfloor n/2\rfloor}
     h(C^{n-2t}/S_d)(-t)^{\oplus m_{n,t}},
 \tag{F08.13}
\]
where the nonnegative integers may be chosen recursively as
\[
 m_{0,0}=d,\quad m_{1,0}=1,\quad
 m_{n,t}=m_{n-1,t}+d\,m_{n-2,t-1},
 \tag{F08.14}
\]
with all unspecified indices zero. For example, the right sides
for $n=2,3$ are respectively
$h(C^2)\oplus\mathbb L^{\oplus d^2}$ and
$h(C^3)\oplus h(C)(-1)^{\oplus(d^2+d)}$.
The retractions are composites and tensor products of the
specified correspondences, so remain defined over the original
base.

\paragraph{A direct relative curve-to-Jacobian splitting.}
We record a construction that avoids needing a separate relative
isomorphism theorem for $h^1(C)$ and $h^1(J)$. The section $e$ has
trivial normal bundle. In the affine chart $x_0=1$, write
$u=x_1,v=x_2$ and $1+u^d+v^d=0$. Along $e$ the derivative
$du^{d-1}=d\epsilon^{d-1}$ is a unit, so the conormal of $e$
is freely generated by $v$. Consequently the standard relative
curve projectors
\[
 \pi_0=[e\times C],\qquad \pi_2=[C\times e],\qquad
 \pi_1=[\Delta_C]-\pi_0-\pi_2
 \tag{F08.15}
\]
are orthogonal idempotents. The only extra term that might occur
over a positive-dimensional base in their composition is the
self-intersection of the section, and it is zero by the trivial
normal bundle. Thus
\[
 h(C/S_d)=\mathbf1\oplus h^1(C/S_d)\oplus\mathbb L.
 \tag{F08.16}
\]

Let $a:C\to J$ send a point $x$ to
$\mathcal O_C(x-e)$, and let $\mathcal P$ be the universal
degree-zero line bundle on $C\times J$, normalized along
$e\times J$. The universal property of the relative Picard
scheme gives
\[
 (1_C\times a)^*\mathcal P
  \simeq\mathcal O_{C\times C}
              (\Delta_C-e\times C-C\times e).
 \tag{F08.17}
\]
Indeed, its restriction to a fibre of the second projection is
$\mathcal O_C(x-e)$, and the right-hand bundle is trivial along
$e\times C$, using the same normal-bundle trivialization. Those
two properties characterize the normalized Poincar\'e family.
Put $\gamma=c_1(\mathcal P)\in
\operatorname{CH}^1(C\times J)_{\mathbf Q}$, a degree-zero
correspondence from $C$ to $J$. Composition with the transpose
graph of $a$ is pullback in the second factor, so (F08.17) says
\[
 a^*\gamma=\pi_1.
 \tag{F08.18}
\]
Therefore
\[
 h^1(C/S_d)
 \underset{\pi_1a^*}{\stackrel{\gamma\pi_1}{\rightleftarrows}}
 h(J/S_d)
\]
is a split injection. The assertion also covers genus zero, when
the displayed primitive curve projector is zero. Combining with
(F08.16), $h(C)$ is a retract of
\[
 U=\mathbf1\oplus h(J)\oplus\mathbb L.
 \tag{F08.19}
\]
Expanding $U^{\otimes k}$ shows that each $h(C^k)(-t)$ is a
retract of a finite direct sum of terms $h(J^s)(-u)$ with
$s,u\geq0$. This turns (F08.13) into (F08.1)--(F08.2).

Finally, tensor products of split injections and retractions are
split injections and retractions. This proves the assertion for
products. Pullback to characteristic-zero fibres and to geometric
special fibres respects every graph, first Chern class, blowup
identity, and correspondence composition used above. Since no
prime outside those dividing $d$ was inverted, the same
construction works at each $p\nmid d$. This proves (F08.3).
\end{proof}

\subsection{Relative primitive projectors and scope of transfer}

\begin{lemma}[C44-F09]\label{C44-F09}
Let $n\geq1$, write $h=c_1(\mathcal O_{F_n}(1))$, and use rational
coefficients.
There is an explicit relative projector
\[
 p_n=[\Delta_{F_n}]
       -\frac1d\sum_{i=0}^{n}h^{n-i}\times h^i.
 \tag{F09.1}
\]
Its realization is the primitive middle cohomology of the Fermat
hypersurface. Algebraic relative correspondences represent the
Lefschetz operator and its lowering operator on $H^*(F_n)$.
For products of Fermat hypersurfaces, their tensor constructions
give the corresponding operators and primitive projectors for
every polarization that is a positive rational combination of the
factor hyperplane classes. These assertions are
unconditional and compatible with all the fibres of F08.
\end{lemma}

\begin{proof}
The $n+1$ coordinate sections $x_0,\ldots,x_n$ of
$\mathcal O_{F_n}(1)$ have no common zero on $F_n$. Thus
$h^{n+1}=0$ in the Chow ring. Moreover the degree of $h^n$ on
each fibre is $d$. The correspondences
\[
 q_i=\frac1d h^{n-i}\times h^i\quad(0\leq i\leq n)
\]
are consequently orthogonal idempotents, with image
$\mathbb L^i$. These assertions follow either by composing the
two factorized cycles or by the degree formula; an exponent
greater than $n$ contributes zero by $h^{n+1}=0$, and an exponent
less than $n$ contributes zero under proper pushforward to the
base. Formula (F09.1) is therefore idempotent.

Weak Lefschetz and Poincar\'e duality for a smooth hypersurface
identify all cohomology outside degree $n$ with these ambient
Tate classes. In degree $n$ they remove the ambient class when
$n$ is even. The remaining subspace is exactly the primitive
middle cohomology. This identification is valid in the complex
fibre and, by smooth proper base change and the relative
hyperplane class, in every geometric good fibre.

On the span of $1,h,\ldots,h^n$, set
\[
 L(h^i)=h^{i+1},\qquad
 \Lambda(h^i)=i(n-i+1)h^{i-1},
 \tag{F09.2}
\]
with the endpoint values zero. Set both operators zero on the
image of $p_n$. The maps between the displayed Tate summands
are algebraic relative correspondences. In realization, this
$L$ equals cup product with $h$: primitive middle classes are
annihilated by $h$. The operator $\Lambda$ is the usual
$\mathfrak{sl}_2$ lowering operator. In particular, inverse
Hard Lefschetz maps are algebraic and extend over $S_d$.

For a product polarization $\sum_j a_jh_j$, with rational
$a_j>0$, take $L=\sum_j a_jL_j$ and
$\Lambda=\sum_j a_j^{-1}\Lambda_j$ on the tensor product.
Their commutator is the sum of the factor degree operators,
so these are the raising and lowering operators for the same
$\mathfrak{sl}_2$ representation. The first realizes cup
product by the specified polarization.
Finite-dimensional $\mathfrak{sl}_2$ representation theory over
$\mathbf Q$ expresses the Lefschetz decomposition projectors
and the inverses of each nonzero Lefschetz map as rational
combinations of these operators and degree projectors. This
supplies algebraic relative correspondences with the stated
realizations.

We do not assert that the diagonal cup correspondence composed
with $p_n$ vanishes in the Chow group. The operators just
constructed have the required cohomological actions; the
possible homologically trivial difference is irrelevant to
the claimed realization statements.
\end{proof}

\begin{corollary}[C44-F10: Fermat positivity under the stated abelian input]
\label{C44-F10}
Let $k$ be an algebraic closure of a finite field of characteristic
$p$, let $p\nmid d$, and let
\[
 X=\prod_{j=1}^{m}F_{n_j,d,k},\qquad n_j\geq1.
\]
Conditional on the unrefereed 032/Milne 2026 inputs specified in
C44-A08, the Hodge standard conjecture holds for $X$ with
respect to every ample class $h=\sum_j a_jh_j$, where
$a_j\in\mathbf Q_{>0}$ and $h_j$ is the pullback of the Fermat
hyperplane class. In particular this proves it for every power
of every positive-dimensional Fermat hypersurface with its product hyperplane
polarization. If every $n_j\geq3$, the conclusion holds for
every ample class of $X$.

This deduction uses C44-A08 and F08--F09. It does not use
Milne2022 Proposition~4.2, a finite-cover positivity transfer,
or a claim that all algebraic cycles of the special fibre lift.
\end{corollary}

\begin{proof}
Choose a closed point of $S_d$ above $p$ and localize at that
point. The result is a mixed-characteristic discrete valuation
ring whose fraction field embeds into $\mathbf C$. Base
extension to its geometric special fibre gives the specified
Fermat schemes over $k$. After multiplying $h$ by a positive
integer, it is the class of a relatively ample line bundle on
this model. Scaling by such an integer changes every relevant
primitive intersection form by a positive scalar and hence
does not affect the asserted sign.

For each even degree $2r\leq\dim X$, take the full primitive
summand for this relative ample class. By F09 its projector
and inverse Lefschetz map are represented by algebraic
correspondences over the same base. Poincar\'e duality and
the inverse Lefschetz map give the inverse of the restricted
primitive pairing, also algebraically. F08 gives a split
injection of the entire cohomology of $X$ into the stated
Tate-twisted cohomology of abelian schemes, and hence an
injection on this primitive summand. All the hypotheses of
C44-A08 hold. Its conclusion is precisely positivity of
\[
 (-1)^r\int_X xyh^{\dim X-2r}
\]
on the rational algebraic primitive classes, with exactly the
conditional input recorded there. This proves the first
assertion for all primitive degrees.

Suppose now that every $n_j\geq3$. Weak Lefschetz gives
$H^1(F_{n_j,d,k})=0$ and
$H^2(F_{n_j,d,k},\mathbf Q_\ell)=\mathbf Q_\ell h_j$.
The N\'eron--Severi cycle map is injective after tensoring
with $\mathbf Q_\ell$, so each factor has rational
N\'eron--Severi group spanned by its hyperplane class.
K\"unneth and $H^1=0$ give the same assertion for the product:
its rational N\'eron--Severi group has basis the $h_j$.
Thus any rational ample class is $\sum_j a_jh_j$ with
rational coefficients. Restricting it to a curve in the
$j$-th factor with all other coordinates fixed shows
$a_j>0$. Such a class is covered by the first assertion.
For factors of dimension two, extra divisor classes can
occur and this last argument is not asserted. For products
of curves, cross terms in the N\'eron--Severi group can also
occur; the first assertion remains as stated.
\end{proof}

\paragraph{What these results do and do not transfer.}
F08 provides a relative split injection into cohomology of
abelian schemes with good reduction. F09 provides the relative
primitive projectors and algebraic inverse Lefschetz maps.
The positive-form comparison in C44-A08 compares the
chosen primitive cup form with the pullback of the abelian
Weil-star form, using positivity of the generic complex Hodge
polarizations and constancy of the comparison endomorphism's
characteristic polynomial. Its use of 032/Milne 2026 remains conditional
on those supplied unrefereed preprints. No implication from an
arbitrary abstract motivic injection to preservation of a
chosen positive form is asserted here.

The distinction between finite covers and this construction is
essential: $f_n$ contracts strata. Pullback of an ample class
can be nef without being ample. The proof of F08 never uses
positivity of that pullback, and F09 uses the original ample
hyperplane classes instead. Thus neither an unproved nef limit
argument nor Milne2022 Proposition~4.2 is a hidden input to the
geometric statements.

\subsection{Primary-source locators}

\begin{itemize}
\item T. Shioda and T. Katsura, \emph{On Fermat varieties},
Tohoku Math. J. \textbf{31} (1979), 97--115:
Lemma~1.1 and formula~(1.3), p.~98, give the degree-$d$ map;
Lemma~1.2, p.~99, resolves it by the specified blowup;
Lemma~1.3, p.~100, gives the finite \'etale open part;
Theorem~1.7 and diagram~(1.25), pp.~101--102, give the full
blowup/quotient/blowdown construction. Most decisively,
Remark~1.8, p.~102, explicitly states that the construction
works over $\mathbf Z[1/d,\zeta_{2d}]$ in the category of
smooth projective schemes. The primary PDF was opened and
Remark~1.8 also checked in the page image:
\url{https://www.jstage.jst.go.jp/article/tmj1949/31/1/31_1_97/_pdf}.
\item C. Vial, \emph{Algebraic cycles and fibrations},
Doc. Math. \textbf{18} (2013), 1521--1553. The author PDF
consulted has 24 pages: Proposition~5.1 on p.~13 is the Chow
blowup formula; Theorem~5.3 on p.~14 verifies the explicit
motivic inverse by Manin's identity principle. It is a
field-base source; the relative use in F08 is via the same
intersection-theoretic identities over the regular base,
not an expansion of Vial's stated hypotheses.
\url{https://www.math.uni-bielefeld.de/~vial/algcycles-fibrations.pdf}.
\item W. Fulton, \emph{Intersection theory}, Proposition~6.7
and \S\S20.1--20.2, for the blowup identities and intersection
theory over a regular base. The applicability of this
relative framework to Chow motives over Dedekind bases is
also stated in J. Laga and A. Shnidman, \emph{Ceresa cycles of
bielliptic Picard curves}, \S4.2, citing K. K\"unnemann,
\emph{Arakelov Chow groups of abelian schemes}, \S1,
Remark~1.1. The supplied Laga--Shnidman source was checked;
its special Picard-type motive isomorphism is not used.
\item V. Guletski\u\i, \emph{Chow motives of abelian type over a
base}, Eur. J. Math. \textbf4 (2018), 1065--1086, \S3
defines relative rational Chow correspondences over a regular
Dedekind base, citing precisely Fulton \S\S20.1--20.2 and
K\"unnemann Remark~1.1, p.~367. The paragraph immediately
before Lemma~3.1 identifies the section self-intersection
correction to the two curve Tate projectors. In F08 this
correction is zero by the explicit trivialization of the
Fermat section's normal bundle. Section~2 gives the normalized
relative Poincar\'e bundle framework.
\url{https://link.springer.com/article/10.1007/s40879-018-0270-9}.
\end{itemize}

% END work/arithmetic/fermat_relative_transfer.tex

% BEGIN work/arithmetic/fermat7_frobenius.tex
\section{C44-F11: a degree-seven Fermat block with extra special classes}
\label{sec:C44-F11}
\paragraph{Statement and dependency flags.}
Let \(X/\mathbb F _2\) be the Fermat fourfold
\[
 X:\quad x_0^7+\cdots+x_5^7=0\subset\mathbb P^5,\qquad T=X\times X,
\]
and put \(K=\mathbb Q(\zeta _7)\). Work over \(\mathbb F _8\) when using
the diagonal automorphisms. There is an explicit relative rational Chow
projector \(e_\alpha\), of cohomological rank six, with
\(\alpha=(1,1,1,3,3,5)\), such that:
\begin{enumerate}
\item In the characteristic-zero Fermat lifting its Hodge types are
\((3,1),(2,2),(1,3)\), each twice, and it has no nonzero rational Hodge
class. On the special fibre Frobenius \(8\) is exactly \(8^2\) on this
whole rank-six block.
\item The fourfold is not supersingular: the character
\(\beta=(1,1,1,1,1,2)\) has normalized Newton slopes \(1,3\).
\item The rank-36 block \(W=e_\alpha\otimes e_\alpha\) in the primitive
middle cohomology of \(T\) splits into six rational rank-six blocks
\(W_r\), \(r=1,\ldots,6\). Only \(W_6\) has rational Hodge classes on the
generic fibre; its entire rank-six space is spanned by explicit
correspondence classes.
\item \emph{Conditional on the supplied 032/Milne-2026 input through
Tate for finite-field abelian varieties}, the special block \(e_\alpha\)
and hence \(W\) are spanned by algebraic cycle classes. Thus \(W\)
contains thirty independent special algebraic directions, in
\(\bigoplus_{r=1}^5W_r\), which cannot be specializations of generic
algebraic cycle classes in the fixed Fermat lifting, even after finite
extension of its base trait.
\end{enumerate}
The first three assertions are unconditional. The last assertion does
not construct equations for the thirty cycles. It proves their
existence and independence under the stated conditional input. It makes
no assertion about lifting those classes in a different deformation.
C44-F13 below strengthens this algebraicity assertion unconditionally
by explicit ruled surfaces.

\begin{proof}
We use the relative Fermat model over
\(\mathbb Z[1/7,\zeta _{14}]\) of C44-F08 and a mixed-characteristic trait
above \(2\). The group \(G=\mu _7^6/\mu _7\) acts diagonally.
Its primitive degree-four characters are the tuples
\(a=(a_0,\ldots,a_5)\), with each \(a_i\in\{1,\ldots,6\}\) and
\(\sum a_i\equiv0\pmod7\). Each occurs once. For a character \(\chi_a\),
the usual group-ring idempotent
\[
 e_a=\frac1{|G|}\sum_{g\in G}\chi_a(g)^{-1}g^*
\]
is defined after adjoining coefficient field \(K\). The orbit sum
\[
 e_{[a]}=\sum_{t=1}^{6}e_{ta}
\]
has rational coefficients and is an actual relative Chow idempotent.
Every such orbit has six elements, since one nonzero coordinate already
forces the stabilizer in \((\mathbb Z/7)^\times\) to be trivial.
The projector is self-adjoint: transposition takes \(e_a\) to \(e_{-a}\).
The group algebra acts on its realization through the field \(K\), and
that realization is a rank-one \(K\)-module.

The classical Fermat character formula is
\[
 q(a)=\frac{\sum_i\langle a_i\rangle}{7}-1,\qquad
 H_a\subset H^{4-q(a),q(a)}.
\]
Its primary locator is Ran, \emph{Cycles on Fermat hypersurfaces},
Compositio Math. 42 (1980), Proposition 1.7, printed p.~125:
\url{https://www.numdam.org/article/CM_1980__42_1_121_0.pdf}.
For \(\alpha=(1,1,1,3,3,5)\) one obtains, in the order \(t=1,\ldots,6\),
\[
 (q(t\alpha))_{t=1}^6=(1,2,1,3,2,3).
\]
A rational Hodge class would generate, under \(\mathbb Q[G]\), a
nonzero rational subrepresentation consisting of Hodge classes.
The orbit representation is irreducible over \(\mathbb Q\), so this
would force its whole realization to have type \((2,2)\), contradicting
the displayed types. Thus its rational Hodge space is zero.

Here is an exact Jacobi-sum calculation. Write
\(\mathbb F _8=\mathbb F _2[\tau]/(\tau^3+\tau+1)\), choose
\(\chi(\tau)=\zeta _7\), and let
\(\psi(x)=(-1)^{\operatorname{Tr}_{\mathbb F _8/\mathbb F _2}(x)}\).
Set \(G_a=\sum_{x\in\mathbb F _8^\times}\chi(x)^a\psi(x)\).
Squaring \(x\) shows that \(G_a\) is constant on
\(H=\{1,2,4\}\) and on \(-H=\{3,5,6\}\). Put
\[
 A=\zeta _7+\zeta _7^2+\zeta _7^4,\quad
 B=\zeta _7^3+\zeta _7^5+\zeta _7^6,\quad \delta=A-B.
\]
The seven nonzero elements have additive-character values \(-1\) at
\(1\), \(+1\) at \(\tau,\tau^2,\tau^4\), and \(-1\) at
\(\tau^3,\tau^5,\tau^6\). Hence
\[
 A+B=-1,\quad AB=2,\quad \delta^2=-7,\qquad
 G_H=-1+\delta,\quad G_{-H}=-1-\delta,\quad G_HG_{-H}=8.
\]
The primitive Frobenius eigenvalue in dimension four is
\[
 j(a)=\frac{\prod_{i=0}^5G_{a_i}}8.
\]
This is the even-dimensional convention in Shioda--Katsura,
\emph{On Fermat varieties}, Tohoku Math. J. 31 (1979),
Section 3, pp.~107--108, formulas (3.1)--(3.10):
\url{https://www.jstage.jst.go.jp/article/tmj1949/31/1/31_1_97/_pdf}.
The five-variable convention and its identification with Frobenius
are also stated in van der Geer--Katsura,
\emph{On the height of Calabi--Yau varieties in positive characteristic},
Documenta Math. 8 (2003), Section 5, printed p.~106:
\url{https://ems.press/content/serial-article-files/25920}.
That paragraph treats arbitrary Fermat dimension and degree; the
Calabi--Yau restriction is imposed only on p.~107 after the Jacobi-sum
and valuation formulas. Equivalently, fix the first nonzero coordinate
to \(1\) and use the five-variable Jacobi sum; the product of the five
remaining characters is nontrivial, and \(G_aG_{-a}=8\) gives the formula.

The tuple \(\alpha\) has three entries in \(H\) and three in \(-H\).
Every multiple \(t\alpha\) retains this balance. Consequently
\[
 j(t\alpha)=8^3/8=64=8^2\qquad(t=1,\ldots,6).
\]
For \(\beta\), all six entries lie in \(H\), so
\[
 j(\beta)=(-1+\delta)^6/8=36-20\delta.
\]
It has polynomial \(z^2-72z+4096\). Its Newton polygon at \(2\), with
vertices \((0,12),(1,3),(2,0)\), gives valuations \(9\) and \(3\).
After division by \(v_2(8)=3\) these are slopes \(3\) and \(1\).
This proves nonsupersingularity.

The characters of \(W\) are all \((t\alpha,s\alpha)\). Their rational
orbits are
\[
 W_r:\quad\{(t\alpha,rt\alpha):t=1,\ldots,6\},
 \qquad r\in(\mathbb Z/7)^\times.
\]
They give six mutually orthogonal relative rational projectors.
Their Hodge second indices are
\[
\begin{array}{c|rrrrrr}
 r\backslash t&1&2&3&4&5&6\\ \hline
 1&2&4&2&6&4&6\\
 2&3&5&4&4&3&5\\
 3&2&5&3&5&3&6\\
 4&4&3&3&5&5&4\\
 5&3&3&2&6&5&5\\
 6&4&4&4&4&4&4
\end{array}
\]
The same irreducibility argument gives no rational generic Hodge
class for \(r\ne6\). The block \(W_6\) is entirely of type \((4,4)\).
Let \(g\) multiply the first coordinate of \(X\) by \(\zeta _7\).
The six projected graph correspondences \(e_\alpha\Gamma_{g^j}\),
\(j=0,\ldots,5\), are supported cohomologically in \(W_6\).
Under Poincare duality they realize \(1,g,\ldots,g^5\) on the rank-one
\(K\)-module. They are independent since \(g\) has minimal polynomial
\(\Phi_7\), of degree six; hence they span \(W_6\).
Each factor \(e_\alpha H^4(X)\) is primitive middle cohomology and is
annihilated by its hyperplane class. Thus \(W\subset P^8(T)\) for
every positive weighted product polarization.

For the last assertion only, impose the supplied conditional
032/Milne-2026 input. C44-F08 gives an actual algebraic injection and
retraction from the Fermat motive into a finite sum of Tate twists
of abelian motives. Tate surjectivity passes to the retract: inject
an invariant class, express its image by algebraic classes, and
apply the algebraic retraction. Since all of \(e_\alpha H^4(X)(2)\)
is Frobenius-\(8\)-invariant, it is spanned by six independent
algebraic classes. Their exterior products span all thirty-six
dimensions of \(W(4)\). The projectors for \(r=1,\ldots,5\) produce
the thirty-dimensional algebraic subspace in the statement.

Specialization commutes with the relative projectors. Every generic
algebraic class has zero projection to those five blocks by the Hodge
calculation. Its specialization therefore has zero projection there.
Every nonzero special algebraic class in their sum consequently
fails to lift as an algebraic cycle class in this Fermat model.
The same argument applies after every finite base extension.
\end{proof}

\paragraph{Status restriction.}
No primary source found here certifies that \(I(X\times X)\) for
this degree-seven example was previously unknown. It is therefore
not labelled an established \emph{B-known/I-unknown} example.
The following result proves a sign statement without claiming
historical priority.

\section{C44-F12: Frobenius symmetry removes the rank-six sign ambiguity}
\label{sec:C44-F12}
\paragraph{Statement and separation of equivalences.}
For \(X/\overline{\mathbb F}_2\) above and \(T=X\times X\):
\begin{enumerate}
\item Unconditionally, \(B(X)\), \(B(T)\), \(I_{\rm num}(X)\), and
\(I_{\rm num}(T)\) hold for every ample class.
\item Unconditionally, numerical and \(13\)-adic homological
equivalence agree on all algebraic cycles of \(X\) and \(T\).
Consequently \(I_{13}(X)\) and \(I_{13}(T)\) hold for every ample class.
\item Conditional on the supplied 032/Milne-2026 input, the same
cohomological statement holds for every \(\ell\ne2\), by C44-F10
(or by the algebraic abelian retraction and conditional \(D_\ell\)
for abelian varieties). C44-F14 below removes this condition for
\(X\), using additional explicit cycles; it does not remove it for \(T\).
\end{enumerate}
The unconditional proof uses the published Ancona--Marmora theorem,
not 032, Milne-2026, or Milne-2022 Proposition 4.2. This proof itself
does not assert Tate surjectivity for the block in C44-F11.

\paragraph{Primary source.}
Ancona--Marmora, \emph{The Hilbert symbol in the Hodge standard conjecture},
Comment. Math. Helv. 100 (2025), 619--662, DOI 10.4171/CMH/591:
\url{https://ems.press/content/serial-article-files/51595?nt=1}.
Section 2.3, pp.~623--624, uses homological motives obtained from
Chow motives. Theorem 3.2 and Remark 3.3, p.~626, give the signature
congruence and zero-or-full-dimension alternative; Lemma 3.4,
pp.~626--627, gives the full-rank comparison with realizations.
The matching author version is arXiv:2210.14001v3, 20 January 2025
(internal date 22 January), printed pp.~5--7:
\url{https://irma.math.unistra.fr/~marmora/2210.14001v3.pdf}.

\paragraph{Rank-six numerical sign lemma.}
Let \(M\) be a relative homological motive over the integers of a finite
extension of \(\mathbb Q_2\), with
\(q:\operatorname{Sym}^2M\to\mathbf1\). Suppose its realization has
dimension six, \(K=\mathbb Q(\zeta _7)\) acts on \(M\), the adjoint
on \(K\) is conjugation, and \(q_B\) polarizes its Hodge structure.
Suppose an algebraic endomorphism \(u\) of the special fibre is a
\(q\)-isometry on numerical classes and is semilinear for
\(\sigma_4:\zeta _7\mapsto\zeta _7^4\). Then \(q\) is positive
definite on numerical classes of the special fibre.

\begin{proof}
Numerical classes have finite dimension and descend to finite residue
extensions. Extend the residue field and trait enough to define a
basis. Put \(V=\operatorname{Hom}_{\rm NUM}(\mathbf1,M_k)\).
If \(V=0\), there is nothing to prove. Otherwise the \(K\)-action and
Ancona--Marmora Remark 3.3 give \(\dim_{\mathbb Q}V=6\), so \(V\)
has rank one over \(K\). Lemma 3.4 identifies \(q|_V\), after realization,
with the nondegenerate form of \(M_k\); it is therefore nondegenerate.

Choose a \(K\)-basis. Trace duality and the adjoint relation give
\[
 q(x,y)=\operatorname{Tr}_{K/\mathbb Q}(c\,x\bar y),
 \qquad c\in K^+\setminus\{0\}.
\]
Indeed represent \(q(x,1)\) by the trace pairing, use the adjoint
relation to recover the form, and use symmetry to obtain \(c=\bar c\).
Nondegeneracy forces \(c\ne0\). Over \(\mathbb R\) this is the orthogonal
sum of three planes indexed by embeddings \(\tau:K^+\hookrightarrow
\mathbb R\), with forms \(2\tau(c)\operatorname{Re}(z\bar w)\).
Each plane is definite, of sign \(\tau(c)\).

The restriction of \(\sigma_4\) to \(K^+\) has order three and acts
transitively on its real embeddings. The semilinear isometry \(u\)
therefore sends these three planes cyclically and their signs agree.
The negative index is \(0\) or \(6\). All four hypotheses of
Ancona--Marmora Theorem 3.2 hold, so this index is divisible by four.
It must be \(0\).
\end{proof}

\paragraph{Separate \(13\)-adic nondegeneracy lemma.}
Let \(U\) be rank one over \(K\otimes\mathbb Q_{13}\), with a
nondegenerate symmetric \(\mathbb Q_{13}\)-form \(q\) whose adjoint
on \(K\) is conjugation. Suppose \(U\) is realized by an algebraic
primitive projector, and let \(A\) be the \(\mathbb Q_{13}\)-span
of its algebraic classes. Then \(q|_A\) is nondegenerate. In particular
every numerically zero projected algebraic class is \(13\)-adically zero.

\begin{proof}
Since \(13\equiv-1\pmod7\), \(K\otimes\mathbb Q_{13}\) is the product
of three unramified quadratic fields, each preserved by conjugation.
For an explicit residue-field check,
\[
 \Phi_7(z)\equiv
 (z^2+6z+1)(z^2+5z+1)(z^2+3z+1)\pmod {13}.
\]
The discriminants \(6,8,5\) are nonsquares modulo \(13\).
The three field idempotents are self-adjoint and split \(U\)
orthogonally into three nondegenerate rank-one field modules.

Algebraic classes are stable under the algebraic \(K\)-action, so
their \(\mathbb Q_{13}\)-span is a \(K\otimes\mathbb Q_{13}\)-submodule.
In a rank-one module over a field a submodule is zero or the whole
module. Hence \(A\) is the sum of a subset of the three orthogonal
nondegenerate pieces and \(q|_A\) is nondegenerate.
A numerically zero algebraic class pairs to zero with all algebraic
classes in \(A\), because \(q\) is intersection with algebraic
polarization powers and projectors. It pairs to zero with their
\(\mathbb Q_{13}\)-span and therefore vanishes. No positivity or
Tate assertion is used in this lemma.
\end{proof}

\paragraph{Application to the fourfold and its square.}
\begin{proof}
C44-F09 supplies relative primitive projectors and Lefschetz operators,
hence \(B(X)\) and \(B(T)\) unconditionally. Write
\[
 H^*(X)=A\oplus P,\qquad
 A=\langle1,h,h^2,h^3,h^4\rangle,\quad P=P^4(X).
\]
The Lefschetz representation \(A\) is the length-five irreducible
string; \(P\) is trivial and concentrated in degree four. Every
primitive character occurs once with a free six-element rational
orbit. Each orbit gives a self-adjoint rank-six relative homological
motive, rank one over \(K\). This uses no Tate surjectivity.

For \(L=a h_1+b h_2\) on \(T\), with \(a,b>0\) rational, the
primitive decomposition is
\[
\begin{array}{c|l}
 \text{source}&\text{primitive summands}\\ \hline
 A\otimes A&\text{one Tate line in each degree }0,2,4,6,8\\
 P\otimes A&P\otimes1\text{ in degree }4\\
 A\otimes P&1\otimes P\text{ in degree }4\\
 P\otimes P&P\otimes P\text{ in degree }8 .
\end{array}
\]
The first row is Clebsch--Gordan for two length-five strings;
equivalently each ambient primitive kernel in degree \(2r\le8\)
has dimension one. In the second row \(h_1P=0\), so its length-five
string starts at \(P\otimes1\); the third row is symmetric. Both
operators vanish in the fourth row. These arguments work for all
\(a,b>0\). C44-F09 makes the projectors algebraic and relative at
the homological-motive level.

Split the non-Tate parts by rational characters of \(G\times G\).
Every nontrivial character in rows two and three occurs once.
Every character pair in row four occurs once with orbit length
six: a nonzero coordinate forces its scalar stabilizer to be trivial.
Thus all non-Tate primitive orbit motives have rank one over \(K\).
Distinct orbits are orthogonal because the cup product pairs a
character only with its negative, in the same rational orbit.
Ambient Tate classes are orthogonal to nontrivial characters.

For a rank-six primitive motive in degree \(2r\), take its Tate
twist \((r)\) and the relative algebraic form
\[
 q_r(x,y)=(-1)^r\int_TxyL^{8-2r}
\]
(and the analogous form on \(X\)). Diagonal automorphisms preserve
\(L\), so the adjoint on \(K\) is conjugation. Hodge--Riemann gives
a Betti polarization, and restriction to an orthogonal rational
sub-Hodge structure preserves this. In degree four on \(T\) the
sign is \(+1\): on \(P\otimes1\) the form is \(7b^4\int_Xuv\);
on \(1\otimes P\) it is \(7a^4\int_Xuv\). Indeed only the term
\(b^4h_2^4\), respectively \(a^4h_1^4\), survives, and
\(\int_Xh^4=7\). These are positive multiples of the generic
primitive degree-four polarization; the two pieces are orthogonal.

Let \(F\) be coordinate-squaring on \(X\) or \(T\), the base change
of Frobenius of its \(\mathbb F _2\)-model, considered over the chosen
residue field. For dimension \(N\),
\[
 F^*L=2L,\qquad \deg F=2^N,\qquad
 q_r(F^*x,F^*y)=2^{2r}q_r(x,y).
\]
The last equality follows by pulling back \(xyL^{N-2r}\) and using
the degree formula. Thus \(u=2^{-r}F^*\) is an algebraic isometry
on numerical classes. For any diagonal order-seven automorphism \(g\),
\(F\circ g=g^2\circ F\), hence
\[
 F^*g^*=(g^4)^*F^*.
\]
The orbit projector is invariant under multiplication by four.
Therefore \(F^*\) preserves each orbit motive and is
\(\sigma_4\)-semilinear on its \(K\)-action. The numerical sign
lemma proves positivity on every non-Tate primitive block.

The ambient primitive Tate lines have hyperplane-generated classes
and intersection numbers unchanged by specialization. Their generic
forms are positive by Hodge--Riemann, hence so are their special forms.
Orthogonality proves \(I_{\rm num}\). The algebraic Lefschetz
decomposition also holds modulo numerical equivalence, since
cohomological identities imply numerical identities. Thus the
primitive numerical space defined by the Lefschetz kernel is exactly
the sum of the spaces just treated, rather than a smaller chosen space.

At \(\ell=13\), apply the separate nondegeneracy lemma to every
rank-six primitive block; the ambient one-dimensional blocks also
have nondegenerate algebraic forms. Project a numerically zero cycle
by the algebraic Lefschetz operators to these primitive pieces,
with the appropriate twists. Each projected class is numerically
zero, hence \(13\)-adically zero. Reassembly proves \(D_{13}\)
in every codimension. Combining this with \(I_{\rm num}\) gives \(I_{13}\).

Finally weak Lefschetz gives \(H^1(X)=0\) and
\(\operatorname{NS}(X)_{\mathbb Q}=\mathbb Qh\). Kunneth gives
\(\operatorname{NS}(T)_{\mathbb Q}=\mathbb Qh_1\oplus\mathbb Qh_2\).
Restriction to a curve in either factor shows that every ample
class on \(T\) has positive coefficients \(a,b\); every ample class
on \(X\) is a positive multiple of \(h\). Thus all ample classes
are covered.

This proof does not transfer \(D_{13}\) to other primes or assert
the existence of classes in a Frobenius-invariant block whose
numerical algebraic space could be zero. The all-\(\ell\) result
for \(T\) remains conditional on 032/Milne-2026 as above;
the additional construction below proves more for \(X\).
\end{proof}

\section{C44-F13: explicit ruled surfaces make the extra classes unconditional}
\label{sec:C44-F13}
\paragraph{Statement.}
The special algebraicity assertion of C44-F11 is unconditional.
The ruled surface
\[
 Z=\overline{\left\{
 [\lambda s:\lambda t:\mu w^2:\mu s^2:\mu t^2:\lambda w]:
 s^7+t^7+w^7=0
 \right\}}\subset X
\]
has nonzero projection to every character line in \(e_\alpha H^4(X)\).
If \(g_0\) scales \(x_0\) by \(\zeta _7\), the six rational cycles
\[
 z_j=e_\alpha(g_0^j)^*[Z],\qquad 0\le j\le5,
\]
form a \(\mathbb Q_\ell\)-basis of this block for every \(\ell\ne2\)
and are defined over \(\mathbb F _8\). Their products and the
projectors \(W_r\) give thirty independent algebraic directions in
\(\bigoplus_{r=1}^5W_r\subset P^8(T)\), none of whose nonzero classes
can specialize from generic algebraic classes in the fixed Fermat
lifting, even after finite base extension. This uses neither
032/Milne-2026 nor a Tate conjecture.

\begin{proof}
Let \(C:s^7+t^7+w^7=0\subset\mathbb P^2\). Partition the coordinates
of \(X\) into \((x_0,x_1,x_5)\) and \((x_2,x_3,x_4)\). The character
\(\alpha\) becomes the concatenation
\[
 u=(1,1,5),\qquad v=(1,3,3).
\]
Both sum to zero modulo seven and occur in \(H^1(C)\).

Over the relative Fermat base of C44-F08 form
\[
 Y=\mathbb P_{C\times C}
 \bigl(\operatorname{pr}_1^*\mathcal O_C(-1)
       \oplus\operatorname{pr}_2^*\mathcal O_C(-1)\bigr),
 \qquad p:Y\longrightarrow C\times C.
\]
The tautological lines lie in complementary coordinate spaces
and define a morphism \(j:Y\to X\), the projective join. The
equation of \(X\) holds because it holds separately on each \(C\).
Thus
\[
 J=j_*p^*:H^1(C)\otimes H^1(C)(-1)\longrightarrow H^4(X)
\]
is an algebraic relative correspondence and sends divisor classes
on \(C\times C\) to codimension-two cycle classes on \(X\).

In characteristic zero this is Ran's operation in Section 4,
printed pp.~138--140. Proposition 4.6 and Corollary 4.7 on p.~140
prove injectivity on primitive cohomology and identify character
pairs with their concatenation:
\url{https://www.numdam.org/article/CM_1980__42_1_121_0.pdf}.
Ran works over \(\mathbb C\); the above correspondence, however,
is defined between smooth projective families over the trait.
It induces a map of lisse \(\mathbb Q_\ell\)-cohomology sheaves.
Injectivity and nonvanishing on each character line therefore
persist to the special fibre by smooth proper base change.
This holds for every \(\ell\ne2\), including \(\ell=7\), since the
coefficient field has characteristic zero. Diagonal equivariance
is compatible with concatenation: the characters under consideration
sum to zero on each coordinate block and so respect the two
blockwise diagonal quotients.

On the special curve define
\[
 \sigma(s,t,w)=(w,s,t),\qquad F(s,t,w)=(s^2,t^2,w^2),
 \qquad f=\sigma\circ F.
\]
Here \(F\) is the base change of the Frobenius of the
\(\mathbb F _2\)-model, viewed as a morphism over \(\mathbb F _8\).
Pullback by \(\sigma\) sends \((a,b,c)\) to \((b,c,a)\),
while \(F^*\) sends \(a\) to \(2a\). Therefore
\[
 f^*H_v=H_{2(3,3,1)}=H_{(6,6,2)}=H_{-u}.
\]
For every \(t=1,\ldots,6\), \(f^*:H_{tv}\to H_{-tu}\) is an
isomorphism: \(\sigma\) is an isomorphism and Frobenius pullback
is invertible on \(\ell\)-adic cohomology.

Let \(\Gamma_f\subset C\times C\) have first coordinate the source.
For \(\eta,\theta\in H^1(C)\),
\[
 \int_{C\times C}[\Gamma_f](\eta\otimes\theta)
   =\int_C\eta f^*\theta.
\]
Choose nonzero \(\eta\in H_{-tu}\), \(\theta\in H_{-tv}\).
Then \(f^*\theta\) is nonzero in \(H_{tu}\) and pairs nontrivially
with \(\eta\), since opposite one-dimensional character spaces
have a nondegenerate Poincare pairing. Hence the
\(H_{tu}\otimes H_{tv}\) component of \([\Gamma_f]\) is nonzero
for every \(t\). The join is nonzero on each such line, so
\(J[\Gamma_f]\) has nonzero component in every \(H_{t\alpha}\).
The two other degree-two Kunneth summands of \(C\times C\)
have a trivial character on one factor; equivariance prevents
them from contributing to or cancelling these components.

The inverse image of \(\Gamma_f\) is
\(\mathbb P_C(\mathcal O_C(-1)\oplus\mathcal O_C(-2))\).
Its map to \(X\) is the displayed parameterization of \(Z\),
homogeneous with these two line-bundle weights. It lies in \(X\)
since, in characteristic two,
\[
 \lambda^7(s^7+t^7+w^7)
 +\mu^7(w^{14}+s^{14}+t^{14})=0.
\]
On \(\lambda\ne0\), the first block recovers \([s:t:w]\) and
the second recovers the fibre coordinate. The map is birational
onto \(Z\), so \(J[\Gamma_f]=[Z]\) with multiplicity one.

The projector is explicitly rational:
\[
 e_\alpha=
 \frac1{|\ker\chi_\alpha|}\sum_{g\in\ker\chi_\alpha}g^*
 -\frac1{|G|}\sum_{g\in G}g^*,
 \qquad |G|=7^5,\quad|\ker\chi_\alpha|=7^4.
\]
It is the sum of the six nontrivial characters of the quotient
\(G/\ker\chi_\alpha\), hence precisely the C44-F11 orbit projector.
Every \(z_j\) is an explicit rational combination of diagonal
translates of \(Z\), defined over \(\mathbb F _8\).
Over a coefficient extension containing \(\zeta _7\), write
\[
 e_\alpha[Z]=\sum_{t=1}^6c_t,\qquad
 0\ne c_t\in H_{t\alpha}.
\]
Since \(\alpha_0=1\), \(z_j=\sum_t\zeta _7^{jt}c_t\).
The six-by-six matrix \((\zeta _7^{jt})_{0\le j\le5,\ 1\le t\le6}\)
is Vandermonde with distinct nodes and is invertible.
Thus the \(z_j\) form a basis after coefficient extension,
and hence over \(\mathbb Q_\ell\), for every \(\ell\ne2\).

Their thirty-six exterior products form a basis of
\(e_\alpha H^4(X)\otimes e_\alpha H^4(X)\).
Applying the relative rational projectors \(W_r\) spans each
of the six blocks and in particular the thirty-dimensional
sum with \(r\ne6\). The generic Hodge obstruction and the
specialization argument of C44-F11 apply unchanged to these
five blocks, proving the claimed nonliftability in this
Fermat model after every finite base extension.
\end{proof}

\section{C44-F14: all Tate classes of the degree-seven fourfold}
\label{sec:C44-F14}
\paragraph{Statement.}
For \(X/\mathbb F _2\), the Tate conjecture holds unconditionally
in every codimension and over every finite extension of
\(\mathbb F _2\). The codimension-two geometric Tate space has
dimension \(2221\): the ambient line and \(2220\) primitive
character lines. Of those primitive lines, \(1860\) come from
linear planes and \(360\) from diagonal translates and coordinate
permutations of the C44-F13 ruled surface.
Consequently \(D_\ell(X)\) and \(I_\ell(X)\) hold unconditionally
for every \(\ell\ne2\) and every ample class.
The \(360\) exceptional directions have zero generic rational
Hodge space in the fixed Fermat lifting and cannot specialize
from generic algebraic classes. This does not establish all-\(\ell\)
\(D\) or \(I\) for \(X\times X\).

\begin{proof}
Use \(G_H=-1+\delta\), \(G_{-H}=-1-\delta\),
\(H=\{1,2,4\}\), and \(\delta^2=-7\) from C44-F11.
For a primitive character \(a\), let \(h(a)\) count its entries in
\(H\). The Jacobi-sum formula gives
\[
 j(a)=\frac{G_H^{h(a)}G_{-H}^{6-h(a)}}8
     =64\,\rho^{h(a)-3},\qquad \rho=G_H/G_{-H}.
\]
The polynomial of \(G_H\) is \(z^2+2z+8\), with \(2\)-adic
Newton valuations \(1,2\). Therefore \(\rho\) has valuation
\(1\) or \(-1\) at a prime above \(2\) and is not a root of
unity. A primitive character is geometric Tate exactly when
\(h(a)=3\), and then its Frobenius-\(8\) eigenvalue is \(64\).
All geometric Tate classes are thus fixed over \(\mathbb F _8\)
after the Tate twist.

Here is a hand classification. Let \(a_1,a_2,a_4\) count the
entries \(1,2,4\), and let \(b_1,b_2,b_4\) count the opposite
entries \(6,5,3\). Both sums equal three. Put \(d_i=a_i-b_i\).
The character-sum condition is
\[
 d_1+d_2+d_4=0,\qquad d_2+3d_4\equiv0\pmod7,\qquad
 -3\le d_i\le3.
\]
For \(d_4=0\), these force all \(d_i=0\).
For \(d_4=1,2,3\), the range of \(d_2\) forces, respectively,
\[
 (d_1,d_2,d_4)=(2,-3,1),\quad(-3,1,2),\quad(-1,-2,3).
\]
Negative \(d_4\) gives their negatives. Each nonzero vector
has positive and negative parts of total three, so it uniquely
determines all the counts \(a_i,b_i\). The zero vector gives
the ten common count vectors \(a_i=b_i\) of sum three.

The ten zero-difference cases are precisely three opposite pairs.
The six exceptional sorted tuples are
\[
\begin{gathered}
 (1,1,1,3,3,5),\quad(1,1,4,5,5,5),\quad(1,2,2,3,3,3),\\
 (2,2,2,3,6,6),\quad(2,4,4,6,6,6),\quad(4,4,4,5,5,6).
\end{gathered}
\]
They are exactly the sorted scalar multiples of \(\alpha\).
Every exceptional character is therefore a coordinate permutation
of a scalar multiple of \(\alpha\). C44-F13, coordinate permutations,
rational orbit projectors, and diagonal translates algebraically
span the entire exceptional subspace.

The opposite-pair characters are spanned by linear-plane classes.
Indeed iterate the projective join of three copies of the
zero-dimensional degree-seven Fermat variety in \(\mathbb P^1\).
Its primitive degree-zero characters are \((b,-b)\), and point
classes span its degree-zero cohomology. Joining three chosen
points gives a plane in \(X\). Ran Proposition 4.6 and Corollary 4.7,
spread out as in C44-F13, show that these primitive join maps
are nonzero on every concatenated character line. Thus linear
planes span all opposite-pair characters in every good fibre.
This uses an actual join correspondence, not a conjectural
specialization-surjectivity assertion.

For the ordered count, three paired patterns have multiplicities
\((3,3)\), six have \((2,2,1,1)\), and one has six distinct entries.
Their contribution is
\[
 3\frac{6!}{3!3!}+6\frac{6!}{2!2!}+6!
 =60+1080+720=1860.
\]
Each exceptional pattern has multiplicities \((3,2,1)\);
the six patterns contribute \(6\cdot6!/(3!2!)=360\).
Thus the primitive Tate space has dimension \(2220\),
and adding \(h^2\) gives \(2221\).

All constructed cycles are defined over \(\mathbb F _8\) and
span the geometric Tate space after every \(\ell\)-adic realization.
Over another finite field \(\mathbb F _{2^m}\), pass to the common
extension containing \(\mathbb F _8\) and average the cycles over
its Galois group. With rational coefficients the normalized norm
projects onto the relevant Frobenius invariants. This proves
codimension-two Tate over every finite extension. The other
codimensions have only ambient one-dimensional cohomology,
spanned by powers of the hyperplane class.

The primitive geometric Tate space is nondegenerate for cup
product: it consists exactly of \(h(a)=3\) character lines,
is closed under \(a\mapsto-a\), and each line pairs perfectly
with its opposite. Since algebraic classes span it, every
numerically zero cycle has zero class for every \(\ell\ne2\).
The ambient pieces are nondegenerate as well. Thus \(D_\ell(X)\)
holds in every codimension, and the unconditional numerical
positivity of C44-F12 gives \(I_\ell(X)\) for every \(\ell\ne2\).
All ample classes are positive multiples of \(h\), as in C44-F12.

Finally a rational generic Hodge character must satisfy
\(\sum_i\langle ta_i\rangle=21\) for every unit \(t\) modulo seven.
Summing over \(t\in H\) gives
\[
 63=7h(a)+14(6-h(a)),
\]
so \(h(a)=3\). Our classification therefore includes every
possible generic rational Hodge character. The opposite-pair
characters have type \((2,2)\) throughout their rational orbit.
Every exceptional character is a coordinate permutation and
scalar multiple of \(\alpha\), whose orbit has mixed Hodge
types by C44-F11 and no nonzero rational Hodge class.
Thus the entire \(360\)-dimensional exceptional special subspace
has zero generic algebraic-class specialization image in
the fixed Fermat lifting.
\end{proof}

% END work/arithmetic/fermat7_frobenius.tex

% BEGIN work/scan_w/contribution.tex
% C44_POS /root/scan_w.  New deductions, 2026-11-22 Europe/Berlin.
% The model is GPT-6 Astra Pro.  Runtime reasoning-effort setting is not exposed.
% This file is a paper fragment.  The root paper owns theorem environments.

\section{An algebraic replacement for Wirtinger in the Weil-plane detector argument}

The statements in this section distinguish an unconditional intersection
estimate from the positivity and algebraicity inputs.  Every invocation of
032, Corollaries 8.3 or 8.4, is conditional on the unrefereed preprints
specified by the user.  The characteristic-zero Theorem W in
\texttt{HODGE\_DIVISION\_ATTACK\_20261121.md}, R4, is not being reproved.
We give its missing positive-characteristic degree estimate and then record
the resulting, slightly differently normalized, detector theorem.

\subsection*{C44-W01. A characteristic-free intersection bound on an abelian variety}

\begin{proposition}[C44-W01]
Let $A$ be an abelian variety of dimension $2n$ over an algebraically
closed field, and let $h=c_1(L)$ be an integral ample class.  Put
\[
 D(Z)=\deg(h^n\cdot Z),\qquad
 V=\frac{\deg h^{2n}}{(2n)!}=\chi(L)\in\mathbf Z_{>0},\qquad
 b_h=\frac{2^n}{n!\sqrt V}.
\]
For effective integral cycles $Z,Z'$ of dimension $n$, including cycles
with multiplicities,
\begin{equation}\label{c44:algebraic-bezout}
 0\leq\deg(Z\cdot Z')\leq b_h^2D(Z)D(Z').
\end{equation}
No standard conjecture and no lifting hypothesis is used.
\end{proposition}

\begin{proof}
It suffices by bilinearity to treat integral subvarieties.  On
$Z\times Z'$ consider $f(z,z')=z-z'$ and $s(z,z')=z+z'$.  Translation
moving, or equivalently refined pullback of a general point of $A$ by
$f$, gives
\[
 f_*[Z\times Z']=e[A],\qquad e=\deg(Z\cdot Z')\geq0.
\]
If $f$ is not dominant, both sides have coefficient $e=0$.  If it is
dominant, it is generically finite and $e$ is its degree, including any
inseparable multiplicity.  To see explicitly that the intersection
coefficient agrees with this degree, a general translate of $Z'$ meets
$Z$ properly and avoids the singular loci of both subvarieties: the
bad incidence loci have dimension at most $2n-1$.  At the resulting
zero-dimensional intersections the two smooth subvarieties are local
complete intersections in the smooth variety $A$, and properness gives
the usual length intersection multiplicity.  The same description
applies to the generic fibre of $f$.  Translation preserves cycle
classes.  This also proves the assertion without separability or
transversality assumptions.

Let $h_1,h_2$ be the pullbacks of $h$ to $Z\times Z'$, and set
$\alpha=f^*h$, $\beta=s^*h$.  Both are nef.  The theorem of the cube,
in its parallelogram consequence, gives the numerical identity
\[
 \alpha+\beta=2(h_1+h_2).
\]
Here $[-1]^*h=h$ numerically; no symmetric representative of the
original line bundle is required.  A convenient exact source is
Milne, \emph{Abelian Varieties}, Corollary 6.5, applied to the maps
$z,z',-z'$; the numerical consequence is the displayed identity.
All mixed intersections of nef divisors with an effective cycle are
nonnegative, by approximation by ample rational classes.  Therefore
\begin{align*}
 e\deg h^{2n}
 &=\deg(\alpha^{2n}\cap[Z\times Z'])\\
 &\leq\deg((\alpha+\beta)^{2n}\cap[Z\times Z'])\\
 &=2^{2n}\binom{2n}{n}D(Z)D(Z').
\end{align*}
The last equality uses $\dim Z=\dim Z'=n$.  Division by
$\deg h^{2n}=(2n)!V$ proves the bound.  The asserted formula and
integrality for $V$ are the abelian Riemann--Roch theorem, Milne,
Theorem 13.3(b); ampleness gives $V>0$.
\end{proof}

\paragraph{Realizations and historical scope.}
Equation \eqref{c44:algebraic-bezout} is an inequality between rational
intersection numbers.  It is valid for classes represented in either
$\ell$-adic or crystalline cohomology because the cycle map preserves
these products and trace.  It does not put an order on $\mathbf Q_\ell$.
This is precisely the kind of degree bound needed below.  In
Grothendieck's framework its operations are proper pushforward, refined
intersection, and the biextension identity; SGA 7 I, Exp. VIII,
4.12.1 and 2.3.11 provide the biextension convention already used by
JAC-T4.3.  It supplies a new quantitative input to the positive
incidence calculation of cote 44, scans 008--009, and to Theorem W,
step 3.  It makes no assertion that a new axis class lifts.

\subsection*{C44-W02. The signed Weil-plane Gram matrix in any characteristic}

Fix an imaginary quadratic field $K$ and an order
$\mathcal O\subseteq K\cap\operatorname{End}(A)$ such that the Rosati
involution of $h$ induces conjugation on $K$.  Choose
$a\in\mathcal O$ with $a/\bar a$ not a root of unity.  Write
\[
 \nu=a\bar a\in\mathbf Z_{>0},\quad q=\nu^n,\quad
 s=a^{2n}+\bar a^{2n}\in\mathbf Z,\quad H=4q^2-s^2>0.
\]
For an invertible prime $\ell$, define
\[
 W_\ell=\bigwedge_{K\otimes\mathbf Q_\ell}^{2n}H^1(A,\mathbf Q_\ell)
 \subset H^{2n}(A,\mathbf Q_\ell).
\]
It has dimension two over $\mathbf Q_\ell$.  Let
\[
 f(t)=(t-a^{2n})(t-\bar a^{2n}),\quad
 g(t)=\prod_{j=1}^{2n-1}(t-a^j\bar a^{2n-j}).
\]
Both polynomials have integer coefficients.  Choose $R>0$ and
$u,v\in\mathbf Z[t]$ with $uf+vg=R$, for example the absolute
resultant with the sign of the Bezout relation adjusted.  Put
\[
 p(t)=v(t)g(t)=\sum_jc_jt^j,\qquad
 \kappa=\sum_j|c_j|\nu^{jn},\qquad e_W=R^{-1}p([a]^*).
\]
These are exactly the projector data of the supplied Theorem W.
We use them, with the same normalization, on the new realization.

\begin{proposition}[C44-W02; conditional on unrefereed 032, Corollary 8.4]
Let $Y\subset A$ be an integral $n$-dimensional subvariety with
$e_W[Y]\ne0$, and let $D=D(Y)$.  Then
\[
 y_1=p([a]^*)[Y],\qquad y_2=[a]^*y_1
\]
are independent integral algebraic cycle classes in $W_\ell$.
The form
\[
 Q(x,y)=(-1)^n\deg(xy)
\]
is positive definite on their rational span.  With $d=Q(y_1,y_1)$,
its Gram matrix is exactly
\begin{equation}\label{c44:weil-gram}
 G=d\begin{pmatrix}1&s/2\\s/2&q^2\end{pmatrix},\qquad
 \det G=d^2H/4\in\mathbf Z_{>0},\qquad
 0<d\leq b_h^2R\kappa D^2\leq b_h^2\kappa^2D^2.
\end{equation}
Moreover
\[
 W_{\rm alg}:=e_W\operatorname{Alg}^n(A)
    =\mathbf Qy_1\oplus\mathbf Qy_2,
 \qquad W_{\rm alg}\otimes\mathbf Q_\ell=W_\ell.
\]
The two classes have integral signed-cycle representatives whose
positive and negative parts have degrees at most $\kappa D$ and
$q\kappa D$, respectively.
\end{proposition}

\begin{proof}
After extending the coefficient field, write $H^1=V_a\oplus V_{\bar a}$.
Both eigenspaces have dimension $2n$, and the Rosati condition says
that $h$ lies in $V_a\otimes V_{\bar a}$.  Each of the two extreme
exterior powers is therefore annihilated by $h$.  Thus $W_\ell$ is
middle primitive.  The eigenspaces of $[a]^*$ in middle cohomology
have eigenvalues $a^j\bar a^{2n-j}$; all are distinct by the choice
of $a$.  Hence the integral polynomial $p$ is $R$ on the extreme
two lines and zero on the other weight spaces.  This verifies the
projector on this realization, rather than importing a Betti lattice.
The exterior-algebra description is valid integrally for etale
cohomology by Milne, Theorem 15.1.

Every $[a^j]$ is a finite flat isogeny, so the pullbacks in
$p([a]^*)[Y]$ are effective integral cycles with their flat
multiplicities.  Rosati gives $[a^j]^*h=\nu^jh$ numerically and
$\deg[a^j]=\nu^{2nj}$.  Projection formula consequently gives
\[
 D([a^j]^*Y)=\nu^{jn}D.
\]
This proves both asserted degree budgets.  The minimal polynomial
of $[a]^*$ on $W_\ell$ is the irreducible quadratic $f$: indeed
$a^{2n}$ is nonreal.  Hence $y_1,y_2$ are rationally independent.

The signed Hodge standard conjecture, used with exactly the sign
$(-1)^n$, is positive definite on every rational subspace of
primitive algebraic classes.  This is the sole input for positivity
here.  It is the even-degree version of the primitive sign in
cote 44, scan 008; JAC-T4.3 establishes the corresponding incidence
polarization by a curve test.  The unsigned form printed in
Milne's report, Section 3.6, is not used.

Let $T=[a]^*|W_\ell$.  Projection formula gives
$Q(Tx,Ty)=q^2Q(x,y)$.  Thus $T^\dagger=q^2T^{-1}$ and
$T+T^\dagger=s$.  It follows that
$Q(y_1,y_2)=sd/2$ and $Q(y_2,y_2)=q^2d$, proving the matrix
formula.  Its entries are integers, since the $y_i$ are integral
cycle classes, and its determinant is a positive integer.
There is a useful improvement over expanding every pair of effective
pullbacks.  On all of $H^{2n}(A)$, the intersection adjoint of
$[a]^*$ is $q^2[a]^{*-1}$.  The pairing therefore matches the
weight indexed by $j$ with the weight indexed by $2n-j$.
The projector $e_W$ selects both extreme weights and annihilates
their orthogonal complement, so it is self-adjoint.  Consequently
$p([a]^*)^\dagger p([a]^*)=R p([a]^*)$ on middle cohomology.
Apply C44-W01 to each term of $[Y]\,p([a]^*)[Y]$ to obtain
\[
 d=|\deg(y_1^2)|=R\left|\deg\bigl([Y]p([a]^*)[Y]\bigr)\right|
 \leq b_h^2R\kappa D^2.
\]
Finally $p(a^{2n})=R$ and $|a^{2n}|=q$ give $R\leq\kappa$
by the ordinary triangle inequality on this explicit algebraic number.
This numerical inequality on projector coefficients is not an
archimedean realization of a cycle class.
This proves the quantitative replacement for Wirtinger.

Finally any finite rationally independent collection of classes in
$W_{\rm alg}$ has a positive definite rational Gram matrix by the
same Hodge-standard input.  Its determinant is nonzero, so those
classes are also $\mathbf Q_\ell$-independent.  There can be at most
two.  The two classes already constructed therefore span
$W_{\rm alg}$ and its realization.
\end{proof}

\paragraph{What ``integral'' means in positive characteristic.}
There is no canonical integral Betti lattice over $\overline{\mathbf F}_p$.
Let $\Lambda$ denote the group of integral algebraic middle classes
modulo numerical equivalence, viewed in cohomology using 032's
numerical--homological equality.  It is a lattice: choose finitely
many integral test cycles giving a nondegenerate rational pairing;
the intersection map embeds $\Lambda$ in a finite-rank free abelian
group.  Write
\[
 \Lambda^\sharp=\{x\in\operatorname{Alg}^n(A):\deg(x\lambda)\in\mathbf Z
                       \text{ for all }\lambda\in\Lambda\}.
\]
The denominator condition needed below is $\delta t\in\Lambda^\sharp$.
The stronger condition $\delta t\in\Lambda$ suffices.  If an integral
compatible realization is specified, its rational classes with integral
pairings also satisfy this condition.  Integrality only at one prime
$\ell$ is insufficient to make the Cramer numerators integers.

\subsection*{C44-W03. A quadratic detector bound in characteristic $p$}

\begin{theorem}[C44-W03; conditional on unrefereed 032, Corollary 8.4]
Retain the preceding notation and suppose $W_{\rm alg}\ne0$.
For $0\ne t\in W_{\rm alg}$ let $\delta\geq1$ be an integer with
$\delta t\in\Lambda^\sharp$, and put $T_0=Q(t,t)>0$.
Explicitly, the detector degree and division complexity are
\begin{align*}
 \delta_W&=\min\{D(Y):Y\subset A\text{ integral of dimension }n,
                                      \ e_W[Y]\ne0\},\\
 c_t&=\min\{\max(N,D(Z^+),D(Z^-)):
       N\in\mathbf Z_{\geq1},\ Z^\pm\text{ effective integral cycles},
                              \ Nt=[Z^+]-[Z^-]\}.
\end{align*}
These are the definitions in the supplied Theorem W; an empty minimum
is interpreted as infinity.  Cycle equalities here are equalities of
classes, and an integral cycle may have several components with
nonnegative integer multiplicities.
Set $B=b_h\sqrt{R\kappa}$.  Then
\begin{equation}\label{c44:weil-quadratic}
 \delta_W\leq c_t\leq
 \delta\max\{HB^2,\,4q\kappa B\sqrt{HT_0}\}\,\delta_W^2.
\end{equation}
In particular one can replace $B$ everywhere by the uniform constant
$B_0=2^n\sqrt{R\kappa}/n!$, independent of the abelian variety, the
characteristic, and the polarization type.  The quartic version
\begin{equation}\label{c44:weil-quartic}
 c_t\leq\delta q^2\max\{B^4,\,4\kappa B^3\sqrt{T_0}\}\,\delta_W^4
\end{equation}
also holds.  Thus the characteristic-$p$ division theorem has both
the requested detector formulation and a quadratic sharpening,
with fully algebraic, explicit constants.
\end{theorem}

\begin{proof}
The hypothesis $W_{\rm alg}\ne0$ supplies a detector: express a
nonzero class as a rational combination of integral effective
cycles and apply $e_W$.  Choose a detector $Y$ of minimal degree
$D=\delta_W$, and form $y_1,y_2$ as in C44-W02.  Put
$b_i=Q(\delta t,y_i)\in\mathbf Z$.  Inverting the exact Gram matrix
\eqref{c44:weil-gram} gives the identity
\begin{equation}\label{c44:exact-clearing}
 \delta Hd\,t=(4q^2b_1-2sb_2)y_1+(4b_2-2sb_1)y_2.
\end{equation}
All coefficients and the positive multiplier $N=\delta Hd$ are
integers.  If $t=t_1y_1+t_2y_2$, Cauchy--Schwarz in this
positive definite rational plane, or the diagonal entries of
$G^{-1}$, gives
\[
 |t_1|\leq\frac{2q\sqrt{T_0}}{\sqrt{dH}},\qquad
 |t_2|\leq\frac{2\sqrt{T_0}}{\sqrt{dH}}.
\]
Therefore the two integral coefficients on the right of
\eqref{c44:exact-clearing} have absolute values at most
$2\delta q\sqrt{HT_0d}$ and $2\delta\sqrt{HT_0d}$.
Using the degree budgets of C44-W02, the positive and negative
parts of the resulting integral cycle each have degree at most
\[
 \kappa D\bigl(2\delta q\sqrt{HT_0d}
                     +2\delta q\sqrt{HT_0d}\bigr)
 \leq4\delta q\kappa B\sqrt{HT_0}\,D^2.
\]
Also $N\leq\delta HB^2D^2$.  These two estimates prove the
upper bound in \eqref{c44:weil-quadratic}.

For completeness the ordinary two-by-two adjugate clearing with
multiplier $\delta\det G$ gives
$N\leq\delta q^2B^4D^4$.  Its two coefficients are bounded by
$2\delta\sqrt{T_0}\|y_1\|\|y_2\|^2$ and
$2\delta\sqrt{T_0}\|y_1\|^2\|y_2\|$, where
$\|y_1\|=\sqrt d\leq BD$ and $\|y_2\|=q\sqrt d$.
Their degree budget is at most
$4\delta q^2\kappa B^3\sqrt{T_0}D^4$, proving
\eqref{c44:weil-quartic}.

Conversely any class equation
$Nt=[Z^+]-[Z^-]$, with $N\geq1$ and effective integral cycles,
has a nonzero Weil projection on at least one integral component
of $Z^+$ or $Z^-$.  That component has degree at most the
corresponding cycle degree.  Hence $\delta_W\leq c_t$.
\end{proof}

\paragraph{Exact change to the supplied Theorem W.}
The original complex constant $|R|/n!$ in the norm estimate is
replaced by $2^n\sqrt{R\kappa}/(n!\sqrt V)$ for the integral numerator
$y_1$.  We do not claim that Wirtinger's sharper constant survives
unchanged.  All other projector, degree, independence, and integer
clearing steps are now valid without an archimedean realization.
The bound is on cohomological or numerical cycle classes, exactly
as in the supplied definition of $c_t$; it is not an equality in
the Chow group modulo rational equivalence.

\subsection*{C44-W04. The precise balanced-signature hypothesis over a finite field}

\begin{proposition}[C44-W04; Tate-algebraicity conclusion conditional on unrefereed 032]
Let $(A,h,K)$ descend to a finite field $\mathbf F_Q$, with the
$K$-action and polarization defined there, and $\dim A=2n$.
If $p$ splits in $K$, assume the two rational $K\otimes\mathbf Q_p$
components of $A[p^\infty]$ both have dimension $n$ (each has
height $2n$).  If $p$ is inert or ramified, no additional determinant
balance is necessary.  Then Frobenius acts with finite order on
$W_\ell(n)$.  Conditional on 032, Corollary 8.3, the whole
$W_\ell$ is algebraic, $W_{\rm alg}$ is its two-dimensional
rational cycle structure, $\delta_W<\infty$, and C44-W03 applies
to every nonzero class in that rational Weil plane.

When $p$ is split and $p$ does not divide the conductor of
$\mathcal O$, the stated $p$-divisible balance is exactly the usual
$(n,n)$ multiplicity condition on $\operatorname{Lie}A$.
It also holds for the special fibre of a balanced polarized
$\mathcal O_K\otimes\mathbf Z_p$-linear abelian scheme in mixed
characteristic.  These hypotheses give the usual meaningful
all-characteristic version of the signature condition.
\end{proposition}

\begin{proof}
Write $\pi$ for the $Q$-power Frobenius endomorphism.  Its
$K$-linear determinant on $H^1$ is an element $\alpha\in K$,
independent of the realization.  One can see its rationality
without choosing a complex lift: the traces of $b\pi^j$ for
$b\in K$ are rational, by the characteristic polynomial theorem
for endomorphisms, and the nondegenerate trace pairing of $K$
recovers the $K$-linear traces.  Newton identities recover the
$K$-linear characteristic polynomial and its determinant.
The same argument in crystalline cohomology gives the same
determinant: in any of these realizations, the exterior-algebra
description of abelian cohomology and the top-degree action give
$\det(m-[b])=\deg([m]-[b])$ for every integer $m$ for which the
map is an isogeny.  Polynomial identity then makes the
characteristic polynomial independent of the realization.
The trace/characteristic-polynomial input is
Milne, \emph{Abelian Varieties}, Proposition 12.9; it is the
degree-defined characteristic polynomial of an endomorphism.

The determinant $\alpha$ is an algebraic integer, since its
eigenvalues are products of Frobenius eigenvalues.  The
polarization pairing gives
\[
 \alpha\bar\alpha=Q^{2n},\qquad |\alpha|=Q^n
\]
at the two complex embeddings.  The absolute-value assertion is
Weil's theorem for abelian varieties, Milne, Theorem 19.1(c).
In particular $\alpha$ is a unit away from primes over $p$.
If $p$ is nonsplit, its unique prime is fixed by conjugation, so
the norm identity implies
$v(\alpha)=n\,v(Q)$ there.  If $p$ splits, the determinant of
Frobenius on each rational Dieudonne summand has normalized
determinant valuation equal to that summand's dimension:
the Newton polygon of a $p$-divisible group of height $h_0$
and dimension $d_0$ ends at $(h_0,d_0)$, so its slopes sum
to $d_0$.  For exact source statements see Oort,
\emph{Abelian Varieties over Finite Fields}, Sections
21.16--21.19 and Corollary 9.3.  Equivalently each
Dieudonne--Manin factor of dimension $d_0$ and height
$d_0+c_0$ has $d_0+c_0$ slopes $d_0/(d_0+c_0)$.
The balance hypothesis therefore gives
$v(\alpha)=n\,v(Q)$ at both primes.

Consequently $\alpha/Q^n$ is a unit in the imaginary quadratic
field $K$ of absolute value one, hence a root of unity.  On the
two Weil lines the Frobenius eigenvalues are $\alpha$ and
$\bar\alpha$.  A finite extension of the finite field kills
their normalized root-of-unity factors.  The Tate conjecture
in the stated 032 input then gives algebraic classes spanning
$W_\ell$: apply the algebraic projector $e_W$ to cycles spanning
the Tate classes.  C44-W02 identifies their rational cycle
span as a plane.  This proves the assertions.

For split $p$ and a $p$-maximal order, the two integral
idempotents act on the $p$-divisible group and on its Lie
algebra, so the dimensions agree with the two Lie
multiplicities.  In an integral balanced family these ranks
are constant.  Isogenies preserve the dimensions and heights
of the rational $p$-divisible summands.
\end{proof}

\begin{proposition}[C44-W05: a bad-order ambiguity in the word signature]
A determinant identity on $\operatorname{Lie}A$ for an order whose
conductor is divisible by a split prime $p$ does not by itself
imply the balance hypothesis of C44-W04.  In particular the
literal naive such identity cannot replace that hypothesis.
\end{proposition}

\begin{proof}
For a completely concrete example take
$E:y^2=x^3-x$ over $\mathbf F_5$, $K=\mathbf Q(i)$.
It has eight rational points, hence Frobenius trace $-2$;
in particular it is ordinary.  Its order-four automorphism
gives the $\mathbf Z[i]$-action, and its Frobenius is one
of $-1\pm2i$.  Thus $p=5$ splits in $K$, as required.
The following construction also works with any ordinary
CM elliptic curve at a split prime.  After a finite
extension all data are defined over $\mathbf F_Q$.
Let $A=E^{2n}$ with the scalar $K$-action and product polarization,
and restrict the action to $\mathcal O=\mathbf Z+p\mathcal O_K$
(or to $\mathbf Z+p\operatorname{End}(E)$ if necessary).
For $a=m+pb$ the action on the Lie algebra is multiplication by
$m\bmod p$, and the reductions of $a$ and $\bar a$ agree.
Thus
\[
 \det(T-a\mid\operatorname{Lie}A)
 =(T-a)^n(T-\bar a)^n
\]
holds after reduction for every $a\in\mathcal O$, although the
two rational $p$-divisible components have dimensions $2n$ and
$0$.  If $\pi$ is the ordinary Frobenius of $E$, the normalized
Weil determinant is $\pi^{2n}/Q^n$; its valuations at the two
primes over $p$ are nonzero and opposite.  It is not a root of
unity, and the two Weil lines are not Tate.  The example
therefore proves the claimed insufficiency.  It is not a
counterexample to C44-W03 or to the usual balanced PEL meaning
of Weil type.
\end{proof}

\section{Geometric positive transfers and special complete intersections}

The abstract abelian-type transfer remains distinct from a transfer that
preserves the actual incidence form. The following statements give
such geometric transfers and identify the extra special-fibre directions
that they do not reach. Their positivity conclusions are conditional
on the unrefereed 032 input whenever that qualification is displayed.

\subsection*{C44-W06. A finite abelian cover gives an explicit inherited polarization}

\begin{theorem}[C44-W06; conditional on unrefereed 032, Corollary 8.4]
Let $X$ be smooth projective of odd dimension $d=2m-1$ over an
algebraically closed field. Suppose that there is a finite surjective
morphism $f:A\to X$ from an abelian variety. For every rational ample
class $h$ on $X$, put $h_A=f^*h$ and $e=\deg f$. The primitive
incidence quotient $B_X=P^m(X)_{\rm prim}$ has the signed rational
polarization of cote 44, Th.~4. More precisely the induced map
\[
 i_f:B_X\longrightarrow B_A=P^m(A)_{\rm prim}
\]
has finite kernel, and its actual incidence classes satisfy
\begin{equation}\label{c44:finite-cover-polarization}
 i_f^*\ell_{B_A}=e\ell_{B_X},\qquad
 \ell_{B_X}=\frac{(-1)^m}{2}\Delta^*d_{B_X}.
\end{equation}
In particular this is a positive transfer on a concrete class of
nonabelian varieties of geometric abelian type. It does not use
Milne's additional transfer assertion from JAC-P20.
\end{theorem}

\begin{proof}
A finite morphism between smooth varieties of the same dimension is
flat; its degree includes inseparable multiplicity. Push--pull gives
$f_*f^*=e$, and $h_A$ is ample. Lefschetz inverses for $X$ are
algebraic: in degree $i\leq d$, an inverse to $L_h^{d-i}$ is
\[
 \frac1e f_* (L_{h_A}^{d-i})^{-1}f^*.
\]
Indeed $f^*$ commutes with the indicated Lefschetz operators and
the formula composes with $L_h^{d-i}$ to the identity. The inverse
on the abelian variety is algebraic. The same pull--push construction
with its Kunneth projectors supplies those of $X$. Thus the middle
primitive projector is algebraic and self-transpose. Saito's
incidence construction and transpose compatibility, used in
JAC-T4.1, give its nondegenerate factor $B_X$.

Pullback takes primitive cycle families on $X$ to primitive cycle
families on $A$, since $h_Af^*=f^*h$. It therefore induces $i_f$.
Pushforward induces a rational left inverse up to $e$, so $i_f$ has
finite kernel. For two normalized families $z,z'$ on parameter
varieties, projection formula gives the equality of divisorial
incidence correspondences
\[
 (p_{12})_*\bigl((1\times f)^*z\cdot(1\times f)^*z'\bigr)
   =e\,(p_{12})_*(z\cdot z').
\]
By the incidence functoriality of JAC-T4.1, this is the equality
$(i_f\times i_f)^*d_{B_A}=e d_{B_X}$. Diagonal pullback gives
\eqref{c44:finite-cover-polarization}, with the same sign and
normalization as JAC-T4.3. That theorem makes $\ell_{B_A}$ ample
conditional on 032. Pullback by a homomorphism with finite kernel
is ample on its source, and division by $e>0$ preserves rational
ampleness. This proves the conclusion. At the motive level
$e^{-1}f_*f^*=1$ exhibits $h(X)$ as an actual Chow-motive summand
of $h(A)$, so the stated abelian-type scope is also explicit.
\end{proof}

\paragraph{Concrete nonabelian examples.}
In characteristic different from two, take
$A=E_0\times E_1\times\cdots\times E_{d-1}$ with $d\geq3$ odd,
choose a nonzero two-torsion point $a\in E_0$, and quotient by
\[
 (x_0,x_1,x_2,\ldots)\longmapsto(x_0+a,-x_1,x_2,\ldots).
\]
The action is free because of its first coordinate. Its quotient
is smooth projective, the cover is finite etale, and the invariant
$H^1(\mathcal O)$ has dimension $d-1$, so the quotient is not an
abelian variety. In characteristic two use a nonzero three-torsion
translation on $E_0$ and an order-three automorphism of
$E_1:y^2+y=x^3$, acting by $x\mapsto\zeta_3x$. The same argument
works. These examples prove that C44-W06 has nonabelian content in
every characteristic.

\subsection*{C44-W07. The fixed part for every smooth complete intersection}

\begin{proposition}[C44-W07; fixed-part positivity conditional on unrefereed 032]
Let $A$ have dimension $N=2m+c-1$, let $h$ be ample, and let
$k:Y\hookrightarrow A$ be a smooth successive intersection of
$c$ divisors of class $h$, with $\dim Y=n=2m-1$. There is no
generic-last-section hypothesis in this statement. Define
\[
 r=(L_h^c)^{-1}k_*,\qquad
 F=k^*\ker L_h^{c+1},\qquad E=\ker k_*\subset H^n(Y),
\]
and let $e_A$ be the primitive projector on $H^n(A)$. Then
\[
 \Pi_F=k^*e_A r,\quad \Pi_E=1-k^*r,\qquad
 H^n(Y)_{\rm prim}=F\mathbin\perp E
\]
are algebraic self-adjoint projectors and an orthogonal decomposition.
They induce a rational isogeny decomposition of primitive incidence
quotients
\[
 P^m(Y)_{\rm prim}\sim B_F\times B_E.
\]
Its signed incidence class is nondegenerate on both factors and is
a rational polarization on $B_F$, conditional on 032. Consequently
the full Th.~4 assertion for this particular $Y$ is equivalent to
ampleness of its signed incidence class on $B_E$. In particular it
holds whenever $B_E=0$.
\end{proposition}

\begin{proof}
The fixed-and-variable projector proof of H9-T13 uses smooth weak
Lefschetz and the identity $[Y]=h^c$, before any pencil monodromy is
introduced. These hypotheses hold here. Explicitly,
$rk^*=1$, and $k^*r$ is the orthogonal projector onto the ambient
middle cohomology. The primitive condition for an ambient class
$k^*u$ is $L_h^{c+1}u=0$: the Gysin map on $H^{n+2}(Y)$ is
injective by dual weak Lefschetz. The same injectivity makes every
element of $\ker k_*$ primitive. The primitive projector of $A$
is algebraic; the projectors of $Y$ outside its middle degree come
from the weak-Lefschetz identifications with $A$ and their duals.
Thus the displayed middle projectors are actual algebraic
correspondences. Self-adjointness and nondegeneracy give the
incidence decomposition exactly as in JAC-T4.1.

For clarity, the positive scalar in the fixed-part test also remains
unchanged. If a normalized fixed family pulled to a pointed curve
$C$ of genus $b\geq1$ has class $\xi$, put $\eta=(1\otimes r)\xi$
and lift its parameter degree-one component to an algebraic class
$w$ on $J(C)\times A$. Then $h^{c+1}w=0$ and, for
$H=\theta+h$,
\[
 H^{b+c}w=0,\qquad
 \int_{J(C)\times A}w^2H^{b+c-1}
   =\frac{(b+c-1)!}{c!}\int_{C\times Y}\xi^2.
\]
This is H9-T13, equation (H9.47); its proof needs no genericity.
The coefficient is positive. Conditional on 032, the signed curve
degree is nonnegative, with equality only for a zero mixed class.
The nondegenerate signed divisor is therefore ample by the
nef-plus-nondegenerate criterion of JAC-T4.3. Orthogonality makes
ampleness of the full signed divisor equivalent to ampleness of
both factors, proving the final assertions.
\end{proof}

\subsection*{C44-W08. Specialization of the inherited incidence polarization}

\begin{theorem}[C44-W08; conditional on unrefereed 032]
Let $\mathscr Y\to U$ be the smooth locus of a genuine Lefschetz
pencil in the setting of H9-T13, and let $s\in U$ be any geometric
point. Write $B_{\bar\eta}=P^m(Y_{\bar\eta})_{\rm prim}$ for the
geometric generic primitive incidence quotient. After a finite
extension of the local generic field, $B_{\bar\eta}$ has good
reduction at $s$, and its special fibre maps with finite kernel to
$P^m(Y_s)_{\rm prim}$. The actual signed incidence form is positive
on this inherited image. Hence Th.~4 holds for $Y_s$ whenever
\begin{equation}\label{c44:no-incidence-jump}
 \dim P^m(Y_s)_{\rm prim}=\dim P^m(Y_{\bar\eta})_{\rm prim}.
\end{equation}
More particularly, if the generic variable incidence factor of the
pencil is nonzero, then Th.~4 holds for every smooth geometric fibre
of the pencil. If that generic factor is zero, only the new
variable incidence directions that occur at a special fibre remain
outside this theorem.
\end{theorem}

\begin{proof}
Work over the strict henselian local trait at $s$, or its completion,
and pass to a finite extension of its fraction field over which an
abelian representative $B$ and a normalized rational Poincare cycle
$z$ are defined. These are finite data; rational descent is the one
fixed in JAC-P00. As in H9-T5, equations (H9.10)--(H9.11), the
mixed class of $z$ injects the Tate module of $B$ into the appropriate
Tate twist of $H^n(Y_{\bar\eta})$. Indeed its self-incidence is
nondegenerate, so a vector killed by the family map would be in
the radical of that incidence pairing. Smooth proper base change
makes the target unramified at this smooth fibre. Thus $B$ has
good reduction by the Neron--Ogg--Shafarevich criterion, just as in
H9-T5. One exact primary author source for the criterion is Milne,
\emph{Complex Multiplication}, Theorem 6.12, p.~51, which cites
Serre--Tate (1968), Theorem 1.

Let $\mathscr B$ denote the resulting abelian scheme over the trait.
Take the closure of a rational-cycle representative of $z$ in
$\mathscr B\times\mathscr Y$. These schemes are smooth over a
regular trait, so the refined fibre restrictions and intersection
products are defined. Its special-fibre mixed class is the
specialization of the generic mixed class. In particular it is
still in parameter degree one and is primitive on $Y_s$; these
conditions are kernels of maps of the smooth-proper cohomology
systems. Rigidification or degree-one projection, if desired,
alters no incidence class.

Form the actual relative self-incidence divisor by the same proper
pushforward of the product of the two cycle families. Clear its
rational denominator. Its signed half-diagonal pullback is a line
bundle $\mathscr L$ on $\mathscr B$, and refined base change says
that its two fibre classes are exactly the corresponding signed
incidence classes, up to the one positive clearing factor. The
generic class is ample by H9-T13.

Here is the specialization step without an assertion that arbitrary
special cycles lift. Choose a relatively ample line bundle
$\mathscr H$ on $\mathscr B$. Flatness and the Hilbert polynomial
make
\[
 \chi(L_t+uH_t)=\frac{(c_1(L_t)+u c_1(H_t))^{\dim B}}{(\dim B)!}
\]
the same real polynomial in $u$ for the generic and special fibres
$t$. On the generic fibre this polynomial is positive for every
$u\geq0$, since $L_{\eta}+uH_{\eta}$ is ample. Thus on the
special fibre $\chi(L_s)\ne0$ and the polynomial has no positive
real root. Mumford's index theorem gives index zero for $L_s$;
the nondegenerate effective criterion makes it ample. These are
exactly the two abelian-variety criteria used in the checked proof
of JAC-T4.3 (Mumford, \emph{Abelian Varieties}, Chapter III,
Section 16, and Chapter II, Section 6, Application 1).

The specialized family induces a homomorphism
$u:B_s\to P^m(Y_s)_{\rm prim}$. Saito's transpose and incidence
compatibility identifies $u^*\ell_s$ with the ample class just
constructed. Therefore $u$ has finite kernel: restriction to any
positive-dimensional connected kernel would make an ample class
zero. This proves positivity on the inherited image. Under
\eqref{c44:no-incidence-jump} the map is an isogeny, and ampleness
descends through that isogeny.

For the final assertion apply the same construction to the generic
variable factor alone. If that factor is nonzero, the irreducibility
after finite-index restriction in H9-T13, or H9-T5 (H9.11), forces
its family map to fill the entire variable local system $E$.
Specialization preserves this rank. The special variable incidence
factor injects into $E_s$ by its own nondegenerate self-incidence,
so its dimension is at most $\operatorname{rank}(E)/2$. The
inherited image already has that dimension. It therefore maps
isogenously onto the special variable incidence factor and gives
its positive signed form. C44-W07 supplies the fixed factor for
every smooth fibre. This proves the stated all-smooth-fibre scope.
\end{proof}

\paragraph{The remaining special-fibre gate.}
When the generic variable incidence quotient vanishes, the
construction specializes only the fixed incidence quotient. A
special fibre may have a new variable incidence factor. Neither
smooth proper base change nor constancy of the rank of $E$ makes
that factor extend as algebraic cycle families through the pencil.
Its signed positivity is the smaller unresolved statement. Thus
C44-W08 does not turn H9-T13 into an unconditional assertion about
every smooth complete intersection; it proves precisely the
rank-preserving specializations and the nonzero-generic-variable
case, with 032 retained as the positivity input.

\subsection*{C44-W09. The full geometric-abelian-type extension with its extra input exposed}

\begin{proposition}[C44-W09; conditional on unrefereed 032 and the separate JAC-P20 transfer]
Assume, in addition to 032, Milne's Proposition 4.2 in
\emph{Grothendieck's standard conjecture of Lefschetz type over
finite fields}, version 14 July 2022, in the total-cohomology form
specified and audited in JAC-P20. Then JAC-T4.3 holds for every
smooth projective odd-dimensional variety of geometric abelian
type, with its actual signed primitive incidence divisor.
\end{proposition}

\begin{proof}
JAC-P16 gives algebraic Lefschetz and primitive projectors. For
every pointed test curve $C$, the product $J(C)\times X$ is still
of geometric abelian type. JAC-P20's actual abelian cohomology
injection and the additional transfer assertion give the Hodge
standard conjecture for this product. The construction of
JAC-T4.2, using the algebraic Abel inverse on the first factor,
now gives an algebraic primitive class $w$ on $J(C)\times X$ and
the exact identity
\[
 \int_{J(C)\times X}w^2(\theta+h)^{b-1}
       =(b-1)!\int_{C\times X}\xi^2.
\]
The Hodge-standard input on this particular product makes its
signed value nonnegative. Saito's nondegenerate incidence factor
and JAC-T4.3's curve criterion then make the signed divisor ample.
This argument proves the stated conditional corollary; it does
not establish the additional transfer assertion itself.
\end{proof}

\paragraph{Framework and precise dependency.}
C44-W06 proves a positive scalar comparison of the actual forms,
not merely an embedding of motives. C44-W07--W08 use smooth
proper base change, SGA 7's fixed and variable projectors, and
the extension of actual cycle correspondences across a regular
trait. Their target is cote 44, scans 008--009, Th.~4 and its
primitive incidence higher Picard variety. The new residual gate
is the positive sign on a special new variable incidence factor,
not the already checked generic pencil sign. For all varieties
of geometric abelian type, JAC-P20 still leaves the comparison
between an abstract transported positive form and the particular
geometric form as an additional dependency. C44-W09 records
the exact implication if that dependency is supplied.


\section{The supersingular Fermat test for cote 16}

\subsection*{C44-S01. The proposed Fermat candidate already has the required sign}

\begin{proposition}[C44-S01; unconditional]
Let $X=X_N^{2r}\subset\mathbf P^{2r+1}$ be the smooth Fermat
hypersurface of degree $N\geq3$ in characteristic $p\nmid N$,
with $r\geq2$.  Suppose $p^\nu\equiv-1\pmod N$ for some
$\nu\geq1$.  Then all even cohomology is algebraic and all odd
cohomology vanishes.  Numerical and homological equivalence
coincide.  The signed form
\[
 Q_V(x,y)=(-1)^r\deg(xy),\qquad
 V=\ker(h:H^{2r}(X)\longrightarrow H^{2r+2}(X)),
\]
is positive definite on its rational algebraic cycle structure.
In particular the Hodge standard conjecture for $X$ is known
unconditionally; it is not a consequence that needs 032.
\end{proposition}

\begin{proof}
Put $q=p^\nu$ and $b=(q+1)/N$.  The Hermitian Fermat
hypersurface $H=X_{q+1}^{2r}$ maps finitely and surjectively to
$X$ by
\[
 [x_0:\cdots:x_{2r+1}]
 \longmapsto[x_0^b:\cdots:x_{2r+1}^b].
\]
Its degree is $b^{2r+1}$: $f^*h_X=bh_H$ and
$\int_Hh_H^{2r}=q+1=bN$ determine the degree by projection
formula.  In particular $f_*f^*=b^{2r+1}$ and $f^*$ is injective.
Shimada, \emph{Lattices of Algebraic Cycles on Fermat Varieties
in Positive Characteristics}, Theorem 1.1(1)--(2) and
Corollary 1.2, proves that linear cycles span the middle
cohomology of $H$ and that its signed primitive intersection
form is positive definite.  Weak Lefschetz gives the other
cohomology groups.  Pushing forward these algebraic classes
shows algebraicity of all cohomology of $X$.

If $x\in V$, then
$h_Hf^*x=b^{-1}f^*(h_Xx)=0$, and
\[
 (-1)^r\int_H(f^*x)^2
  =b^{2r+1}(-1)^r\int_Xx^2.
\]
The injectivity of $f^*$ and Shimada's positive form prove the
claim.  In all other degrees $X$ has the ambient projective
space Lefschetz strings, whose sign is immediate.  The
intersection pairing is nondegenerate on the rational
algebraic classes because these span cohomology; thus
numerical and homological equivalence agree.  Since
$\operatorname{NS}(X)_\mathbf Q=\mathbf Qh$, this also covers
every ample polarization.
\end{proof}

\paragraph{Exact status of the proposed test.}
Shioda--Katsura, \emph{On Fermat Varieties}, Theorem II,
printed p.~97, says for even dimension and $N\geq4$ that
algebraic generation of cohomology is equivalent to the
displayed congruence.  Thus, in that range, the entire stated
supersingular Fermat candidate class is covered by C44-S01.
The new argument above is only the finite-cover descent
of Shimada's already known sign.  It neither attributes a
new theorem to Grothendieck nor treats algebraicity alone as
a positivity theorem.  The square needed in scan 049 is
handled explicitly below.

\subsection*{C44-S02. An explicit smooth Segre hyperplane}

\begin{proposition}[C44-S02; unconditional]
Let $X=X_N^d$ be a smooth Fermat hypersurface with $N\geq3$.
Choose a finite field $k_0$ containing every $(N-2)$-nd root
of unity in $\overline{\mathbf F}_p$.  Let $M>N(d+1)$ and
choose $\beta$ of degree $M$ over $k_0$ in its extension of
degree $M$.  Define $c_i=\beta^{(N-2)i}$ for
$0\leq i\leq d+1$.  Then
\[
 W=\left\{\sum_{i=0}^{d+1}c_ix_iy_i=0\right\}
       \subset X\times X
\]
is a smooth hyperplane section for the Segre polarization
$L=h_1+h_2$.  Also $Y=X\cap\{x_{d+1}=0\}$ is a smooth
hyperplane section of $X$.
\end{proposition}

\begin{proof}
The assertion for $Y$ follows from $p\nmid N$.  If $W$ had
a singular point, the differential of its equation would
annihilate both projective tangent spaces.  Since the
Fermat gradients are $N(x_i^{N-1})$ and $N(y_i^{N-1})$,
there would be scalars $\lambda,\mu$ with
\[
 c_iy_i=\lambda x_i^{N-1},\qquad
 c_ix_i=\mu y_i^{N-1}\quad\text{for every }i.
\]
Both scalars are nonzero, since all $c_i$ are nonzero and
neither projective coordinate vector is zero.  The nonzero
coordinates of $x$ and $y$ have the same nonempty support
$S$.  Substitution gives, for $i\in S$,
\[
 (x_i^N)^{N-2}=(\mu\lambda^{N-1})^{-1}c_i^N.
\]
Choose $\rho$ with $\rho^{N-2}=(\mu\lambda^{N-1})^{-1}$.
It follows that
$x_i^N=\rho\zeta_i\beta^{Ni}$ for roots of unity
$\zeta_i\in k_0$ with $\zeta_i^{N-2}=1$.
This remains valid if $p\mid(N-2)$, since only the distinct
roots are used.  The Fermat equation for $x$ now gives
\[
 \sum_{i\in S}\zeta_i\beta^{Ni}=0.
\]
This is a nonzero polynomial over $k_0$ of degree at most
$N(d+1)<M$ vanishing at $\beta$, a contradiction.  The
Jacobian criterion therefore proves smoothness.
\end{proof}

\subsection*{C44-S03. The exact square sign and all three corrected scan-050 conditions}

\begin{theorem}[C44-S03; unconditional]
For $X=X_N^{2r}$ of C44-S01, put $d=2r$, $T=X\times X$,
$L=h_1+h_2$, and take $Y,W$ as in C44-S02.  Then
\[
 C_{\rm G}(X),\quad C_{\rm G}(Y),\quad R_0(T,W),\quad
 G_{L,0}(T,W),\quad A_0(T),\quad G_0(T,W),\quad A_0(W)
\]
all hold with the meanings fixed by SIGN\_NEXT SN-C03--SN-C04.
The Hodge standard conjecture $I(T)$ holds.  In its middle
degree the primitive space and form are explicitly
\begin{gather}
 P^{2d}(T)=\mathbf Qw\ \perp\ (V\otimes V),\qquad
 w=\sum_{i=0}^{d}(-1)^ih_1^ih_2^{d-i},\label{c44:fermat-square-primitive}\\
 \int_T(cw+z)^2=(d+1)N^2c^2+(Q_V\otimes Q_V)(z,z)>0
 \quad\text{if }(c,z)\ne(0,0).\label{c44:fermat-square-sign}
\end{gather}
The ring of degree-zero rational correspondence classes
modulo homological equivalence is
\[
 \prod_{i=0}^{d}\operatorname{End}_{\mathbf Q}
       \bigl(\operatorname{Alg}^i(X)\bigr)
 \simeq \mathbf Q^d\times M_{b_d(X)}(\mathbf Q),
\]
and is therefore semisimple.  These are the actual sign and
semisimplicity conclusions of cote 16, scan 049, for this
nonabelian family.
\end{theorem}

\begin{proof}
Weak Lefschetz gives the orthogonal decomposition
\[
 H^{\rm even}(X)=U\perp V,\qquad
 U=\mathbf Q\langle1,h,\ldots,h^d\rangle,
\]
where $V$ is in cohomological degree $d$ and $hV=0$.
The rational structures here are given by cycle classes;
they become the full cohomology after scalar extension.
The middle cohomology of $T$ consequently consists of the
ambient terms, the two terms $V\otimes h^r$ and
$h^r\otimes V$, and $V\otimes V$.  On the ambient
terms the kernel of multiplication by $L$ is the single
line displayed in \eqref{c44:fermat-square-primitive}:
successive coefficients satisfy $c_i+c_{i-1}=0$.
Multiplication by $L$ is injective on each of the two
cross terms and their images have different Kunneth
degrees.  It is zero on $V\otimes V$.  Orthogonality
follows from $\int_Xvh^r=0$ for $v\in V$.

In $w^2$ the $d+1$ contributing pairs have indices
$i,d-i$ and sign $(-1)^d=1$, giving
$\int_Tw^2=(d+1)N^2$.  The pairing on $V\otimes V$
is the tensor product of the two raw pairings; their
two factors $(-1)^r$ cancel.  It is therefore
$Q_V\otimes Q_V$, a positive definite form.  This proves
the exact critical sign.

For the full Hodge standard statement, decompose
$H^*(T)$ into the four $L$-stable blocks
$U\otimes U$, $V\otimes U$, $U\otimes V$, and
$V\otimes V$.  They are orthogonal.  The first has
the Lefschetz module and intersection forms of
$\mathbf P^d\times\mathbf P^d$, multiplied by $N^2$;
its signed forms are positive by characteristic-zero
Hodge--Riemann, with identical rational matrices in
every characteristic.  In $V\otimes U$ the primitive
classes are exactly the base $V\otimes1$, in
codimension $r$.  Their signed form is $NQ_V$;
the same is true for $U\otimes V$.  The last block
is primitive in middle degree and was just treated.
This proves $I(T)$ in all degrees.  Every ample
class on $T$ is a positive linear combination of
$h_1,h_2$; the same argument, with positive rescalings,
also handles these polarizations if required.

For a smooth hypersurface $Z$ of degree $N$ and
dimension $e$, the Lefschetz inverses below its middle
degree are already explicit.  For even $i<e$ the
correspondence
\[
 \frac1N h_Z^{i/2}\times h_Z^{i/2}
\]
induces the inverse of
$L_Z^{e-i}:H^i(Z)\longrightarrow H^{2e-i}(Z)$;
for odd $i<e$ both spaces vanish.  The middle inverse
is the identity.  Thus $C_{\rm G}(X)=B(X)$ and
$C_{\rm G}(Y)=B(Y)$.  The product stability of $B$
(Kleiman 1968, Proposition 2.12), or the cycle-basis
matrix units below, gives $B(T)$ and in particular
the algebraicity of $(L^2)^{-1}$ in the critical
degrees of SN-C03.

Write $j:W\hookrightarrow T$.  If $j^*x$ is algebraic
for $x\in H^{2d-2}(T)$, then
$L j_*j^*x=L^2x$ is algebraic; applying the algebraic
inverse proves $R_0(T,W)$.  For
$\gamma\in\operatorname{Alg}^{d+1}(T)$, set
$x=(L^2)^{-1}\gamma$ and $y=j^*x$.  Both are
algebraic and $Lj_*y=\gamma$, proving $G_{L,0}$.
The critical $A_0(T)$ was used here; the equivalence
in SN-C03 supplies $G_0$ and $A_0(W)$ as well.
This verifies the maps with the required extra $L$;
no unweighted lower-Gysin replacement is made.

Finally choose rational cycle bases for each
$\operatorname{Alg}^i(X)$ and the intersection-dual
bases in degree $d-i$.  Exterior products of these
classes give all rational matrix units on
$H^{2i}(X)$.  Every degree-zero correspondence
acts degreewise; Kunneth and Poincare duality make
this action faithful modulo homological equivalence.
It is therefore exactly the product of the stated
full matrix algebras.  Since all nonmiddle Betti
numbers in even degree are one, the displayed
description follows.
\end{proof}

\paragraph{What task 4 does and does not close.}
C44-S01--C44-S03 rule out the suggested supersingular Fermat
family as an example whose square-sign status is unknown.
They verify the corrected scan-050 proposition and the
scan-049 sign and correspondence ring explicitly and
unconditionally on that family.  A different smooth
projective $X$ with a genuinely unestablished $I(X\times X)$
is still required for the adversarial part of task 4; absence
of a source for one named equation would not certify that
status.  None of these findings closes the special-axis
sign for a complete intersection in an abelian variety.

\paragraph{Framework and programme crosswalk.}
For C44-S01--S03 the relevant archive locators are
\texttt{ARCHIVE\_16\_scan049.tex}, the semisimple
correspondence paragraph, and
\texttt{ARCHIVE\_16\_scan050.tex}, the proposition and
final weighted diagram.  They instantiate SIGN\_NEXT
SN-C03--SN-C04 and its successor question 4.  Their
Lefschetz inverses also give the $C_{\rm G}$ input to
cote 44, Th.~6, scan 004, through JAC-P02.  Weak
Lefschetz and Gysin are the Weil-cohomology operations
in SGA 5, Exp. VII, Section 7 (complete intersections).
The finite power maps and explicit hyperplanes are
scheme morphisms over finite extensions of $\mathbf F_p$;
the identities survive every Weil cohomology realization.
At the level of homological motives the Fermat object
is a direct sum of Tate objects with the positive
Lefschetz forms just computed.  No equality with its
Chow motive modulo rational equivalence is silently
deduced.  C44-W01--W05 use the exterior-power, hence
$\lambda$-operation, of $h^1(A)$, the algebraic
endomorphism projector, and rational cycle intersections.
No additional $K$-theory assertion, Selmer injection,
or motivic $L$-function nonvanishing is inferred.

% END work/scan_w/contribution.tex

% BEGIN work/scan_w/relative_th4.tex
% C44_POS /root/scan_w; GPT-6 Astra Pro; effort not exposed.
\section{Theorem 4 from an extending odd-degree abelian correspondence}

\begin{theorem}[C44-W10: a relative criterion for the actual incidence polarization]
Let $R$ be a complete discrete valuation ring of mixed characteristic
$(0,p)$ with algebraically closed residue field $k$ and fraction field
admitting an embedding into $\mathbf C$. Fix $\ell\ne p$.
Let $\mathscr X/R$ be smooth projective of relative dimension
$n=2m-1$, with relatively ample class $h$, and write $X=\mathscr X_k$.
Suppose a relative rational algebraic correspondence $e$ realizes the
orthogonal primitive projector
\[
 P^n(\mathscr X)=\ker(L_h:H^n(\mathscr X)\longrightarrow
                                      H^{n+2}(\mathscr X)),
 \tag{C44.W10.1}
\]
extended by zero in other degrees. Suppose also that relative
algebraic correspondences give an injection on every geometric fibre
\[
 \alpha:P^n(\mathscr X)\longrightarrow
       \bigoplus_\nu H^{2s_\nu-1}(\mathscr A_\nu)(s_\nu-m),
 \tag{C44.W10.2}
\]
where the $\mathscr A_\nu/R$ are abelian schemes and $s_\nu\geq1$.

Then, conditional on the unrefereed 032/Milne 2026 input in C44-A08,
cote~44, Theorem~4 holds for $X$ with this primitive polarization.
Precisely, on Saito's primitive incidence quotient
$B=P^m(X)_{\rm prim}$ the actual rational incidence divisor is ample:
\[
       \ell_B=\frac{(-1)^m}{2}\,\Delta_B^*d_B
                  \in\operatorname{NS}(B)_{\mathbf Q}.
 \tag{C44.W10.3}
\]
After multiplication by a positive integer it defines a polarization,
including when $p=2$. No extension to characteristic zero of a test
map $C\to B$ or of its cycle family is assumed. Milne 2022,
Proposition~4.2, is not an additional input.
\end{theorem}

\begin{proof}
The correspondence $e$ induces an idempotent on $P^m(X)$ by the
incidence-detection result JAC-P01: $e^2-e$ has zero relevant odd
cohomological action. In middle degree, primitive cohomology is the
orthogonal complement of $hH^{n-2}(X)$, and hard Lefschetz makes both
summands nondegenerate. Thus $e$ is self-transpose on the relevant
cohomology. Saito's transpose compatibility for the canonical
incidence duality makes its image $B$ a nondegenerate incidence
factor. This is exactly the argument of JAC-T4.1; that argument uses
the algebraicity and self-transpose property just supplied, not that
$X$ itself is abelian. It gives a normalized rational cycle family
$z$ on $B\times X$, with class entirely in
\[
                H^1(B)\otimes P^n(X),
\]
whose self-incidence is $d_B$. The associated symmetric rational
homomorphism is an isogeny.

Let $f:C\to B$ be the normalization of an arbitrary integral curve
in $B$, translated so that a chosen point maps to the origin. Its
genus $b$ is positive, since an abelian variety has no rational
curves. The pullback $\xi=(f\times1)^*z$ is a rational algebraic
class in
\[
                  H^1(C)\otimes P^n(X).
 \tag{C44.W10.4}
\]
We will prove the actual geometric sign
\[
                         (-1)^m\int_{C\times X}\xi^2\geq0,
 \tag{C44.W10.5}
\]
with strict inequality if its cohomology class is nonzero.

\paragraph{Lifting the test curve alone.}
Every smooth proper curve over $k$ lifts to a smooth proper curve
$\mathscr C/R$. Here is a precise way to use the curve deformation
theorem over this particular trait. The stack of smooth proper
curves is smooth over $\operatorname{Spec}\mathbf Z$ (Stacks
Project, Tag~0E1N, Lemma~109.16.2). Choose a scheme atlas near the
point represented by $C$. Since $k$ is algebraically closed, after
passing to an atlas point there is a $k$-point representing this
curve. The atlas is smooth over $R$. A smooth scheme over a
henselian local ring has an $R$-point through any given residue-field
point: use an etale coordinate map to affine space, lift its
coordinates, and apply the henselian lifting property. Pulling back
the universal curve gives $\mathscr C$. The chosen point of $C$
lifts by the same argument to a section. Its divisor has degree one
on each fibre and is relatively ample, so this curve model is
projective. Write $u$ for that relative ample divisor class.

The relative Jacobian $\mathscr J=\operatorname{Pic}^0_{\mathscr C/R}$
is an abelian scheme. Its Abel and Poincare correspondences identify
$H^1(\mathscr C)$ with $H^1(\mathscr J)$ on every geometric fibre,
by algebraic correspondences over $R$. Concretely, if $a$ is the
relative Abel map and $\theta$ the canonical principal polarization,
an inverse to $a^*$ on degree one is
\[
        j_C=(b-1)!\,(L_\theta^{b-1})^{-1}a_*.
 \tag{C44.W10.6}
\]
The inverse Lefschetz correspondence is relative algebraic by the
explicit abelian operators in C44-A08. The formula follows fibrewise
from $a_*[C]=\theta^{b-1}/(b-1)!$ and the projection formula, as in
JAC-T4.2. A relatively ample representative of the polarization can
be chosen over this complete trait after a finite extension if
necessary; the argument is unchanged by that extension.

\paragraph{A primitive middle summand with its actual duality.}
On $\mathscr C\times_R\mathscr X$ put
\[
                  V=H^1(\mathscr C)\otimes P^n(\mathscr X).
\]
Its degree is $1+n=2m$, the middle degree of this $2m$-fold.
The product of the relative degree-one curve projector and $e$
is an algebraic projector with image $V$. Both factors are
self-transpose for their middle Poincare dualities, so the product
is self-transpose for middle Poincare duality on the product.
Its restricted pairing is nondegenerate. The inverse duality map
is represented algebraically by this projector together with the
middle Poincare duality correspondence. Thus the algebraic-inverse
hypothesis of C44-A08 holds here, rather than being inferred from
an abstract choice of positive form.

For the product ample class $u+h$ we have
\[
                       (u+h)V=0.
 \tag{C44.W10.7}
\]
Indeed $uH^1(\mathscr C)=0$ by the degree on a curve and
$hP^n(\mathscr X)=0$ by (C44.W10.1). Hence $V$ is primitive.
Tensoring (C44.W10.2) with (C44.W10.6), and then using the Kunneth
summand of a product of abelian schemes, gives an extending
algebraic injection
\[
 V\longrightarrow
 \bigoplus_\nu H^{2s_\nu}(\mathscr J\times_R\mathscr A_\nu)
                                  (s_\nu-m).
 \tag{C44.W10.8}
\]
It is injective on every geometric fibre. All hypotheses of
C44-A08 are now verified with $r=m$ and dimension $2m$.
Its signed primitive form is exactly
$(-1)^m\int_{C\times X}\xi^2$, with no additional Lefschetz
factor or comparison scalar. It proves (C44.W10.5).
The cycle $\xi$ itself, and the map $f$, have been used only in the
special fibre; neither was lifted.

\paragraph{The incidence divisor.}
The normalized incidence identity and refined pullback give
\[
                  \deg_C f^*\ell_B
                         =\frac{(-1)^m}{2}
                                       \int_{C\times X}\xi^2\geq0.
 \tag{C44.W10.9}
\]
The degree-one projection of the parameter family removes the
degree-zero and degree-two parameter components in this formula.
Since every integral curve in $B$ has been tested, $\ell_B$ is nef.
Its associated symmetric homomorphism is an isogeny by the
nondegeneracy of the primitive incidence factor proved above.
A nef nondegenerate rational divisor on an abelian variety is ample,
by the criterion used in JAC-T4.3 (Mumford's index theorem and the
effective nondegenerate criterion). This proves (C44.W10.3).
The zero-dimensional case for $B$ is immediate. All normalizations
are rational; division by $2$ therefore creates no exception at
characteristic $2$.
\end{proof}

\begin{corollary}[C44-W11: odd Fermat varieties and their product families]
Let $k=\overline{\mathbf F}_p$, let $d\geq2$ with $p\nmid d$, and let
\[
             X=\prod_{j=1}^{a}F_{n_j,d,k},\qquad
             n_j\geq1,\qquad \sum_jn_j=2m-1.
\]
For every polarization $h=\sum_j b_jh_j$ with $b_j\in\mathbf Q_{>0}$
and factor hyperplane classes $h_j$, the actual primitive incidence
divisor (C44.W10.3) is ample, conditional on the unrefereed
032/Milne 2026 input. In particular cote~44, Theorem~4 holds for
every odd-dimensional smooth Fermat hypersurface in every good
characteristic, with every rational ample class. For the displayed
products it holds for every rational ample class if all $n_j\geq3$;
when curve or surface factors occur, only the displayed product
polarizations are asserted.
\end{corollary}

\begin{proof}
C44-F08 gives an actual relative split Chow-motive injection of
the Fermat product into a sum of Tate twists of powers of the
relative Fermat Jacobian, over
$\operatorname{Spec}\mathbf Z[1/d,\zeta_{2d}]$.
Restriction to odd middle degree gives (C44.W10.2). C44-F09 gives
relative Lefschetz operators and the primitive projector for every
displayed $h$. In middle degree its primitive image is orthogonal
to the Lefschetz image, so this projector is self-transpose as
required in C44-W10. For a single odd Fermat hypersurface this can
also be read directly from
\[
 [\Delta]-\frac1d\sum_{i=0}^{n}h^{n-i}\times h^i;
\]
transposition permutes the terms and its image is the primitive
middle cohomology.

Choose a prime of the cyclotomic base above $p$. Its local ring is
a discrete valuation ring. Complete its maximal unramified
extension and choose the algebraic closure of the residue field
to identify that residue field with $k$. This gives a complete
mixed-characteristic discrete valuation ring to which all the
relative correspondences and abelian schemes pull back. Its
fraction field has cardinality at most that of $\mathbf C$ and
characteristic zero, so it admits an abstract embedding into
$\mathbf C$: choose a transcendence basis and extend across the
remaining algebraic extension. A common such embedding is used
for the whole construction. An embedding is only needed to apply
complex Hodge--Riemann to this common generic fibre; no continuity
with a $p$-adic topology is required. Consequently C44-W10 applies.

For one odd-dimensional hypersurface of dimension at least three,
weak Lefschetz gives rational Neron--Severi rank one, generated by
the hyperplane class. In dimension one, numerical divisor classes
are likewise determined by their degree, so every rational ample
class has the same Lefschetz action as a positive rational multiple
of the hyperplane class. These observations prove the single-factor
assertion for all ample classes. For products with all $n_j\geq3$,
the Neron--Severi computation in C44-F10 shows that every ample
class is a positive rational combination of the factor hyperplane
classes. No such computation is invoked for products containing
curves or surfaces.
\end{proof}

\paragraph{Scope and the remaining geometric-type gate.}
C44-W10 is a sufficient relative condition, and C44-W11 supplies a
nonabelian family meeting it. It strengthens C44-W09 on this family
by replacing its additional Milne2022 transfer input with actual
extending correspondences and an explicit geometric-form proof.
It does not show that an arbitrary variety of geometric abelian
type has such a lift or extending injection. The input from 032
is still conditional; the constructions F08--F09 and the test-curve
lifting are unconditional. The result concerns the incidence
quotient of cote~44, scans~008--009, with the primitive qualification
in scan~008, and supplies a positive sign on its actual divisor,
not a separately transported polarization.


% END work/scan_w/relative_th4.tex

% BEGIN work/scan_w/qt4_abelian_type_transfer_reaudit.tex
% C44_POS /root/scan_w. Focused source audit; no new theorem or gate closure.
\section{Focused re-audit of the geometric-abelian-type transfer in Theorem 4}
\label{sec:qt4-transfer-reaudit}

\paragraph{Outcome.}
Conditional on the unrefereed preprints 032 and Milne's 2026 report,
their accepted inputs provide
the full Tate conjecture over finite fields and the Hodge standard
conjecture for abelian varieties, including all their products.
The checked arguments do not independently identify an arbitrary
geometric-abelian-type variety's own primitive intersection form
with the positive form transported from that category.
Consequently C44-W09 retains its separate Milne 2022,
Proposition~4.2, transfer input. This audit neither refutes that
assertion nor promotes a new conditional formulation as progress.

\paragraph{The precise positive inputs.}
In the supplied 032 CM manuscript, Corollary~8.3, printed
pp.~49--50, states algebraic spanning of Tate classes and
numerical--homological equality for abelian varieties over finite
fields. Corollary~8.4, printed p.~50, states positivity of
\[
       (-1)^r\deg(xyh^{g-2r})
\]
on the actual primitive algebraic classes of an abelian variety,
for every rational ample class. Its proof invokes Milne 2002,
Theorems~2.1 and 3.3, with the stated addendum and correction.
Products of abelian varieties are still abelian varieties, so
these inputs include all powers and all relevant Hom spaces.
Neither corollary asserts the same geometric-form statement for
every variety whose motive belongs to the generated category.

Milne's supplied report, dated 9 November 2026,
\emph{On the Hodge, Tate, and standard conjectures for abelian
varieties}, Theorems~3.1 and 3.5, printed pp.~6 and 8, has the
same abelian-variety scope. The table in Section~3.4, p.~7,
defines its category $\operatorname{Mot}(\mathbb F)$ using
abelian varieties and algebraic correspondences.
Section~3.8, p.~8, explicitly requires a category polarization
\emph{for which the geometric pairings are positive};
Section~3.10, pp.~8--9, describes the compatible CM reduction
polarization. Thus the existence of an abstract category
polarization is not the entire geometric conclusion even in
the report's own formulation.

\paragraph{What the published polarization theorem checks.}
Milne 2002, Proposition~3.2, printed pp.~606--607, distinguishes
existence of a polarization from positivity of the specified
geometric forms $\varphi^r$ and $\vartheta^r$.
Theorem~2.1, pp.~604--605, establishes compatibility with CM
reduction. In the proof of Theorem~3.3, pp.~607--608, the actual
abelian geometric form is obtained by lifting the polarized
abelian variety up to isogeny and specializing that form.
Corollary~3.4, p.~608, retains the category based on abelian
varieties. The addendum concerns the faithful character-module
calculation in Theorem~2.1; it adds no arbitrary-variety transfer.

These passages were read in the primary journal-version author
PDF, including the displayed proofs of Lemma~1.2,
Proposition~1.4, Theorem~2.1, Proposition~3.2 and Theorem~3.3:
\url{https://arxiv.org/pdf/math/0103175}, version~3.
The author's complete addendum was read at
\url{https://www.jmilne.org/math/articles/2002a.html}.

\paragraph{The actual algebraic transfer and the unproved sign step.}
Let $X$ be of geometric abelian type. JAC-P16 already proves
numerical--homological equality, algebraic Kunneth and primitive
projectors, and the algebraic hard-Lefschetz inverse on $X$
and its powers. JAC-P20 constructs an actual algebraic injection
of total cohomology into the cohomology of an abelian variety;
Tate twists and multiple summands are separated by auxiliary
elliptic Kunneth factors. There is no remaining assertion here
that only an abstract isomorphism of realizations was produced.

Fix a primitive, Tate-normalized motive
$M=P^{2r}h(X)(r)$. The positive polarization of the abelian
category restricts to a positive form $Q_0$ on this summand.
One can choose $Q_0$ rationally by pulling back the appropriate
direct sum of rational geometric abelian forms along the
rational algebraic inclusion.
The geometric pairing is a separate algebraic morphism
\[
 Q_X:M\otimes M\longrightarrow\mathbf1,
 \qquad Q_X(x,y)=(-1)^r\int_Xxyh^{\dim X-2r}.
\]
Its nondegeneracy follows from the already established
Lefschetz theory and equality of equivalence relations.
There is therefore an invertible $Q_0$-self-adjoint algebraic
endomorphism $a$ such that
\[
                       Q_X(x,y)=Q_0(ax,y).
\]
On the rational algebraic space $\operatorname{Hom}(\mathbf1,M)$,
the required sign is exactly the positivity of this particular
operator's real eigenvalues. Semisimplicity supplies its
algebraic inverse and the category supplies $Q_0$; neither
assertion selects those signs.

Full Tate algebraicity does not resolve this distinction by
itself. It makes invariant morphisms available, but both a
form and its negative remain algebraic morphisms. Moreover
$Q_0$ and $-Q_0$ have exactly the same adjoint involution on
$\operatorname{End}(M)$ and hence the same adjoint trace pairing.
Thus even positivity of that endomorphism trace pairing cannot
select the scalar sign of the form. This is a formal limitation
of the proposed inference, not an example of a geometric
variety violating Hodge standard. Higher multiplicity summands
have the corresponding self-adjoint-operator ambiguity.

The same issue persists in the actual Th.~4 curve test.
For $J(C)\times X$, the algebraic injection transports a positive
form from an abelian target; it does not identify its value
with the signed self-intersection of the primitive class used
in C44-W09. The latter is the value whose exact factorial
comparison gives the incidence divisor on the higher Picard
quotient. An unrelated ample polarization on that quotient
does not determine the sign of its given incidence divisor.

\paragraph{The exact 2022 transfer assertion.}
Milne, \emph{Grothendieck's standard conjecture of Lefschetz type
over finite fields}, 14 July 2022, v1.1,
Proposition~4.2, printed pp.~7--8, explicitly states the broader
transfer from an algebraic total-cohomology injection.
Its complete proof is the citation of Kleiman 1968, Theorem~3.11,
and Saavedra Rivano, VI.4.4.2. The adjacent Theorem~4.5 uses
Proposition~4.2 and is consequently not an independent route
around this step. These statements and proofs were reread in
\url{https://www.jmilne.org/math/articles/LFF.pdf}.

Kleiman's cited theorem and its whole proof were reread in the
primary scan, printed pp.~380--381. Algebraicity of the
geometric adjoint uses $B(X)$ and $B(Y)$; its strict trace
positivity separately assumes $I(X\times Y)$ for the product
polarization. Using that positivity for the unknown source
therefore does not prove the proposed transfer.
The primary scan is
\url{https://download.uni-mainz.de/mathematik/Algebraische%20Geometrie/Lehre/WS23.Padische.1966.kleiman.weil.pdf}.
The earlier JAC-P20 visual audit of Saavedra VI.4.2.0, p.~363,
VI.4.4.2, p.~373, and Appendix A3, pp.~400--401, recorded the
standing standard-conjecture hypothesis. That source audit is
carried forward here; it is not represented as a new complete
reading of Saavedra in this focused interval.

\paragraph{Ledger consequence.}
The numerical and Lefschetz conclusions of JAC-P16 remain
proved under the accepted 032 input. The explicit finite-cover
comparison of C44-W06 and the extending-correspondence proof
of C44-W10--W11 remain the established geometric form transfers
in their respective scopes. The unrestricted
QT4-abeliantype/C44-W09 row is unchanged: the extra content is
that $Q_X$, not merely some transported form on its motive,
belongs to the positive compatibility class. No new sign
application or progress-clock event is claimed.

% END work/scan_w/qt4_abelian_type_transfer_reaudit.tex

% BEGIN work/scan_w/sharp_cycle_bound.tex
% C44_POS /root/scan_w; GPT-6 Astra Pro; effort not exposed.
\section{Sharp complementary-cycle degrees and the refined detector constant}

\begin{proposition}[C44-W12: a sharp algebraic complementary-cycle estimate]
Let $A$ be an abelian variety of dimension $g\geq1$ over an
algebraically closed field of arbitrary characteristic, and let
$h=c_1(L)$ be an integral ample class. Put
\[
                   V=\frac{\int_Ah^g}{g!}=\chi(L)>0.
\]
If $Z,Z'$ are effective cycles of dimensions $r$ and $g-r$ with
nonnegative integer coefficients, then
\[
 0\leq\deg(Z\cdot Z')
   \leq\frac{\deg(h^rZ)\deg(h^{g-r}Z')}{r!(g-r)!V}.
 \tag{C44.W12.1}
\]
The coefficient is sharp. Neither the standard conjectures nor a
characteristic-zero lift is an input.
\end{proposition}

\begin{proof}
By bilinearity it suffices to treat integral subvarieties. Their
product $T=Z\times Z'$ is an integral projective variety, which may
be singular. Set $e=\deg(Z\cdot Z')$. The translation-moving
argument in C44-W01, with complementary dimensions $r,g-r$, gives
\[
 f_*[T]=e[A],\qquad f(z,z')=z-z',\qquad e\geq0.
\]
Inseparable generic degrees are included. Let $s(z,z')=z+z'$.
The degree of this map, interpreted as zero when it is not dominant,
is $\deg(Z\cdot[-1]_*Z')$. In a Weil cohomology of characteristic
different from the coefficient prime, $[-1]^*$ acts by $(-1)^i$
on $H^i(A)$, and therefore acts identically on every cycle class.
It follows that
\[
                         s_*[T]=e[A].
 \tag{C44.W12.2}
\]
This is an equality of top-dimensional pushforward cycles, whose
coefficients have just been computed as intersection numbers; no
algebraic equivalence between $Z'$ and $[-1]_*Z'$ is asserted.

Write $h_1,h_2$ for the pullbacks to $T$ and put
$\alpha=f^*h$, $\beta=s^*h$. The two classes are nef and the
numerical parallelogram identity gives
\[
 \alpha+\beta=2(h_1+h_2),\qquad
 \int_T\alpha^g=\int_T\beta^g=e\int_Ah^g.
\]
The Khovanskii--Teissier inequality for two nef divisors gives
\[
 \left(\int_T(\alpha+\beta)^g\right)^{1/g}
       \geq\left(\int_T\alpha^g\right)^{1/g}
            +\left(\int_T\beta^g\right)^{1/g}.
 \tag{C44.W12.3}
\]
It applies to integral possibly singular projective varieties in
arbitrary characteristic: see Lazarsfeld, \emph{Positivity in
Algebraic Geometry I}, Corollary~1.6.3(iii), pp.~90--91, and
Remark~1.6.5, p.~91. Equivalently, Cutkosky,
\emph{Teissier's problem on inequalities of nef divisors over an
arbitrary field}, Corollary~6.3, formula~(20), gives precisely this
statement for a complete integral variety over any field.
In particular no smoothness of $Z$ or $Z'$ is being assumed here.

Substitute the preceding equalities into (C44.W12.3) and cancel
the factor $2$. Raising to the $g$-th power gives
\[
 e\int_Ah^g\leq\int_T(h_1+h_2)^g
       =\binom gr\deg(h^rZ)\deg(h^{g-r}Z').
\]
Using $\int_Ah^g=g!V$ proves (C44.W12.1).

For sharpness, take a product of principally polarized abelian
varieties of dimensions $r$ and $g-r$, with its product principal
polarization, and take the two complementary factors as $Z,Z'$.
Then $e=1$, $V=1$, and their degrees are $r!$ and $(g-r)!$,
respectively. Equality holds. The cases $r=0$ or $r=g$ are given
by a point and the whole variety.
\end{proof}

\begin{theorem}[C44-W13: refined Gram and division constants]
Retain all the data and hypotheses of C44-W02--C44-W03, including
the explicit integral polynomial $p$, the dual numerical-lattice
condition $\delta t\in\Lambda^\sharp$, and the conditional
unrefereed 032 positivity input. Define
\[
 \widehat B=\frac{\sqrt{R\kappa/2}}{n!\sqrt V},\qquad
 \varepsilon=
 \begin{cases}1,&s\text{ even},\\2,&s\text{ odd}.
 \end{cases}
 \tag{C44.W13.1}
\]
For an integral detector $Y$ of degree $D$, its Gram entry obeys
\[
                  0<d=Q(y_1,y_1)\leq\widehat B^2D^2.
 \tag{C44.W13.2}
\]
The exact Gram matrix remains
$d\left(\begin{smallmatrix}1&s/2\\s/2&q^2\end{smallmatrix}\right)$.
For every nonzero rational Weil class $t$ under the stated lattice
condition, with $T_0=Q(t,t)$, the detector complexity satisfies
\[
 \delta_W\leq c_t\leq
 \delta\varepsilon\max\left\{
          \frac H4\widehat B^2,
          \frac{q\kappa}{2}\widehat B\sqrt{HT_0}\right\}
                                        \delta_W^2.
 \tag{C44.W13.3}
\]
The corresponding quartic bound is
\[
 c_t\leq\delta\max\left\{
          \frac H4\widehat B^4,
          \frac{q\kappa}{2}\widehat B^3\sqrt{HT_0}\right\}
                                        \delta_W^4.
 \tag{C44.W13.4}
\]
Both bounds remain valid with the variety-independent constant
$\widehat B_0=\sqrt{R\kappa/2}/n!$ in place of $\widehat B$.
\end{theorem}

\begin{proof}
The middle $K$-weight is $a^n\bar a^n=q$, so $g(q)=0$ and hence
$p(q)=0$. Use the actual effective cycle representatives
\[
 Y^+=\sum_{c_j>0}c_j[a^j]^*Y,\qquad
 Y^-=\sum_{c_j<0}(-c_j)[a^j]^*Y.
\]
They have difference representing $y_1$. The degree formula from
C44-W02 and $p(q)=0$ imply the exact identities
\[
       D(Y^+)=D(Y^-)=\frac{\kappa D}{2}.
 \tag{C44.W13.5}
\]
Thus $\kappa$ is even, and $[a]^*Y^\pm$ represent $y_2$ with
each part of degree $q\kappa D/2$. These are explicit signed
representatives; cancellation of common components is unnecessary.

The self-adjoint projector identity from C44-W02 is
$p(T)^\dagger p(T)=Rp(T)$. Consequently
\[
 d=R\left|\deg\bigl(Y\cdot(Y^+-Y^-)\bigr)\right|.
\]
By C44-W12 with $g=2n,r=n$, the two numbers
$I_\pm=\deg(Y\cdot Y^\pm)$ are nonnegative and satisfy
\[
              I_\pm\leq\frac{\kappa D^2}{2(n!)^2V}.
\]
Use $|I_+-I_-|\leq\max(I_+,I_-)$, rather than adding the two
upper bounds. This proves (C44.W13.2). All intersection estimates
in this paragraph are unconditional; 032 is used for the positive
Gram form and the subsequent Cauchy--Schwarz inequality.

Choose $D=\delta_W$ and put $b_i=Q(\delta t,y_i)\in\mathbf Z$.
Exact inversion of the Gram matrix gives
\[
 \frac{\delta\varepsilon Hd}{4}\,t
   =\varepsilon\left(q^2b_1-\frac{s}{2}b_2\right)y_1
        +\varepsilon\left(b_2-\frac{s}{2}b_1\right)y_2.
 \tag{C44.W13.6}
\]
Every displayed coefficient, including the multiplier on the left,
is an integer. If $s$ is even then $4\mid H$. If $s$ is odd,
integrality of the Gram entry $sd/2$ forces $2\mid d$, and
$\varepsilon=2$ makes both the multiplier and coefficients integral.
The multiplier is positive.

Writing $t=t_1y_1+t_2y_2$, positive definiteness gives
\[
 |t_1|\leq\frac{2q\sqrt{T_0}}{\sqrt{dH}},\qquad
 |t_2|\leq\frac{2\sqrt{T_0}}{\sqrt{dH}}.
\]
The two integer coefficients in (C44.W13.6) therefore have
absolute values at most
$\delta\varepsilon q\sqrt{HT_0d}/2$ and
$\delta\varepsilon\sqrt{HT_0d}/2$, respectively. With the exact
part-degrees (C44.W13.5), each positive or negative part of their
cycle combination has degree at most
\[
 \frac{\delta\varepsilon q\kappa}{2}
            D\sqrt{HT_0d}
    \leq\frac{\delta\varepsilon q\kappa}{2}
            \widehat B\sqrt{HT_0}\,D^2.
\]
The multiplier is at most
$\delta\varepsilon H\widehat B^2D^2/4$. These give the upper
bound in (C44.W13.3). The lower bound is exactly the integral
component argument of C44-W03.

For (C44.W13.4), use the integer multiplier
$N=\delta\det G=\delta d^2H/4$. The two adjugate coefficients
are $Nt_1,Nt_2$ and are integers by the integral Gram matrix.
Their absolute values are at most
$\delta qd^{3/2}\sqrt{HT_0}/2$ and
$\delta d^{3/2}\sqrt{HT_0}/2$. The degree of either signed part
is consequently at most
\[
        \frac{\delta q\kappa}{2}D d^{3/2}\sqrt{HT_0}
       \leq\frac{\delta q\kappa}{2}
                         \widehat B^3\sqrt{HT_0}D^4,
\]
whereas $N\leq\delta H\widehat B^4D^4/4$.
Finally $V\geq1$ proves the uniform-constant assertion.
\end{proof}

\paragraph{Exact relation to the earlier constants.}
C44-W12 sharpens the valid elementary estimate C44-W01, and
C44-W13 uses additionally the vanishing of the middle weight
$p(q)$ and the parity of the integral Gram matrix. No alteration
of the projector, detector degree, denominator convention, or
class-level meaning of $c_t$ is involved. In the characteristic-$p$
Theorem W deduction, the strongest constant established here for
$\sqrt{Q(y_1,y_1)}$ is $\widehat B D$. An identification with the
original analytic Wirtinger constant is not asserted.

The primary nef-inequality sources inspected were
\url{https://www.math.stonybrook.edu/~roblaz/Reprints/Laz.PAG.Chapt1.pdf},
Theorem~1.6.1 and Corollary~1.6.3 with their proofs and
Remark~1.6.5, and
\url{https://arxiv.org/pdf/1304.1218}, Theorem~6.1 and
Corollary~6.3, especially formula~(20) and its proof.


% END work/scan_w/sharp_cycle_bound.tex

% BEGIN work/consolidated_th5.tex
\section{The corrected fifth theorem, with all domains visible}

\begin{theorem}[C44-T05: synthesis of H9-T1, H9-T2, H9-T7,
H9-T17 and H9-T18]\label{thm:final-th5}
Use rational Chow groups and the signed geometric intersection height
of JAC1964 H9-T1. The following is the consolidated form of cote 44,
scan 004, Th.~5. The formulas below are unconditional under their
stated admissibility and geometric hypotheses. Each assertion that
uses the abelian Hodge standard conjecture is explicitly conditional
on the unrefereed 032/Milne inputs.

\begin{enumerate}
\item Let $f:\mathcal X\to C$ be projective with $\mathcal X$ smooth
of dimension $2m$, $C$ a smooth projective curve, and smooth generic
fibre $Y$. Let $x_0,y_0$ extend numerically trivial cycles on $Y$.
Choose a finite-dimensional vertical correction space, supported in
finitely many relevant fibres. For each such fibre let $W_s$ be its vertical numerical cycle
space, $R_s$ its radical, and $M_s$ the nonsingular matrix on a chosen
basis of $W_s/R_s$. Suppose the extensions annihilate $R_s$, and
that these spaces account for all vertical pairings. Write $a_s,b_s$
for their pairing vectors. Then admissible corrections exist and
\[
 h_m(x_0|_Y,y_0|_Y)=(-1)^m
 \left(\deg(x_0y_0)-\sum_s a_s^{t}M_s^{-1}b_s\right).
 \tag{C44.T05.1}
\]
If only a smaller vertical space has been used, this expression is
the corrected intersection, and is a height only after the remaining
vertical orthogonality is verified. Failure to annihilate a vertical
radical is an obstruction to correction, not a term one may omit.

\item For a Lefschetz pencil on a smooth $2m$-fold $X$, if
$H^{2m}(Y_s)\to H^{2m}(Y_{\bar\eta_s})$ is injective at every
singular fibre, all corrections vanish for ambient cycles whose
generic restrictions are homologically trivial:
\[
                  h_m(j^*z,j^*z')=(-1)^m\int_Xzz'.
 \tag{C44.T05.2}
\]
The same injectivity makes any global extension of a homologically
trivial generic cycle admissible. In the ordinary Picard--Lefschetz
setting the condition is equivalent to nonzero vanishing cycles;
SGA 7 condition $(A)$ is the characteristic-independent formulation
used in the predecessor. For genuine abelian pencils it holds by
JAC-P07. If $X$ is abelian, conditional on the unrefereed 032/Milne
preprints, (C44.T05.2) is positive definite
on primitive ambient algebraic cohomology.

\item If $X$ is abelian, let $B$ be its smooth axis and use
$\mathcal H,\mathcal H_0,A_X,A_B$ of H9-T7, (H9.17). Then
\[
 \mathcal H/\mathcal H_0\simeq A_X\mathbin\perp A_B,
 \qquad h_m\simeq
 \left((-1)^m\int_Xzz'\right)\mathbin\perp
 \left((-1)^{m-1}\int_Bww'\right).
 \tag{C44.T05.3}
\]
In particular $h_Bw=0$ for $w\in A_B$. Conditional on 032/Milne,
semipositivity on the whole $\mathcal H$ is equivalent to the weak
axis sign on $A_B$. If the axis form is positive definite, the
radical is exactly $\mathcal H_0$. If it is only semidefinite,
the quotient of the radical by $\mathcal H_0$ is the radical of
that axis form; equality with $\mathcal H_0$ is not inferred.
The axis sign is known for $m=2$, for characteristic zero, and in
the generic-axis vanishing cases H9-T8, H9-T8a and JAC-P09.
C44-A04 and C44-A06 give additional explicit positive subspaces,
not a sign theorem for every special axis.

\item For a smooth threefold generic fibre, on the admissible
homologically trivial domain let
\[
 \mathcal H_h^{\rm Alb}=\{z:\operatorname{alb}(hz)=0\},\qquad
 \psi_h(D)=\operatorname{alb}(h^2D),\qquad
 D_z=\psi_h^{-1}\operatorname{alb}(hz).
\]
The rational isogeny $\psi_h:\Pic^0(Y)\to\Alb(Y)$ gives the
orthogonal decomposition
\[
                    z=(z-hD_z)+hD_z.
 \tag{C44.T05.4}
\]
The first term lies in $\mathcal H_h^{\rm Alb}$; the second has
negative definite height modulo the constant trace. In
characteristic zero the first height is semidefinite by Zhang,
Theorem 1, with the precise admissible-model equality criterion
of H9-T17, (H9.77)--(H9.80). This sign statement does not extend
to the whole homologically trivial group after an arbitrary
finite cover: H9-T17 gives the exact negative value $-12d^3$ on
the displayed $E^4$ example, in every characteristic.

\item For an abelian-fourfold pencil and a finite separable cover
$C\to\mathbf P^1$ unramified over the critical values, let
$k:\mathcal X\hookrightarrow A\times C$ be the smooth incidence
fourfold. All homologically trivial generic cycles are admissible.
On $\mathcal H_h^{\rm Alb}$ the primitive correction of H9-T18
gives
\[
              h_2(z,z)=\int_{A\times C}a^2\lambda
                              +\int_{\mathcal X}v^2,
 \quad \lambda^2a=0,\quad k_*v=0.
 \tag{C44.T05.5}
\]
Conditional on 032/Milne the first term is positive definite on
its algebraic domain. The remaining assertion is precisely the
nonnegativity of the algebraic primitive form on $\ker k_*$.
C44-C01 splits off its positive old surface-axis part; the new
trace-kernel part is the remaining special-cover gate. The
generic vanishing results H9-T19--T26 and C44-C02--C04 give
the norm identities in their explicitly stated parameter scopes.
They do not establish this sign for every closed special cover.
\end{enumerate}
\end{theorem}

\begin{proof}
This is an assembly of the specified predecessor results, not a new
proof of them. Part (1) is H9-T1, (H9.3)--(H9.5), equal to
SIGN\_NEXT SN-44-01 after its normalization, with vertical blocks
separated by disjoint fibres. Part (2) is H9-T2, (H9.6)--(H9.7),
and its full-extension argument at the beginning of H9-T7;
the Picard--Lefschetz equivalence is SN-44-03. Part (3) is the
isometry H9-T7, (H9.18), with the elementary radical statement
of an orthogonal sum and the conditional ambient sign. Part (4)
is H9-T17, (H9.79)--(H9.84), including both of Zhang's hypotheses
and Kahn's trace quotient. Part (5) is H9-T12 admissibility and
H9-T18, (H9.86)--(H9.91), followed by the independently checked
old/new decomposition C44-C01. No step asserts algebraic
equivalence from homological triviality.
\end{proof}

\paragraph{The incidence higher Picard variety and the arithmetic reading.}
Grothendieck's $P^m$ in cote 44, scans 002 and 004, is attached
to algebraically trivial cycles via incidence. A homologically
trivial class in the displayed larger Chow groups is not thereby
a point of $P^m$. The Neron--Tate interpretation on a primitive
incidence quotient requires the actual regular homomorphism and
the height comparison; these are the separate inputs of
JAC-T4.1--T4.5 and SN-44-01. The full-middle-factor polarization
is refuted by JAC-P03, and the unrestricted covered height is
refuted by H9-T17. Neither counterexample refutes the primitive
axis sign asked here. The normalization and the Albanese kernel
implement the primitive-decomposition margin of scan 008.

In cote 16, scans 049--050, the algebraicity of Lefschetz operators
is condition $C=C_G=B$, while the height signs use $I$ on the
specified primitive algebraic spaces. Restoring the Lefschetz
factor in $A''$ (SN-C03--C04) is a separate issue from these
signs. Theorem C44-T05 does not deduce $I$ from $C$ or identify
cohomological semisimplicity with injectivity on rational Chow
groups. In the SGA framework, the local input is SGA 7 II,
XVIII 5.3.2, 5.3.4--5.3.5 and 6.2; the variable irreducibility
is XVIII 6.7.2. In the Chow-motive framework the global terms
are the blowup summand $h(B)(1)$ and the old/new finite-cover
projectors. These locators and their roles are retained from
the checked predecessor arguments.

% END work/consolidated_th5.tex

\part{Unchecked candidates and unresolved work}

% BEGIN work/final_research_boundary.tex
\section{Final unpromoted research routes and their exact limits}
\label{sec:final-research-limits}
This section records the final state of the A31u/A31v route. It is not
an additional ramification theorem. The complete route calculations,
including rejected intermediate inferences, are retained in the
accompanying research notes and their whole-note reading audits.

\paragraph{Locally exact corrections.}
An integral supported one-form that is genuinely rationally exact has
special reduction in a higher Cartier boundary. The calculation writes
its primitive with a bounded power of $p$ in the denominator, decomposes
coefficients over $W(k)$ in a ramified DVR power basis, and applies
Illusie's higher-boundary criterion (Ast\'erisque 63, printed p.~95,
Lemma~2.4.3). At a fixed finite pole bound, sufficiently many Cartier
operations annihilate every supported normal weight. Together with
A31u this permits ordinary boundary corrections inside the higher-boundary
submodule. It says nothing comparable about an arbitrary ambient
ordinary cohomology class.

At characteristic-$p$ Artin precision $n\leq e$, top supported coherent
cohomology can lift a prescribed ordinary ambient image without first
lifting a global ambient supported section. Finite Serre duality and
Krull intersection give the required special-fibre surjectivity; the
vanishing of $H^5_S$ permits induction in precision. This route has no
retained pole bound compatible with A31n on the fixed axis, and does
not handle arbitrary mixed-characteristic precision.

\paragraph{Quantitative conormal attempt.}
The power-pullback construction in A31p has a linear ampleness margin.
Keeping its Frobenius-pullback cokernel fixed and varying the six-generator
surjection puts the remaining kernel in a bounded family. A finite
Cartier calculation can be organized from the bottom through
$B^1,Z^1,B^2,Z^2,B^3,Z^3$, using cokernel and extension bounds.
An earlier proposed regularity bound for a kernel from its two right-hand
terms was rejected. Even the proposed uniform stabilization of the Ext
system would not establish the required closed-period saturation or
the bounded global representatives. Full canonical restriction to a
fixed-order neighbourhood is not a uniform substitute: its global
section count grows in order $t^6$, whereas the fixed-order jet count
on the fixed surface has order $t^2$. This comparison does not determine
the rank of the primitive closed-period image.

\paragraph{The constrained test that remains unproved.}
Top Cartier gives the distinguished triangle
\[
 Z_B^3[-3]\longrightarrow F_*\Omega_B^{\geq3}
       \longrightarrow\Omega_{B^{(1)}}^4[-4].
\]
The preceding hypercohomology term of the last complex is zero.
Thus $H^1(Z_B^3)$ injects into
$\mathbb H^4(\Omega_B^{\geq3})$ and is the kernel of its top Cartier
map. Under the A29 degeneration hypotheses, this is the annihilator
inside $F^3H^4_{\mathrm{dR}}(B)$ of the ordinary closed-function
periods. Cartier wedge duality identifies the coefficient bundle
$E^*\otimes\omega$ in A31v with $Z_B^3$. A possible direct replacement
for full closed-period saturation would therefore prove surjectivity of
\[
 H^1(Z_B^3)_{\mathrm{primitive}}
       \longrightarrow H^1(\widehat B_S,\widehat Z_B^3).
\]
\emph{This surjectivity is not proved.} Lifting a finite closed jet to
a global $F^3$ class can leave a nonzero conjugate-top canonical section.
The mixed-characteristic ambient correction also remains unproved.
Neither finite-dimensionality nor the bounded regularity route supplies
these two missing steps.

These limits concern task~1(a), JAC1964 QH9-axis-lift, and the axis term
of cote~44 Th.~5 (scan~004, with the primitive warning on scans~008--009).
They make no additional assertion about cote~16 scans~049--050 and use
no 032 or Milne2026 input.

\paragraph{Research stopping decision.}
Scheduled Review~9 was recorded at 10:14:05 UTC on 22 November 2026.
Reviews~8 and~9 both found no additional qualifying movement; the
user's stopping rule therefore ended research. The quaternion extensions
had been counted once at Review~5, and the A31f onward lifting
extensions once at Review~7. No unresolved goal is converted to a
proved conclusion by that stopping decision.

% END work/final_research_boundary.tex

\section{Unresolved goals}
The fixed-axis sign is not proved for every special axis, and no
wrong-sign algebraic class is constructed. The universal lifting
question for classes remains open. The special-cover sign remains
open at arbitrary closed points with new algebraic trace-kernel
classes. Fermat Abel--Jacobi injectivity remains open, with the
exact gates specified above. Full Th.~4 for arbitrary new special
incidence directions is not proved. The proposed supersingular
Fermat family does not meet the requested initially-unknown-$I$
status. The ledger and question register give the current precise
state of every treated row.
\end{document}
